\documentclass[11pt,a4paper]{article}
\usepackage[T1]{fontenc}
\usepackage{lmodern,amsmath,amssymb,amsthm,mathtools,mathrsfs,booktabs,longtable,enumitem}
\usepackage[margin=27mm]{geometry}
\usepackage{microtype,tabularx,array,xurl}
\usepackage[hidelinks]{hyperref}
\hypersetup{pdftitle={Theta Prime Gamma Energy and its Spectral Structure},pdfauthor={Jongmin Choi},pdfsubject={Integrated closed energy, verified response gates, spectral synthesis, sparse and contiguous zero bands, and fixed lattice positivity},pdfkeywords={theta kernel, prime powers, Gamma form, verified response gate, spectral synthesis, zero jets, sparse windows, fixed lattice positivity}}
\setlength{\parindent}{1.2em}\setlength{\parskip}{0.2em}
\setlength{\emergencystretch}{2em}\allowdisplaybreaks[2]
\numberwithin{equation}{section}
\newtheorem{theorem}{Theorem}[section]
\newtheorem{proposition}[theorem]{Proposition}
\newtheorem{lemma}[theorem]{Lemma}
\newtheorem{corollary}[theorem]{Corollary}
\newtheorem{hypothesis}[theorem]{Hypothesis}
\theoremstyle{definition}\newtheorem{definition}[theorem]{Definition}
\theoremstyle{remark}\newtheorem{remark}[theorem]{Remark}
\newcommand{\op}{\mathrm{op}}\newcommand{\lo}{\mathrm{lo}}\newcommand{\up}{\mathrm{up}}\newcommand{\eps}{\varepsilon}
\newcommand{\rhip}[2]{\langle#1,#2\rangle_H}
\newcommand{\R}{\mathbb R}\newcommand{\C}{\mathbb C}\newcommand{\N}{\mathbb N}
\newcommand{\E}{\mathcal E}\newcommand{\D}{\mathcal D}
\newcommand{\He}{L^2_{\mathrm{even}}(\mu)}\newcommand{\Le}{L^2_{\mathrm{even}}(dx)}
\newcommand{\Hh}{\mathcal H}\newcommand{\Hw}{\mathcal H}
\newcommand{\HH}{\mathcal H_W}\newcommand{\XX}{X_W}\newcommand{\EE}{\mathcal E_{N,\beta}}
\newcommand{\tg}{t_\Gamma}\newcommand{\cg}{c_\Gamma}
\newcommand{\qD}{\mathfrak q_D}\newcommand{\qth}{\mathfrak q_\theta}
\newcommand{\dom}{\operatorname{Dom}}\newcommand{\Var}{\operatorname{Var}}
\newcommand{\rank}{\operatorname{rank}}\newcommand{\ess}{\operatorname{ess}}\newcommand{\supp}{\operatorname{supp}}
\newcommand{\ip}[2]{\langle#1,#2\rangle}\newcommand{\norm}[1]{\lVert#1\rVert}\newcommand{\abs}[1]{\lvert#1\rvert}
\newcommand{\dd}{\mathop{}\!\mathrm d}\newcommand{\one}{\mathbf1}
\newcommand{\Zc}{\mathcal Z_\delta}\newcommand{\Rcal}{\mathcal R}
\newcommand{\RR}{\mathscr R}\newcommand{\SSS}{\mathcal S_{\!\exp}}\newcommand{\CE}{\mathcal C_{\exp}}
\newcommand{\Eloc}{\mathcal E_{\mathrm{loc}}}\newcommand{\Def}{\mathfrak D}
\newcommand{\PP}{P}\newcommand{\J}{J}\newcommand{\V}{V}\newcommand{\wo}{\mathcal W}
\newcommand{\Nn}{\mathcal N}\newcommand{\Ii}{\mathcal I}\newcommand{\Zz}{\mathcal Z}
\newcommand{\Kk}{\mathcal K}\newcommand{\Aa}{\mathcal A}\newcommand{\Ff}{\mathcal F}
\newcommand{\cl}{\overline{\operatorname{span}}}\newcommand{\Ran}{\operatorname{Ran}}\newcommand{\HS}{\mathfrak S_2}
\DeclareMathOperator{\Rea}{Re}\DeclareMathOperator{\Ima}{Im}
\DeclareMathOperator{\adj}{adj}
\title{Theta Prime Gamma Energy\\and its Spectral Structure}
\author{Jongmin Choi\\[3pt]\small Independent Researcher, Seoul, Korea\\\small ORCID iD: \href{https://orcid.org/0009-0008-7448-514X}{0009-0008-7448-514X}}
\date{November 2026\\Version 3 (working integration)}
\begin{document}\maketitle
\begin{abstract}
We develop an integrated closed-form and spectral framework for the even theta-weighted energy containing the complete Gamma difference integral and every prime power. The maximal energy domain is identified with its weighted physical realization. Exact radicals, compact spectral reductions, weighted synthesis, and full-multiplicity zero-jet form and graph cores isolate the remaining sign problem. Quantitative variance bounds are derived from clearly separated local-sign and complete-response premises. A dimension-independent residual gate and full-norm continuation retain actual source coordinates, all tails, and the mean correction. Retaining the original fourth-stage witness margin gives a stronger conditional comparison. A completed fourth-stage recomputation uses repaired integer-power bounds and fresh exact acceptance; historical same-code replays retain a separate arithmetic qualification. We prove positivity for arbitrarily many sufficiently separated logarithmic-complexity windows and for the entire actual zero-jet space in a single band of width eleven times the logarithm of height. A complete interaction estimate also permits sparse families of entire growing bands under the stated stronger spacing rule. For a fixed low space with a finite theta exponential-polynomial representation, the exact response cost of the entire sparse high family tends to zero in its actual weighted Gram. Separately replayed ball-arithmetic certificates supply positive margins on the first four and five exact zero modes, with their respective weighted coupling constants. The fifth-mode proof retains a derivative Gram and verifies its complete shifted Schur cost. Derivative-Gram graph bounds transfer that five-mode certificate, and then an incremental six-mode certificate, to their Gamma-low projections together with the complete complementary Gamma-high form domain under explicit sector and operator-row inputs, while leaving the other low directions unresolved. Under the stated sparse estimates, the four-mode low space attaches to an infinite selected tail above an explicit, very large height. A separate fixed-real-lattice theorem permits unrestricted finite or infinite mode count, and a signed block criterion identifies the additional coupling requirements. A restricted Gamma spectral decomposition on the actual zero-jet sector gives projected graph cores and an exact low-space Schur criterion. Fixed inverse filters, complete residual optimization, exact null-response repair, and ordered tail transfer clarify the further obligations of a cofinal certificate. Uniform two-sided operator brackets separate response accuracy from source dimension, while an asymptotic obstruction shows why a fixed graph debit cannot certify an exhaustive raw-source family. Each result states its norm, domain, geometry, and verification boundary. The qualitative prime-number-theorem refinement supplies no effective final height here. A later complete single-band estimate covers widths $o(\log T\sqrt{\log\log T})$ in ordinary $L^2$, and a conservative fixed multiple of this scale gives eventual positivity; Appendix~\ref{app:growing-zero-bands} retains the precise hypotheses and full proof. These restricted positivity results do not establish positivity on arbitrary unions or the full zero family. The global half inequality and the Riemann hypothesis remain unproved.
\end{abstract}
\noindent\textbf{Keywords.} Theta kernel; prime powers; Gamma difference form; verified response gate; Schur complement; spectral synthesis; zero jets; sparse frequency windows; fixed lattice positivity.
\clearpage\setcounter{tocdepth}{1}\tableofcontents\clearpage

\section{Introduction and logical structure}\label{sec:intro}
The object of this paper is one nonnegative jump energy, rather than a family of truncated prime or finite-matrix substitutes. Its centered variance gap can be studied in three complementary coordinates: profiles in $L^2(\mu)$, physical numerators in a reciprocal-theta weighted space, and an ordinary-$L^2$ confining pencil. The distinction between these spaces is essential. A form congruence is not an ordinary-$L^2$ unitary equivalence, a convergent response calculation does not determine a Schur sign, and approximation by full zero jets does not prove positivity on those jets.

The analytic conclusions are the full-domain physical bridge, exact spectral and radical identities, parameter-scaled finite complements, weighted spectral synthesis, and stronger-topology zero-jet cores. Positivity is established at fixed exponential complexity, on arbitrarily many sparse logarithmic-complexity unit windows, on one entire actual-zero band of width eleven times the logarithm of height, and on the stated sparse families of entire growing bands. The last geometry uses exponent-$12$ spacing, with the ordinary norm of each whole block. A fixed real lattice gives a different infinite-dimensional synthesis theorem. Section~\ref{sec:four-zero-certificate} supplies one concrete positive low seed in the admissible global exponential-polynomial class: the first four zero modes, with an actual weighted coupling bound and a separately conditional sparse attachment. These geometries must not be conflated. A further theorem bounds the complete response of a sparse high family against a fixed low space with a finite theta--exponential-polynomial representation; its low-space constant and weighted Gram are retained. Quantitative global comparisons also have different computational premises: a whole-domain local sign, historical complete-response contracts, and newer finite response gates with full infinite-complement and residual enclosures. No finite pivot computation alone proves a whole-domain inequality.

\subsection{Common notation}
All physical spaces are complex and even unless a real profile is explicitly required. Inner products are linear in their first argument. The positive even kernel $\kappa$ is constructed in Section~\ref{sec:foundations}. Set
\begin{equation}\label{a:eq:weights}
 c(x)=\cosh(x/2),\qquad d\mu=2c\kappa\,dx,\qquad
 W=\frac{\kappa}{2c},\qquad \Hh=L^2_{\mathrm{even}}(\R,dx/W).
\end{equation}
The measure $\mu$ has mass one. Write $\norm{u}_{\Hh}^2=\int|u|^2/W$, and reserve $\norm{\cdot}_2$ for ordinary Lebesgue $L^2$. With $\Lambda$ the von Mangoldt function, define
\begin{equation}\label{eq:common-shifts}
 \rho(s)=\frac{e^{-s/2}}{1-e^{-2s}},\qquad
 a_n=\frac{\Lambda(n)}{\sqrt n},\qquad s_n=\log n.
\end{equation}
Thus $a_n$ retains every prime power, not merely primes. The exact numerator forms are
\begin{align}
 \tg[u]&=\int_0^\infty\rho(s)\norm{u(\cdot+s)-u}_2^2\,ds,\label{a:eq:physicalgamma}\\
 p[u]&=2\operatorname{Re}\sum_{n\ge2}a_n\int_\R u(x+s_n)\overline{u(x)}\,dx,\label{a:eq:physicalprime}\\
 q_0[u]&=\tg[u]-\cg\norm{u}_2^2-p[u],\qquad
 \cg=\log\pi-\psi(1/4),\quad\psi=\Gamma'/\Gamma.\label{a:eq:qzero}
\end{align}
The complete prime form is bounded on $\Hh$. The common numerator domain and centered space are
\begin{equation}\label{a:eq:X}
 \begin{gathered}
 X=\{u\in\Hh:\tg[u]<\infty\},\qquad
 \norm{u}_X^2=\norm{u}_{\Hh}^2+\tg[u],\\
 m(u)=\int_\R cu,\quad \Hh_c=\ker m,\quad X_c=X\cap\Hh_c.
 \end{gathered}
\end{equation}
The original energy $\E$ and its maximal domain $\D$ are defined in~\eqref{a:eq:Eg}--\eqref{a:eq:Edomain}. Their bridge is
\begin{equation}\label{eq:intro-bridge}
 \kappa\D=X,\qquad q_0[\kappa f]=\E[f]-\tfrac12\norm f_\mu^2,
 \qquad \mu(f)=2m(\kappa f).
\end{equation}
Multiplication $Uf=\kappa f$ is unitary from $\He$ to $\Hh$, but not to ordinary $\Le$. We use $U$ only for this unitary; the letters $J$ occurring in interface kernels or zero jets have their separately displayed meanings. The entire transform is $F_u(z)=\int u(x)\cos(zx)\,dx$; on even functions it equals $\mathcal Fu(z)=\int u(x)e^{izx}\,dx$. The ordinary unitary transform, used in multiplier estimates, is $\widehat u(t)=(2\pi)^{-1/2}\int u(x)e^{-itx}\,dx$.

\subsection{Results and their dependency classes}
\begin{center}\small
\begin{tabular}{@{}>{\raggedright\arraybackslash}p{0.17\textwidth}>{\raggedright\arraybackslash}p{0.49\textwidth}>{\raggedright\arraybackslash}p{0.26\textwidth}@{}}
\toprule Class & Conclusion & Additional premise or limitation\\\midrule
Analytic & Closed domains, actual operator, radicals, synthesis and zero-jet cores & Classical explicit formula and standard analysis\\[3pt]
Analytic & Fixed-complexity, sparse-window, single and sparse growing-band, and fixed-lattice positivity & The stated geometry, complexity and actual form domain\\[3pt]
Analytic & Vanishing exact low-to-sparse response cost & Fixed admissible low space and sparse high family; no positive low margin supplied\\[3pt]
Computed low spaces & Margins $1/32$ on four exact zero modes and $1/25$ on five & Arb/Acb replay; actual weighted Grams; four-mode conditional sparse attachment\\[3pt]
Local sign & $\E\ge(49307/100000)\Var_\mu$ & Accepted whole-local-domain sign at radius $4/5$\\[3pt]
Historical response & $\E\ge(2861149/5760000)\Var_\mu$ & Complete residual and centered-response contracts\\[3pt]
Historical replay & Third-stage coefficient $\lambda_+$ in Section~\ref{vrg:sec:generalized} & Complete enclosure contract; shared integer-power exposure remains unresolved\\[3pt]
Recomputed response & Fourth-stage base $\lambda_4$ and retained-margin $\lambda_{4,\mathrm{ret}}$ in Corollary~\ref{vrg:cor:retained4} & Repaired full generation and fresh exact gates; analytic and implementation premises retained\\[3pt]
Open & $\E\ge\tfrac12\Var_\mu$; every finite exact full-zero-jet matrix positive semidefinite & No proof is supplied\\\bottomrule
\end{tabular}
\end{center}
The computationally independent analytic seed remains useful even though it is not the largest coefficient. The historical coefficients $0.49307$ and approximately $0.496727$ have not been relabeled as the newest bounds. The stronger third and fourth stages give, respectively,
\[
 \lambda_+=\frac{1381250815636970598719595}{2767005669802701449752576}
 \quad\hbox{and}\quad
 \lambda_4=\frac{542006223062224701317434549}{1084868104425248291110985728}.
\]
Their numerical values are approximately $0.499186116859478389$ and $0.4996056394794406$. The unchanged fourth-stage matrices and original witness additionally give $\lambda_{4,\mathrm{ret}}\approx0.49964024316343157254$ by retaining the strict margin, as proved in Corollary~\ref{vrg:cor:retained4}. This is not a new source stage. Sections~\ref{sec:four-zero-certificate} and~\ref{sec:five-zero-certificate} are separate finite actual-zero computations: their positive unshifted margins use their own weighted Grams and do not improve a global coefficient. These are implications from the complete stated enclosure contracts. The fourth stage has now been regenerated with the repaired integer-power routine, and both gates and the retained margin were accepted afresh. The earlier third-stage same-code replay remains a reproduction with an unresolved shared-primitive enclosure obligation; it is not retroactively validated by the fourth-stage repair. Appendix~\ref{app:current-evidence} separates these histories. Neither run is a second complete implementation or a formal proof. Every remaining loss is strictly positive. A separate incremental certificate in Section~\ref{sx:sec:six}
gives $q_0\ge\norm{\cdot}_{\Hh}^2/40$ on six exact raw zero functions.
At the same fixed Gamma split, its matrix response bound gives a projected
six-mode Schur margin $1/50$ and complete-high attachment margin $1/51$,
under the stated admission and enclosure inputs. These are restricted
weighted-domain bounds, not improvements of a global comparison.

Sections~\ref{sec:foundations}--\ref{sec:spectral} construct the common domains and actual spectral objects. Section~\ref{sec:bounds} retains the historical comparison proofs; Section~\ref{sec:verified-response} derives the newer response gate and its exact continuation. Section~\ref{sec:schur-unified} provides general finite reductions, and Section~\ref{sec:obstructions} gives local obstructions. Sections~\ref{sec:synthesis-unified}--\ref{sec:actual-jets} identify the full zero-jet core and exact remaining finite sign. Section~\ref{gb:sec:bridge} constructs the restricted Gamma operator and projected-jet response criterion on the actual reducing sector, while keeping its low spectral space potentially infinite dimensional. Sections~\ref{sec:complexity-unified}--\ref{sec:fixed-lattice} establish the restricted positivity regimes, the fixed-low-space sparse response estimate, computed four- and five-zero positive low spaces, the conditional five-mode graph-to-Schur transfer and its sixth-direction extension at the same split, both with complete high attachment, the uniform response brackets and fixed-debit obstruction, and the signed coupling criterion. Appendices~\ref{rtc:sec:tools} and~\ref{ext:sec:exterior} integrate fixed-resolvent cores, complete residual optimization, exact null repair, source-image leakage bounds, cofinal logic and exterior equality directions. All proof appendices are internal to this source. Computational evidence is identified separately; no earlier manuscript is used as a mathematical premise.

\section{Common analytic foundations}\label{sec:foundations}


\subsection{The theta kernel and the original closed energy}\label{a:sec:energy}

For $x\ge0$, define
\begin{equation}\label{a:eq:kappa}
 \kappa(x)=e^{x/2}\sum_{j\ge1}(4\tau_j(x)^2-6\tau_j(x))e^{-\tau_j(x)},
 \qquad \tau_j(x)=\pi j^2e^{2x},
\end{equation}
and extend evenly. We record the normalization and decay needed later.

\begin{lemma}[Smoothness and normalization]\label{a:lem:theta}
The kernel $\kappa$ is smooth, strictly positive, and even. There is a finite $C_\kappa>0$ such that
\begin{equation}\label{a:eq:envelopek}
 \kappa(x)\le C_\kappa\exp\!\left(\frac92|x|-\pi e^{2|x|}\right),
 \qquad
 C_\kappa=4\pi^2\sum_{j\ge1}j^4e^{-\pi(j^2-1)}.
\end{equation}
Furthermore $\int_\R2c\kappa\,dx=1$, and
\begin{equation}\label{a:eq:envelopeW}
 0<W(x)\le C_\kappa\exp(4|x|-\pi e^{2|x|}).
\end{equation}
\end{lemma}
\begin{proof}
The series in~\eqref{a:eq:kappa}, together with every derivative, converges uniformly on compact subsets of $[0,\infty)$. Each summand is positive because $\tau_j\ge\pi>3/2$. Dropping the negative polynomial term and using $e^{2x}\ge1$ gives~\eqref{a:eq:envelopek} for $x\ge0$.

To justify smooth evenness at zero, rather than assume it from reflection, set
\[
 H(x)=e^{x/2}\sum_{j\ge1}e^{-\pi j^2e^{2x}},\qquad x\in\R.
\]
The Jacobi transformation \cite[Eq.~20.7.32]{DLMF} implies
$H(-x)=H(x)+\sinh(x/2)$. Applying $D^2-1/4$ shows that $H''-H/4$ is even; direct differentiation identifies it with~\eqref{a:eq:kappa} on the positive half-line. Thus it gives the claimed smooth extension. Differentiating the transformation at zero yields $H'(0)=-1/4$. Since $c''=c/4$, integration by parts and the rapid decay of $H,H'$ at $+\infty$ give
\[
 \int_0^\infty c(x)\kappa(x)\,dx
  =[cH'-c'H]_0^\infty=\frac14.
\]
Evenness proves the normalization. Finally $2c(x)\ge e^{|x|/2}$ gives~\eqref{a:eq:envelopeW}.
\end{proof}

Put $\rho(s)=e^{-s/2}/(1-e^{-2s})$ for $s>0$, and $a_n=\Lambda(n)/\sqrt n$. The full positive energy is
\begin{align}
 \E_\Gamma[f]&=\int_0^\infty\rho(s)\int_\R
       \kappa(x)\kappa(x+s)|f(x+s)-f(x)|^2\,dx\,ds,\label{a:eq:Eg}\\
 \E_P[f]&=\sum_{n\ge2}a_n\int_\R
       \kappa(x)\kappa(x+s_n)|f(x+s_n)-f(x)|^2\,dx,\label{a:eq:Ep}\\
 \E[f]&=\E_\Gamma[f]+\E_P[f],\qquad
 \D=\{f\in\He:\E[f]<\infty\}.\label{a:eq:Edomain}
\end{align}
These are nonnegative integrals and sums, interpreted by Tonelli's theorem; no cancellation defines their domain. Write
\[
 \mu(f)=\int_\R f\,d\mu,\qquad
 \Var_\mu(f)=\norm{f}_\mu^2-|\mu(f)|^2.
\]

\begin{proposition}[Maximal domain and compact smooth form core]\label{a:prop:core}
The form $\E$ on $\He$ is densely defined, nonnegative, and closed. Its maximal integral domain $\D$ is the closure of even $C_c^\infty(\R)$ in the norm $(\norm{f}_\mu^2+\E[f])^{1/2}$. Constants belong to $\D$, and $\E[f]=0$ holds precisely for almost-everywhere constant $f$.
\end{proposition}
\begin{proof}
The theta envelope implies that for every $A_0>0$ there is a constant $C_{A_0}$ such that $\kappa(x)\le C_{A_0}e^{-A_0|x|}$. Consequently
\begin{equation}\label{a:eq:edgefinite}
 \int_\R\kappa(x)\kappa(x+s)\,dx
 \le C_{A_0}^2(s+A_0^{-1})e^{-A_0s},\qquad s\ge0.
\end{equation}
Since $a_n\le n$, the complete prime edge measure has finite mass, by taking $A_0>2$ in~\eqref{a:eq:edgefinite}. For every $L<\infty$,
\begin{equation}\label{a:eq:rhointegrable}
 \int_0^\infty\rho(s)\min\{1,L^2s^2\}\,ds<\infty.
\end{equation}
These observations admit compact smooth functions and constants to the domain.

For closedness, let $(f_j)$ be Cauchy in the form norm, with $L^2(\mu)$ limit $f$. A subsequence converges almost everywhere for $\mu$, hence for Lebesgue measure, because the density is strictly positive. A countable union of translated null sets is null, so all prime differences converge almost everywhere. Fubini's theorem gives convergence for the Gamma edge measure as well. For each fixed $k$, Fatou's lemma gives
\[
 \E[f-f_k]\le\liminf_{j\to\infty}\E[f_j-f_k].
\]
It follows that $f\in\D$ and $f_k\to f$ in form norm.

We prove the core assertion first for real profiles. Let $T_M(t)=\max\{-M,\min\{t,M\}\}$. The remainder $r_M(t)=t-T_M(t)$ is $1$-Lipschitz, tends pointwise to zero, and satisfies $|r_M(t)|\le|t|$. Thus dominated convergence in both edge measures and in $L^2(\mu)$ gives $T_M(f)\to f$ in form norm. It suffices to approximate bounded even profiles.

Choose smooth even cutoffs $0\le\chi_R\le1$ equal to one on $[-R,R]$, supported in $[-2R,2R]$, and with uniformly bounded Lipschitz constants. Equation~\eqref{a:eq:rhointegrable}, the finite prime-edge mass, and dominated convergence give $\E[\chi_R]\to0$. For bounded $f$, the identity for a product difference yields
\begin{align*}
 |[f(1-\chi_R)](y)-[f(1-\chi_R)](x)|^2
 &\le2(1-\chi_R(y))^2|f(y)-f(x)|^2\\
 &\quad+2\norm{f}_\infty^2|\chi_R(y)-\chi_R(x)|^2.
\end{align*}
Hence $f\chi_R\to f$ in form norm. It remains to smooth a compactly supported profile $g\in\D$.

Suppose $\supp g\subset[-K,K]$. Its ordinary $L^2$ norm is finite, since the density of $\mu$ is bounded below on this interval. For $0<s\le1$, a nonzero difference of $g$ has both endpoints in $[-K-1,K+1]$. The positive lower bound of $\kappa$ there controls the small-shift part of $\tg[g]$ by $\E_\Gamma[g]$. The part with $s>1$ is bounded by $4\norm{g}_2^2\int_1^\infty\rho(s)\,ds$. Thus $\tg[g]<\infty$.

Let $g_\epsilon=\zeta_\epsilon*g$, where $\zeta_\epsilon$ is an even compactly supported smooth approximate identity. Translations commute with convolution, and
\[
 \norm{\Delta_s(g_\epsilon-g)}_2\le2\norm{\Delta_sg}_2,
 \qquad\Delta_sg=g(\cdot+s)-g.
\]
Dominated convergence gives $\tg[g_\epsilon-g]\to0$. Since $\kappa$ is bounded,
$\E_\Gamma[g_\epsilon-g]\le\norm{\kappa}_\infty^2\tg[g_\epsilon-g]\to0$.
All the smoothed profiles have support in one fixed compact interval $J$. For any $h$ supported there,
\begin{equation}\label{a:eq:compactprimerate}
 \E_P[h]\le2\int_J\kappa(x)|h(x)|^2
     \sum_{n\ge2}a_n\bigl(\kappa(x-s_n)+\kappa(x+s_n)\bigr)\,dx
 \le C_J\norm{h}_2^2.
\end{equation}
Indeed, the complete rate in parentheses is uniformly bounded on $J$: use $\kappa(x\pm\log n)\le C_{A_0}e^{A_0\sup_J|x|}n^{-A_0}$ and $A_0>2$. Ordinary $L^2$ convergence therefore proves prime-energy convergence. The local upper bound of the density gives $L^2(\mu)$ convergence. Applying the argument to real and imaginary parts proves the complex core assertion.

Finally $\E[1]=0$. If $\E[f]=0$, the strictly positive Gamma density implies $f(y)=f(x)$ for almost every pair $x<y$. Fubini's theorem forces $f$ to be almost everywhere constant. Density of the core in $\He$ is immediate from the positive smooth density of $\mu$.
\end{proof}

The representation theorem for closed nonnegative forms \cite[Chapter~VI, Section~2]{Kato} defines $A$ by
\begin{equation}\label{a:eq:Adef}
 \begin{split}
 f\in\dom(A),\ Af=h\quad\Longleftrightarrow\quad
 f\in\D,\ h\in\He,\quad
 \E(f,g)=\ip{h}{g}_\mu\ \text{for all }g\in\D.
 \end{split}
\end{equation}
The operator is self-adjoint, $A1=0$, and $\ker A=\C1$. The centered space
$H_0=\{f\in\He:\mu(f)=0\}$ reduces $A$; its restriction is $A_E$. This construction uses only the nonnegative integral energy, before any identification with the physical form.


\subsection{Explicit bounds for the complete form}\label{a:sec:prime}

The bounds in this section account for all theta summands and all prime powers. A finite head is used only together with its complete remainder. We write $w_{\max}=\norm W_\infty$; the symbol $J$ will later denote a Neumann iteration index.

\begin{lemma}[Elementary theta bounds]\label{a:lem:explicitW}
For $r=|x|$,
\begin{equation}\label{a:eq:explicit_envelope}
 0<\kappa(x)\le42e^{(9/2)r-\pi e^{2r}},\qquad
 0<W(x)\le42e^{4r-\pi e^{2r}},\qquad w_{\max}<\frac12.
\end{equation}
For every integer $R\ge1$,
\begin{equation}\label{a:eq:rationalOmega}
 \Omega_R:=\sup_{|x|>R}W(x)
 <\overline\Omega(R):=42\,81^R(3/8)^{3\cdot7^R}.
\end{equation}
\end{lemma}
\begin{proof}
We use $3<\pi<22/7$, $\pi^2<10$, and the elementary exponential bounds
\begin{equation}\label{a:eq:elementarye}
 \frac{163}{60}=\sum_{j=0}^5\frac1{j!}<e
 <\frac{163}{60}+\frac1{720}\sum_{k\ge0}7^{-k}
 =\frac{11743}{4320}<\frac{87}{32}<\frac{11}{4}.
\end{equation}
In particular $e>27/10>8/3$, $e^2>7$, and $e^3>20$, since
$163^3=4330747>4320000=20\cdot60^3$.
The theta coefficient in Lemma~\ref{a:lem:theta} is less than $42$. Indeed, the tail of
$\sum_{k\ge1}k^4e^{-\pi(k^2-1)}$ starts below $16e^{-9}$, and its successive ratios, for $k\ge2$, are at most
$(3/2)^4e^{-15}<1/2$. Thus that sum is below
$1+32e^{-9}<26/25$, and
$4(22/7)^2(26/25)<42$.
This proves the two envelopes in~\eqref{a:eq:explicit_envelope}.

For the sharper supremum bound put $F(t)=(4t^2-6t)e^{-t}$. Since
\[
 F'(t)=-2(2t-1)(t-3)e^{-t},\qquad
 W(x)=\frac{\sum_{k\ge1}F(\pi k^2e^{2r})}{1+e^{-r}},
\]
$F$ decreases for $t\ge3$. The complete tail with $k\ge2$ satisfies
\[
 H(r):=\sum_{k\ge2}F(\pi k^2e^{2r})
 <\sum_{k\ge2}40k^4e^{-3k^2}<1280e^{-12}<\frac1{125}.
\]
The majorant starts at $640e^{-12}$ and has successive ratio at most
$(3/2)^4e^{-15}<1/2$; the last inequality uses $e^{12}>20^4=160000$.
Consider three fixed radial intervals. If $0\le r\le1/6$, then
$F(\pi e^{2r})\le F(3)=18e^{-3}<9/10$ and
$1+e^{-r}\ge2-r\ge11/6$, giving
\[
 W(x)<\frac{9/10+1/125}{11/6}=\frac{681}{1375}<\frac12.
\]
If $1/6\le r\le1/2$, then $\pi e^{2r}>4$,
$F(\pi e^{2r})<F(4)=40e^{-4}<20/27$, since $e^4>54$,
and $1+e^{-r}\ge3/2$. Therefore
\[
 W(x)<\frac{20/27+1/125}{3/2}=\frac{5054}{10125}<\frac12.
\]
Here $3/4-20/27=1/108>1/125$. Finally, if $r\ge1/2$, then
$\pi e^{2r}>3e>8$ and $e^8=e^6e^2>2800$, so
\[
 W(x)<208e^{-8}+\frac1{125}
 <\frac{13}{175}+\frac1{125}=\frac{72}{875}<\frac12.
\]
These uniform rational margins prove $w_{\max}<1/2$.

For~\eqref{a:eq:rationalOmega}, the function $4r-3e^{2r}$ decreases for $r\ge0$. Use $\pi>3$, $e^{4R}<81^R$, $e^{2R}>7^R$, and $e^{-1}<3/8$ in the envelope.
\end{proof}

\begin{lemma}[Complete weighted prime tail]\label{a:lem:prime}
For $N\ge2$, define
\begin{equation}\label{a:eq:rN}
 r_N=96(N+1)^3(27/512)^{N+1}.
\end{equation}
The full prime form is a bounded Hermitian form on $\Hw$, and
\begin{equation}\label{a:eq:headtail}
 |p[u]-p_N[u]|\le r_N\norm u_{\Hh}^2,
 \qquad |p_N[u]|\le2\sum_{2\le n\le N}a_n\norm u_2^2.
\end{equation}
In particular,
\begin{equation}\label{a:eq:head4}
 r_4=\frac{5380840125}{1099511627776}<\frac1{128},
 \qquad 0<\cg<6,
 \qquad \cg\norm u_2^2+|p_4[u]|\le9\norm u_2^2.
\end{equation}
Consequently the form
\begin{equation}\label{a:eq:Dform}
 d(u,z)=\cg\ip uz_2+p(u,z)
\end{equation}
is represented by a bounded self-adjoint operator $D$ on $\Hw$ with
\begin{equation}\label{a:eq:C0}
 \norm D\le C_0:=9+r_4<10.
\end{equation}
\end{lemma}
\begin{proof}
Let $\beta=\sqrt W$ and let $\tau_sF(x)=F(x+s)$ on ordinary $L^2$. With $z=|x|+|x+s|\ge s$, the envelope and the arithmetic--geometric mean inequality give
\begin{align*}
 \beta(x)\beta(x+s)
 &\le42\exp\left(2z-\frac\pi2(e^{2|x|}+e^{2|x+s|})\right)\\
 &\le42e^{2z-\pi e^z}\le42e^{2s-\pi e^s}.
\end{align*}
The last exponent decreases for $z\ge0$. Hence the self-adjoint series
\begin{equation}\label{a:eq:KP}
 K_P=\sum_{n\ge2}a_nM_\beta(\tau_{s_n}+\tau_{-s_n})M_\beta
\end{equation}
converges absolutely in operator norm: its $n$th term has norm at most
$84\Lambda(n)n^{3/2}e^{-\pi n}$.
Reflection interchanges the two shifts, so the sum preserves even functions. Substituting $u=\beta F$ gives
\begin{equation}\label{a:eq:prepresentation}
 p[u]=\ip{K_P(u/\beta)}{u/\beta}_2.
\end{equation}
This also proves absolute convergence of the original scalar series.

Using $\Lambda(n)\le\log n\le n$ and $n^{5/2}\le n^3$, its tail is bounded by
$84\sum_{n>N}n^3e^{-3n}$. For $n\ge3$ the successive ratio is below
$(4/3)^3(3/8)^3=1/8$. Summing the geometric majorant and using
$e^{-3}<27/512$ yields~\eqref{a:eq:rN}. Translation invariance and Cauchy--Schwarz give the finite-head estimate.

For the small head, the elementary bounds
\[
 \log2<7/10,\quad \log3<11/10,\quad
 \sqrt2>7/5,\quad \sqrt3>22/13
\]
give $a_2<1/2$, $a_3<13/20$, $a_4<7/20$, and thus
$2(a_2+a_3+a_4)<3$. The logarithm inequalities follow, for example, by keeping the first five terms of $e^{7/10}$ and the first six terms of $e^{11/10}$; these exceed $2$ and $3$, respectively. The square-root comparisons follow by squaring.
The digamma reflection and duplication formulas, together with $\psi(1)=-\gamma$, give
\[
 \cg=\log\pi+\gamma+\frac\pi2+3\log2.
\]
These classical formulas are recorded in \cite[Sections 5.4--5.5]{DLMF}. In particular $\cg>0$. For the upper bound, use $0<\gamma<1$, $\pi/2<11/7$, $\log2<7/10$, and $\log\pi<6/5$. The last inequality follows from $\pi<22/7<e^{6/5}$, using the first five terms of the exponential series. Their sum gives
$\cg<6/5+1+11/7+21/10<6$.
Substitution $N=4$ in~\eqref{a:eq:rN} yields the fraction in~\eqref{a:eq:head4}; comparison with $1/128$ is immediate by cross multiplication.
Finally~\eqref{a:eq:head4} and $w_{\max}<1$ imply
$|d[u]|\le(9+r_4)\norm u_{\Hh}^2$. Polarization and the bounded-form representation theorem give $D$ and~\eqref{a:eq:C0}.
\end{proof}

The operator $K_P$ is not claimed to be compact. A decaying multiplication operator composed with a translation need not be compact on a nonatomic $L^2$ space. The bounded operator $D$ acts in the weighted space; its existence does not assert that the full unweighted prime row belongs to ordinary $L^2$.


\subsection{The explicit formula and the compact physical bridge}\label{a:sec:bridge}

We now prove the identity that connects the positive energy to the physical form. The classical explicit formula is used as a theorem, without any assumption about the location of nontrivial zeta zeros inside the critical strip. We first fix its normalization, then justify its use with the noncompact theta kernel, and finally perform the ground-state calculation with all endpoints retained.

In this section use the nonunitary Fourier transform
\[
 \mathcal Fu(z)=\int_\R u(x)e^{izx}\,dx,
 \qquad m(u)=\int_\R c(x)u(x)\,dx.
\]
Let $\rho_\zeta$ range over the distinct nontrivial zeros of $\zeta$, with multiplicity $\nu_{\rho_\zeta}$, and put $w_{\rho_\zeta}=(\rho_\zeta-1/2)/i$. Thus $|\operatorname{Im}w_{\rho_\zeta}|<1/2$; no reality of $w_{\rho_\zeta}$ is asserted. We use the classical zero-counting bound
\begin{equation}\label{a:eq:zerocount}
 \sum_{|\operatorname{Im}\rho_\zeta|\le T}\nu_{\rho_\zeta}
       =O(T\log(T+2))\qquad(T\ge2).
\end{equation}
This weak form of the classical Riemann--von Mangoldt estimate \cite{vonMangoldt} is sufficient for every zero sum below.

\begin{theorem}[Classical Guinand explicit formula]\label{a:thm:guinand}
For $g\in C_c^\infty(\R)$,
\begin{align}\label{a:eq:guinand}
 \sum_{\rho_\zeta}\nu_{\rho_\zeta}\mathcal Fg(w_{\rho_\zeta})
 &=\mathcal Fg(i/2)+\mathcal Fg(-i/2)\notag\\
 &\quad+\frac1{2\pi}\int_\R
   \bigl(\operatorname{Re}\psi(1/4+it/2)-\log\pi\bigr)
       \mathcal Fg(t)\,dt\notag\\
 &\quad-\sum_{n\ge2}a_n\bigl(g(s_n)+g(-s_n)\bigr).
\end{align}
The zero sum is absolutely convergent. This is the classical formula in the Fourier convention displayed here \cite{Guinand}, with prime powers collected by the von Mangoldt function.
\end{theorem}

The statement is quoted, not proposed as a new result. Absolute convergence follows from the rapid decay of $\mathcal Fg$ in the closed strip $|\operatorname{Im}z|\le1/2$ and~\eqref{a:eq:zerocount}.

\begin{lemma}[Bilinear physical normalization]\label{a:lem:bilinear}
For real even $u,v\in C_c^\infty(\R)$,
\begin{equation}\label{a:eq:bilinearzeros}
 q_0(u,v)+2m(u)m(v)
   =\sum_{\rho_\zeta}\nu_{\rho_\zeta}
       \mathcal Fu(w_{\rho_\zeta})\mathcal Fv(w_{\rho_\zeta}).
\end{equation}
\end{lemma}
\begin{proof}
The digamma integral \cite[Eq.~5.9.16]{DLMF} gives, after subtracting its value at $1/4$ and setting its integration variable equal to $2s$,
\begin{equation}\label{a:eq:digammamultiplier}
 \operatorname{Re}\psi(1/4+it/2)-\psi(1/4)
    =2\int_0^\infty\rho(s)(1-\cos(ts))\,ds=:d_\Gamma(t).
\end{equation}
The difference integral is convergent at zero and infinity. Plancherel and polarization identify
\[
 \tg(u,v)=\frac1{2\pi}\int_\R
       d_\Gamma(t)\mathcal Fu(t)\mathcal Fv(t)\,dt
\]
for real even tests. Apply Theorem~\ref{a:thm:guinand} to $g=u*v$. Its Fourier transform is the product of the two transforms, and $g(0)=\int uv$. Evenness gives
\[
 \mathcal Fu(i/2)=\mathcal Fu(-i/2)=m(u),\qquad
 g(s_n)+g(-s_n)
  =\int_\R\bigl(u(x+s_n)v(x)+u(x)v(x+s_n)\bigr)\,dx.
\]
These are respectively the pole and the polarized prime term. The constant in the Gamma term is $\psi(1/4)-\log\pi=-\cg$, proving~\eqref{a:eq:bilinearzeros}.
\end{proof}

For complex even tests the zero-side product is
$\mathcal Fu(w)\overline{\mathcal Fv(\overline w)}$. It is not generally $|\mathcal Fu(w)|^2$ when $u=v$ and $w$ is nonreal. We need only the real bilinear identity and its conjugate-linear extension in the second argument; no positivity is inferred from this zero sum.

\begin{lemma}[Entire theta transform]\label{a:lem:thetaxi}
Let
\[
 \xi(s)=\tfrac12s(s-1)\pi^{-s/2}\Gamma(s/2)\zeta(s).
\]
Then the entire continuation of this function satisfies
\begin{equation}\label{a:eq:thetaxi}
 \mathcal F\kappa(z)=\xi(1/2+iz),\qquad z\in\C.
\end{equation}
In particular, $\mathcal F\kappa(w_{\rho_\zeta})=0$ for every nontrivial zero, whether or not it is on the critical line.
\end{lemma}
\begin{proof}
Write $\Theta(t)=\sum_{j\ge1}e^{-\pi j^2t}$, so that the function in Lemma~\ref{a:lem:theta} is $H(x)=e^{x/2}\Theta(e^{2x})$. For $\operatorname{Re}s>1$, termwise integration gives
\[
 \pi^{-s/2}\Gamma(s/2)\zeta(s)
   =\int_0^\infty\Theta(t)t^{s/2-1}\,dt.
\]
The Jacobi transformation is
$\Theta(t)=t^{-1/2}\Theta(1/t)+(t^{-1/2}-1)/2$.
Split the integral at one and substitute $t\mapsto1/t$ in its lower part. The elementary remainder integrates to $1/[s(s-1)]$, and hence
\begin{align}\label{a:eq:mellintheta}
 \pi^{-s/2}\Gamma(s/2)\zeta(s)
 &=\frac1{s(s-1)}+I(s),\notag\\
 I(s)&=\int_1^\infty\Theta(t)
       \bigl(t^{s/2-1}+t^{(1-s)/2-1}\bigr)\,dt.
\end{align}
The integral $I$ is entire by local uniform double-exponential domination after $t=e^{2x}$. Thus
$\xi(s)=1/2+s(s-1)I(s)/2$ is entire. For $s=1/2+iz$, the same substitution gives
\begin{equation}\label{a:eq:xiH}
 \xi(1/2+iz)=\frac12
       -2(z^2+1/4)\int_0^\infty H(x)\cos(zx)\,dx.
\end{equation}
Because $\kappa=H''-H/4$ on the positive half-line, evenness and two integrations by parts give
\[
 \mathcal F\kappa(z)
  =-2H'(0)-2(z^2+1/4)\int_0^\infty H(x)\cos(zx)\,dx.
\]
All boundary terms at infinity vanish for each complex $z$. The value $H'(0)=-1/4$ from Lemma~\ref{a:lem:theta} proves~\eqref{a:eq:thetaxi}. The defining prefactors of $\xi$ are finite and nonzero at a nontrivial zero of $\zeta$, so the last assertion follows.
\end{proof}

\begin{proposition}[Weak theta source identity]\label{a:prop:thetanull}
For every even $h\in C_c^\infty(\R)$,
\begin{equation}\label{a:eq:thetanull}
 q_0(\kappa,h)=-\overline{m(h)}.
\end{equation}
All terms on the left exist separately as convergent form pairings. Equivalently, $(D_\Gamma-\cg-P)\kappa=-c$ on even compact smooth tests. This is a weak identity and does not assert that $c$ belongs to ordinary $L^2$.
\end{proposition}
\begin{proof}
We first take $h$ real. Let $\chi_R(x)=\chi(x/R)$, $R\ge1$, where $\chi$ is even and smooth, equals one on $[-1,1]$, has support in $[-2,2]$, and satisfies $0\le\chi\le1$. Put $k_R=\kappa\chi_R$. Differentiating the theta series on each half-line shows that $\kappa$ and its derivatives decrease faster than every exponential. Consequently $k_R\to\kappa$ in $H^1(\R)$ and in $L^1(e^{|x|/2}dx)$. For $v\in H^1(\R)$,
\[
 \tg[v]\le\norm{v'}_2^2\int_0^1\rho(s)s^2\,ds
              +4\norm{v}_2^2\int_1^\infty\rho(s)\,ds.
\]
It follows that the Gamma pairing, the ordinary $L^2$ pairing, and the moment in~\eqref{a:eq:bilinearzeros} converge as $R\to\infty$.

The complete prime pairing also converges without a spatially uniform incoming-rate assumption. If $\supp h\subset[-K,K]$, changing variables in one of its two terms gives
\[
 p(k_R,h)=\sum_{n\ge2}a_n\int_{-K}^K h(x)
           \bigl(k_R(x+s_n)+k_R(x-s_n)\bigr)\,dx.
\]
For any fixed $A_0>2$,
\begin{equation}\label{a:eq:bridgeprimerate}
 \sup_{|x|\le K}\sum_{n\ge2}a_n
       \bigl(\kappa(x+s_n)+\kappa(x-s_n)\bigr)
 \le2C_{A_0}e^{A_0K}\sum_{n\ge2}a_nn^{-A_0}<\infty.
\end{equation}
Here $a_n\le n$ suffices. Since $|k_R|\le\kappa$, dominated convergence proves the assertion and absolute convergence of the limiting pairing.

For the zero side, uniformly in $R$ and in $|\operatorname{Im}z|\le1/2$,
\[
 |\mathcal Fk_R(z)|\le\int_\R\kappa(x)e^{|x|/2}\,dx<\infty.
\]
Repeated integration by parts in the transform of the fixed compact test gives
$|\mathcal Fh(z)|\le C_{h,N}(1+|\operatorname{Re}z|)^{-N}$ throughout this strip. Equation~\eqref{a:eq:zerocount} therefore supplies a summable majorant for the products in~\eqref{a:eq:bilinearzeros}, for example with $N=3$. Pointwise transform convergence and Lemma~\ref{a:lem:thetaxi} make their limiting sum zero. Thus
\[
 q_0(\kappa,h)+2m(\kappa)m(h)=0.
\]
Lemma~\ref{a:lem:theta} gives $m(\kappa)=1/2$, proving~\eqref{a:eq:thetanull} for real $h$. Splitting an arbitrary $h$ into real and imaginary parts and using conjugate linearity in the second argument proves the stated version.
\end{proof}

\begin{theorem}[Compact physical bridge]\label{a:thm:compactbridge}
For every even $f\in C_c^\infty(\R)$,
\begin{equation}\label{a:eq:provedbridge}
 q_0[\kappa f]=\E[f]-\frac12\norm{f}_\mu^2.
\end{equation}
\end{theorem}
\begin{proof}
Set $h=\kappa|f|^2$, which is a real even compact smooth test. For $a=\kappa(x)$, $b=\kappa(x+s)$, $F=f(x)$, and $G=f(x+s)$, the pointwise identity
\[
 |bG-aF|^2-(b-a)(b|G|^2-a|F|^2)=ab|G-F|^2
\]
gives, after integration,
\begin{equation}\label{a:eq:gammaground}
 \tg[\kappa f]-\tg(\kappa,h)=\E_\Gamma[f].
\end{equation}
This is an equality between finite quantities: $\kappa$, $h$, and $\kappa f$ are in $H^1$, their Gamma pairings are finite by the estimate in the preceding proof, and $\E_\Gamma[f]$ is finite by Proposition~\ref{a:prop:core}.

For the primes the two endpoint terms are both essential. Their full expansion gives
\begin{align}\label{a:eq:primeground}
 p[\kappa f]-p(\kappa,h)
 &=\sum_{n\ge2}a_n\int_\R\kappa(x)\kappa(x+s_n)
       \bigl(2\operatorname{Re}(f(x+s_n)\overline{f(x)})\notag\\
 &\hspace{44mm}-|f(x+s_n)|^2-|f(x)|^2\bigr)\,dx\notag\\
 &=-\E_P[f].
\end{align}
The compact-profile prime energy is finite, the compact correlation $p[\kappa f]$ has only finitely many nonzero shifts, and $p(\kappa,h)$ is absolutely convergent by~\eqref{a:eq:bridgeprimerate}. Thus no difference of divergent prime expressions is used. Finally,
$\norm{\kappa f}_2^2=\ip{\kappa}{h}_2$ cancels the constant term. Combining~\eqref{a:eq:gammaground}, \eqref{a:eq:primeground}, and~\eqref{a:eq:thetanull} yields
\[
 q_0[\kappa f]=\E[f]-m(h)
             =\E[f]-\int_\R c\kappa|f|^2\,dx,
\]
which is~\eqref{a:eq:provedbridge}.
\end{proof}

The theorem uses the classical explicit formula and the theta transform, not a zero-location hypothesis. It establishes an identity of forms, not a positivity statement for $q_0$ or for the centered energy shifted by $1/2$.

\subsection{The exact numerator domain}\label{a:sec:domain}

We first analyze $X$ directly from its two defining nonnegative terms. Because $W$ is positive and bounded,
\begin{equation}\label{a:eq:L2bound}
 \norm{u}_2^2\le\norm{W}_\infty\norm{u}_{\Hh}^2.
\end{equation}
The Gamma form is closed on ordinary $L^2$, for example by the same difference-and-Fatou argument used for $\E$.

\begin{lemma}\label{a:lem:Xcore}
The space $X$ of~\eqref{a:eq:X} is Hilbert, and even $C_c^\infty(\R)$ is dense in $X$.
\end{lemma}
\begin{proof}
A sequence Cauchy in $X$ has a weighted $L^2$ limit and an ordinary Gamma-form limit. Equation~\eqref{a:eq:L2bound} identifies their ordinary $L^2$ limits, proving completeness.

Use the cutoffs $\chi_R$ from Proposition~\ref{a:prop:core}. For $h_R=(1-\chi_R)u$ the product identity gives
\begin{align*}
 |h_R(x+s)-h_R(x)|^2
 &\le2(1-\chi_R(x+s))^2|u(x+s)-u(x)|^2\\
 &\quad+2|u(x)|^2|\chi_R(x+s)-\chi_R(x)|^2.
\end{align*}
The first term tends to zero in the Gamma integral by dominated convergence. The second is dominated by
$2|u(x)|^2\min\{1,L^2s^2\}$, which is integrable by~\eqref{a:eq:rhointegrable}. Weighted convergence is another application of dominated convergence. Thus compactly supported vectors are dense in $X$.

For a fixed compactly supported vector, even smooth mollification converges in ordinary $L^2$ and in $\tg$, by the translation argument in Proposition~\ref{a:prop:core}. The supports remain in one compact interval, where $1/W$ is bounded. Ordinary $L^2$ convergence consequently gives weighted convergence as well.
\end{proof}

For original profiles and physical numerators the coordinates are
\begin{equation}\label{a:eq:coordinates}
 u=\kappa f,\qquad v=2cf=\frac uW,\qquad
 F=\frac u{\sqrt W}=\sqrt{2c\kappa}\,f.
\end{equation}
Only $f\mapsto F$ is the displayed unitary map from $\He$ to ordinary $\Le$. The map $f\mapsto u$ is an isometry to the weighted space with norm $\norm{\cdot}_{\Hh}$, since
\begin{equation}\label{a:eq:weightedidentity}
 \norm{\kappa f}_{\Hh}=\norm{f}_\mu.
\end{equation}

\begin{proposition}[Extension of the compact bridge]\label{a:prop:domain}
Multiplication by $\kappa$ gives a bijection $\D\to X$, with equivalent graph norms, and
\begin{equation}\label{a:eq:extendedbridge}
 q_0(\kappa f,\kappa g)=\E(f,g)-\tfrac12\ip{f}{g}_\mu,
 \qquad f,g\in\D.
\end{equation}
\end{proposition}
\begin{proof}
Put $C_{\mathrm b}=|\cg|\norm{W}_\infty+\norm{K_P}$. On the compact smooth core, Theorem~\ref{a:thm:compactbridge}, Lemma~\ref{a:lem:prime}, and~\eqref{a:eq:weightedidentity} imply
\begin{align}
 \tg[\kappa f]&\le\E[f]+(C_{\mathrm b}+\tfrac12)\norm{f}_\mu^2,\label{a:eq:graph1}\\
 \E[f]&\le\tg[\kappa f]+(C_{\mathrm b}+\tfrac12)\norm{f}_\mu^2.\label{a:eq:graph2}
\end{align}
Taking compact-profile approximants from Proposition~\ref{a:prop:core} and using completeness of $X$ proves that $\kappa\D\subset X$.

Conversely, for $u\in X$ choose even $u_j\in C_c^\infty$ converging in $X$, by Lemma~\ref{a:lem:Xcore}. Since $\kappa$ is smooth and strictly positive, $f_j=u_j/\kappa$ are legitimate even compact smooth profiles. Equation~\eqref{a:eq:graph2} makes $(f_j)$ Cauchy in the original energy graph norm. Its weighted limit is $f=u/\kappa$, so closedness of $\E$ gives $f\in\D$. This proves surjectivity and graph-norm equivalence. The Gamma form and the prime form are continuous in $X$, so the quadratic identity passes to the limit; polarization gives~\eqref{a:eq:extendedbridge}.
\end{proof}

Recall $m(u)=\int_\R c(x)u(x)\,dx$. The integral exists absolutely for finite weighted norm, because
\begin{equation}\label{a:eq:meanbound}
 |m(u)|\le\left(\int_\R c^2W\,dx\right)^{1/2}\norm{u}_{\Hh}
       =\tfrac12\norm{u}_{\Hh}.
\end{equation}
Here $\int c^2W=\tfrac12\int c\kappa=1/4$. Thus $\mu(f)=2m(\kappa f)$, and the full pole form obeys
\begin{equation}\label{a:eq:polebridge}
 q[u]:=q_0[u]+2|m(u)|^2
       =\E[f]-\tfrac12\Var_\mu(f),\qquad u=\kappa f\in X.
\end{equation}
The coefficient $2$ in the pole is fixed by normalization. It is not an additional positive contribution available after centering.


\subsection{Complex theta continuation and bounded Fourier evaluations}\label{d:sec:theta}
The real-line normalization has already been proved. We record the additional complex estimates needed for spectral synthesis.


\begin{lemma}[Theta continuation and tails]\label{d:lem:theta-tails}
The function $\kappa$ has an even holomorphic continuation to $|\operatorname{Im}z|<\pi/4$.
There are constants $0<A_0\le A_1<\infty$ such that
\begin{equation}\label{d:eq:W-two-sided}
 A_0 e^{4x-\pi e^{2x}}\le W(x)\le A_1 e^{4x-\pi e^{2x}}
 \qquad(x\ge0).
\end{equation}
For each nonnegative integer $j$ and each $\beta_0<\pi/4$,
\begin{equation}\label{d:eq:strip-upper}
 |\kappa^{(j)}(x+i\beta)|\le C_{j,\beta_0}
 \exp\!\left((9/2+2j)|x|-\pi\cos(2\beta_0)e^{2|x|}\right)
 \quad(|\beta|\le\beta_0).
\end{equation}
For fixed $b$ in the strip, as $x\to+\infty$,
\begin{equation}\label{d:eq:complex-asymptotic}
 \kappa(x+b)=4\pi^2e^{9(x+b)/2}e^{-\pi e^{2x}e^{2b}}
 \bigl(1+O(e^{-2x})\bigr).
\end{equation}
The last error estimate is uniform when $b$ ranges over a compact subset of the strip.
\end{lemma}
\begin{proof}
Let
\[
 \Psi(t)=\sum_{n\ge1}e^{-\pi n^2t}\quad(\operatorname{Re}t>0),\qquad
 h(z)=e^{z/2}\Psi(e^{2z}).
\]
The Jacobi transformation~\cite[Eq.~20.7.32]{DLMF} gives, first on the real line,
$h(-x)=h(x)+\sinh(x/2)$. Both sides are holomorphic in the strip, so the identity holds there. Direct differentiation gives
$\kappa=h''-h/4$ on $x\ge0$. Since $(D^2-1/4)\sinh(x/2)=0$, this formula defines an even holomorphic extension and agrees with the stipulated reflection. All derivatives of the defining series converge locally uniformly in the strip.

For a lower bound in~\eqref{d:eq:W-two-sided}, keep the $n=1$ term and use
$4\tau^2-6\tau\ge(4-6/\pi)\tau^2$ for $\tau\ge\pi$, together with
$e^{x/2}\le2c(x)\le2e^{x/2}$. For the upper bound, discard the negative term and use
\[
 \sum_{n\ge1} n^4e^{-\pi n^2e^{2x}}
 \le e^{-\pi e^{2x}}\sum_{n\ge1}n^4e^{-\pi(n^2-1)}.
\]
The series on the right converges. Termwise differentiation introduces only fixed polynomial powers of $n^2e^{2z}$ and fixed exponential factors. The same estimate with $\pi\cos(2\beta_0)$ in the exponential proves~\eqref{d:eq:strip-upper}; evenness handles negative $x$.

For~\eqref{d:eq:complex-asymptotic}, the lower-degree part of the $n=1$ polynomial has relative size $O(e^{-2x})$. The terms $n\ge2$ have, on a compact set of $b$, an additional factor bounded by $C\exp(-c e^{2x})$ with $c>0$. This proves the stated relative estimate. In particular,
\begin{equation}\label{d:eq:W-asymptotic}
 W(x)=4\pi^2e^{4x-\pi e^{2x}}\bigl(1+O(e^{-x})\bigr)
 \quad(x\to+\infty).
\end{equation}
\end{proof}


\begin{lemma}[Entire transform and bounded evaluations]\label{d:lem:Fourier}
For $u\in\Hh$, the function $F_u$ is even entire, the map $u\mapsto F_u$ is injective, and
\begin{equation}\label{d:eq:eval-bound}
 |F_u^{(k)}(z)|\le\norm u_{\Hh}
 \left(\int_\R W(x)|x|^{2k}e^{2B|x|}\,dx\right)^{1/2}
 \quad (|\operatorname{Im}z|\le B).
\end{equation}
Moreover
\begin{equation}\label{d:eq:Xi-xi}
 \Xi(z):=F_\kappa(z)=\xi(1/2+iz),\qquad
 \Xi(0)>0,\quad \Xi(i/2)=\Xi(-i/2)=\tfrac12.
\end{equation}
Its zero set is discrete and invariant under $w\mapsto-w$ and $w\mapsto\overline w$, preserving multiplicity.
\end{lemma}
\begin{proof}
Cauchy--Schwarz gives~\eqref{d:eq:eval-bound}; the two-sided theta estimate makes every displayed integral finite and justifies locally uniform differentiation. It also puts $u$ in $L^1(e^{B|x|}dx)$ for every $B$. Ordinary Fourier uniqueness proves injectivity. Lemma~\ref{a:lem:thetaxi} already proves the exact completed transform, including its Mellin calculation. Positivity, evenness, real coefficients and normalization give the remaining assertions.
\end{proof}


\subsection{The exact Gamma form and its weighted representative}\label{d:sec:Gamma-form}

We now construct the weighted Gamma representative and its ambient graph core. These facts will serve both the physical realization and the later stronger-topology synthesis theorem.

Retain the same theta weight $W$ and put
\begin{align}
 \rho(s)&=\frac{e^{-s/2}}{1-e^{-2s}},\qquad s>0,\label{d:eq:Gamma-rho}\\
 \tg[u]&=\int_0^\infty\rho(s)
       \norm{u(\cdot+s)-u}_{L^2(\R)}^2\,ds,\label{d:eq:Gamma-form}\\
 {X}&=\{u\in\Hh:\tg[u]<\infty\},\qquad
 \norm{u}_{{X}}^2=\norm{u}_{\Hh}^2+\tg[u].
 \label{d:eq:Gamma-X}
\end{align}
The form is polarized with the linear-first convention. Although its vectors are even, its translation differences are taken on the whole real line. All positive shifts in~\eqref{d:eq:Gamma-form} are retained.

Write $M=\norm{W}_\infty<\infty$. The inclusion
$\norm{u}_2\le\sqrt M\norm{u}_{\Hh}$ is continuous. The kernel satisfies $\rho(s)=O(s^{-1})$ at zero and $O(e^{-s/2})$ at infinity, so
\begin{equation}\label{d:eq:Gamma-J}
 J_j=\int_0^\infty s^j\rho(s)\,ds<\infty\qquad(j=1,2).
\end{equation}
Use the unitary ordinary Fourier transform in the next two sections. Plancherel and Tonelli give, including infinite values,
\begin{equation}\label{d:eq:Gamma-multiplier}
 \tg[u]=\int_\R d_\Gamma(\xi)|\widehat u(\xi)|^2\,d\xi,
 \qquad d_\Gamma(\xi)=2\int_0^\infty\rho(s)(1-\cos(\xi s))\,ds.
\end{equation}
The inequalities $1-\cos y\le |y|$ and
$\norm{u(\cdot+s)-u}_2\le s\norm{u'}_2$ imply
\begin{equation}\label{d:eq:Gamma-first-order}
 0\le d_\Gamma(\xi)\le2J_1|\xi|,\qquad
 \tg[u]\le J_2\norm{u'}_2^2
 \quad(u\in H^1(\R)).
\end{equation}
Here $H^1(\R)$ always means the ordinary unweighted Sobolev space. These are upper estimates, not replacements of the exact Gamma form by a first-order form.

\begin{proposition}[Closed form and compact smooth form core]\label{d:prop:Gamma-form-core}
The form~\eqref{d:eq:Gamma-form} is densely defined, nonnegative, and closed in $\Hh$. The even compact smooth functions are dense in ${X}$. It has a nonnegative self-adjoint representative $T_\Gamma$ in $\Hh$.
\end{proposition}
\begin{proof}
The domain and compact smooth form-core assertion are Lemma~\ref{a:lem:Xcore}. The multiplier identity above proves nonnegativity and closedness equivalently. The representation theorem gives $T_\Gamma$; no graph-core conclusion is inferred from form density.
\end{proof}

Let $D_\Gamma$ be the ordinary nonnegative Fourier multiplier with symbol $d_\Gamma$ and its maximal ordinary $L^2$ multiplier domain. Fix any bounded self-adjoint operator $\mathsf D$ on $\Hh$, and define
\begin{equation}\label{d:eq:Gamma-graph}
 B_{\mathsf D}=T_\Gamma-\mathsf D,\qquad
 \mathcal G_{\mathsf D}=\operatorname{Dom}T_\Gamma,
 \qquad
 \norm{u}_{\mathcal G_{\mathsf D}}^2
 =\norm{u}_{\Hh}^2+\norm{B_{\mathsf D}u}_{\Hh}^2.
\end{equation}
The bounded self-adjoint perturbation is self-adjoint on $\operatorname{Dom}T_\Gamma$, and the two graph norms are equivalent, since
\[
 \norm{B_{\mathsf D}u}_{\Hh}\le\norm{T_\Gamma u}_{\Hh}
       +\norm{\mathsf D}\norm{u}_{\Hh},\qquad
 \norm{T_\Gamma u}_{\Hh}\le\norm{B_{\mathsf D}u}_{\Hh}
       +\norm{\mathsf D}\norm{u}_{\Hh}.
\]
No relation between $\mathsf D$ and $\Kk$ is assumed.

\begin{lemma}[Actual graph admission from one ordinary derivative]\label{d:lem:Gamma-H1}
If $u\in\Hh\cap H^1(\R)$ is even, then $u\in\operatorname{Dom}T_\Gamma$ and
\begin{equation}\label{d:eq:Gamma-H1}
 T_\Gamma u=W D_\Gamma u,\qquad
 \norm{B_{\mathsf D}u}_{\Hh}
 \le2\sqrt M J_1\norm{u'}_2+\norm{\mathsf D}\norm{u}_{\Hh}.
\end{equation}
Consequently, for either $E={X}$ or $E=\mathcal G_{\mathsf D}$,
\begin{equation}\label{d:eq:Gamma-E-upper}
 \norm{u}_{\Hh}\le\norm{u}_E
 \le C_E\bigl(\norm{u}_{\Hh}+\norm{u}_{H^1(\R)}\bigr)
 \qquad(u\in\Hh\cap H^1(\R)).
\end{equation}
\end{lemma}
\begin{proof}
Equation~\eqref{d:eq:Gamma-first-order} first gives $u\in{X}$ and $D_\Gamma u\in L^2$. The vector $g=WD_\Gamma u$ belongs to $\Hh$ because
\[
 \norm{g}_{\Hh}^2=\int_\R W|D_\Gamma u|^2
 \le M\norm{D_\Gamma u}_2^2.
\]
For every $z\in{X}$, the ordinary multiplier form identity yields
$\tg(u,z)=\ip{D_\Gamma u}{z}_2=\ip{g}{z}_{\Hh}$.
The representation criterion therefore admits $u$ to the actual weighted operator domain and identifies $T_\Gamma u=g$. Now~\eqref{d:eq:Gamma-first-order} and boundedness of $\mathsf D$ prove~\eqref{d:eq:Gamma-H1}. The form and graph estimates imply~\eqref{d:eq:Gamma-E-upper}. This is an admission argument, not an assumed maximal-domain identity for a formal product.
\end{proof}

\subsection{A direct compact smooth graph core}\label{d:sec:Gamma-graph-core}

We need compact smooth density in the graph norm itself. The form core from Proposition~\ref{d:prop:Gamma-form-core} alone is insufficient.

\begin{lemma}[Weighted cutoff commutator]\label{d:lem:Gamma-commutator}
If $\chi$ is real, bounded, even, and Lipschitz, multiplication by $\chi$ preserves ${X}$ and $\operatorname{Dom}T_\Gamma$. There is a bounded operator $C_\chi$ on $\Hh$ such that
\begin{align}
 \tg(\chi u,z)-\tg(u,\chi z)
 &=\ip{C_\chi u}{z}_{\Hh},\label{d:eq:Gamma-comm-form}\\
 T_\Gamma(\chi u)&=\chi T_\Gamma u+C_\chi u,
 \qquad
 \norm{C_\chi}\le2M\operatorname{Lip}(\chi)J_1.
 \label{d:eq:Gamma-comm-row}
\end{align}
The first identity holds for $u,z\in{X}$, and the second for $u\in\operatorname{Dom}T_\Gamma$.
\end{lemma}
\begin{proof}
The product rule gives
\[
 \tg[\chi u]
 \le2\norm{\chi}_\infty^2\tg[u]
       +2\operatorname{Lip}(\chi)^2J_2\norm{u}_2^2,
\]
proving form-domain preservation. Put $\beta=\sqrt W$ and use the isometry $U_Wu=u/\beta$ to ordinary even $L^2$. Expansion of the combined form difference in~\eqref{d:eq:Gamma-comm-form} gives the kernel
\begin{equation}\label{d:eq:Gamma-comm-kernel}
 (U_W C_\chi U_W^{-1}f)(x)
 =\int_\R\beta(x)\beta(y)\rho(|x-y|)
                 (\chi(x)-\chi(y))f(y)\,dy.
\end{equation}
Its absolute row and column integrals are at most
$2M\operatorname{Lip}(\chi)\int_0^\infty s\rho(s)\,ds$.
The elementary Schur estimate gives the asserted norm bound. For precision, first delete $|x-y|<\epsilon$; the omitted operator norm is at most
$2M\operatorname{Lip}(\chi)\int_0^\epsilon s\rho(s)\,ds\to0$.
The form difference converges to~\eqref{d:eq:Gamma-comm-form} by the difference-product identity and form Cauchy--Schwarz. The separately divergent diagonal rate is never split off. The sign in~\eqref{d:eq:Gamma-comm-kernel} is the sign of $[T_\Gamma,\chi]=T_\Gamma\chi-\chi T_\Gamma$.

For $u\in\operatorname{Dom}T_\Gamma$ and all $z\in{X}$,
\[
 \tg(\chi u,z)
 =\ip{T_\Gamma u}{\chi z}_{\Hh}+\ip{C_\chi u}{z}_{\Hh}
 =\ip{\chi T_\Gamma u+C_\chi u}{z}_{\Hh}.
\]
The representation criterion proves operator-domain preservation and~\eqref{d:eq:Gamma-comm-row}, rather than presupposing them.
\end{proof}

For the scaled cutoffs from the preceding section, every $u\in\operatorname{Dom}T_\Gamma$ therefore satisfies
\begin{align}
 \norm{u-\chi_Ru}_{\Hh}&\le\norm{\boldsymbol1_{|x|>R}u}_{\Hh},\notag\\
 \norm{T_\Gamma u-T_\Gamma(\chi_Ru)}_{\Hh}
 &\le\norm{\boldsymbol1_{|x|>R}T_\Gamma u}_{\Hh}
       +\frac{2M\norm{\chi'}_\infty J_1}{R}\norm{u}_{\Hh}.
 \label{d:eq:Gamma-cutoff-graph}
\end{align}
Both right sides tend to zero for the fixed graph vector. No uniform tail rate on a graph unit ball is claimed.

\begin{lemma}[Compact graph vectors and their noncompact rows]\label{d:lem:Gamma-compact-row}
If $u\in\operatorname{Dom}T_\Gamma$ is supported in $[-A,A]$, then $u\in\operatorname{Dom}D_\Gamma$ in ordinary $L^2$. For $h>0$ and $|x|>A+h$,
\begin{equation}\label{d:eq:Gamma-exterior}
 D_\Gamma u(x)=-\int_{-A}^A\rho(|x-y|)u(y)\,dy,
 \qquad
 |D_\Gamma u(x)|\le
 \frac{e^{A/2}\norm{u}_1e^{-|x|/2}}{1-e^{-2h}}.
\end{equation}
\end{lemma}
\begin{proof}
For $u\in{X}$, define $D_\Gamma u$ as a tempered distribution through the multiplier; the at-most-linear bound~\eqref{d:eq:Gamma-first-order} suffices. For $u\in\operatorname{Dom}T_\Gamma$, compact testing gives
\begin{equation}\label{d:eq:Gamma-local-row}
 D_\Gamma u=(T_\Gamma u)/W\quad\hbox{in }L^2_{\rm loc}.
\end{equation}
This follows by comparing the distributional form pairing with the representative pairing. Division by $W$ is legitimate locally only.

On a test support at positive distance from $[-A,A]$, the diagonal term vanishes. Expanding the form there and applying Fubini gives the exterior convolution in~\eqref{d:eq:Gamma-exterior}. The distance $h$ controls the denominator of $\rho$, while $|x-y|\ge|x|-A$ controls its numerator. If
$w_{A+h}=\min_{|x|\le A+h}W(x)>0$, the local row and exterior bound imply
\[
 \norm{D_\Gamma u}_2^2
 \le w_{A+h}^{-1}\norm{T_\Gamma u}_{\Hh}^2
       +\frac{2e^{-h}\norm{u}_1^2}{(1-e^{-2h})^2}<\infty.
\]
Thus the global distributional row is ordinary $L^2$, which is exactly the ordinary multiplier-domain criterion. Compact support of the input was essential; the output row is not asserted to be compactly supported.
\end{proof}

\begin{theorem}[Compact smooth graph core]\label{d:thm:Gamma-graph-core}
The even compact smooth functions form a graph core for $T_\Gamma$ and for every $B_{\mathsf D}$ in~\eqref{d:eq:Gamma-graph}.
\end{theorem}
\begin{proof}
First use~\eqref{d:eq:Gamma-cutoff-graph}. For a fixed compact graph vector $u$ supported in $[-A,A]$, let $u_\epsilon=\eta_\epsilon*u$, where $\eta$ is a nonnegative even compact smooth mollifier of mass one supported in $[-1,1]$. Lemma~\ref{d:lem:Gamma-compact-row} allows ordinary multiplier commutation:
\[
 D_\Gamma u_\epsilon=\eta_\epsilon*(D_\Gamma u).
\]
Both input and row converge in ordinary $L^2$. For $\epsilon\le1$ the input difference has support in $[-A-1,A+1]$, so
\[
 \norm{u_\epsilon-u}_{\Hh}
 \le w_{A+1}^{-1/2}\norm{u_\epsilon-u}_2,\qquad
 \norm{T_\Gamma(u_\epsilon-u)}_{\Hh}
 \le\sqrt M\norm{D_\Gamma(u_\epsilon-u)}_2.
\]
Each $u_\epsilon$ is already admitted by Lemma~\ref{d:lem:Gamma-H1}; the second estimate compares actual rows. A diagonal cutoff--mollifier sequence proves the graph core. Equivalent graph norms under bounded perturbation prove the assertion for $B_{\mathsf D}$. This theorem concerns the full ambient space $\Hh$, not a restriction to $\Kk$.
\end{proof}


\section{Actual realization and the spectral threshold}\label{sec:spectral}
The weighted realization and the ordinary-$L^2$ pencil are distinct operators with a common form numerator. The first is unitarily equivalent to the original energy operator; the second is a confining parameter family used to count signs.


\subsection{The weighted physical representative}\label{sec:weighted-physical}
Let $D_b$ be the bounded self-adjoint correction defined by
\begin{equation}\label{eq:Db-physical}
 \ip{D_bu}{v}_{\Hh}=\cg\ip uv_2+p(u,v),\qquad
 B=T_\Gamma-D_b,\quad \dom B=\dom T_\Gamma.
\end{equation}
Lemma~\ref{a:lem:prime} proves boundedness, including every prime power; in its notation $D_b=D$ and $\norm{D_b}<10$. The closed-form representation theorem and~\eqref{eq:intro-bridge} imply
\begin{equation}\label{eq:actual-unitary}
 UAU^{-1}=B+\tfrac12I,\qquad Uf=\kappa f.
\end{equation}
Indeed both sides represent the same closed form on $X$, so uniqueness of the self-adjoint representative proves the operator and domain identity. This does not identify the ordinary-$L^2$ representative of the same algebraic form.


\subsection{Derivative radicals and finite interpolation}\label{a:sec:radical}

The scaled response uses actual radicals of $q_0$, meaning vectors that annihilate the entire form domain under polarization. A vector with only $q_0[u]=0$ would not suffice. We also need ordinary-$L^2$ density, but not weighted density; the two statements differ for this kernel.

\begin{lemma}[Full-domain derivative radicals]\label{a:lem:radical}
For $j\ge1$, put
\begin{equation}\label{a:eq:radical}
 r_j=\kappa^{(2j)}-4^{-j}\kappa,\qquad
 \Rcal=\operatorname{span}\{r_j:j\ge1\}.
\end{equation}
Then $r_j\in X_c$ and
\begin{equation}\label{a:eq:radicalidentity}
 q_0(r_j,z)=0\qquad(z\in X).
\end{equation}
The family $(r_j)_{j\ge1}$ is linearly independent.
\end{lemma}
\begin{proof}
For each integer $m\ge0$, differentiating the theta series on the positive half-line gives
\[
 \kappa^{(m)}(x)=e^{x/2}\sum_{n\ge1}P_m(\pi n^2e^{2x})e^{-\pi n^2e^{2x}},
\]
where $P_0(t)=4t^2-6t$ and
$P_{m+1}(t)=(1/2-2t)P_m(t)+2tP_m'(t)$.
Thus $\deg P_m\le m+2$. Positivity and the first theta term also give
\begin{align}
 |\kappa^{(m)}(x)|&\le C_m
        e^{(9/2+2m)|x|-\pi e^{2|x|}},\label{a:eq:derivativeupper}\\
 \kappa(x)&\ge(4\pi^2-6\pi)
        e^{(9/2)|x|-\pi e^{2|x|}}.\label{a:eq:thetalower}
\end{align}
The upper bound follows by summing
$n^{2m+4}e^{-\pi(n^2-1)}$; the lower bound uses
$4\pi^2e^{4x}-6\pi e^{2x}\ge(4\pi^2-6\pi)e^{4x}$ for $x\ge0$.
Evenness extends both estimates to the full line. Therefore
\begin{equation}\label{a:eq:derivativeweighted}
 \frac{|\kappa^{(m)}(x)|^2}{W(x)}
 \le C_m'e^{(5+4m)|x|-\pi e^{2|x|}}.
\end{equation}
Every derivative has finite weighted norm and belongs to $H^1(dx)$, hence has finite Gamma form. In particular $d_j=\kappa^{(2j)}\in X$.

Integration by parts gives
\[
 \mathcal Fd_j(z)=(-z^2)^j\xi(1/2+iz),\qquad
 m(d_j)=4^{-j}m(\kappa)=4^{-j}/2.
\]
Use $d_j\chi_R$ in the bilinear explicit formula exactly as in the proof of Proposition~\ref{a:prop:thetanull}. For clarity, all three cutoff controls persist:~\eqref{a:eq:derivativeupper} gives convergence in $H^1$ and in $L^1(e^{|x|/2}dx)$; on each compact test support the complete shifted prime row is dominated by a summable $n^{-A_0}$ majorant with $A_0>2$; and the zero side is dominated by
\[
 \left(\int_\R|d_j(x)|e^{|x|/2}\,dx\right)
 C_{h,N}(1+|\operatorname{Re}w|)^{-N}.
\]
The zero-counting bound makes this summable. Every limiting product vanishes because $\mathcal Fd_j(w_{\rho_\zeta})=0$. Hence, on compact even tests,
\[
 q_0(d_j,h)=-4^{-j}\overline{m(h)}.
\]
Continuity of $q_0$ on $X\times X$ and the compact smooth form core extend the identity to all $h\in X$. Subtracting $4^{-j}$ times the source identity for $\kappa$ proves~\eqref{a:eq:radicalidentity}; the moment formula proves $r_j\in X_c$.

If $\sum_{j=1}^Jb_jr_j=0$, Fourier transformation gives
\[
 \left(\sum_{j=1}^Jb_j\bigl((-z^2)^j-4^{-j}\bigr)\right)
                 \xi(1/2+iz)=0\qquad(z\in\R).
\]
The second factor is nonzero on a neighborhood of zero, since
$\mathcal F\kappa(0)=\int\kappa>0$. The polynomial must vanish identically. Comparing its highest degrees successively gives all $b_j=0$.
\end{proof}

\begin{lemma}[Holomorphic translations]\label{a:lem:strip}
The real kernel has an even holomorphic continuation to
$\mathcal S=\{z:|\operatorname{Im}z|<\pi/4\}$. For every $m\ge0$ and every $0\le\vartheta<\pi/4$,
\begin{equation}\label{a:eq:stripbound}
 |\kappa^{(m)}(r+iy)|\le C_{m,\vartheta}
 e^{(9/2+2m)|r|}
 e^{-\pi\cos(2\vartheta)e^{2|r|}}
 \quad(|y|\le\vartheta).
\end{equation}
The map $b\mapsto\kappa(\cdot-b)$ is holomorphic from $\mathcal S$ into ordinary $L^2(dx)$.
\end{lemma}
\begin{proof}
The continuation and derivative estimate are Lemma~\ref{d:lem:theta-tails}. It remains to prove ordinary-$L^2$-valued holomorphy. Put $d_\vartheta=\pi\cos(2\vartheta)>0$.

On a compact set of translation parameters choose $B<\infty$ and $\vartheta<\pi/4$ with $|\operatorname{Re}b|\le B$ and $|\operatorname{Im}b|\le\vartheta$. Formula~\eqref{a:eq:stripbound} supplies the common square-integrable envelope
\[
 C_{m,\vartheta}e^{(9/2+2m)(|x|+B)}
       e^{-d_\vartheta e^{2\max\{|x|-B,0\}}}.
\]
Taylor's integral remainder, with the bound for the second derivative, proves complex differentiability of $b\mapsto\kappa(\cdot-b)$ in $L^2$. Repetition gives derivative $(-1)^m\kappa^{(m)}(\cdot-b)$ of every order.
\end{proof}

\begin{proposition}[Ordinary density and weighted incompleteness]\label{a:prop:density}
The closure of $\Rcal$ in ordinary $\Le$ is all of $\Le$. Its closure in $\Hh_c$ with the weighted norm is a proper subspace.
\end{proposition}
\begin{proof}
Suppose $h\in\Le$ is orthogonal to every $r_j$, and define the bilinear convolution pairing
\[
 F(b)=\int_\R h(x)\kappa(x-b)\,dx,\qquad b\in\mathcal S.
\]
Lemma~\ref{a:lem:strip} makes $F$ holomorphic. Since $h$ and $\kappa$ are even, $F$ is even. Orthogonality to the real functions $r_j$ gives
$F^{(2j)}(0)=4^{-j}F(0)$ for all $j\ge1$.
All Taylor coefficients therefore agree with those of $F(0)\cosh(b/2)$, and the identity theorem yields
\[
 F(b)=F(0)\cosh(b/2)\qquad(b\in\mathcal S).
\]
For real $b$, Cauchy--Schwarz bounds $|F(b)|$ by
$\norm h_2\norm\kappa_2$, independently of $b$. Thus $F(0)=0$, so $F=0$. Since $\kappa\in L^1$, the convolution belongs to $L^2$ and Plancherel gives
$\mathcal F\kappa(t)\widehat h(t)=0$ almost everywhere. The first factor is a nonzero entire function, whose real zeros are discrete and hence have measure zero. It follows that $h=0$, proving ordinary density.

For the different weighted assertion, choose any nontrivial zero of $\zeta$, and write it as $1/2+iw$. Such zeros exist by the classical counting theorem \cite{vonMangoldt}; their position on or off the critical line is irrelevant. Put
\[
 u_w(x)=W(x)\cos(wx).
\]
The theta envelope gives $u_w\in\Hw$. The entire transform identity and integration by parts imply
\[
 m(u_w)=\frac12\mathcal F\kappa(w)=0,\qquad
 \ip{u_w}{r_j}_{\Hh}
 =\int_\R\cos(wx)r_j(x)\,dx
 =\bigl((-w^2)^j-4^{-j}\bigr)\mathcal F\kappa(w)=0.
\]
Thus $u_w\in \Hh_c$ is orthogonal to $\Rcal$, and it is nonzero because $u_w(0)=W(0)>0$. This proves weighted incompleteness. It does not identify $u_w$ as an eigenstate or as a counterexample to any variance inequality.
\end{proof}

The ordinary density in Proposition~\ref{a:prop:density} will be applied only to finitely many ordinary-$L^2$ continuous moment functionals. It gives no bound on the weighted norm of an interpolant. In particular it cannot be substituted for density in the original centered profile space under $u=\kappa f$.


\subsection{A source map and its actual invariant domain}\label{b:sec:source}
In this subsection let $a(x)=(2c(x))^{-1}$, $L=-A$, and
\begin{equation}\label{b:eq:nu}
 d\nu(s)=\rho(s)\,ds+\sum_{n\ge2}a_n\delta_{s_n}(ds),\qquad s>0.
\end{equation}
The function $a(x)$ is distinct from the coefficients $a_n$ and from the complex scalar $a=i/2$ used later only in the synthesis sections.


\begin{equation}\label{b:eq:source-main}
 Q\kappa=-c,\qquad Q=D_\Gamma-\cg-P,
\end{equation}
where, on smooth rapid inputs,
\begin{align}
 D_\Gamma h(x)&=\int_0^\infty\rho(s)
       [2h(x)-h(x+s)-h(x-s)]\,ds,\label{b:eq:Drow}\\
 Ph(x)&=\sum_{n\ge2}a_n[h(x+s_n)+h(x-s_n)].\label{b:eq:Prow}
\end{align}
The compensated difference in~\eqref{b:eq:Drow} is never split into divergent diagonal rates.

Let $\RR$ be the space of smooth rapid-exponential functions with seminorms
\[
 p_{B,m}(h)=\sup_x e^{B|x|}|h^{(m)}(x)|<\infty
 \quad(B>0,\ m\ge0).
\]
Let $\SSS$ be the smooth small-exponential class with seminorms
\[
 q_{\epsilon,m}(v)=\sup_xe^{-\epsilon|x|}|v^{(m)}(x)|<\infty
 \quad(\epsilon>0,\ m\ge0).
\]
Finally, $\CE$ consists of all smooth functions satisfying, for each fixed $m$, a bound $|f^{(m)}(x)|\le C_m e^{B_m|x|}$ for some finite $C_m,B_m$. The exponent may depend on $m$. This is an algebra, $\SSS\subset\CE$, and the theta derivative estimates~\eqref{a:eq:derivativeupper} imply $\kappa\CE\subset\RR$.

\begin{lemma}[Continuous full source map]\label{b:lem:T}
The map $T=aQ$ sends $\RR$ continuously into $\SSS$. The maps $h\mapsto WQh$, $h\mapsto W(D_\Gamma-\cg)h$, and $h\mapsto WPh$ send $\RR$ continuously into $\RR$.
\end{lemma}
\begin{proof}
For $0<s\le1$ the $m$th differentiated Gamma difference is bounded by $s^2\norm{h^{(m+2)}}_\infty$, and for $s>1$ by $4\norm{h^{(m)}}_\infty$. These bounds are integrable against $\rho(s)\,ds$ and justify every differentiation. Thus every derivative of $D_\Gamma h$ is bounded.

For $B>1/2$, the elementary inequalities $\Lambda(n)\le\log n$ and $|x\pm\log n|\ge\log n-|x|$ give
\begin{equation}\label{b:eq:primepoint}
 |(Ph)^{(m)}(x)|\le
  2p_{B,m}(h)e^{B|x|}\sum_{n\ge2}(\log n)n^{-B-1/2}.
\end{equation}
The scalar sum is finite by the integral test, and includes every prime power. It proves locally uniform absolute convergence of the series and every differentiated series. Since $|a^{(j)}(x)|\le C_j e^{-|x|/2}$, choosing $B=1/2+\epsilon$ and applying Leibniz's rule proves the $\SSS$ bounds. The same argument gives the full tail estimate
\begin{equation}\label{b:eq:primetailpoint}
 q_{\epsilon,m}\bigl(a(P-P_N)h\bigr)
 \le C_{m,\epsilon}\max_{0\le j\le m}p_{1/2+\epsilon,j}(h)
           \sum_{n>N}(\log n)n^{-1-\epsilon}\longrightarrow0.
\end{equation}
Here $P_N$ retains the terms with $n\le N$ in~\eqref{b:eq:Prow}. This is a weighted-growth seminorm estimate, not a uniform unweighted incoming-rate estimate. Every derivative of $Qh$ grows at most as $e^{(1/2+\epsilon)|x|}$. Multiplication by the rapid theta factor $W$ gives all target $\RR$ seminorms. The same proof applies separately to its Gamma and prime parts.
\end{proof}

\begin{proposition}[Actual generator and invariant algebra]\label{b:prop:generator}
Every even $f\in\CE$ belongs to $\dom(A)$, and
\begin{equation}\label{b:eq:generator-source}
 Af=\tfrac12f+T(\kappa f),\qquad
 Lf=-\tfrac12f-T(\kappa f).
\end{equation}
The right sides lie in $\CE$, so every such profile belongs to $\dom(A^m)$ for each finite $m$. Pointwise,
\begin{equation}\label{b:eq:generator-jump}
 Lf(x)=a(x)\int_{s>0}\!
 \left\{\kappa(x+s)[f(x+s)-f(x)]
       +\kappa(x-s)[f(x-s)-f(x)]\right\}\,d\nu(s).
\end{equation}
\end{proposition}
\begin{proof}
Put $u=\kappa f\in\RR$. The theta bounds imply $u\in H^1(\R)$ and
$\int |u|^2/W=\int |f|^2\,d\mu<\infty$. Thus $u\in X$, and the full bridge admits $f$ to $\D$. Against an even compact smooth $h$, the compensated Gamma row can be paired with $\kappa h$ by the Taylor estimates above; the prime row is locally absolutely convergent by~\eqref{b:eq:primepoint}. Hence
\[
 q_0(u,\kappa h)=\int Q(u)\kappa\overline h\,dx
                =\ip{T(u)}h_\mu.
\]
Lemma~\ref{b:lem:T} gives $T(u)\in\SSS\subset\He$, because $\mu$ has all exponential moments. The bridge yields
$\E(f,h)=\ip{f/2+T(u)}h_\mu$ on the compact smooth form core. Both sides are continuous in the form norm of $h$, so the equality extends to all $h\in\D$. The representation criterion~\eqref{a:eq:Adef} proves~\eqref{b:eq:generator-source}. Its right side is in $\CE$, proving invariance and iteration.

The weak equation~\eqref{b:eq:source-main} upgrades to a pointwise equality because Lemma~\ref{b:lem:T} makes $Q\kappa$ smooth and the residual is even. In the jump row, near zero the profile difference is $O(s)$ and $\rho(s)=O(1/s)$; its Gamma tail is integrable. The prime row converges locally absolutely by~\eqref{b:eq:primepoint}, applied to $\kappa f$ and $\kappa$. Subtracting the two compensated source rows gives
\[
 Q(\kappa f)-fQ\kappa
  =-\int_{s>0}\!\left\{\kappa(x+s)[f(x+s)-f(x)]
             +\kappa(x-s)[f(x-s)-f(x)]\right\}\,d\nu(s).
\]
Now use $Q\kappa=-c$ and $ac=1/2$. This proves~\eqref{b:eq:generator-jump} with the coefficient $1/(2c)$.
\end{proof}


For every $\delta>0$, set
\begin{equation}\label{eq:ell-unified}
 \ell_\delta[u]=q_0[u]+\delta\norm u_{\Hh}^2,\qquad \dom\ell_\delta=X.
\end{equation}


\subsection{The closed pencil and compact resolvent}\label{a:sec:pencil}

\begin{proposition}[Closed ordinary-$L^2$ realization]\label{a:prop:closedpencil}
For every $\delta>0$, the form $\ell_\delta$ of~\eqref{eq:ell-unified} is closed and lower bounded on $\Le$. Its self-adjoint representative is characterized by
\begin{equation}\label{a:eq:Ldef}
 \begin{split}
 u\in\dom(L_\delta),\ L_\delta u=h\quad\Longleftrightarrow\quad
 u\in X,\ h\in\Le,\quad
 \ell_\delta(u,z)=\ip{h}{z}_2\ \text{for all }z\in X.
 \end{split}
\end{equation}
The construction here uses the direct numerator estimates.
\end{proposition}
\begin{proof}
Choose an integer $N$ for which $r_N\le\delta/2$ and set
\[
 C_N=|\cg|+2\sum_{2\le n\le N}a_n.
\]
The complete head--tail estimate gives
\begin{equation}\label{a:eq:coercive}
 \ell_\delta[u]\ge\tg[u]+\frac\delta2\norm{u}_{\Hh}^2-C_N\norm{u}_2^2.
\end{equation}
Its upper absolute bound is a finite multiple of $\norm{u}_X^2$, by Lemma~\ref{a:lem:prime} and~\eqref{a:eq:L2bound}. Therefore
$\ell_\delta[u]+(C_N+1)\norm{u}_2^2$ defines a positive norm squared equivalent to $\norm{u}_X^2$. Completeness and density follow from Lemma~\ref{a:lem:Xcore}. The closed-form representation theorem gives the self-adjoint operator and~\eqref{a:eq:Ldef}.
\end{proof}

The choice of $N$ in this proof is existential and analytic. The form remains the complete form for every such choice. In distributional notation the operator is
\begin{equation}\label{a:eq:formalL}
 L_\delta=D_\Gamma-\cg-P+\frac\delta W,
 \qquad
 Pu=\sum_{n\ge2}a_n\bigl(u(\cdot+s_n)+u(\cdot-s_n)\bigr),
\end{equation}
where $D_\Gamma$ is represented by $\tg$. For $u\in X$, the prime row converges in $L^2$ on every fixed compact interval: each shifted local norm is bounded by the local supremum of $\sqrt W$ times $\norm{u}_{\Hh}$, and the theta envelope makes the resulting tails summable. Also $u/W$ is locally $L^2$. The Gamma term is defined distributionally by its form. Operator-domain membership requires that the \emph{sum} in~\eqref{a:eq:formalL} be represented by an ordinary $L^2$ vector; it does not assert that each term is globally $L^2$. Compact smooth functions are a form core. No operator-core claim for this ordinary-$L^2$ realization is being made.

\begin{lemma}[Logarithmic Fourier coercivity]\label{a:lem:fourier}
For the unitary Fourier transform,
\begin{equation}\label{a:eq:multiplier}
 \tg[u]=\int_\R d_\Gamma(\xi)|\widehat u(\xi)|^2\,d\xi,
 \quad
 d_\Gamma(\xi)=2\int_0^\infty\rho(s)(1-\cos(\xi s))\,ds
 \ge\frac12\log(1+4\xi^2).
\end{equation}
\end{lemma}
\begin{proof}
Plancherel's theorem and Tonelli's theorem give the multiplier identity, allowing infinite values. The inequality $1-e^{-2s}\le2s$ gives
\[
 d_\Gamma(\xi)\ge\int_0^\infty e^{-s/2}\frac{1-\cos(\xi s)}s\,ds.
\]
For general $a>0$, differentiating this last integral with $e^{-as}$ in place of $e^{-s/2}$ yields $\int_0^\infty e^{-as}\sin(\xi s)\,ds=\xi/(a^2+\xi^2)$. Integrating back from $\xi=0$ gives $\tfrac12\log(1+\xi^2/a^2)$. Taking $a=1/2$ proves the bound.
\end{proof}

\begin{proposition}[Compactness]\label{a:prop:compact}
The embedding $X\hookrightarrow\Le$ is compact. Consequently $L_\delta$ has compact resolvent for every $\delta>0$.
\end{proposition}
\begin{proof}
The unit ball of $X$ is uniformly spatially tight, since
\begin{equation}\label{a:eq:spatialtight}
 \int_{|x|>R}|u(x)|^2\,dx
 \le\sup_{|x|>R}W(x)\,\norm{u}_{\Hh}^2\longrightarrow0
\end{equation}
uniformly on that ball. Lemma~\ref{a:lem:fourier} gives uniform Fourier tightness:
\begin{equation}\label{a:eq:fouriertight}
 \int_{|\xi|>M}|\widehat u(\xi)|^2\,d\xi
 \le\frac{\tg[u]}{\tfrac12\log(1+4M^2)}\longrightarrow0.
\end{equation}
Let $P_M$ be the Fourier projection onto $[-M,M]$. The operator
$\one_{[-R,R]}P_M$ is Hilbert--Schmidt: its space--frequency kernel has constant modulus $(2\pi)^{-1/2}$ on a finite rectangle and vanishes elsewhere. Moreover
\[
 \norm{u-\one_{[-R,R]}P_Mu}_2
 \le\norm{\one_{\{|x|>R\}}u}_2+\norm{(I-P_M)u}_2.
\]
Thus the inclusion is a uniform limit on the $X$ unit ball of compact operators. The cutoffs and Fourier projection preserve evenness.

By~\eqref{a:eq:coercive}, a resolvent $(L_\delta+aI)^{-1}$ with $a>C_N$ maps ordinary $L^2$ boundedly into $X$. Composing this map with the compact inclusion proves compactness of the resolvent. Hence the eigenvalues are real, have finite multiplicity, are bounded below, and tend to $+\infty$.
\end{proof}

This compactness concerns $L_\delta$ on ordinary $L^2$, not the original operator $A_E$ on $L^2(\mu)$. The two Hilbert norms are different.


\subsection{Actual eigenstates and the negative index}\label{f:sec:spectrum}

Proposition~\ref{a:prop:domain} gives, for $\lambda=1/2-\delta$,
\begin{equation}\label{f:eq:congruence}
 \ell_\delta(\kappa f,\kappa g)
       =\E(f,g)-\lambda\ip{f}{g}_\mu,
       \qquad f,g\in\D.
\end{equation}
This is a form congruence through a bijection of form domains, not a unitary equivalence of operators on their displayed Hilbert spaces.

\begin{proposition}[Fixed-distance spectral finiteness]\label{f:prop:finite}
For every $t<1/2$, the spectral projection $\one_{(-\infty,t]}(A)$ has finite rank. The centered restriction satisfies $A_E\ge\varepsilon_0 I$ for some $\varepsilon_0>0$ and
\begin{equation}\label{f:eq:essential}
 \inf\sigma_{\ess}(A_E)\ge\frac12,
\end{equation}
where the infimum of the empty essential spectrum is $+\infty$.
\end{proposition}
\begin{proof}
For $t<0$ the assertion follows from $A\ge0$. Let $0\le t<1/2$, and take $\delta=1/2-t$. Every $f$ in the range of $\one_{[0,t]}(A)$ lies in $\dom(A)$ and satisfies $\E[f]\le t\norm{f}_\mu^2$. Equation~\eqref{f:eq:congruence} sends this range injectively into a subspace of $X$ on which $\ell_\delta\le0$.

A lower-bounded compact-resolvent operator has only finitely many nonpositive eigenvalues. Its nonpositive spectral subspace is therefore finite dimensional. Any subspace of its form domain on which the form is nonpositive has dimension at most that finite number: otherwise a nonzero vector in the subspace would be orthogonal to the entire nonpositive spectral subspace, where the form is strictly positive. Applying this observation to $L_\delta$ proves finite rank.

The centered restriction has no zero eigenvector, by Proposition~\ref{a:prop:core}. Apply the finite-rank statement at any one $t_0\in(0,1/2)$, for example $t_0=1/4$. On its centered finite-dimensional spectral range, if nonzero, the smallest eigenvalue is strictly positive; on the complementary range the operator is bounded below by $t_0$. This proves a positive lower bound. Finally, finite rank below every $t<1/2$ excludes essential spectrum below $1/2$ and permits accumulation of subthreshold eigenvalues only at $1/2$.
\end{proof}

The proposition is a consequence of the compact pencil and the bridge. It does not require an additional estimate for primes in short intervals. The next result establishes equality in~\eqref{f:eq:essential}.


\begin{proposition}[Actual half eigenstates and essential bottom]\label{f:prop:halfeigen}
The profiles
\begin{equation}\label{f:eq:halfprofiles}
 g_j=\kappa^{(2j)}/\kappa-4^{-j}\quad(j\ge1)
\end{equation}
belong to $\dom A_E$, satisfy $A_Eg_j=g_j/2$, and are linearly independent. Hence
\begin{equation}\label{f:eq:exactessential}
 \dim\ker(A_E-\tfrac12I)=\infty,\qquad
 \inf\sigma_{\mathrm{ess}}(A_E)=\tfrac12.
\end{equation}
\end{proposition}
\begin{proof}
Lemma~\ref{a:lem:radical} proves $r_j\in X_c$ and $q_0(r_j,z)=0$ for every $z\in X$, with linear independence. The representation criterion for $B$ gives $r_j\in\dom B$, $Br_j=0$. Apply~\eqref{eq:actual-unitary} and $U^{-1}r_j=g_j$. The resulting infinite-dimensional eigenspace puts $1/2$ in the essential spectrum. Proposition~\ref{f:prop:finite} excludes essential spectrum below it.
\end{proof}
The equality states do not exclude lower discrete spectrum. No completeness of the derivative family in the centered weighted space has been asserted; its complement is the subject of Section~\ref{sec:synthesis-unified}.


The compact-resolvent eigenvalues are enumerated increasingly, with multiplicity, as $\eta_1(\delta)\le\eta_2(\delta)\le\cdots$. For $0<\delta<1/2$ and $\lambda=1/2-\delta$, the kernel correspondence to be proved is
\begin{equation}\label{eq:pencil-kernel}
 \ker L_\delta=\{\kappa f:f\in\dom A_E,\ A_Ef=\lambda f\}.
\end{equation}


\begin{proposition}[Kernel correspondence]\label{f:prop:kernel}
For $0<\delta<1/2$ and $\lambda=1/2-\delta$, equation~\eqref{eq:pencil-kernel} holds as an exact linear correspondence, including multiplicities.
\end{proposition}
\begin{proof}
If $u\in\ker L_\delta$, then $u\in X$ and $f=u/\kappa\in\D$. Testing~\eqref{a:eq:Ldef} against every $z=\kappa g\in X$ and using~\eqref{f:eq:congruence} gives
\[
 \E(f,g)=\lambda\ip{f}{g}_\mu\qquad(g\in\D).
\]
The representation criterion~\eqref{a:eq:Adef} proves $f\in\dom(A)$ and $Af=\lambda f$. Since $\lambda>0$ and $A1=0$, self-adjointness implies $\mu(f)=0$. Thus $f\in\dom(A_E)$ is a genuine centered eigenstate.

Conversely, for such an eigenstate $f$, the vector $u=\kappa f$ belongs to $X$ and~\eqref{f:eq:congruence} vanishes against every $z\in X$. Its ordinary-$L^2$ representing vector is zero. Hence $u\in\dom(L_\delta)$ and $L_\delta u=0$.
\end{proof}

No pointwise eigen-row or additional regularity of an eigenstate is used here. The correspondence identifies kernels at the indicated parameter; it does not identify nonzero eigenvalues of $L_\delta$ with eigenvalues of $A_E$.

\begin{proposition}[Exact form index]\label{f:prop:index}
For $0<\delta<1/2$ and $\lambda=1/2-\delta$,
\begin{align}
 n_-(L_\delta)&=1+\rank\one_{(0,\lambda)}(A_E),\label{f:eq:index}\\
 \dim\ker L_\delta&=\dim\ker(A_E-\lambda).\label{f:eq:nullity}
\end{align}
For $u\in X$, put $a=2m(u)$ and $u_0=u-a\kappa$. Then
\begin{align}
 m(u_0)&=0,\qquad
 \norm{u}_{\Hh}^2=|a|^2+\norm{u_0}_{\Hh}^2,\label{f:eq:decompnorm}\\
 \ell_\delta[u]&=-\lambda|a|^2+
              \E[u_0/\kappa]-\lambda\norm{u_0}_{\Hh}^2.\label{f:eq:decompform}
\end{align}
\end{proposition}
\begin{proof}
Write $f=u/\kappa=a+g$, where $a=\mu(f)$ and $\mu(g)=0$. Constants have zero energy, and polarization gives $\E(1,g)=0$. Also $\norm{1}_\mu=1$. The orthogonal decomposition in $L^2(\mu)$ and~\eqref{f:eq:congruence} prove~\eqref{f:eq:decompnorm}--\eqref{f:eq:decompform}.

A bijective linear change of form domain preserves the maximal dimension of a strictly negative form subspace. On the constant line the right side of~\eqref{f:eq:decompform} is strictly negative. On the centered subspace its negative index is the rank of the spectral projection of $A_E$ below $\lambda$, by the spectral theorem. Proposition~\ref{f:prop:finite} and the positive gap make this the finite rank in~\eqref{f:eq:index}. For a compact-resolvent operator its form negative index equals its number of negative eigenvalues. The nullity identity is Proposition~\ref{f:prop:kernel}.
\end{proof}

\begin{corollary}[Failure of the half bound]\label{f:cor:failure}
The global inequality $\E[f]\ge\tfrac12\Var_\mu(f)$ on $\D$ fails if and only if there exists an actual eigenpair
\[
 A_Ef=\lambda f,\qquad f\in\dom(A_E)\setminus\{0\},\qquad0<\lambda<1/2.
\]
\end{corollary}
\begin{proof}
Subtract the mean of a violating profile to obtain $g\ne0$ with $\E[g]/\norm{g}_\mu^2<1/2$. Choose $t$ strictly between this quotient and $1/2$. The spectral projection $\one_{(-\infty,t]}(A_E)$ is nonzero, since otherwise the spectral theorem would give $\E[g]\ge t\norm{g}_\mu^2$. It is finite rank by Proposition~\ref{f:prop:finite}, so its invariant range contains an eigenvector. Its eigenvalue is positive by the gap and at most $t<1/2$. The converse follows by evaluating the energy on the eigenvector.
\end{proof}

\subsection{The mean and the operator-domain obstruction}\label{f:sec:mean}

The decomposition in Proposition~\ref{f:prop:index} is orthogonal in the weighted norm, but not in ordinary physical $L^2$. Two concrete facts make the distinction unavoidable.

\begin{proposition}[The mean is unbounded in ordinary $L^2$]\label{f:prop:mean}
The functional $m$ is continuous on $X$ and on the weighted space, with bound~\eqref{a:eq:meanbound}, but is not bounded in the ordinary $L^2$ norm, even on even $C_c^\infty(\R)$. Its zero-mean subspace in $X$ is closed in $X$ and dense in ordinary $\Le$.
\end{proposition}
\begin{proof}
Only unboundedness and density remain to be shown. Fix a nonzero, nonnegative, even $\phi\in C_c^\infty(\R)$ and put $M=m(\phi)>0$. For sufficiently large $R$, the two translates in
\[
 \psi_R(x)=\phi(x-R)+\phi(x+R),\qquad
 b_R(x)=\frac{\psi_R(x)}{2\cosh(R/2)M}
\]
have disjoint supports. The addition formula for the hyperbolic cosine and evenness of $\phi$ give $m(\psi_R)=2\cosh(R/2)M$. Hence
\begin{equation}\label{f:eq:unboundedmean}
 m(b_R)=1,\qquad
 \norm{b_R}_2^2=\frac{\norm{\phi}_2^2}{2\cosh^2(R/2)M^2}\longrightarrow0.
\end{equation}
Every $b_R$ is in the compact smooth numerator core. Thus neither $m$ nor the quadratic pole $2|m|^2$ has an ordinary-$L^2$ bound. For any even compact smooth $u$, the corrected function $u-m(u)b_R$ has zero mean and converges to $u$ in ordinary $L^2$. Density of the compact smooth core in $\Le$ proves the last claim. No smallness of $\norm{b_R}_X$ or $\norm{b_R}_{\Hh}$ is asserted.
\end{proof}

\begin{proposition}[The theta kernel is generally not in the operator domain]\label{f:prop:kdomain}
For every $\delta>0$, $\kappa\in X$. If $\delta\ne1/2$, then $\kappa\notin\dom(L_\delta)$. At $\delta=1/2$,
\begin{equation}\label{f:eq:halfkernel}
 L_{1/2}\ge0,\qquad \ker L_{1/2}=\C\kappa.
\end{equation}
For $0<\delta<1/2$, $\kappa$ is nevertheless a strictly negative form direction.
\end{proposition}
\begin{proof}
Since $1\in\D$, Proposition~\ref{a:prop:domain} admits $\kappa$ to $X$ and gives $\norm{\kappa}_{\Hh}=1$, $m(\kappa)=1/2$. Put $\lambda=1/2-\delta$, now allowing either sign. Equation~\eqref{f:eq:congruence} gives the exact polarized identity
\begin{equation}\label{f:eq:sourcerow}
 \ell_\delta(\kappa,z)=-2\lambda\overline{m(z)},\qquad z\in X.
\end{equation}
If $\kappa\in\dom(L_\delta)$, then compact smooth even tests in~\eqref{a:eq:Ldef} force its ordinary-$L^2$ representative to equal $-2\lambda c$ as an even distribution. For $\lambda\ne0$ this is impossible, since $c\notin L^2(\R)$. When $\lambda=0$, the form is $\E[u/\kappa]\ge0$ and its zero vectors are precisely multiples of $\kappa$, by Proposition~\ref{a:prop:core}. Finally $\ell_\delta[\kappa]=-\lambda<0$ for $0<\delta<1/2$.
\end{proof}

Thus a negative \emph{form direction} is not necessarily an operator-domain vector, much less an ordinary-$L^2$ ground eigenvector. The mean-zero space is not an ordinary-$L^2$ closed codimension-one hyperplane. Moreover $\int u\overline\kappa\,dx=0$ is different from $m(u)=0$: the two functionals on compact tests would be proportional only if $\kappa$ were a scalar multiple of $c$, contrary to their decay. Removing $\kappa$ by an ordinary-$L^2$ orthogonal projection therefore does not implement centering of the original problem.

\subsection{The second eigenvalue and the remaining sign}\label{f:sec:sign}

The compact-resolvent construction gives a precise family of sign questions, but it supplies no positive sign for the second eigenvalue throughout the required parameter interval.

\begin{proposition}[Equivalent sign conditions]\label{f:prop:sign}
The following statements are equivalent:
\begin{enumerate}
\item $\E[f]\ge\tfrac12\Var_\mu(f)$ for every $f\in\D$.
\item For every $0<\delta<1/2$, $\eta_2(\delta)>0$.
\item For every $0<\delta<1/2$ and every nonzero $u\in X$ with $m(u)=0$, one has $\ell_\delta[u]>0$.
\item $\ker L_\delta=\{0\}$ for every $0<\delta<1/2$.
\end{enumerate}
\end{proposition}
\begin{proof}
If the first statement holds, then on centered vectors~\eqref{f:eq:congruence} gives $\ell_\delta[u]\ge\delta\norm{u}_{\Hh}^2>0$. It also gives no centered spectrum of $A_E$ at or below $\lambda=1/2-\delta$. Equations~\eqref{f:eq:index}--\eqref{f:eq:nullity} then show that $L_\delta$ has exactly one negative eigenvalue and no zero eigenvalue; compact resolvent gives $\eta_2(\delta)>0$.

Conversely, if $\eta_2(\delta)>0$, the index and nullity identities exclude centered spectrum at or below $\lambda$, proving strict positivity on the centered form subspace. If the third statement holds for all $\delta$, letting $\delta\downarrow0$ for each fixed centered $u$ gives $q_0[u]\ge0$. By~\eqref{a:eq:polebridge} and subtraction of the original mean, this is the first statement. Finally, the equivalence of the first and fourth statements follows from Corollary~\ref{f:cor:failure} and Proposition~\ref{f:prop:kernel}.
\end{proof}

Two valid estimates illustrate the limits of the construction. First, min--max on the common domain $X$ and~\eqref{a:eq:L2bound} give
\begin{equation}\label{f:eq:monotonicity}
 \eta_j(\delta_2)\ge\eta_j(\delta_1)
       +\frac{\delta_2-\delta_1}{\norm{W}_\infty},
       \qquad\delta_2>\delta_1>0.
\end{equation}
At $\delta=1/2$, Proposition~\ref{f:prop:kdomain} and compact resolvent give $\eta_2(1/2)>0$. However, decreasing $\delta$ decreases the form. Inequality~\eqref{f:eq:monotonicity} does not prohibit a second eigenvalue from crossing zero at a smaller positive parameter.

Second, suppose $u\in\ker L_\delta$ with $0<\delta<1/2$ and normalize $\norm{u}_{\Hh}=1$. Then
\begin{equation}\label{f:eq:stateidentity}
 \tg[u]+\delta=\cg\norm{u}_2^2+p[u].
\end{equation}
Choose $N$ with $r_N\le\delta/2$ as in Proposition~\ref{a:prop:closedpencil}. The same complete-tail estimate yields
\begin{equation}\label{f:eq:statebound}
 \tg[u]+\frac\delta2\le C_N\norm{u}_2^2,
 \qquad \norm{u}_2^2\ge\frac{\delta}{2C_N}.
\end{equation}
Thus a normalized fixed-gap candidate cannot lose all its ordinary physical $L^2$ mass. The estimate is not contradictory: reducing the tail by enlarging $N$ also enlarges the head constant $C_N$. No estimate reversing this balance is proved. Testing the kernel equation against $\kappa$ gives only $m(u)=0$, by~\eqref{f:eq:sourcerow}; this recovers centering and provides no second independent sign constraint.


\section{Global variance comparisons and exact mean repair}\label{sec:bounds}
There are three comparison levels. The first is wholly analytic. The second uses only a whole-local-domain sign. The third has additional response hypotheses. None of the proofs replaces the complete energy by its finite head.


\begin{theorem}[Analytic variance and inertia bounds]\label{a:thm:mainseed}
For every $f\in\D$,
\begin{equation}\label{a:eq:mainseed}
 \E[f]\ge c_{\log 2}\Var_\mu(f),\qquad
 c_{\log 2}=\frac4{85}\left(\frac{63}{128}-7238e^{-12}\right)
 >\frac{8939}{425000}>\frac1{48}.
\end{equation}
For each $\delta>0$, the form $\ell_\delta$ on $X$ has a lower-bounded self-adjoint representative $L_\delta$ in ordinary $\Le$, with compact resolvent. If
\[
 \frac12-c_{\log 2}<\delta<\frac12,
\]
then $L_\delta$ has exactly one negative eigenvalue, counted with multiplicity, and zero operator nullity. At $\delta_3=23/48$, its exact mean Schur scalar is $-1/48$, and on $X_c=\{u\in X:\int_\R cu=0\}$,
\begin{equation}\label{a:eq:mainmargin}
 \ell_{\delta_3}[u]\ge\frac{509}{2550000}\norm{u}_{\Hh}^2.
\end{equation}
The strongest open endpoint furnished by this particular bound is $1/2-c_{\log 2}$, not its rational rounding $1/2-8939/425000$.
\end{theorem}


\subsection{A radius-sensitive analytic variance bound}\label{a:sec:seed}

We now prove the first quantitative conclusion. The estimate uses a fixed small prime head, but the weighted remainder from Lemma~\ref{a:lem:prime} remains present throughout. Combining the full-domain bridge with~\eqref{a:eq:head4} gives, for every $f\in\D$,
\begin{equation}\label{a:eq:masscontrol}
 \E[f]\ge\tg[\kappa f]+\frac{63}{128}\norm f_\mu^2
                         -9\norm{\kappa f}_2^2.
\end{equation}
Indeed the coefficient of $\norm f_\mu^2$ before rounding is $1/2-r_4>63/128$. The head through $4$ includes the prime square $4$; every later prime correlation is charged to $r_4$.

\subsubsection{Even folding and the global mean}

For $r>0$, write $I_r=[-r,r]$ and set
\begin{equation}\label{a:eq:radiusdata}
 \varepsilon_r=\mu(I_r^c),\qquad
 \Omega_r=\sup_{I_r^c}W,\qquad
 \alpha_r=\frac{\rho(r)}{4c(r)}
 =\frac1{2(e^r+1)(1-e^{-2r})}.
\end{equation}

\begin{lemma}[Folded local variance with global centering]\label{a:lem:fold}
If $f\in\D$ and $\mu(f)=0$, put $T=\norm f_\mu^2$ and
$A_r=\int_{I_r}|f|^2\,d\mu$. Then
\begin{equation}\label{a:eq:foldbound}
 \E[f]\ge\alpha_r(A_r-\varepsilon_rT),\qquad
 A_r\le\frac{\E[f]}{\alpha_r}+\varepsilon_rT.
\end{equation}
\end{lemma}
\begin{proof}
On $[0,r]$, let $d\nu(a)=4c(a)\kappa(a)\,da$. Evenness gives
\[
 \nu([0,r])=\mu(I_r),\qquad
 \int_0^r|f|^2\,d\nu=A_r,\qquad
 \int_0^rf\,d\nu=\int_{I_r}f\,d\mu.
\]
Restrict the full ordered-pair Gamma integral to $I_r^2$. The two same-sign regions and the opposite-sign region fold to
\begin{equation}\label{a:eq:foldidentity}
 2\int_{0<a<b<r}[\rho(b-a)+\rho(b+a)]
       \kappa(a)\kappa(b)|f(b)-f(a)|^2\,da\,db.
\end{equation}
For $0<a<b$, the complete positive exponential series gives
\begin{align*}
 \rho(b-a)+\rho(b+a)
 &=2\sum_{j\ge0}e^{-(2j+1/2)b}\cosh((2j+1/2)a)\\
 &\ge2c(a)\rho(b).
\end{align*}
Relative to $d\nu(a)d\nu(b)$, the coefficient in~\eqref{a:eq:foldidentity} is therefore at least
\[
 \frac{\rho(b-a)+\rho(b+a)}{8c(a)c(b)}
 \ge\frac{\rho(b)}{4c(b)}\ge\alpha_r,
\]
since $\rho$ decreases and $c$ increases on the positive half-line. Tonelli's theorem justifies both the restriction and the positive series. The ordered-pair variance identity yields
\[
 \E[f]\ge\alpha_r\left((1-\varepsilon_r)A_r
                  -\left|\int_{I_r}f\,d\mu\right|^2\right).
\]
The original global mean is retained:
\[
 \left|\int_{I_r}f\,d\mu\right|^2
 =\left|\int_{I_r^c}f\,d\mu\right|^2
 \le\varepsilon_r(T-A_r).
\]
Substitution proves~\eqref{a:eq:foldbound}. This proof applies to complex profiles on the entire maximal domain; it makes no local-mean substitution.
\end{proof}

Because $w_{\max}<1/2$,
\begin{equation}\label{a:eq:masswindow}
 \norm{\kappa f}_2^2=\int_\R W|f|^2\,d\mu
 \le\frac12A_r+\Omega_rT.
\end{equation}
Dropping only the nonnegative $\tg$ term in~\eqref{a:eq:masscontrol}, and then using Lemma~\ref{a:lem:fold}, gives the general template
\begin{equation}\label{a:eq:seedtemplate}
 \E[f]\ge
 \frac{63/128-9\Omega_r-(9/2)\varepsilon_r}
      {1+9/(2\alpha_r)}\,T
\end{equation}
whenever the numerator is positive. The restriction of a nonnegative integral supplies a lower bound; it does not replace the underlying energy by a local model.

\subsubsection{The radius \texorpdfstring{$\log2$}{log 2}}

At $r_0=\log2$, formula~\eqref{a:eq:radiusdata} becomes
\begin{equation}\label{a:eq:alphalog2}
 \alpha_{r_0}=\frac29,\qquad
 A_{r_0}\le\frac92\E[f]+\varepsilon_{r_0}T.
\end{equation}
The complete theta envelope holds at every nonnegative radius. Since $4r-3e^{2r}$ decreases,
\begin{equation}\label{a:eq:Omegalog2}
 \Omega_{r_0}\le42e^{4r_0-3e^{2r_0}}=672e^{-12}.
\end{equation}
For the mass tail, use $2c(r)\le2e^{r/2}$ and the envelope for $\kappa$ itself. Evenness and $t=e^{2r}$ give
\begin{align}
 \varepsilon_{r_0}
 &\le168\int_{r_0}^\infty e^{5r-3e^{2r}}\,dr
 =84\int_4^\infty t^{3/2}e^{-3t}\,dt\notag\\
 &\le42\int_4^\infty t^2e^{-3t}\,dt
 =\frac{2380}{9}e^{-12}.\label{a:eq:epslog2}
\end{align}
Here $t^{3/2}\le t^2/2$ for $t\ge4$, and integration by parts gives
\[
 \int_4^\infty t^2e^{-3t}\,dt
 =e^{-12}\left(\frac{16}{3}+\frac89+\frac2{27}\right)
 =\frac{170}{27}e^{-12}.
\]
Equations~\eqref{a:eq:Omegalog2}--\eqref{a:eq:epslog2} are full exterior estimates, not values sampled at the boundary.

\begin{proof}[Proof of the variance bound in Theorem~\ref{a:thm:mainseed}]
For centered $f$, combine~\eqref{a:eq:masscontrol}, \eqref{a:eq:masswindow}, and~\eqref{a:eq:alphalog2} to get
\[
 \E[f]\ge\left(\frac{63}{128}-9\Omega_{r_0}
                       -\frac92\varepsilon_{r_0}\right)T
                         -\frac{81}{4}\E[f].
\]
Move the last term to the left. The coefficient of $e^{-12}$ in the tail cost is
$9\cdot672+(9/2)(2380/9)=6048+1190=7238$. Therefore
\[
 \E[f]\ge\frac4{85}\left(\frac{63}{128}-7238e^{-12}\right)T
 =c_{\log2}T.
\]
By~\eqref{a:eq:elementarye}, $e^{12}>160000$, whence
\begin{align*}
 c_{\log2}&>\frac4{85}\left(\frac{63}{128}-\frac{7238}{160000}\right)
 =\frac4{85}\frac{8939}{20000}=\frac{8939}{425000}=:c_3,\\
 c_3-\frac1{48}
 &=\frac{8939\cdot48-425000}{425000\cdot48}
 =\frac{509}{2550000}>0.
\end{align*}
For general $f\in\D$, replace $f$ by $f-\mu(f)$. Constants belong to $\D$, the energy is unchanged, and the squared norm becomes $\Var_\mu(f)$. This proves~\eqref{a:eq:mainseed} on the full original domain.
\end{proof}

No optimality is asserted for $c_{\log2}$. It is the explicit coefficient proved by this particular radius and these estimates. In particular the coarse template~\eqref{a:eq:seedtemplate}, with~\eqref{a:eq:masscontrol} fixed, is bounded above by $63/128<1/2$: its loss terms are nonnegative and its denominator is at least one. That observation limits this template, not other possible estimates for the original energy.

\subsubsection{The exact mean and the inertia interval}

Let $\Hh_c=\{u\in\Hw:m(u)=0\}$ and $X_c=X\cap \Hh_c$. Since
$\norm\kappa_{\Hh}=1$ and $\ip u\kappa_{\Hh}=2m(u)$, the map
\begin{equation}\label{a:eq:centering}
 \mathsf C u=u-2m(u)\kappa
\end{equation}
is the exact weighted orthogonal projection onto $\Hh_c$, and it preserves $X$. The full-domain bridge gives, for $u=a\kappa+w$ with $w\in X_c$ and $a=2m(u)$,
\begin{align}
 \norm u_{\Hh}^2&=|a|^2+\norm w_{\Hh}^2,\label{a:eq:meansplitnorm}\\
 \ell_\delta[u]&=(\delta-1/2)|a|^2+\ell_\delta[w].\label{a:eq:meansplit}
\end{align}
This is weighted centering, not ordinary-$L^2$ projection away from $\kappa$.

\begin{proof}[Proof of the inertia conclusion in Theorem~\ref{a:thm:mainseed}]
Write $g_\delta=c_{\log2}+\delta-1/2$. The variance bound implies
\begin{equation}\label{a:eq:seedcoercivity}
 \ell_\delta[w]\ge g_\delta\norm w_{\Hh}^2\qquad(w\in X_c).
\end{equation}
For $1/2-c_{\log2}<\delta<1/2$, this margin is positive. The head--tail bound and $w_{\max}<1/2$ also give
\[
 \ell_\delta[w]\ge\tg[w]-C\norm w_{\Hh}^2,
 \qquad C=9/2+r_4.
\]
Combine this estimate with~\eqref{a:eq:seedcoercivity}, assigning weight
$\lambda=g_\delta/[2(g_\delta+C)]$ to the first one. Then
\begin{equation}\label{a:eq:fullseedcoercivity}
 \ell_\delta[w]\ge\lambda\tg[w]
                         +\frac{g_\delta}{2}\norm w_{\Hh}^2.
\end{equation}
Thus the centered infinite complement is coercive in the full form norm. By~\eqref{a:eq:meansplit}, its exact elimination leaves the scalar $\delta-1/2<0$. The maximal strictly negative form dimension is one and the bilinear form kernel is zero. For a lower-bounded compact-resolvent operator, projection onto its negative spectral subspace is injective on every strictly negative form subspace, while that spectral subspace attains the bound. Thus the form negative index equals the number of negative eigenvalues with multiplicity. By~\eqref{a:eq:Ldef}, a form-kernel vector is precisely an operator-domain vector with representing row zero, so the nullities also agree.

At $\delta_3=23/48$, the scalar is $-1/48$ and
$g_{\delta_3}>c_3-1/48=509/2550000$, proving~\eqref{a:eq:mainmargin}.
The endpoint $1/2-c_{\log2}$ is not certified by the non-strict variance inequality alone. The rational endpoint $1/2-c_3$ lies strictly inside the certified interval, with margin $c_{\log2}-c_3>0$.
\end{proof}

Zero operator nullity does not mean that an indefinite quadratic form has no isotropic vectors. Likewise, the mean Schur scalar $-1/48$ is a congruence entry, not the lowest ordinary-$L^2$ eigenvalue. Indeed the source row gives
$\ell_\delta(\kappa,z)=(2\delta-1)\overline{m(z)}$ on $X$; for $\delta\ne1/2$ its distributional representative is $(2\delta-1)c\notin L^2(dx)$. Hence $\kappa$ is a negative form direction for $0<\delta<1/2$, but not an operator-domain eigenvector of $L_\delta$.


\subsection{The local interval and the two computational premises}\label{c:sec:premises}
Set $B_0=4/5$, $I=(-B_0,B_0)$ and $O=\R\setminus[-B_0,B_0]$. This scalar radius is distinct from the weighted operator $B$ in~\eqref{eq:Db-physical}. The local minimum domain $\D_I$ is the closure of even $C_c^\infty(I)$, zero extended, in $(\tg+\norm{\cdot}_2^2)^{1/2}$. The form
\[
 q[u]=q_0[u]+2|m(u)|^2
\]
has restriction $q_I=q|_{\D_I}$ and ordinary physical $L^2(I)$ representative $L_I$. The sharp-cut and exterior minimum-domain statements used below are proved in Appendix~\ref{c:app:cuts}; they are not additional computational premises.


\begin{hypothesis}[Accepted local input theorem]\label{c:hyp:local}
On the whole domain $\D_I$, $q_I[u]\ge0$.
\end{hypothesis}
The accepted computer-assisted input states the stronger inequality $q_I[u]\ge10^{-22}\norm u_2^2$. Only its sign is used in the first theorem. Its physical $L^2$ coefficient is not transferred unchanged to a weighted global gap. Appendix~\ref{c:app:provenance} identifies the accepted evidence; it does not claim a fresh replay.

\begin{theorem}[Local-sign branch]\label{c:thm:localglobal}
Under Hypothesis~\ref{c:hyp:local}, every even $f\in\D$ satisfies
\begin{equation}\label{c:eq:49307}
 \E[f]\ge\frac{49307}{100000}\Var_\mu(f).
\end{equation}
The same comparison permits the endpoint
\begin{equation}\label{c:eq:lambdaendpoint}
 \begin{gathered}
 \lambda_{\rm cmp}=\frac12-\frac12\left(a_1+\sqrt{a_1^2+4L_1^2}\right)
 >\frac{49307}{100000},\\
 a_1=\frac{39703}{56448000},\qquad L_1=\frac{2099}{320000},\qquad
 \E[f]\ge\lambda_{\rm cmp}\Var_\mu(f).
 \end{gathered}
\end{equation}
No local inverse-response enclosure is needed. The displayed rational coefficient improves $493/1000$ by exactly $7/100000$; it does not approach $1/2$ through a parameter family.
\end{theorem}

\begin{theorem}[Additional-response branch]\label{c:thm:conditionalglobal}
Under the complete-operator and centered-response contracts in Hypothesis~\ref{c:hyp:response}, every even $f\in\D$ satisfies
\begin{equation}\label{c:eq:newglobal}
 \E[f]\ge\frac{2861149}{5760000}\Var_\mu(f).
\end{equation}
These contracts imply local positivity but are not implied by Hypothesis~\ref{c:hyp:local}. Their numerical enclosure assertions are inputs, not computations established here.
\end{theorem}
The latter coefficient exceeds $71501077/144000000$, obtained from a larger source cost under the same contracts, by exactly $24/125000$. This comparison concerns conditional implications, not the endpoint $1/2$. Arguments are given for real-even functions; decomposition into real and imaginary parts gives the complex statements. Inner products are linear in the first argument. 


\subsection{Analytic estimates and the two source coordinates}\label{c:sec:estimates}
Write $\alpha=\sqrt W$, and set
\[
 A(x)=\frac{(4t^2-6t)e^{-t}}{1+e^{-x}}+128t^2e^{-4t},
 \quad t=\pi e^{2x},\qquad A=A({B_0}).
\]
Appendix~\ref{c:app:scalar} proves
\begin{gather}
 A<1/9216,\quad W({B_0}+h)\le Ae^{-25h},\quad
 \int_{B_0}^\infty W<1/230400,\label{c:eq:envelope}\\
 \sup W<1/2,\quad \int_0^\infty W\le1/8,\quad
 p_*:=\mu(O)<3/49000,\quad \cg<383/70<11/2,\label{c:eq:basiccaps}\\
 W(1/10)<21/50,\quad W(1/4)<1/4,\quad W(9/20)<3/40,\label{c:eq:pointW}\\
 \kappa(2/5)<1/4,\quad \kappa(1/2)<7/100,\quad
 \kappa(3/5)<1/50,\quad \kappa(7/10)<3/1000,\quad
 \kappa(3/4)<1/1000.\label{c:eq:pointkappa}
\end{gather}
Both $\kappa$ and $W$ decrease on $[1/10,\infty)$. Also,
\begin{equation}\label{c:eq:rhoelementary}
 \rho(s)\le\frac1{2s}+\frac14,\qquad
 \frac56<\rho({B_0})<\frac9{10},\qquad \pi^2<10.
\end{equation}
These are full-series analytic inequalities.

There are two source norms. The weighted profile interface $\mathcal J=\mathcal J_p+\mathcal J_\Gamma$ is the inside-to-outside operator in
\begin{equation}\label{c:eq:energycross}
 \E[g+h]=\E[g]+\E[h]-2\ip{\mathcal Jg}{h}_\mu,
 \qquad \supp g\subset\overline I,\quad\supp h\subset\overline O.
\end{equation}
Under the unitary map
\begin{equation}\label{c:eq:foldU}
 (U_+ f)(x)=\sqrt{4c(x)\kappa(x)}f(x),\qquad x>0,
\end{equation}
its Gamma kernel is $\sqrt{W(x)W(y)}[\rho(x-y)+\rho(x+y)]$, $x>{B_0}$, $0<y<{B_0}$.

The physical exterior source in the second branch is
\begin{equation}\label{c:eq:physicalsource}
 (Ju)(y)=\int_O\rho(|x-y|)u(x)\,dx+
 \sum_{n\ge2}a_n[u(y+s_n)+u(y-s_n)],\qquad y\in I,
\end{equation}
where exterior inputs are extended by zero. Its input norm is
$\wo(u)=\int_O|u|^2/W\,dx$, and its output norm is physical $L^2(I)$.
With $v(x)=\sqrt2u(x)/\alpha(x)$ on $x>{B_0}$, the input norm becomes $\norm v_2^2=\wo(u)$ and the output fold is multiplication by $\sqrt2$. The extra factor $\sqrt{W(y)}$ belongs only to the profile interface.

\begin{lemma}[Carleman bound]\label{c:lem:carleman}
The kernel $(s+t)^{-1}$ on the positive half-line has norm at most $\pi$, as does every rectangular restriction.
\end{lemma}
\begin{proof}
For the Schur weight $w(s)=s^{-1/2}$,
\[
 \int_0^\infty\frac{t^{-1/2}}{s+t}\,dt
 =2s^{-1/2}\int_0^\infty\frac{dr}{1+r^2}=\pi s^{-1/2}.
\]
The transposed integral is identical. Weighted Cauchy--Schwarz proves the assertion on compact tests away from zero, and density completes it.
\end{proof}

\subsection{Four weighted prime channels and all their overlaps}\label{c:sec:weightedprime}
For the folded adjoint of the profile interface the branches are
\begin{align*}
 (P_nv)(y)&=a_n\sqrt{W(y)W(s_n+y)}v(s_n+y)\one_{s_n+y>{B_0}},\\
 (M_nv)(y)&=a_n\sqrt{W(y)W(s_n-y)}v(s_n-y)\one_{s_n-y>{B_0}}.
\end{align*}
If $y-s_n>0$, it is still inside; a negative exterior input gives exactly the reflected-minus branch. Retain $H_p=P_2+M_3+M_4+M_5$. Its domains, up to null endpoints, are
\begin{center}\small
\begin{tabular}{@{}lll@{}}\toprule
Channel & Inside output $y$ & Exterior input $x$\\\midrule
$P_2$ & $({B_0}-\log2,{B_0})$ & $({B_0},{B_0}+\log2)$\\
$M_3$ & $(0,\log3-{B_0})$ & $({B_0},\log3)$\\
$M_4$ & $(0,\log4-{B_0})$ & $({B_0},\log4)$\\
$M_5$ & $(0,{B_0})$ & $(\log5-{B_0},\log5)$\\\bottomrule
\end{tabular}
\end{center}
Here $M_2=0$ and $a_6=0$. Appendix~\ref{c:app:scalar} gives
\begin{equation}\label{c:eq:coefficientcaps}
 a_2<99/200,\quad a_3<16/25,\quad a_4<7/20,\quad a_5<29/40,
 \qquad a_n\le1.
\end{equation}
In particular $a_4=(\log2)/2$, not $(\log4)/2$.

The four channel-square weights, in units of $A$, are bounded by
$21/50,1/4,3/40,1/100$. For $P_2$, $y>{B_0}-\log2>1/10$.
For $M_3$ split at $y=1/4$: below it $x-{B_0}>6/125$ gives
$W(x)W(y)<Ae^{-6/5}/2<A/4$.
For $M_4$ split at $9/20$: below it $x-{B_0}>17/125$ gives $Ae^{-17/5}/2<A/40$.
For $M_5$ split at $3/5$: below it $x-{B_0}>209/1000$ gives $A/280<A/100$.
Above each split use \eqref{c:eq:pointW}--\eqref{c:eq:pointkappa}.
The common input diagonal is at most
\begin{align}
 A\bigl[(99/200)^2(21/50)&+(16/25)^2/4
 +(7/20)^2(3/40)+(29/40)^2/100\bigr]\notag\\
 &=A\frac{879017}{4000000}<\frac{11A}{50}.\label{c:eq:weighteddiag}
\end{align}
No output orthogonality is assumed.

A cross operator has the form $Cv(x)=\one_D(x)w(x)v(\phi(x))$,
where $\phi$ is injective and $|\phi'|=1$, so $\norm C\le\sup_D|w|$.
The pair $C+C^*$ costs twice this bound, even when $D$ is not reflection invariant.
Every pair is listed below; the last column bounds $\sqrt{W(x)W(x')}/A$.
\begin{center}\small
\begin{tabular}{@{}llll@{}}\toprule
Pair & First input domain & Second input $x'$ & Cap\\\midrule
$P_2,M_3$ & $({B_0},\log6-{B_0})$ & $\log6-x$ & $1/10$\\
$P_2,M_4$ & $({B_0},\log8-{B_0})$ & $\log8-x$ & $1/300$\\
$P_2,M_5$ & $({B_0},{B_0}+\log2)$ & $\log10-x$ & $1/6000$\\
$M_3,M_4$ & $({B_0},\log3)$ & $x+\log(4/3)$ & $1/32$\\
$M_3,M_5$ & $({B_0},\log3)$ & $x+\log(5/3)$ & $1/500$\\
$M_4,M_5$ & $({B_0},\log4)$ & $x+\log(5/4)$ & $1/15$\\\bottomrule
\end{tabular}
\end{center}
These rational comparisons follow from the envelope and are proved in Appendix~\ref{c:app:scalar}. In the first three rows use $W(y)<21/50$, and in the next two use $W(y)<1/2$.
In the last row the sharper full product is below $A/200$:
for $y\ge9/20$ use $(3/40)(1/15)$; otherwise the sum of exterior distances is
$\log20-2{B_0}-2y>495/1000$, its exponent exceeds $6$, and the product is below $A/800$.
The six paired costs in units of $A$ are at most
\[
 27/1000,\quad1/2000,\quad3/50000,\quad7/1000,\quad1/1000,\quad13/5000.
\]
Their sum is $1908/50000<1/25$. Thus $\norm{H_p}^2<(13/50)A<(51/100)^2A$.

The exact complement is $\sum_{n\ge3}P_n+\sum_{n\ge7}M_n$.
From $e^{B_0}<9/4$, $e^{2{B_0}}<5$ and $25/2>12$,
\begin{gather}
 \norm{P_n}\le\sqrt{A/2}(9/(4n))^{12},\qquad
 \norm{M_n}\le\sqrt{A/2}(5/n)^{12},\notag\\
 \sum_{n\ge3}(9/(4n))^{12}\le\frac{14}{11}(3/4)^{12}<1/24,\qquad
 \sum_{n\ge7}(5/n)^{12}\le\frac{18}{11}(5/7)^{12}<3/100.\label{c:eq:bothtails}
\end{gather}
The last comparisons are $178564176<184549376$ and
$439453125000<456762477633$. Both complete tails converge in operator norm.
Their triangle inequalities pay all head--tail and tail--tail products.
Since $\sqrt2>7/5$,
\begin{equation}\label{c:eq:weightedprimenorm}
 \norm{\mathcal J_p}<\sqrt A\left(\frac{51}{100}+\frac{43}{840}\right)
 <\frac9{16}\sqrt A<\frac3{512}<\frac{59}{10000}.
\end{equation}
The scalar slack before multiplication by $\sqrt A$ is $11/8400$.

\subsection{The exact mean repair and six Gamma strips}\label{c:sec:meanrepair}
For $x>{B_0}$ and $0<y<{B_0}$ define
\begin{equation}\label{c:eq:remainderkernel}
 B_\Gamma(x,y)=\rho(x-y)+\rho(x+y)-2c(y)\rho(x).
\end{equation}
The full positive series for $\rho$ gives
\[
 B_\Gamma(x,y)=2\sum_{j\ge1}e^{-(2j+1/2)x}
 [\cosh((2j+1/2)y)-\cosh(y/2)]\ge0.
\]
It decreases in $x$ and increases in $y$. Only the $j=0$ term cancels.
Let $\mathcal T_\Gamma$ have kernel $\sqrt{W(x)W(y)}B_\Gamma(x,y)$ and set
\[
 e(y)=U_+ 1(y)=2\sqrt2c(y)\sqrt{W(y)},\qquad
 z(x)=\frac{\sqrt{W(x)}\rho(x)}{\sqrt2}.
\]
For an inside profile $g$ the exact identity is
\begin{equation}\label{c:eq:rankrepair}
 U_+ \mathcal J_\Gamma g=\mathcal T_\Gamma U_+ g+z\mu(g),\qquad
 \mu(g)=\ip{U_+ g}{e}.
\end{equation}
Thus the rank-one term has not been discarded.

For $y\le b<{B_0}$, monotonicity and \eqref{c:eq:rhoelementary} give
\[
 0\le B_\Gamma(x,y)\le B_\Gamma({B_0},b)
 <\frac{{B_0}}{{B_0}^2-b^2}-\frac76=:F(b).
\]
The subtraction is made at $x={B_0}$, not by imposing a false lower bound for $\rho(x)$ in the exterior.
Divide the input at $2/5,1/2,3/5,7/10,3/4$.
The first five strips give valid Hilbert--Schmidt estimates:
\begin{center}\small
\begin{tabular}{@{}llll@{}}\toprule
Strip & $\int W(y)\,dy$ cap & $F(b)$ & Squared norm cap\\\midrule
$(0,2/5)$ & $1/8$ & $1/2$ & $1/7372800$\\
$(2/5,1/2)$ & $1/80$ & $23/26$ & $529/12460032000$\\
$(1/2,3/5)$ & $7/2000$ & $71/42$ & $5041/116121600000$\\
$(3/5,7/10)$ & $1/1000$ & $25/6$ & $625/8294400000$\\
$(7/10,3/4)$ & $3/40000$ & $1703/186$ & $2900209/106278912000000$\\\bottomrule
\end{tabular}
\end{center}
Each row uses the exterior integral $1/230400$.
The squared norms are respectively below $(140,43,44,76,28)/10^9$.

On the touching strip $(3/4,4/5)$, $W(y)<1/2000$ and
\[
 B_\Gamma(x,y)\le\rho(x-y)+\rho(x+y)
 \le\frac1{2(x-y)}+\frac{51}{62}.
\]
The singular term is half a Carleman kernel after $s=x-{B_0}$, $t={B_0}-y$.
Its norm is below
$\frac\pi2(9216\cdot2000)^{-1/2}<37/100000$.
Its squared cap is $1/7372800$, and $1369\cdot7372800>10^{10}$.
The bounded term has norm at most
\[
 \frac{51}{62}(230400\cdot40000)^{-1/2}
 =\frac{17}{1984000}<\frac9{1000000}.
\]
Thus the last strip norm is below $379/10^6<19/50000$, whose square is below $145/10^9$.
No Hilbert--Schmidt claim is made for its touching singularity.
The input projections are orthogonal, giving
$\mathcal T_\Gamma \mathcal T_\Gamma^*=\sum_iT_iT_i^*$, without output orthogonality.
Consequently
\begin{equation}\label{c:eq:gammaremainnorm}
 \norm{\mathcal T_\Gamma}^2<\frac{140+43+44+76+28+145}{10^9}
 =\frac{476}{10^9}<\left(\frac7{10000}\right)^2.
\end{equation}
Both physical boundaries are already included in the even fold.

\subsection{The global comparison from local sign}\label{c:sec:firstproof}
For an exterior profile $h$, the physical bridge gives
\[
 \E[h]=\tfrac12\norm h_\mu^2+\tg[\kappa h]
 -\cg\norm{\kappa h}_2^2-p[\kappa h].
\]
The full-prime signed identity in Appendix~\ref{c:sec:signed} is analytic and needs neither response data nor an exterior moment constraint. With $u=\kappa h$ it gives
\begin{equation}\label{c:eq:exteriorenergy}
 \E[h]\ge\left(\tfrac12-\frac{31}{46080}\right)\norm h_\mu^2
 +\tg[\kappa h]+D_p(\kappa h/\alpha),
\end{equation}
where $D_p$ is the complete nonnegative prime slack defined in \eqref{c:eq:slack}.
Thus the cap $31/46080$, not the earlier coarser $7/10240$, is available in the local-sign branch.

For a real centered profile $v=g+h$ cut on $I$ and $O$, one has
$\mu(g)=-\mu(h)$, not $\mu(g)=0$. Put
\[
 H=U_+ h,\qquad e_O=U_+ \one_O=2\sqrt2c\alpha\big|_{({B_0},\infty)},\qquad
 p_*:=\norm{e_O}^2=\mu(O).
\]
Then $\mu(h)=\ip H{e_O}$.
The removed term in \eqref{c:eq:rankrepair} enters the energy with sign
\[
 -2\mu(g)\ip zH=2\mu(h)\ip zH.
\]
This product of distinct exterior functionals is not assigned a sign on its own. Local positivity simultaneously supplies the penalty
$-\tfrac12|\mu(g)|^2=-\tfrac12|\mu(h)|^2$.

\begin{lemma}[The combined mean form]\label{c:lem:coupledmean}
Let $e_O,z$ be real vectors in the folded exterior space with
$0\le z\le e_O/4$ pointwise and $p_*=\norm{e_O}^2$.
For every real exterior vector $H$,
\begin{equation}\label{c:eq:combinedmean}
 -\tfrac12|\ip H{e_O}|^2+2\ip H{e_O}\ip zH
 \ge-\tfrac12p_*\norm H^2.
\end{equation}
More precisely, the optimal nonnegative loss for this form is
$\norm{e_O}\norm w-\ip{e_O}w$, where $w=z-e_O/4$.
\end{lemma}
\begin{proof}
With $(a\otimes b)H=\ip Hb a$ on the real space, the combined operator is
\[
 K=e_O\otimes z+z\otimes e_O-\tfrac12e_O\otimes e_O
   =e_O\otimes w+w\otimes e_O.
\]
Its two possible nonzero eigenvalues are
$\ip{e_O}w\pm\norm{e_O}\norm w$; the formula includes dependent vectors by continuity and the remaining spectrum is zero.
Since $-e_O/4\le w\le0$, one has
$\norm w\le\norm{e_O}/4$ and $\ip{e_O}w\ge-p_*/4$.
Consequently $K\succeq-p_*I/2$, proving the inequality and the stated exact loss.
\end{proof}
Here the pointwise hypothesis is automatic:
\[
 \frac{z(x)}{e_O(x)}=\frac{\rho(x)}{4c(x)}<\frac14\qquad(x>{B_0}),
\]
since $\rho$ decreases, $\rho({B_0})<9/10$ and $c(x)\ge1$.
The estimate removes the need for the separate absolute rank charge; it does not discard the local mean penalty.

Local sign gives $\E[g]\ge\tfrac12\norm g_\mu^2-\tfrac12|\mu(g)|^2$.
Use \eqref{c:eq:exteriorenergy}, Lemma~\ref{c:lem:coupledmean},
the complete weighted prime bound $\norm{\mathcal J_p}<3/512$, and
$\norm{\mathcal T_\Gamma}<7/10000$ to obtain
\begin{equation}\label{c:eq:twobytwo}
 \E[v]\ge\tfrac12\norm g_\mu^2+(\tfrac12-a_1)\norm h_\mu^2
 -2L_1\norm g_\mu\norm h_\mu,
\end{equation}
where
\begin{equation}\label{c:eq:improvedconstants}
 a_1=\frac{31}{46080}+\frac3{98000}=\frac{39703}{56448000},
 \qquad L_1=\frac3{512}+\frac7{10000}=\frac{2099}{320000}.
\end{equation}
Only $p_*<3/49000$ was used for the combined mean loss. The nonnegative terms in \eqref{c:eq:exteriorenergy} were dropped only after the signed identity had been applied.

For $\delta_1=693/100000$, the residual comparison matrix is
\[
 M_{\delta_1}=\begin{pmatrix}\delta_1&-L_1\\-L_1&\delta_1-a_1\end{pmatrix}.
\]
The exact principal-minor checks are
\begin{equation}\label{c:eq:margin}
 \delta_1-a_1=\frac{8787041}{1411200000}>0,\qquad
 \det M_{\delta_1}=\frac{2244433}{17920000000000}>0.
\end{equation}
For the determinant, on denominator $1128960000000000$, the numerators of
$\delta_1^2$, $\delta_1a_1$ and $L_1^2$ are respectively
$54218191104$, $5502835800$ and $48573956025$.
Their signed difference is $141399279=63\cdot2244433$.
Thus \eqref{c:eq:twobytwo} proves the coefficient
$1/2-\delta_1=49307/100000$.
More generally $M_\delta\succeq0$ precisely when
\[
 \delta\ge\frac12\left(a_1+\sqrt{a_1^2+4L_1^2}\right),
\]
which gives \eqref{c:eq:lambdaendpoint}. This threshold is strictly positive.
It is an exact endpoint for this comparison, not an assertion of the optimal constant of the actual energy, and not a family of signs tending to zero loss.

First apply the argument to $v=f-\mu(f)$ with $f$ even compact smooth.
Constants and sharp cuts belong to the corresponding minimum domains by Appendix~\ref{c:app:cuts}.
The bounded interface operators and norm-convergent tails justify the cross identity on these cuts.
Finally use the global smooth form core, along which energy and variance converge.
This proves Theorem~\ref{c:thm:localglobal} on the full maximal domain.
Real and imaginary parts give the complex statement without assigning a sign to the rank-one cross term separately.

\subsection{A joint Gram estimate for the physical source}\label{c:sec:joint}
The output coordinate in this section is physical $L^2(I)$, not weighted profile space.
For $v(x)=\sqrt2u(x)/\alpha(x)$ on $x>{B_0}$ the folded Gamma source is
\[
 (\mathcal T_\Gamma^{\rm phys}v)(y)=\int_{B_0}^\infty
 [\rho(x-y)+\rho(x+y)]\alpha(x)v(x)\,dx,\qquad0<y<{B_0}.
\]
There is no further factor $\sqrt2$.
Because $x+y>{B_0}$, \eqref{c:eq:rhoelementary} implies the positive majorant
\[
 0\le\rho(x-y)+\rho(x+y)\le\frac1{2(x-y)}+b,\qquad b=\frac{23}{20}.
\]
Define
\[
 (Cv)(y)=\int_{B_0}^\infty\frac{\alpha(x)v(x)}{2(x-y)}\,dx,\qquad
 (Rv)(y)=b\int_{B_0}^\infty\alpha(x)v(x)\,dx.
\]
Pointwise $|\mathcal T_\Gamma^{\rm phys}v|\le(C+R)|v|$.
This is domination, not equality with the actual Gamma source.
The touching kernel is half a Carleman kernel, hence
\[
 \norm C\le\frac\pi2\sqrt A.
\]
Set
\[
 I_0=\int_{B_0}^\infty\alpha(x)^2\,dx,\qquad
 I_2=\int_{B_0}^\infty\alpha(x)^2\log^2\frac{x}{x-{B_0}}\,dx.
\]
Then $I_0\le A/25$ and
\begin{equation}\label{c:eq:regulargram}
 \norm R^2=b^2B_0I_0\le\frac{529}{12500}A,\qquad
 \norm{C^*R}=\frac b2\sqrt{I_2I_0}.
\end{equation}
The second identity follows from the exact rank-one formula
\[
 (C^*Rv)(x)=\frac b2\alpha(x)\log\frac{x}{x-{B_0}}
                  \int_{B_0}^\infty\alpha(t)v(t)\,dt.
\]

\subsubsection{The logarithmic cross integral}
With $t=25(x-{B_0})$,
\[
 I_2/A\le\frac1{25}\int_0^\infty e^{-t}\log^2(1+20/t)\,dt.
\]
Take $a=61/20>\log21$, since $e^3>20$ and $e^{1/20}>21/20$.
Writing $r(t)=\max\{-\log t,0\}$, one has
$\log(1+20/t)\le a+r(t)$.
The inequality $e^{-t}\le1-t+t^2/2$ on $[0,1]$ and
\[
 \int_0^1t^k(-\log t)\,dt=(k+1)^{-2},\qquad
 \int_0^1t^k\log^2t\,dt=2(k+1)^{-3}
\]
give
\[
 \int_0^1e^{-t}(-\log t)\,dt\le29/36,\qquad
 \int_0^1e^{-t}\log^2t\,dt\le193/108.
\]
Consequently
\begin{equation}\label{c:eq:I2}
 I_2/A\le\frac1{25}\left(\frac{3721}{400}+
       \frac{1769}{360}+\frac{193}{108}\right)
 =\frac{172837}{270000}<\left(\frac{81}{100}\right)^2.
\end{equation}
Thus $2\norm{C^*R}<(1863/10000)A$ and
\begin{equation}\label{c:eq:gammajoint}
 \norm{J_\Gamma}_{\wo\to L^2(I)}^2/A
 <\frac{121}{49}+\frac{529}{12500}+\frac{1863}{10000}
 <\frac{529}{196}.
\end{equation}
The last margin is $9524/9800000$.
This proves the standalone bound
$\norm{J_\Gamma}_{\wo\to L^2(I)}<(23/14)\sqrt A$.
The calculation concerns $C^*R$, not the square integral of a touching Gamma kernel.

\subsubsection{The physical prime head and complete tails}
In the same folded spaces the physical channels are
\[
 (P_nv)(y)=a_n\alpha(s_n+y)v(s_n+y)\one_{s_n+y>{B_0}},\quad
 (M_nv)(y)=a_n\alpha(s_n-y)v(s_n-y)\one_{s_n-y>{B_0}}.
\]
In this section $H_p=P_2+M_3+M_4+M_5$ refers to these physical channels.
They have the same domains and pair maps as in Section~\ref{c:sec:weightedprime}, but without $\sqrt{W(y)}$.
Their common input diagonal is at most $(422091/400000)A$.
Indeed below $x_5=\log5-{B_0}$ the square sum is at most $(6217/8000)A$;
above $x_5$, $\alpha(x)/\sqrt A<e^{-9/80}<9/10$, giving
\[
 D(x)/A\le\frac{81}{100}\bigl[(99/200)^2+(16/25)^2
 +(7/20)^2+(29/40)^2\bigr]=\frac{422091}{400000}.
\]
The six cross products use the last column of the table in Section~\ref{c:sec:weightedprime}.
Their paired sum is
\[
 \frac{198}{3125}+\frac{231}{200000}+\frac{957}{8000000}
 +\frac7{500}+\frac{116}{62500}+\frac{203}{6000}
 =\frac{109751}{960000}.
\]
Hence
\begin{equation}\label{c:eq:physicalhead}
 \norm{H_p}^2\le\frac{5613847}{4800000}A<\frac{169}{144}A.
\end{equation}
The exact complement $T_p=\sum_{n\ge3}P_n+\sum_{n\ge7}M_n$ obeys
\begin{equation}\label{c:eq:physicaltail}
 \norm{T_p}<\left(\frac1{24}+\frac3{100}\right)\sqrt A
\end{equation}
by the same complete integer tails \eqref{c:eq:bothtails}, now without $1/\sqrt2$.
Thus no prime power or head--tail cross term is omitted.

\subsubsection{Every singular Gamma and prime-head product}
Let $\sigma=25/2$. Direct substitution gives the \emph{input}-Gram kernels
\begin{align}
 (C^*P_n)(z,t)&=\frac{a_n\alpha(z)\alpha(t)}
                      {2(z-t+s_n)},&
 z>{B_0},\quad\max({B_0},s_n)<t<{B_0}+s_n,\label{c:eq:inputgramplus}\\
 (C^*M_n)(z,t)&=\frac{a_n\alpha(z)\alpha(t)}
                      {2(z+t-s_n)},&
 z>{B_0},\quad\max({B_0},s_n-{B_0})<t<s_n.\label{c:eq:inputgramminus}
\end{align}
They vanish outside the indicated intervals. They are not the output-Gram products $CP_n^*$ or $CM_n^*$.

For $P_2$, put $\ell=\log2$, $s=z-{B_0}$, $r=t-{B_0}\in(0,\ell)$.
The kernel is bounded by
$(a_2A/2)e^{-\sigma(s+r)}/(s+\ell-r)$.
On $r\le\ell/2$ the denominator is at least $\ell/2$, so an actual Hilbert--Schmidt majorant gives
\[
 \norm{(C^*P_2)_{\rm low}}/A
 \le\frac{a_2}{2\sigma\ell}<\frac{297}{10000}.
\]
On $r>\ell/2$, set $h=\ell-r$. The denominator is $s+h$, so use Carleman, retaining the factor $1/2$:
\[
 \norm{(C^*P_2)_{\rm high}}/A
 \le\frac{a_2\pi}{2}\,2^{-25/4}<\frac{1089}{89600}.
\]
This corner $(z,t)=({B_0},{B_0}+\log2)$ is not Hilbert--Schmidt.

For $M_3,M_4$, let $d_n=2{B_0}-s_n$. The bounds
$d_3>1/2$, $d_4>1/5$ genuinely separate their kernels.
The exponential square integrals give
\[
 \norm{C^*M_n}/A\le\frac{a_n}{4\sigma d_n},\qquad
 \norm{C^*M_3}/A<16/625,\quad \norm{C^*M_4}/A<7/200.
\]
For $M_5$, however, set $d=s_5-2{B_0}>0$,
$s=z-{B_0}$ and $h=t-(s_5-{B_0})$.
The kernel majorant is
\[
 \frac{a_5A}{2}e^{-\sigma d}\frac{e^{-\sigma(s+h)}}{s+h}.
\]
Although its two input locations are separated, the corner is touching.
Carleman, not a denominator gap or Hilbert--Schmidt estimate, gives
$\norm{C^*M_5}/A\le a_5\pi/2<319/280$.
All four products are therefore paid by
\begin{equation}\label{c:eq:CHp}
 \norm{C^*H_p}/A<
 \frac{297}{10000}+\frac{1089}{89600}+\frac{16}{625}
 +\frac7{200}+\frac{319}{280}
 =\frac{2781497}{2240000}<\frac54.
\end{equation}

\subsubsection{The regular-majorant product and the full source}
On the common exterior input coordinate,
\begin{align*}
 H_p^*1(t)&=\alpha(t)w(t),\\
 w(t)&=a_2\one_{({B_0},{B_0}+\log2)}(t)+a_3\one_{({B_0},\log3)}(t)
 +a_4\one_{({B_0},\log4)}(t)+a_5\one_{(\log5-{B_0},\log5)}(t).
\end{align*}
Since $w<221/100$ and $R=b\,1\otimes\alpha$, the exact rank-one product satisfies
\begin{equation}\label{c:eq:HpR}
 H_p^*R=b(H_p^*1)\otimes\alpha,\qquad
 2\norm{H_p^*R}\le2b(221/100)I_0\le\frac{5083}{25000}A.
\end{equation}
Expanding the bounded-operator Gram of $C+R+H_p$ now charges all three cross pairs:
\begin{align}
 \norm{C+R+H_p}^2/A
 &<\frac{121}{49}+\frac{169}{144}+\frac52+
   \frac{529}{12500}+\frac{1863}{10000}+\frac{5083}{25000}\notag\\
 &=\frac{121}{49}+\frac{169}{144}+\frac52+\frac{21597}{50000}
 <\frac{324}{49}.\label{c:eq:jointgram}
\end{align}
For the last comparison use $21597/50000<9/20$ and
$169/144+59/20=2969/720<203/49$, whose cross products are
$145481<146160$.
Positive-kernel domination gives
$|Jv|\le(C+R+H_p+T_p)|v|$.
Combining \eqref{c:eq:physicaltail} and \eqref{c:eq:jointgram} proves
\begin{theorem}[Complete analytic source bound]\label{c:thm:source}
The physical source extends from compact exterior tests to every even finite-$\wo$ input and satisfies
\begin{equation}\label{c:eq:sourcej}
 j:=\norm J_{\wo\to L^2(I)}
 <\left(\frac{18}{7}+\frac1{24}+\frac3{100}\right)\sqrt A
 =\frac{11101}{4200}\sqrt A<\frac{11101}{403200}<\frac{69}{2500}.
\end{equation}
This theorem needs neither local sign nor response enclosures.
\end{theorem}
The extension follows by density and boundedness; both full prime tails converge in operator norm. The domination statement is not a Loewner order between rectangular operators.

\subsection{The additional response contracts and inverse cost}\label{c:sec:response}
Let $e_0=(2{B_0})^{-1/2}$ and
$e_n={B_0}^{-1/2}\cos(n\pi x/{B_0})$, $n\ge1$, on $I$.
Let $E:\R^{64}\to L^2_{\rm even}(I)$ insert $e_0,\ldots,e_{63}$ and set
\[
 H_{64}=L_I+9EE^*,\qquad \mathsf U_{64}=H_{64}^{-1}E.
\]
Appendix~\ref{c:sec:lift} proves analytically that
\begin{equation}\label{c:eq:Hcoercivity}
 H_{64}\succeq61I/28,\qquad H_{64}^{-1}\preceq D_{\rm old},
\end{equation}
where $D_{\rm old}$ is the explicit cosine-diagonal majorant defined there.
These facts are not numerical enclosures.

Choose fixed real trial coefficients $C_{\rm trial}\in\R^{1024\times64}$,
defining columns $\widetilde{\mathsf U}_{64}$ in the first 1024 local cosine modes.
These columns belong to the local operator domain: their killed Gamma row has logarithmic endpoint singularities in $L^2$, and their remaining rows are bounded.
Set
\[
 R_{64}=E-H_{64}\widetilde{\mathsf U}_{64},\qquad
 F=I_{64}/9-E^*H_{64}^{-1}E.
\]
The regular-majorant operator $R$ in Section~\ref{c:sec:joint} is unrelated to $R_{64}$.

\begin{hypothesis}[Full residual and centered-response contracts]\label{c:hyp:response}
For the fixed trial matrix there are symmetric matrices $F_*,V_{\rm up}$,
and its positive-diagonal Cholesky factor $M$ with $F_*=MM^*$, such that:
\begin{enumerate}[label=\textup{(\alph*)}]
\item The actual complete-operator inequalities hold:
\begin{equation}\label{c:eq:contracta}
 0\prec F_*\preceq F,\qquad
 R_{64}^*H_{64}^{-1}R_{64}
 \preceq R_{64}^*D_{\rm old}R_{64}
 \preceq V_{\rm up}\preceq F_*/100.
\end{equation}
All residual rows, including those beyond 1024 and the infinite remainder, are included.
\item Set $V=\widetilde{\mathsf U}_{64} M^{-*}$ and $Z=C_{\rm trial}M^{-*}$.
For one fixed rational centering vector $a\in\R^{64}$, the exact finite-row, spatial-tail, Gamma-tail, and prime-tail quantities specified in Appendix~\ref{c:app:rows} have enclosures satisfying
\begin{equation}\label{c:eq:contractb}
 \sqrt{I_{\rm fin}}+e_X<\frac{81}{2500},\qquad
 e_\Gamma<\frac{39}{10000},\qquad
 b_{16}\norm{Z_{<64}}_{\mathrm{HS}}<10^{-21},\qquad
 \norm{Z_{64:1023}}_{\mathrm{HS}}<\frac13.
\end{equation}
Here $b_{16}=500/20^{25}$ bounds the entire omitted prime source after 16.
The rows and their centering refer to the same $V$, not to independently chosen response vectors.
\end{enumerate}
\end{hypothesis}
The theorem is valid for any fixed data satisfying these precise contracts.
An accepted finite matrix sign or a finite replay alone does not establish their full analytic enclosure interpretation.
Neither \eqref{c:eq:contracta} nor \eqref{c:eq:contractb} is inferred from Hypothesis~\ref{c:hyp:local}.
No new enclosure run is asserted.

To see local coercivity from (a), put $f_*=\lambda_{\min}(F_*)>0$ and
$K_0=3H_{64}^{-1/2}E$. Then
$K_0^*K_0=I-9F\preceq(1-9f_*)I$, so the same upper bound holds for
$K_0K_0^*$. Therefore
\begin{equation}\label{c:eq:Lcoercive}
 L_I=H_{64}^{1/2}(I-K_0K_0^*)H_{64}^{1/2}
 \succeq9f_*H_{64}\succeq(549f_*/28)I>0.
\end{equation}
The mean pole remains in $L_I$ throughout.

For $m_O(u)=\int_Ocu=0$, the fixed centered-row decomposition in Appendix~\ref{c:app:rows} gives
\begin{equation}\label{c:eq:responsebound}
 \norm{V^*Ju}\le r\sqrt{\wo(u)},\qquad
 r<\frac{363}{10000}+\frac j3+10^{-21}<\frac{91}{2000}.
\end{equation}
It is essential to insert the sharper fraction from \eqref{c:eq:sourcej} before rounding:
\[
 \frac{91}{2000}-\frac{363}{10000}-\frac{11101}{1209600}
 =\frac{683}{30240000}>10^{-21}.
\]
The coarser bound $j<69/2500$ alone leaves no displayed room for this last tail.

Woodbury gives the actual inverse
\[
 T=L_I^{-1}=H_{64}^{-1}+\mathsf U_{64}F^{-1}\mathsf U_{64}^*.
\]
For $s\in L^2_{\rm even}(I)$ set $h=\ip s{H_{64}^{-1}s}$.
Since $\mathsf U_{64}=\widetilde{\mathsf U}_{64}+H_{64}^{-1}R_{64}$, inverse order and (a) imply
\[
 \ip s{Ts}\le h+
 \left[\norm{V^*s}+\frac1{10}\sqrt h\right]^2.
\]
Indeed
$\norm{M^{-1}R_{64}^*H_{64}^{-1}s}\le\sqrt h/10$
follows by conjugating \eqref{c:eq:contracta}.
Young's inequality with parameter $1/20$ yields
\[
 \ip s{Ts}\le\frac{121}{100}h+\frac{21}{20}\norm{V^*s}^2.
\]
For $s=Ju$, $h\le(28/61)j^2\wo(u)$.
Using \eqref{c:eq:responsebound} and then $j<69/2500$,
\begin{align}
 \frac{\ip{Ju}{TJu}}{\wo(u)}
 &<\frac{847}{1525}\left(\frac{69}{2500}\right)^2
 +\frac{21}{20}\left(\frac{91}{2000}\right)^2\notag\\
 &=\frac{4032567}{9531250000}+\frac{173901}{80000000}
 <\frac{17}{40000}+\frac{87}{40000}=\frac{13}{5000}.\label{c:eq:inversecost}
\end{align}
For the first comparison use $847/1525<5/9$ and
$17/40000-4761/11250000=81/45000000>0$; the second uses $173901<174000$.
Including $u=0$, the universal statement is
$\ip{Ju}{L_I^{-1}Ju}\le(13/5000)\wo(u)$ on centered exterior inputs.

\subsection{The signed exterior form and the conditional global bound}\label{c:sec:conditionalproof}
Let $\D_O^0$ be the minimum closure of even $C_c^\infty(O)$ with
$m_O(u)=0$ in $(\tg+\wo)^{1/2}$.
For $u\in\D_O^0$ define
\[
 S[u]=\tg[u]-\cg\norm u_2^2-p[u]-\ip{Ju}{L_I^{-1}Ju}.
\]
This is the actual signed Schur form of the mean-inclusive local operator.
Appendix~\ref{c:sec:signed} proves, without a numerical premise,
\begin{equation}\label{c:eq:signedcapmain}
 \cg\norm u_2^2+p[u]\le\frac{31}{46080}\wo(u)-D_p(u/\alpha),
\end{equation}
where the full nonnegative slack is
\[
 D_p(g)=\sum_{n\ge2}a_n\int_{\{x,x+s_n\in O\}}
  \alpha(x)\alpha(x+s_n)|g(x+s_n)-g(x)|^2\,dx.
\]
Every prime power and both row orientations are retained.
Combining \eqref{c:eq:inversecost} and \eqref{c:eq:signedcapmain},
\begin{equation}\label{c:eq:Slower}
 S[u]\ge\tg[u]+D_p(u/\alpha)-\varepsilon_8\wo(u),\qquad
 \varepsilon_8=\frac{31}{46080}+\frac{13}{5000}
 =\frac{18851}{5760000}.
\end{equation}

For $f\in\D$ use the exterior-mean gauge
\[
 \beta_O=\mu(O)^{-1}\int_Of\,d\mu,\qquad
 u=\kappa\one_O(f-\beta_O),\quad
 v=\kappa\one_I(f-\beta_O).
\]
Appendix~\ref{c:app:cuts} puts $u$ in $\D_O^0$ and $v$ in $\D_I$.
Moreover
\[
 \wo(u)=\int_O|f-\beta_O|^2\,d\mu
 \le\int_O|f-\mu(f)|^2\,d\mu\le\Var_\mu(f).
\]
The full bridge and local square completion give
\begin{equation}\label{c:eq:squaremain}
 \E[f]-\tfrac12\Var_\mu(f)=q_I[v-L_I^{-1}Ju]+S[u].
\end{equation}
The local square is nonnegative by \eqref{c:eq:Lcoercive}.
Dropping only that square, $\tg$ and $D_p$ in \eqref{c:eq:Slower} proves
Theorem~\ref{c:thm:conditionalglobal}.
The exact comparison is
\[
 \frac{2861149}{5760000}-\frac{71501077}{144000000}
 =\frac{24}{125000}>0.
\]
The full-domain claim uses the proved equality of minimum and maximal energy domains, not a change of realization.

\subsection{The remaining sign and the absence of a uniform weighted gap}\label{c:sec:limits}
The exact remaining target is a spatially faithful inverse comparison, not a uniform estimate obtained by replacing its two costs by constants. Assume here that $L_I$ is strictly positive, as supplied by the accepted stronger local theorem $q_I\ge10^{-22}\norm{\cdot}_2^2$ or by Hypothesis~\ref{c:hyp:response}. Then the global half bound on $\D$ is equivalent to
\begin{equation}\label{c:eq:remainingS}
 S[u]\ge0\quad\hbox{for every }u\in\D_O^0.
\end{equation}
Sufficiency follows from \eqref{c:eq:squaremain}.
For necessity, take any $u\in\D_O^0$, set $v=L_I^{-1}Ju\in\D_I$ and
$f=(u+v)/\kappa$. The domain identification in Section~\ref{a:sec:domain} gives $f\in\D$; its exterior mean is zero and its local square vanishes, so its half-defect is $S[u]$.
If one assumes only $L_I\succeq0$, an ordinary inverse requires an additional coercivity or generalized-inverse and range argument.

By \eqref{c:eq:signedq0}, the target \eqref{c:eq:remainingS} is equivalently
\begin{equation}\label{c:eq:spatialtarget}
 \int_O(\cg W+V_p)|u/\alpha|^2\,dx+\ip{Ju}{L_I^{-1}Ju}
 \le\tg[u]+D_p(u/\alpha),\qquad u\in\D_O^0.
\end{equation}
Here $V_p$ is the exact complete row potential \eqref{c:eq:primepotential}.
The following analytic obstruction explains why the spatial potential and actual inverse response in this inequality cannot simply be replaced by the fixed upper cost $\varepsilon_8\wo$ and absorbed by a uniform weighted gap.

\begin{proposition}[No uniform weighted gap for the full positive slack]\label{c:prop:nogap}
On the original centered exterior minimum domain,
\begin{equation}\label{c:eq:nogap}
 \inf_{0\ne u\in\D_O^0}
 \frac{\tg[u]+D_p(u/\alpha)}{\wo(u)}=0.
\end{equation}
The same holds with $D_p$ replaced by its $n=2$ summand $D_2$.
Neither local positivity nor response contracts are required for this proposition.
\end{proposition}
\begin{proof}
Choose a nonnegative $\eta\in C_c^\infty(0,1)$ with ordinary $L^2$ norm one.
For $R>{B_0}$, translate it to $(R,R+1)$ and $(R+2,R+3)$, reflect each evenly to the negative half-line, and divide each reflected pair by $\sqrt2$.
Denote the resulting full-line orthonormal functions by $\phi_{1,R},\phi_{2,R}$.
Their derivative norms are independent of $R$.
Put
\[
 m_i=\int_O c\alpha\phi_{i,R}\,dx>0,\qquad
 g_R=\frac{m_2\phi_{1,R}-m_1\phi_{2,R}}{\sqrt{m_1^2+m_2^2}},
 \qquad u_R=\alpha g_R.
\]
Then $\norm{g_R}_2=1$, $m_O(u_R)=0$, and $\wo(u_R)=1$ exactly.
The coefficient vector has Euclidean norm one, so $\norm{g_R'}_2$ is bounded uniformly even when the moments are small.
Every $u_R$ is itself in the centered compact smooth exterior core.

We justify decay of $\alpha'$ from the positive theta series, not from an upper bound for $W$ alone.
On the positive exterior write
\[
 W=\sum_{j\ge1}w_j,\qquad
 w_j=\frac{(4t_j^2-6t_j)e^{-t_j}}{1+e^{-x}},\qquad
 t_j=\pi j^2e^{2x}>3.
\]
Termwise differentiation gives
\[
 (\log w_j)'=2+\frac{4t_j}{2t_j-3}-2t_j+
 \frac{e^{-x}}{1+e^{-x}},\qquad
 |(\log w_j)'|\le2t_j+7.
\]
Cauchy--Schwarz in this positive series yields
\begin{equation}\label{c:eq:alphaprimedecay}
 |\alpha'|^2=\frac{|W'|^2}{4W}
 \le\frac14\sum_{j\ge1}w_j(2t_j+7)^2.
\end{equation}
The polynomial-times-double-exponential series on the right has a summable majorant at any fixed exterior radius and tends uniformly to zero on escaping tails. The same holds for $\alpha$; evenness supplies the negative half-line.
Consequently
\[
 \norm{u_R}_2\le\sup_{|x|>R}\alpha\longrightarrow0,\qquad
 \norm{u_R'}_2\le\sup_{|x|>R}|\alpha'|
 +\sup_{|x|>R}\alpha\,\norm{g_R'}_2\longrightarrow0.
\]
The full nonlocal Gamma estimate
\[
 \tg[u]\le\norm{u'}_2^2\int_0^1\rho(s)s^2\,ds
 +4\norm u_2^2\int_1^\infty\rho(s)\,ds
\]
proves $\tg[u_R]\to0$.

For completeness, the full prime potential also vanishes uniformly on escaping tails.
For $x>{B_0}$, each existing same-side product in \eqref{c:eq:primepotential} is at most $An^{-25/2}$; the reflected product is at most $Ae^{25{B_0}}n^{-25/2}$.
Since $a_n\le1$, a uniform summable row majorant is
\[
 A(2+e^{25{B_0}})\sum_{n\ge2}n^{-25/2}.
\]
For each fixed shift, the row product tends uniformly to zero as $x\to+\infty$; its reflected branch eventually disappears.
A finite-head and summable-tail split, followed by evenness, therefore gives
$\sup_{|x|>R}V_p(x)\to0$.
The complete edge-square inequality gives
\[
 0\le D_p(g_R)\le2\int_OV_p|g_R|^2\,dx
 \le2\sup_{|x|>R}V_p(x)\longrightarrow0.
\]
All prime powers and both physical half-lines are retained.
Together with $\wo(u_R)=1$, this proves \eqref{c:eq:nogap}.
Finally $0\le D_2\le D_p$ proves its dyadic version.
\end{proof}
The proposition is compatible with the trivial kernel of $D_p$ proved in Appendix~\ref{c:sec:signed}: injectivity is not a positive uniform gap.
It rules out $\tg+D_p\ge\varepsilon\wo$ for every fixed $\varepsilon>0$, hence in particular absorption of $\varepsilon_8\wo$ by that route.
It gives no negative direction of $S$ and does not obstruct the relative estimate \eqref{c:eq:spatialtarget}.

\subsubsection{What a regularized sign would and would not imply}
Let $A_2[u]=\tg[u]+D_2(u/\alpha)$ on the same original centered exterior domain.
Under the additional response contracts, \eqref{c:eq:Slower} implies only
\[
 A_2[u]\le S[u]+\varepsilon_8\wo(u).
\]
It does not imply $S[u]+A_2[u]\ge0$.
If that sign were independently proved on the original domain, averaging the two inequalities would give
$S\ge-(\varepsilon_8/2)\wo$ and the coefficient $5741149/11520000$.
That number has an unproved premise and is not an achieved bound.
More generally a separate sign $S+\eta A_2\ge0$ with $\eta>0$ yields only
\[
 S\ge-\frac{\eta\varepsilon_8}{1+\eta}\wo.
\]
One case $\eta=1$ cannot be iterated to zero loss by reusing the same inequalities.
A family with $\eta\downarrow0$ on the same original domain, with all residual, source-tail and domain errors controlled, would prove \eqref{c:eq:remainingS}; no such family is established here.
A weaker $A_2^{1/2}$ completion would require its own faithful-domain and continuous-extension proofs. The full slack is not extended to a weaker topology by assertion.
A rank-seven lift, if used in a future sign analysis, would be $H_7=L_I+7E_7E_7^*$, not the rank-64 lift $H_{64}=L_I+9EE^*$ used here.
Neither the improved fixed comparison nor the conditional response cost supplies a uniform continuation to zero loss.
The global half bound and the Riemann hypothesis remain unproved.

\subsubsection{A positive loss interval for a congruent global pencil}
For $0<\delta<1/2$ define the profile form on $\D$ by
\[
 Q_\delta[f]=\E[f]-(\tfrac12-\delta)\norm f_\mu^2.
\]
If $\E\ge\lambda\Var_\mu$, then on the centered subspace
\[
 Q_\delta[f]\ge(\lambda-\tfrac12+\delta)\norm f_\mu^2,
 \qquad\mu(f)=0.
\]
Thus Theorem~\ref{c:thm:localglobal} excludes centered negativity and nullity whenever
\begin{equation}\label{c:eq:pencilinterval}
 \frac{693}{100000}<\delta<\frac12.
\end{equation}
Under the response contracts, Theorem~\ref{c:thm:conditionalglobal} also gives
$18851/5760000<\delta<1/2$.
The constant direction has value $Q_\delta[1]=-(1/2-\delta)<0$ and is form-orthogonal to the centered subspace. Hence the profile form has negative index one and zero nullity on either applicable strict interval.

The closed form congruence proved in Sections~\ref{a:sec:pencil} and~\ref{f:sec:spectrum} transfers this negative index and nullity to the actual ordinary-$L^2$ pencil $L_\delta$.
This $L_\delta$ is distinct from the fixed local operator $L_I$ inverted above.
The transfer requires the stated congruence and domain identification; it is not inferred from a finite matrix.
These are sufficient positive-loss intervals. Neither endpoint inequality, a single lifted calculation, nor the location of an essential spectral threshold supplies positivity for every $\delta>0$ or a cofinal certified sequence tending to zero.


\section{Complete response gates and full domain continuation}\label{sec:verified-response}
The general Schur machinery is specialized here to finite induced-source coordinates and complete residual enclosures. The formulas below use the actual physical energy already constructed above.

\subsection{Result and scope}
The fourth reported stage uses
\[
a_4=\frac1{2048},\qquad B_4=\frac{\log8}{2},\qquad k_4=256.
\]
In its actual induced-source coordinates, both the historical report and the completed corrected recomputation give the same exact sufficient matrix
\begin{equation}\label{vrg:eq:headline}
M_4=4T_{\lo}-\frac{5a_4}{4}G_{\up}
 +\frac{a_4}{4}F_{\lo}I_{256}
 -\left(\frac4{\nu_4}+\frac{5a_4}{\nu_4^2}\right)K_{\up}
 \succeq\frac1{2000}I_{256}.
\end{equation}
Both this exact penalty and the conservative penalty \(26240\) pass fresh
strict shifted-matrix sign tests after the corrected full source generation. The base analytic composition in
Section~\ref{vrg:sec:fourth} gives, conditionally on the stated full-domain and
enclosure premises,
\begin{equation}\label{vrg:eq:lambda4}
\E[f]\geq\lambda_4\Var_\mu(f),\qquad
\lambda_4=\frac{542006223062224701317434549}{1084868104425248291110985728}
 \approx0.4996056394794406.
\end{equation}
This conclusion concerns every even profile in the full maximal energy
domain, not only the trial space. The simpler \(21487/43008\) also follows. Corollary~\ref{vrg:cor:retained4} strengthens this same fourth-stage comparison to $\lambda_{4,\mathrm{ret}}\approx0.49964024316343157254$ by retaining the original witness margin, without new source enclosures.

The earlier \(a_1=1/288\), \(a_2=1/512\), and \(a_3=1/1024\) stages remain.
The third-stage comparison \(M_3-(3/20000)\mathsf U_3\succ0\), giving
about \(0.499186116859478389\) under its enclosure contract, has a completed
historical separate-environment replay and a fresh exact witness using the
same implementation. Its source contains the integer-power primitive
subsequently found not universally outward. That replay establishes
reproduction, not resolution of this arithmetic obligation.

The fourth-stage producer run, later static review, and exact stored-data
checks are historical layers. A separate recomputation has now completed
the full source, energy, Gram, Gamma, model and residual chain using the
repaired integer-power routine, followed by both fresh gates. A further
exact check reuses only the original witness rows and recomputes the
retained-margin acceptance on the corrected outputs. Intermediate
cell, error-ledger and shifted-pivot enclosures changed; final sufficient
matrices and their exact signs were retained. The repair does not validate
historical earlier stages. Appendix~\ref{app:current-evidence} gives the
per-stage evidence and remaining implementation assumptions. The scoped
full-text analytical review supports the custom source/domain bridge;
neither that review nor the computation is formal verification.


\subsection{The shared form and analytic realization}\label{vrg:sec:objects}
Use the real even subspace of the common $\Hh$ and $X$, and write $H=\Hh$ in this section. Complex inequalities follow by separating real and imaginary parts. Write $\omega_n=\Lambda(n)/\sqrt n$ for the prime weights, to distinguish them from the positive stage parameters $a_1,a_2,a_3,a_4$, and put
\begin{equation}\label{vrg:eq:forms}
 \ell_a[u]=q_0[u]+a\norm u_H^2.
\end{equation}
The kernel, complete forms and maximal energy domain are those of Section~\ref{sec:foundations}, not truncated replacements. The identity $c_\Gamma=\log\pi+\gamma_E+\pi/2+3\log2$ follows from the standard digamma reflection and duplication identities.

\begin{proposition}[Analytic realization used by the response gate]\label{vrg:hyp:realization}
The kernel is positive and even, $\mu(\R)=1$, $\kappa\in X$, and $W(x)\le42e^{4|x|-\pi e^{2|x|}}<2$. The complete $q_0$ is closed and semibounded on $X$. Its self-adjoint representative $B_{\op}$ has the same operator domain as the weighted Gamma representative, with bounded complete prime and constant perturbation, and $B_{\op}\kappa=-\kappa/2$. The smooth form core, minimum--maximal energy-domain equality and $\kappa\D=X$ hold, with
\[
 q_0[\kappa f]=\E[f]-\tfrac12\norm f_\mu^2.
\]
Sharp cuts at every fixed finite radius are continuous in the form norm and belong to the original zero-extended minimum domains.
\end{proposition}
\begin{proof}
The theta normalization, full prime bound and exact domain are proved in Section~\ref{sec:foundations}; the weighted representative and source identity are established in Section~\ref{sec:weighted-physical} and Proposition~\ref{a:prop:thetanull}. Its compact-test source identity extends to $X$ by Lemma~\ref{a:lem:Xcore} and continuity of $q_0$ and the moment; closed-form representation then gives $B_{\op}\kappa=-\kappa/2$. The cut proof in Appendix~\ref{c:app:cuts} uses only a fixed positive finite endpoint, the uniform Carleman crossing bound and strict positivity of the continuous weight on bounded intervals. The same proof applies with that endpoint replaced by each radius used here; its endpoint-strip Gamma error remains $O(\delta(1+|\log\delta|))$. Thus none of these properties is inferred from a finite matrix sign.
\end{proof}
Since $\norm\kappa_H=1$, $h(u)=\rhip u\kappa=2m(u)$. Orthogonal centering is $C=I-\kappa\otimes\kappa$, and $X_c=X\cap\kappa^\perp$. The uncentered bridge uses the norm, not variance; adding $2|m(\kappa f)|^2$, or centering $f$, gives the variance identity. The scoped analytical audit of the numerical stages checks these custom dependencies and their uses, not every theorem of the present monograph.

\subsection{Finite geometry and the changed local bound}\label{vrg:sec:geometry}
For \(I=(-B,B)\), \(O=\R\setminus I\), use the physical even cosine basis
\[
e_0=(2B)^{-1/2},\qquad e_n=B^{-1/2}\cos(n\pi x/B)\quad(1\leq n<k).
\]
Put \(b_n=W1_Ie_n\), \(v_n=Cb_n\),
\(V=\operatorname{span}\{v_n\}_{n<k}\), and \(Z=X_c\cap V^{\perp_H}\).
For centered \(u\), \(\rhip u{v_n}=\int_Iu e_n\).
The representers belong to \(X\); a vanishing combination first forces its
centering term to vanish outside \(I\), then all cosine coefficients to
vanish inside. Hence
\[
X_c=V\mathbin{\oplus_H}Z,\qquad \dim V=k.
\]
The notation \(P_Z\) means projection onto \(\overline{Z}^{H}\), not onto a
possibly nonclosed form domain. There are \(k\) moment constraints after
centering and \(k+1\) constraints on the full space.

The cut-admission argument works at nonrational \(B_3\) as well:
endpoint-strip errors have squared \(L^2\) size \(O(\delta)\),
translation increment \(O(\min(\delta,s))\), and Gamma cost
\(O(\delta(1+|\log\delta|))\). The two folded Carleman branches control the
touching interface. No Neumann boundary condition is imposed; the Neumann
resolvent is only a comparison for a cosine high-mode lower bound.
Write \(q_I[i]=q_0[i]+2|m(i)|^2\) on the local minimum domain.

At \(B=4/5,k=64\), the retained analytic comparisons give
\[
q_I[i]\geq\frac{61}{28}\norm i_2^2
\]
on the first-64-moment kernel, using
\(d_{64}>9\), \(c_\Gamma<383/70\), and complete local prime norm \(<27/20\).
This does not assume positivity on the whole local space.

\paragraph{New radius prime geometry.}
At \(2B_3=\log7\), a compressed shift pair is a direct integral of finite
path adjacency matrices. A path with \(v\) vertices has norm
\(2\cos(\pi/(v+1))\). The \(n=2\) fibers have at most three vertices;
\(n=3,4,5\) have at most two. These norms already include both orientations:
\[
\norm{P_{I_3}}\leq\sqrt2\,\omega_2+\omega_3+\omega_4+\omega_5
 <\frac34+\frac{16}{25}+\frac7{20}+\frac{29}{40}
 =\frac{493}{200}<\frac52.
\]
Here \(\omega_4=\log2/\sqrt4\), \(\omega_6=0\), and \(n=7\) has zero
\emph{local} overlap except null endpoints. Its raw-trial, global,
cross-source, and exterior contributions are not deleted.

\paragraph{New high-mode Gamma floor.}
For \(t=256\pi/B_3>768\), retaining two Gamma terms and integrating the
remaining decreasing mesh gives
\[
d(t)\geq\frac{4t^2}{t^2+1/4}
 +\frac{(4/5)t^2}{t^2+25/4}
 +\frac12\log\left(1+\frac{4t^2}{81}\right).
\]
Since \(\log(512/3)>392/81+2/7\),
\[
h_{256}>\frac{478}{100}+\frac{392}{81}+\frac27
 -\frac{383}{70}-\frac52
 =\frac{54823}{28350}>\frac{19}{10}.
\]
The first surviving mode is 256. The nonnegative mean pole remains in
\(q_I\), although it may be dropped for this lower bound.


\subsection{Full source tails and complement coercivity}\label{vrg:sec:complement}
Assume the stated local floor \(q_I[i]\geq h_k\norm i_2^2\) on the
moment kernel, and complete uncentered bounds
\[
\norm{J_Bo}_2\leq j_B\norm o_H,\quad
|c_\Gamma\norm o_2^2+p[o]|\leq d_B\norm o_H^2,\quad
\mu(O)\leq\bar p_B,
\]
where
\[
J_Bo(y)=\int_O\rho(|x-y|)o(x)\,dx
 +\sum_{n\geq2}\omega_n\{o(y+\log n)+o(y-\log n)\}.
\]
All tails and both orientations are retained. These source/exterior
estimates and local comparison constants are analytical inputs traced to
the audited uniform theta-envelope argument.

\begin{lemma}\label{vrg:lem:complement}
Set \(\eps_{B,k}=d_B+\bar p_B/2+j_B^2/h_k\). For \(a>\eps_{B,k}\),
the full complement representative satisfies
\(A_{a,Z}\succeq(a-\eps_{B,k})I\).
\end{lemma}
\begin{proof}
For \(u=i+o\in Z\), global centering gives \(m(i)=-m(o)\), and
\[
q_0[u]=q_I[i]+q_0[o]-2\Re\langle J_Bo,i\rangle_2-2|m(o)|^2.
\]
Since \(\int_Oc^2W=\mu(O)/4\), the mean costs \(\bar p_B/2\).
Drop nonnegative exterior Gamma energy and complete the interior square.
This gives
\(\ell_a[u]\geq a\norm i_H^2+(a-\eps_{B,k})\norm o_H^2\).
\end{proof}
At \(B=4/5,k=64\), the caps
\(j_B\leq11101/403200\), \(d_B\leq31/46080\),
\(\bar p_B\leq3/49000\), and \(h_k=61/28\) give
\[
\eps_{64}=\frac{2606381167}{2479196160000},\quad
\nu_1=\frac{6001938833}{2479196160000}>\frac{2a_1}{3},\quad
\nu_2=\frac{2235798833}{2479196160000}<\frac{2a_2}{3}.
\]
The old \(2a/3\) weakening is valid only for the first stage.

At \(B_3=\log7/2,k=256\), the uniform \(N=7\) envelopes are
\begin{align*}
M_7&=42\cdot49e^{-7\pi},&
J_7&=2\sqrt{42}\,e^{-7\pi/2}(49+7\pi+7/\pi+1/\pi^2),\\
D_7&=20M_7+96\cdot8^3(27/512)^8,&
P_7&=\frac{168\cdot7^{5/2}e^{-7\pi}}{14\pi-5}.
\end{align*}
Here \(P_7\) bounds full mass, \(J_7\) includes touching Gamma and all prime
source tails, and \(D_7\) is an absolute exterior correction bound.
With \(r=27/512\), \(A=661/9\), safe rational caps are
\[
\mathcal J=168A^2r^7,\qquad
\mathcal D=20(42\cdot49r^7)+96\cdot8^3r^8,\qquad
\mathcal P=\frac{84\cdot343}{37}r^7.
\]
They bound \(j_B^2,d_B,\bar p_B/2\), respectively. Hence
\begin{align*}
\eps_3&=\mathcal D+\mathcal P+\frac{\mathcal J}{19/10}
 =\frac{958695282839627433}{1621007635477226848256},\\
\nu_3&=a_3-\eps_3=\frac{624319986181101911}{1621007635477226848256}
 >\frac{a_3}{3}.
\end{align*}
The coupled reference \(P_a[u]=q_I[i]+t_\Gamma[o]+a\norm u_H^2\) obeys
\(P_{a_3}/5\preceq\ell_{a_3}\preceq9P_{a_3}/5\), not the old half
comparison. Its dual factor would be 5. The certificates here use the
full weighted row route instead.


\subsection{A dimension independent residual gate}\label{vrg:sec:gate}
Let \(F_0=(\rhip{v_j}{v_i})\) and \(Lf=(\int_Ife_i)_{i<k}\).
For a stage's graph-admitted raw columns \(u_j\), set \(f_j=Cu_j\) and
\(J=(Lf_j)_j\). Certified nonsingularity of \(J\) gives
\[
v'_j=P_Vf_j,\qquad v'=VA,\quad A=F_0^{-1}J,\qquad
F=J^{\mathsf T}F_0^{-1}J.
\]
This is a basis of exactly \(V\), with its actual Gram \(F\).
The trial lift \(\widetilde e(v'_j)=f_j\) satisfies
\(P_V\widetilde e=I_V\) exactly. Rounded inverse interpolation and
identity-Gram substitutions are unnecessary and are not used.

For \(A_{a,Z}\succeq\nu I>0\), coercivity gives the harmonic lift
\(e_a(v)\in v+Z\), characterized by \(\ell_a(e_a(v),z)=0\) for all \(z\in Z\).
Define
\[
S_a[v]=\ell_a[e_a(v)],\quad G_a[v]=\norm{e_a(v)}_H^2,\quad
T[v]=\ell_a[\widetilde e(v)],\quad G_{\rm tr}[v]=\norm{\widetilde e(v)}_H^2.
\]
Every matrix uses the same induced column coordinates within its stage.

\begin{proposition}\label{vrg:prop:gate}
Suppose \(T_{\lo}\preceq T\), \(G_{\rm tr}\preceq G_{\up}\),
\(F\succeq F_{\lo}I_k\), and \(K_{\up}\) bounds a full weighted
residual-row Gram representing \(z\mapsto\ell_a(\widetilde e(v),z)\)
on all of \(Z\). Then
\[
4S_a-aG_a\succeq
4T_{\lo}-\frac{5a}{4}G_{\up}+\frac a4F_{\lo}I_k
 -\mathcal C(a,\nu)K_{\up},\qquad
\mathcal C(a,\nu)=\frac4\nu+\frac{5a}{\nu^2}.
\]
\end{proposition}
\begin{proof}
For \(d=\widetilde e-e_a\in Z\), the exact error-energy Gram satisfies
\[
R[v]=\ell_a[d(v)]
 =\sup_{z\in Z\setminus\{0\}}\frac{|\ell_a(\widetilde e(v),z)|^2}{\ell_a[z]},
\qquad R\preceq K_{\up}/\nu,\qquad S_a=T-R.
\]
Also \(\norm{d(v)}_H^2\leq K_{\up}[v]/\nu^2\).
Writing \(\widetilde e(v)=v+\widetilde w(v)\), Young's inequality with
parameter \(1/4\), used only on \(Z\), gives
\[
G_a[v]\leq\tfrac54G_{\rm tr}[v]-\tfrac14F[v]
                 +5K_{\up}[v]/\nu^2.
\]
Combining energy and norm errors gives the stated penalty.
\end{proof}
An arbitrary witness
\((B_{\op}+a)u-\kappa\mu-\sum b_n\lambda_n\) is allowed:
its projection onto \(\overline Z^{H}\) is the actual residual.
Projection contraction suffices; the witness need not itself have zero
moments. Its complete norm must be bounded for every column combination,
including the noncompact mean term, not merely at sampled points or on a
finite Galerkin subspace.

Stage 1 uses the conservative valid \(\nu=2a_1/3\), giving 4968.
Stages 2 and 3 use their exact \(\nu_2,\nu_3\); in particular
\[
\mathcal C_2=\frac{82195505835368325120000}{4998796421644161889},\qquad
\mathcal C_3=
\frac{16878509049120604936352249196300346589184}
 {389775445145171281112984394167851921}.
\]
Neither the old 4968 penalty nor an unavailable reference-dual comparison
is transferred. The margins \(1/250,1/1000,1/2000\) are Euclidean,
basis-dependent floors. They are not \(H\)-metric floors without comparison
to \(F\); congruence preserves positivity, not the same scalar margin.


\subsection{Full norm continuation and the original coefficients}\label{vrg:sec:continuation}
\begin{theorem}\label{vrg:thm:global}
Assume the full realization, finite geometry, and complete complement
inputs above, with \(\nu_*=a-\eps>0\).
If the actual harmonic matrices satisfy \(S_a\succeq\alpha G_a\) for
\(\alpha>0\), then on the full even maximal energy domain
\[
\E[f]\geq\left(\frac12-a+h_*\right)\Var_\mu(f),
\qquad h_*=\frac{\alpha\nu_*}{\alpha+\nu_*}.
\]
\end{theorem}
\begin{proof}
For \(u=e_a(v)+z\in X_c\), harmonic \emph{form} orthogonality gives
\(\ell_a[u]\geq\alpha\norm{e_a(v)}_H^2+\nu_*\norm z_H^2\).
No Hilbert orthogonality is assumed: for arbitrary \(e,z\),
\[
\alpha\norm e_H^2+\nu_*\norm z_H^2-h_*\norm{e+z}_H^2
 =\frac{\norm{\alpha e-\nu_*z}_H^2}{\alpha+\nu_*}\geq0.
\]
Set \(b_*=a-h_*\), \(u_f=\kappa(f-\mu(f))\). The unchanged bridge gives
\[
\norm{u_f}_H^2=\Var_\mu(f),\qquad
q_0[u_f]=\E[f]-\tfrac12\Var_\mu(f)\geq-b_*\Var_\mu(f).
\]
Real and imaginary parts extend the argument to complex even profiles.
\end{proof}

The independent Schur reconciliation also retains the norm cross term.
Write \(A=A_{a,Z}\), \(w_a=P_Ze_a\), \(b=a-h\), \(0<h<\nu_*\). Then
\[
e_b=e_a+h(A-hI)^{-1}w_a,
\]
\begin{align*}
S_b&=S_a-hG_a-h^2\langle w_a,(A-hI)^{-1}w_a\rangle_H\\
&\succeq S_a-\left(h+\frac{h^2}{\nu_*-h}\right)G_a
   +\frac{h^2}{\nu_*-h}F .
\end{align*}
At \(h=h_*\), \(0<h_*<\nu_*\) and
\(h_*+h_*^2/(\nu_*-h_*)=\alpha\), hence
\[
S_{b_*}\succeq\frac{\alpha^2}{\alpha+\nu_*}F\succ0,
\qquad A_{b_*,Z}\succeq(\nu_*-h_*)I\succ0.
\]
There is no endpoint inverse or new solve at \(b_*\).

For each original certificate, \(4S_a-aG_a\succeq0\) gives
\(\alpha=a/4\). Thus the first three coefficients remain
\[
\lambda_1=\frac{1556700237103}{3131143231872},\qquad
\lambda_2=\frac{3517506508493}{7058112169984},\qquad
\lambda_3=\frac{2085655174209946102779}{4178222298875020275712}.
\]
They are approximately \(0.49716672851541954\),
\(0.49836364509080533\), and \(0.4991728598958235\).
Stage 1 also implies the simpler \(175/352\).
For the original third stage,
\[
h_3=\frac{624319986181101911}{4178222298875020275712},\qquad
b_3=\frac{3455975227564035077}{4178222298875020275712},
\]
and \(S_{b_3}\succeq(96619584290816/1020073803436284247)F\).
These statements do not establish infinite continuation, the half bound,
RH, or any unreviewed odd-sector conclusion.


\subsection{A stronger exact comparison of the stored matrices}\label{vrg:sec:generalized}
Use the unchanged third-stage data \(a=a_3,\nu=\nu_3\) and source coordinates.
The response estimate in Proposition~\ref{vrg:prop:gate} yields the \emph{matrix} upper bound
\[
G_a\preceq\mathsf U:=
 \frac54G_{\up}-\frac14F_{\lo}I+\frac5{\nu^2}K_{\up}.
\]
Here \(\mathsf U\) is not the scalar residual-model norm bound \(U\) in
Section~\ref{vrg:sec:numerics}. The negative source term uses the actual
induced-source floor, and \(K_{\up}\) retains every inherited tail and error.

The additional finite certificate proves
\[
M-\theta\mathsf U\succ0,\qquad \theta=\frac3{20000}.
\]
Consequently,
\[
4S_a-aG_a\succeq M\succeq\theta\mathsf U\succeq\theta G_a,
\qquad
S_a\succeq\alpha G_a,\quad
\alpha=\frac{a+\theta}{4}=\frac{721}{2560000}.
\]
No matrices are assumed to commute. This comparison retains their relative
directions rather than dividing separate scalar spectral estimates.
Theorem~\ref{vrg:thm:global} gives
\[
\lambda_+=\frac{1381250815636970598719595}{2767005669802701449752576}
 \approx0.499186116859478389,
\]
with \(S_{a-h}\succeq
(50226821317322080256/422211558502609474144375)F\succ0\).
The exact improvement is
\[
\lambda_+-\lambda_3=
\frac{1169326335435513843338953182503555763}
 {88204687428918903519291432810551217389696}
 \approx1.3256963654883\,10^{-5}>0.
\]
The scalar estimate \(\mathsf U\preceq5944I\) gives a smaller improvement.

\paragraph{Floating proposal and exact acceptance.}
A floating generalized-eigenvalue estimate near \(1.52657938\,10^{-4}\)
suggested the rational \(\theta\); it is not a certified eigenvalue.
After exact assembly \(M-\theta\mathsf U=A_{\rm num}/D\), \(D>0\),
floating Cholesky only proposes \(Q=Q_{\rm num}/2^{40}\).
Acceptance uses arbitrary-precision integers:
\[
H=Q_{\rm num}^{\mathsf T}A_{\rm num}Q_{\rm num}=H^{\mathsf T},
\qquad H_{ii}-\sum_{j\ne i}|H_{ij}|>0\quad(1\leq i\leq256).
\]
Thus \(H\succ0\); this forces the square witness \(Q_{\rm num}\) to be
nonsingular, and proves \(M-\theta\mathsf U\succ0\).
Neither a floating condition estimate nor a sampled eigenvalue is a premise.

A separate standard-library verifier checks input hashes and reconstructs
all 65,536 entries of both the original \(M\) and the generalized matrix
from exact \(T,G,K,F_{\lo}\), then checks the integer witness.
A further local checker reconstructs the common-integer comparison and
all 256 positive congruence margins. These checks validate the stored-matrix
combination, not the source/enclosure generators or full-domain assumptions.

\paragraph{A distinct valid comparison using the original witness.}
For the same stored matrices, an exact check using the original congruence witness also gives
\[
 M-mI-\beta\mathsf U\succ0,\qquad
 m=\frac1{2000},\quad \beta=\frac{87990877}{10^{12}},\quad
 \mathsf U\preceq5944I.
\]
Consequently $S_a\succeq\alpha_{\rm alt}G_a$ with
\[
 \alpha_{\rm alt}=\frac{a+\beta+m/5944}{4}
 =\frac{791025659111}{2972000000000000},
\]
and the same continuation gives
\[
 \lambda_{\rm alt}=
 \frac{1566298472325898306724681382188145}
 {3137737632207571874164167580860416}
 \approx0.499180828966863220.
\]
This is smaller than $\lambda_+$. The two sufficient matrix right-hand sides are not asserted to be Loewner-ordered: their difference has opposite signs on two coordinate vectors in exact stored-matrix checks. The earlier selected-file check of this comparison did not execute its full numerical test suite; that historical limitation does not undo the later complete separate-environment replay of the third stage. The alternate comparison introduces no new source enclosure and does not alter the fixed-complement ceiling below.

\paragraph{Limit of this continuation, not of the operator.}
For fixed \(a,\nu\), the harmonic mean always obeys \(h<\nu\).
Even allowing \(\alpha\) to increase without bound, this method only
approaches
\[
\frac12-a+\nu=\frac12-\eps_3
 =\frac{809545122455773796695}{1621007635477226848256}
 \approx0.4994085806495554.
\]
This is a ceiling for this fixed-complement argument alone. It is not an
upper bound on the optimal operator gap and does not disprove the half bound.


\subsection{The fourth radius and its sharper complement}\label{vrg:sec:fourth}
At \(2B_4=\log8=3\log2\), the compressed \(n=2\) shift pair has
three-vertex fibers almost everywhere, hence norm \(\sqrt2\).
A possible fourth vertex lies on a null endpoint fiber. The \(n=3,4,5,7\)
pairs have norm one, so the full local prime norm is
\[
\norm{P_{I_4}}\leq\sqrt2\,\omega_2+\omega_3+\omega_4+\omega_5+\omega_7
 <\frac{643}{200}<\frac{13}{4}.
\]
The \(n=8\) \emph{local} overlap has measure zero. Its raw-trial,
cross-source and exterior terms remain; no global prime term is deleted.

For \(t=256\pi/B_4>5120/7\), use the two-term-plus-tail Gamma bound
of Section~\ref{vrg:sec:geometry}. The rational terms exceed \(399/100\) and \(79/100\);
the log tail exceeds
\(\log(10240/63)>392/81+34/143\). Thus the independently audited
earlier analytical preflight already gives
\[
h_{256}>
 \frac{478}{100}+\frac{392}{81}+\frac{34}{143}
 -\frac{383}{70}-\frac{13}{4}
 =\frac{9209503}{8108100}>\frac{11}{10}.
\]
Its preflight used the weaker value \(1\); selecting \(11/10\) introduces
no new assumption. The producer's separate shorter bound
\(63521/56700>11/10\) was also checked in the later static comparison.

With \(r=27/512\), the complete \(N=8\) source, exterior, and half-mass caps are
\[
\mathcal J_8=168(5791/63)^2r^8,\qquad
\mathcal D_8=22\cdot42\cdot64r^8+96\cdot9^3r^9,\qquad
\mathcal P_8=\frac{84\cdot8^3}{43}r^8.
\]
Both source orientations and every tail are retained.
Lemma~\ref{vrg:lem:complement} pays the global mean once and gives
\begin{align*}
\eps_4&=\mathcal D_8+\mathcal P_8+\frac{10}{11}\mathcal J_8
 =\frac{20262735541658121715431}{250172086796502324840759296},\\
\nu_4&=a_4-\eps_4
 =\frac{101891603714446529085721}{250172086796502324840759296}
 >\frac{4a_4}{5}.
\end{align*}
Also \(\eta_4=\mathcal D_8+\mathcal P_8+(20/11)\mathcal J_8<a_4/2\).
Conditional on the retained reference-form realization, this gives
\(P_{a_4}/2\preceq\ell_{a_4}\preceq3P_{a_4}/2\) on the \emph{same}
centered moment complement. The numerical certificate uses the full row,
not a reference-dual solve.

The exact residual penalty is
\[
C_4=\frac{254759770327879362689077749075670381093185203994624}
 {10381898907501813723094317124159982832166089841}<26240.
\]
Substitution of \(\nu_0=4a_4/5\) gives \(C(a_4,\nu_0)=26240\).
Since valid \(K_{\up}\) dominates a Gram, it is positive semidefinite;
hence \(M_{\rm exact}-M_{\rm coarse}=(26240-C_4)K_{\up}\succeq0\).
The report separately checks both shifted matrices. Their coordinate
margins are not physical spectral gaps.

Taking \(\alpha=a_4/4\) in Theorem~\ref{vrg:thm:global} gives
\[
h_4=\frac{101891603714446529085721}{1084868104425248291110985728},
\qquad
b_4=\frac{427829150399444238058315}{1084868104425248291110985728},
\]
and \(\lambda_4=1/2-b_4\) in \eqref{vrg:eq:lambda4}.
The coarse floor gives \(b=17/43008\), coefficient \(21487/43008\).
The stronger result exceeds the \emph{third-stage} fixed-floor ceiling
because both the radius and complement changed; that is no contradiction.


\subsection{Retaining the fourth-stage witness margin}\label{vrg:sec:retained4}
The fourth-stage comparison admits a further extraction from the same
stored matrices and the unchanged original congruence witness. It uses no
new source stage. Set $a=a_4$, $\nu=\nu_4$, and retain the matrix upper
bound for the actual harmonic-lift Gram from Proposition~\ref{vrg:prop:gate}:
\[
 0\preceq G_a\preceq\mathsf U
 =\frac54G_{\up}+\frac5{\nu^2}K_{\up}-\frac14F_{\lo}I.
\]
The stored-data checks give the two finite inequalities
\begin{equation}\label{vrg:eq:retained4-gates}
 M_4-mI-\beta\mathsf U\succ0,\qquad \mathsf U\prec gI,
 \qquad m=\frac1{2000},\quad
 \beta=\frac{10508819}{40000000000},\quad g=4060.
\end{equation}
These are statements in the original induced-source coordinates, not
physical spectral-gap claims.

\begin{corollary}[Same-data retained-margin comparison]\label{vrg:cor:retained4}
Under all full-domain, complement and complete-enclosure premises of
Proposition~\ref{vrg:prop:gate} and the fourth-stage construction, the
finite inequalities~\eqref{vrg:eq:retained4-gates} give
\begin{equation}\label{vrg:eq:retained4-alpha}
 S_a\succeq\alpha_{4,\mathrm{ret}}G_a,\qquad
 \alpha_{4,\mathrm{ret}}=\frac{a+\beta+m/g}{4}
 =\frac{6099134007}{32480000000000}.
\end{equation}
Consequently, on the full even maximal energy domain,
\begin{align}\label{vrg:eq:retained4-lambda}
 \E[f]&\ge\lambda_{4,\mathrm{ret}}\Var_\mu(f),\qquad
 \lambda_{4,\mathrm{ret}}=\frac12-a+
       \frac{\alpha_{4,\mathrm{ret}}\nu}{\alpha_{4,\mathrm{ret}}+\nu},\\[-2pt]
 \lambda_{4,\mathrm{ret}}
 &=\frac{2415896663121887366359985794094352519}
         {4835272370827926282795053961934979072}
 \approx0.49964024316343157254.\notag
\end{align}
In particular the rounded-down coefficient $0.499640243163$ is valid
under the same premises and is strictly larger than the base
$\lambda_4$ in~\eqref{vrg:eq:lambda4}.
\end{corollary}
\begin{proof}
The actual response gate gives $4S_a-aG_a\succeq M_4$. Since
$0\preceq G_a\preceq\mathsf U\preceq gI$, the two retained margins yield
\[
 4S_a\succeq aG_a+\beta\mathsf U+mI
       \succeq(a+\beta+m/g)G_a.
\]
This uses order comparisons, not commutation or a substitution of the
trial/source Gram for the actual lift Gram. Apply
Theorem~\ref{vrg:thm:global} with $\alpha=\alpha_{4,\mathrm{ret}}$.
Its harmonic-mean identity pays the complete norm cross term; no Hilbert
orthogonality between the harmonic lift and complement is assumed.
Substitution of the exact fractions gives~\eqref{vrg:eq:retained4-lambda}.
\end{proof}

\paragraph{Exact acceptance and its boundary.}
The original square integer witness $Q$ is not refitted or rescaled.
After clearing positive denominators, the full product
$Q^{\mathsf T}(M_4-mI-\beta\mathsf U)Q$ has 256 strictly positive
symmetric diagonal-dominance margins. Every one of its 65,536 entries
also agrees with the separately assembled difference of the original
congruence and $\beta Q^{\mathsf T}\mathsf UQ$. Strict positivity forces
$Q$ to be nonsingular, proving the first finite inequality. The exact
symmetric row upper bound for $\mathsf U$ is less than
$4059.659<4060$; all 256 row gaps for $4060I-\mathsf U$ are strictly
positive. A separate exact checker reconstructs the complete retained
congruence and the row bound from the original stored entries. Both
checks agree with the scalar fraction above. They verify finite stored
data, not the source-enclosure generator or the analytic premises.

The improvement over the base coefficient is approximately
$3.4603683990985266\,10^{-5}$. The complement did not change, so the same
method ceiling is
\[
 \lambda_{4,\mathrm{ret}}<\frac12-a_4+\nu_4
 =\frac{125065780662709504298664217}{250172086796502324840759296}
 \approx0.49991900481064404098<\frac12.
\]
This is a limitation of the fixed-complement continuation, not an upper
bound for the actual operator. The extraction gives no positive
unshifted $q_0$ margin, no identification with the fixed low space of
Section~\ref{sec:low-sparse}, and no proof of the global half inequality
or the Riemann hypothesis. The separate fourth-stage source-replay
status remains the one stated in Appendix~\ref{app:current-evidence}.


\subsection{New trial coordinates and the exact nonrational interface}\label{vrg:sec:interface}
All four reported candidates are raw even continuous P1 families on
\([-6/5,6/5]\), zero at the support endpoints. The first two have
\(k=64,N=768,h=1/640\) on the positive half-line; the third and fourth have
\(k=256,N=3072,h=1/2560\), with distinct fresh candidates. Continuity and zero extension give \(H^1\),
not \(H^2\), admission. The Gamma multiplier bound, bounded \(W\), and
source eigenline give weighted operator admission.

Every moment matrix uses its true radius. Directed logarithms locate
\(qh<B<(q+1)h\): \(q=2490\) at \(B_3\), while the fourth report gives
\(q=2661\) at \(B_4\).
For frequency \(t_n=n\pi/B\), \(n>0\), the positive-side terminal
hat integrals before cosine normalization are
\[
c_q=\frac{2\cos(qt_nh)-\cos((q-1)t_nh)-(-1)^n}{ht_n^2},
\qquad
c_{q+1}=\frac{(-1)^n-\cos(qt_nh)}{ht_n^2}.
\]
For \(n=0\), with \(\delta=B-qh\),
\[
c_q=h/2+\delta-\delta^2/(2h),\qquad c_{q+1}=\delta^2/(2h).
\]
Both terminal nodes are retained. Interior hats use
\(h\,\operatorname{sinc}(t_nh/2)^2\cos(jt_nh)\), with half for the
central positive-side hat; here \(\operatorname{sinc}(z)=\sin z/z\),
with value 1 at zero. The full-line factor is 2.
The exact identities \(\sin(t_nB)=0\), \(\cos(t_nB)=(-1)^n\) preserve
the true endpoint, rather than classifying it by rounded support.

The fresh inverse witness and interval \(J\) use norm-conversion factor
\(\sqrt{256}=16\), not 8. With \(F_0\preceq2I\), a lower bound for the
smallest singular value of \(J\) gives
\(F=J^{\mathsf T}F_0^{-1}J\succeq\sigma_{\min}(J)^2I/2\).
The third-stage reported \(F_{\lo}\) is approximately \(4.421639094332094\times10^{-9}\),
distinct from both earlier bounds near \(2.21\times10^{-6}\).
The fourth report gives a fresh source floor about \(2.2536263\times10^{-10}\).

The full physical residual is
\[
r=D_\Gamma u-c_\Gamma u-Pu+au/W-\eta c-1_Ie\beta,\qquad
K=\int_\R Wr^{\mathsf T}r,
\]
where \(D_\Gamma u=\int_0^\infty\rho(s)[2u-u(\cdot+s)-u(\cdot-s)]\,ds\)
and \(Pu=\sum_{n\geq2}\omega_n[u(\cdot+\log n)+u(\cdot-\log n)]\).
The free mean-witness coefficient is \(\eta\). Its \(H\)-row is \(Wr\), and \(Wc=\kappa/2\), \(W1_Ie=b\);
thus it is a valid residual witness with the correct units.

\paragraph{Model cut and paid sliver.}
Only the residual integration model uses a 120-bit dyadic \(B_0\);
frequencies and source moments use true \(B\).
The straddling cell carries a separate cosine polynomial, subtracted only
on interior subsegments in the square and both logarithmic cross terms.
The exact \(B_0\) phase is a breakpoint; float conversions have Ball radii.
Since \(W<2\), the full-line cutoff error is at most
\[
\sqrt{4|B-B_0|}\,\norm\beta_F\sqrt{k/B}.
\]
The third-stage enclosure is below \(2.690\times10^{-16}\); the fourth
reports below \(1.480\times10^{-16}\).
It is included once in the row error \emph{before} Gram inflation.
Analytical cut admission alone would not pay this numerical discrepancy.


\subsection{Complete residual enclosures and finite sign checks}\label{vrg:sec:numerics}
For P1 derivative jumps \(J_\ell\),
\[
D_\Gamma u(x)=-\sum_\ell J_\ell\Phi(|x-t_\ell|),\qquad
\Phi(r)=\sum_{j\geq0}\frac{e^{-(2j+1/2)r}}{(2j+1/2)^2}.
\]
This uses the atomic second derivative and includes all Gamma labels.
Each 256-column candidate has 6,145 full-line jumps; both cancellation identities,
all norms, and all model radii were recomputed. Near zero the third-stage scalar
enclosure retains 64 Bernoulli terms with complete tail starting at 65.
Five-prime NTT/CRT at length 16,384 pays 60-bit kernel/jump quantization
and conversion errors. No old jump cap, mode-63 bound, or rank-64 constant
is imported.

Primal \(T_{\lo}\) freshly retains 1024 positive Gamma labels and drops the
positive remainder; this is not exact evaluation of the infinite primal
energy. All compact raw prime correlations through 11, including 7 and 8,
are included. Exact signed integer convolution handles the large contractions.
The residual retains all 13 compact prime powers through 23, both
orientations, H1 displacement errors, and the complete outer tail.

After keeping endpoint logarithms explicit, the degree-32 polynomial
residual is combined before outer products. Positive panel weights dominate
\(W\). A verified 33-node Gauss rule integrates degree-64 polynomial
squares; logarithmic products use separately enclosed analytic moments.
The third stage has 5,120 positive compact cells and 43 exact subsegments
per cell, including the boundary-model breakpoint.

For compact model Gram \(K_0\), midpoint \(\widehat K\), row error \(\eps_r\),
trace-derived \(\norm{K_0}^{1/2}\leq U\), arithmetic radius \(\delta_G\),
and complete outer Gram bound \(\tau_{\rm out}I\),
\[
K_{\rm full}\preceq\widehat K+
 (\delta_G+2U\eps_r+\eps_r^2+\tau_{\rm out})I.
\]
All cross terms and the outer tail are paid once. The table records the
historically reported display bounds; these rounded summaries do not
replace exact endpoints or discharge an arithmetic premise. The corrected
fourth-stage acceptance uses its freshly generated exact inputs, as
recorded in Appendix~\ref{app:current-evidence}.
\begin{center}\small
\begin{tabular}{@{}lllll@{}}
\toprule
Quantity & \(a_1=1/288\) & \(a_2=1/512\) & \(a_3=1/1024\) & \(a_4=1/2048\)\\
\midrule
Model trace & \(<0.001\) & \(<0.0008632\) & \(<0.001986\) & \(<0.001425\)\\
\(U\) & \(4/125\) & \(4/125\) & \(<0.044563\) & \(<0.037737\)\\
\(\eps_r\) & \(<7.862\,10^{-9}\) & \(<7.628\,10^{-9}\) & \(<8.024778\,10^{-9}\) & \(<7.073\,10^{-9}\)\\
\(\delta_G\) & \(<2.082\,10^{-15}\) & \(<2.087\,10^{-15}\) & \(<4.889172\,10^{-15}\) & \(<4.021\,10^{-15}\)\\
\(\tau_{\rm out}\) & \(<1.508\,10^{-67}\) & \(<1.507\,10^{-67}\) & \(<5.110\,10^{-67}\) & \(<4.864\,10^{-67}\)\\
Total debit & \(<5.032\,10^{-10}\) & \(<4.882\,10^{-10}\) & \(<7.152139\,10^{-10}\) & \(<5.338\,10^{-10}\)\\
Matrix floor & \(I_{64}/250\) & \(I_{64}/1000\) & \(I_{256}/2000\) & \(I_{256}/2000\)\\
\bottomrule
\end{tabular}
\end{center}
The fourth report also recomputes its jump, trace, row, and tail bounds.
Its rounded summary components need not add to a separately rounded total;
they cannot certify \(U^2\) against a rounded trace. Exact raw endpoints
are required for such a check. They are available in the mathematical
records used by the subsequent stored-data and corrected-source checks;
their availability is distinct from validation of their enclosure meaning.
The third \(U\) is freshly derived, not a reused \(4/125\).
The historical records report directed interval LDL and exact integer
congruence acceptance for each shifted matrix. The fourth stage has now
passed both kinds of acceptance using corrected arithmetic and fresh data. The latter proves positive diagonal and
strict diagonal dominance of \(Q^{\mathsf T}(M-\delta I)Q\), forcing \(Q\)
nonsingular and the original matrix positive. Floating Cholesky only
proposes \(Q\); no sampled eigenvalue or minimum pivot is a sign premise.


\subsection{Review scope and retained inconclusive attempts}\label{vrg:sec:review}
The first three primary reviews used their producers' executor.
All three earlier stages also have completed historical separate-environment
reviews or replays. These use the same implementation and libraries, not
second complete residual implementations or formal proof checkers. Such
reproduction is not a validation of a shared arithmetic primitive. The
corrected fourth-stage run does not retroactively revalidate earlier runs.

The separate checkpoint 3 audit reran the full residual/sign chain, not
its inherited primal bounds. Stage 2 reran the whole fresh chain, including
uncached primal bounds. Both used independent integer Bareiss checks:
all 64 leading principal minors were positive. Stage 2 used its fresh
jump-norm cap below \(92783.881\), not the old \(92204\).

\paragraph{Completed separate third-stage replay.}
All 31 tests pass. The accepted candidate, true-radius source geometry,
uncached \(T\), \(G\), complete Gamma/model, \(K\), and gate were regenerated.
All compared numerical fields and binary arrays agree; metadata-dependent
paths, times and hash links are checked against their own fresh artifacts.
The auditor separately reconstructed all 65,536 entries of each \(T,G,K,M\):
\(T\) by even-folded Toeplitz sums, \(G\) by a cell formula, and \(K\) by its
full debit ledger. Freshly replayed scalar enclosure inputs are reused;
this is not a second special-function implementation.

Fresh eigenvector-proposed 48-bit dyadic witnesses accept
\(M_3-I/2000\succ0\) and \(M_3-(3/20000)\mathsf U_3\succ0\) by exact
integer diagonal dominance and separate standard-library product checks.
The latter does not reuse the earlier Cholesky witness.
The scalar comparison \(\mathsf U_3\preceq5944I\) also passes.
These checks establish the unchanged stored-matrix signs and scalar
compositions. Their third-stage operator consequences retain the full
enclosure premise, including the unresolved integer-power qualification.

\paragraph{Fourth-stage report and independent static review.}
The fourth producer reports a complete fresh replay in the same executor
and implementation, 31 passing tests, 400-digit shifted LDL, exact integer
congruence, and pure-Python reconstruction for both penalties.
No separate reviewer was commissioned during that producer turn.
A later independent analytical/static audit reconciles the full received
proof, 32 exact rational values, 22 strict margins and ten ledger quantities.
That static review did not receive raw matrices, regenerate enclosures or certify their signs. The later exact stored-matrix check described in Appendix~\ref{app:current-evidence} is a distinct verification layer. The producer's missing earlier-preflight file is a historical
provenance issue: separate local analytical and received derivations both justify
the selected Gamma floor, so it is not a remaining analytic mismatch.

\paragraph{Failures are preserved, not reclassified.}
For stage 3, the initial 1536-cell diagnostic had a negative sufficient
estimate, not a certified negative direction of the harmonic gate.
The accepted 3072-cell candidate is new. Overwide trigonometric and
48-term scalar enclosures were replaced without changing that frozen candidate.
The 70-digit LDL attempt was inconclusive at pivot index 61, the 62nd pivot.
The separate audit reproduced it. At 400 digits, all 256 shifted pivots
pass for the same exact matrix and gap; no margin or candidate changed.
The fourth producer report recorded no numerical failure at its original
snapshot, before the later integer-power finding and corrected recomputation.

High-precision moment, full-jump Gamma and boundary/log integration probes
are diagnostic. Agreement of point values or quadratures is not an error
bound. Exact enclosures, full tail/error budgets and integer sign acceptance
are the proof components. A pivot magnitude is not a physical eigenvalue bound.

\paragraph{Integer-power correction and analytic trust.}\label{vrg:par:arithmetic}
An exact-rational audit found a genuine defect in the original interval
integer-power routine on the observed runtime. At working precision $70$,
for the exact finite-decimal input $x=1+10^{-40}$, its returned upper
endpoint for $x^2$ is $1+2\cdot10^{-40}$, whereas
\[
 x^2=1+2\cdot10^{-40}+10^{-80}.
\]
Thus the endpoint is strictly inward. This is a synthetic primitive
counterexample, not an identified final-operator counterexample or a
proof that a particular historical production operand failed. The
historical third-stage source uses the same endpoint-power mechanism;
its one-sided enclosure obligation remains unresolved.

The corrected routine uses only directed nonnegative multiplications
and repeated squaring. Each product is monotone in its nonnegative
factors, so induction preserves the appropriate lower or upper bound.
For a negative endpoint and odd exponent, exact sign reversal exchanges
the two rounding directions. An even power on a zero-crossing interval
has lower bound zero and upper bound from the largest endpoint magnitude;
negative integer powers first use the zero-excluding reciprocal.
Invalid or nonfinite results fail closed. This operation-level argument
is distinct from the adversarial exact-rational regression tests.

The complete fourth-stage chain was regenerated using this repair.
Genuine intermediate cell and error-ledger bounds, and shifted-pivot
enclosures, changed outwards. Final sufficient matrices and exact
integer acceptance were retained, and the retained-margin inequality
was recomputed against the corrected outputs. Thus equality of final
matrix entries is not a claim that the old arithmetic or every old
intermediate enclosure was valid. The remaining numerical trust includes
Python integers/Fraction, directed Decimal multiplication and the other
stated elementary enclosures, IEEE binary64/subnormals/nextafter, and
single-thread full-precision BLAS with Ball errors. Neither exact finite
acceptance nor this repair formally verifies the entire runtime or
implementation. The scoped analytical review checks the custom bridge
and standard-theorem applications, not all foundational or RH-equivalence
claims. The independent Arb/Acb low-space calculations below and the
separate local-input history do not inherit this specific Decimal defect.



\section{Delta scaled finite Schur reductions}\label{sec:schur-unified}
Under their respective local-sign or complete-enclosure premises, the
preceding variance bounds give actual signs on specified parameter intervals. At every positive parameter one can also construct a positive infinite complement and exact scaled response, but the finite sign is then a separate question. The radical interpolation cost must remain in the reference form.


\begin{theorem}[Scaled complement and exact Schur response]\label{a:thm:mainschur}
For every $\delta>0$ and $0<\theta<1$, finitely many explicitly prescribed interior moments define a weighted orthogonal decomposition $X_c=V\oplus Z$ such that
\begin{equation}\label{a:eq:maincomparison}
 (1-\theta)P_\delta[z]\le\ell_\delta[z]
 \le(1+\theta)P_\delta[z],\qquad
 P_\delta[z]=\tg[z]+\delta\norm{z}_{\Hh}^2.
\end{equation}
There is a finite-dimensional radical lift $H:V\to X_c$ satisfying $P_VHv=v$ and $q_0(Hv,z)=0$ for all $z\in X$. Write
\[
 h(v)=Hv,\quad p(v)=h(v)-v\in Z,\quad G[v]=\norm{h(v)}_{\Hh}^2.
\]
If $A_{\delta,Z}$ is the positive weighted operator represented by $\ell_\delta|_Z$, then the exact Schur form is
\begin{equation}\label{a:eq:mainschur}
 S_\delta[v]=\delta G[v]
    -\delta^2\ip{p(v)}{A_{\delta,Z}^{-1}p(v)}_{\Hh}.
\end{equation}
Finite Neumann sums $z_J$ of exact positive reference solves satisfy
\begin{equation}\label{a:eq:mainerror}
 0\le\ell_\delta[h(v)-z_J(v)]-S_\delta[v]
 \le\frac{\delta\theta^{2J+2}}{1-\theta}G[v].
\end{equation}
They also have a weighted residual bound $C_0\theta^J\sqrt{G[v]}$ with $C_0<10$. The norm of $H$ is charged through $G$; no uniform conditioning or weighted completeness is assumed.
\end{theorem}

The proof of Theorem~\ref{a:thm:mainseed} is in Section~\ref{a:sec:seed}. Sections~\ref{a:sec:radical}--\ref{a:sec:response} prove Theorem~\ref{a:thm:mainschur}. On the interval in Theorem~\ref{a:thm:mainseed}, the exact finite form also satisfies
$S_\delta[v]\ge(c_{\log 2}+\delta-1/2)\norm v_{\Hh}^2$.
This finite sign follows from the full-domain variance estimate. It is not a uniform continuation theorem as $\delta$ tends to zero.


\subsection{Finite moment complements for positive parameters}\label{a:sec:moments}

The full correction $d=\cg\ip\cdot\cdot_2+p$ is bounded in the weighted norm, but need not be small globally. Interior moments make its ordinary-$L^2$ part small relative to the Gamma form, while the theta weight controls spatial leakage. The construction below quantifies both effects.

Write
\[
 d(T)=\frac12\log(1+4T^2),\qquad
 \epsilon_m=\frac{2RT}{\pi}\frac{(RT)^{4m}}{((2m)!)^2}.
\]
The scalar function $d(T)$ is distinct from the bilinear correction $d(u,z)$.

\begin{theorem}[Explicit relative complement]\label{a:thm:moment}
Fix $\delta>0$ and $0<\theta<1$. Choose $N\ge2$, $C>0$, $T>0$, an integer $R\ge1$, and an integer $m\ge1$ so that
\begin{align}
 r_N&\le\theta\delta/2,&
 C&\ge\cg+2\sum_{2\le n\le N}a_n,\label{a:eq:choices1}\\
 d(T)&\ge2C/\theta,&
 \Omega_R&\le\theta\delta/(8C),&
 \epsilon_m&\le1/4.\label{a:eq:choices2}
\end{align}
If $u\in X$ and
\begin{equation}\label{a:eq:moments}
 \int_{-R}^Rx^{2k}u(x)\,dx=0\qquad(0\le k<m),
\end{equation}
then
\begin{align}
 \norm u_2^2&\le\frac2{d(T)}\tg[u]+4\Omega_R\norm u_{\Hh}^2,
                   \label{a:eq:Fouriermoment}\\
 |d[u]|&\le\theta P_\delta[u],\qquad
 (1-\theta)P_\delta[u]\le\ell_\delta[u]
                \le(1+\theta)P_\delta[u],\label{a:eq:relative}
\end{align}
where $P_\delta[u]=\tg[u]+\delta\norm u_{\Hh}^2$.
\end{theorem}
\begin{proof}
Put $I=[-R,R]$, $u_i=\one_Iu$, and $u_o=\one_{I^c}u$. Evenness and~\eqref{a:eq:moments} remove the first $m$ terms of the cosine Taylor expansion. For real arguments the remainder is bounded by $|\xi x|^{2m}/(2m)!$. Hence the unitary transform satisfies, for $|\xi|\le T$,
\begin{align*}
 |\widehat{u_i}(\xi)|
 &\le(2\pi)^{-1/2}\int_I|u(x)|\frac{|\xi x|^{2m}}{(2m)!}\,dx\\
 &\le\sqrt{R/\pi}\,\norm{u_i}_2\frac{(RT)^{2m}}{(2m)!}.
\end{align*}
If $F_T$ denotes Fourier projection onto $[-T,T]$, then
$\norm{F_Tu_i}_2^2\le\epsilon_m\norm{u_i}_2^2$.
Only ordinary $L^2$ membership of the cut piece is used; no Gamma-domain assertion about $u_i$ is required. The high-frequency part of the original $u$ is controlled by Lemma~\ref{a:lem:fourier}. Thus
\begin{align*}
 \norm u_2^2
 &\le2\epsilon_m\norm{u_i}_2^2+2\norm{u_o}_2^2
                     +\tg[u]/d(T)\\
 &\le2\epsilon_m\norm u_2^2+2\Omega_R\norm u_{\Hh}^2
                     +\tg[u]/d(T).
\end{align*}
Absorbing $2\epsilon_m\le1/2$ proves~\eqref{a:eq:Fouriermoment}.
The complete head--tail estimate then gives
\begin{align*}
 |d[u]|&\le C\norm u_2^2+r_N\norm u_{\Hh}^2\\
 &\le\frac{2C}{d(T)}\tg[u]+(4C\Omega_R+r_N)\norm u_{\Hh}^2
 \le\theta P_\delta[u].
\end{align*}
Since $\ell_\delta=P_\delta-d$, the two-sided estimate follows.
\end{proof}

\begin{corollary}[A terminating explicit prescription]\label{a:cor:choices}
The choices in Theorem~\ref{a:thm:moment} can be made without a spectral sign assumption. One sufficient prescription is:
\begin{enumerate}
\item Choose an integer $N\ge2$ with $r_N\le\theta\delta/2$.
\item Put $C=2N+4$, $L=\lceil2C/\theta\rceil$, and $T=(11/4)^L/2$.
\item Choose an integer $R\ge1$ with
$42\,81^R(3/8)^{3\cdot7^R}\le\theta\delta/(8C)$.
\item Put $m=\lceil2(RT)^2\rceil$.
\end{enumerate}
For rational $\delta$ and $\theta$, all sufficient tests in this prescription use integer and rational operations.
\end{corollary}
\begin{proof}
The bound $a_n\le\log n/\sqrt n\le2/e<1$ follows by maximizing $\log x/\sqrt x$ for $x\ge1$. Together with $\cg<6$ this gives the stated $C$. Since $\log(11/4)>1$ by~\eqref{a:eq:elementarye},
$d(T)>\log(2T)=L\log(11/4)>L\ge2C/\theta$.
The remainder $r_N$ decreases to zero geometrically: its successive ratio for $N\ge2$ is at most $1/8$. The rational envelope in~\eqref{a:eq:rationalOmega} tends to zero, so the choices of $N$ and $R$ terminate.

Write $A_0=RT\ge1$. Since $(2m)!\ge m^m$, $m\ge2A_0^2$, and $4^m\ge4m$,
\[
 \epsilon_m<\frac{2A_0}{3}\left(\frac{A_0^2}{m}\right)^{2m}
 \le\frac{2A_0}{3}4^{-m}
 \le\frac{A_0}{6m}\le\frac1{12A_0}<\frac14.
\]
This verifies the final test.
\end{proof}

The prescription is deliberately coarse. It proves finite existence with explicit sufficient inequalities, not a practically small matrix dimension. No parameter selection or finite matrix is evaluated here.

\subsubsection{Exact weighted moment representatives}

With $I=[-R,R]$ and the choices above, define
\begin{equation}\label{a:eq:momentbasis}
 b_k=W\one_Ix^{2k},\qquad
 v_k=\mathsf Cb_k=b_k-2m(b_k)\kappa\quad(0\le k<m),
\end{equation}
and let
\begin{equation}\label{a:eq:VZ}
 V=\operatorname{span}\{v_0,\ldots,v_{m-1}\},\qquad
 Z=\{u\in X_c:\text{\eqref{a:eq:moments} holds}\}.
\end{equation}

\begin{lemma}[Admitted basis and exact decomposition]\label{a:lem:VZ}
Each $v_k$ belongs to $X_c$, the family is independent, and
\begin{equation}\label{a:eq:orthdecomp}
 X_c=V\oplus Z,\qquad V\perp_{\Hh} Z.
\end{equation}
The subspace $Z$ is form closed and has codimension $m$ in $X_c$. Its weighted closure is
$\Zc=\Hh_c\cap V^{\perp_{\Hh}}$, and $X\cap\Zc=Z$.
\end{lemma}
\begin{proof}
Each $b_k$ is bounded, compactly supported, and piecewise $C^1$, with at most two boundary jumps. For $0<s\le1$ the squared translation increment is $O(s)$: the boundary strips have length $O(s)$ and bounded values, while the smooth common interior contributes $O(s^2)$. Since $\rho(s)=O(1/s)$ at zero and is integrable away from zero, $\tg[b_k]<\infty$. Its weighted norm is finite on the compact interval. Since $\kappa\in X$, centering gives $v_k\in X_c$.

For $u\in X_c$,
\[
 \ip u{v_k}_{\Hh}=\ip u{b_k}_{\Hh}
                    =\int_Ix^{2k}u(x)\,dx.
\]
If $\sum\alpha_kv_k=0$, then $W\one_IP=a\kappa$ for an even polynomial $P$. Outside $I$, positivity of $\kappa$ forces $a=0$. Inside $I$, positivity of $W$ gives $P=0$, so all coefficients vanish. Thus the weighted Gram matrix is positive definite. Projection onto its finite span preserves $X_c$ and proves~\eqref{a:eq:orthdecomp}. The moment functionals are weighted-norm continuous, so their common kernel is form closed and has the asserted codimension.

Finally $X_c$ is dense in $\Hh_c$: apply $\mathsf C$ to a weighted-norm dense sequence in $X$. Applying $I-P_V$ to this approximation proves density of $Z$ in $\Zc$, and the characterization $X\cap\Zc=Z$ follows from the moment identities.
\end{proof}

In particular the infinite complement has genuine full-form coercivity,
\begin{equation}\label{a:eq:Zcoercive}
 \ell_\delta[z]\ge(1-\theta)P_\delta[z]
 \ge(1-\theta)\min\{1,\delta\}\norm z_X^2,
\end{equation}
not merely positivity on a finite sample. Its weighted gap is at least $(1-\theta)\delta$.

\subsection{Exact finite Schur reduction and radical lifts}\label{a:sec:schur}

For $v\in V$, define the usual exact response $z_v\in Z$ by
\begin{equation}\label{a:eq:usualresponse}
 \ell_\delta(z_v,z)=\ell_\delta(v,z)\quad(z\in Z),\qquad
 S_\delta[v]=\ell_\delta[v-z_v].
\end{equation}
Lax--Milgram in the full form norm gives its existence and uniqueness by~\eqref{a:eq:Zcoercive}. Completing the square gives
\begin{equation}\label{a:eq:usualsquare}
 \ell_\delta[v+z]=S_\delta[v]+\ell_\delta[z+z_v]
                         \qquad(v\in V,\ z\in Z).
\end{equation}
Consequently the negative index and bilinear nullity of $\ell_\delta|_{X_c}$ equal those of the finite form $S_\delta$. Together with~\eqref{a:eq:meansplit}, the full form is congruent to the mean scalar, the finite Schur form, and the positive infinite complement. For $0<\delta<1/2$, positivity of $S_\delta$ is exactly the remaining finite condition for negative index one and nullity zero.

The bare finite block has the useful identity
\begin{equation}\label{a:eq:bareblock}
 \ell_\delta(v_i,v_j)=\ell_\delta(b_i,b_j)
       +(2-4\delta)m(b_i)\overline{m(b_j)}.
\end{equation}
It follows by expanding the weighted centering and using
$q_0(\kappa,z)=-\overline{m(z)}$ and
$\ip\kappa z_{\Hh}=2\overline{m(z)}$.
For compact $b_i,b_j$, prime shifts with $\log n>2R$ have disjoint supports and contribute exactly zero to this bare block. This support fact does not truncate the full noncompact response; all its tails remain controlled by Lemma~\ref{a:lem:prime}.

\begin{lemma}[Finite radical interpolation]\label{a:lem:lift}
There are a finite-dimensional subspace $F\subset\Rcal$ and a linear map $H:V\to F$ such that $P_VHv=v$. Writing
\begin{equation}\label{a:eq:liftdata}
 h(v)=Hv=v+p(v),\qquad p(v)=P_{\Zc}h(v)\in Z,
\end{equation}
the finite form
\begin{equation}\label{a:eq:G}
 G[v]=\norm{h(v)}_{\Hh}^2
     =\norm v_{\Hh}^2+\norm{p(v)}_{\Hh}^2
\end{equation}
is positive definite.
\end{lemma}
\begin{proof}
If a nonzero linear combination of the $m$ interior moments annihilated $\Rcal$, its ordinary-$L^2$ representative
$\one_I\sum_{k<m}\alpha_kx^{2k}$ would be orthogonal to the ordinary-dense space $\Rcal$. Proposition~\ref{a:prop:density} would force that representative to vanish, and then all coefficients would be zero. Therefore the moment map $\Rcal\to\C^m$ is onto. Its image already has full rank on a finite span $F$ of derivative radicals. On $V$, the same moment map is the invertible weighted Gram matrix. A linear right inverse on $F$ followed by this matrix defines $H$. The moment identities give $P_VHv=v$ and $p(v)\in Z$. Orthogonality gives~\eqref{a:eq:G}; $P_VHv=v$ makes $H$ injective, so $G$ is positive definite.
\end{proof}

Only finite interpolation has been proved. There is no uniform bound on $\norm H$ as $\delta$, $\theta$, or the moment choices change. Nor has a particular full-rank derivative minor been located and enclosed. An implementation must find and certify that minor and charge its conditioning. The reference form $G$ explicitly retains this cost, rather than replacing it silently by $\norm v_{\Hh}^2$.

\subsection{Delta-scaled response and uniform preconditioning}\label{a:sec:response}

All operators in this section act in the weighted complement Hilbert space $\Zc$, not in ordinary $L^2$. Let $T_Z$ be the nonnegative self-adjoint representative of the Gamma form restricted to $Z$, whose form domain is $Z$. Define
\begin{equation}\label{a:eq:weightedops}
 P_{\delta,Z}=T_Z+\delta I,\qquad
 D_Z=P_{\Zc}D|_{\Zc},\qquad
 A_{\delta,Z}=P_{\delta,Z}-D_Z.
\end{equation}
The Gamma form on $\Zc$ is densely defined and closed: its norm after adding the weighted squared norm is precisely the restricted $X$ norm, and Lemma~\ref{a:lem:VZ} supplies density. Since $D_Z$ is bounded with norm at most $C_0$, the bounded-perturbation theorem gives
\begin{equation}\label{a:eq:opdomains}
 \dom A_{\delta,Z}=\dom P_{\delta,Z}=\dom T_Z.
\end{equation}
The relative complement estimate gives
$A_{\delta,Z}\ge(1-\theta)\delta I$ and hence
$\norm{A_{\delta,Z}^{-1}}\le[(1-\theta)\delta]^{-1}$.
These statements concern a weighted constrained operator. They construct no ordinary-$L^2$ endpoint operator at $\delta=0$.

\begin{theorem}[Exact scaled Schur formula]\label{a:thm:scaled}
For the lift in Lemma~\ref{a:lem:lift}, put
\begin{equation}\label{a:eq:liftresponse}
 z_*(v)=\delta A_{\delta,Z}^{-1}p(v),\qquad
 e(v)=h(v)-z_*(v).
\end{equation}
Then
\begin{equation}\label{a:eq:scaledschur}
 S_\delta[v]=\ell_\delta[e(v)]
 =\delta G[v]-\delta^2\ip{p(v)}{A_{\delta,Z}^{-1}p(v)}_{\Hh}.
\end{equation}
The usual response in~\eqref{a:eq:usualresponse} is $z_v=z_*(v)-p(v)$.
\end{theorem}
\begin{proof}
The full-domain radical identity gives, for every $z\in Z$,
\[
 \ell_\delta(h(v),z)=\delta\ip{h(v)}z_{\Hh}
                    =\delta\ip{p(v)}z_{\Hh}.
\]
Thus $z_*$ is exactly the response from the affine representative $h(v)$. The affine sets $h(v)+Z$ and $v+Z$ coincide. Moreover $q_0[h(v)]=0$, so $\ell_\delta[h(v)]=\delta G[v]$. Completing the square gives~\eqref{a:eq:scaledschur}. Substituting $h(v)=v+p(v)$ into the response equation gives the last assertion.
\end{proof}

The factor $\delta$ in the forcing is exact. It does not arise by subtracting two separately uncontrolled terms of order one. Nevertheless the second term in~\eqref{a:eq:scaledschur} can be comparable to or larger than the first, and its sign consequence needs a separate estimate.

Regard $Z$ as the energy Hilbert space with inner product $P_\delta$. Set
\begin{equation}\label{a:eq:Kx}
 K_\delta=P_{\delta,Z}^{-1}D_Z,\qquad
 x(v)=\delta P_{\delta,Z}^{-1}p(v).
\end{equation}
The relative bound~\eqref{a:eq:relative} implies that $K_\delta$ is self-adjoint in this energy inner product and
\begin{equation}\label{a:eq:contraction}
 \norm{K_\delta}_{P_\delta\to P_\delta}\le\theta,
 \qquad \norm{x(v)}_{P_\delta}^2\le\delta\norm{p(v)}_{\Hh}^2.
\end{equation}
Indeed $P_\delta(K_\delta u,z)=d(u,z)$, so the quadratic bound yields the operator bound by the spectral theorem for bounded self-adjoint operators. Testing the reference equation for $x$ against $x$, and using
$\norm x_{\Hh}\le\delta^{-1/2}\norm x_{P_\delta}$, proves the second estimate.

\begin{theorem}[Exact-reference response series and charged errors]\label{a:thm:neumann}
For each integer $J\ge0$, define
\begin{equation}\label{a:eq:iterates}
 z_J(v)=\sum_{k=0}^JK_\delta^kx(v),\qquad
 e_J(v)=h(v)-z_J(v),\qquad T_J[v]=\ell_\delta[e_J(v)].
\end{equation}
Then $z_J\to z_*$ in the full form norm and
\begin{equation}\label{a:eq:Schurerror}
 0\le T_J[v]-S_\delta[v]
 \le\frac{\delta\theta^{2J+2}}{1-\theta}\norm{p(v)}_{\Hh}^2
 \le\frac{\delta\theta^{2J+2}}{1-\theta}G[v].
\end{equation}
For the residual
\begin{equation}\label{a:eq:residual}
 r_J(v;z)=\delta\ip{p(v)}z_{\Hh}-\ell_\delta(z_J(v),z),
 \qquad z\in Z,
\end{equation}
one has
\begin{align}
 r_J(v;z)&=P_\delta(K_\delta^{J+1}x(v),z),\label{a:eq:energyrow}\\
 \norm{r_J(v;\cdot)}_{P_\delta^*}
 &\le\sqrt\delta\,\theta^{J+1}\norm{p(v)}_{\Hh},\label{a:eq:energyresidual}\\
 r_J(v;z)&=\ip{D_ZK_\delta^Jx(v)}z_{\Hh},\label{a:eq:weightedrow}\\
 |r_J(v;z)|&\le C_0\theta^J\sqrt{G[v]}\norm z_{\Hh}.\label{a:eq:weightedresidual}
\end{align}
For the usual response, the corresponding approximate vector is
$\widetilde z_v=z_J(v)-p(v)$; its residual is exactly~\eqref{a:eq:residual}.
\end{theorem}
\begin{proof}
The response equation is $(I-K_\delta)z_*=x$. Since $\norm K_\delta\le\theta<1$ in the energy norm, the Neumann series converges to $z_*$ in that norm, which is equivalent to the $X$ norm for fixed $\delta>0$. Telescoping gives~\eqref{a:eq:energyrow} and~\eqref{a:eq:energyresidual}.

Put $r=K_\delta^{J+1}x$. Then
$z_*-z_J=(I-K_\delta)^{-1}r$. Exact completion of the square gives
\begin{align*}
 T_J-S_\delta
 &=\ell_\delta[z_*-z_J]
 =P_\delta(r,(I-K_\delta)^{-1}r)\\
 &\le\frac1{1-\theta}\norm r_{P_\delta}^2.
\end{align*}
Use~\eqref{a:eq:contraction} and~\eqref{a:eq:G} to obtain~\eqref{a:eq:Schurerror}.

For the weighted residual, a domain argument is essential. Every application of $P_{\delta,Z}^{-1}$ maps $\Zc$ into $\dom T_Z$, and $D_Z$ is bounded on $\Zc$. Hence $x$ and all finite terms $K_\delta^kx$ belong to the common operator domain in~\eqref{a:eq:opdomains}. The geometric identity now holds as an actual weighted operator equation:
\[
 A_{\delta,Z}z_J=\delta p-D_ZK_\delta^Jx.
\]
This proves~\eqref{a:eq:weightedrow}. Finally,
\[
 \norm{D_ZK_\delta^Jx}_{\Hh}
 \le C_0\delta^{-1/2}\norm{K_\delta^Jx}_{P_\delta}
 \le C_0\theta^J\norm p_{\Hh}
 \le C_0\theta^J\sqrt{G[v]}.
\]
For $\widetilde z_v=z_J-p$, the radical identity gives
\[
 \ell_\delta(v,z)-\ell_\delta(\widetilde z_v,z)
 =\delta\ip pz_{\Hh}-\ell_\delta(z_J,z)=r_J(v;z),
\]
which verifies the stated relation to the usual response.
\end{proof}

This weighted residual estimate is not inferred from mere form membership of an arbitrary trial vector. It uses the bounded correction $D_Z$ and equality of the two operator domains. No form-core-to-operator-core upgrade is made.

\subsubsection{Residual budgets and the implementation boundary}

For $\theta=1/2$, the energy error is
\begin{equation}\label{a:eq:halferror}
 0\le T_J-S_\delta\le2\delta\,4^{-J-1}G.
\end{equation}
For any $\eta>0$, choosing $J$ so that
$\theta^{2J+2}/(1-\theta)\le\eta$ makes the error at most $\eta\delta G$, uniformly in $\delta$. The weighted-dual estimate has a different parameter budget. If
\begin{equation}\label{a:eq:weightedbudget}
 C_0^2\theta^{2J}\le(1-\theta)\eta\delta^2,
\end{equation}
then its residual constant $\tau=C_0\theta^J$ obeys
$\tau^2\le(1-\theta)\eta\delta^2$.
To see the relevance, a residual satisfying
$|r(v;z)|\le\tau\sqrt{G[v]}\norm z_{\Hh}$ and a complement gap
$\ell_\delta[z]\ge b\delta\norm z_{\Hh}^2$ gives Schur error at most
$\tau^2G/(b\delta)$. This follows by testing the exact response error equation against the error and applying Cauchy--Schwarz. With $b=1-\theta$, condition~\eqref{a:eq:weightedbudget} gives the same target $\eta\delta G$.

These formulas are analytic constructions involving finitely many \emph{exact infinite-dimensional} positive reference solves $P_{\delta,Z}^{-1}$. They are not computed finite-matrix certificates. Replacing a reference solve by a Galerkin vector requires additional certified residual or enclosure estimates. Form density does not itself provide those estimates. Neither an evaluated reference inverse nor a quantitative enclosure of the interpolation matrix is supplied here.

\subsection{Fixed-interval sign and signed continuation}\label{a:sec:continuation}

\subsubsection{The actual finite sign supplied by the variance bound}

\begin{theorem}[Finite Schur sign on the proved interval]\label{a:thm:finiteseed}
Fix $1/2-c_{\log2}<\delta<1/2$ and any one of the moment decompositions above. With $g_\delta=c_{\log2}+\delta-1/2$,
\begin{align}
 S_\delta[v]&\ge g_\delta\norm v_{\Hh}^2,\label{a:eq:Sseed}\\
 \delta^2\ip p{A_{\delta,Z}^{-1}p}_{\Hh}
 &\le\delta G[v]-g_\delta\norm v_{\Hh}^2.\label{a:eq:signedresponse}
\end{align}
In particular $S_{23/48}\ge(509/2550000)I_V$ in the inherited weighted metric. For the same fixed lift, put
\begin{equation}\label{a:eq:sigmaH}
 \sigma_H=\inf_{v\ne0}\frac{\norm v_{\Hh}^2}{G[v]}>0.
\end{equation}
Every exact-reference iterate satisfies
\begin{equation}\label{a:eq:Tseed}
 T_J[v]\ge g_\delta\norm v_{\Hh}^2\ge g_\delta\sigma_HG[v].
\end{equation}
\end{theorem}
\begin{proof}
Since $z_*\in Z$, weighted projection gives $P_Ve(v)=v$. Thus
$\norm{e(v)}_{\Hh}^2\ge\norm v_{\Hh}^2$. Apply the full-domain bound~\eqref{a:eq:seedcoercivity} to $e(v)$ and use~\eqref{a:eq:scaledschur}. The same projection argument applies to $e_J(v)=h(v)-z_J(v)$. Finally $\sigma_H>0$ follows from finite dimensionality and injectivity of $H$.
\end{proof}

This is a sign theorem for the actual finite Schur form with its full infinite response retained. Its source is the global energy bound, not residual smallness. It gives no lower bound for $\sigma_H$ uniform over changing moment spaces or parameters.
For this fixed $V,Z,H,\delta$, choose
$\eta=g_\delta\sigma_H/(2\delta)$. Then~\eqref{a:eq:Tseed} is
$T_J\ge2\eta\delta G$. A finite $J$ meets both
\begin{equation}\label{a:eq:seedbudgets}
 \frac{\theta^{2J+2}}{1-\theta}\le\frac{g_\delta\sigma_H}{2\delta},
 \qquad C_0^2\theta^{2J}\le(1-\theta)\eta\delta^2.
\end{equation}
This existence statement does not compute $\sigma_H$, locate the interpolation minor, or turn the exact positive solves into numerical matrices.

\begin{corollary}[The stronger comparison branches in the same Schur coordinates]\label{cor:stronger-schur}
Suppose a branch of Section~\ref{sec:bounds} proves $\E\ge c_*\Var_\mu$. For every $\delta>1/2-c_*$ and every moment decomposition at that parameter,
\[
 S_\delta[v]\ge(c_*+\delta-1/2)\norm v_{\Hh}^2.
\]
If also $\delta<1/2$, the ordinary-$L^2$ pencil has negative index one and zero nullity. In particular the local-sign input gives this conclusion for $693/100000<\delta<1/2$, and the additional response contracts give it for $18851/5760000<\delta<1/2$.
\end{corollary}
\begin{proof}
The bridge gives $\ell_\delta[u]\ge(c_*+\delta-1/2)\norm u_{\Hh}^2$ on $X_c$. The exact response vector has projection $v$, so its weighted norm is at least $\norm v_{\Hh}$. Evaluating this inequality at that vector proves the Schur estimate. The weighted mean decomposition then leaves one strictly negative scalar and a positive centered form, proving the index and nullity statements exactly as in Theorem~\ref{a:thm:mainseed}.
\end{proof}
This corollary transfers already established branch bounds; it adds no new signed estimate and does not extend any branch to $\delta=0$.

\subsubsection{The signed loss under a parameter decrease}

Fix one $V,Z,H$ on an interval where the complement remains positive, and write
\begin{equation}\label{a:eq:renormalized}
 B_Z=T_Z-D_Z,\qquad A_{\delta,Z}=B_Z+\delta I,
 \qquad R_\delta=S_\delta/\delta.
\end{equation}
The same weighted constrained operator $B_Z$ occurs at each parameter, with a common operator domain. It is not an ordinary-$L^2$ endpoint realization.

\begin{proposition}[Signed renormalized continuation]\label{a:prop:continuation}
Suppose $A_{a,Z}\ge\nu I$, $0<h<\nu$, and $b=a-h>0$. Then
\begin{align}
 R_b[v]-R_a[v]
 &=h\ip{p(v)}{B_Z(B_Z+b)^{-1}(B_Z+a)^{-1}p(v)}_{\Hh},
                  \label{a:eq:signedcontinuation}\\
 R_b[v]&\ge R_a[v]
 -\frac{h(a-\nu)_+}{\nu(\nu-h)}\norm{p(v)}_{\Hh}^2.
                  \label{a:eq:continuationloss}
\end{align}
If the complement was built at $a$ with relative parameter $\theta$, then for $b=qa$ and $\theta<q<1$,
\begin{equation}\label{a:eq:ratio_loss}
 R_{qa}[v]\ge R_a[v]
 -\frac{(1-q)\theta}{(1-\theta)(q-\theta)}\norm{p(v)}_{\Hh}^2.
\end{equation}
\end{proposition}
\begin{proof}
Divide~\eqref{a:eq:scaledschur} by $\delta$ and use the resolvent identity
\[
 a(B_Z+a)^{-1}-b(B_Z+b)^{-1}
 =(a-b)B_Z(B_Z+b)^{-1}(B_Z+a)^{-1}.
\]
It is a bounded operator identity by the spectral theorem and proves~\eqref{a:eq:signedcontinuation}. The spectrum of $B_Z$ lies in $[\nu-a,\infty)$. Nonnegative spectral values contribute nonnegatively. At a negative value $-d$, where $0\le d\le a-\nu$, the negative magnitude is
\[
 \frac{hd}{(a-d)(b-d)}.
\]
It increases with $d$ for $0\le d<b$. Its endpoint is
$h(a-\nu)/[\nu(\nu-h)]$, proving~\eqref{a:eq:continuationloss}. For the relative complement use $\nu=(1-\theta)a$ and substitute $h=(1-q)a$.
\end{proof}

Differentiating the same spectral identity gives
\begin{equation}\label{a:eq:derivativeR}
 \frac{d}{d\delta}R_\delta[v]
 =-\frac{q_0[z_*(v)]}{\delta^2}.
\end{equation}
The response belongs to $\dom B_Z$, so the expression is defined. Formula~\eqref{a:eq:derivativeR} is not a monotonicity theorem: no favorable sign of $q_0[z_*]$ is supplied beyond the proved interval. The continuation loss is charged only to the possible negative spectrum of $B_Z$, but that signed contribution has not been removed.

For example, if at one parameter $a$ the full centered form obeys
$\ell_a\ge g\norm{\cdot}_{\Hh}^2$, let
\begin{equation}\label{a:eq:beta}
 \beta_H=\sup_{v\ne0}\frac{\norm{p(v)}_{\Hh}^2}{G[v]}=1-\sigma_H<1.
\end{equation}
The same projection argument gives
$R_a\ge(g/a)(1-\beta_H)G$. Therefore a sufficient one-step sign criterion is
\begin{equation}\label{a:eq:onestep}
 \frac ga(1-\beta_H)
 >\frac{(1-q)\theta}{(1-\theta)(q-\theta)}\beta_H.
\end{equation}
Every quantity must refer to the same $V,Z,H$. The inequality has not been verified here at any parameter below $1/2-c_{\log2}$.

\subsection{The unresolved sign and the scope of the construction}\label{a:sec:scope}

At a general positive parameter, the positive complement and the scaled response formula give only
\begin{equation}\label{a:eq:anglebound}
 S_\delta[v]\ge\delta\left(G[v]
                 -\frac{\norm{p(v)}_{\Hh}^2}{1-\theta}\right).
\end{equation}
Its sufficient angle condition is $\beta_H<1-\theta$. Ordinary density does not establish that condition. Decreasing $\theta$ changes the required moments and generally changes $H$ and $\beta_H$; one cannot keep a favorable interpolation cost while independently improving the complement parameter.

A finite-dimensional model makes the logical limitation explicit. On $\C^2$ take
\[
 \begin{gathered}
 B=\operatorname{diag}(0,-d),\qquad d>0,\qquad0<\phi<\pi/2,\\
 V=\operatorname{span}\{(\cos\phi,\sin\phi)\},\qquad
 Z=\operatorname{span}\{(-\sin\phi,\cos\phi)\}.
 \end{gathered}
\]
The first coordinate is an exact radical and supplies a lift. For $B+\delta I$ the complement is $\delta-d\cos^2\phi$, while the unit-vector Schur scalar is
\begin{equation}\label{a:eq:model}
 S_\delta=\frac{\delta(\delta-d)}{\delta-d\cos^2\phi}.
\end{equation}
This follows either by the $2\times2$ Schur formula or by dividing the determinant $\delta(\delta-d)$ by the complement entry. When
$d\cos^2\phi<\delta<d$, the complement is positive but the Schur form is negative. Choosing $\cos^2\phi$ sufficiently small also gives a prescribed relative complement bound on a nonempty interval. Exact reference iterations converge there. The example is not the theta operator and is not a counterexample to its desired half bound. It shows why the abstract complement, radical, and residual algebra alone cannot assign the missing sign.

There is a corresponding conditional cost statement for the actual weighted form. Suppose that a normalized centered weighted eigenstate $u$ satisfies
$q_0(u,z)=-d\ip uz_{\Hh}$ for all $z\in X_c$, with $d>0$. It is weighted-orthogonal to $\Rcal$ by~\eqref{a:eq:radicalidentity}. If
$\ell_\delta|_Z\ge b\delta\norm{\cdot}_{\Hh}^2$ with $0<b<1$, put
$v=P_Vu$ and $w=u-Hv\in Z$. Since $u\perp_{\Hh} Hv$,
\[
 \ell_\delta[w]=(\delta-d)\norm u_{\Hh}^2
                   +\delta\norm{Hv}_{\Hh}^2,
 \qquad
 \norm w_{\Hh}^2=\norm u_{\Hh}^2+\norm{Hv}_{\Hh}^2.
\]
Complement coercivity forces
\begin{equation}\label{a:eq:conditionalcost}
 \norm{Hv}_{\Hh}^2\ge
 \left(\frac{d}{(1-b)\delta}-1\right)\norm u_{\Hh}^2.
\end{equation}
Thus a hypothetical negative state is compatible with the complement estimate: its cost can move into the interpolation norm as $\delta$ shrinks. This conditional inequality is not an exclusion of such a state.


\section{Curvature and positive local row obstructions}\label{sec:obstructions}
The generator of Section~\ref{b:sec:source} is actual on its invariant smooth algebra. On real profiles write $L=-A$ and
\begin{equation}\label{b:eq:carredef}
 \Gamma(f,h)=\tfrac12\{L(fh)-fLh-hLf\},\quad
 \Gamma(f)=\Gamma(f,f),\quad \Gamma_2(f)=\tfrac12L\Gamma(f)-\Gamma(f,Lf).
\end{equation}
For $f\in\dom A$ the integrated half defect is
\begin{equation}\label{b:eq:defect}
 \Def[f]=\norm{Af}_\mu^2-\tfrac12\E[f].
\end{equation}
Let $\Eloc=C^\infty_{\mathrm{even}}(\R;\R)$ have the compact-open smooth seminorms
\begin{equation}\label{b:eq:localtopology}
 p^{\mathrm{loc}}_{R,m}(f)=\max_{0\le q\le m}\sup_{|x|\le R}|f^{(q)}(x)|.
\end{equation}
A local row is a continuous real-linear map $\Eloc\to\C$, equivalently an even compactly supported distribution of finite order. A common row family is required in a factorization; its support or order need not be uniform over the family.


\begin{theorem}[Curvature obstruction]\label{b:thm:maincurvature}
For the actual operator $A$, the profile
\[
 g=\kappa''/\kappa-1/4
\]
belongs to $\dom(A)$, satisfies $Ag=g/2$ and $\mu(g)=0$, and obeys
\begin{equation}\label{b:eq:mainquarter}
 \frac{\Gamma_2(g)(x)}{\Gamma(g)(x)}\longrightarrow\frac14
 \quad(|x|\to\infty).
\end{equation}
The denominator is positive in both sufficiently distant tails. For every $K>1/4$ there are a real even $f\in C_c^\infty(\R)$ with $\mu(f)=0$ and a nonempty open interval on which $\Gamma_2(f)<K\Gamma(f)$. Nevertheless,
\[
 \int_\R[\Gamma_2(g)-\tfrac12\Gamma(g)]\,d\mu=0.
\]
\end{theorem}

\begin{theorem}[Local density and positive local rows]\label{b:thm:mainlocal}
The real span $\mathcal G$ of
\begin{equation}\label{b:eq:halfprofiles}
 g_j=\frac{\kappa^{(2j)}}\kappa-4^{-j},\qquad j\ge1,
\end{equation}
is dense in $\Eloc$. Every $g_j$ is an actual centered half-eigenstate, and has exactly centered even compact smooth approximants converging in the graph norm of $A$.

Let $(Z,\eta)$ be a positive measure space, and suppose that outside one null set $\ell_z:\Eloc\to\C$ is a continuous real-linear row, with $z\mapsto\ell_z(f)$ measurable for the profiles in question. There is no identity
\begin{equation}\label{b:eq:mainnoidentity}
 \Def[f]=\int_Z|\ell_z(f)|^2\,d\eta(z)
\end{equation}
for every real even $f\in C_c^\infty(\R)$ with $\mu(f)=0$, using one common row family and measure. If only the lower bound
$\Def[f]\ge\int_Z|\ell_z(f)|^2\,d\eta\ge0$ is asserted on that class, all rows must vanish on $\Eloc$ almost everywhere. Finite or countable collections of rows are included.
\end{theorem}

Sections~\ref{b:sec:curvature}--\ref{b:sec:cutoffs} prove Theorem~\ref{b:thm:maincurvature}. Sections~\ref{b:sec:density}--\ref{b:sec:nogo} prove Theorem~\ref{b:thm:mainlocal}, and Section~\ref{b:sec:paths} treats positive semidefinite finite-path matrices without discarding compatibility relations. The common foundation, explicit-formula bridge and actual half-eigenstates were proved in Sections~\ref{sec:foundations}--\ref{sec:spectral}. The classical explicit formula is used as a theorem, not as a hypothesis about zero locations.


\subsection{Carr\'e du champ and all mixed jump compositions}\label{b:sec:carre}

Define the symmetric displacement measure $J$ by
\[
 \int F(h)\,J(dh)=\int_{s>0}[F(s)+F(-s)]\,d\nu(s),
 \qquad r(x,h)=a(x)\kappa(x+h).
\]
The measure $d\mu(x)r(x,h)J(dh)=\kappa(x)\kappa(x+h)\,dx\,J(dh)$ is invariant under edge reversal $(x,h)\mapsto(x+h,-h)$. Thus
\begin{equation}\label{b:eq:energy-symmetric}
 \E[f]=\tfrac12\int_\R\int
       \kappa(x)\kappa(x+h)|f(x+h)-f(x)|^2\,J(dh)\,dx.
\end{equation}
The factor $1/2$ accounts for the two orientations.

\begin{proposition}[Pointwise and integrated identities]\label{b:prop:carre}
For real even $f,h\in\CE$, all terms in~\eqref{b:eq:carredef} belong to $\CE$, and
\begin{align}
 \Gamma(f,h)(x)&=\tfrac12\int r(x,z)
   [f(x+z)-f(x)][h(x+z)-h(x)]\,J(dz),\label{b:eq:carre-jump}\\
 \Gamma(f,h)&=\tfrac14fh+\tfrac12
   \{fT(\kappa h)+hT(\kappa f)-T(\kappa fh)\},\label{b:eq:carre-source}\\
 \int\Gamma(f)\,d\mu&=\E[f],\qquad
 \int\Gamma_2(f)\,d\mu=\norm{Af}_\mu^2.\label{b:eq:integrated}
\end{align}
\end{proposition}
\begin{proof}
The jump product rule gives~\eqref{b:eq:carre-jump}; its integral is~\eqref{b:eq:energy-symmetric}. Substitution of~\eqref{b:eq:generator-source} gives~\eqref{b:eq:carre-source} and membership in $\CE$. Proposition~\ref{b:prop:generator} admits every product and every subsequent application of $L$. Since $A1=0$ and $\Gamma(f)\in\dom(A)$,
$\int L\Gamma(f)\,d\mu=0$. Also $Lf\in\D$ and edge polarization gives
\[
 \int\Gamma(f,Lf)\,d\mu=\E(f,Lf)=-\norm{Af}_\mu^2.
\]
The bilinear edge integral is absolutely integrable by Cauchy--Schwarz in the edge measure; the other terms lie in $L^2(\mu)\subset L^1(\mu)$. This proves~\eqref{b:eq:integrated}.
\end{proof}

The formulas use the full operator, not its Gamma component alone. To record explicitly the two-step terms, put $G_\Gamma=D_\Gamma-\cg$. For $h\in\RR$,
\begin{align}\label{b:eq:mixed}
 T(\kappa Th)=aQ(WQh)
 &=aG_\Gamma(WG_\Gamma h)-aG_\Gamma(WPh)\notag\\
 &\quad-aP(WG_\Gamma h)+aP(WPh).
\end{align}
Every intermediate $WG_\Gamma h$ and $WPh$ belongs to $\RR$ by Lemma~\ref{b:lem:T}. Thus Gamma--Gamma, both mixed orders, and prime--prime are separately defined smooth $\SSS$ functions. Constants in $G_\Gamma$ and their products are retained.

There is also locally uniform absolute prime--prime summability. Let
\[
 U_h(x)=\sum_{n\ge2}a_n\bigl(|h(x+s_n)|+|h(x-s_n)|\bigr).
\]
Equation~\eqref{b:eq:primepoint} bounds it by $C_\epsilon e^{(1/2+\epsilon)|x|}$. Therefore $WU_h$ decreases faster than every exponential, and the absolute double sum is bounded by
\[
 \sum_{m\ge2}a_m\{W(x+s_m)U_h(x+s_m)+W(x-s_m)U_h(x-s_m)\},
\]
which converges locally uniformly by the same bound. Differentiating the inner prime series is legitimate; multiplication by $W$ places it in $\RR$ before the outer Gamma or prime operation.

For clarity, the corresponding jump integrals are absolutely controlled with their small-shift cancellations retained. For $\sigma\in\{1,-1\}$ define
\[
 I_f(z,t)=a(z)\kappa(z+\sigma t)[f(z+\sigma t)-f(z)],\qquad
 J_f(z,t)=\tfrac12a(z)\kappa(z+\sigma t)[f(z+\sigma t)-f(z)]^2.
\]
For $f\in\CE$, the $\nu$-integral of the absolute value of every fixed $z$ derivative has an exponential envelope in $z$. On $0<t\le1$ the two differentiated expressions are respectively $O(t)$ and $O(t^2)$ with such envelopes. Above one, expansion reduces the estimate to shifted rapid functions $\kappa f^j$ and exponential factors in $z$; the Gamma tail and~\eqref{b:eq:primepoint} apply to each channel.

In $L\Gamma(f)$ retain the difference $J_f(x+h,t)-J_f(x,t)$. Its absolute $t$ integral is $O(|h|)$ for $|h|\le1$ on compact $x$ intervals. In $\Gamma(f,Lf)$ retain
$[f(x+h)-f(x)][I_f(x+h,t)-I_f(x,t)]$, whose absolute $t$ integral is $O(h^2)$. For $|h|>1$, the outer factor $\kappa(x+h)$ controls the exponential envelope and the complete prime and Gamma tails. Fubini therefore applies to these cancelled integrands for every pair of channels, including large opposite shifts. It does not authorize separating the infinite Gamma diagonal rate.

For finite jump cutoffs in both instances of the generator, put
\[
 p=f(x+h)-f(x),\quad q=f(x+k)-f(x),\quad
 z=f(x+h+k)-f(x+h)-f(x+k)+f(x),
\]
\[
 C_{h,k}=a(x)\kappa(x+h)a(x+h)\kappa(x+h+k),\qquad
 D_{h,k}=a(x)^2\kappa(x+h)\kappa(x+k).
\]
Direct expansion gives
\begin{align}\label{b:eq:doublejump}
 4\Gamma_2^{\mathrm{cut}}(f)(x)
 =\iint\bigl\{&C_{h,k}(q+z)^2-D_{h,k}q^2\notag\\
 &-2p[C_{h,k}(q+z)-D_{h,k}q]\bigr\}
     J_{\mathrm{cut}}(dh)J_{\mathrm{cut}}(dk).
\end{align}
The preceding estimates give the full limit with the braces kept together. The first two terms cannot be independently integrated over the infinite activity in $h$. No general nonnegative-square interpretation of the bracket is asserted.

\subsection{The exact half state and the quarter curvature limit}\label{b:sec:curvature}

Proposition~\ref{f:prop:halfeigen} proves, using the full form domain, that every $g_j$ in~\eqref{f:eq:halfprofiles} satisfies
\begin{equation}\label{b:eq:actualhalf}
 g_j\in\dom(A),\qquad Ag_j=\tfrac12g_j,\qquad \mu(g_j)=0.
\end{equation}
The upper and lower theta estimates~\eqref{a:eq:derivativeupper}--\eqref{a:eq:thetalower} imply
\begin{equation}\label{b:eq:ratio}
 |\kappa^{(m)}(x)/\kappa(x)|\le C_me^{2m|x|}.
\end{equation}
Every derivative of a quotient in~\eqref{f:eq:halfprofiles} is a finite sum of products of these ratios; hence $g_j\in\CE$. In particular all ensuing identities concern actual operator-domain vectors, and~\eqref{b:eq:generator-source} gives the pointwise equality
\begin{equation}\label{b:eq:halfsource}
 T(\kappa g_j)=0.
\end{equation}
The smooth residual of the Hilbert-space identity vanishes pointwise because the density of $\mu$ is strictly positive.

For the first state $g=g_1$, let $t=\pi e^{2x}$ for $x\ge0$. The theta derivative recurrence is
\begin{equation}\label{b:eq:recurrence}
 \kappa^{(m)}(x)=e^{x/2}\sum_{n\ge1}P_m(n^2t)e^{-n^2t},\qquad
 P_{m+1}(t)=(1/2-2t)P_m(t)+2tP_m'(t),
\end{equation}
with $P_0(t)=4t^2-6t$. In particular,
\[
 P_2(t)=16t^4-112t^3+165t^2-\tfrac{75}{2}t.
\]
For $t\ge\pi$,
\begin{equation}\label{b:eq:thetatail}
 \sum_{n\ge2}n^{2m+4}e^{-(n^2-1)t}
 \le e^{-3t}\sum_{n\ge2}n^{2m+4}e^{-\pi(n^2-4)}.
\end{equation}
The last series is finite. Since $P_0(t)$ is bounded below by a positive multiple of $t^2$ on this range,
\[
 \frac{\kappa''(x)}{\kappa(x)}
  =\frac{P_2(t)}{P_0(t)}+O(t^2e^{-3t})=4t^2+O(t).
\]
Evenness gives the two-sided asymptotic
\begin{equation}\label{b:eq:gasymptotic}
 g(x)=4\pi^2e^{4|x|}+O(e^{2|x|}).
\end{equation}
No numerical theta truncation enters this estimate.

Define the full source corrections
\begin{equation}\label{b:eq:HK}
 h=\kappa g^2\in\RR,\qquad H=Th\in\SSS,\qquad
 K=T(\kappa H)=aQ(WQh)\in\SSS.
\end{equation}
Formula~\eqref{b:eq:mixed} expands $K$ with all four channel combinations.

\begin{proposition}[Exact signed identities]\label{b:prop:quarteridentities}
Pointwise on $\R$,
\begin{align}
 \Gamma(g)&=\tfrac14g^2-\tfrac12H,\label{b:eq:Gg}\\
 L\Gamma(g)&=-\tfrac18g^2+\tfrac12K,\label{b:eq:LGg}\\
 \Gamma_2(g)&=\tfrac1{16}g^2-\tfrac14H+\tfrac14K,\label{b:eq:G2g}\\
 \Gamma_2(g)-\tfrac12\Gamma(g)&=-\tfrac1{16}g^2+\tfrac14K.\label{b:eq:halfcurvature}
\end{align}
Consequently, as $r=|x|\to\infty$,
\begin{align}
 \Gamma(g)(x)&=4\pi^4e^{8r}(1+o(1)),\label{b:eq:Gtail}\\
 \Gamma_2(g)(x)&=\pi^4e^{8r}(1+o(1)),\label{b:eq:G2tail}\\
 \Gamma_2(g)(x)-\tfrac12\Gamma(g)(x)&=-\pi^4e^{8r}(1+o(1)).\label{b:eq:deftail}
\end{align}
\end{proposition}
\begin{proof}
Equations~\eqref{b:eq:carre-source} and~\eqref{b:eq:halfsource} give~\eqref{b:eq:Gg}. The source row gives $L(g^2)=-g^2/2-H$ and $LH=-H/2-K$, so the two $H$ terms cancel in $L\Gamma(g)$ and yield~\eqref{b:eq:LGg}. Since $Lg=-g/2$, bilinearity gives
$\Gamma_2(g)=L\Gamma(g)/2+\Gamma(g)/2$, proving~\eqref{b:eq:G2g}--\eqref{b:eq:halfcurvature}. Every application of $L$ is permitted by Proposition~\ref{b:prop:generator}. For every $\epsilon>0$, $H,K=O(e^{\epsilon r})$. Combining this with~\eqref{b:eq:gasymptotic} gives the three asymptotics.
\end{proof}

The quotient in~\eqref{b:eq:mainquarter} follows. More generally, for each fixed $K_0>1/4$,
\begin{equation}\label{b:eq:anycurvature}
 \Gamma_2(g)-K_0\Gamma(g)
  =(1/16-K_0/4)g^2+(K_0/2-1/4)H+K/4<0
\end{equation}
in both sufficiently distant tails. The number $1/4$ is an upper restriction on a possible pointwise lower bound; the calculation does not establish $\Gamma_2\ge\Gamma/4$.

\subsection{Compact witnesses and the integrated cancellation}\label{b:sec:cutoffs}

Fix real even $\chi\in C_c^\infty$ with $0\le\chi\le1$, equal to one on $[-1,1]$ and supported in $[-2,2]$, and put $\chi_R(x)=\chi(x/R)$ for $R\ge1$. Initially let $g_R=\chi_Rg$. These profiles and their squares have derivativewise exponential bounds uniform in $R$. The rapid theta decay therefore gives
\begin{equation}\label{b:eq:rapidcutoffs}
 \kappa g_R\longrightarrow\kappa g,\qquad
 \kappa g_R^2\longrightarrow\kappa g^2
 \quad\hbox{in }\RR.
\end{equation}
Indeed the profiles agree with their limits on every fixed compact set for large $R$, and outside that set stronger rapid seminorms control all cutoff derivative terms.

Set $F_R=T(\kappa g_R)$ and $H_R=T(\kappa g_R^2)$. Lemma~\ref{b:lem:T} gives
\begin{equation}\label{b:eq:sourcecutoffs}
 F_R\longrightarrow0,\qquad H_R\longrightarrow H\quad\hbox{in }\SSS.
\end{equation}
Using~\eqref{b:eq:carre-source}, write
\[
 B_R=\Gamma(g_R)=g_R^2/4+g_RF_R-H_R/2,
 \qquad B=\Gamma(g).
\]
Then $B_R\to B$ locally smoothly and $\kappa B_R\to\kappa B$ in $\RR$. For the product $g_RF_R$, use the uniform exponential bounds on $g_R$, convergence in one specified small-exponential seminorm at each derivative order, and the theta decay. Equation~\eqref{b:eq:generator-source} now yields $LB_R\to LB$ locally uniformly. Also
\[
 \Gamma(g_R,F_R)=g_RF_R/4+\tfrac12\{g_RT(\kappa F_R)
                  +F_RT(\kappa g_R)-T(\kappa g_RF_R)\}\longrightarrow0
\]
locally uniformly, since the two $T$ arguments containing $F_R$ converge to zero in $\RR$. Finally $Lg_R=-g_R/2-F_R$, so
\[
 \Gamma_2(g_R)-\tfrac12\Gamma(g_R)
       =\tfrac12LB_R+\Gamma(g_R,F_R).
\]
Thus
\begin{equation}\label{b:eq:curvaturecutoffs}
 \Gamma(g_R)\to\Gamma(g),\qquad \Gamma_2(g_R)\to\Gamma_2(g)
 \quad\hbox{locally uniformly}.
\end{equation}
This argument avoids interchanging an uncontrolled double jump sum.

Choose a fixed real even $\psi\in C_c^\infty$ with $\mu(\psi)=1$ and set
\begin{equation}\label{b:eq:centeredcutoff}
 f_R=\chi_Rg-\mu(\chi_Rg)\psi.
\end{equation}
Theta domination gives $\chi_Rg\to g$ in $\He$ and $\mu(\chi_Rg)\to0$. These centered compact profiles have the same uniform derivative envelopes and rapid-weighted limits. The proof above therefore also gives~\eqref{b:eq:curvaturecutoffs} with $f_R$ in place of $g_R$. Fix $K_0>1/4$ and a finite $x_0$ satisfying the strict tail inequality~\eqref{b:eq:anycurvature}. For sufficiently large finite $R$, that strict inequality persists at $x_0$ for $f_R$. Smoothness makes it hold on a nonempty open interval, hence not merely on a null set. This proves the compact-test assertion of Theorem~\ref{b:thm:maincurvature}.

There is actual graph convergence for these particular cutoffs: $g_R\to g$ and $Ag_R\to Ag$ in $\He$, by~\eqref{b:eq:sourcecutoffs} and theta domination. Likewise $B_R\to B$ and $AB_R\to AB$ in $\He$, using the uniform exponential envelopes and $\kappa B_R\to\kappa B$ in $\RR$. No assertion that compact smooth functions are an operator core for arbitrary vectors is required.

\begin{proposition}[Exact integrated cancellation]\label{b:prop:cancellation}
The curvature defect of $g$ integrates to zero, and
\begin{align}
 \int\Gamma(g)\,d\mu&=\tfrac12\norm g_\mu^2,
 &\int\Gamma_2(g)\,d\mu&=\tfrac14\norm g_\mu^2,\label{b:eq:halfintegrals}\\
 \int H\,d\mu&=-\tfrac12\norm g_\mu^2,
 &\int K\,d\mu&=\tfrac14\norm g_\mu^2.\label{b:eq:HKintegrals}
\end{align}
\end{proposition}
\begin{proof}
All terms are integrable by the admitted-class estimates. Equation~\eqref{b:eq:integrated} and $Ag=g/2$ prove~\eqref{b:eq:halfintegrals}; integrating~\eqref{b:eq:Gg} and~\eqref{b:eq:halfcurvature} proves~\eqref{b:eq:HKintegrals}.
\end{proof}

The negative tails in~\eqref{b:eq:deftail} must therefore be balanced by a positive part elsewhere. In particular $K$ cannot be discarded in an integrated argument. Pointwise curvature failure is not a counterexample to the integrated half inequality.

\subsection{Local smooth density of the half eigenstates}\label{b:sec:density}

\begin{theorem}[Local density]\label{b:thm:density}
The closure of $\mathcal G=\operatorname{span}_{\R}\{g_j:j\ge1\}$ in $\Eloc$ is all of $\Eloc$.
\end{theorem}
\begin{proof}
By Hahn--Banach separation it suffices to show that every continuous linear functional $\ell$ annihilating all $g_j$ is zero. Extend it to all smooth functions by even symmetrization, and then complexify. The extension is an even compactly supported distribution. Set $w=\ell/\kappa$, meaning multiplication by the smooth positive reciprocal on a neighborhood of its compact support.

By Lemma~\ref{d:lem:theta-tails}, $\kappa$ has an even holomorphic continuation to the strip $\mathcal S=\{z:|\operatorname{Im}z|<\pi/4\}$ with all the derivative tail bounds.

The function
\[
 F(t)=\ip{w(x)}{\kappa(x+t)},\qquad t\in\mathcal S,
\]
where brackets denote distributional pairing without conjugation, is holomorphic. Compact support and finite order justify differentiation under this pairing. Evenness gives $F(-t)=F(t)$, and annihilation of the $g_j$ gives
\[
 F(0)=\ell(1),\qquad F^{(2j)}(0)=4^{-j}\ell(1)\quad(j\ge1).
\]
Its Taylor series is therefore the series of $\ell(1)\cosh(t/2)$. The identity theorem on the connected strip gives
\begin{equation}\label{b:eq:Fcosh}
 F(t)=\ell(1)\cosh(t/2),\qquad t\in\mathcal S.
\end{equation}
On the real axis $F$ decreases faster than every exponential. Indeed, for $\supp(w)\subset[-R,R]$ and order at most $m$, the distribution bound with a fixed auxiliary cutoff gives
\[
 |F(t)|\le C\max_{0\le q\le m}\sup_{|x|\le R+1}|\kappa^{(q)}(x+t)|
       \le C_Be^{-B|t|}\quad(B>0).
\]
The derivative theta envelope justifies the last step. Hence~\eqref{b:eq:Fcosh} forces $\ell(1)=0$ and $F=0$.

On the real line, evenness identifies $F$ with $w*\kappa$. This convolution is Schwartz, because $w$ is compactly supported and $\kappa$ is Schwartz. In the convention $\widehat h(\xi)=\int e^{-i\xi x}h(x)\,dx$, Fourier transformation gives
\[
 \widehat w(\xi)\widehat\kappa(\xi)=0.
\]
The transform of a compact distribution is entire, directly by pairing it with $e^{-izx}$. Since $\widehat\kappa(0)=\int\kappa>0$, continuity makes $\widehat\kappa$ nonzero on a real interval around zero. Therefore $\widehat w$ vanishes on that interval, and the identity theorem makes it identically zero. Fourier injectivity for distributions yields $w=0$, hence $\ell=0$. Hahn--Banach proves the assertion.
\end{proof}

This proof does not divide by the completed zeta function at its zeros. Only the nonvanishing of the theta transform near zero is needed in its last step. Weighted centering is not continuous in the local smooth topology, so there is no contradiction between $\mu(g_j)=0$ and local approximation of the constant profile.

\begin{corollary}[Exact finite data interpolation]\label{b:cor:finiteimage}
For every continuous real-linear map $B:\Eloc\to V$ into a finite-dimensional space,
\[
 B(\mathcal G)=B(\Eloc).
\]
\end{corollary}
\begin{proof}
Density makes $B(\mathcal G)$ dense in $B(\Eloc)$. The former is a linear subspace of a finite-dimensional space and is closed, proving equality.
\end{proof}

Thus finite endpoint and jet data can be matched exactly by a finite linear combination of the half states, whenever the data are compatible with evenness and any repeated endpoints. This statement is local, not weighted.

\subsection{Graph cutoffs and a strictly positive compact defect}\label{b:sec:graph}

Use the fixed $\chi_R$ and $\psi$ of Section~\ref{b:sec:cutoffs}. For each fixed $j\ge1$ set
\begin{equation}\label{b:eq:fjR}
 e_{j,R}=\mu(\chi_Rg_j),\qquad
 f_{j,R}=\chi_Rg_j-e_{j,R}\psi.
\end{equation}
These profiles are real, even, compact smooth, and exactly centered.

\begin{lemma}[Special graph approximation]\label{b:lem:graph}
For each fixed $j$, as $R\to\infty$,
\[
 f_{j,R}\to g_j\text{ in }\Eloc\text{ and }\He,
 \qquad Af_{j,R}\to Ag_j\text{ in }\He,
 \qquad\Def[f_{j,R}]\to0.
\]
\end{lemma}
\begin{proof}
Theta domination and~\eqref{b:eq:actualhalf} give $e_{j,R}\to0$. Local smooth and weighted Hilbert convergence follow. Put $r_j=\kappa g_j=\kappa^{(2j)}-4^{-j}\kappa\in\RR$. Then
\[
 \kappa(f_{j,R}-g_j)=(\chi_R-1)r_j-e_{j,R}\kappa\psi\longrightarrow0
 \quad\hbox{in }\RR.
\]
For every seminorm this follows from Leibniz's rule and stronger rapid decay on the tail supporting the cutoff derivatives. Lemma~\ref{b:lem:T} gives convergence of its $T$ image in $\SSS$, hence in $\He$. Equation~\eqref{b:eq:generator-source} proves graph convergence. Finally~\eqref{b:eq:defect} is graph-continuous and $\Def[g_j]=0$.
\end{proof}

This is an approximation theorem for the specified half states. It does not assert a general compact smooth operator core.

\begin{lemma}[Logarithmic Gamma coercivity]\label{b:lem:fourier}
For the unitary Fourier transform,
\begin{equation}\label{b:eq:multiplier}
 \tg[u]=\int_\R d_\Gamma(\xi)|\widehat u(\xi)|^2\,d\xi,
 \qquad d_\Gamma(\xi)\ge\tfrac12\log(1+4\xi^2).
\end{equation}
\end{lemma}
\begin{proof}
This is Lemma~\ref{a:lem:fourier}, including its Plancherel--Tonelli identity for possibly infinite form values.
\end{proof}

\begin{proposition}[The defect is not identically zero]\label{b:prop:positive}
There is a real even $f\in C_c^\infty(\R)$ with $\mu(f)=0$ and $\Def[f]>0$.
\end{proposition}
\begin{proof}
Choose nonzero real even $\phi\in C_c^\infty$. For a real frequency $R>0$, put
\[
 m_R=\mu(\phi\cos(Rx)),\qquad h_R=\phi\cos(Rx)-m_R\psi,
 \qquad u_R=\kappa h_R,
\]
and write $b=\kappa\phi$, $d=\kappa\psi$. These are centered compact smooth profiles with fixed compact support. Riemann--Lebesgue gives
\begin{equation}\label{b:eq:HFmass}
 m_R\to0,\quad
 \norm{h_R}_\mu^2\to\tfrac12\int\phi^2\,d\mu>0,\quad
 \norm{u_R}_2^2\to\tfrac12\norm b_2^2>0.
\end{equation}
In the unitary Fourier convention,
\[
 \widehat u_R(\xi)=\tfrac12\{\widehat b(\xi-R)+\widehat b(\xi+R)\}
                         -m_R\widehat d(\xi).
\]
Because $\widehat b$ and $\widehat d$ are Schwartz, the Fourier mass in $[-M,M]$ tends to zero for every fixed $M$. A positive amount of mass therefore escapes to arbitrarily high frequencies. Lemma~\ref{b:lem:fourier} gives
\[
 \liminf_{R\to\infty}\tg[u_R]
 \ge\tfrac12\log(1+4M^2)\,\tfrac12\norm b_2^2.
\]
Letting $M\to\infty$ proves $\tg[u_R]\to\infty$.

The complete weighted prime estimate in Lemma~\ref{a:lem:prime} gives
$|p[u]|\le C_P\norm u_{\Hh}^2$ for a finite $C_P$, and
$\norm u_2^2\le\norm W_\infty\norm u_{\Hh}^2$. Since $\norm{\kappa h}_{\Hh}=\norm h_\mu$, the full bridge yields
\[
 \E[h_R]\ge\tg[u_R]-C_0\norm{h_R}_\mu^2\longrightarrow\infty,
 \qquad C_0=|\cg|\norm W_\infty+C_P.
\]
No nonzero prime tail is discarded. Proposition~\ref{b:prop:generator} places each $h_R$ in $\dom(A)$. Cauchy--Schwarz then gives
\[
 \Def[h_R]\ge\E[h_R]
       \left(\frac{\E[h_R]}{\norm{h_R}_\mu^2}-\frac12\right)>0
\]
for all sufficiently large $R$, by~\eqref{b:eq:HFmass}. Choose any such finite frequency.
\end{proof}

The witness is strictly positive. It is not a negative direction for the integrated inequality and does not construct an actual subthreshold state.

\subsection{No nontrivial positive local row factorization}\label{b:sec:nogo}

\begin{proof}[Proof of the factorization assertions in Theorem~\ref{b:thm:mainlocal}]
Suppose~\eqref{b:eq:mainnoidentity} holds on the exactly centered compact class. Fix $j$ and take integer $R\to\infty$ in~\eqref{b:eq:fjR}. Local smooth convergence and row continuity give
$\ell_z(f_{j,R})\to\ell_z(g_j)$ wherever the row is continuous. Fatou's lemma and Lemma~\ref{b:lem:graph} imply
\begin{align}\label{b:eq:Fatou}
 \int_Z|\ell_z(g_j)|^2\,d\eta
 &\le\liminf_{R\to\infty}\int_Z|\ell_z(f_{j,R})|^2\,d\eta\notag\\
 &=\liminf_{R\to\infty}\Def[f_{j,R}]=0.
\end{align}
Thus $\ell_z(g_j)=0$ almost everywhere. There are countably many $j$, so removal of their exceptional null sets and the common continuity exceptional set leaves one full-measure set on which every $g_j$ is annihilated. Theorem~\ref{b:thm:density} and continuity force $\ell_z=0$ on all of $\Eloc$ at every point of this set. The right side of~\eqref{b:eq:mainnoidentity} is then zero for every test, contradicting Proposition~\ref{b:prop:positive}.

For the lower-bound assertion, replace the equality in~\eqref{b:eq:Fatou} by the upper bound $\Def[f_{j,R}]\to0$. The same reasoning forces all rows to vanish, but gives no contradiction to $\Def\ge0$ itself: its lower bound may simply be zero. A countable row collection is handled by one further countable intersection of full-measure sets, or by the product with counting measure. Neither Fatou's lemma nor the zero-integral argument requires sigma-finiteness.
\end{proof}

The graph-cutoff step is essential. An identity asserted only on compact tests has not been substituted at an unbounded eigenstate without justification. The row family and positive measure must be fixed for all inputs; choosing rows separately for each $f$ is outside the theorem. Supports and distributional orders may vary without any uniform bound. Consequently the statement includes integrations over arbitrarily large Gamma and prime-power displacements, provided each individual row has compact dependence.

\subsection{Finite paths and the actual compatible data space}\label{b:sec:paths}

For each parameter $z$, let $B_z$ collect finitely many endpoint values, derivatives, or edge differences of a profile in $\Eloc$. Its components are continuous local rows. The image
\[
 V_z=B_z(\Eloc)
\]
is regarded as a real finite-dimensional space, even when the coordinates are complex. For edge differences it already contains all cycle relations, together with evenness and repeated-endpoint relations. It need not be the full redundant coordinate space.

\begin{corollary}[Positive finite-path matrices]\label{b:cor:paths}
Suppose $M_z$ is a finite Hermitian matrix whose quadratic form is nonnegative on the actual image $V_z$. Assume a common measurable family of these quadratic expressions and a positive measure $\eta$. There is no identity
\begin{equation}\label{b:eq:matrixidentity}
 \Def[f]=\int_Z(B_zf)^*M_z(B_zf)\,d\eta(z)
\end{equation}
for every real even exactly centered $f\in C_c^\infty(\R)$. If equality is replaced by a nonnegative lower bound, the quadratic form must vanish on $V_z$ almost everywhere.
\end{corollary}
\begin{proof}
Apply Lemma~\ref{b:lem:graph} and Fatou to each fixed $g_j$. On one common full-measure set,
$(B_zg_j)^*M_z(B_zg_j)=0$ for all $j$. On a real vector space a nonnegative quadratic form has zero real polarized pairing between any zero-norm vector and every other vector: expand its value on $v+tw$ for all real $t$. Hence the form vanishes on the real span of $\{B_zg_j:j\ge1\}$. Corollary~\ref{b:cor:finiteimage} identifies this span with $V_z$. The integrand is consequently zero for every even smooth profile, contradicting Proposition~\ref{b:prop:positive} in the equality case. The lower-bound conclusion follows by the same Fatou argument.
\end{proof}

This proof uses real polarization $\operatorname{Re}(v^*M_zw)$; it does not assume that $V_z$ is a complex subspace and does not require a measurable square-root selection. The support points can be
\[
 x,\quad x+h_1,\quad x+h_1+h_2,\quad\ldots,\quad x+h_1+\cdots+h_m
\]
for any finite $m$, with no bound on the individual shifts. The measure may include all Gamma and prime-power paths, both orientations, every mixed channel, and countably many finite path lengths. The obstruction comes from the finite local dependence of each nonnegative integrand, not from a bound on the total range or number of jumps.

\subsection{Weighted limits and the remaining integrated problem}\label{b:sec:scope}

Local density gives no control of the weighted cost of an approximation. A precise distinction is useful.

\begin{proposition}[Weighted cost of local approximation]\label{b:prop:blowup}
Let $f\in\He\cap\Eloc$ lie outside $\ker(A-\tfrac12I)$. If $v_n\in\mathcal G$ converges to $f$ in $\Eloc$, then $\norm{v_n}_\mu\to\infty$.
\end{proposition}
\begin{proof}
Otherwise there is a bounded subsequence in $\He$, with a further weakly convergent subsequence. The half-eigenspace is a closed linear subspace of $\He$, hence weakly closed, so the weak limit lies in it. Pairing with any compact smooth even test and using local convergence identifies the limit with $f$. Such tests are dense in $\He$, giving a contradiction. Thus no bounded subsequence exists, which is precisely divergence of the norms to infinity.
\end{proof}

The constant profile and the compact witness of Proposition~\ref{b:prop:positive} are examples. In particular even a centered compact target can have this divergent approximation cost. This local argument does not prove weighted completeness or identify the entire half-eigenspace. The weighted derivative closure, a different space from the full centered Hilbert space, is analyzed in Section~\ref{sec:synthesis-unified}.

For an admitted real profile, write $t=T(\kappa f)$. The actual source row gives the exact integrated identity
\begin{equation}\label{b:eq:signedsource-defect}
 \Def[f]=\norm t_\mu^2+\tfrac12\ip t f_\mu.
\end{equation}
The square is nonnegative, but no universal sign has been proved for the sum. Its global nature is visible without guessing a topology: for every fixed $x$, the row $f\mapsto T(\kappa f)(x)$ annihilates every $g_j$, while
\[
 T(\kappa)(x)=-\tfrac12.
\]
By Theorem~\ref{b:thm:density} it cannot be continuous in $\Eloc$. Thus this genuinely global row is outside the local-square obstruction, despite being evaluated at a single spatial point.

The integrated edge formulation is also exact but does not supply a sign. Equip edge functions with the norm
\[
 \norm F_{\mathrm{edge}}^2=\tfrac12\int_\R\int
     \kappa(x)\kappa(x+h)|F(x,h)|^2\,J(dh)\,dx,
 \qquad df(x,h)=f(x+h)-f(x).
\]
The difference map $d$ with domain $\D$ is closed, by the same almost-everywhere and Fatou argument used for the energy; its norm squared is $\E[f]$. The closed-form representation gives $A=d^*d$. For $f\in\CE$, invariance implies $Af\in\D$, so $df\in\dom(dd^*)$ and
\begin{equation}\label{b:eq:edge-defect}
 \Def[f]=\ip{df}{(dd^*-\tfrac12I)df}_{\mathrm{edge}}.
\end{equation}
This is a domain-correct reformulation, not a positivity proof. Corollary~\ref{b:cor:paths} shows that imposing the finite cycle relations alone cannot produce a nonzero positive finite-path density for this defect.

The obstruction permits genuinely global integral or frequency-space rows, nonlocal inverses, nonlinear constructions, and signed local identities with cancellation after integration. It also leaves open estimates using additional information from an actual candidate eigen-equation. No such method is proved positive here. In particular the explicit negative pointwise tails and the positive compact integrated witness establish neither a negative integrated direction nor a global half gap. The integrated half inequality and the Riemann hypothesis remain unresolved by the results of this paper.


\section{Weighted synthesis and full multiplicity zero jets}\label{sec:synthesis-unified}
The local density just proved is not weighted completeness. This section determines the weighted derivative closure through the exact theta transform, without a zero-location or simplicity hypothesis. The proofs of the actual Abel bound and its source recovery are included in Appendix~\ref{app:abel}; compact division and root-space structure are recorded in Appendix~\ref{app:compactmodel}.
Put $a=i/2$ and define
\begin{equation}\label{d:eq:center-intro}
 m(u)=F_u(a),\qquad Cu=u-2m(u)\kappa,\qquad \Hh_c=\ker m.
\end{equation}
Thus $C$ is exactly the weighted centering projection used earlier. Set
\begin{equation}\label{d:eq:N-intro}
 \Nn=\cl_{\Hh}\{r_j:j\ge1\},\qquad \Kk=\Hh_c\ominus\Nn,
 \qquad r_j=\kappa^{(2j)}-4^{-j}\kappa.
\end{equation}
If $\nu(w)$ is the full multiplicity of a zero $w$ of $\Xi$, define
\begin{align}
 \Ii&=\{u\in\Hh_c:F_u=\Xi Q\text{ for an even entire }Q\},\label{d:eq:I-intro}\\
 \Zz&=\cl_{\Hh}\{W(x)\partial_z^k\cos(zx)|_{z=w}:\Xi(w)=0,\ 0\le k<\nu(w)\}.
 \label{d:eq:Z-intro}
\end{align}
The algebraic span in the last line is denoted by $\mathcal J_{\mathrm{fin}}$. Its members are full zero jets, not only zero-value representatives.
The exact defect decomposition, proved below, is
\begin{equation}\label{d:eq:main-defect}
 \Kk=\Zz\oplus(\Ii\cap\Nn^\perp).
\end{equation}
The open translation lens is
\begin{equation}\label{d:eq:lens-intro}
 \Omega=\{b:|\operatorname{Im}b|<\pi/4,\quad
 2e^{-2|\operatorname{Re}b|}\cos(2\operatorname{Im}b)>1\}.
\end{equation}


\subsection{The full ideal and its annihilator}\label{d:sec:ideal}

The Fourier image $\Aa=\{F_u:u\in\Hh\}$ is a Hilbert space with the transported norm $\norm{F_u}_{\Aa}=\norm{u}_{\Hh}$. Injectivity makes this definition unambiguous. Equation~\eqref{d:eq:eval-bound} makes every derivative evaluation continuous. Since
\[
 \ip{u}{\kappa}_{\Hh}=2m(u),\qquad \norm{\kappa}_{\Hh}^2=1,
\]
$C$ is the orthogonal projection asserted above. Integration by parts, justified by the theta tails, gives
\begin{equation}\label{d:eq:derivative-transforms}
 F_{\kappa^{(2j)}}(z)=(-z^2)^j\Xi(z),\qquad
 F_{r_j}(z)=\bigl((-z^2)^j-4^{-j}\bigr)\Xi(z).
\end{equation}
Every derivative lies in $\Hh$, again by the tails.

It is useful to distinguish the raw wave from its conjugated Riesz representative:
\begin{equation}\label{d:eq:waves}
 J_{w,k}(x)=W(x)\partial_z^k\cos(zx)|_{z=w},\qquad
 E_{w,k}=J_{\overline w,k}.
\end{equation}
They belong to $\Hh$ for every fixed $w,k$, and the linear-first convention gives
\begin{equation}\label{d:eq:Riesz}
 F_u^{(k)}(w)=\ip{u}{E_{w,k}}_{\Hh}.
\end{equation}
The raw and Riesz families have the same closed span when $w$ ranges over all zeros. For $k<\nu(w)$,
\begin{align}
 m(J_{w,k})&=\tfrac12\Xi^{(k)}(w)=0,\label{d:eq:wave-center}\\
 \ip{J_{w,k}}{r_j}_{\Hh}
 &=\left.\partial_z^k\left[\bigl((-z^2)^j-4^{-j}\bigr)\Xi(z)\right]\right|_{z=w}=0.
 \label{d:eq:wave-radical}
\end{align}
Thus $\Zz\subset\Kk$.

\begin{proposition}[Closed ideal and exact defect]\label{d:prop:ideal}
The space $\Ii$ is closed and
\begin{equation}\label{d:eq:I-jets}
 \Ii=\{u\in\Hh_c:F_u^{(k)}(w)=0\text{ whenever }\Xi(w)=0,\ 0\le k<\nu(w)\}
 =\Hh_c\cap\Zz^\perp.
\end{equation}
It contains $\Nn$. For $u\in\Ii$, the entire quotient $Q=F_u/\Xi$ is even and satisfies $Q(a)=0$. Conversely these conditions on a quotient characterize $\Ii$. In particular~\eqref{d:eq:main-defect} holds and
\begin{equation}\label{d:eq:equivalence}
 \Zz=\Kk\quad\Longleftrightarrow\quad\Ii=\Nn.
\end{equation}
\end{proposition}
\begin{proof}
At each zero of $\Xi$, the conditions in~\eqref{d:eq:I-jets} are exactly the removable-singularity conditions for $F_u/\Xi$. This gives an entire quotient; its evenness follows off the zero set and then everywhere by continuation. The converse is immediate. Since $\Xi(a)=1/2$, centering is equivalent to $Q(a)=0$. The bounded evaluations make the intersection of all the specified kernels closed. Equations~\eqref{d:eq:Riesz} and~\eqref{d:eq:derivative-transforms} prove the orthogonal-complement description and $\Nn\subset\Ii$.

In $\Hh_c$, the decomposition $\Hh_c=\Zz\oplus\Ii$ is orthogonal. As $\Nn\subset\Ii$, taking its orthogonal complement within $\Hh_c$ gives~\eqref{d:eq:main-defect} and~\eqref{d:eq:equivalence}.
\end{proof}

This is full-multiplicity divisibility, not merely vanishing of the values at distinct zeros. Nothing in Proposition~\ref{d:prop:ideal} supplies a growth bound on $Q$. The full approximation statement proved in Theorem~\ref{d:thm:full-synthesis} is
\begin{equation}\label{d:eq:ideal-synthesis}
 \overline{\{\Xi p:p\text{ even polynomial},\ p(a)=0\}}^{\Aa}
 =\{\Xi Q\in\Aa:Q\text{ even entire},\ Q(a)=0\}.
\end{equation}

\subsection{The exact complex translation lens}\label{d:sec:translations}

For $|\operatorname{Im}b|<\pi/4$ define
\begin{equation}\label{d:eq:H-R}
 H_b(x)=\frac{\kappa(x+b)+\kappa(x-b)}2,\qquad
 R_b(x)=H_b(x)-\cosh(b/2)\kappa(x).
\end{equation}
These are even functions of $x$. Membership will restrict the parameter to the lens~\eqref{d:eq:lens-intro}.

\begin{theorem}[Exact admission and holomorphy]\label{d:thm:lens}
Within the continuation strip $|\operatorname{Im}b|<\pi/4$, each of the conditions $H_b\in\Hh$ and $R_b\in\Hh$ holds if and only if $b\in\Omega$. On $\Omega$ the maps $b\mapsto H_b$ and $b\mapsto R_b$ are $\Hh$-valued holomorphic. They satisfy
\begin{equation}\label{d:eq:translation-transform}
 F_{H_b}(z)=\cos(bz)\Xi(z),\qquad
 F_{R_b}(z)=\bigl(\cos(bz)-\cosh(b/2)\bigr)\Xi(z).
\end{equation}
In particular $R_b\in\Ii$.
\end{theorem}
\begin{proof}
Write $b=u+iv$. On a compact subset of $\Omega$,
$2e^{-2|u|}\cos(2v)-1$ has a positive lower bound. The theta upper bounds and the lower bound for $W$ show that $|H_b|^2/W$ and the corresponding expressions for every fixed parameter derivative are bounded by
\begin{equation}\label{d:eq:lens-majorant}
 C\exp(C'|x|-\epsilon e^{2|x|}),\qquad \epsilon>0,
\end{equation}
with constants uniform on that compact set. Differentiation under the integral and dominated difference quotients prove $\Hh$-holomorphy. Contour shifting within the theta strip is justified by the same rapid tails. Applying it first to the exponential Fourier transform and then using evenness gives~\eqref{d:eq:translation-transform}. Evaluation at $a$ proves centering.

For the converse, first suppose $u\ne0$ and set $A=e^{-2|u|}\cos(2v)>0$. In~\eqref{d:eq:H-R}, one summand has strictly slower exponential decay as $x\to+\infty$ than the other. Equation~\eqref{d:eq:complex-asymptotic} therefore gives
\begin{equation}\label{d:eq:nonadmission-u}
 \frac{|H_b(x)|^2}{W(x)}
 \sim C_b e^{5x-\pi(2A-1)e^{2x}},\qquad C_b>0.
\end{equation}
If $2A\le1$, the right side is not integrable. This proves nonadmission when the real part is nonzero.

If $u=0$, the leading terms have equal magnitude and must be treated together. Realness of the theta series gives
\begin{equation}\label{d:eq:oscillatory-tail}
 H_{iv}(x)=4\pi^2e^{9x/2-\pi\cos(2v)e^{2x}}
 \left[\cos\bigl(9v/2-\pi\sin(2v)e^{2x}\bigr)+O(e^{-2x})\right].
\end{equation}
Outside $\Omega$ but within the strip, $|v|\ge\pi/6$ and $\sin(2v)\ne0$. Set $t=e^{2x}$. On a fixed positive-length subinterval of every sufficiently late period in $t$, the absolute value of the cosine is at least $3/4$ and the error is at most $1/4$. The weighted integral over these intervals is bounded below by a positive constant times
\begin{equation}\label{d:eq:oscillatory-divergence}
 \int_{\mathcal P} t^{3/2}\exp\bigl(\pi[1-2\cos(2v)]t\bigr)\,dt,
\end{equation}
where $\mathcal P$ is that union of periodic subintervals. This integral diverges even when $\cos(2v)=1/2$. Hence the boundary points $b=\pm i\pi/6$, as well as the remaining excluded imaginary parameters, are not admitted. No pointwise positive lower bound for the oscillatory factor has been used.

Finally $\kappa\in\Hh$, so subtracting its scalar multiple cannot turn a nonmember into a member. This proves the identical admission criterion for $R_b$.
\end{proof}

\begin{theorem}[Translation span and derivative span]\label{d:thm:span}
For every real interval $U$ containing zero in its interior and contained in $\Omega$,
\begin{equation}\label{d:eq:span-equality}
 \Nn=\cl_{\Hh}\{R_b:b\in U\}
     =\cl_{\Hh}\{R_b:b\in\Omega\}.
\end{equation}
\end{theorem}
\begin{proof}
The lens is connected: its defining condition can be written
$2|u|-\log\cos(2v)<\log2$, a convex sublevel condition in $|v|<\pi/4$.
Holomorphic differentiation at zero gives
\[
 \left.\frac{d^{2j}}{db^{2j}}R_b\right|_{b=0}=r_j,
\]
while all odd derivatives and the zeroth derivative vanish. Hilbert-valued finite differences along $U$ therefore place every $r_j$ in the first closed span in~\eqref{d:eq:span-equality}.

Conversely, let a bounded linear functional on $\Hh$ annihilate all the $r_j$. Its value on $R_b$ is holomorphic on $\Omega$ with all Taylor coefficients at zero equal to zero. The scalar identity theorem makes it vanish throughout $\Omega$. Taking annihilators proves the reverse inclusion.
\end{proof}

\begin{lemma}[The nearest boundary]\label{d:lem:radius}
The distance from zero to $\partial\Omega$ is
\begin{equation}\label{d:eq:a-star}
 a_*=\frac{\log2}{2}.
\end{equation}
In particular every closed disk $|b|\le R<a_*$ lies in $\Omega$.
\end{lemma}
\begin{proof}
By symmetry it suffices to consider the first-quadrant boundary. Set $t=2v$, so $2u=\log(2\cos t)$ and $0\le t\le\pi/3$. Four times the squared distance is
$f(t)=\log^2(2\cos t)+t^2$. For $0<t<\pi/3$ its derivative has the sign of
\[
 h(t)=t\cot t-\log(2\cos t).
\]
Here $h(0+)=1-\log2>0$, and
\[
 h'(t)=\frac{\sin t-t\cos t}{\sin^2t\cos t}>0.
\]
Thus $f$ is increasing, and its minimum is at $t=0$.
\end{proof}

\begin{remark}[No global weighted convolution pairing]\label{d:rem:pairing}
The restriction on real translations is a genuine norm issue. Set
$u_*(x)=\sqrt{W(x)}/(1+|x|)$ and $v_*=u_*/W$. Then $u_*\in\Hh$. For any real $b\ge a_*$, the positive-tail product $v_*(x)\kappa(x-b)$ is comparable to
\[
 (1+x)^{-1}\exp\!\left(\tfrac52x-\pi[e^{-2b}-1/2]e^{2x}\right),
\]
so its integral diverges. Replacing $u_*$ by $Cu_*$ changes $v_*$ by a multiple of $2c$; that correction has a finite pairing with the translated theta function and does not remove the divergence. This is not a counterexample to~\eqref{d:eq:ideal-synthesis}: $Cu_*$ is not asserted to annihilate $\Nn$.
\end{remark}

\subsection{Strict type synthesis and the one ray barrier}\label{d:sec:type}

An entire function has exponential type strictly below $a_*$ here if
$|Q(z)|\le M e^{\tau|z|}$ on $\C$ for some $M<\infty$ and $0\le\tau<a_*$. This is a global bound, not merely a bound on one ray.

\begin{theorem}[Strict type partial synthesis]\label{d:thm:type}
Let $Q(z)=\sum_{j\ge0}a_{2j}z^{2j}$ be even entire of exponential type strictly below $a_*$. Then
\begin{equation}\label{d:eq:derivative-series}
 u_Q=\sum_{j\ge0}(-1)^j a_{2j}\kappa^{(2j)}
\end{equation}
converges absolutely in $\Hh$ and satisfies $F_{u_Q}=\Xi Q$. Its centered part belongs to $\Nn$ and is the $\Hh$-limit of
\begin{equation}\label{d:eq:centered-series}
 \sum_{j=1}^J(-1)^j a_{2j}r_j.
\end{equation}
Consequently $Q(a)=0$ implies $u_Q\in\Nn$. In particular every $u\in\Ii$ whose quotient has this strict global type bound belongs to $\Nn$.
\end{theorem}
\begin{proof}
Choose $0<\tau<R<a_*$ with $|Q(z)|\le M e^{\tau|z|}$, increasing an available exponent if necessary. Lemma~\ref{d:lem:radius} and the Banach-valued Cauchy estimate applied to $H_b$ give
\begin{equation}\label{d:eq:Banach-Cauchy}
 \norm{\kappa^{(2j)}}_{\Hh}\le M_R(2j)!R^{-2j},\qquad
 M_R=\sup_{|b|=R}\norm{H_b}_{\Hh}<\infty.
\end{equation}
Scalar Cauchy estimates at radius $n/\tau$, together with the positive-real Stirling estimate~\cite[\S5.11]{DLMF}, give
\begin{equation}\label{d:eq:scalar-Cauchy}
 |a_n|n!\le M'\sqrt{n+1}\,\tau^n\qquad(n\ge0).
\end{equation}
The $n=0$ term is included by enlarging $M'$. Combining the estimates yields a summable bound
$M'M_R\sqrt{2j+1}(\tau/R)^{2j}$ for the norms of the terms in~\eqref{d:eq:derivative-series}.

Bounded transform evaluations and~\eqref{d:eq:derivative-transforms} identify the transform of its sum with $\Xi Q$. Applying $C$ to its partial sums gives~\eqref{d:eq:centered-series}; equivalently their transforms are
$\Xi[P_J(z)-P_J(a)]$, where $P_J=\sum_{j=0}^J a_{2j}z^{2j}$. Since $C$ is bounded, they converge to $Cu_Q$. If $Q(a)=0$, then $Cu_Q=u_Q$. Fourier injectivity gives the last assertion.
\end{proof}

The same constant appears in a necessary bound for every divisible transform, but with no strict margin and only on the imaginary axis.

\begin{proposition}[Imaginary axis growth bound]\label{d:prop:ray}
As $y\to+\infty$,
\begin{equation}\label{d:eq:gamma-saddle}
 \int_\R W(x)\cosh^2(yx)\,dx
 =(1+o(1))\pi^{-y}\Gamma(y+2),
\end{equation}
and
\begin{equation}\label{d:eq:gamma-ratio}
 \frac{\sqrt{\pi^{-y}\Gamma(y+2)}}{\Xi(iy)}
 =(\sqrt2+o(1))\frac{2^{y/2}}{y}.
\end{equation}
Thus there is a finite constant $C$ such that every $u\in\Ii$, with $F_u=\Xi Q$, satisfies
\begin{equation}\label{d:eq:ray-bound}
 |Q(iy)|\le C\norm{u}_{\Hh}\frac{2^{y/2}}{y}
\end{equation}
for all sufficiently large $y$. The same bound holds at $-iy$.
\end{proposition}
\begin{proof}
By evenness and $\cosh^2(yx)=(e^{2yx}+2+e^{-2yx})/4$, the only leading term is
$\frac12\int_0^\infty W(x)e^{2yx}\,dx$. Substitute~\eqref{d:eq:W-asymptotic} and then $t=e^{2x}$. The leading expression becomes
\[
 \pi^2\int_1^\infty t^{y+1}e^{-\pi t}\,dt
 \sim\pi^{-y}\Gamma(y+2).
\]
Here the factor $1/2$ from evenness and $\cosh^2$, the factor $4\pi^2$ in $W$, and $dx=dt/(2t)$ give exactly $\pi^2$. To justify the asymptotic, fix a large $X$. On $x\ge X$ the relative error in~\eqref{d:eq:W-asymptotic} is uniformly small; on $[0,X]$ the integral is $O_X(e^{2yX})$, negligible against $\Gamma(y+2)\pi^{-y}$. The two remaining terms of $\cosh^2$ are also negligible. This proves~\eqref{d:eq:gamma-saddle}.

For $s=y+1/2>1$, Lemma~\ref{d:lem:Fourier} and its symmetry give
\[
 \Xi(iy)=\xi(s)=\frac{s(s-1)}2\pi^{-s/2}\Gamma(s/2)\zeta(s)>0.
\]
The Dirichlet series shows $\zeta(s)\to1$. Positive-real Stirling applied to the two Gamma factors gives~\eqref{d:eq:gamma-ratio}. Finally Cauchy--Schwarz in the defining transform gives
$|F_u(iy)|\le\norm{u}_{\Hh}(\int W\cosh^2(yx))^{1/2}$; division by $\Xi(iy)$ proves~\eqref{d:eq:ray-bound}.
\end{proof}

Equation~\eqref{d:eq:ray-bound} permits growth at the critical exponential rate $a_*$. It neither controls all directions nor characterizes the norm of the quotient. In particular it cannot replace the strict global bound in Theorem~\ref{d:thm:type}. The next section establishes a full directional lens bound, but not a strict circular type bound for every quotient.

\subsection{Global quotient growth and the closed lens}\label{d:sec:global-growth}

The one-ray estimate of Proposition~\ref{d:prop:ray} is only a necessary consequence of membership. We now obtain a global necessary estimate with the exact support of the closed translation lens. Set
\begin{equation}\label{d:eq:closed-lens}
 L=\overline\Omega,\qquad
 h_L(w)=\sup_{b\in L}\operatorname{Re}(bw),\qquad
 H_L(z)=h_L(iz)=\sup_{b\in L}\operatorname{Re}(ibz).
\end{equation}
The bilinear convention in $h_L$, with no hidden conjugation, will also fix the rotation in the Borel representation.

\begin{lemma}[Preliminary order and local quotient bounds]\label{d:lem:quotient-order}
For $u\in\Ii$ and $Q=F_u/\Xi$,
\begin{equation}\label{d:eq:quotient-order}
 \log^+M_Q(R)=O_u(R\log(R+2)),
 \qquad M_Q(R)=\max_{|z|=R}|Q(z)|.
\end{equation}
For every fixed $K<\infty$ there is a constant $C_K$, independent of $u$, such that
\begin{equation}\label{d:eq:quotient-local}
 \sup_{|z|\le K}|Q(z)|\le C_K\norm{u}_{\Hh}.
\end{equation}
\end{lemma}
\begin{proof}
The zero source is trivial. Cauchy--Schwarz, the theta upper bound, and the substitution $t=e^{2|x|}$ bound the squared evaluation norm for $|z|\le R$, $R\ge1$, by a constant times $\pi^{-R}\Gamma(R+2)$. Positive-real Stirling therefore gives
$\log^+M_{F_u}(R)=O_u(R\log(R+2))$. The same holds for $\Xi=F_\kappa$.

For an entire function $f$, let
\[
 m_R^\pm(f)=\frac1{2\pi}\int_0^{2\pi}\log^\pm|f(Re^{it})|\,dt.
\]
Jensen's formula, using $\Xi(0)\ne0$, gives
$m_R^-(\Xi)\le m_R^+(\Xi)-\log|\Xi(0)|$. Zeros on the circle cause only integrable logarithmic singularities; the same inequality follows by a limit through radii avoiding them. Almost everywhere on a circle,
\[
 \log^+|Q|\le\log^+|F_u|+\log^-|\Xi|.
\]
Consequently $m_R^+(Q)=O_u(R\log(R+2))$. Since $Q$ is entire, $\log^+|Q|$ is subharmonic. The Poisson kernel for the disk of radius $2R$ is at most $3$ at points of modulus at most $R$, so its Poisson majorant yields
\[
 \log^+M_Q(R)\le3m_{2R}^+(Q).
\]
This proves~\eqref{d:eq:quotient-order}. No lower bound for $\Xi$ near its moving zeros was used.

For~\eqref{d:eq:quotient-local}, choose one radius $R>K$ whose circle contains no zero of $\Xi$, possible by discreteness. Bounded Fourier evaluation and the positive minimum of $|\Xi|$ on that fixed circle give
$|F_u/\Xi|\le C_K\norm{u}_{\Hh}$ there. The maximum principle for the entire quotient gives the bound inside, including at the removable zero points.
\end{proof}

The order bound~\eqref{d:eq:quotient-order} is not yet finite exponential type. A uniform estimate away from the horizontal zero strip and a sector maximum-principle argument provide the additional information.

\begin{lemma}[Uniform off-strip estimate]\label{d:lem:offstrip}
Write $z=x+iy$, $X=|x|$, $Y=|y|$, and $R=(X^2+Y^2)^{1/2}$. For $Y>0$ set
\begin{equation}\label{d:eq:G-function}
 G(X,Y)=\frac Y2\log\frac{2Y}{R}+\frac X2\arctan\frac XY.
\end{equation}
There is a constant $C$ independent of $u\in\Ii$ such that, uniformly for $Y\ge1$,
\begin{equation}\label{d:eq:offstrip}
 |Q(x+iy)|\le C\norm{u}_{\Hh}Y^{3/4}R^{-7/4}e^{G(X,Y)}
 \le C\norm{u}_{\Hh}R^{-1}e^{G(X,Y)}.
\end{equation}
\end{lemma}
\begin{proof}
Since $|\cos((x+iy)t)|\le\cosh(Yt)$, the numerator estimate follows from the same Gamma integral used in Proposition~\ref{d:prop:ray}:
\begin{equation}\label{d:eq:offstrip-numerator}
 |F_u(x+iy)|\le C\norm{u}_{\Hh}
 Y^{Y/2+3/4}\pi^{-Y/2}e^{-Y/2}\qquad(Y\ge1).
\end{equation}
This is a uniform upper bound, with constants enlarged on bounded $Y$.

For the denominator, the symmetries of $\xi$ permit evaluation of its modulus at $s=v+1/2$, where $v=Y+iX$. Thus
\[
 |\Xi(x+iy)|
 =\left|\frac12(v^2-1/4)\pi^{-v/2-1/4}
       \Gamma(v/2+1/4)\zeta(v+1/2)\right|.
\]
The Dirichlet series and the absolutely convergent Euler product in $\operatorname{Re}s>1$ give
\begin{equation}\label{d:eq:zeta-uniform}
 \zeta(3/2)^{-1}\le|\zeta(s)|\le\zeta(3/2)
 \qquad(\operatorname{Re}s\ge3/2).
\end{equation}
Indeed the upper bound follows from the series, while
$|1/\zeta(s)|\le\prod_p(1+p^{-\operatorname{Re}s})
 \le\prod_p(1-p^{-\operatorname{Re}s})^{-1}
 =\zeta(\operatorname{Re}s)$.
Absolute convergence justifies these products and their nonvanishing; no prime enumeration is involved. These standard formulas are recorded in~\cite[Eq.~25.2.11]{DLMF}.

The uniform complex Stirling formula~\cite[\S5.11]{DLMF} gives the shifted form
\[
 \Gamma(v/2+1/4)
 =\sqrt{2\pi}(v/2)^{v/2-1/4}e^{-v/2}
       \bigl(1+O(R^{-1})\bigr),
\]
uniformly as $R\to\infty$ with $|\arg v|\le\pi/2$. Taking moduli and absorbing the remaining compact region, where the Gamma function is continuous and nonzero, yields
\[
 |\Gamma(v/2+1/4)|\asymp
 R^{-1/4}(R/2)^{Y/2}
 \exp\left(-Y/2-\frac X2\arctan(X/Y)\right).
\]
Also $|v^2-1/4|\asymp R^2$ for $Y\ge1$. Therefore
\begin{equation}\label{d:eq:offstrip-denominator}
 |\Xi(x+iy)|\ge c R^{7/4}
 \left(\frac R{2\pi}\right)^{Y/2}
 \exp\left(-Y/2-\frac X2\arctan(X/Y)\right).
\end{equation}
Dividing~\eqref{d:eq:offstrip-numerator} by this estimate gives~\eqref{d:eq:offstrip}. The final inequality uses $Y\le R$. The reflection and conjugation steps apply to the denominator only; no real-coefficient assumption on $u$ or $Q$ is needed.
\end{proof}

\begin{lemma}[Support function of the lens]\label{d:lem:lens-support}
For $X=|x|$ and $Y=|y|$,
\begin{equation}\label{d:eq:lens-support}
 H_L(x+iy)=
 \begin{cases}
 G(X,Y),&X\le\sqrt3Y,\ Y>0,\\
 \pi X/6,&X\ge\sqrt3Y.
 \end{cases}
\end{equation}
The two formulas agree at their common boundary, and the second includes $Y=0$.
\end{lemma}
\begin{proof}
Coordinate-reflection symmetry of $L$ and its first-quadrant boundary give
\[
 H_L(x+iy)=\max_{0\le\beta\le\pi/6}
 \left[\frac Y2\log(2\cos(2\beta))+X\beta\right].
\]
For $Y>0$, the interior critical point satisfies $\tan(2\beta)=X/Y$. It is available precisely when $X\le\sqrt3Y$, and substitution gives $G(X,Y)$. Otherwise the maximum is attained at $\beta=\pi/6$, where the real coordinate of the lens is zero. The case $Y=0$ follows from the same maximum.
\end{proof}

\begin{theorem}[Global closed-lens bound]\label{d:thm:global-growth}
There is a constant $C$, independent of $u\in\Ii$, such that its entire quotient satisfies
\begin{equation}\label{d:eq:global-growth}
 |Q(z)|\le C\norm{u}_{\Hh}\frac{e^{H_L(z)}}{1+|z|}
 \qquad(z\in\C).
\end{equation}
In particular every such quotient has finite exponential type.
\end{theorem}
\begin{proof}
In the vertical cones $X\le\sqrt3Y$, $R\ge2$ implies $Y\ge1$, so Lemmas~\ref{d:lem:offstrip} and~\ref{d:lem:lens-support} prove the assertion there. Lemma~\ref{d:lem:quotient-order} handles the compact disk.

For the right horizontal sector $|\arg z|\le\pi/6$, put
$f(z)=zQ(z)e^{-\pi z/6}$. On its boundary rays, the same off-strip estimate gives
$|f(z)|\le C\norm{u}_{\Hh}$ outside a fixed disk; the local quotient bound covers the remaining segments. By~\eqref{d:eq:quotient-order},
$\log^+|f(z)|=O_u(R\log(R+2))$.
For $\epsilon>0$ consider
\[
 f_\epsilon(z)=f(z)e^{-\epsilon z^2}.
\]
In this sector $\operatorname{Re}z^2\ge|z|^2/2$. Thus on the outer circular arc of a truncated sector the damping dominates the preliminary growth and the maximum tends to zero as the radius tends to infinity. On the two radial boundaries the damping has modulus at most one, preserving the bound $C\norm{u}_{\Hh}$. The maximum principle and then $\epsilon\downarrow0$ give that bound for $f$ throughout the sector. This proves~\eqref{d:eq:global-growth} there for $|z|\ge1$, because $H_L(z)=\pi|\operatorname{Re}z|/6$. Evenness covers the left horizontal sector, and~\eqref{d:eq:quotient-local} covers the remaining compact disk.
\end{proof}

The horizontal value $\pi|x|/6$ is essential: the raw off-strip expression $G$ is not the lens support near the real axis. The explicit damping argument supplies the needed improvement without a minimum-modulus assumption or an unproved quotient-order assertion.

\subsection{Directional Borel continuation and a sharp common carrier}\label{d:sec:Borel}

Set $P(w)=Q(-iw)$. Then~\eqref{d:eq:global-growth} becomes
\begin{equation}\label{d:eq:P-growth}
 |P(w)|\le C\norm{u}_{\Hh}\frac{e^{h_L(w)}}{1+|w|}.
\end{equation}
For $P(w)=\sum_{n\ge0}p_nw^n$, define its exponential Borel transform initially by
\begin{equation}\label{d:eq:B-series}
 B_P(b)=\sum_{n\ge0}n!p_n b^{-n-1}.
\end{equation}
The series converges outside a sufficiently large disk: scalar Cauchy estimates as in~\eqref{d:eq:scalar-Cauchy} apply to every finite exponential-type bound. A compact convex set is called a \emph{carrier} here when this Borel transform continues as a single-valued holomorphic function to its complement and the contour inversion below holds around that set.

\begin{theorem}[Borel continuation and uniform boundary estimate]\label{d:thm:Borel}
The function $B_P$ has a single-valued holomorphic continuation to $\C\setminus L$. For any positively oriented rectifiable Jordan contour $\Gamma$ in $\C\setminus L$ enclosing $L$ with winding number one,
\begin{equation}\label{d:eq:B-inversion}
 P(w)=\frac1{2\pi i}\int_\Gamma B_P(b)e^{bw}\,db.
\end{equation}
With $\delta=\operatorname{dist}(b,L)>0$,
\begin{equation}\label{d:eq:B-bound}
 |B_P(b)|\le C\norm{u}_{\Hh}
       \int_0^\infty\frac{e^{-\delta t}}{1+t}\,dt
 \le C'\norm{u}_{\Hh}\log(1+\delta^{-1}).
\end{equation}
The set $L$ is the smallest common closed convex carrier for all the rotated quotients $P$. Individual carriers may be smaller.
\end{theorem}
\begin{proof}
For a real direction $\theta$, the integral
\begin{equation}\label{d:eq:directional-Laplace}
 B_\theta(b)=e^{i\theta}\int_0^\infty
     P(te^{i\theta})e^{-bte^{i\theta}}\,dt
\end{equation}
converges and is holomorphic in the half-plane
\begin{equation}\label{d:eq:separation-halfplane}
 \operatorname{Re}(be^{i\theta})>h_L(e^{i\theta}).
\end{equation}
This follows directly from~\eqref{d:eq:P-growth}, with locally uniform majorants for derivatives. Far enough into that half-plane, termwise integration of the power series is absolutely justified and gives~\eqref{d:eq:B-series}.

Convex separation shows that these half-planes cover $\C\setminus L$. To verify compatibility, fix $b$ in an overlap. The separating directions $v$ satisfy
$\operatorname{Re}((b-k)v)>0$ for every $k\in L$. They form an open convex cone with no opposite pair. Two allowable unit directions can therefore be joined by an angular arc of length less than $\pi$ within that cone. On the compact arc the separation gap is uniformly positive. Applying Cauchy's theorem to the integrand $P(w)e^{-bw}$ between the two rays, its outer-arc integral is $O(e^{-cT})$ as the radius $T$ tends to infinity, by~\eqref{d:eq:P-growth}; the inner arc tends to zero with its radius. The two Laplace integrals agree. Thus the local functions patch to one holomorphic continuation.

On a large circle, multiplying~\eqref{d:eq:B-series} by the exponential series and integrating selects $p_nw^n$, proving~\eqref{d:eq:B-inversion}. Cauchy deformation through $\C\setminus L$ then proves it for the stated contours.

At a nearest point of the convex set $L$, choose a unit separating direction whose gap in~\eqref{d:eq:separation-halfplane} equals $\delta$. Formula~\eqref{d:eq:directional-Laplace} gives the first bound in~\eqref{d:eq:B-bound}. For $0<\delta\le1$, split the integral at $1$ and $\delta^{-1}$ to obtain $C(1+\log\delta^{-1})$. For $\delta\ge1$, it is at most $\delta^{-1}$. These are bounded by a constant times $\log(1+\delta^{-1})$, proving the second estimate.

For sharpness, take any nonzero $b_0\in\Omega$. The quotient of $R_{b_0}$ rotates to
$P_{b_0}(w)=\cosh(b_0w)-\cosh(b_0/2)$. Its Borel transform is
\[
 \frac{b}{b^2-b_0^2}-\frac{\cosh(b_0/2)}b,
\]
which has nonzero poles at both $b_0$ and $-b_0$. Any common closed convex carrier must contain all these pairs and therefore their closed convex hull $L$. Conversely the continuation and inversion just proved show that $L$ is a carrier for every quotient. A polynomial quotient can have carrier $\{0\}$, so no individual equality with $L$ is asserted.
\end{proof}

The rotation is important. In the original variable $Q(z)=P(iz)$, the exponential in~\eqref{d:eq:B-inversion} is $e^{ibz}$. Thus the translation parameters lie in $L$, whereas an exponential carrier expressed directly for $Q$ is rotated to $iL$.

\subsection{Compact interior Abel regularization}\label{d:sec:Abel-maps}

For $0<r<1$ the desired centered dilation is
$\Xi(z)[Q(rz)-Q(ra)]$. The Borel representation now proves its actual weighted admission for every $u\in\Ii$, not just for small circular type.

\begin{theorem}[Interior Abel maps]\label{d:thm:Abel-maps}
For every $0<r<1$ there is a unique linear operator $A_r:\Ii\to\Nn$ such that
\begin{equation}\label{d:eq:Abel-transform}
 F_{A_ru}(z)=\Xi(z)\bigl[Q(rz)-Q(ra)\bigr],
 \qquad F_u=\Xi Q.
\end{equation}
It is bounded, compact, and injective, and
\begin{equation}\label{d:eq:Abel-norm}
 \norm{A_r}_{\Ii\to\Hh}
 \le C(1-r)^{-5/4}\log\frac e{1-r}.
\end{equation}
Moreover
\begin{equation}\label{d:eq:Abel-semigroup}
 A_rA_s=A_{rs},\qquad A_rr_j=r^{2j}r_j,\qquad
 \overline{\Ran A_r}^{\Hh}=\Nn.
\end{equation}
The bound~\eqref{d:eq:Abel-norm} is a coarse interior estimate. Theorem~\ref{d:thm:uniform-endpoint} will replace its endpoint loss by a uniform bound using actual weighted admission.
\end{theorem}
\begin{proof}
Let $P_r(w)=P(rw)$, whose Borel transform is
$B_r(b)=r^{-1}B_P(b/r)$ outside $rL$. Put $\rho=(1+r)/2$ and let
$\Gamma_r=\rho\partial L$ with positive orientation. This is a piecewise smooth Jordan contour inside $\Omega$ and outside $rL$. Inversion and the evenness of $P_r$ give
\[
 Q(rz)=\frac1{2\pi i}\int_{\Gamma_r}B_r(b)\cos(bz)\,db.
\]
Define
\begin{equation}\label{d:eq:Abel-Bochner}
 A_ru=\frac1{2\pi i}\int_{\Gamma_r}B_r(b)R_b\,db.
\end{equation}
The contour is compactly contained in $\Omega$, so the integrand is continuous in the $\Hh$ norm and the Bochner integral exists. Every $R_b$ belongs to the closed linear subspace $\Nn$, so the integral does too. Bounded transform evaluations pass under the integral. Formula~\eqref{d:eq:translation-transform} then gives~\eqref{d:eq:Abel-transform}, including its exact centering. Fourier injectivity proves uniqueness and independence of the chosen intermediate contour.

We give quantitative estimates before proving compactness. For $b=u+iv\in\Omega$, put
\[
 \epsilon(b)=2e^{-2|u|}\cos(2v)-1.
\]
Uniform theta estimates on the compact closed lens give
\begin{equation}\label{d:eq:translation-blowup}
 \norm{R_b}_{\Hh}\le C\bigl(1+\epsilon(b)^{-5/4}\bigr).
\end{equation}
Indeed the squared positive-tail integrand is bounded by
$C\exp(5x-\pi\epsilon(b)e^{2x})$. With $t=e^{2x}$ its integral is at most a constant times
$\int_1^\infty t^{3/2}e^{-\pi\epsilon(b)t}\,dt
 \le C'\epsilon(b)^{-5/2}$.
The compact $x$ region and the centering term are uniformly bounded. Taking square roots proves~\eqref{d:eq:translation-blowup}.

The convex function $f(u+iv)=2|u|-\log\cos(2v)$ has $f(0)=0$ and equals $\log2$ on $\partial L$. Therefore, for $b\in\rho\partial L$,
\begin{equation}\label{d:eq:radial-epsilon}
 \epsilon(b)\ge 2^{1-\rho}-1\ge(1-\rho)\log2.
\end{equation}
At $b=\rho b_0$, choose a unit supporting direction at $b_0\in\partial L$. Since $L$ contains the closed disk of radius $a_*$, its support in that direction is at least $a_*$. The separation from $rL$ is consequently at least
\begin{equation}\label{d:eq:contour-separation}
 \delta=\operatorname{dist}(b,rL)\ge a_*(\rho-r).
\end{equation}

Applying the directional Laplace estimate directly to $P_r$ gives
\begin{equation}\label{d:eq:scaled-Borel}
 |B_r(b)|\le C\norm{u}_{\Hh}
       \int_0^\infty\frac{e^{-\delta t}}{1+rt}\,dt.
\end{equation}
For $r\ge1/2$,~\eqref{d:eq:contour-separation} bounds this by
$C\norm{u}_{\Hh}\log(e/(1-r))$. For $r\le1/2$, $\delta\ge a_*/4$, and dropping the denominator gives a uniform bound. This separate estimate avoids a spurious small-$r$ loss from $r^{-1}B_P(b/r)$. The contour lengths are uniformly bounded. Substituting~\eqref{d:eq:translation-blowup}--\eqref{d:eq:scaled-Borel} into~\eqref{d:eq:Abel-Bochner} proves~\eqref{d:eq:Abel-norm}.

For fixed $r$, partition the contour into short arcs and replace $R_b$ on each arc by one representative $R_{b_j}$. The resulting operator has finite rank: its coefficients are integrals of the bounded linear functionals $u\mapsto B_r(b)$. The uniform bound~\eqref{d:eq:scaled-Borel} and uniform $\Hh$ continuity of $R_b$ on the contour make the error tend to zero in operator norm. Thus $A_r$ is compact, not merely pointwise approximated by finite-rank maps.

For the composition law, the quotient of $A_su$ is
$Q_s(z)=Q(sz)-Q(sa)$. Hence
\[
 Q_s(rz)-Q_s(ra)=Q(sr z)-Q(sr a),
\]
and Fourier injectivity gives $A_rA_s=A_{rs}$. If $A_ru=0$, then $Q(rz)$ is constant. Since $r>0$, $Q$ itself is constant, and $Q(a)=0$ forces $u=0$. Finally~\eqref{d:eq:derivative-transforms} gives $A_rr_j=r^{2j}r_j$. The range is contained in $\Nn$ and contains every finite radical direction, proving its stated closure.
\end{proof}

For every $u\in\Ii$,~\eqref{d:eq:Abel-transform} tends compact-open to $F_u$ as $r\uparrow1$. Neither fixed-parameter compactness nor the semigroup law alone upgrades this to boundedness in $\Hh$. The actual-weight argument in Sections~\ref{d:sec:source-recovery}--\ref{d:sec:criterion} supplies the stronger estimate.

\subsection{A sharp boundary family and strong Abel convergence}\label{d:sec:boundary-family}

The strict-type theorem does not exhaust the synthesized quotients. A family whose carriers reach the boundary of $L$ has a sharp weighted admission threshold.

For $b_0\in\partial L$ and real $\alpha>0$, set
\begin{align}
 Q_{\alpha,b_0}(z)
 &=\int_0^1(1-t)^{\alpha-1}
       \bigl[\cos(tb_0z)-\cosh(tb_0/2)\bigr]\,dt,\label{d:eq:boundary-quotient}\\
 u_{\alpha,b_0}(x)
 &=\int_0^1(1-t)^{\alpha-1}R_{tb_0}(x)\,dt.\label{d:eq:boundary-source}
\end{align}
The second formula is initially pointwise, not an assumed $\Hh$ integral. For fixed $x$, the parameter segment lies inside the theta analytic strip, including its endpoint, so it defines a continuous integrand apart from the integrable scalar weight. The theta bounds uniformly on this compact segment also show that $u_{\alpha,b_0}$ is in $L^1(\R,e^{B|x|}dx)$ for every finite $B$. Indeed the slowest real exponential coefficient is $1/2$, still giving double-exponential ordinary decay. Fubini therefore proves
\begin{equation}\label{d:eq:boundary-transform}
 F_{u_{\alpha,b_0}}=\Xi Q_{\alpha,b_0},\qquad m(u_{\alpha,b_0})=0
\end{equation}
even before weighted admission is decided.

\begin{theorem}[Sharp boundary threshold]\label{d:thm:boundary-family}
For every $b_0\in\partial L$ and $\alpha>0$,
\begin{equation}\label{d:eq:boundary-threshold}
 u_{\alpha,b_0}\in\Hh\quad\Longleftrightarrow\quad\alpha>\frac54.
\end{equation}
In the admitted range, $u_{\alpha,b_0}\in\Nn$ and
\begin{equation}\label{d:eq:boundary-Abel}
 A_ru_{\alpha,b_0}\longrightarrow u_{\alpha,b_0}
 \quad\hbox{in }\Hh\quad(r\uparrow1).
\end{equation}
In the rotated variable $P_{\alpha,b_0}(w)=Q_{\alpha,b_0}(-iw)$, the segment $[-b_0,b_0]$ is a carrier and its endpoints are genuine boundary contacts.
\end{theorem}
\begin{proof}
By~\eqref{d:eq:translation-blowup} and~\eqref{d:eq:radial-epsilon},
\begin{equation}\label{d:eq:boundary-majorant}
 \norm{R_{tb_0}}_{\Hh}\le C(1-t)^{-5/4}\qquad(0\le t<1).
\end{equation}
When $\alpha>5/4$, multiplication by $(1-t)^{\alpha-1}$ gives an integrable norm majorant. Thus~\eqref{d:eq:boundary-source} is an $\Hh$ Bochner integral of vectors in $\Nn$, so it belongs to $\Nn$. The centered dilation of~\eqref{d:eq:boundary-quotient} is the same integral with $R_{rtb_0}$ in place of $R_{tb_0}$. Since $1-rt\ge1-t$, the same majorant is valid uniformly for $r\le1$. For fixed $t<1$, holomorphy gives $R_{rtb_0}\to R_{tb_0}$ in $\Hh$. Bochner dominated convergence and~\eqref{d:eq:Abel-transform} prove~\eqref{d:eq:boundary-Abel}.

For necessity, use $R_b=R_{-b}$ to reduce first to $b_0=u+iv$ with $u>0$. Put $q_0=e^{-2b_0}$. The boundary equation implies
\begin{equation}\label{d:eq:boundary-q}
 \operatorname{Re}q_0=\frac12,\qquad
 \operatorname{Re}(b_0q_0)
 =e^{-2u}\bigl[u\cos(2v)+v\sin(2v)\bigr]>0.
\end{equation}
In the branch $\kappa(x-tb_0)$ of~\eqref{d:eq:boundary-source}, set $T=e^{2x}$ and $\delta=1-t$. Uniformly near $\delta=0$,
\[
 e^{-2tb_0}=q_0+2b_0q_0\delta+O(\delta^2),\qquad
 e^{-9tb_0/2}=e^{-9b_0/2}\bigl(1+O(\delta)\bigr).
\]
The real part of $e^{-2tb_0}$ strictly decreases for $0\le t\le1$, since its derivative is
$-2e^{-2tu}[u\cos(2tv)+v\sin(2tv)]<0$.
Thus near the endpoint its increment above $\operatorname{Re}q_0$ is at least $c\delta$, and away from the endpoint it has a fixed positive gap.

Apply~\eqref{d:eq:complex-asymptotic}, split the integral at a small fixed $\delta$, and rescale $\delta=s/T$ in the endpoint piece. Dominated convergence against $s^{\alpha-1}e^{-cs}$ gives
\begin{equation}\label{d:eq:boundary-tail}
 u_{\alpha,b_0}(x)
 =C_{\alpha,b_0}e^{(9/2-2\alpha)x}
       e^{-\pi q_0e^{2x}}\bigl(1+o(1)\bigr),
\end{equation}
where
\begin{equation}\label{d:eq:boundary-constant}
 C_{\alpha,b_0}
 =2\pi^2 e^{-9b_0/2}
      \frac{\Gamma(\alpha)}{(2\pi b_0q_0)^\alpha}\ne0.
\end{equation}
The power is defined in the right half-plane from~\eqref{d:eq:boundary-q}. The Laplace identity
$\int_0^\infty s^{\alpha-1}e^{-As}\,ds
 =\Gamma(\alpha)A^{-\alpha}$ for $\operatorname{Re}A>0$
follows from the real Gamma integral by holomorphy.

No omitted term cancels~\eqref{d:eq:boundary-tail}. For the opposite branch, the real coefficient $e^{2tu}\cos(2tv)$ stays strictly above $1/2$ on the entire compact interval: for $t<1$ this follows from $tb_0\in\Omega$, and at $t=1$ it equals $e^{4u}/2>1/2$. The centered unshifted kernel is faster still. Consequently~\eqref{d:eq:W-asymptotic} yields
\begin{equation}\label{d:eq:boundary-weighted-tail}
 \frac{|u_{\alpha,b_0}(x)|^2}{W(x)}
 \sim c_{\alpha,b_0}e^{(5-4\alpha)x},
 \qquad c_{\alpha,b_0}>0.
\end{equation}
This is not integrable for $\alpha\le5/4$, including equality.

It remains to treat $b_0=\pm i\pi/6$. The two endpoint branches now have conjugate coefficients of equal magnitude. The same rescaling is valid because $\operatorname{Re}(b_0e^{-2b_0})>0$. Their sum has the form
\begin{equation}\label{d:eq:boundary-tip}
 2|C|e^{(9/2-2\alpha)x}e^{-\pi T/2}
 \left[\cos\left(\frac{\sqrt3\pi}{2}T+\phi\right)+o(1)\right],
 \qquad C\ne0.
\end{equation}
This is an additive error inside the bracket, not a relative estimate at the zeros of the cosine. On a fixed positive proportion of each sufficiently late $T$-period the cosine has magnitude at least $1/2$ and the error is at most $1/4$. After division by $W$ and $dx=dT/(2T)$, the weighted integral over these subintervals dominates a constant times
\[
 \int_{\mathcal P} T^{3/2-2\alpha}\,dT.
\]
It diverges for $\alpha\le5/4$, with logarithmic divergence at equality. This proves the threshold at the interference tips as well.

Finally, $P_{\alpha,b_0}$ is an integral of $\cosh(tb_0w)$ and a constant with an integrable scalar weight. Its Borel transform can therefore be represented, outside $[-b_0,b_0]$, by integrating the corresponding rational transforms. This proves that segment is a carrier. In every fixed direction with $\operatorname{Re}(b_0w)>0$, endpoint rescaling gives a nonzero leading term
\[
 \frac{\Gamma(\alpha)}2\,\frac{e^{b_0w}}{(b_0w)^\alpha}
 \quad(|w|\to\infty),
\]
where the opposite exponential and the centered constant are smaller. Symmetry gives the analogous term at the other endpoint. The directional exponential growth therefore attains the support of the segment, so its endpoints cannot be removed from the convex carrier. In particular, taking $b_0=i\pi/6$ gives type $\pi/6>a_*$ in the original variable, beyond the strict-type subclass of Theorem~\ref{d:thm:type}.
\end{proof}

\begin{proposition}[Carrier growth is not an admission criterion]\label{d:prop:sinc-warning}
There exists an even entire $Q_*$ with $Q_*(a)=0$ and
\begin{equation}\label{d:eq:sinc-growth}
 |Q_*(z)|\le C\frac{e^{H_L(z)}}{1+|z|}
\end{equation}
whose rotated carrier lies in $L$, but $\Xi Q_*$ is not the transform of any vector in $\Hh$.
\end{proposition}
\begin{proof}
Take $b_0=a_*$ and $\alpha=1$ in~\eqref{d:eq:boundary-quotient}. Then
\begin{equation}\label{d:eq:sinc-quotient}
 Q_*(z)=\frac{\sin(a_*z)}{a_*z}
        -\frac{\sinh(a_*/2)}{a_*/2},
\end{equation}
with its removable value at zero. Its carrier is the segment $[-a_*,a_*]$ in the rotated variable. The sinc term satisfies~\eqref{d:eq:sinc-growth} by
$|\sin(a_*z)|\le e^{a_*|\operatorname{Im}z|}$ and
$a_*|\operatorname{Im}z|\le H_L(z)$, with a compact-disk adjustment. The constant term is also absorbed: Lemma~\ref{d:lem:radius} implies $H_L(z)\ge a_*|z|$.

The unique ordinary $L^1$ source in~\eqref{d:eq:boundary-source} is not in $\Hh$, since~\eqref{d:eq:boundary-weighted-tail} gives squared weighted tail comparable to $e^x$. If an $\Hh$ source existed, it would also be in $L^1$ and have the same Fourier transform, contradicting Fourier uniqueness.

For each $r<1$, the centered function $Q_*(rz)-Q_*(ra)$ has strict exponential type $ra_*<a_*$, so Theorem~\ref{d:thm:type} supplies an $\Nn$ source. These sources cannot have a bounded sequence as $r\uparrow1$: a weakly convergent subsequence and bounded evaluations would provide the forbidden $\Hh$ inverse of $\Xi Q_*$.
\end{proof}

This example disproves only sufficiency of the carrier and growth condition for weighted admission. It is not an element of $\Ii$ and not a counterexample to $\Ii=\Nn$ or to endpoint boundedness for sources already in $\Ii$.


\subsection{Uniform Abel bounds and full weighted synthesis}\label{d:sec:criterion}

\begin{theorem}[The actual Hilbert endpoint]\label{d:thm:uniform-endpoint}
The centered Abel operators satisfy
\begin{equation}\label{d:eq:uniform-endpoint}
 \sup_{3/4\le r<1}\norm{A_r}_{\Ii\to\Hh}<\infty.
\end{equation}
\end{theorem}
\begin{proof}
For $u\in\Ii$, use its fixed scalar functional $\mu$ to define the uncentered source
\[
 v_r(x)=\mu(\kappa(x+rb)).
\]
The parameter set $r\Gamma$, $3/4\le r<1$, remains compactly inside the theta strip. Hence $v_r$ is even and has all ordinary exponential moments, by the same argument as Lemma~\ref{d:lem:scalar-source}. Its transform is $\Xi(z)Q(rz)$ and its leading Laplace source is exactly $f_r$ from~\eqref{d:eq:fr-definition}. The exact theta identity~\eqref{d:eq:U-H-identity} applies with $f_r$.

Theorem~\ref{d:thm:weighted-kernel}, Theorem~\ref{d:thm:recovery}, and Lemma~\ref{d:lem:forward-reconstruction} give
\begin{equation}\label{d:eq:vr-bound}
 \norm{v_r}_{\Hh}
 \le C\norm{f_r}_{E_\pi}
 \le C'\norm{f}_E
 \le C''\norm{u}_{\Hh}.
\end{equation}
In particular this proves actual weighted admission of $v_r$; it was not inferred from a formal transform or from an exterior $\Hh$-valued contour integral.

Centering is exact:
\[
 2m(v_r)=Q(ra),\qquad Cv_r=v_r-Q(ra)\kappa.
\]
The scalar coefficient is bounded uniformly by $C\norm{u}_{\Hh}$ from~\eqref{d:eq:quotient-local}, since $|ra|\le1/2$, and $\norm{\kappa}_{\Hh}=1$. Thus~\eqref{d:eq:vr-bound} pays for the centering correction and bounds $Cv_r$ uniformly. Its transform is precisely~\eqref{d:eq:Abel-transform}, so Fourier injectivity identifies it with the already admitted $A_ru\in\Nn$. This proves~\eqref{d:eq:uniform-endpoint}.

The argument used actual membership $u\in\Hh$ in~\eqref{d:eq:actual-weight-comparison} and Theorem~\ref{d:thm:recovery}. It assumed neither positivity of $u$ nor a positive measure representation of its Borel functional. Only the auxiliary fractional-power kernel is positive.
\end{proof}

\begin{theorem}[Full multiplicity weighted synthesis]\label{d:thm:full-synthesis}
The full divisible ideal equals the derivative span closure:
\begin{equation}\label{d:eq:full-synthesis}
 \Ii=\Nn.
\end{equation}
For every $u\in\Ii$,
\begin{equation}\label{d:eq:strong-Abel}
 A_ru\longrightarrow u\quad\hbox{strongly in }\Hh
 \qquad(r\uparrow1).
\end{equation}
Consequently
\begin{equation}\label{d:eq:zerojet-completeness}
 \Kk=\Zz
 =\cl_{\Hh}\{J_{w,k}:\Xi(w)=0,\ 0\le k<\nu(w)\}.
\end{equation}
Equivalently, the generalized nullspace of $T|_{\Ii}$ is dense in $\Ii$, and its quotient $\Ii/\Nn$ is zero.
\end{theorem}
\begin{proof}
Fix $u\in\Ii$ and any sequence $r_n\uparrow1$ with $r_n\ge3/4$. Theorem~\ref{d:thm:uniform-endpoint} makes $A_{r_n}u$ bounded in $\Hh$. A subsequence has a weak limit $v$. Since $\Nn$ is a closed linear subspace, it is weakly closed, so $v\in\Nn$. Bounded transform evaluations and~\eqref{d:eq:Abel-transform} give
\[
 F_v(z)=\lim_n\Xi(z)[Q(r_nz)-Q(r_na)]
       =\Xi(z)[Q(z)-Q(a)]=F_u(z).
\]
Fourier injectivity yields $v=u$, so $\Ii\subset\Nn$. The reverse inclusion was proved in Proposition~\ref{d:prop:ideal}; this establishes~\eqref{d:eq:full-synthesis}.

On a finite radical combination $p$, the identity $A_rr_j=r^{2j}r_j$ gives $A_rp\to p$ in $\Hh$. This finite span is dense in $\Ii=\Nn$. With the uniform bound already proved,
\[
 \norm{A_ru-u}_{\Hh}
 \le C\norm{u-p}_{\Hh}
       +\norm{A_rp-p}_{\Hh}+\norm{p-u}_{\Hh}.
\]
First let $r\uparrow1$, then let $p$ approach $u$. This proves strong convergence, not merely compact-open convergence of transforms.

Taking orthogonal complements in $\Hh_c$ in Proposition~\ref{d:prop:ideal} gives~\eqref{d:eq:zerojet-completeness}. The generalized-nullspace statement follows from Theorem~\ref{d:thm:backward}, and the quotient assertion is immediate.
\end{proof}

The uniform bound precedes both limiting arguments. Neither fixed-$r$ compactness, the semigroup identity, nor a carrier estimate was substituted for it. Proposition~\ref{d:prop:sinc-warning} is not contradicted: its source is outside $\Hh$, so the actual-weight recovery hypothesis fails. The carrier condition is still not a sufficient admission criterion.


\subsection{Stronger-topology duals and the same closed carrier}\label{d:sec:Gamma-dual-carrier}

Let
\begin{equation}\label{d:eq:jet-fin}
 \mathcal J_{\rm fin}=\operatorname{span}\{J_{w,k}:\Xi(w)=0,\ 0\le k<\nu(w)\},
 \qquad \Nn_0=\C\kappa\oplus\Nn.
\end{equation}
The sum is $\Hh$-orthogonal and closed, and $\Kk=\Nn_0^\perp$. Fix either $E={X}$ or $E=\mathcal G_{\mathsf D}$, with $\mathsf D$ fixed as above.

The theta tails and smooth positivity on compact intervals give
\begin{equation}\label{d:eq:W-derivative-bound}
 W^2\le C W,\qquad |W'|^2\le C W\quad\hbox{on }\R.
\end{equation}
For the second inequality, differentiating $W=\kappa/(2c)$ adds at most a fixed exponential factor to the theta derivative bound; dividing its square by the lower bound for $W$ leaves
$C\exp(C'|x|-\pi e^{2|x|})$ outside a compact interval. It is bounded there, and compact positivity handles the rest. A bounded logarithmic derivative of $W$ is neither used nor asserted.

For $z\in\C$ define
\[
 e_z(x)=W(x)\cos(zx),\qquad
 M_y^2=\int_\R W(x)\cosh^2(yx)\,dx\quad(y\ge0).
\]
Since $|\cos(zx)|,|\sin(zx)|\le\cosh(|\operatorname{Im}z|x)$, equations~\eqref{d:eq:W-derivative-bound} and~\eqref{d:eq:Gamma-E-upper} give
\begin{equation}\label{d:eq:dual-evaluation}
 \norm{e_z}_{\Hh}\le M_{|\operatorname{Im}z|},\qquad
 \norm{e_z}_{H^1}\le C(1+|z|)M_{|\operatorname{Im}z|},\qquad
 \norm{e_z}_E\le C_E(1+|z|)M_{|\operatorname{Im}z|}.
\end{equation}
The map $z\mapsto e_z$ is entire into $E$: on compact parameter sets, all fixed parameter derivatives add only fixed powers of $x$, dominated by the same theta tails in $\Hh$ and ordinary $H^1$. Difference quotients and~\eqref{d:eq:Gamma-E-upper} give $E$ holomorphy. In particular every jet in~\eqref{d:eq:jet-fin} belongs to both stronger spaces.

Let $\ell$ be a continuous conjugate-linear functional on $E$ annihilating $\mathcal J_{\rm fin}$, and put $L_\ell=\norm{\ell}_{E^*}$. Define the even tempered distribution
\begin{equation}\label{d:eq:dual-distribution}
 v(\varphi)=\ell(W\overline{\varphi_{\rm even}}),\qquad
 \varphi_{\rm even}(x)=\frac{\varphi(x)+\varphi(-x)}2,
 \quad\varphi\in\mathcal S(\R).
\end{equation}
It is linear in the test function. Multiplication by $W$ maps Schwartz functions continuously into $\Hh\cap H^1$, and~\eqref{d:eq:Gamma-E-upper} gives
$\norm{W\varphi}_E\le C_E(\sup|\varphi|+\sup|\varphi'|)$.
Thus~\eqref{d:eq:dual-distribution} genuinely defines a tempered distribution; no weighted function representative is yet assumed.

Set
\begin{equation}\label{d:eq:dual-entire}
 F(z)=\ell(e_{\overline z}).
\end{equation}
Conjugate linearity and the entire $E$-valued family make $F$ even entire. Its derivatives obey
$F^{(k)}(w)=\ell(J_{\overline w,k})$.
Conjugation invariance of the zero set, with multiplicities, therefore implies
\begin{equation}\label{d:eq:dual-quotient}
 F=\Xi Q,\qquad Q\text{ even entire},\qquad
 |F(z)|\le C_E L_\ell(1+|z|)M_{|\operatorname{Im}z|}.
\end{equation}
There is no centering assumption: $Q(a)$ may be nonzero, and the $\kappa$ direction must be retained in the eventual annihilator.

\begin{theorem}[Closed-lens bound for stronger-topology duals]\label{d:thm:dual-lens}
The quotient in~\eqref{d:eq:dual-quotient} satisfies
\begin{equation}\label{d:eq:dual-lens}
 |Q(z)|\le C_E L_\ell e^{H_L(z)}\qquad(z\in\C).
\end{equation}
This has the same exponential carrier as~\eqref{d:eq:global-growth}, but not its reciprocal-linear prefactor. No weighted admission of $v$ is used.
\end{theorem}
\begin{proof}
The zero functional is trivial. The theta evaluation bound and~\eqref{d:eq:dual-quotient} give
$\log^+M_F(R)=O_\ell(R\log(R+2))$, with the same bound for $\Xi$.
Jensen, using $\Xi(0)\ne0$, bounds the circle average of $\log^-|\Xi|$ by its positive average minus $\log|\Xi(0)|$. Since
$\log^+|Q|\le\log^+|F|+\log^-|\Xi|$, Poisson majorization from radius $2R$ to $R$ yields
\begin{equation}\label{d:eq:dual-order}
 \log^+M_Q(R)=O_\ell(R\log(R+2)).
\end{equation}
Circle zeros have integrable logarithmic singularities; limiting radii avoiding them give the same result. On any fixed compact disk, choose a larger zero-free circle for $\Xi$ and divide~\eqref{d:eq:dual-quotient} by its positive minimum there. The maximum principle gives
\begin{equation}\label{d:eq:dual-local}
 \sup_{|z|\le R_0}|Q(z)|\le C_{E,R_0}L_\ell.
\end{equation}

For $z=x+iy$, write $X=|x|$, $Y=|y|\ge1$, and $R=(X^2+Y^2)^{1/2}$. The Gamma integral used in Proposition~\ref{d:prop:ray} gives
\[
 M_Y\le C Y^{Y/2+3/4}\pi^{-Y/2}e^{-Y/2}.
\]
The same symmetry, reciprocal Euler product in $\operatorname{Re}s\ge3/2$, and uniform complex Stirling calculation used in~\eqref{d:eq:offstrip-denominator} give
\[
 |\Xi(z)|\ge cR^{7/4}\left(\frac R{2\pi}\right)^{Y/2}
       \exp\left(-\frac Y2-\frac X2\arctan\frac XY\right).
\]
These estimates do not require an $\Hh$ representative for $F$; the numerator bound needed here is exactly~\eqref{d:eq:dual-quotient}. With $G(X,Y)$ defined in~\eqref{d:eq:G-function}, division gives
\[
 |Q(z)|\le C_E L_\ell(1+R)Y^{3/4}R^{-7/4}e^{G(X,Y)}
 \le C'_E L_\ell e^{G(X,Y)}.
\]
The last inequality uses $Y\le R$ and $(1+R)/R\le2$. Thus the single derivative loss removes the reciprocal-linear gain, without enlarging the exponential support.

The support calculation in Lemma~\ref{d:lem:lens-support} gives $H_L=G$ in the vertical cones and $H_L(z)=\pi|\operatorname{Re}z|/6$ in the horizontal cones. For the right sector $|\arg z|\le\pi/6$, put $f(z)=Q(z)e^{-\pi z/6}$. The preceding estimate and~\eqref{d:eq:dual-local} bound $f$ on both radial boundaries by $C_E L_\ell$. Apply the maximum principle on truncated sectors to $f(z)e^{-\epsilon z^2}$. There $\operatorname{Re}z^2\ge|z|^2/2$, so~\eqref{d:eq:dual-order} makes the outer circular boundary tend to zero. On the rays the damping has modulus at most one. Let the outer radius tend to infinity and then $\epsilon\downarrow0$. Evenness covers the left sector, and~\eqref{d:eq:dual-local} covers the compact remainder. This proves~\eqref{d:eq:dual-lens} with a constant linear in $L_\ell$.
\end{proof}

\begin{lemma}[Ordinary Sobolev source from the exterior contour]\label{d:lem:dual-source}
There is a smooth even ordinary $H^1$ function $v_0$, with all fixed ordinary exponential moments, such that
\begin{equation}\label{d:eq:dual-source}
 F_{v_0}=F=\Xi Q,\qquad
 \norm{v_0}_{H^1(\R)}\le C_E L_\ell.
\end{equation}
This conclusion initially makes no assertion that $v_0\in\Hh$.
\end{lemma}
\begin{proof}
Set $P(w)=Q(-iw)$, so~\eqref{d:eq:dual-lens} gives
$|P(w)|\le C_E L_\ell e^{h_L(w)}$.
The exponential Borel series $B_\ell(b)=\sum_{n\ge0}n!p_nb^{-n-1}$ converges outside a disk by scalar Cauchy estimates. Its directional continuation is
\[
 e^{i\theta}\int_0^\infty
       P(te^{i\theta})e^{-bte^{i\theta}}\,dt,
 \qquad\operatorname{Re}(be^{i\theta})>h_L(e^{i\theta}).
\]
Expansion sufficiently far out recovers the series. The half-planes cover $\C\setminus L$ by convex separation. At a point in an overlap, separating directions form an open convex cone without opposite rays; a compact angular arc joining two such directions has a uniform positive gap. Contour rotation is therefore justified by an outer-arc bound $O(e^{-\delta T})$, and the inner arc tends to zero. The local continuations patch to a single holomorphic function. Choosing a normal at a nearest point of $L$ gives
\begin{equation}\label{d:eq:dual-Borel}
 |B_\ell(b)|\le\frac{C_E L_\ell}{\operatorname{dist}(b,L)}
 \qquad(b\notin L).
\end{equation}
The inverse-distance bound replaces the earlier logarithmic one; a fixed exterior contour only requires its finiteness.

Choose a fixed positively oriented rectifiable Jordan contour $\Gamma\subset\C\setminus L$ with $L\subset\operatorname{int}\Gamma$, winding number one, and contour and interior contained in the theta strip. Series inversion on a large circle, followed by deformation in $\C\setminus L$, yields
\[
 P(w)=\frac1{2\pi i}\int_\Gamma B_\ell(b)e^{bw}\,db.
\]
Define
\begin{equation}\label{d:eq:dual-contour-source}
 v_0(x)=\frac1{2\pi i}\int_\Gamma B_\ell(b)H_b(x)\,db,
 \qquad H_b(x)=\frac{\kappa(x+b)+\kappa(x-b)}2.
\end{equation}
The theta tails and their fixed derivatives are uniform on this compact parameter contour. Consequently this is an ordinary $H^1$-valued Bochner integral, with norm bounded as in~\eqref{d:eq:dual-source}; it is smooth, even, and has all ordinary exponential moments. It is not an $\Hh$-valued contour integral. Ordinary Fourier interchange and contour shifting in the strip give $F_{H_b}(z)=\Xi(z)\cos(bz)$. Since $P$ is even, integrating $\cos(bz)$ gives $P(iz)=Q(z)$, proving the transform identity in~\eqref{d:eq:dual-source}.
\end{proof}

\subsection{Dual recovery and full zero-jet closure in both topologies}\label{d:sec:Gamma-jet-core}

The ordinary source must now be identified with the dual distribution, and only then admitted to the actual weighted space.

\begin{lemma}[Distribution identity and weighted dual recovery]\label{d:lem:dual-recovery}
The distribution $v$ in~\eqref{d:eq:dual-distribution} equals the ordinary function $v_0$ in~\eqref{d:eq:dual-contour-source}. Moreover
\begin{equation}\label{d:eq:dual-H-recovery}
 v_0\in\Hh,\qquad \norm{v_0}_{\Hh}\le C_E L_\ell,\qquad
 \ell(u)=\ip{v_0}{u}_{\Hh}\quad(u\in E).
\end{equation}
\end{lemma}
\begin{proof}
For this test-level identification write the nonunitary Fourier transform explicitly as
$\widehat\varphi_0(t)=\int_\R\varphi(x)e^{-itx}\,dx$.
For even Schwartz $\varphi$, its conjugate inversion formula multiplied by $W$ is the $E$-valued identity
\[
 W\overline\varphi
 =\frac1{2\pi}\int_\R\overline{\widehat\varphi_0(t)}\,e_t\,dt.
\]
The Bochner integral converges because~\eqref{d:eq:dual-evaluation} on the real axis is $O(1+|t|)$ and $\widehat\varphi_0$ is Schwartz. Its $\Hh$ component identifies it with the displayed pointwise product by ordinary inversion. Applying the conjugate-linear $\ell$ gives
\begin{equation}\label{d:eq:dual-test-Fourier}
 v(\varphi)=\frac1{2\pi}\int_\R\widehat\varphi_0(t)F(t)\,dt.
\end{equation}
The ordinary $L^1$ function $v_0$ has the same identity by Fubini and~\eqref{d:eq:dual-source}. For non-even tests both sides equal their even part. Thus $v=v_0$ as tempered distributions, with no weighted admission premise.

Take the real even compact smooth cutoffs $\chi_R$ above, for $R\ge1$, and define the finite nonnegative number
\[
 b_R=\int_\R\chi_R(x)\frac{|v_0(x)|^2}{W(x)}\,dx.
\]
The compact smooth test $\varphi=\chi_R\overline{v_0}/W$ is legitimate because $v_0$ is smooth and $W$ is smooth and positive on its support. Equation~\eqref{d:eq:dual-distribution}, now identified with $v_0$, gives exactly
\begin{equation}\label{d:eq:dual-cutoff-test}
 b_R=\ell(\chi_Rv_0),\qquad
 \norm{\chi_Rv_0}_{\Hh}^2\le b_R.
\end{equation}
The inequality uses $\chi_R^2\le\chi_R$, not equality of the two cutoff masses. The ordinary estimate already proved in Lemma~\ref{d:lem:dual-source} gives, uniformly in $R\ge1$,
\[
 \norm{\chi_Rv_0}_{H^1}\le C\norm{v_0}_{H^1}\le C_E L_\ell.
\]
Every $\chi_Rv_0$ is compact smooth, so~\eqref{d:eq:Gamma-E-upper} applies before any assertion about weighted admission of $v_0$. Thus
\begin{equation}\label{d:eq:dual-quadratic}
 b_R\le L_\ell\norm{\chi_Rv_0}_E
 \le C_E L_\ell(b_R+L_\ell^2)^{1/2}.
\end{equation}
For the form norm this uses the full Gamma energy bound in~\eqref{d:eq:Gamma-first-order}; for the graph norm it uses the actual operator admission~\eqref{d:eq:Gamma-H1}. The bounded perturbation costs only $\norm{\mathsf D}\sqrt{b_R}$ and is absorbed by the same inequality.

If $L_\ell=0$ the conclusion is immediate. Otherwise divide~\eqref{d:eq:dual-quadratic} by $L_\ell^2$, square, and solve the resulting quadratic inequality to obtain $b_R\le C'_E L_\ell^2$. Since $\chi_R\to1$ pointwise, Fatou's lemma gives $v_0\in\Hh$ and the norm bound in~\eqref{d:eq:dual-H-recovery}. The ordinary $H^1$ estimate was independent of this recovery, so the argument is noncircular.

For any even compact smooth $u$, the test $\varphi=\overline u/W$ in~\eqref{d:eq:dual-distribution} yields
$\ell(u)=\int v_0\overline u/W=\ip{v_0}{u}_{\Hh}$.
The compact smooth form core, or respectively graph core, now extends this equality to every $u\in E$, because both sides are $E$-continuous. This last extension uses the graph core in the graph case, not mere form density.
\end{proof}

\begin{theorem}[Full zero jets in the exact Gamma and ambient graph norms]\label{d:thm:Gamma-jet-core}
For every bounded self-adjoint $\mathsf D$ on $\Hh$,
\begin{align}
 \overline{\mathcal J_{\rm fin}}^{{X}}
   &={X}\cap\Kk,\label{d:eq:Gamma-form-jet-core}\\
 \overline{\mathcal J_{\rm fin}}^{\mathcal G_{\mathsf D}}
   &=\operatorname{Dom}T_\Gamma\cap\Kk.
 \label{d:eq:Gamma-graph-jet-core}
\end{align}
More precisely, for either $E={X}$ or $E=\mathcal G_{\mathsf D}$, its continuous conjugate-linear annihilator is
\begin{equation}\label{d:eq:Gamma-annihilator}
 \{\ell\in E^*: \ell(\mathcal J_{\rm fin})=0\}
 =\{u\mapsto\ip{n_0}{u}_{\Hh}:n_0\in\Nn_0\}.
\end{equation}
The graph closure is in the ambient domain norm of $T_\Gamma-\mathsf D$. It does not assert that $\Kk$ reduces this operator or that its restriction to $\Kk$ is self-adjoint.
\end{theorem}
\begin{proof}
For an annihilating $\ell$, Lemma~\ref{d:lem:dual-recovery} gives an actual $\Hh$ representative $v_0$ with $F_{v_0}=\Xi Q$. Only at this stage apply the Hilbert synthesis theorem. Since $\Xi(a)=1/2$,
\[
 Cv_0=v_0-Q(a)\kappa,\qquad
 F_{Cv_0}(z)=\Xi(z)[Q(z)-Q(a)].
\]
Thus $Cv_0\in\Ii=\Nn$ by Theorem~\ref{d:thm:full-synthesis}, and $v_0\in\Nn_0$. This proves one inclusion in~\eqref{d:eq:Gamma-annihilator}. Conversely, each $n_0\in\Nn_0$ represents a continuous functional on $E$ because $E\hookrightarrow\Hh$. It annihilates every full jet by the $\kappa$ and radical pairings in~\eqref{d:eq:wave-center}--\eqref{d:eq:wave-radical}, followed by $\Hh$ continuity. This proves the reverse inclusion.

The common kernel of the functionals on the right of~\eqref{d:eq:Gamma-annihilator} is $E\cap\Nn_0^\perp=E\cap\Kk$. The closed-linear-span form of Hahn--Banach therefore gives the two closure identities. Explicitly, an element of $E\cap\Kk$ outside the $E$-closure of $\mathcal J_{\rm fin}$ would admit a continuous separating functional; its conjugate-linear form is classified by~\eqref{d:eq:Gamma-annihilator} and must vanish on that element, a contradiction. The reverse inclusion follows from $\mathcal J_{\rm fin}\subset\Kk$ and continuous inclusion $E\hookrightarrow\Hh$.

The representer in~\eqref{d:eq:Gamma-annihilator} is unique because compact smooth vectors make $E$ dense in $\Hh$. Its $\Hh$ norm and its $E$-dual norm are equivalent on this particular annihilator by~\eqref{d:eq:dual-H-recovery} and the continuous inclusion. No equivalence of the stronger and weaker vector norms on $\Kk$ is asserted.
\end{proof}

The theorem is a topology-lifting argument, not an inference from $\Hh$ density alone or from individual jet smoothness. The distributional quotient, its weakened prefactor, the ordinary exterior-contour source, compact weighted recovery, and the appropriate compact smooth core are all needed. It supplies existence of finite-jet approximations in these two precise norms, but no uniform rate, coefficient bound, unconditional basis, finite-section conditioning, or algebraic invariance of the finite jet span.


The preceding graph closure is in the ambient graph norm for an arbitrary bounded self-adjoint perturbation. It does not assert that $\Kk$ reduces such a perturbation. The actual physical operator has additional radical identities, which provide precisely that missing reduction in the next section.


\section{Actual physical reduction and the exact finite sign}\label{sec:actual-jets}
The stronger-topology theorem applies to the actual physical correction $D_b$ from~\eqref{eq:Db-physical}. The conclusion here is more specific than ambient graph density: the radical identities make the actual complementary sector reducing. No corresponding claim is made for an arbitrary bounded self-adjoint perturbation.

\begin{theorem}[Actual reducing sector and full zero jet cores]\label{thm:actual-core}
Let $B=T_\Gamma-D_b=UAU^{-1}-I/2$ on $\Hh$, and let $\Nn$ and $\Kk$ be as in~\eqref{d:eq:N-intro}. Then
\begin{equation}\label{eq:actual-radical}
 \Nn\subset\dom B,\qquad B|_{\Nn}=0,\qquad B\kappa=-\kappa/2.
\end{equation}
The orthogonal decomposition
\begin{equation}\label{eq:actual-decomposition}
 \Hh=\C\kappa\oplus\Nn\oplus\Kk
\end{equation}
reduces $B$. In particular $B_{\Kk}=B|_{\dom B\cap\Kk}$ is self-adjoint in $\Kk$, and
\begin{equation}\label{eq:actual-cores}
 \overline{\mathcal J_{\mathrm{fin}}}^{X}=X\cap\Kk,
 \qquad
 \overline{\mathcal J_{\mathrm{fin}}}^{\,\operatorname{graph}(B)}
       =\dom B\cap\Kk.
\end{equation}
Thus the restriction of $B_{\Kk}$ to the algebraic full-zero-jet span is essentially self-adjoint and has closure $B_{\Kk}$. Its finite span need not be invariant under $B$.
\end{theorem}
\begin{proof}
Every finite radical $r$ belongs to $\dom B$ and satisfies $Br=0$, by Lemma~\ref{a:lem:radical} and the representative criterion. If $r_n\to n\in\Nn$ in $\Hh$, then $(r_n,Br_n)\to(n,0)$ in $\Hh\oplus\Hh$. Closedness of $B$ proves $n\in\dom B$ and $Bn=0$. The source identity $q_0(\kappa,z)=-\overline{m(z)}$ and $\ip\kappa z_{\Hh}=2\overline{m(z)}$ give $B\kappa=-\kappa/2$ by the same criterion.

For $u\in\dom B$, self-adjointness gives
\[
 \ip{Bu}{n}_{\Hh}=0\quad(n\in\Nn),\qquad
 \ip{Bu}{\kappa}_{\Hh}=-\tfrac12\ip u\kappa_{\Hh}.
\]
The projections onto $\Nn$ and $\C\kappa$ map $\dom B$ into itself, since both target spaces are in the domain, and commute with $B$ by these identities. Their complementary projection onto $\Kk$ therefore also preserves the domain and commutes with $B$. This proves reduction and self-adjointness of the part.

We also record form-domain preservation explicitly. For $n\in\Nn$,
\[
 \tg[n]=\ip{D_bn}{n}_{\Hh}\le\norm{D_b}\norm n_{\Hh}^2,
 \qquad \norm n_X\le\sqrt{1+\norm{D_b}}\norm n_{\Hh}.
\]
Hence $P_{\Nn}:\Hh\to X$ is bounded. The one-dimensional projection onto $\kappa\in X$ is bounded into $X$ as well, so $P_{\Kk}=I-P_{\Nn}-P_{\kappa}$ preserves $X$ continuously. This does not assert that $T_\Gamma$ alone reduces these spaces.

Theorem~\ref{d:thm:Gamma-jet-core}, with $\mathsf D=D_b$, gives both closures in~\eqref{eq:actual-cores}; its hypotheses hold by the common Gamma and prime-domain construction. Every jet is in the actual operator domain by Lemma~\ref{d:lem:Gamma-H1}: the theta derivative estimates give $W\partial_z^k\cos(zx)|_{z=w}\in\Hh\cap H^1(\R)$. The graph closure is therefore an operator core for the self-adjoint part $B_{\Kk}$. This proves essential self-adjointness of its symmetric core restriction, without assuming invariance of the finite span.
\end{proof}

Under $U^{-1}$, the finite profile family is
\begin{equation}\label{eq:actual-profile-jets}
 \frac{J_{w,k}(x)}{\kappa(x)}
 =\frac{\partial_z^k\cos(zx)|_{z=w}}{2\cosh(x/2)},
 \qquad \Xi(w)=0,\quad 0\le k<\nu(w).
\end{equation}
It is a form core and an $A$-graph core for the reducing profile sector $U^{-1}\Kk$. Unitarity and the bounded graph shift by $I/2$ preserve these closure statements. No backward-division operator has been commuted through $B$ or through the prime and Gamma rows.

\subsection{The exact Hermitian matrices}
For a finite list of actual zero jets indexed by $\alpha=(w,k)$, put $j_\alpha=J_{w,k}$ and define
\begin{align}\label{eq:actual-matrix}
 M_{\alpha\beta}
 ={}&\int_0^\infty\rho(s)\int_\R
 [j_\alpha(x+s)-j_\alpha(x)]
 \overline{[j_\beta(x+s)-j_\beta(x)]}\,dx\,ds\notag\\
 &-\cg\int_\R j_\alpha(x)\overline{j_\beta(x)}\,dx\notag\\
 &-\sum_{n\ge2}a_n\left\{\int_\R j_\alpha(x+s_n)\overline{j_\beta(x)}\,dx
       +\int_\R j_\alpha(x)\overline{j_\beta(x+s_n)}\,dx\right\}.
\end{align}
All entries exist. Gamma pairings follow by Cauchy--Schwarz from admitted finite Gamma energies; prime pairings are absolutely convergent by Lemma~\ref{a:lem:prime}. The matrix is Hermitian, and
\begin{equation}\label{eq:actual-quadratic}
 q_0\!\left[\sum_\alpha c_\alpha j_\alpha\right]
 =\sum_{\alpha,\beta}c_\alpha\overline{c_\beta}M_{\alpha\beta}.
\end{equation}
No Gamma shift cutoff, prime head, reality hypothesis on $w$, or simplicity assumption appears in this definition. Redundant columns can make the matrix singular without affecting the sign criterion.

\begin{theorem}[Exact remaining sign equivalence]\label{thm:actual-sign}
The following four assertions are equivalent:
\begin{enumerate}[label=\textup{(\roman*)}]
\item $q_0[u]\ge0$ for every $u\in X_c$.
\item $q_0[u]\ge0$ for every $u\in X\cap\Kk$.
\item Every finite exact full-zero-jet matrix~\eqref{eq:actual-matrix} is positive semidefinite.
\item $\E[f]\ge\tfrac12\Var_\mu(f)$ for every $f\in\D$.
\end{enumerate}
The theorem proves the equivalence, not any one of these assertions.
\end{theorem}
\begin{proof}
The full-domain physical bridge after subtracting the original mean gives (i)$\Leftrightarrow$(iv). For $u\in X_c$, Theorem~\ref{thm:actual-core} gives $u=n+k$ with $n=P_{\Nn}u\in\Nn$ and $k=P_{\Kk}u\in X\cap\Kk$. Since $Bn=0$, the operator-to-form identity yields $q_0(n,z)=\ip{Bn}{z}_{\Hh}=0$ for every $z\in X$. Thus $q_0[u]=q_0[k]$, proving (i)$\Leftrightarrow$(ii).

The form $q_0$ is continuous on $X\times X$: the Gamma term is controlled by Cauchy--Schwarz in its form norm, and the complete correction is bounded on $\Hh$. The form-core equality in~\eqref{eq:actual-cores}, together with~\eqref{eq:actual-quadratic}, therefore proves (ii)$\Leftrightarrow$(iii). This passage is in the exact form norm, not merely the Hilbert norm.
\end{proof}

In particular, if a centered negative-form vector exists, some finite full-zero-jet combination already has strictly negative $q_0$. The same is true if an actual centered eigenstate of $A_E$ has eigenvalue in $(0,1/2)$, by the bridge. This follows from form convergence and strict negativity; it supplies no bound on the number, height, multiplicity or conditioning of the necessary jets. The constant ground state is not in the centered sector. Establishing all the finite signs remains open.


% Body-only addition. Insert after the exact actual-zero-jet sign equivalence.
% Existing macros: \Hh, \Kk, \norm, \ip, \tg, \dom, \Nn.
% New bibliography key: gbTeschl. No source manuscript is a bibliography premise.
\section{Gamma spectral blocks on the actual zero jet sector}
\label{gb:sec:bridge}

The actual reduction and graph core give a spectral decomposition in which
the complete high response has a uniform bound. The remaining condition is
the sign of a bounded Schur operator on the entire low spectral space.
The low spectral space is not assumed finite dimensional. This construction
uses the restricted Gamma form rather than a restriction of the ambient
Gamma operator.

Only in this section write $\mathcal L=T_\Gamma-D_b$ for the ambient
physical operator, reserving $B$ for an off-diagonal block. Retain
$\mathcal J_{\rm fin}$, $\Kk$, and $X$ from
Theorem~\ref{thm:actual-core}. Thus $\Kk$ reduces $\mathcal L$,
its self-adjoint part $\mathcal L_\Kk$ has domain
$\dom\mathcal L\cap\Kk$, and $\mathcal J_{\rm fin}$ is its graph
core and a form core in $X_\Kk:=X\cap\Kk$.
Physical norms and pairings in this section are those of $\Hh$ unless
another space is indicated.

\subsection{The restricted form and projected graph cores}

\begin{lemma}[The Gamma representative on the reducing sector]
\label{gb:lem:restricted}
The restriction $\tg|_{X_\Kk}$ is a densely defined closed nonnegative
form in $\Kk$. Denote its representative by $\Gamma_\Kk$, and put
$D_\Kk=P_\Kk D_b|_\Kk$. Then
\begin{equation}\label{gb:eq:restricted}
 \Gamma_\Kk=\mathcal L_\Kk+D_\Kk,\qquad
 \dom\Gamma_\Kk=\dom\mathcal L\cap\Kk.
\end{equation}
The finite full-zero-jet span is a graph core for $\Gamma_\Kk$.
\end{lemma}
\begin{proof}
The embedding $X\hookrightarrow\Hh$ is continuous and $\Kk$ is closed
in $\Hh$. Therefore $X_\Kk$ is closed in $X$, so the restricted form
is closed. Its domain is dense in $\Kk$, since it contains the
Hilbert-total family $\mathcal J_{\rm fin}$.
The form representative exists by the closed-form representation theorem;
see, for example, \cite[Theorem 2.13]{gbTeschl}.

As $\Kk$ reduces the semibounded self-adjoint operator $\mathcal L$,
the restriction $q_0|_{X_\Kk}$ is the closed form of its part
$\mathcal L_\Kk$. For $u,v\in X_\Kk$ the exact identity is
\[
 \tg(u,v)=q_0(u,v)+\ip{D_\Kk u}{v}_{\Hh}.
\]
The bounded self-adjoint perturbation $\mathcal L_\Kk+D_\Kk$ has
the same operator domain as $\mathcal L_\Kk$, and its form is the
displayed sum on $X_\Kk$. Uniqueness of the representative proves
\eqref{gb:eq:restricted}. Equivalently, the form-domain assertion follows
by spectral calculus for the reducing projection and the operator-domain
assertion by bounded perturbation. The two graph norms are equivalent:
\[
 \norm{\Gamma_\Kk u}\le\norm{\mathcal L_\Kk u}
       +\norm{D_\Kk}\norm u,\qquad
 \norm{\mathcal L_\Kk u}\le\norm{\Gamma_\Kk u}
       +\norm{D_\Kk}\norm u.
\]
Consequently their graph cores coincide.
\end{proof}
Neither $T_\Gamma$ nor $D_b$ is assumed to reduce $\Kk$ separately.
On the common domain the identity
$\Gamma_\Kk u=P_\Kk T_\Gamma u$ follows from
\eqref{gb:eq:restricted}; this does not identify the restricted-form
representative with the naive operator part of $T_\Gamma$.

Fix $d\ge\norm{D_b}$ and $R>d$. All subsequent spectral projections
are those of this same $\Gamma_\Kk$:
\begin{gather}
 P=\mathbf1_{[0,R)}(\Gamma_\Kk),\qquad Q=I-P,\qquad
 G=\Gamma_\Kk|_{Q\Kk},\qquad g=R-d>0,\notag\\
 A=P\Gamma_\Kk P-PD_\Kk P,\qquad B=-QD_\Kk P,\notag\\
 C=G-QD_\Kk Q,\qquad
 E=-G^{-1/2}QD_\Kk QG^{-1/2}.\label{gb:eq:blocks}
\end{gather}
Here $A$ is bounded self-adjoint on $P\Kk$, $B:P\Kk\to Q\Kk$
is bounded, and $C$ is self-adjoint on $\dom G$. Its form domain is
$QX_\Kk$. Spectral calculus gives $P\Kk\subset\dom\Gamma_\Kk$
and $X_\Kk=P\Kk\oplus QX_\Kk$.

\begin{theorem}[Actual high bounds and projected jet cores]
\label{gb:thm:high}
The complete blocks satisfy
\begin{align}
 E&\succeq-\frac dR I,&
 C&\succeq(1-d/R)G\succeq gI,\label{gb:eq:high}\\
 M_R:=B^*G^{-1}B&\preceq\frac{d^2}{R}I,&
 F_R:=B^*C^{-1}B&\preceq(1-d/R)^{-1}M_R
                     \preceq\frac{d^2}{g}I.\label{gb:eq:response}
\end{align}
The space $P\mathcal J_{\rm fin}$ is Hilbert dense in $P\Kk$, and
$Q\mathcal J_{\rm fin}$ is a graph core for $G$ and $C$.
These conclusions also hold for the graph core of
$G-QD'Q$ for any bounded self-adjoint $D'$ on $\Kk$.
\end{theorem}
\begin{proof}
The spectral theorem gives $G\succeq RI$ and
$\norm{G^{-1/2}}\le R^{-1/2}$. Thus $\norm E\le d/R$.
For $h\in QX_\Kk$,
\[
 \ip{D_\Kk h}{h}\le d\norm h^2
       \le\frac dR\norm{G^{1/2}h}^2,
\]
which proves the high-form inequalities. Since $\norm B\le d$,
$M_R\preceq d^2I/R$. Inverse order gives
$C^{-1}\preceq(1-d/R)^{-1}G^{-1}$ and proves the remaining bounds.

For $h\in\dom G$, choose $j_m\in\mathcal J_{\rm fin}$ such that
$j_m\to h$ and $\Gamma_\Kk j_m\to\Gamma_\Kk h$ in $\Hh$.
The spectral projection $Q$ preserves $\dom\Gamma_\Kk$ and commutes
with its action. Hence $Qj_m\to h$ and
$GQj_m=Q\Gamma_\Kk j_m\to Gh$. This proves the graph-core assertion
for $G$. Bounded perturbations give equivalent graph norms on the same
domain, proving both remaining graph-core assertions. Applying $P$ to
Hilbert approximations by jets proves the low-space density.
\end{proof}

If $Q\Kk=\bigoplus_j\mathcal H_j$ reduces $G$, let $G_j$ denote
its parts and $B_j$ the rows of $B$. Orthogonality gives
\[
 \sum_j B_j^*G_j^{-1}B_j=M_R
\]
in the positive strong sense. Thus the high-bound premise of
Theorem~\ref{lat:thm:blocks} is available with $\kappa=d/R$ for this
spectral decomposition. Its additional relative premise would be
\begin{equation}\label{gb:eq:relative}
 M_R\preceq(1-d/R)A.
\end{equation}
The absolute estimate \eqref{gb:eq:response} does not prove this premise
or even $A\succeq0$.

\subsection{The exact Schur sign and complete residual debit}

\begin{theorem}[Exact sign on the whole low spectral space]
\label{gb:thm:schur}
For the fixed splitting \eqref{gb:eq:blocks},
\begin{equation}\label{gb:eq:sign}
 q_0\ge0\text{ on }X_\Kk
 \quad\Longleftrightarrow\quad
 S_R:=A-B^*C^{-1}B\succeq0\text{ on }P\Kk.
\end{equation}
\end{theorem}
\begin{proof}
For $l\in P\Kk$ and $h\in QX_\Kk$ the exact block form is
\[
 q_0[l+h]=\ip{Al}{l}+2\operatorname{Re}\ip{Bl}{h}
                       +\norm{C^{1/2}h}^2.
\]
The lower bound $C\succeq gI$ ensures that
$C^{-1}Bl\in\dom C\subset QX_\Kk$. Completing the square on this
full form domain gives
\begin{equation}\label{gb:eq:square}
 q_0[l+h]=\norm{C^{1/2}(h+C^{-1}Bl)}^2+\ip{S_Rl}{l}.
\end{equation}
This proves sufficiency. Choosing $h=-C^{-1}Bl$ proves necessity.
\end{proof}
By Theorem~\ref{thm:actual-sign}, \eqref{gb:eq:sign} is also equivalent
to the centered global half inequality. This is an equivalence, not an
established sign. Positivity would force $A\succeq0$ and
$B(\ker A)=0$. Absolute smallness of $B$ cannot replace that exact
kernel compatibility.

The scalar relative criterion can be strictly stronger than the exact
Schur criterion. For $A=1$, $G=I_2$,
$E=\operatorname{diag}(-1/2,1)$, and $B=(0,1)^{\mathsf T}$, the full
block matrix
\[
 \begin{pmatrix}1&0&1\\0&1/2&0\\1&0&2\end{pmatrix}
\]
is positive definite and $B^*C^{-1}B=1/2$. Its leading principal
minors are $1,1/2,1/2$. But the original scalar pair requires both
$\kappa\ge1/2$ and $1\le1-\kappa$, which are incompatible.
This is a finite structural example, not a theta counterexample.

For the prime truncation keep $P,Q,\Gamma_\Kk$ fixed. Let $D_{b,N}$
retain the constant term and every prime power through $N$, and set
$D_{\Kk,N}=P_\Kk D_{b,N}|_\Kk$.
Choose a common $d$ with
\begin{equation}\label{gb:eq:truncation}
 \norm{D_b},\ \sup_{N\ge4}\norm{D_{b,N}}\le d,
 \qquad\norm{D_b-D_{b,N}}\le r_N\longrightarrow0.
\end{equation}
One may take $d=10$ and the already proved rational $r_N$ in
\eqref{a:eq:rN}. Indeed the head through four costs at most $9$ in
the weighted norm, and the absolute operator-series majorant beyond
four is at most $r_4$; the same majorant bounds each partial tail.
Thus both the complete and every retained correction have norm at most
$9+r_4<10$. This uses Lemma~\ref{a:lem:prime}, without another prime
enumeration.

Use $D_{\Kk,N}$ in \eqref{gb:eq:blocks} to define $A_N,B_N,C_N$ and
$S_{R,N}=A_N-B_N^*C_N^{-1}B_N$. Every $C_N\succeq gI$ on the common
domain $\dom G$, and $Q\mathcal J_{\rm fin}$ is its graph core.
The retained correction need not reduce $\Kk$ before compression.

\begin{proposition}[Complete prime tail and residual lower matrix]
\label{gb:prop:residual}
Put
\begin{equation}\label{gb:eq:schur-tail}
 e_N=\left(1+\frac{2d}{g}+\frac{d^2}{g^2}\right)r_N.
 \qquad \norm{S_R-S_{R,N}}\le e_N.
\end{equation}
Let $J:\mathbb C^k\to P\Kk$ be an injective finite synthesis and set
$G_V=J^*J$. For an admitted response
$Y:\mathbb C^k\to\dom C_N$ define
\begin{gather*}
 \mathcal R=B_NJ-C_NY,\\
 H_N(Y)=(B_NJ)^*Y+Y^*(B_NJ)-Y^*C_NY.
\end{gather*}
All adjoints use the actual weighted Hilbert space and the ordinary
coefficient inner product. The full Schur compression has the lower bound
\begin{equation}\label{gb:eq:lower}
 \Lambda_{J,N,Y}:=J^*A_NJ-H_N(Y)-g^{-1}\mathcal R^*\mathcal R-e_NG_V
       \preceq J^*S_RJ.
\end{equation}
\end{proposition}
\begin{proof}
Compression gives
$\norm{A-A_N},\norm{B-B_N},\norm{C-C_N}\le r_N$,
where the last difference is its bounded extension on the common domain.
Replace the two $B$ factors at cost $2dr_N/g$. The resolvent identity
\[
 C^{-1}-C_N^{-1}=C^{-1}Q(D_\Kk-D_{\Kk,N})QC_N^{-1}
\]
then costs at most $r_N/g^2$, and multiplication by the two remaining
$B_N$ factors costs $d^2r_N/g^2$. Including $A-A_N$ proves
\eqref{gb:eq:schur-tail}; all response cross terms are retained.

Writing $Z=B_NJ$ and expanding $\mathcal R=Z-C_NY$ gives the exact
identity
\begin{equation}\label{gb:eq:residual-identity}
 Z^*C_N^{-1}Z=H_N(Y)+\mathcal R^*C_N^{-1}\mathcal R.
\end{equation}
Since $0\preceq C_N^{-1}\preceq g^{-1}I$, subtracting
$H_N(Y)+g^{-1}\mathcal R^*\mathcal R$ from $J^*A_NJ$ is a lower
bound for $J^*S_{R,N}J$. The operator error in
\eqref{gb:eq:schur-tail} becomes $e_NG_V$ under congruence by $J$,
proving \eqref{gb:eq:lower}.
\end{proof}

The residual norm includes every output row in $Q\Kk$. Compressing
it to the response trial space can discard an arbitrarily large omitted
component and is not sufficient for \eqref{gb:eq:lower}.
For a numerical evaluation, a certified Hermitian error must also be
subtracted from the displayed lower matrix. A bound
$-\varepsilon G_V\preceq\Delta\preceq\varepsilon G_V$ controls its
effect in the represented norm. If only $|\Delta_{ab}|\le e$ is known,
the crude sufficient relative cost is $ke/\lambda_{\min}(G_V)$.
Redundant columns instead require passage to the represented quotient,
with exact kernel compatibility; an unsupported pseudoinverse is not a
substitute.

\subsection{A cofinal criterion with projected jet responses}

\begin{theorem}[Cofinal low spectral certificates]
\label{gb:thm:cofinal}
Choose nested finite spaces $L_n\subset P\mathcal J_{\rm fin}$ with
Hilbert-dense union in $P\Kk$, injective syntheses $J_n$ onto $L_n$,
and any sequence $\delta_n>0$ tending to zero. Put $G_n=J_n^*J_n$.
Then $S_R\succeq0$ if and only if, for every $n$, there exist a prime
cutoff $N_n\ge4$ and a response $Y_n$ whose columns belong to
$Q\mathcal J_{\rm fin}$ such that
\begin{equation}\label{gb:eq:cofinal}
 \Lambda_{J_n,N_n,Y_n}\succeq-\delta_nG_n.
\end{equation}
This theorem concerns exact matrices and supplies no uniform effective
bound on $N_n$, response dimension, coefficients, or projection accuracy.
\end{theorem}
\begin{proof}
If \eqref{gb:eq:cofinal} holds, Proposition~\ref{gb:prop:residual}
gives $\ip{S_Rl}{l}\ge-\delta_n\norm l^2$ for $l\in L_n$.
Every fixed vector in the union belongs to all sufficiently late $L_n$.
Letting $n\to\infty$ makes its Schur form nonnegative. The operator
$S_R$ is bounded on $P\Kk$, so density proves $S_R\succeq0$.

Conversely assume $S_R\succeq0$ and fix $n$.
First choose $N\ge4$ so that $2e_N<\delta_n/2$.
For each of the finitely many columns of $J_n$, the exact response
$C_N^{-1}B_NJ_n$ belongs to $\dom C_N$. Approximate each column by
$Q\mathcal J_{\rm fin}$ in the $C_N$ graph norm, using
Theorem~\ref{gb:thm:high}. On this fixed finite coefficient space this
makes $\norm{\mathcal R G_n^{-1/2}}$ arbitrarily small.
Equations~\eqref{gb:eq:schur-tail} and
\eqref{gb:eq:residual-identity} imply
\[
 \Lambda_{J_n,N,Y}
 =J_n^*S_{R,N}J_n+\mathcal R^*C_N^{-1}\mathcal R
       -g^{-1}\mathcal R^*\mathcal R-e_NG_n
 \succeq-2e_NG_n-g^{-1}\mathcal R^*\mathcal R.
\]
Choose the responses so that
$g^{-1}\mathcal R^*\mathcal R\preceq(\delta_n/2)G_n$.
Then \eqref{gb:eq:cofinal} follows. The choices are made separately
at every finite stage, with no exchange of the dimension and accuracy
quantifiers.
\end{proof}

For an added shift $\delta\ge0$, the same positive high gap gives
\[
 0\preceq B^*C^{-1}B-B^*(C+\delta I)^{-1}B
       \preceq\frac{\delta d^2}{g^2}I.
\]
Indeed the difference is
$\delta B^*C^{-1}(C+\delta I)^{-1}B$, a positive function of $C$
between the two $B$ factors. If a complete residual satisfies
$\mathcal R^*\mathcal R\preceq r^2G_V$, its response debit is at most
$(r^2/g)G_V$. Thus $r=o(\sqrt\delta)$ suffices for an $o(\delta)$
debit with this fixed gap. Graph density provides finite approximants,
not a uniform rate in growing source dimension.

The projections $P,Q$ are Gamma spectral projections on the actual
reducing sector. They are not zero-height windows, Fourier-band
projections, coordinate deletions, raw compact sources, or images under
the fixed inverse filters of Appendix~\ref{rtc:sec:tools}.
Projected jets generally are not finite combinations of the original
jets. No finite-rank assertion for $P$ and no computable enclosure of
these projections is proved here.

The signs in \eqref{gb:eq:sign} or all stages of
\eqref{gb:eq:cofinal} remain unproved for the actual theta operator.
Sending $R$ to infinity shrinks the absolute response bound but enlarges
the low sector and supplies no lower bound for $A$. This section adds
an exact domain and approximation connection, not an actual theta sign
computation, a new global numerical coefficient, or a proof of the global
half inequality or the Riemann hypothesis.

\paragraph{Raw positive low spaces and spectral projections.}
The raw four- and five-zero spaces in
Sections~\ref{sec:four-zero-certificate} and~\ref{sec:five-zero-certificate}
are not the low Gamma spectral space $P\Kk$. A positive bound on a raw
space $V$ does not automatically pass to $PV$: the actual operator
$\mathcal L_\Kk=\Gamma_\Kk-D_\Kk$ need not commute with $P$.
Its low/high cross terms therefore prevent inferring the sign on a
projection from the sign on the original vectors. Those finite
certificates alone neither provide the low Schur sign in
\eqref{gb:eq:sign} nor identify a finite Gram with the Gram of the entire
low spectral sector. Any such use requires a proved transfer with the
appropriate form and Gram errors. Their global exponential-polynomial
representation is instead compatible with the separate fixed-low-space,
unit-window sparse setting of Section~\ref{sec:low-sparse}, where all
coupling constants and positive margins must still be paid for the
chosen low space. Section~\ref{gs:sec:transfer} supplies a specific conditional transfer
using the full graph bound of the raw five-mode source. It certifies its
five-dimensional projected low space together with the complete high
Gamma form domain. Section~\ref{sx:sec:six} extends the projected source
to six dimensions at the same split, using a complete matrix response
cost. Neither certificate covers the other low spectral directions.

\section{Uniform coercivity at fixed exponential complexity}\label{sec:complexity-unified}
The finite-jet sign equivalence concerns every finite complexity. A complementary analytic result controls the high-frequency region when that complexity is fixed. We state it for general rapidly decreasing positive weights and then specialize to the same theta form. In this section and its cluster appendix only, $\HH=L^2(\R,dx/W)$ is the full space: neither weights nor functions need be even unless this is stated. The physical theta specialization then restricts to the even sector. The proof uses compact ODE control through collisions and finite clusters at arbitrary separations, not a lower bound on frequency spacing.


\begin{definition}[Admissible weight]\label{e:def:weight}
A function $W:\R\to(0,\infty)$ is admissible if it is $C^1$ and, for every
$a>0$, there is $C_a<\infty$ such that
\begin{equation}\label{e:eq:weight}
 W(x)+|W'(x)|\le C_a e^{-a|x|}\qquad(x\in\R).
\end{equation}
Set $\HH=L^2(\R,dx/W)$, with
$\norm{u}_W^2=\int_\R |u(x)|^2/W(x)\,dx$.
\end{definition}

The assumptions imply that $W$ is bounded, has a positive minimum on every
compact interval, and tends to zero at both ends. In particular,
\begin{equation}\label{e:eq:embedding}
 \norm{u}_2^2\le \norm W_\infty\norm{u}_W^2\qquad(u\in\HH).
\end{equation}
The converse below is specific to a nonlinear class, not to all of $\HH$.

\begin{definition}[Exponential complexity]\label{e:def:complexity}
For $N\in\N$ and $\beta\ge0$, let $\EE$ consist of zero and all functions
\begin{equation}\label{e:eq:exp-poly}
 P(x)=\sum_{j=1}^r p_j(x)e^{i\lambda_jx},\qquad
 \sum_{j=1}^r(\deg p_j+1)\le N,\qquad |\Ima\lambda_j|\le\beta,
\end{equation}
where the $\lambda_j\in\C$ are distinct and the $p_j$ are nonzero
polynomials. Repeated frequencies are merged first. Let $\EE(T)$ be the
subclass for which $|\Rea\lambda_j|\ge T$ for every displayed frequency.
\end{definition}

The cost of a degree-$k$ polynomial at one frequency is $k+1$, regardless
of how many of its coefficients vanish. Distinct exponential-polynomial
terms are linearly independent, as recalled in Section~\ref{e:sec:cosine},
so the merged nonzero representation is unique. These classes need not be
closed under addition at fixed $N$.

For $s>0$ define
\begin{equation}\label{e:eq:gamma-form}
 \rho(s)=\frac{e^{-s/2}}{1-e^{-2s}},\qquad
 \tg[u]=\int_0^\infty\rho(s)\norm{u(\cdot+s)-u}_2^2\,ds,
 \qquad \XX=\{u\in\HH:\tg[u]<\infty\}.
\end{equation}
The nonnegative integral is interpreted without cancellation. If $D$ is
any bounded self-adjoint operator on $\HH$, set
\begin{equation}\label{e:eq:q}
 \qD(u,v)=\tg(u,v)-\ip{Du}{v}_W,\qquad
 \qD[u]=\qD(u,u).
\end{equation}
Here $\tg(u,v)$ is the polarized difference form. Every $WP$ in
Definition~\ref{e:def:complexity} belongs to $\XX$, as proved below.

\begin{theorem}[Uniform finite-complexity coercivity]\label{e:thm:main}
Fix an admissible weight $W$, $N\in\N$, and $\beta\ge0$.
\begin{enumerate}[label=\textup{(\roman*)},leftmargin=2.1em]
\item There is $C=C(N,\beta,W)<\infty$ such that
\begin{equation}\label{e:eq:comparison}
 \norm{WP}_W^2\le C\norm{WP}_2^2\qquad(P\in\EE).
\end{equation}
\item For every fixed $0\le R<\infty$,
\begin{equation}\label{e:eq:escape}
 \varepsilon_{N,\beta,W}(T;R):=
 \sup_{0\ne P\in\EE(T)}
 \frac{\norm{\mathbf1_{[-R,R]}\widehat{WP}}_2}{\norm{WP}_2}
 \longrightarrow0\qquad(T\to\infty).
\end{equation}
\item For every bounded self-adjoint $D$ on $\HH$,
\begin{equation}\label{e:eq:coercivity}
 \inf_{0\ne P\in\EE(T)}\frac{\qD[WP]}{\norm{WP}_2^2}
 \longrightarrow+\infty,\qquad
 \inf_{0\ne P\in\EE(T)}\frac{\qD[WP]}{\norm{WP}_W^2}
 \longrightarrow+\infty.
\end{equation}
\end{enumerate}
No lower bound on frequency separation or bound on coefficient
conditioning is required. The constants and height thresholds may depend
on $N$, $\beta$, $W$, and, in \textup{(iii)}, $\norm D$.
\end{theorem}

In particular, for every prescribed $L>0$ there is $T_L<\infty$ such that
$\qD[WP]\ge L\norm{WP}_W^2$ for every $P\in\EE(T_L)$.
This is uniform in the represented function norm, not in the Euclidean
norm of raw exponential coefficients. The theorem gives no effective
value of $T_L$ and no uniformity as $N\to\infty$.

Section~\ref{e:sec:gamma} fixes the exact Gamma form and its domain.
Sections~\ref{e:sec:ode}--\ref{e:sec:proof} prove the theorem.
Section~\ref{e:sec:theta} verifies the theta and full-prime specialization
without importing an energy-realization theorem. The remaining sections
record finite cosine-jet consequences and the limitations of the
complexity hypothesis. No historical-priority claim is made.


\subsection{The exact Gamma form}\label{e:sec:gamma}

\begin{lemma}[Multiplier and domain]\label{e:lem:gamma}
For $u\in\HH$,
\begin{equation}\label{e:eq:multiplier}
 \tg[u]=\int_\R d_\Gamma(t)|\widehat u(t)|^2\,dt,
 \qquad
 d_\Gamma(t)=2\int_0^\infty\rho(s)(1-\cos(ts))\,ds.
\end{equation}
The multiplier is nonnegative and satisfies
\begin{equation}\label{e:eq:logbound}
 d_\Gamma(t)\ge\frac12\log(1+4t^2).
\end{equation}
Moreover, for a constant $C_\Gamma<\infty$,
\begin{equation}\label{e:eq:h1bound}
 \tg[u]\le C_\Gamma\bigl(\norm u_2^2+\norm{u'}_2^2\bigr)
 \qquad(u\in H^1(\R)).
\end{equation}
Thus $WP\in\HH\cap H^1(\R)\subset\XX$ for every exponential polynomial
$P$.
\end{lemma}
\begin{proof}
By~\eqref{e:eq:embedding}, $u\in L^2$. Plancherel for each shift and Tonelli
for the nonnegative integrand give~\eqref{e:eq:multiplier}, including when
either side is infinite. Near zero, $\rho(s)=O(s^{-1})$; at infinity,
$\rho(s)=O(e^{-s/2})$. Thus the integral defining $d_\Gamma(t)$ converges
for each $t$. The estimates
\[
 \norm{u(\cdot+s)-u}_2\le s\norm{u'}_2,
 \qquad \norm{u(\cdot+s)-u}_2\le2\norm u_2
\]
prove~\eqref{e:eq:h1bound} by splitting at $s=1$.

Since $1-e^{-2s}\le2s$,
\[
 d_\Gamma(t)\ge I(t):=\int_0^\infty
 \frac{e^{-s/2}(1-\cos(ts))}{s}\,ds.
\]
Differentiation under the integral on compact $t$ intervals gives
$I'(t)=\int_0^\infty e^{-s/2}\sin(ts)\,ds=t/(t^2+1/4)$.
Integrating from $0$ yields $I(t)=\tfrac12\log(1+4t^2)$.
Finally $P$ and $P'$ grow at most like a polynomial times an exponential.
Equation~\eqref{e:eq:weight} makes
$\int W|P|^2$, $\int W^2|P|^2$, and
$\int |W'P+WP'|^2$ finite. This proves admission.
\end{proof}

For reference, the exact digamma normalization is
\begin{equation}\label{e:eq:digamma}
 d_\Gamma(t)=\Rea\psi\left(\frac14+\frac{it}{2}\right)-\psi(1/4),
 \qquad \psi=\Gamma'/\Gamma.
\end{equation}
Indeed, subtract the integral representations
\[
 \psi(z)+\gamma=\int_0^\infty
 \frac{e^{-r}-e^{-zr}}{1-e^{-r}}\,dr\qquad(\Rea z>0)
\]
at $z=1/4+it/2$ and $z=1/4$, take real parts, and set $r=2s$
\cite[Eq.~5.9.16]{DLMF}. This identity names the form; only the directly
proved bound~\eqref{e:eq:logbound} is needed for coercivity.

\begin{proposition}[Closed form]\label{e:prop:closed}
The space $\XX$, with norm
$\norm u_X^2=\norm u_W^2+\tg[u]$, is complete. The form $\tg$ is densely
defined and nonnegative on $\HH$. The form $\qD$ is closed and bounded
below by $-\norm D\norm u_W^2$.
\end{proposition}
\begin{proof}
If $u_n$ is Cauchy in $\norm{\cdot}_X$, it converges to $u$ in $\HH$,
hence in ordinary $L^2$. In Fourier variables,
$d_\Gamma^{1/2}\widehat u_n$ is Cauchy in $L^2$. A subsequence of the two
convergent $L^2$ sequences converges almost everywhere, which identifies
the second limit as $d_\Gamma^{1/2}\widehat u$. Therefore $u\in\XX$ and
$u_n\to u$ in $\norm{\cdot}_X$. Compactly supported smooth functions lie
in $\XX$ by Lemma~\ref{e:lem:gamma} and are dense in $\HH$: truncate first,
then use the positive continuous weight bounds on a compact interval and
ordinary smooth approximation. Finally
$|\ip{Du}{u}_W|\le\norm D\norm u_W^2$ and
\[
 \norm u_X^2\le \qD[u]+(\norm D+1)\norm u_W^2
 \le \tg[u]+(2\norm D+1)\norm u_W^2.
\]
The shifted form norm is equivalent to the complete $X$ norm.
\end{proof}

No ordinary Fourier projection is asserted to preserve $\HH$. Fourier
cutoffs below are used only to estimate integrals in ordinary $L^2$.


The complete proof of Theorem~\ref{e:thm:main} is in Appendix~\ref{app:clusters}. Its constants are qualitative and may depend on the total confluent complexity.


\subsection{Specialization to the actual theta form}\label{e:sec:theta}
Lemmas~\ref{a:lem:theta} and~\ref{d:lem:theta-tails}, together with Leibniz's rule for $W=\kappa/(2c)$, show that $W$ and every fixed derivative decrease faster than every exponential. Thus the theta weight is admissible. On the full space $\HH$, put $K_su(x)=W(x)u(x+s)$ and
\[
 D_\theta=\cg M_W+\sum_{n\ge2}a_n(K_{s_n}+K_{s_n}^*),\qquad
 (K_s^*u)(x)=W(x)u(x-s).
\]
Weighted Cauchy--Schwarz gives $\norm{K_s}\le\sup_x\sqrt{W(x)W(x+s)}$. The complete norm-summable majorant proved in Lemma~\ref{a:lem:prime} did not require evenness, so this is a bounded self-adjoint operator on the full $\HH$. Reflection interchanges $K_s$ and $K_s^*$, hence the sum preserves the even sector and its restriction is $D_b$. Define $\qth=\tg-\ip{D_\theta\,\cdot}{\cdot}_W$. Its even restriction is precisely $q_0$ on $X$; no physical term is discarded.
\begin{corollary}[Complete theta specialization]\label{e:cor:theta}
For each fixed $N$ and $\beta$, both Rayleigh quotients in~\eqref{e:eq:coercivity} for $\qth$ tend uniformly to $+\infty$. After restriction to even functions, this is the same conclusion for the actual form $q_0$.
\end{corollary}
\begin{proof}
Apply Theorem~\ref{e:thm:main} with $D=D_\theta$. Restricting the class cannot decrease the infimum.
\end{proof}


\subsection{Confluent cosine families and finite matrices}\label{e:sec:cosine}

Let $W$ be even and admissible. For $w\in\C$ and $k\ge0$, put
\begin{equation}\label{e:eq:jets}
 J_{w,k}(x)=W(x)\left.\partial_z^k\cos(zx)\right|_{z=w}.
\end{equation}
These are even functions; here the word jet means only a parameter
derivative and does not impose a zero equation. Suppose that
$w_1,\ldots,w_r$ are distinct, $\Rea w_j\ge T>0$,
$|\Ima w_j|\le\beta$, and $m_j\ge1$. Set $d=\sum_jm_j$ and
\begin{equation}\label{e:eq:cosine-combination}
 u=\sum_{j=1}^r\sum_{k=0}^{m_j-1}b_{j,k}J_{w_j,k}.
\end{equation}
The identity
\begin{equation}\label{e:eq:jet-exp}
 \partial_z^k\cos(zx)=\frac12\bigl((ix)^ke^{izx}+(-ix)^ke^{-izx}\bigr)
\end{equation}
shows that $u/W\in\mathcal E_{2d,\beta}(T)$.

\begin{corollary}[Uniform finite cosine matrices]\label{e:cor:cosine}
Fix $N$, $\beta$, and a bounded self-adjoint $D$ on $\HH$. For every
$L>0$ there is $T_N(L)<\infty$ such that every combination
\eqref{e:eq:cosine-combination} with $d\le N$ and
$\min_j\Rea w_j\ge T_N(L)$ satisfies
\begin{equation}\label{e:eq:cosine-coercivity}
 \qD[u]\ge L\norm u_W^2.
\end{equation}
For the nonredundant prefix list displayed above, the form matrix is
positive definite at those heights. More precisely it dominates $L$ times
the weighted Gram matrix as a Hermitian form on the coefficient space.
\end{corollary}
\begin{proof}
The inequality is Theorem~\ref{e:thm:main} at complexity $2N$. To justify
strictness of the coefficient form, recall exponential-polynomial
independence. If $\sum_jp_j(x)e^{i\lambda_jx}=0$ with distinct
$\lambda_j$, apply
$\prod_{h\ne j}(\partial_x-i\lambda_h)^{\deg p_h+1}$.
Every term except the $j$th vanishes. After removing $e^{i\lambda_jx}$,
the remaining operator on $p_j$ is a product of powers of
$\partial_x+i(\lambda_j-\lambda_h)$.
Each factor is invertible on the finite polynomial space: in the monomial
basis its matrix is triangular with a nonzero constant diagonal.
Therefore $p_j=0$ for every $j$.

The frequencies $w_1,\ldots,w_r,-w_1,\ldots,-w_r$ are all distinct
because the first group has positive real parts and the $w_j$ are distinct.
Equation~\eqref{e:eq:jet-exp} and independence imply that the prefix jets
are linearly independent. Hence a nonzero coefficient vector represents a
nonzero $u$, and~\eqref{e:eq:cosine-coercivity} gives $\qD[u]>0$.
\end{proof}

For repeated or redundant columns the same inequality holds, but the
coefficient matrix can be singular. For a sparse selection at $w_j$, the
correct complexity cost is one plus the largest derivative order selected,
not the number of selected columns. Near collisions the unnormalized Gram
matrix can be ill-conditioned; no uniform lower eigenvalue in raw
Euclidean coordinates is asserted.

This corollary applies to arbitrary frequencies meeting its hypotheses.
It neither uses nor proves that a given frequency is a zero of a special
function. In particular, no conclusion about all actual zeta-zero jets,
a centered half inequality, or the Riemann hypothesis follows from it.

\subsection{A fixed width cosine packet}\label{e:sec:packet}

Although the general theorem has no stated rate, a common bounded offset
window permits a useful explicit form of estimate with unevaluated
constants. Fix $R_0\ge0$ and $\beta\ge0$. Let $P_+$ have complexity at most $N$ and frequencies in
$|\Rea\lambda|\le R_0$, $|\Ima\lambda|\le\beta$. Put $A=WP_+$ and
\begin{equation}\label{e:eq:packet}
 u_T(x)=\frac12\bigl(e^{iTx}A(x)+e^{-iTx}A(-x)\bigr),\qquad T>0,
\end{equation}
where $W$ is even. By Lemma~\ref{e:lem:ode}, choose a constant $C\ge1$
depending only on $N,R_0,\beta,W$ such that
$\norm A_W,\norm{A'}_2\le C\norm A_2$.

\begin{proposition}[Carrier estimate]\label{e:prop:packet}
For $T\ge2C$,
\begin{equation}\label{e:eq:packet-bound}
 \qD[u_T]\ge\left\{
 \frac12\log(1+T^2)\left(1-\frac{16C^2}{T^2}\right)
 -4\norm D C^2\right\}\norm{u_T}_2^2.
\end{equation}
\end{proposition}
\begin{proof}
Integration by parts gives
\[
 \left|\int_\R e^{2iTx}A(x)\overline{A(-x)}\,dx\right|
 \le\frac{\norm{A'}_2\norm A_2}{T}
 \le\frac C T\norm A_2^2.
\]
Expanding~\eqref{e:eq:packet} therefore yields
$\norm{u_T}_2^2\ge\norm A_2^2/4$ for $T\ge2C$.
The weighted triangle inequality gives
$\norm{u_T}_W^2\le C^2\norm A_2^2\le4C^2\norm{u_T}_2^2$.
Using~\eqref{e:eq:low-band} with the two carriers $\pm T$ and band
$[-T/2,T/2]$, including their factors $1/2$, gives
\[
 \norm{\mathbf1_{[-T/2,T/2]}\widehat{u_T}}_2
 \le\frac{2\norm{A'}_2}{T}
 \le\frac{4C}{T}\norm{u_T}_2.
\]
Now apply~\eqref{e:eq:logbound} outside that band and subtract
$\norm D\norm{u_T}_W^2$.
\end{proof}

The right side tends to a positive and then arbitrarily large coefficient
as $T\to\infty$ with the window and complexity fixed. The constants have
not been evaluated, so this is not an effective numerical threshold.
If the displayed lower bound is negative at a smaller carrier, that fact
is not a negative value of the form.

\subsection{Why complexity cannot be removed}\label{e:sec:limits}

The quantifiers in Theorem~\ref{e:thm:main} have the order
\[
 \text{for each fixed }N,\beta,L\text{ there is a height }T(N,\beta,L).
\]
They do not provide a single height for all complexities. Taking the span
of all high-frequency jets allows unbounded complexity and is not covered
by the theorem. The following constructions demonstrate that two key
uniform estimates genuinely fail on the larger arbitrary-frequency class.

\begin{proposition}[Failure of complexity independent norm comparison]\label{e:prop:norm-counter}
For any admissible $W$, let $P_m(x)=x^{2m}$. Then
\begin{equation}\label{e:eq:norm-counter}
 \frac{\norm{WP_m}_W^2}{\norm{WP_m}_2^2}\longrightarrow\infty.
\end{equation}
\end{proposition}
\begin{proof}
Define the probability measure
\[
 d\nu_m(x)=\frac{W(x)^2x^{4m}\,dx}{\int_\R W(y)^2y^{4m}\,dy}.
\]
For each fixed $R\ge1$,
\[
 \nu_m([-R,R])\le
 \left(\frac R{R+1}\right)^{4m}
 \frac{\int_{-R}^R W(x)^2\,dx}{\int_{R+1}^{R+2}W(x)^2\,dx}
 \longrightarrow0.
\]
The denominator in the fixed ratio is strictly positive. Since
$W(x)\to0$ as $|x|\to\infty$, given $B>0$ choose $R$ so that
$1/W\ge B$ outside $[-R,R]$. Then
\[
 \frac{\norm{WP_m}_W^2}{\norm{WP_m}_2^2}
 =\int_\R\frac1W\,d\nu_m
 \ge B\bigl(1-\nu_m([-R,R])\bigr).
\]
Take the lower limit and then let $B\to\infty$.
\end{proof}

\begin{proposition}[Two high carrier locations with growing degree]\label{e:prop:escape-counter}
Suppose $W$ is even and admissible. There exist $T_n\to\infty$ and real
even exponential polynomials $P_n$ whose only frequency locations are
$\pm T_n$, but whose complexities are unbounded, such that
\begin{equation}\label{e:eq:taylor-limit}
 WP_n\longrightarrow W\quad\text{in }\HH\text{ and in }H^1(\R),
 \quad\text{hence in }\XX.
\end{equation}
Consequently uniform low-frequency escape fails when no complexity bound
is imposed. For every fixed bounded self-adjoint $D$,
\begin{equation}\label{e:eq:q-limit}
 \frac{\qD[WP_n]}{\norm{WP_n}_2^2}
 \longrightarrow\frac{\qD[W]}{\norm W_2^2}<\infty.
\end{equation}
\end{proposition}
\begin{proof}
For fixed $T>0$, put
\[
 S_{m,T}(x)=\sum_{k=0}^m\frac{(-iTx)^k}{k!},\qquad
 P_{m,T}(x)=\frac12\bigl(e^{iTx}S_{m,T}(x)
 +e^{-iTx}S_{m,T}(-x)\bigr).
\]
These functions are real and even on $\R$, have frequencies $\pm T$, and
have complexity $2(m+1)$. For fixed $T$, the Taylor series and its derivative
converge locally uniformly, so $P_{m,T}\to1$ and $P_{m,T}'\to0$ pointwise.
For all $m$,
\[
 |S_{m,T}(x)|\le e^{T|x|},\qquad
 |S_{m,T}'(x)|\le T e^{T|x|},
\]
and hence $|P_{m,T}|\le e^{T|x|}$ and
$|P_{m,T}'|\le2T e^{T|x|}$. Equation~\eqref{e:eq:weight} supplies integrable
dominating functions for the weighted norm, ordinary norm, and derivative
of $W(P_{m,T}-1)$. Dominated convergence proves
$WP_{m,T}\to W$ in $\HH\cap H^1$ at fixed $T$.

Choose any $T_n\to\infty$ and then choose $m_n$ so large that the sum of
the two norm distances is below $1/n$. Set $P_n=P_{m_n,T_n}$.
Equation~\eqref{e:eq:h1bound} proves~\eqref{e:eq:taylor-limit} in $\XX$.
Because $W\in L^1$, $\widehat W$ is continuous and
$\widehat W(0)=(2\pi)^{-1/2}\int W>0$.
Thus some fixed low Fourier band has positive limiting mass, contradicting
any complexity independent escape statement. Finally $\qD$ is continuous
on $\XX$: Cauchy--Schwarz holds for $\tg$, and $D$ is bounded.
This proves~\eqref{e:eq:q-limit}. The complexities must be unbounded,
otherwise Theorem~\ref{e:thm:main}\textup{(ii)} would contradict the
nonzero limiting low-band mass.
\end{proof}

These are auxiliary arbitrary-frequency examples. The first has frequency
zero; the second prescribes arbitrarily high polynomial order at two
chosen frequencies. Neither is asserted to use actual zeta zeros or their
actual multiplicities. Neither proves that $\qth$ is negative at any
vector, and~\eqref{e:eq:q-limit} does not determine the sign of $\qth[W]$.
They rule out deleting the complexity hypothesis from this general
mechanism, not a possible theorem using additional zero-specific structure.


\section{Complexity uniform sparse windows and contiguous zero bands}\label{sec:sparse-bands}
The fixed-complexity theorem does not itself permit a growing number of jets. This section proves additional estimates under explicit window geometry or one contiguous-band restriction.

\subsection{The form and the two positivity regimes}\label{sw:sec:main}
Two complementary estimates address the cross terms of a prime--Gamma form. On a sparse family, we bound the full normalized interaction rows, uniformly in the number of windows and their total prefix complexity. Within one growing contiguous band, we instead retain the aggregate positive Gram and compare the complete form to it. The first method controls many separated blocks; the second retains all internal interactions without an internal spacing condition.

All spaces and coefficients are complex, all logarithms are natural, and inner products are linear in their first argument. Write $\norm{\cdot}_2$ for ordinary Lebesgue $L^2(\R)$ and $\tau_su(x)=u(x+s)$. The weight $W=\kappa/(2c)$ and forms $t_\Gamma,p,q_0$ are the common complete forms of Section~\ref{sec:intro}. In this section $\rho_\Gamma=\rho$. The polarization of the prime form is
\begin{equation}\label{sw:eq:p}
 p(u,v)=\sum_{n\ge2}a_n
 \bigl(\langle\tau_{s_n}u,v\rangle_2+\langle u,\tau_{s_n}v\rangle_2\bigr),
\end{equation}
so both orientations and every prime power are retained. The actual theta series is equivalently
\begin{equation}\label{sw:eq:W}
 W(x)=\frac{1}{1+e^{-|x|}}\sum_{j\ge1}(4\tau_j^2-6\tau_j)e^{-\tau_j},
 \qquad\tau_j=\pi j^2e^{2|x|}.
\end{equation}

For domain statements we use the full spaces and their even parts:
\begin{equation}\label{sw:eq:spaces}
 \begin{split}
 \mathscr H&=L^2(\R,dx/W),\qquad
 \mathscr X=\{u\in\mathscr H:\tg[u]<\infty\},\\
 \mathcal H&=\mathscr H_{\rm even},\qquad
 X=\mathscr X_{\rm even},\qquad
 \norm u_{\mathscr X}^2=\norm u_{\mathscr H}^2+\tg[u].
 \end{split}
\end{equation}
Intermediate modulated summands need not be even and are treated in the full spaces. The final paired packets belong to $X$. Section~\ref{sw:sec:domain} proves the precise continuity facts needed for the closure conclusions, without a graph-core or operator-realization assertion.

An exponential polynomial of prefix complexity at most $N$ has a representation
\begin{equation}\label{sw:eq:prefix}
 P(x)=\sum_{\ell=1}^J e^{i\lambda_\ell x}p_\ell(x),\qquad
 \deg p_\ell<m_\ell,\qquad \sum_{\ell=1}^Jm_\ell\le N.
\end{equation}
Equal frequencies may be merged. A selected degree $k$ costs its entire prefix of length $k+1$, even if some lower coefficients vanish. No separation of frequencies or bound on the coefficient norm is imposed.

At height $T$, a paired unit-window packet is
\begin{equation}\label{sw:eq:packet}
 u_T(x)=\frac{e^{iTx}A(x)+e^{-iTx}A(-x)}2,
 \qquad A=WQ,
\end{equation}
where $Q$ has frequencies in $[0,1]+i[-1/2,1/2]$. A declared integer budget $N\ge1$ must bound both the prefix budget of $Q$ and the total prefix budget of the paired representation. Twice the displayed prefix budget of $Q$ always suffices. Each window in a finite sum contains one aggregate packet satisfying that window's budget.

Define, for $T\ge e^{18}$,
\begin{equation}\label{sw:eq:Phi}
 \Phi(T)=\int_{e^{18}}^T\frac{du}{(\log u)^8}.
\end{equation}
\begin{theorem}[Sparse window positivity]\label{sw:thm:sparse}
Fix $C_0>0$. There is an effectively specifiable $T_0=T_0(C_0,W)\ge e^{18}$ such that the following holds. Let $T_j\ge T_0$ be any finite or countably infinite increasing sequence with
\begin{equation}\label{sw:eq:spacing}
 \Phi(T_{j+1})-\Phi(T_j)\ge1.
\end{equation}
For each $j$, let $u_j=u_{T_j}$ be a packet~\eqref{sw:eq:packet} whose total declared paired budget obeys
\begin{equation}\label{sw:eq:log-budget}
 1\le N_j\le C_0\log T_j.
\end{equation}
Then every finite selection satisfies
\begin{equation}\label{sw:eq:main}
 q_0\!\left[\sum_j u_j\right]
       \ge\frac12\sum_j\log T_j\,\norm{u_j}_2^2.
\end{equation}
The estimate has no bound on the number of selected windows or their total complexity. A sufficient, stronger spacing condition is
\begin{equation}\label{sw:eq:sufficient-spacing}
 T_{j+1}-T_j\ge(\log T_{j+1})^8.
\end{equation}
If the partial sums of a fixed series converge in $X$, then~\eqref{sw:eq:main} holds for its sum, with the full infinite series on the right; that series is necessarily finite. Nonnegativity also holds on the $X$-closure of the set of finite admissible sums.
\end{theorem}

A variable-frequency class with bounded prefix complexity is not generally a vector space. In the closure statement, admissibility is retained before taking limits; unrestricted linear combinations inside one window could exceed its budget. If each window instead has a fixed finite frequency-prefix list meeting the total paired cap, every vector in its fixed block space is admissible. In that case Theorem~\ref{sw:thm:sparse} gives nonnegativity on the $X$-closure of the algebraic sum of those fixed spaces. The actual full zero-jet spaces considered in Section~\ref{sw:sec:zero} have such fixed lists.

The second regime concerns one whole contiguous band, without internal spacing. Let $V_{T,L}$ denote the complex span of all actual cosine jets~\eqref{sw:eq:jets} associated with nontrivial zeta zeros in the closed band $[T,T+L]$, with every multiplicity prefix retained. No critical-line, real-coefficient, or simple-zero assumption is made.

\begin{theorem}[Full logarithmic zero bands]\label{sw:thm:pnt-bands}
For every fixed $A>0$,
\begin{equation}\label{sw:eq:pnt-band-liminf}
 \liminf_{T\to\infty}\ \inf_{0\ne u\in V_{T,A\log T}}
 \frac{q_0[u]}{\log T\,\norm u_2^2}
 \ge1-\sqrt{\frac{A\log(4e)}{2\pi^2}}.
\end{equation}
The assertion is vacuous for an empty space. In particular, every fixed $0<A<2\pi^2/\log(4e)$ gives eventual strict positivity on that entire single band. There exists $T_8$ such that for every real $T\ge T_8$,
\begin{equation}\label{sw:eq:eight-band}
 q_0[u]\ge\frac1{1000}\log T\,\norm u_2^2,
 \qquad u\in V_{T,8\log T}.
\end{equation}
All internal couplings are included; the norm is that of the whole represented function.
\end{theorem}
The proof is given in Sections~\ref{sw:sec:pnt-local}--\ref{sw:sec:pnt-bands}. Its qualitative PNT input gives existence of the final threshold, not an evaluated or effective final height in this paper. It does not authorize adding separately positive bands and discarding their cross terms.

For a complementary finite-constant branch, Theorem~\ref{sw:thm:contiguous} supplies an explicit symbolic threshold rule under total paired complexity at most $\delta_*(\log T)^2$, with $\delta_*=1/39200$. Corollary~\ref{sw:cor:contiguous-jets} applies it to actual bands of width at most $\eta_*\log T$, with $\eta_*=1/7840000$, and proves a two-sided aggregate Gram asymptotic at sublogarithmic widths. Those results do not use PNT and remain valid with their original constants.

\subsection{Analytical estimates used in the proof}\label{sw:sec:inputs}
The following estimates are established in Appendices~\ref{sw:sec:theta}--\ref{sw:sec:carrier}; they are stated here to keep the cross-window argument readable. They concern the same actual weight and weight~\eqref{sw:eq:W} and common complete forms, and all norms are ordinary $L^2$.

\begin{proposition}\label{sw:prop:analytic}
There are common effectively specifiable constants $K\ge e^2$, $K_1,K_2>0$ depending only on $W$ with the following properties. For every integer $N\ge1$, set
\begin{equation}\label{sw:eq:analytic-scales}
 M_N=K(N+1)\log(N+2),\qquad R_N=\tfrac12\log M_N,
 \qquad b_N=K_2(N+1)^4\log(N+2).
\end{equation}
\begin{enumerate}
\item If $P$ has prefix complexity at most $N$ and $|\Ima\lambda_\ell|\le1/2$, with no restriction on its real carriers, then
\begin{align}
 \norm{\one_{\{|x|>r\}}WP}_2&\le e^{-e^{2r}}\norm{WP}_2
                  &&(r\ge R_N),\label{sw:eq:input-tail}\\
 |p[WP]|&\le K_1\sqrt{(N+1)\log(N+2)}\,\norm{WP}_2^2.\label{sw:eq:prime-main}
\end{align}
\item If $Q$ additionally has $|\Rea\lambda_\ell|\le1$, then
\begin{equation}\label{sw:eq:deriv-main}
 \norm{(WQ)'}_2\le b_N\norm{WQ}_2.
\end{equation}
\item For a nonzero packet~\eqref{sw:eq:packet}, with budget $N$ as specified there, if $T\ge2b_N$, then
\begin{equation}\label{sw:eq:packet-norm}
 \norm{A}_2\le2\norm{u_T}_2,
\end{equation}
and
\begin{equation}\label{sw:eq:lower}
 \frac{q_0[u_T]}{\norm{u_T}_2^2}
 \ge\frac12\log(1+T^2)\left(1-\frac{16b_N^2}{T^2}\right)
       -\cg-K_1\sqrt{(N+1)\log(N+2)}.
\end{equation}
\end{enumerate}
\end{proposition}
The constants are independent of coefficients, multiplicities, and frequency collisions. Appendix~\ref{sw:sec:weighted} specifies one common $K$ by a scalar inequality, and explicit symbolic $K_1,K_2$ in terms of it. No numerical values are needed for this general theorem. Section~\ref{sec:four-zero-certificate} later instantiates one conservative sparse-family constant recipe; it does not evaluate the qualitative-PNT threshold.

We use the unitary Fourier transform
\begin{equation}\label{sw:eq:fourier}
 \widehat u(\xi)=(2\pi)^{-1/2}\int_\R u(x)e^{-i\xi x}\,dx.
\end{equation}
The exact Gamma multiplier is
\begin{equation}\label{sw:eq:gamma-symbol}
 \tg(u,v)=\int_\R d_\Gamma(\xi)\widehat u(\xi)
                  \overline{\widehat v(\xi)}\,d\xi,
 \quad d_\Gamma(\xi)=2\int_0^\infty\rho_\Gamma(s)(1-\cos(\xi s))\,ds.
\end{equation}
Its lower bound $d_\Gamma(\xi)\ge\frac12\log(1+4\xi^2)$ is proved in Appendix~\ref{sw:sec:theta}. For cross terms we also need an upper bound on the full, unbounded multiplier.

\begin{lemma}\label{sw:lem:gamma-upper}
For every real $\xi$,
\begin{equation}\label{sw:eq:gamma-upper}
 0\le d_\Gamma(\xi)\le4+\tfrac12\log(1+4\xi^2).
\end{equation}
\end{lemma}
\begin{proof}
The inequality $(1-e^{-t})^{-1}=1+(e^t-1)^{-1}\le1+t^{-1}$ for $t>0$ gives
$\rho_\Gamma(s)\le e^{-s/2}(1+(2s)^{-1})$. Therefore
\[
 d_\Gamma(\xi)\le
 2\int_0^\infty e^{-s/2}(1-\cos(\xi s))\,ds
 +\int_0^\infty\frac{e^{-s/2}(1-\cos(\xi s))}{s}\,ds.
\]
The first term is $4-(1/4+\xi^2)^{-1}\le4$. Differentiating the second in $\xi$ and using its value at zero gives $\frac12\log(1+4\xi^2)$. This proves the bound without truncating the original Gamma integral.
\end{proof}

\subsection{The actual cross window form}\label{sw:sec:cross}
Let $A=WQ_A$ and $B=WQ_B$ have bounded-offset prefix budgets $N,M\ge1$, with frequencies in $[-1,1]+i[-1/2,1/2]$. Write $b_A=b_N$, $b_B=b_M$, and consider
\[
 f(x)=e^{i\alpha x}A(x),\qquad g(x)=e^{i\beta x}B(x),
\]
with real carriers $\alpha,\beta$ and
\begin{equation}\label{sw:eq:cross-hyp}
 d=|\alpha-\beta|>0,\qquad |\alpha|,|\beta|\le R,
 \qquad R\ge2,\qquad b_A,b_B\le R/2.
\end{equation}
These intermediate functions need not be even. The theta bounds place them in $\mathscr H\cap H^1(\R)\subset\mathscr X$, so the full-space formulas apply.

\subsubsection{The unbounded Gamma tails}
Set
\begin{equation}\label{sw:eq:GR}
 G_R=4+\log(1+8R^2).
\end{equation}
For $\eta=\xi-\alpha$, inequality~\eqref{sw:eq:gamma-upper} and $4(\alpha+\eta)^2\le8\alpha^2+8\eta^2$ give, by concavity of the logarithm and Plancherel,
\begin{align}\label{sw:eq:gamma-full}
 \tg[f]&\le\left[4+\tfrac12\log(1+8\alpha^2+8b_A^2)\right]\norm A_2^2
 \le G_R\norm A_2^2.
\end{align}
The first step uses the probability measure $|\widehat A(\eta)|^2d\eta/\norm A_2^2$ for $A\ne0$. The second uses $b_A\le R/2$ and $1+10R^2\le(1+8R^2)^2$. The zero envelope is trivial. The same argument applies to $g$.

An ordinary Fourier-tail mass estimate cannot simply be multiplied by a global bound for $d_\Gamma$, since that multiplier is unbounded. Instead, use
\[
 1+4(\alpha+\eta)^2\le(1+8\alpha^2)(1+8\eta^2).
\]
For $r>0$, the function $\log(1+8r^2)/r^2$ is decreasing: its derivative has the required sign because $\log(1+x)\ge x/(1+x)$ for $x\ge0$. With $h=d/2$, this gives, when $|\eta|\ge h$,
\[
 d_\Gamma(\alpha+\eta)
 \le\frac{\eta^2}{h^2}
       \left[4+\tfrac12\log(1+8\alpha^2)+\tfrac12\log(1+8h^2)\right].
\]
Integrating against $|\widehat A(\eta)|^2$ and using $d\le2R$ proves
\begin{equation}\label{sw:eq:gamma-tail}
 \int_{|\xi-\alpha|\ge d/2}d_\Gamma(\xi)|\widehat f(\xi)|^2\,d\xi
 \le\frac{4b_A^2}{d^2}G_R\norm A_2^2.
\end{equation}
Only the first derivative moment has been used; no second derivative or bounded-multiplier assumption is hidden.

Partition the Fourier line into $E=\{\xi:|\xi-\alpha|\ge d/2\}$ and its complement. On $E^c$, $|\xi-\beta|\ge d/2$. Weighted Cauchy--Schwarz on each part, combining one tail estimate~\eqref{sw:eq:gamma-tail} with the opposite full estimate~\eqref{sw:eq:gamma-full}, gives
\begin{equation}\label{sw:eq:gamma-cross}
 |\tg(f,g)|\le\frac{2G_R(b_A+b_B)}d\norm A_2\norm B_2.
\end{equation}

\subsubsection{Mass and both prime orientations}
The product $A\overline B$ belongs to $W^{1,1}(\R)$, has zero limits at infinity, and its derivative has $L^1$ norm at most
$(b_A+b_B)\norm A_2\norm B_2$. Integration by parts against $e^{i(\alpha-\beta)x}$ yields
\begin{equation}\label{sw:eq:mass-cross}
 |\langle f,g\rangle_2|\le\frac{b_A+b_B}{d}\norm A_2\norm B_2.
\end{equation}
The same argument holds after any real translation of either envelope, because translations preserve its ordinary norm and derivative norm.

For each shift $s$, the two polarized correlations in~\eqref{sw:eq:p} are
\begin{align*}
 \langle\tau_sf,g\rangle_2
   &=e^{i\alpha s}\int_\R e^{i(\alpha-\beta)x}A(x+s)\overline{B(x)}\,dx,\\
 \langle f,\tau_sg\rangle_2
   &=e^{-i\beta s}\int_\R e^{i(\alpha-\beta)x}A(x)\overline{B(x+s)}\,dx.
\end{align*}
The outside phases have modulus one. The elementary prime-head estimate proved in Appendix~\ref{sw:sec:prime} is
\begin{equation}\label{sw:eq:prime-head-input}
 \sum_{2\le n\le Y}\frac{\Lambda(n)}{\sqrt n}\le6\sqrt Y\qquad(Y\ge1).
\end{equation}
Thus, for an integer $X\ge\max(M_N,M_M,2)$, integration by parts for both orientations gives
\begin{equation}\label{sw:eq:prime-cross-head}
 |p_{n\le X}(f,g)|\le12\sqrt X\,\frac{b_A+b_B}{d}\norm A_2\norm B_2.
\end{equation}

For $n>X$ set $s=\log n$. Except on a null set, either $|x+s|>s/2$ or $|x|>s/2$. Splitting on these sets, Cauchy--Schwarz and the normalized tail~\eqref{sw:eq:input-tail} give
\[
 \int_\R |A(x+s)B(x)|\,dx\le2e^{-n}\norm A_2\norm B_2,
\]
because $s/2\ge\max(R_N,R_M)$. The other orientation has the same bound. Also $\Lambda(n)/\sqrt n\le\log n/\sqrt n\le1$: the latter function has maximum $2/e<1$ on the positive real line. Therefore the entire remaining prime series obeys
\begin{equation}\label{sw:eq:prime-cross-tail}
 |p_{n>X}(f,g)|\le4\sum_{n>X}e^{-n}\norm A_2\norm B_2
 =\frac{4e^{-X}}{e-1}\norm A_2\norm B_2
 \le4e^{-X}\norm A_2\norm B_2.
\end{equation}
This is a linear sum of polarized correlations, not a squared prime sum; no additional head--tail products exist.

Choose the common integer cutoff
\begin{equation}\label{sw:eq:Xcut}
 X=\left\lceil\max\{M_N,M_M,2\log(2+R),2\}\right\rceil.
\end{equation}
Combining~\eqref{sw:eq:prime-cross-head} and~\eqref{sw:eq:prime-cross-tail} gives
\begin{equation}\label{sw:eq:prime-cross}
 |p(f,g)|\le\left[12\sqrt X\,\frac{b_A+b_B}d
                     +\frac4{(2+R)^2}\right]\norm A_2\norm B_2.
\end{equation}
Every prime power and the infinite remainder have been included.

\subsubsection{All four paired carriers}
Equations~\eqref{sw:eq:gamma-cross}, \eqref{sw:eq:mass-cross}, and~\eqref{sw:eq:prime-cross} prove
\begin{equation}\label{sw:eq:unpaired-cross}
 |q_0(f,g)|\le\mathcal L(R,d;N,M)\norm A_2\norm B_2,
\end{equation}
where
\begin{equation}\label{sw:eq:Lcross}
 \mathcal L(R,d;N,M)=
 (2G_R+|\cg|+12\sqrt X)\frac{b_N+b_M}d+\frac4{(2+R)^2}.
\end{equation}

\begin{proposition}[Actual paired cross bound]\label{sw:prop:cross}
Let $u_T,v_S$ be packets~\eqref{sw:eq:packet} with declared budgets $N,M$, respectively, and let $T,S\ge2$, $T\ne S$. If $T\ge2b_N$ and $S\ge2b_M$, then, with $R=\max(T,S)$ and $D=|T-S|$,
\begin{equation}\label{sw:eq:paired-cross}
 |q_0(u_T,v_S)|\le4\mathcal L(R,D;N,M)\norm{u_T}_2\norm{v_S}_2.
\end{equation}
\end{proposition}
\begin{proof}
Evenness of $W$ gives $A(-x)=W(x)Q_A(-x)$ and the corresponding identity for $B$. Reflection preserves prefix length, the symmetric offset rectangle, derivative norm, and symmetric tail norm. The four carrier pairs in the sesquilinear expansion have differences $T-S,T+S,-T-S,-T+S$, whose absolute values lie between $D$ and $2R$. The assumptions give $b_N,b_M\le R/2$. Thus~\eqref{sw:eq:unpaired-cross} applies to all four terms with the same $X,G_R$, and replacing each denominator by $D$ only enlarges the bound. Their four coefficients have modulus $1/4$, with sum one. Finally~\eqref{sw:eq:packet-norm} gives $\norm A_2\norm B_2\le4\norm{u_T}_2\norm{v_S}_2$.
\end{proof}
The factor four in~\eqref{sw:eq:paired-cross} comes from the two envelope-to-packet norm comparisons, not from an additional count of the four coefficients. No coefficient Gram inverse or reciprocal-theta norm comparison enters.

\subsection{Explicit constants and the diagonal margin}\label{sw:sec:constants}
Fix $C_0>0$ as in Theorem~\ref{sw:thm:sparse}, and retain the common $K,K_1,K_2$ of Proposition~\ref{sw:prop:analytic}. Define
\begin{equation}\label{sw:eq:constants}
 \begin{split}
 A_0&=C_0+1,\qquad B_0=C_0+2,\qquad L_0=\log B_0+1,\\
 B_1&=K_2A_0^4L_0,\qquad C_X=KA_0L_0+5+2\log3,\\
 G_0&=6+\log9,\qquad E_0=2G_0+|\cg|+12\sqrt{C_X},\\
 C_*&=8E_0B_1+16.
 \end{split}
\end{equation}
These deliberately conservative constants depend only on $C_0$ and the fixed theta constants. They remain symbolic here; Section~\ref{sec:four-zero-certificate} evaluates one recipe for $C_0=8$.

Choose $t_0\ge18$ to satisfy three scalar conditions:
\begin{align}
 16B_1^2t_0^{10}e^{-2t_0}&\le\tfrac18,\label{sw:eq:threshold-a}\\
 K_1\sqrt{\frac{A_0\log(B_0t_0)}{t_0}}
       +\frac{|\cg|}{t_0}&\le\tfrac18,\label{sw:eq:threshold-b}\\
 E_8(t_0)&\le\tfrac14,\label{sw:eq:threshold-c}
\end{align}
where
\begin{equation}\label{sw:eq:E8}
 E_8(t)=C_*\left[
       \frac{2^{17}+4/3}{t^2}+\frac{2^{19/2}}{t^3}
       +(2+2^{11/2})t^5e^{-t}\right].
\end{equation}
Set $T_0=e^{t_0}$. All three left sides tend to zero and decrease on $t_0\ge18$: $t^{10}e^{-2t}$ decreases for $t\ge5$; $\log(B_0t)/t$ decreases when $\log(B_0t)>1$; and each term of~\eqref{sw:eq:E8} decreases for $t\ge18$. Thus these are effective whole-ray conditions, not an isolated successful check at one height.

For $t=\log T\ge t_0$ and $N\le C_0t$,
\begin{equation}\label{sw:eq:budget-reduce}
 N+1\le A_0t,\quad N+2\le B_0t,\quad
 \log(B_0t)\le L_0t,\quad b_N\le B_1t^5.
\end{equation}
Condition~\eqref{sw:eq:threshold-a} implies $16b_N^2/T^2\le1/8$, and in particular $T\ge2b_N$. Condition~\eqref{sw:eq:threshold-b}, multiplied by $t$, bounds the prime and constant losses in~\eqref{sw:eq:lower} by $t/8$. Since $\frac12\log(1+T^2)\ge t$, we obtain the uniform diagonal margin
\begin{equation}\label{sw:eq:diagonal}
 q_0[u_T]\ge\tfrac34\log T\,\norm{u_T}_2^2.
\end{equation}

For two distinct heights $T,S$, put $R=\max(T,S)$, $D=|T-S|$, and $v=\log R$. Both budgets are at most $C_0v$. The cutoff~\eqref{sw:eq:Xcut} obeys
\begin{align}\label{sw:eq:Xbound}
 X&\le KA_0v\log(B_0v)+2\log(2+e^v)+3
 \le C_Xv^2,
\end{align}
since a ceiling costs at most one, a maximum is at most the sum of its nonnegative entries, and $\log(2+e^v)\le v+\log3$. Also
\[
 b_N+b_M\le2B_1v^5,\qquad
 G_R\le4+\log9+2v\le G_0v.
\]
Proposition~\ref{sw:prop:cross} therefore gives the actual pair estimate
\begin{align}\label{sw:eq:reduced-cross}
 |q_0(u_T,v_S)|
 &\le\left[\frac{8E_0B_1v^6}{D}+\frac{16}{(2+R)^2}\right]
                    \norm{u_T}_2\norm{v_S}_2\\
 &\le\frac{C_*v^6}{D}\norm{u_T}_2\norm{v_S}_2.\nonumber
\end{align}
The last line uses $0<D<R<(2+R)^2$ and $v\ge1$. It absorbs the already bounded full prime tail, rather than dropping it. No multiplicative gap assumption was used.

For $t_j=\log T_j$, normalize against the matching block energies by setting
\begin{equation}\label{sw:eq:kij}
 k_{ij}=\frac{C_*\max(t_i,t_j)^6}
                {|T_i-T_j|\sqrt{t_it_j}},\qquad i\ne j.
\end{equation}
These coefficients are symmetric and nonnegative. The normalization is ordinary physical $L^2$, exactly as in~\eqref{sw:eq:diagonal}.

\subsection{Packing and the complete row sum}\label{sw:sec:packing}
We prove $\sup_i\sum_{j\ne i}k_{ij}\le E_8(t_0)$. The estimate includes both directions and arbitrarily many windows, not merely nearest neighbors.

\begin{lemma}[Packing in the increasing coordinate]\label{sw:lem:packing}
Suppose points $S_j\in[U,V]\subset[e^{18},\infty)$ satisfy~\eqref{sw:eq:spacing}. For every nonnegative decreasing function $f$ on this interval,
\begin{equation}\label{sw:eq:packing}
 \sum_j f(S_j)\le f(U)+\int_U^V\frac{f(s)}{(\log s)^8}\,ds.
\end{equation}
The same statement holds for $V=\infty$ when the integral is finite. Moreover,
\begin{equation}\label{sw:eq:gap-sandwich}
 \frac{V-U}{(\log V)^8}\le\Phi(V)-\Phi(U)
                       \le\frac{V-U}{(\log U)^8}\qquad(V>U\ge e^{18}).
\end{equation}
\end{lemma}
\begin{proof}
Write $y_j=\Phi(S_j)$ and $F=f\circ\Phi^{-1}$. Apart from the first point, attach the disjoint intervals $[y_j-1,y_j]$. They lie after the first coordinate and within $[\Phi(U),\Phi(V)]$. Since $F$ decreases, the integral over each such interval is at least $F(y_j)$. Bound the first value by $f(U)$ and change variables back. Finite partial sums and monotone convergence prove the infinite version. The gap sandwich follows directly from the decreasing density $(\log s)^{-8}$.
\end{proof}
In particular, condition~\eqref{sw:eq:spacing} implies
\begin{equation}\label{sw:eq:gap-lower}
 T_{j+1}-T_j\ge(\log T_j)^8,
\end{equation}
while~\eqref{sw:eq:sufficient-spacing} implies~\eqref{sw:eq:spacing}. The stronger sufficient condition is not asserted to be necessary.

Fix a selected height $T=e^t$, $t\ge t_0$, and abbreviate the coefficient~\eqref{sw:eq:kij} by $k(T,S)$. We split the remaining heights into three disjoint regions.

\subsubsection{Near heights on both sides}
If $S\in[T/2,2T]$, both logarithms lie between $t/2$ and $2t$, and hence
\begin{equation}\label{sw:eq:near-k}
 k(T,S)\le\frac{128C_*t^5}{|T-S|}.
\end{equation}
Every consecutive selected pair inside this interval has height gap at least
$\delta=(t/2)^8$ by~\eqref{sw:eq:gap-lower}. Starting at the selected anchor $T$, the $k$th point on either side has distance at least $k\delta$. All such distances are at most $T$. If that side has $m\ge1$ points, then $m\le T/\delta$ and
\[
 \sum_{k=1}^m\frac1{k\delta}
 \le\frac{1+\log m}{\delta}
 \le\frac{1+\log(1+T/\delta)}{\delta}
 \le\frac{2t}{\delta},
\]
since $\delta\ge1$ and $t\ge18$. An empty side contributes zero. Adding both sides in~\eqref{sw:eq:near-k} gives
\begin{equation}\label{sw:eq:near-row}
 \sum_{\substack{S\in[T/2,2T]\\S\ne T}}k(T,S)
 \le128C_*t^5\frac{4t}{\delta}
 =\frac{2^{17}C_*}{t^2}.
\end{equation}
The number of near heights is not assumed bounded; its harmonic cost has been retained.

\subsubsection{Every farther lower height}
If $T_0\le S<T/2$, then
\begin{equation}\label{sw:eq:lower-k}
 k(T,S)\le\frac{2C_*t^{11/2}}{T\sqrt{\log S}}.
\end{equation}
If $T/2<T_0$, the sum is empty. Otherwise apply Lemma~\ref{sw:lem:packing} to $f(S)=(\log S)^{-1/2}$ on $[T_0,T/2]$. Since $\log S\ge t_0\ge18>17$,
\[
 \frac{d}{dS}\left(\frac{S}{(\log S)^{17/2}}\right)
 =\frac{1}{(\log S)^{17/2}}\left(1-\frac{17}{2\log S}\right)
 \ge\frac{1}{2(\log S)^{17/2}}.
\]
Dropping the nonnegative lower-end subtraction in the resulting integral bound gives
\begin{align*}
 \sum_{S<T/2}\frac1{\sqrt{\log S}}
 &\le\frac1{\sqrt{t_0}}+\int_{T_0}^{T/2}\frac{dS}{(\log S)^{17/2}}\\
 &\le\frac1{\sqrt{t_0}}+\frac{T}{(t-\log2)^{17/2}}
 \le\frac1{\sqrt{t_0}}+\frac{2^{17/2}T}{t^{17/2}}.
\end{align*}
Together with~\eqref{sw:eq:lower-k}, this proves
\begin{align}\label{sw:eq:lower-row}
 \sum_{S<T/2}k(T,S)
 &\le\frac{2C_*t^{11/2}e^{-t}}{\sqrt{t_0}}+\frac{2^{19/2}C_*}{t^3}\\
 &\le2C_*t_0^5e^{-t_0}+\frac{2^{19/2}C_*}{t_0^3}.\nonumber
\end{align}
The last step uses monotonicity of $t^{11/2}e^{-t}$ on $t\ge18$.

\subsubsection{The entire farther upper tail}
If $S>2T$ and $v=\log S$, then
\begin{equation}\label{sw:eq:upper-k}
 k(T,S)\le\frac{2C_*}{\sqrt t}\frac{(\log S)^{11/2}}S.
\end{equation}
The function $f(S)=(\log S)^{11/2}/S$ decreases when $\log S>11/2$. Its density-weighted packing integral converges, so Lemma~\ref{sw:lem:packing} gives
\begin{align*}
 \sum_{S>2T}f(S)
 &\le\frac{(t+\log2)^{11/2}}{2T}
       +\int_{2T}^{\infty}\frac{dS}{S(\log S)^{5/2}}\\
 &=\frac{(t+\log2)^{11/2}}{2T}
       +\frac23(t+\log2)^{-3/2}.
\end{align*}
Since $t+\log2\le2t$, equation~\eqref{sw:eq:upper-k} yields
\begin{align}\label{sw:eq:upper-row}
 \sum_{S>2T}k(T,S)
 &\le2^{11/2}C_*t^5e^{-t}+\frac{4C_*}{3t^2}\\
 &\le2^{11/2}C_*t_0^5e^{-t_0}+\frac{4C_*}{3t_0^2}.\nonumber
\end{align}
This is an estimate of the infinite upper-frequency tail, not an inference from decay of each individual pairing.

Combining~\eqref{sw:eq:near-row}, \eqref{sw:eq:lower-row}, and~\eqref{sw:eq:upper-row}, and using $t\ge t_0$, gives the promised complete row bound
\begin{equation}\label{sw:eq:row}
 \sup_i\sum_{j\ne i}k_{ij}\le E_8(t_0)\le\tfrac14.
\end{equation}
Missing windows, finite subsequences, or additional irregular gaps preserve the same estimate.

\subsection{Proof of the theorem and form closure}\label{sw:sec:domain}
For a finite selection set $z_j=\sqrt{\log T_j}\norm{u_j}_2$. The exact sesquilinear expansion, the diagonal estimate~\eqref{sw:eq:diagonal}, and~\eqref{sw:eq:reduced-cross} give
\begin{align*}
 q_0\!\left[\sum_j u_j\right]
 &\ge\frac34\sum_j z_j^2-2\sum_{i<j}k_{ij}z_iz_j\\
 &\ge\frac34\sum_j z_j^2-\sum_i z_i^2\sum_{j\ne i}k_{ij}
 \ge\frac12\sum_j z_j^2.
\end{align*}
The middle step is $2z_iz_j\le z_i^2+z_j^2$, and the last is~\eqref{sw:eq:row}. This proves~\eqref{sw:eq:main}, uniformly in the number of blocks, without block orthogonality or a coefficient-matrix condition. It remains to justify the exact topology used in the infinite extension.

\begin{lemma}[Continuity in the stated form norm]\label{sw:lem:domain}
The space $\mathscr X$ in~\eqref{sw:eq:spaces} is complete. The full prime form is bounded on $\mathscr H\times\mathscr H$, and $q_0$ is continuous on $\mathscr X\times\mathscr X$. The same holds after restriction to the even spaces $\mathcal H,X$.
\end{lemma}
\begin{proof}
Appendix~\ref{sw:sec:theta} gives
$W(x)\le C_4e^{4|x|-\pi e^{2|x|}}$. This envelope decreases for $|x|\ge0$, since $4-2\pi e^{2|x|}<0$, so $\norm W_\infty\le C_4e^{-\pi}$. For $s\ge0$, at least one of $|x+s|,|x|$ is at least $s/2$. Consequently
\begin{equation}\label{sw:eq:weighted-shift}
 \sup_x\sqrt{W(x+s)W(x)}
 \le C_4e^{s-(\pi/2)e^s-\pi/2}
 \le C_4e^{s-(\pi/2)e^s}.
\end{equation}
Write $u=\sqrt W F$ and $v=\sqrt W G$. Cauchy--Schwarz and translation invariance in ordinary $L^2$ show
\[
 |\langle\tau_su,v\rangle_2|
 \le\sup_x\sqrt{W(x+s)W(x)}\,\norm u_{\mathscr H}\norm v_{\mathscr H}.
\]
The other orientation has the same bound. Thus the actual infinite prime form is absolutely bounded by
\begin{equation}\label{sw:eq:H-prime}
 |p(u,v)|\le D_P\norm u_{\mathscr H}\norm v_{\mathscr H},
 \quad D_P=2C_4\sum_{n\ge2}\sqrt n\log n\,e^{-\pi n/2}<\infty.
\end{equation}
This majorant includes all integers, hence all prime powers. Its use here establishes continuity only; the much sharper ordinary-$L^2$ estimates remain those used in the positivity proof. The mass term is also $\mathscr H$-bounded, because
$\norm u_2\le\sqrt{\norm W_\infty}\norm u_{\mathscr H}$.

For completeness, let $u_n$ be Cauchy in $\mathscr X$ and let $u_n\to u$ in $\mathscr H$, hence in ordinary $L^2$. For each fixed $s$, translation is continuous in $L^2$, so Fatou in the positive $s$ integral gives
\[
 \tg[u_n-u]\le\liminf_{m\to\infty}\tg[u_n-u_m].
\]
It follows that $u\in\mathscr X$ and $u_n\to u$ in $\mathscr X$. Finally
$|\tg(u,v)|\le\tg[u]^{1/2}\tg[v]^{1/2}$ by Cauchy--Schwarz in the difference integral. Together with~\eqref{sw:eq:H-prime} and the mass bound, this proves continuity of $q_0$. The even subspaces are closed, giving the asserted restrictions.
\end{proof}

If the partial sums of a fixed block series converge in $X$, continuity makes the left side of~\eqref{sw:eq:main} converge. Its right side increases to the full nonnegative series, which the finite inequalities force to be finite. Passing to the limit proves the infinite version. Continuity also gives nonnegativity at every $X$-limit of finite admissible sums. This does not admit arbitrary abstract block-$\ell^2$ series into $X$, nor extend the conclusion to bare $\mathcal H$ or ordinary-$L^2$ closure. The distinction between fixed block spaces and variable-frequency bounded-complexity classes stated after Theorem~\ref{sw:thm:sparse} remains essential.

\begin{corollary}[Fixed multiplicative gaps]\label{sw:cor:geometric}
For every $C_0>0$ and $\theta>1$, there is an effectively specifiable height after which~\eqref{sw:eq:main} holds for any sequence satisfying $T_{j+1}\ge\theta T_j$, with the same budget and domain qualifications.
\end{corollary}
\begin{proof}
For $T\ge e^{18}$,
\[
 \Phi(\theta T)-\Phi(T)
 \ge\frac{(\theta-1)T}{(\log(\theta T))^8}\longrightarrow\infty.
\]
The displayed ratio increases once $\log(\theta T)>8$. Choose a whole-ray threshold where it is at least one, and enlarge the threshold of Theorem~\ref{sw:thm:sparse}. Any larger next height only increases the $\Phi$ gap.
\end{proof}

\subsection{Moving anchor localization and the sharper prime bound}\label{sw:sec:moving}
The sparse argument above remains valid with its original conservative constants. We now sharpen physical mass localization, without changing its derivative estimates. To keep the two scales distinct, retain $M_N,R_N$ from~\eqref{sw:eq:analytic-scales} for the derivative proof and introduce
\begin{equation}\label{sw:eq:linear-scale}
 \widetilde M_N=8(N+1),\qquad
 \widetilde R_N=\tfrac12\log\widetilde M_N.
\end{equation}
Only mass localization and the prime estimate use this new linear cutoff.

\begin{theorem}[Linear cutoff and complete prime bound]\label{sw:thm:linear-prime}
Let $P$ have prefix complexity at most the integer $N\ge1$, with $|\Ima\lambda_\ell|\le1/2$ and arbitrary real carriers. Put $u=WP$. For $r\ge\widetilde R_N$,
\begin{equation}\label{sw:eq:linear-tail}
 \norm{\one_{\{|x|>r\}}u}_2\le e^{-e^{2r}}\norm u_2.
\end{equation}
The complete prime form is absolutely convergent and satisfies
\begin{equation}\label{sw:eq:prime35}
 |p[u]|\le35\sqrt{N+1}\,\norm u_2^2.
\end{equation}
The constants are uniform over coefficients, multiplicities, collisions, and nonminimal admissible prefix budgets. Neither $P$ nor $u$ need be even.
\end{theorem}
\begin{proof}
The theta constants of Appendix~\ref{sw:sec:theta} satisfy the useful exact bounds
\begin{equation}\label{sw:eq:theta-numeric-majorants}
 c_W>9,\qquad C_4<42.
\end{equation}
Indeed, $c_W=2\pi^2-3\pi>9$ follows from $\pi>3$. In the sum defining $C_4$, the $j=2$ term is less than $16/2^9=1/32$. For $j\ge2$, the ratio of consecutive terms is at most
\[
 (3/2)^4e^{-15}<\frac{81}{16\cdot2^{15}}<\frac12.
\]
The full sum over $j\ge2$ is therefore less than $1/16$, and $\pi<22/7$ gives
$C_4<4(22/7)^2(17/16)=2057/49<42$. These are full-series majorants, not numerical evaluations of $W$.

For a fixed $P$, polynomial-exponential growth and the theta envelopes imply $u\in L^2$ before any normalization. If $u=0$ the assertions are immediate. Otherwise put $U=\norm u_2>0$.
Fix $r>0$ and apply Lemma~\ref{sw:lem:growth} to $F(t)=P((r/2)t)$ at the fixed scaled point $t=2$. The exponent real parts are bounded in modulus by $\beta=r/4$. The lemma allows every $\beta\ge0$; no lower bound on $r/2$ is required. Its weaker exponential factor $e^{\beta(t+1)}$ gives
\begin{equation}\label{sw:eq:moving-P}
 |P(r)|\le2\sqrt N\,e^{3r/4}(8e)^N(r/2)^{-1/2}
                        \norm P_{L^2[0,r/2]}.
\end{equation}
Scaling preserves each prefix degree. The confluent passage in that lemma is local uniform convergence on compact intervals, so it is valid on this moving finite anchor as well.

The lower envelope $e^{4x-\pi e^{2x}}$ decreases for $x\ge0$. Thus
\[
 \inf_{0\le x\le r/2}W(x)\ge c_We^{2r-\pi e^r},\qquad
 \norm P_{L^2[0,r/2]}\le c_W^{-1}e^{-2r+\pi e^r}U.
\]
Combining with~\eqref{sw:eq:moving-P} and the upper theta envelope yields
\begin{equation}\label{sw:eq:moving-majorant}
 \frac{|u(\pm r)|^2}{U^2}\le
 E_N(r):=\frac{8NC_4^2}{rc_W^2}(8e)^{2N}
        \exp\left\{\frac{11}2r+2\pi e^r-2\pi e^{2r}\right\}.
\end{equation}
For the negative sign apply the same argument to $P(-(r/2)t)$, whose anchor is $[-r/2,0]$. Evenness of $W$, not of $P$, supplies the same lower bound.

We verify absorption on the entire ray. Let $M=e^{2r}\ge8(N+1)$, so $M\ge16$ and $r>1$. We use $2<e<3$, $3<\pi<22/7$, and the exact inequality $\log2<7/10$, which follows from
\[
 e^{7/10}>1+\frac7{10}+\frac{(7/10)^2}{2}+\frac{(7/10)^3}{6}
 =\frac{12013}{6000}>2.
\]
By~\eqref{sw:eq:theta-numeric-majorants},
\[
 \frac{8NC_4^2}{rc_W^2}\le\frac{196M}{9r}<22M.
\]
Since $\log22<4$, and $\log M/M$ decreases on $M\ge16$,
\[
 \log\frac{8NC_4^2}{rc_W^2}\le4+\log M\le\frac{17}{40}M.
\]
The other positive terms obey
\[
 2N\log(8e)\le\frac{31}{40}M,\qquad
 \frac{11}2r\le\frac{77}{160}M,\qquad
 2\pi e^r\le\frac{11}{7}M.
\]
Here $N\le M/8$, $\log M/M\le\log16/16<7/40$, and $\sqrt M\ge4$ suffice. Consequently
\begin{equation}\label{sw:eq:sharp-absorb}
 \log E_N(r)\le
 \left(\frac{17}{40}+\frac{31}{40}+\frac{77}{160}+\frac{11}{7}-6\right)M
 =-\frac{3077}{1120}M\le-2M.
\end{equation}
Every inequality holds for all $M\ge8(N+1)$. Integrating both pointwise tails gives the stronger squared estimate
\begin{equation}\label{sw:eq:sharp-integrated}
 \frac{\norm{\one_{\{|x|>r\}}u}_2^2}{U^2}
 \le2\int_r^\infty e^{-2e^{2x}}\,dx
 =\int_{e^{2r}}^\infty\frac{e^{-2v}}v\,dv
 \le\frac{e^{-2e^{2r}}}{2e^{2r}}\le e^{-2e^{2r}},
\end{equation}
which proves~\eqref{sw:eq:linear-tail}.

To estimate the complete prime form, put $Y=8(N+1)$, an integer. The correlation split~\eqref{sw:eq:correlation} gives
$|C_u(s)|\le2U\norm{\one_{\{|x|>s/2\}}u}_2$, while always $|C_u(s)|\le U^2$. The entire head is at most $12\sqrt Y\,U^2$ by~\eqref{sw:eq:prime-head-input}, including the two orientations. For every integer $n>Y$, equation~\eqref{sw:eq:linear-tail} at $s/2=(\log n)/2$ gives $|C_u(\log n)|\le2e^{-n}U^2$. Since $a_n\le1$, the full remaining series is at most
\[
 4\sum_{n>Y}e^{-n}U^2=\frac{4e^{-Y}}{e-1}U^2
 <4\,2^{-16}U^2<U^2.
\]
This proves absolute convergence. Finally $12\sqrt8<34$, because $1152<1156=34^2$, and $\sqrt{N+1}\ge1$ give~\eqref{sw:eq:prime35}.
\end{proof}

The bound~\eqref{sw:eq:prime35} improves the complexity order of the prime loss from $\sqrt{N\log N}$ to $\sqrt N$. It does not claim an optimal exponent or a better numerical constant for every small $N$. In particular, the new mass cutoff does not replace $M_N$ in the derivative-tail estimates of Appendix~\ref{sw:sec:weighted}.

\subsubsection{A budget dependent reciprocal theta comparison}
The same moving anchor also supplies a useful, explicitly nonuniform norm comparison.
\begin{proposition}\label{sw:prop:weighted-comparison}
Under the hypotheses of Theorem~\ref{sw:thm:linear-prime},
\begin{equation}\label{sw:eq:DN}
 \norm{WP}_{\mathscr H}^2\le D_N\norm{WP}_2^2,
 \qquad D_N=1+\frac{e^{8\pi(N+1)}}{576(N+1)^2}.
\end{equation}
For even functions the weighted norm is that of $\mathcal H$. The constant depends on the budget and is not a uniform weighted gap.
\end{proposition}
\begin{proof}
Weighted membership follows directly from the theta envelope and polynomial-exponential growth. Suppose $U=\norm{WP}_2>0$. Before multiplying~\eqref{sw:eq:moving-P} by $W^2$, the moving-anchor argument gives
\[
 \frac{|P(\pm r)|^2}{U^2}\le\frac{8N}{rc_W^2}(8e)^{2N}
               e^{-5r/2+2\pi e^r}.
\]
Multiplication by only one factor $W$ bounds the weighted density by
\begin{equation}\label{sw:eq:weighted-density}
 \frac{W(\pm r)|P(\pm r)|^2}{U^2}
 \le H_N(r):=\frac{8NC_4}{rc_W^2}(8e)^{2N}
                   e^{3r/2+2\pi e^r-\pi e^{2r}}.
\end{equation}
For $M=e^{2r}\ge8(N+1)\ge16$, the prefactor is at most
$14M/(27r)<M$. The same elementary bounds as above give
\begin{align*}
 \log H_N(r)
 &\le\left(\frac7{40}+\frac{31}{40}+\frac{21}{160}+\frac{11}{7}-3\right)M\\
 &=-\frac{389}{1120}M\le-\frac13M.
\end{align*}
With $M_0=8(N+1)$ and $R=\frac12\log M_0$, integrating both weighted tails gives
\[
 \int_{|x|>R}\frac{|u(x)|^2}{W(x)}\,dx
 \le U^2\int_{M_0}^\infty\frac{e^{-v/3}}v\,dv
 \le\frac3{M_0}e^{-M_0/3}U^2<U^2.
\]
On the central interval the decreasing lower theta envelope yields
$W(x)\ge c_WM_0^2e^{-\pi M_0}>9M_0^2e^{-\pi M_0}$.
Thus its weighted mass is at most $e^{\pi M_0}U^2/(9M_0^2)$. Adding the two regions proves~\eqref{sw:eq:DN}, since $9\cdot8^2=576$.
\end{proof}
For $N\le C_0\log T$, the loss in~\eqref{sw:eq:DN} is at most
$1+(e^{8\pi}/576)T^{8\pi C_0}$. A fixed-budget physical lower margin can therefore be divided by $D_N$ to obtain a budget-dependent weighted margin. This observation supplies no additional coupling theorem and no positive gap uniform in increasing complexity or height.

\subsection{Growing offsets with the same derivative constant}\label{sw:sec:wide-derivative}
The derivative estimate extends beyond fixed-width offsets. Its proof still uses the original derivative scale $M_N=K(N+1)\log(N+2)$, not the linear mass cutoff of Section~\ref{sw:sec:moving}.

\begin{proposition}\label{sw:prop:wide-derivative}
Let $Q$ have prefix complexity at most $N\ge1$, with
$|\Ima\lambda_\ell|\le1/2$ and $|\Rea\lambda_\ell|\le N$. Then, with exactly the same $K_2$ and $b_N$ as in Proposition~\ref{sw:prop:analytic},
\begin{equation}\label{sw:eq:wide-derivative}
 \norm{(WQ)'}_2\le b_N\norm{WQ}_2.
\end{equation}
\end{proposition}
\begin{proof}
Write $z_\ell=i\lambda_\ell$, so $|\Rea z_\ell|\le1/2$ and $|\Ima z_\ell|\le N$. We first prove the same local inequality as~\eqref{sw:eq:local-derivative}, with $D_0=10\,18^{3/2}$, on this enlarged root rectangle.
For an ordinary sum with at most $N$ roots, put $a=16(N+1)$ and $g(t)=e^{-at}Q(t)$. The growth lemma depends only on the real parts of $z_\ell$, so its bound with $\beta=2$ remains valid. The squared damped tail relative to $\norm Q_{L^2[0,1]}^2$ is at most
\[
 4Ne^{4t+4}(4et)^{2N}e^{-32(N+1)t},\qquad t\ge1.
\]
Its logarithm is at most $-16Nt-32t$, exactly as in Appendix~\ref{sw:sec:derivative}. Hence $\norm g_{L^2[0,\infty)}\le\sqrt2\norm Q_{L^2[0,1]}$.

Use the exact complex half-line multipliers~\eqref{sw:eq:Bmu} and~\eqref{sw:eq:endpoint}, obtained from Erd\'elyi's weighted Theorem 2.9.1 and Lemma 3.4~\cite{Erdelyi}. With $\mu_\ell=z_\ell-a$,
\[
 \max_\ell|\mu_\ell+1/2|\le a+N+1,\qquad
 2\sqrt{\sum_{j<k}\Rea\mu_j\Rea\mu_k}\le2N(a+1/2).
\]
It follows that
\begin{equation}\label{sw:eq:wide-B}
 B(\mu)\le(2N+1)a+2N+\tfrac32\le4N(a+2),\qquad
 \sqrt{-2\sum_\ell\Rea\mu_\ell}\le\sqrt{2N(a+2)}.
\end{equation}
The original damping has already paid for imaginary roots of size $N$. Since $Q'(0)=g'(0)+ag(0)$, these inequalities yield
\[
 |Q'(0)|\le\sqrt{2N(a+2)}[4N(a+2)+a]\sqrt2\norm Q_{L^2[0,1]}
 \le D_0(N+1)^3\norm Q_{L^2[0,1]}.
\]

For confluence use $p_\ell(D_\varepsilon(t))$ with $D_\varepsilon$ as in~\eqref{sw:eq:D}. Real root parts are unchanged, and for $t\ge0$,
$|D_\varepsilon(t)|\le t$, $|D_\varepsilon'(t)|=1$. The damped functions and their first derivatives have, for fixed $Q$, an $\varepsilon$-independent polynomial-times-$e^{-(a-1/2)t}$ majorant. Its square is integrable. A boundary root's imaginary part may slightly exceed $N$ during approximation; apply the exact continuous source multipliers first, pass to the repeated-root limit, and only then use~\eqref{sw:eq:wide-B}. Dominated convergence therefore gives the same local constant. Scaling $F(t)=Q(x+\delta t)$ for $0<\delta\le1$ preserves the enlarged rectangle and proves
\begin{equation}\label{sw:eq:wide-local}
 |Q'(x)|\le D_0(N+1)^3\delta^{-3/2}\norm Q_{L^2[x,x+\delta]}.
\end{equation}

For the weighted passage, use the original $\delta=M_N^{-1}$ and $R_N=\frac12\log M_N$. The same theta oscillation bound $C_{\rm osc}=e^{L_We^2}$ and the same Fubini overlap factor give
\[
 \norm{\one_{|x|\le R_N}WQ'}_2
 \le D_0C_{\rm osc}(N+1)^3M_N\norm{WQ}_2,
 \quad
 \norm{\one_{|x|\le R_N}W'Q}_2\le L_WM_N\norm{WQ}_2.
\]
Also~\eqref{sw:eq:wide-local} with $\delta=1$ supplies the noncircular anchor
$\norm{Q'}_{L^2[-1,0]}\le D_0m_0^{-1}(N+1)^3\norm{WQ}_2$.
The growth bounds for $Q,Q'$ depend only on their root real parts. Thus the reflected and translated derivative anchors, and both exterior majorants in Appendix~\ref{sw:sec:weighted}, are unchanged. The same common constant in~\eqref{sw:eq:common-B} and the same $K$ absorb them, giving each exterior norm at most $\norm{WQ}_2$. Combining all four pieces gives
\[
 \norm{(WQ)'}_2\le(D_0C_{\rm osc}+L_W+2)(N+1)^3M_N\norm{WQ}_2
 =K_2(N+1)^4\log(N+2)\norm{WQ}_2.
\]
No derivative-tail radius or derivative constant has been improved or replaced.
\end{proof}

\subsection{The aggregate Gram on one contiguous band}\label{sw:sec:aggregate}
Let $Q$ have frequencies in $[0,L]+i[-1/2,1/2]$, where $0\le L\le N$, and let $N$ bound the full paired prefix budget as well as the positive envelope budget. Form the single aggregate packet
\begin{equation}\label{sw:eq:aggregate-packet}
 A=WQ,\qquad
 u(x)=\frac{e^{iTx}A(x)+e^{-iTx}A(-x)}2.
\end{equation}
Its norm is the norm of the entire represented sum, not a sum of separate unit-window norms. Every internal pairing is retained in this aggregate.

\begin{theorem}[Whole-band comparison]\label{sw:thm:aggregate}
Suppose $T\ge\max(1,2b_N)$, and put
\begin{equation}\label{sw:eq:aggregate-constants}
 \begin{split}
 t&=\log T,\qquad h_T=\tfrac12\log(1+T^2),\qquad
 \epsilon_{T,N}=16b_N^2/T^2,\\
 C_\Gamma&=4+\tfrac12\log37,\qquad
 E^*_{T,N}=|\cg|+35\sqrt{N+1}+\max(C_\Gamma,\epsilon_{T,N}h_T).
 \end{split}
\end{equation}
Then
\begin{align}
 q_0[u]&\ge\left[h_T(1-\epsilon_{T,N})-\cg-35\sqrt{N+1}\right]\norm u_2^2,\label{sw:eq:aggregate-lower}\\
 |q_0[u]-t\norm u_2^2|&\le E^*_{T,N}\norm u_2^2.\label{sw:eq:aggregate-error}
\end{align}
For any fixed common frequency-prefix list fitting the budget and any finite spanning list $v_1,\ldots,v_d$ of its packet space, including redundant lists, define
\begin{equation}\label{sw:eq:aggregate-matrices}
 G_{ij}=\langle v_j,v_i\rangle_2,\qquad
 (Q_0)_{ij}=q_0(v_j,v_i).
\end{equation}
The actual matrices satisfy the Loewner inequalities
\begin{equation}\label{sw:eq:loewner}
 -E^*_{T,N}G\preceq Q_0-tG\preceq E^*_{T,N}G.
\end{equation}
\end{theorem}
\begin{proof}
Proposition~\ref{sw:prop:wide-derivative} supplies the same derivative bound used in Appendix~\ref{sw:sec:carrier}. Integration by parts in the reflected carrier cross term and Plancherel consequently give
\[
 \norm u_2^2\ge\tfrac12(1-b_N/T)\norm A_2^2,
 \qquad \norm A_2\le2\norm u_2,
 \qquad
 \norm{\one_{[-T/2,T/2]}\widehat u}_2\le\frac{4b_N}{T}\norm u_2.
\]
The full nonnegative Gamma multiplier therefore obeys
\begin{equation}\label{sw:eq:aggregate-gamma}
 \tg[u]\ge h_T(1-\epsilon_{T,N})\norm u_2^2.
\end{equation}
Subtract the exact constant term and apply~\eqref{sw:eq:prime35} to the total paired polynomial; this proves~\eqref{sw:eq:aggregate-lower}.

For the upper bound, differentiation of both carriers gives
$\norm{u'}_2\le(T+b_N)\norm A_2\le2(T+b_N)\norm u_2$.
For $u\ne0$, Jensen's inequality for the complete multiplier upper bound~\eqref{sw:eq:gamma-upper} yields
\begin{align*}
 \frac{\tg[u]}{\norm u_2^2}
 &\le4+\tfrac12\log\left(1+4\frac{\norm{u'}_2^2}{\norm u_2^2}\right)\\
 &\le4+\tfrac12\log(1+16(T+b_N)^2)
 \le t+C_\Gamma,
\end{align*}
since $T\ge1$ and $b_N\le T/2$ give $T^{-2}+16(1+b_N/T)^2\le37$. On the other hand, $h_T\ge t$ in~\eqref{sw:eq:aggregate-gamma} gives
$\tg[u]-t\norm u_2^2\ge-\epsilon_{T,N}h_T\norm u_2^2$.
Combining the upper and lower Gamma errors, then adding the absolute prime and constant costs, proves~\eqref{sw:eq:aggregate-error}. The zero function needs no division.

For $u=\sum_jc_jv_j$, the linear-first convention makes $c^*Gc=\norm u_2^2$ and $c^*Q_0c=q_0[u]$. Thus~\eqref{sw:eq:aggregate-error} for every coefficient vector is exactly~\eqref{sw:eq:loewner}. If $c\in\ker G$, the represented function is zero, so all its pairings vanish for both forms. On the represented quotient,~\eqref{sw:eq:loewner} is a relative operator-norm estimate. No coefficient Gram inverse or additional compatibility assumption is required.
\end{proof}
The positive principal Gram here is the complete Gamma form measured against the whole ordinary-$L^2$ Gram. The argument does not sum absolute near-window Gamma entries and does not assert a new arithmetic cancellation or a new carrier-kernel identity. Nearly dependent columns are handled through their actual aggregate Gram.

\begin{theorem}[Explicit sufficient contiguous complexity]\label{sw:thm:contiguous}
Let
\begin{equation}\label{sw:eq:delta-eta}
 \delta_* =\frac1{39200},\qquad
 \eta_* =\frac{\delta_*}{200}=\frac1{7840000},\qquad B_*=48K_2.
\end{equation}
Choose $t_*\ge10$ so that
\begin{equation}\label{sw:eq:contiguous-threshold}
 \delta_*t_*^2\ge1,\qquad
 |\cg|\le t_*/8,\qquad
 16B_*^2t_*^{18}e^{-2t_*}\le1/8.
\end{equation}
For $T\ge e^{t_*}$, $0\le L\le N$, and an integer total paired budget
$1\le N\le\delta_*(\log T)^2$, every packet~\eqref{sw:eq:aggregate-packet} satisfies
\begin{equation}\label{sw:eq:contiguous-positive}
 q_0[u]\ge\tfrac12\log T\,\norm u_2^2.
\end{equation}
The threshold is effectively specified, not numerically evaluated.
\end{theorem}
\begin{proof}
All three conditions persist for $t\ge t_*$. In particular,
$d\log(t^{18}e^{-2t})/dt=18/t-2\le-1/5$ for $t\ge10$, and the other two conditions are increasing inequalities. Such a $t_*$ exists.
For $t=\log T$, the budget gives
\[
 N\le t^2,\qquad N+1\le2\delta_*t^2\le2t^2,\qquad
 \log(N+2)\le3\log t.
\]
Thus $b_N\le48K_2t^8\log t\le B_*t^9$, so~\eqref{sw:eq:contiguous-threshold} gives $\epsilon_{T,N}\le1/8$ and $T\ge2b_N$. The sharper prime loss is
\[
 35\sqrt{N+1}\le35\sqrt{2\delta_*}\,t=t/4,
\]
because $35^2\cdot2\delta_*=1/16$. In~\eqref{sw:eq:aggregate-lower}, use $h_T\ge t$ and $|\cg|\le t/8$ to obtain
$(7/8-1/4-1/8)t=t/2$.
\end{proof}
These sufficient constants are intentionally small and are not asserted sharp. For families with $L\le N=N(T)=o((\log T)^2)$, the stronger uniform two-sided conclusion follows from~\eqref{sw:eq:aggregate-error}:
\begin{equation}\label{sw:eq:aggregate-asymptotic}
 E^*_{T,N}/\log T\longrightarrow0,\qquad
 Q_0=(1+o(1))(\log T)G
 \quad\text{in the relative form sense of~\eqref{sw:eq:loewner}}.
\end{equation}
Indeed, eventually $N\le t^2$ and $b_N\le B_*t^9$; hence the leakage term divided by $t$ tends to zero exponentially, while $35\sqrt{N+1}/t\to0$ and the fixed constants divided by $t$ vanish. The conclusion is uniform over every represented coefficient vector in each fixed common prefix list.

\subsection{Actual zero jets in the two regimes}\label{sw:sec:zero}
We use a large-height analytic zero-counting estimate, with no finite-range zero-table extension. Let $N_\zeta(Y)$ count nontrivial zeta zeros with positive ordinate at most $Y$, including multiplicity, and put
\begin{equation}\label{sw:eq:RvM}
 F(Y)=\frac{Y}{2\pi}\log\frac{Y}{2\pi e},\qquad
 Y_*=30610046000.
\end{equation}
Hasanalizade, Shen, and Wong~\cite[Theorem 1.1 and Table 2]{HSW}, on the branch $Y\ge Y_*$, bound $|N_\zeta(Y)-F(Y)+1/8|$. The first table row has constants $0.103787$, $0.257297$, and $9.367419$. Explicitly charging the additional $1/8$, since $9.367419+1/8=9.492419<9.5$, gives
\begin{equation}\label{sw:eq:HSW-large}
 |N_\zeta(Y)-F(Y)|
 \le0.104\log Y+0.258\log\log Y+9.5
 \le10\log Y\qquad(Y\ge Y_*).
\end{equation}
The last step uses $\log\log Y\le\log Y$, $\log Y\ge1$, and $0.104+0.258+9.5=9.862<10$. We use the published large-height theorem as a primary analytical premise; we neither recompute its published constants nor claim they were obtained without computation. The finite zero-database branch used for the finite-range extension in that paper is not invoked here.

Suppose $T-1\ge Y_*$, $t=\log T\ge1$, and $0\le L\le T$. If $M$ is the full multiplicity count in the closed band $[T,T+L]$, then
\[
 M\le N_\zeta(T+L+1)-N_\zeta(T-1).
\]
The larger endpoint is at most $3T$, so its logarithm is at most $3t$, while the lower endpoint has logarithm at most $t$. Both endpoints are in the range of~\eqref{sw:eq:HSW-large}, and their total remainder cost is at most $40t$. Throughout the intervening interval,
$F'(Y)=\log(Y/(2\pi))/(2\pi)\le t$. Its length is $L+2$, hence
\begin{equation}\label{sw:eq:band-count}
 M\le(L+42)t\le42(L+1)t.
\end{equation}
The widened endpoints handle zeros at either target endpoint. This proof needs the full counting main term and its logarithmic remainder, not merely a weak bound of order $Y\log Y$.

For distinct nontrivial zero locations $\rho=\beta+i\gamma$, let $\nu_\rho$ denote multiplicity and define
\begin{equation}\label{sw:eq:jets}
 \lambda_\rho=\gamma+i(1/2-\beta),\qquad
 J_{\rho,k}(x)=W(x)\left.\partial_\lambda^k\cos(\lambda x)
                           \right|_{\lambda=\lambda_\rho},
 \quad0\le k<\nu_\rho.
\end{equation}
Let $V_{T,L}$ be the complex span of all these jets with $T\le\gamma\le T+L$. The critical strip supplies $|\Ima\lambda_\rho|<1/2$ without a critical-line or simplicity assumption. The exact identity
\[
 \partial_\lambda^k\cos(\lambda x)
 =\tfrac12\bigl((ix)^ke^{i\lambda x}+(-ix)^ke^{-i\lambda x}\bigr)
\]
puts every vector of $V_{T,L}$ into~\eqref{sw:eq:aggregate-packet}, with offsets in $[0,L]+i[-1/2,1/2]$ and total paired prefix count at most $2M$. Reflection gives the factor $(-ix)^k$ exactly because $W$ is even.

For the whole actual band use the common integer budget
\begin{equation}\label{sw:eq:padded-band-budget}
 N=\left\lceil\max\{1,L,2M\}\right\rceil
 \le85(L+1)t\le100(L+1)t.
\end{equation}
Indeed, every entry of the maximum is at most $84(L+1)t$, by~\eqref{sw:eq:band-count}, and the ceiling costs less than one, which is at most $(L+1)t$. Padding explicitly ensures $L\le N$, even if the band is empty or contains very few zeros; no lower zero-count bound is assumed.

\begin{corollary}[Finite-constant actual-band branch]\label{sw:cor:contiguous-jets}
Increase the threshold $t_*$ in~\eqref{sw:eq:contiguous-threshold} so that
\begin{equation}\label{sw:eq:actual-threshold}
 e^{t_*}-1\ge Y_*,\qquad \eta_*t_*\ge1.
\end{equation}
For every $T\ge e^{t_*}$, $0\le L\le\eta_*\log T$, and $u\in V_{T,L}$,
\begin{equation}\label{sw:eq:actual-contiguous}
 q_0[u]\ge\tfrac12\log T\,\norm u_2^2.
\end{equation}
Every internal coupling is included, with no internal spacing hypothesis. If $L(T)=o(\log T)$, then uniformly on the whole actual band,
\begin{equation}\label{sw:eq:actual-asymptotic}
 q_0[u]=(1+o(1))\log T\,\norm u_2^2,
\end{equation}
in the two-sided aggregate Gram sense~\eqref{sw:eq:loewner}, for any spanning list of $V_{T,L(T)}$.
\end{corollary}
\begin{proof}
If $L\le\eta_*t$ and $\eta_*t\ge1$, then $L\le t\le T$ and $L+1\le2\eta_*t$. Equation~\eqref{sw:eq:padded-band-budget} gives
$N\le200\eta_*t^2=\delta_*t^2$, so Theorem~\ref{sw:thm:contiguous} applies. For $L=o(t)$, the same padded budget satisfies $N/t^2\le100(L+1)/t\to0$. Equation~\eqref{sw:eq:aggregate-asymptotic} proves the two-sided statement. Empty bands and zero vectors are included without normalization.
\end{proof}
At the fixed endpoint $L=\eta_*\log T$ the estimate proves a positive margin, not a vanishing relative error from this argument alone. The constants are sufficient, and the required heights can be very large; no numerical threshold has been evaluated.

\begin{corollary}[Sparse full unit-window spaces]\label{sw:cor:zeros}
Let $\mathcal J_T=V_{T,1}$. Choose $C_0\ge168$, and enlarge the threshold of Theorem~\ref{sw:thm:sparse} so that $T_0-1\ge Y_*$. Then every finite choice $u_j\in\mathcal J_{T_j}$ along a sequence satisfying~\eqref{sw:eq:spacing} obeys~\eqref{sw:eq:main}, with its stated actual-$X$ convergence and fixed-block closure conclusions.
\end{corollary}
\begin{proof}
At $L=1$, equation~\eqref{sw:eq:band-count} gives $M\le84\log T$. The doubled prefix count is at most $168\log T$. Every full actual unit-window space has a fixed root-prefix list, so it fits the hypotheses of Theorem~\ref{sw:thm:sparse}. If a window is empty the assertion is trivial.
\end{proof}

\subsection{Optimized localization for asymptotic budgets}\label{sw:sec:pnt-local}
The finite estimate~\eqref{sw:eq:prime35} holds for every declared budget. We now obtain a smaller asymptotic leading constant, using a variable anchor closer to the evaluation point and the unconditional prime number theorem. The original derivative scale and constant remain unchanged. In this and the next section a generic polynomial has budget $N$; for a paired packet, $n$ will instead bound only its positive envelope. Its total paired budget can be $2n$.

\begin{lemma}[Optimized mass tail]\label{sw:lem:optimized-tail}
Fix $0<\theta<1$ and
\begin{equation}\label{sw:eq:optimized-parameters}
 a_\theta=\log(4e/\theta),\qquad c>a_\theta/\pi,
 \qquad d=\pi-a_\theta/c>0.
\end{equation}
There is an effectively specifiable integer $N_0=N_0(\theta,c,W)$ such that every $P$ of prefix budget at most $N\ge N_0$, with $|\Ima\lambda_\ell|\le1/2$, obeys
\begin{equation}\label{sw:eq:optimized-tail}
 \norm{\one_{\{|x|>r\}}WP}_2
 \le e^{-\gamma e^{2r}}\norm{WP}_2,
 \qquad e^{2r}\ge cN,\qquad \gamma=d/2.
\end{equation}
The bound is uniform in all real carriers, coefficients, collisions, and confluent prefixes, on the entire stated radius ray.
\end{lemma}
\begin{proof}
Put $U=\norm{WP}_2$, with membership established by the theta envelope as before. The case $U=0$ is immediate. Apply Lemma~\ref{sw:lem:growth} to $F(t)=P(\theta rt)$ at $t=1/\theta$, with exponent real-part bound $\beta=\theta r/2$. This gives
\[
 |P(r)|\le2\sqrt N\,e^{(1+\theta)r/2}(4e/\theta)^N
          (\theta r)^{-1/2}\norm P_{L^2[0,\theta r]}.
\]
The decreasing lower envelope for $W$, rather than monotonicity of $W$ itself, gives
$\norm P_{L^2[0,\theta r]}\le c_W^{-1}e^{-4\theta r+\pi e^{2\theta r}}U$.
For $M=e^{2r}$, the upper envelope therefore yields
\begin{equation}\label{sw:eq:optimized-point}
 \frac{|WP(\pm r)|^2}{U^2}\le
 E_{N,\theta}(r):=
 \frac{4NC_4^2}{\theta r c_W^2}(4e/\theta)^{2N}
 e^{(9-7\theta)r+2\pi M^\theta-2\pi M}.
\end{equation}
The negative sign follows by applying the growth lemma to $P(-\theta rt)$ and using evenness of $W$. The two-sided strip gives the same exponent bound. Compact-interval confluent convergence applies for each fixed anchor, with no real-root or parity assumption.

For $M\ge cN$ and $M\ge e$, put
\[
 H=\max\left\{1,\frac{8C_4^2}{\theta c c_W^2}\right\},
 \qquad D_\theta=\frac{11-7\theta}{2}.
\]
Taking logarithms in~\eqref{sw:eq:optimized-point}, using $N\le M/c$ and discarding $-\log\log M\le0$, gives
\[
 \log E_{N,\theta}(r)
 \le-2dM+\log H+D_\theta\log M+2\pi M^\theta.
\]
Choose $M_0\ge\max(e,1/d)$ with
\begin{equation}\label{sw:eq:optimized-threshold}
 \frac{\log H+D_\theta\log M_0}{M_0}
                  +2\pi M_0^{\theta-1}\le d,
 \qquad N_0=\max\{1,\lceil M_0/c\rceil\}.
\end{equation}
Each term on the left decreases to zero on $M\ge M_0$, so the inequality persists on the whole ray. For $N\ge N_0$ and $M\ge cN$, both point bounds are at most $e^{-dM}U^2$. Integrating both sides gives
\[
 \frac{\norm{\one_{\{|x|>r\}}WP}_2^2}{U^2}
 \le2\int_r^\infty e^{-de^{2x}}\,dx
 =\int_M^\infty\frac{e^{-dv}}v\,dv
 \le\frac{e^{-dM}}{dM}\le e^{-dM}.
\]
Taking square roots proves the claim, including its uniformity.
\end{proof}

For a fixed rational choice used below, take
\begin{equation}\label{sw:eq:rational-anchor}
 \theta=\frac{39}{40},\qquad c=\frac{31}{40}.
\end{equation}
The elementary bounds $\pi>157/50$, $\log2<7/10$, and $\log(40/39)<1/39$ give
\begin{equation}\label{sw:eq:anchor-slack}
 \pi c-a_\theta>
 \frac{157}{50}\frac{31}{40}
       -\left(1+\frac{14}{10}+\frac1{39}\right)
 =\frac{613}{78000}>0.
\end{equation}
Thus $d>d_0:=613/60450$ and one may use $\gamma_0=613/120900$. For complete whole-ray specificity, set
\[
 H_0=\max\left\{1,\frac{1568}{9\theta c}\right\},\qquad
 M_0\ge\max(e,1/d_0),
\]
and require
\begin{equation}\label{sw:eq:rational-anchor-threshold}
 \frac{\log H_0+(167/80)\log M_0}{M_0}
          +\frac{44}{7}M_0^{-1/40}\le d_0.
\end{equation}
Here $C_4<42$, $c_W>9$, and $\pi<22/7$ justify every enlargement. The slow power $M_0^{-1/40}$ is retained; no uniformity as $\theta\uparrow1$ is asserted.

\subsection{The exact compressed prime head}\label{sw:sec:schur}
Write
\[
 \Psi(y)=\sum_{2\le k\le y}\Lambda(k),\qquad
 S(y)=\sum_{2\le k\le y}\frac{\Lambda(k)}{\sqrt k},\qquad
 S_1(y)=\sum_{2\le k\le y}\frac{\Lambda(k)}k.
\]
The sole new arithmetic input is the unconditional prime number theorem
\begin{equation}\label{sw:eq:PNT}
 \Psi(y)=y+o(y)\qquad(y\longrightarrow\infty).
\end{equation}
Both the all-prime-power definition and this asymptotic are recorded in the NIST Digital Library of Mathematical Functions~\cite[25.16.1 and 25.16.3]{DLMF}. The separate RH-strength remainder is not used. We invoke~\eqref{sw:eq:PNT} qualitatively and do not supply an effective remainder.

For an integer $X\ge2$, set
\begin{equation}\label{sw:eq:head-operators}
 \begin{gathered}
 R=\tfrac12\log X,\qquad I=(-R,R),\qquad \Pi f=\one_I f,\\
 H_X=\sum_{2\le k\le X}a_k(\tau_{\log k}+\tau_{-\log k}),
 \qquad B_X=\Pi H_X\Pi.
 \end{gathered}
\end{equation}
The operators are bounded and self-adjoint on ordinary $L^2$, and
$\norm{H_X}\le2S(X)$. For $R\le s<2R$, the two pieces of
$\Pi(\tau_s+\tau_{-s})\Pi f$ have disjoint output intervals $(-R,R-s)$ and $(-R+s,R)$, and, after translation, disjoint input intervals. Hence their squared norms add to at most $\norm f_2^2$. The compressed pair has norm at most one in this range and is zero for $s\ge2R$, including the null endpoint. For smaller $s$ its norm is at most two. These facts alone give the useful preliminary bound
$\norm{B_X}\le S(X)+S(\sqrt X)$, with both orientations included. The following positive weighted-row argument improves its leading constant.

\begin{proposition}[Sharp local compressed norm]\label{sw:prop:schur}
For every $\varepsilon>0$, there are finite nonnegative constants $C_\varepsilon,C'_\varepsilon$ such that
\begin{equation}\label{sw:eq:schur-finite}
 \norm{B_X}\le(1+\varepsilon)(\sqrt X+\log X)
                         +C_\varepsilon+C'_\varepsilon
 \qquad(X\ge2).
\end{equation}
Moreover,
\begin{equation}\label{sw:eq:schur-asymptotic}
 \norm{B_X}=(1+o(1))\sqrt X,
 \qquad
 \norm{B_X|_{L^2_{\rm even}(I)}}=(1+o(1))\sqrt X.
\end{equation}
This sharpness concerns the local compressed operator, not the final prefix-complexity constant.
\end{proposition}
\begin{proof}
The finite translation-edge measure of $B_X$ is nonnegative and symmetric. With the strictly positive weight $w(x)=\cosh(x/2)$ on $I$, apply
\[
 2|f(x)f(y)|\le\frac{w(y)}{w(x)}|f(x)|^2
                    +\frac{w(x)}{w(y)}|f(y)|^2
\]
to each undirected retained edge and integrate. Symmetry, or equivalently adding the two orientations and dividing by two, proves
\begin{equation}\label{sw:eq:atomic-schur}
 |\langle B_Xf,f\rangle_2|
 \le\int_I|f(x)|^2\frac{B_Xw(x)}{w(x)}\,dx,
 \qquad \norm{B_X}\le\mathop{\rm ess\,sup}_{x\in I}\frac{B_Xw(x)}{w(x)}.
\end{equation}
There is no additional factor two. Self-adjointness turns the absolute diagonal bound into the operator norm bound; positivity of the kernel does not assert positivity of its quadratic form.

By reflection it suffices to take $0\le x<R$. Set $a=e^{R-x}$ and $b=e^{R+x}$. A positive shift stays in $I$ when $k<a$, and a negative shift when $k<b$. Expansion of the hyperbolic cosine gives the exact row with strict cutoffs. With inclusive cutoffs it gives the upper bound
\begin{equation}\label{sw:eq:schur-row}
 \frac{B_Xw(x)}{w(x)}\le
 \frac{e^{x/2}\{\Psi(a)+S_1(b)\}
       +e^{-x/2}\{\Psi(b)+S_1(a)\}}
      {e^{x/2}+e^{-x/2}}.
\end{equation}
The equality exceptions for inclusive cutoffs form a finite null set for each fixed $X$. Every prime power remains in the row.

By~\eqref{sw:eq:PNT}, choose $C_\varepsilon\ge0$ such that
$\Psi(y)\le(1+\varepsilon)y+C_\varepsilon$ for every real $y\ge1$: the compact initial interval is covered by enlarging the constant. Partial summation yields
\[
 S_1(y)=\frac{\Psi(y)}y+\int_1^y\frac{\Psi(v)}{v^2}\,dv
 \le(1+\varepsilon)\log y+C'_\varepsilon,
 \qquad C'_\varepsilon=1+\varepsilon+C_\varepsilon.
\]
The main terms in~\eqref{sw:eq:schur-row} simplify exactly:
\begin{align*}
 \frac{e^{x/2}a+e^{-x/2}b}{e^{x/2}+e^{-x/2}}&=e^R=\sqrt X,\\
 \frac{e^{x/2}\log b+e^{-x/2}\log a}{e^{x/2}+e^{-x/2}}
 &=R+x\tanh(x/2)\le2R=\log X.
\end{align*}
Each pair of constant errors is one convex weighted average. Equation~\eqref{sw:eq:schur-finite} follows. Fixing $\varepsilon$ before taking $X\to\infty$, then letting it decrease to zero, proves the upper limiting coefficient one.

For the matching lower bound, PNT likewise supplies $C^-_\varepsilon\ge0$ with $\Psi(y)\ge(1-\varepsilon)y-C^-_\varepsilon$ on all $y\ge1$. The exact row almost everywhere, nonnegative $S_1$ terms, and the same identity give
\[
 \frac{B_Xw(x)}{w(x)}\ge(1-\varepsilon)\sqrt X-C^-_\varepsilon.
\]
Test the Rayleigh quotient on the even function $w\in L^2(I)$. This proves the matching lower bound on the even subspace and hence on the full space. Reflection commutes with $B_X$, so the restriction is legitimate.
\end{proof}
This proof treats the exact atomic operator. PNT does not provide, and the argument does not use, operator-norm approximation by a continuous-density model. No phase cancellation, Hilbert-transform cancellation, or arithmetic rank subtraction is invoked.

\subsection{Complete prime asymptotics and logarithmic zero bands}\label{sw:sec:pnt-bands}
\subsubsection{From physical concentration to the full prime form}
Suppose a nonzero $u\in L^2(\mathbb R)$ satisfies, with fixed $D\ge1$, $\gamma>0$, $c>0$, and integer $m$,
\begin{equation}\label{sw:eq:abstract-tail}
 \norm{\one_{\{|x|>r\}}u}_2\le D e^{-\gamma e^{2r}}U,
 \qquad e^{2r}\ge cm,\qquad U=\norm u_2.
\end{equation}
Let $X=\lceil cm\rceil\ge2$, and take $\Pi$ from~\eqref{sw:eq:head-operators}. Write $i=\Pi u$, $o=u-i$. Then $\norm i_2\le U$ and $\norm o_2\le De^{-\gamma X}U$. The leading compressed head is charged to the energy of $i$, not to an envelope energy. The exact difference identity and the full-line head norm give
\begin{align*}
 \langle H_Xu,u\rangle_2-\langle H_Xi,i\rangle_2
 &=\langle H_Xo,u\rangle_2+\langle H_Xi,o\rangle_2,\\
 |\langle H_Xu,u\rangle_2-\langle H_Xi,i\rangle_2|
 &\le4D S(X)e^{-\gamma X}U^2.
\end{align*}
It would be incorrect to use the smaller compressed norm to bound these uncompressed errors. The hard-cut $i$ is used only for bounded ordinary-$L^2$ operators; no Gamma or $q_0$ form is applied to it.

The correlation split~\eqref{sw:eq:correlation}, for every integer $k>X$, gives
$|\langle\tau_{\log k}u,u\rangle_2|\le2De^{-\gamma k}U^2$.
Since $0\le a_k\le1$, all remaining prime powers and both orientations cost at most
$4De^{-\gamma(X+1)}U^2/(1-e^{-\gamma})$. This also proves absolute convergence. Thus, for every fixed $\varepsilon>0$,
\begin{equation}\label{sw:eq:pnt-complete-bound}
 \begin{split}
 \frac{|p[u]|}{U^2}\le{}&(1+\varepsilon)(\sqrt X+\log X)
                  +C_\varepsilon+C'_\varepsilon\\
 &+4D S(X)e^{-\gamma X}
       +\frac{4De^{-\gamma(X+1)}}{1-e^{-\gamma}}.
 \end{split}
\end{equation}
PNT and partial summation give
\begin{equation}\label{sw:eq:S-asymptotic}
 S(y)=\frac{\Psi(y)}{\sqrt y}
       +\frac12\int_1^y\frac{\Psi(v)}{v^{3/2}}\,dv
     =2\sqrt y+o(\sqrt y).
\end{equation}
To justify the integrated remainder, for any tolerance bound $|\Psi(v)-v|$ by that tolerance times $v$ on one entire eventual ray; its integral there is at most the tolerance times $\sqrt y$. The compact initial integral is fixed and becomes negligible. Taking $m\to\infty$ with $c,D,\gamma,\varepsilon$ fixed in~\eqref{sw:eq:pnt-complete-bound}, and only then letting $\varepsilon\downarrow0$, yields the uniform leading bound
\begin{equation}\label{sw:eq:concentration-prime}
 \limsup_{m\to\infty}\ \sup_u
       \frac{|p[u]|}{\sqrt m\,\norm u_2^2}\le\sqrt c,
\end{equation}
where the supremum is over functions satisfying~\eqref{sw:eq:abstract-tail} with these common constants.

\begin{theorem}[Uniform asymptotic prefix bounds]\label{sw:thm:pnt-prime}
Put $a_0=\log(4e)$. For every $\zeta>0$, all sufficiently large generic budgets $N$ satisfy
\begin{equation}\label{sw:eq:pnt-generic}
 |p[WP]|\le\left(\sqrt{a_0/\pi}+\zeta\right)\sqrt N\,\norm{WP}_2^2
\end{equation}
for every strip polynomial of prefix at most $N$. The same coefficient holds with the positive-envelope budget $n$ for every paired packet
\begin{equation}\label{sw:eq:positive-budget-packet}
 u(x)=\tfrac12\{e^{iTx}A(x)+e^{-iTx}A(-x)\},\qquad A=WQ,
\end{equation}
where $Q$ has prefix at most $n$, frequencies in $[0,L]+i[-1/2,1/2]$,
$0\le L\le n$, and $T\ge\max(1,2b_n)$:
\begin{equation}\label{sw:eq:pnt-paired}
 |p[u]|\le\left(\sqrt{a_0/\pi}+\zeta\right)\sqrt n\,\norm u_2^2.
\end{equation}
Each eventual budget threshold is uniform over all these coefficients, frequencies, widths, and carriers. These are asymptotic statements, not replacements for~\eqref{sw:eq:prime35} at every finite budget.
\end{theorem}
\begin{proof}
For generic $WP$, Lemma~\ref{sw:lem:optimized-tail} supplies~\eqref{sw:eq:abstract-tail} with $m=N,D=1$. For a paired packet, Proposition~\ref{sw:prop:wide-derivative} uses the positive envelope's own budget $n$ and gives $\norm{A'}_2\le b_n\norm A_2$. The reflected-carrier integration by parts of Appendix~\ref{sw:sec:carrier} then gives
\[
 \norm u_2^2\ge\tfrac12(1-b_n/T)\norm A_2^2,
 \qquad \norm A_2\le2\norm u_2.
\]
Apply the optimized tail to $A$ with budget $n$. Reflection preserves its two-sided tail norm, real carriers have modulus one, and the factors $1/2$ sum to one. Consequently
\[
 \norm{\one_{\{|x|>r\}}u}_2
 \le2e^{-\gamma e^{2r}}\norm u_2,
 \qquad e^{2r}\ge cn,\qquad n\ge N_0.
\]
Thus~\eqref{sw:eq:concentration-prime} applies with $m=n,D=2$. The envelope-to-packet factor two occurs only in exponentially small errors in~\eqref{sw:eq:pnt-complete-bound}; the leading term uses $\norm{\Pi u}_2^2\le\norm u_2^2$. We have not applied a generic $2n$-budget estimate and then dropped a factor.

For any fixed $c>a_0/\pi$, choose a fixed $\theta<1$ close enough to one that $c>a_\theta/\pi$. With those choices fixed, apply~\eqref{sw:eq:concentration-prime}. Finally infimize over $c>a_0/\pi$. This proves both tolerance statements. Neither $\theta$ nor $\varepsilon$ varies with the budget inside a limit, and no uniform-in-$\theta$ threshold is needed.
\end{proof}

\subsubsection{The full actual zero band and its positive budget}
The Gamma comparison also needs only the positive-envelope budget. The same proof as Theorem~\ref{sw:thm:aggregate} gives, for~\eqref{sw:eq:positive-budget-packet},
\begin{equation}\label{sw:eq:positive-budget-gamma}
 h_T(1-16b_n^2/T^2)\norm u_2^2
 \le\tg[u]\le(\log T+C_\Gamma)\norm u_2^2,
 \qquad h_T=\tfrac12\log(1+T^2).
\end{equation}
Indeed, Plancherel bounds the low-frequency norm by $4b_n\norm u_2/T$, and differentiation gives $\norm{u'}_2\le2(T+b_n)\norm u_2$. The lower full multiplier and Jensen upper bound apply exactly as before. For each fixed $C>0$, uniformly on $1\le n\le C(\log T)^2$ and $L\le n$, one has $b_n=O_C((\log T)^8\log\log T)$, so eventually $T\ge2b_n$ and
\begin{equation}\label{sw:eq:positive-gamma-asymptotic}
 \tg[u]=(1+o(1))\log T\,\norm u_2^2.
\end{equation}
This includes the complete unbounded multiplier and all internal pairings.

Fix $A>0$, put $t=\log T$, $L=At$, and let $M=M(T,L)$ count all positive-height zeros in the closed band with multiplicities. Use the same large-height HSW premise~\eqref{sw:eq:HSW-large} and widened endpoints as in Section~\ref{sw:sec:zero}. Since $F'(Y)=\log(Y/(2\pi))/(2\pi)$, uniformly between these endpoints,
\begin{align}
 F(T+L+1)-F(T-1)
 &=\frac{(L+2)(t-\log(2\pi))}{2\pi}+O_A(t^2/T)\notag\\
 &=\frac{A}{2\pi}t^2+O_A(t).
\end{align}
The two logarithmic counting remainders are $O_A(t)$, so
\begin{equation}\label{sw:eq:sharp-band-count}
 M(T,At)\le\frac{A}{2\pi}t^2+O_A(t).
\end{equation}
The leading coefficient counts positive-height zeros only; there is no extra negative-height multiplicity.

For any complex linear combination of the jets~\eqref{sw:eq:jets}, choose
\[
 Q(x)=\sum_{\rho,k}d_{\rho,k}(ix)^k
                    e^{i(\lambda_\rho-T)x}.
\]
Then~\eqref{sw:eq:positive-budget-packet} reproduces the whole combination exactly. Its reflected partner uses $Q(-x)$, not its complex conjugate. The positive envelope prefix is at most $M$, although the total paired prefix can be $2M$. Declare the common positive budget
\begin{equation}\label{sw:eq:positive-padded-budget}
 n(T)=\left\lceil\max\{1,L,M(T,L)\}\right\rceil,
 \qquad
 At\le n(T)\le\frac{A}{2\pi}t^2+O_A(t).
\end{equation}
The padding guarantees $L\le n$ without a lower zero-density premise, and $n(T)\to\infty$. It supplies a fixed common prefix space for the whole actual band.

\begin{proof}[Proof of Theorem~\ref{sw:thm:pnt-bands}]
Fix $\theta,c$ as in~\eqref{sw:eq:optimized-parameters}. Apply~\eqref{sw:eq:concentration-prime} to the paired class with $m=n(T),D=2$, then use~\eqref{sw:eq:positive-padded-budget}. Uniformly on the nonzero vectors of the full actual space,
\[
 \limsup_{T\to\infty}\ \sup_{0\ne u\in V_{T,A\log T}}
 \frac{|p[u]|}{\log T\,\norm u_2^2}
 \le\sqrt{\frac{cA}{2\pi}}.
\]
The Gamma coefficient tends uniformly to one by~\eqref{sw:eq:positive-gamma-asymptotic}, while $\cg/\log T\to0$. Taking the fixed-parameter infimum over $c>a_0/\pi$ proves~\eqref{sw:eq:pnt-band-liminf}. This means that for every tolerance the lower coefficient minus that tolerance holds for every sufficiently large real $T$ and every vector in its band, not merely along a sequence. Empty spaces are interpreted vacuously.

For the concrete width $A=8$, use the fixed rational parameters~\eqref{sw:eq:rational-anchor}, whose admissibility is proved in~\eqref{sw:eq:anchor-slack}. The prime limiting coefficient is at most
\[
 \sqrt{\frac{c\,8}{2\pi}}=\sqrt{\frac{31}{10\pi}}<\frac{199}{200}.
\]
Indeed, $\pi>157/50$ gives $31/(10\pi)<155/157<(199/200)^2$, and the last comparison is the exact integer inequality
$157\cdot39601-155\cdot40000=17357>0$.
Thus the lower limiting coefficient exceeds $1/200$, so the conservative margin $1/1000$ holds after one whole-ray threshold. Alternatively fix $\varepsilon=1/1000$ in~\eqref{sw:eq:pnt-complete-bound}: its limiting prime coefficient is below
$(1001/1000)(199/200)=199199/200000<997/1000$, leaving the stated margin after the Gamma and constant errors vanish.
\end{proof}
The constants $C_\varepsilon,C'_\varepsilon$ and their eventual PNT thresholds are only existential here. Accordingly the eight-logarithmic-band theorem supplies no evaluated or effective final height. Its elementary localization threshold alone is specified by~\eqref{sw:eq:rational-anchor-threshold}. This differs from the explicitly specified, unevaluated threshold rules of the earlier finite-constant and sparse branches.

For any spanning list of a fixed common prefix space, including a redundant one, every displayed lower bound is equivalent to the corresponding Loewner inequality against the ordinary represented-function Gram. In particular,~\eqref{sw:eq:eight-band} gives $Q_0\succeq(\log T)G/1000$ for the full actual band. A vector in $\ker G$ represents zero, so every form pairing vanishes; no coefficient inverse or conditioning assumption is required.

\subsection{Flexible anchors and eleven logarithmic zero bands}\label{bl:sec:anchor}
The interior anchor in the growth estimate can vary while the arithmetic and Gamma estimates remain unchanged. This improves the sufficient band width and permits a sparse family of entire growing bands. We retain the ordinary norm, full forms and positive-envelope packets of Section~\ref{sw:sec:pnt-bands}. In particular, $b_n=K_2(n+1)^4\log(n+2)$, a packet has real offsets in $[0,L]$ with $L\le n$, and its reflected partner is $Q(-x)$, not $\overline{Q(x)}$.

\begin{lemma}[Flexible interior anchor]\label{bl:anchorlemma}
Let $F$ have prefix budget at most $N$ and exponent real parts bounded in absolute value by $\beta$. For $0<\delta<1/2$ and $t\ge1$,
\begin{equation}\label{bl:flex}
 |F(t)|\le e^{\beta t}\sqrt{N/\delta}
 \left(\frac{2e(t-\delta)}{1-2\delta}\right)^N\norm F_{L^2[0,1]}.
\end{equation}
\end{lemma}
\begin{proof}
For each $y\in[\delta,1-\delta]$, \eqref{sw:eq:erd-interior} bounds $|F(y)|$ by $\sqrt{N/\delta}\norm F_{L^2[0,1]}$. The stronger first inequality of Erd\'elyi's Lemma~3.2~\cite{Erdelyi} is
\[
 |G(t)|\le\left(\frac{2e(t-a)}{d}\right)^N
 \norm G_{L^\infty[a,a+d]},\qquad
 a,d>0,\quad t\ge a+d,\quad \Rea z_j\le0.
\]
Apply it to $G(t)=e^{-\beta t}F(t)$ with $a=\delta,d=1-2\delta$. Damping is at most one on the anchor; undamping at the evaluation point gives $e^{\beta t}$. For a confluent prefix, replace $x^k$ by $((e^{i\varepsilon x}-1)/(i\varepsilon))^k$ and let real $\varepsilon\to0$. A prefix of length $m$ becomes at most $m$ exponentials with unchanged exponent real parts, and compact-interval uniform convergence passes the bound to the limit. The constants are independent of all coefficients and frequency collisions.
\end{proof}

\begin{lemma}[Improved normalized physical tail]\label{bl:newtail}
Fix $0<\theta<1$, $0<\delta<1/2$ and set
\begin{equation}\label{bl:anew}
 a_{\theta,\delta}=\log\left(\frac{2e(\theta^{-1}-\delta)}{1-2\delta}\right).
\end{equation}
If $c>a_{\theta,\delta}/\pi$, then there are fixed $\gamma>0$ and $N_0$ such that every strip polynomial $P$ of budget at most $N\ge N_0$ satisfies
\begin{equation}\label{bl:newtailformula}
 \norm{\one_{|x|>r}WP}\le e^{-\gamma e^{2r}}\norm{WP},\quad e^{2r}\ge cN.
\end{equation}
The statement is uniform in all real carriers, coefficients and confluences.
\end{lemma}
\begin{proof}
Let $U=\norm{WP}>0$ and apply Lemma \ref{bl:anchorlemma} to $F(t)=P(\theta rt)$ at $t=\theta^{-1}$, with $\beta=\theta r/2$. The decreasing lower theta envelope $c_We^{4r-\pi e^{2r}}$, with $c_W>9$ and upper-envelope constant $C_4<42$, gives
\begin{align}
 |P(r)|&\le e^{r/2}\sqrt{\frac{N}{\delta\theta r}}
 e^{Na_{\theta,\delta}}\norm P_{L^2[0,\theta r]},\\
 \norm P_{L^2[0,\theta r]}&\le c_W^{-1}e^{-4\theta r+\pi e^{2\theta r}}U.
\end{align}
Consequently, putting $M=e^{2r}$,
\begin{equation}\label{bl:newdensity}
 \frac{|WP(\pm r)|^2}{U^2}
 \le\frac{C_4^2N}{c_W^2\delta\theta r}
 \exp\{2Na_{\theta,\delta}+(9-8\theta)r+2\pi M^\theta-2\pi M\}.
\end{equation}
Reflection proves the negative tail, using the other side of the strip. Let $d=\pi-a_{\theta,\delta}/c>0$. For $M\ge\max(e,cN)$, the logarithm of the right side is at most
\begin{equation}
 -2dM+\log H+\frac{11-8\theta}{2}\log M+2\pi M^\theta,
 \quad H=\max\left(1,\frac{2C_4^2}{c_W^2\delta\theta c}\right).
\end{equation}
Choose $M_0\ge\max(e,1/d)$ such that the last three terms, divided by $M_0$, sum to at most $d$. Each term decreases on this ray, so the same holds for every $M\ge M_0$. Take $N_0\ge\lceil M_0/c\rceil$. Each pointwise squared tail is at most $e^{-dM}U^2$, and integrating both tails gives
\begin{equation}
 \frac{\norm{\one_{|x|>r}WP}^2}{U^2}
 \le\int_M^\infty\frac{e^{-dv}}v\,dv
 \le\frac{e^{-dM}}{dM}\le e^{-dM}.
\end{equation}
Thus $\gamma=d/2$ suffices. Membership and the zero case follow directly from the theta envelope. This proof supplies a symbolic whole-ray localization threshold for fixed parameters.
\end{proof}

\begin{theorem}[Improved asymptotic budget and actual-band bounds]\label{bl:improvement}
For every $\zeta>0$ and sufficiently large budgets, the generic strip-polynomial prime bound is
\begin{equation}
 |p[WP]|\le\left(\sqrt{\frac{\log(2e)}\pi}+\zeta\right)\sqrt N\norm{WP}^2.
\end{equation}
For packets \eqref{sw:eq:positive-budget-packet} with positive budget $n$, $L\le n$ and $T\ge\max(1,2b_n)$, the same bound holds with $N$ replaced by $n$ and $WP$ replaced by $u_T$. Moreover, for every fixed $A>0$,
\begin{equation}\label{bl:newband}
 \liminf_{T\to\infty}\inf_{0\ne u\in V_{T,A\log T}}
 \frac{q_0[u]}{\log T\norm u^2}
 \ge\kappa_{\rm band}(A):=1-\sqrt{\frac{A\log(2e)}{2\pi^2}}.
\end{equation}
Thus every fixed $0<A<A_*=2\pi^2/\log(2e)\approx11.658295$ gives eventual strict positivity on the whole band.
\end{theorem}
\begin{proof}
For every fixed $c>\log(2e)/\pi$, choose fixed $\theta<1$ close to one and fixed $\delta>0$ close to zero so that $c>a_{\theta,\delta}/\pi$. This is possible because $a_{\theta,\delta}\to\log(2e)$. Apply Lemma \ref{bl:newtail} and the exact prime estimate \eqref{sw:eq:pnt-complete-bound}. For generic polynomials use $m=N,D=1$. For packets, the positive envelope has its own budget $n$ and $\norm A_2\le2\norm u_2$ for $T\ge2b_n$. Its tail and its reflection therefore give $m=n,D=2$. The leading compressed-prime term is charged to $\norm{\Pi u}_2^2\le\norm u_2^2$; the factor two appears only in the exponentially small head-replacement and prime-tail errors. No generic $2n$ estimate has been divided by two. Take the budget limit with these constants fixed, then infimize over $c>\log(2e)/\pi$. This proves both prime inequalities with any tolerance. For the actual band, use the same common padded budget
\[
 n(T)=\lceil\max\{1,A\log T,M(T,A\log T)\}\rceil
 \le\frac{A}{2\pi}(\log T)^2+O_A(\log T).
\]
This enforces $L\le n$ even if the actual zero count is small, includes all multiplicity prefixes, and tends to infinity. Equation~\eqref{sw:eq:positive-budget-gamma} applies uniformly on these budgets, so the full Gamma coefficient tends to one while $\cg/\log T\to0$. The prime coefficient is at most $\sqrt{cA/(2\pi)}$ in the height limit. Taking the fixed-parameter infimum proves~\eqref{bl:newband}. No limit requires a threshold uniform in $\theta$ or $\delta$. No claim is made at $A=A_*$.
\end{proof}

\begin{corollary}[A concrete eleven-logarithmic band]\label{bl:eleven}
There exists a height $T_{11}$ such that, for every real $T\ge T_{11}$,
\begin{equation}\label{bl:eleveneq}
 q_0[u]\ge\frac1{100}\log T\norm u^2,\qquad u\in V_{T,11\log T}.
\end{equation}
\end{corollary}
\begin{proof}
One can use fixed rational parameters $\theta=99/100$, $\delta=1/100$, $c=11/20$. Here
\begin{equation}
 a_{\theta,\delta}=1+\log2+\log(9901/9702).
\end{equation}
Using $\pi>157/50$, $\log2<7/10$ and $\log(1+x)<x$ gives the exact positive margin
\begin{equation}
 \pi c-a_{\theta,\delta}>
 \frac{157}{50}\frac{11}{20}-\frac{17}{10}-\frac{199}{9702}
 =\frac{31477}{4851000}>0.
\end{equation}
The limiting prime coefficient is at most
\begin{equation}
 \sqrt{\frac{121}{40\pi}}<\sqrt{\frac{605}{628}}<\frac{99}{100},
\end{equation}
where $628\cdot9801-605\cdot10000=105028>0$. The Gamma coefficient tends to one and $c_\Gamma/\log T\to0$, so the lower limit strictly exceeds $1/100$. This proves the eventual whole-ray assertion. The final PNT-dependent height remains unevaluated here.
\end{proof}

\subsection{The residual form across an internal band cut}\label{bl:sec:adjacent}
Let $I_1,I_2$ be adjacent or overlapping frequency intervals contained in $[T,T+A\log T]$, and let $U,V$ be their actual zero-jet spaces. All multiplicity prefixes are retained. Fix $0<A<A_*$ and $0<\mu<\kappa_{\rm band}(A)$. For sufficiently large $T$, Theorem \ref{bl:improvement} applies to $U+V$.

\begin{proposition}[Residual contraction across an internal cut]\label{bl:adjacent}
Set $r(u,v)=q_0(u,v)-\mu\log T\ip uv$. Then, for $u\in U$ and $v\in V$,
\begin{equation}\label{bl:csresidual}
 \boxed{|q_0(u,v)-\mu\log T\ip uv|^2
 \le(q_0[u]-\mu\log T\norm u^2)
      (q_0[v]-\mu\log T\norm v^2).}
\end{equation}
For arbitrary spanning lists in the two spaces, write $R_{ij}=Q_{ij}-\mu\log T\,G_{ij}$. Then
\begin{equation}\label{bl:blockresidual}
 \begin{pmatrix}R_{11}&R_{12}\\R_{12}^*&R_{22}\end{pmatrix}\succeq0,
 \qquad R_{12}=R_{11}^{1/2}K R_{22}^{1/2},\quad\norm K\le1,
\end{equation}
on the represented supports. If $R_{22}$ is positive definite, this includes the operator inequality
\begin{equation}\label{bl:schurresidual}
 R_{12}R_{22}^{-1}R_{12}^*\preceq R_{11}.
\end{equation}
\end{proposition}
\begin{proof}
The merged-band theorem says $r[w]\ge0$ for every $w\in U+V$. Apply it to $u+zv$ for every $z\in\C$. Minimizing the resulting quadratic polynomial proves \eqref{bl:csresidual}; if a diagonal term is zero, its mixed term must vanish. The same argument for arbitrary coefficient vectors gives the block positivity and contraction factorization. Completing the square proves \eqref{bl:schurresidual}. More explicitly, block positivity gives $\ker R_{22}\subseteq\ker R_{12}$ and $\ker R_{11}\subseteq\ker R_{12}^*$. On the two ranges, residual Cauchy--Schwarz defines a contraction between the square-root ranges, which gives the stated factorization. In finite dimension, the singular Schur statement is
\[
 R_{12}R_{22}^{\dagger}R_{12}^*\preceq R_{11},\qquad
 \ker R_{22}\subseteq\ker R_{12},
\]
where $\dagger$ denotes the Moore--Penrose inverse. Both the kernel condition and the inequality are needed; an ordinary inverse of a singular Gram is never used.
\end{proof}

This is a control of the \emph{actual cross form}, with the mass cross term retained. It is stronger than separately asserting the two diagonal blocks are positive, but it is not a small off-diagonal estimate in the sum of separate block norms. No orthogonality of neighboring theta packets is asserted.

For example, intervals $[T,T+5\log T]$ and $[T+5\log T,T+10\log T]$ fit in one merged band and satisfy
\begin{equation}
 q_0[u+v]\ge\frac7{100}\log T\norm{u+v}^2
\end{equation}
eventually, since $\kappa_{\rm band}(10)\approx0.0738475>0.07$. Alternatively, any two pieces whose union lies in the eleven-logarithmic band enjoy \eqref{bl:eleveneq} and Proposition \ref{bl:adjacent} with $\mu=1/100$. Theorem~\ref{sw:thm:pnt-bands} already merges two four-logarithmic pieces into an eight-logarithmic band. Two full eight-logarithmic pieces have width $16\log T$ and lie beyond both sufficient ranges; a negative sufficient coefficient at $16$ is \emph{not} evidence of actual negativity.

\subsection{Sparse families of entire growing bands}\label{bl:sec:sparse}
We now allow a whole growing band in each sparse slot. The following retains all internal couplings and all cross-slot prime and Gamma terms.

\begin{lemma}[Cross bound for growing packets]\label{bl:cross}
Let $u_T,v_S$ have positive-envelope budgets $n,m$ and real offsets in $[0,L_T]$, $[0,L_S]$, with $L_T\le n$ and $L_S\le m$. Suppose $T,S\ge2$ and $T\ge2b_n,S\ge2b_m$. Let $R=\max(T,S)\ge2$, $D=|T-S|>0$, and
\begin{equation}
 G_R=4+\log(1+8R^2),\quad
 X=\left\lceil\max\{8(n+1),8(m+1),2\log(2+R),2\}\right\rceil.
\end{equation}
Then
\begin{equation}\label{bl:cross-explicit}
 |q_0(u_T,v_S)|\le4\left[(2G_R+|c_\Gamma|+12\sqrt X)
 \frac{b_n+b_m}{D}+\frac4{(2+R)^2}\right]\norm{u_T}\norm{v_S}.
\end{equation}
\end{lemma}
\begin{proof}
For two unpaired carriers $e^{i\alpha x}A,e^{i\beta x}B$, split the Fourier line at distance $|\alpha-\beta|/2$. The full Gamma bound and its derivative-weighted tail give the mixed cost $2G_R(b_n+b_m)/|\alpha-\beta|$. Integration by parts gives mass cost $(b_n+b_m)/|\alpha-\beta|$. The same argument for each translated correlation and $\sum_{k\le X}a_k\le6\sqrt X$ gives the two-oriented prime-head cost $12\sqrt X(b_n+b_m)/|\alpha-\beta|$.

For $k>X$, at least one translated envelope lies beyond $\frac12\log k$. Equation \eqref{sw:eq:linear-tail}, applied to each envelope's own budget, gives the full prime tail at most $4e^{-X}\norm A\norm B$. Each of the four paired carrier differences has modulus at least $D$. Their expansion coefficients have absolute values summing to one, and $\norm A\norm B\le4\norm{u_T}\norm{v_S}$. Finally $e^{-X}\le(2+R)^{-2}$. This is precisely \eqref{bl:cross-explicit}; the growing-offset derivative estimate is Proposition~\ref{sw:prop:wide-derivative}.
\end{proof}

Assume $n\le C(\log T)^2$, $m\le C(\log S)^2$ and $\log T,\log S\ge1$. Define
\begin{align}
 B&=K_2(C+1)^4(\log(C+2)+2),\\
 C_X&=8(C+1)+5+2\log3,\quad G_0=6+\log9,\\
 E&=2G_0+|c_\Gamma|+12\sqrt{C_X},\qquad C_*=8EB+16.
\end{align}
Then $b_n,b_m\le B(\log R)^9$, $X\le C_X(\log R)^2$, and \eqref{bl:cross-explicit} implies
\begin{equation}\label{bl:cross-simplified}
 |q_0(u_T,v_S)|\le C_*\frac{(\log R)^{10}}{|T-S|}\norm{u_T}\norm{v_S}.
\end{equation}
The last tail is absorbed using $D\le R$. The constants are deliberately conservative.

\begin{theorem}[Sparse families of entire logarithmic bands]\label{bl:sparsebands}
Fix $0<A<A_*$ and $0<\mu<\kappa_{\rm band}(A)$. There exists $T_0\ge e^{26}$ such that the following holds. Put
\begin{equation}
 \Phi_{12}(T)=\int_{e^{26}}^T\frac{ds}{(\log s)^{12}}.
\end{equation}
If $T_j\ge T_0$ are increasing and $\Phi_{12}(T_{j+1})-\Phi_{12}(T_j)\ge1$, then every finite choice $u_j\in V_{T_j,A\log T_j}$ satisfies
\begin{equation}\label{bl:sparse-result}
 q_0\left[\sum_j u_j\right]\ge\frac\mu2\sum_j\log T_j\norm{u_j}^2.
\end{equation}
The same estimate passes to a fixed series whose partial sums actually converge in the form norm
$\norm u_X^2=\norm u_{\mathcal H}^2+\tg[u]$. Nonnegativity also holds on the $X$-closure of the algebraic sum of these fixed actual block spaces. In particular, using only Theorem~\ref{sw:thm:pnt-bands}, one may take $A=8,\mu=1/1000$ and obtain the coefficient $1/2000$ on the right.
\end{theorem}
\begin{proof}
The padded positive budget in~\eqref{sw:eq:positive-padded-budget} supplies a common $C$ with $n_j\le C(\log T_j)^2$ for every whole block. After increasing $T_0$, its diagonal is at least $\mu\log T_j\norm{u_j}^2$ and $T_j\ge2b_{n_j}$. Normalize $z_j=\sqrt{\log T_j}\norm{u_j}$ and define
\begin{equation}
 k(T,S)=\frac{C_*\max(\log T,\log S)^{10}}
 {|T-S|\sqrt{\log T\log S}}.
\end{equation}
The packing argument of Section~\ref{sw:sec:packing} applies with the new exponent. We give the full row estimates to include the infinite far tail.

Fix $T=e^t$, $t\ge t_0=\log T_0\ge26$. Near heights $S\in[T/2,2T]$ have logarithms in $[t/2,2t]$, so $k(T,S)\le2^{11}C_*t^9/|T-S|$. Successive heights there are separated by at least $\Delta=(t/2)^{12}$. The two harmonic sums cost at most $4t/\Delta$, giving
\begin{equation}
 \sum_{T/2\le S\le2T,\ S\ne T}k(T,S)\le2^{25}C_*t^{-2}.
\end{equation}
For a nonnegative decreasing function $f$, unit spacing in $\Phi_{12}$ gives
$\sum f(S)\le f(U)+\int_U^V f(s)(\log s)^{-12}\,ds$.
For the lower far tail, use $f(s)=(\log s)^{-1/2}$ and
\begin{equation}
 \frac{d}{ds}\frac{s}{(\log s)^{25/2}}
 =\frac{1-25/(2\log s)}{(\log s)^{25/2}}
 \ge\frac1{2(\log s)^{25/2}}\quad(s\ge e^{26}).
\end{equation}
It follows that
\begin{equation}
 \sum_{S<T/2} k(T,S)
 \le\frac{2C_*t^{19/2}e^{-t}}{\sqrt{t_0}}+2^{27/2}C_*t^{-3}.
\end{equation}
For the upper far tail use the decreasing $f(s)=(\log s)^{19/2}/s$. Its density-weighted integral is
\begin{equation}
 \int_{2T}^\infty\frac{ds}{s(\log s)^{5/2}}
 =\tfrac23(\log(2T))^{-3/2},
\end{equation}
which gives
\begin{equation}
 \sum_{S>2T} k(T,S)\le2^{19/2}C_*t^9e^{-t}+\tfrac43 C_*t^{-2}.
\end{equation}
All three bounds are uniform in the number of blocks. Therefore
\begin{equation}\label{bl:E12}
 \sup_i\sum_{j\ne i}k(T_i,T_j)\le\mathcal E_{12}(t_0):=
 C_*\left[\frac{2^{25}+4/3}{t_0^2}+\frac{2^{27/2}}{t_0^3}
 +(2+2^{19/2})t_0^9e^{-t_0}\right].
\end{equation}
This decreases to zero for $t_0\ge26$. Choose $\mathcal E_{12}(t_0)\le\mu/2$. Expanding the actual form and using $2z_iz_j\le z_i^2+z_j^2$ proves \eqref{bl:sparse-result}. Actual $X$ convergence and Lemma~\ref{sw:lem:domain} pass the finite inequalities to the limit; the right side is monotone and therefore finite. Continuity also gives nonnegativity on the stated $X$-closure. This is not a claim that every abstract square-summable coefficient series converges in $X$.
\end{proof}

The exponent $12$ is sufficient, not optimized. A stronger easily checked spacing rule is $T_{j+1}-T_j\ge(\log T_{j+1})^{12}$. These are widely separated growing bands, not an adjacent covering. The proof connects the contiguous estimate to the old sparse-row mechanism, but cannot fill the gaps between its slots.

\subsection{Limits of adjacent and weighted assembly}\label{bl:sec:limits}
Positive neighboring pieces do not, on their own, prove positivity of a long chain. For example,
\begin{equation}\label{bl:counterexample}
 M=\begin{pmatrix}1&-3/4&0\\-3/4&1&-3/4\\0&-3/4&1\end{pmatrix}
\end{equation}
has positive adjacent $2\times2$ principal blocks, but has eigenvalue $1-3\sqrt2/4<0$. This is a logical counterexample to the inference, not a theta-matrix counterexample. Even a bound of one on each normalized pair contraction is insufficient for a full row of interactions.

For a finite family of represented block spaces, let $Q=(Q_{ij})$ be the exact matrix, including every cross entry. If each diagonal block is positive definite on its represented quotient, set $D=\operatorname{diag}(Q_{ii})$ and
\begin{equation}
 \mathcal C=D^{-1/2}(Q-D)D^{-1/2}.
\end{equation}
Then
\begin{equation}\label{bl:global-condition}
 Q\succeq0\quad\Longleftrightarrow\quad I+\mathcal C\succeq0.
\end{equation}
The stronger sufficient condition $\norm{\mathcal C}\le1$ or a normalized absolute block-row sum at most one is useful but need not be necessary. The sparse theorem proves such a row condition in its own geometry. No such estimate has been obtained here for an adjacent covering of all high zeros.

For a target weighted lower bound, replace $Q$ by the \emph{actual shifted matrix} $Q-\lambda H$, with $H$ the reciprocal-theta Gram, before taking blocks. An ordinary-$L^2$ margin cannot silently become a uniform $H$ margin. The comparison in Proposition~\ref{sw:prop:weighted-comparison} grows exponentially with the budget and supplies no uniform positive weighted gap. The physical bridge~\eqref{eq:intro-bridge} does not change this norm distinction or supply a missing full-family sign.

Finally, for a finite low/high split
\begin{equation}
 Q=\begin{pmatrix}A&B\\B^*&D_H\end{pmatrix},\quad D_H>0,
\end{equation}
completing the square gives
\begin{equation}
 Q\succeq0\quad\Longleftrightarrow\quad A-BD_H^{-1}B^*\succeq0.
\end{equation}
This is the exact additional response-cost condition. In an infinite setting, the inverse, domains, uniform bounds and passage to the form limit need their own justification. They are not supplied by the existence of positive logarithmic bands.


For an infinite direct sum, the displayed matrix congruence is initially a finite-support form identity only. To extend it by a bounded normalized operator, one must construct the cross form on the completion of the diagonal-energy space and prove its boundedness there. The sparse estimate gives this control: since $Q_{jj}\succeq\mu(\log T_j)G_{jj}$ on represented blocks, each diagonal-energy-normalized cross block has norm at most $k(T_i,T_j)/\mu$, and its full row sum is at most $\mathcal E_{12}(t_0)/\mu\le1/2$. It does not by itself identify that completion with $X$, justify an unbounded inverse on an arbitrary high space, or imply convergence in the actual form norm. These are separate obligations.

The new bounds are analytic consequences of the cited exponential-sum estimates, the unconditional PNT and the full zero-counting bound. They do not certify a numerical theta matrix or an effective final PNT height. All actual zeros, complete multiplicity prefixes, prime powers and Gamma shifts remain present. Arbitrary adjacent chains, a uniform reciprocal-theta gap, and the unrestricted global half inequality remain unresolved.

\subsection{Interpretation and remaining boundary}\label{sw:sec:limits}
The two regimes have different metrics and different coupling scope. Sparse-window positivity~\eqref{sw:eq:main} allows arbitrarily many selected unit windows and uses the sum of their individual ordinary-$L^2$ block energies. The contiguous theorems~\eqref{sw:eq:actual-contiguous} and~\eqref{sw:eq:eight-band} remove internal spacing within one growing band, but use the ordinary norm of the entire represented sum. Replacing that aggregate norm by a sum of separate window norms would discard near-cancellation and is not proved.

The sparse coordinate satisfies $\Phi(T)\sim T/(\log T)^8$. Thus it permits order $T/(\log T)^8$ allowed window slots up to height $T$, rather than only order $\log T$ under a fixed multiplicative gap. These are allowed slots, not assertions of zero occupancy. Adjacent unit heights have $\Phi$ gaps tending to zero and do not meet the sparse hypothesis.

For the contiguous theorem, fixed common prefix lists define genuine vector spaces, including the full actual zero-jet spaces above. A union of separately budgeted variable-frequency classes is not automatically one admissible vector space. By Lemma~\ref{sw:lem:domain}, inequalities pass to actual $X$-limits of admissible approximants with the same stated lower coefficient; this does not admit arbitrary coefficient series, authorize unrestricted linear spans, or replace form convergence by bare $\mathcal H$ convergence. It does not combine arbitrary distinct growing bands. Theorem~\ref{bl:sparsebands} supplies the further complete cross estimate for its stated exponent-$12$ sparse geometry, retaining the sum of whole-band ordinary energies.

The earlier PNT argument reaches $L=8\log T$ with the sufficient range $A<2\pi^2/\log(4e)$. The flexible-anchor refinement extends the sufficient range to $A<2\pi^2/\log(2e)$ and yields the explicit margin in the eleven-logarithmic band. The stronger sufficient boundary itself is not covered. Neither argument proves a vanishing relative Gram error for fixed positive logarithmic width, and neither sufficient boundary is asserted optimal. Only the local compressed-head constant one is proved sharp. The PNT branches supply no effective final height here, whereas the earlier sparse and finite-constant branches retain their symbolic threshold rules.

Proposition~\ref{sw:prop:weighted-comparison} supplies only an exponentially budget-dependent reciprocal-theta norm conversion. It does not produce a uniform weighted coercivity gap as the complexity grows. Redundant coefficient lists may have singular Grams; all matrix statements are about represented functions or their quotient by the Gram kernel.

The positivity arguments in this section assert no low-frequency block, general low/high coupling, arbitrary union of contiguous bands, completeness of a sparse family, or operator-reducing property. The new sparse growing-band result is separate from the unit-window low-to-sparse response theorem in Section~\ref{sec:low-sparse}; applying that response argument to growing bands requires a further mixed estimate and is not asserted here. They use $q_0$ directly and do not depend on the centering, radical, energy-bridge or operator-core results elsewhere in the monograph. These restricted results do not establish the unrestricted global half inequality or the Riemann hypothesis. These analytic arguments require no numerical zero data, theta matrix, response inverse or evaluated threshold, and make no claim of historical priority.




\section{Vanishing coupling from a fixed low space to sparse windows}\label{sec:low-sparse}
The sparse theorem controls interactions within the high family. We now
bound its complete interaction with a fixed low space in the actual
weighted Gram. The estimate is uniform in the number of selected sparse
windows, but its low-space constant depends on the full low basis.

We retain the complete forms, full spaces $\mathscr H,\mathscr X$ and
linear-first inner product convention of Section~\ref{sec:sparse-bands}.
The final low and high functions are even. Intermediate carrier
modulations need not be even. The completeness and continuity required
below are established in Section~\ref{sw:sec:domain}; no additional
operator-core or bounded-inverse assertion is inferred.

Fix $C_0>0$. Use the constants $K,K_2$ and scales
$M_N=K(N+1)\log(N+2)$ and $b_N=K_2(N+1)^4\log(N+2)$ of
Proposition~\ref{sw:prop:analytic}. The envelope derivative and tail bounds
hold for offsets in $[-1,1]+i[-1/2,1/2]$. The full Gamma upper bound is
Lemma~\ref{sw:lem:gamma-upper}. The high packets are
\eqref{sw:eq:packet}, with a fixed frequency-prefix list in each block,
\emph{total} paired budget $N_j\le C_0\log T_j$, and spacing
\eqref{sw:eq:spacing}. Here ``total'' includes the reflected paired
representation, not merely the unpaired envelope. Theorem~\ref{sw:thm:sparse}
provides
\[
 q_0\!\left[\sum_j u_j\right]\ge\tfrac12
       \sum_j\log T_j\norm{u_j}_2^2
\]
for every finite selection. Enlarge its threshold $T_0\ge e^{18}$ so
that $T\ge2b_N$ whenever $T\ge T_0$ and $N\le C_0\log T$.
The envelope-to-packet bound is then
$\norm A_2\le2\norm{u_T}_2$. A union of varying frequency lists is not
silently treated as a fixed vector space of the same budget.

\subsection{A low carrier against a high packet}
The constants in this subsection are local to this argument. They remain symbolic in this general argument;
Section~\ref{sec:four-zero-certificate} later verifies one conservative
instance for a particular four-dimensional low space.

Set
\begin{align*}
 A_0&=C_0+1,&B_0&=C_0+2,&L_0&=1+\log B_0,\\
 B_1&=K_2A_0^4L_0,& C_X&=KA_0L_0+5+2\log3,\\
 G_0&=6+\log9,& E_0&=2G_0+|c_\Gamma|+12\sqrt{C_X},
 &C_*&=8E_0B_1+16.
\end{align*}

The complete unpaired cross estimate is already proved in
Section~\ref{sw:sec:cross}, including the unbounded Gamma tail, both
prime orientations and every prime power. In particular, for
$f=e^{i\alpha x}A$, $g=e^{i\beta x}B$ satisfying
\eqref{sw:eq:cross-hyp}, equations~\eqref{sw:eq:unpaired-cross}--\eqref{sw:eq:Lcross}
give
\begin{equation}\label{ls:eq:unpaired}
 |q_0(f,g)|\le
 \left[(2G_R+|c_\Gamma|+12\sqrt X)\frac{b_N+b_M}{|\alpha-\beta|}
       +\frac4{(R+2)^2}\right]\norm A_2\norm B_2,
\end{equation}
where $G_R=4+\log(1+8R^2)$ and
$X=\lceil\max(M_N,M_M,2\log(2+R),2)\rceil$.
This is a bound for the full form, not a finite-prime or bounded-multiplier
substitute. We now specialize it to one fixed low space.

\begin{lemma}[Low carrier against one high packet]\label{ls:lem:mixed}
Let $f=e^{i\alpha x}A$, $|\alpha|\le L$, with envelope budget at most
$N_V$. If
\[
 T\ge T_0,\quad T\ge2L,\quad N_V\le C_0\log T,\quad b_{N_V}\le T/2,
\]
and $u_T$ is an admissible high packet, then
\begin{equation}\label{ls:eq:one-cross}
 |q_0(f,u_T)|\le C_*\frac{(\log T)^6}{T}\norm A_2\norm{u_T}_2.
\end{equation}
\end{lemma}
\begin{proof}
Put $t=\log T\ge18$ and $R=T$. Each carrier separation from $\alpha$
to $\pm T$ is at least $T/2$. For both budgets at most $C_0t$,
\[
 b_N\le B_1t^5,\quad X\le C_Xt^2,\quad
 G_R\le G_0t,\quad 2G_R+|c_\Gamma|+12\sqrt X\le E_0t.
\]
Each bracket in \eqref{ls:eq:unpaired} is at most
$4E_0B_1t^6/T+4/(T+2)^2$. The two coefficients of the high packet
have total modulus one, and its envelope norm is at most twice the
packet norm. This gives $8E_0B_1t^6/T+8/(T+2)^2$, bounded by
$C_*t^6/T$. Reflection preserves the envelope bounds.
\end{proof}

\subsection{The actual low Gram and the packing sum}

Let $V$ be a fixed finite-dimensional even subspace with basis
$v_1,\ldots,v_d$. Suppose each basis vector has a fixed finite representation with real
carriers $\alpha_{\ell r}\in\mathbb R$
\begin{equation}\label{ls:eq:low-decomp}
 v_\ell=\sum_r e^{i\alpha_{\ell r}x}A_{\ell r},\qquad
 |\alpha_{\ell r}|\le L,
\end{equation}
where $A_{\ell r}=WQ_{\ell r}$ has offsets in
$[-1,1]+i[-1/2,1/2]$ and prefix budget at most $N_V$.
Thus each envelope satisfies \eqref{sw:eq:input-tail} and
\eqref{sw:eq:deriv-main}. This is a restricted class of fixed low
spaces, not an assertion for an arbitrary finite-dimensional subspace of
$\mathscr X$.
Define the true weighted Gram by
$c^*G_Vc=\norm{\sum_\ell c_\ell v_\ell}_{\mathscr H}^2$, and put
\begin{equation}\label{ls:eq:KV}
 a_\ell=\sum_r\norm{A_{\ell r}}_2,\qquad
 K_V=\sum_{\ell=1}^d a_\ell\sqrt{(G_V^{-1})_{\ell\ell}}.
\end{equation}
Then $G_V\succ0$ and
\begin{equation}\label{ls:eq:coordinate}
 \sum_\ell |c_\ell|a_\ell\le K_V\norm v_{\mathscr H},
 \qquad v=\sum_\ell c_\ell v_\ell.
\end{equation}
To see this, apply Cauchy--Schwarz in the $G_V$ metric to each coordinate:
\[
|c_\ell|\le\sqrt{(G_V^{-1})_{\ell\ell}}\sqrt{c^*G_Vc}.
\]
Thus a basis condition number is not silently replaced by one. $K_V$ is
finite for this fixed space, but need not remain bounded for growing spaces.

Choose $U\ge\max(T_0,2L)$ with $N_V\le C_0\log U$ and
$b_{N_V}\le U/2$. Let the higher windows all start at $T_j\ge U$ and
satisfy \eqref{sw:eq:spacing}. Define
\begin{equation}\label{ls:eq:theta}
 \Theta(U)=\frac{(\log U)^{11}}{U^2}
 +\frac{(\log U)^3+3(\log U)^2+6\log U+6}{U}.
\end{equation}

\begin{lemma}[Complete squared mixed-tail sum]\label{ls:lem:packing}
For any finite or countable sequence as above,
\begin{equation}\label{ls:eq:pack-sum}
 \sum_j\frac{(\log T_j)^{11}}{T_j^2}\le\Theta(U),
 \qquad \Theta(U)\sim\frac{(\log U)^3}{U}\longrightarrow0.
\end{equation}
\end{lemma}
\begin{proof}
For a nonnegative decreasing function $F$, the change of coordinate
$y=\Phi(T)$ and separation by at least one give
\[
 \sum_j F(T_j)\le F(U)+\int_U^\infty
                    \frac{F(s)}{(\log s)^8}\,ds.
\]
For instance, after the first point each value is paid by the integral on
the preceding unit interval in the $y$ coordinate; these intervals are
disjoint. Take $F(s)=(\log s)^{11}/s^2$, decreasing for $\log s>11/2$.
The remaining integral is
\[
 \int_U^\infty\frac{(\log s)^3}{s^2}\,ds
 =\frac{(\log U)^3+3(\log U)^2+6\log U+6}{U}.
\]
The asymptotic follows directly. This estimates the entire sum, not just
individual pairings.
\end{proof}

\subsection{Vanishing Schur response on the sparse high family}

\begin{theorem}\label{ls:thm:main}
Under the packet and sparse-family estimates recalled above and the stated low-space conditions, put
\[
 \epsilon_V(U)=2C_*^2K_V^2\Theta(U).
\]
For every finite higher selection, $h=\sum_j u_j$, let
$Z^2=\sum_j\log T_j\norm{u_j}_2^2$. Then
\begin{align}
 |q_0(v,h)|&\le C_*K_V\sqrt{\Theta(U)}\norm v_{\mathscr H}Z,\label{ls:eq:mixed-sum}\\
 q_0[v+h]&\ge q_0[v]-\epsilon_V(U)\norm v_{\mathscr H}^2,\label{ls:eq:lower}\\
 q_0[v+h]&\ge q_0[v]-2\epsilon_V(U)\norm v_{\mathscr H}^2
                         +\tfrac14Z^2.\label{ls:eq:retained}
\end{align}
Choose an independent basis in each fixed frequency-prefix high block.
For these coordinates, let $C_F$ be the exact combined
high form matrix, $B_F$ the exact low--high coupling matrix, and $A_V$ the
low form matrix. The convention is
$q_0[v+h]=c^*A_Vc+2\Re(d^*B_Fc)+d^*C_Fd$.
Then $C_F\succ0$ and
\begin{equation}\label{ls:eq:response}
 \boxed{0\preceq B_F^*C_F^{-1}B_F
          \preceq\epsilon_V(U)G_V,\qquad
 A_V-B_F^*C_F^{-1}B_F\succeq A_V-\epsilon_V(U)G_V.}
\end{equation}
All estimates are uniform in the number of selected higher windows and
their total complexity.
\end{theorem}
\begin{proof}
Lemma~\ref{ls:lem:mixed} and \eqref{ls:eq:coordinate} give
\[
 |q_0(v,u_j)|\le C_*K_V\frac{(\log T_j)^6}{T_j}
                         \norm v_{\mathscr H}\norm{u_j}_2.
\]
Apply Cauchy--Schwarz in $j$, putting $\sqrt{\log T_j}\norm{u_j}_2$
in the second factor. The squared first factor is exactly the sum in
\eqref{ls:eq:pack-sum}. This proves \eqref{ls:eq:mixed-sum}.
Set $k=C_*K_V\sqrt{\Theta(U)}\norm v_{\mathscr H}$ and use
$q_0[h]\ge Z^2/2$. The scalar identities
\[
 \tfrac12 Z^2-2kZ=\tfrac12(Z-2k)^2-2k^2,\qquad
 \tfrac12 Z^2-2kZ=\tfrac14Z^2+\tfrac14(Z-4k)^2-4k^2
\]
prove \eqref{ls:eq:lower} and \eqref{ls:eq:retained}.

The sparse inequality implies that no nonzero block tuple synthesizes to
zero, and that its exact matrix $C_F$ is positive definite. Moreover
\[
 |q_0(v,h)|^2\le\epsilon_V(U)\norm v_{\mathscr H}^2q_0[h].
\]
Taking the supremum of this quotient over the finite high space gives
$c^*B_F^*C_F^{-1}B_Fc\le\epsilon_V(U)c^*G_Vc$. This is the exact Schur
response, not a residual approximation; the finite rule is also recorded
in \cite[Appendix A.5]{Boyd}. The second matrix inequality follows.
\end{proof}

\begin{proposition}[Infinite variational response]\label{ls:prop:infinite}
Let $Y_U$ be the $\mathscr X$-closure of the algebraic sum of the fixed selected
high block spaces. Then $q_0[h]\ge0$ on $Y_U$ and, for $v\in V$,
\begin{equation}\label{ls:eq:variational}
 \mathfrak R_U[v]:=\sup_{h\in Y_U}
 \{2\Re q_0(v,h)-q_0[h]\}
 \le\epsilon_V(U)\norm v_{\mathscr H}^2.
\end{equation}
Along increasing finite high selections that exhaust these blocks, the
positive response matrices increase and converge in finite-dimensional
operator norm to the matrix of $\mathfrak R_U$.
\end{proposition}
\begin{proof}
Continuity in $\mathscr X$ extends both $q_0[h]\ge0$ and
$|q_0(v,h)|^2\le\epsilon_V(U)\norm v_{\mathscr H}^2q_0[h]$ to $Y_U$.
Complete the scalar square in $\sqrt{q_0[h]}$ to obtain
\eqref{ls:eq:variational}. In particular a null vector of the high form has
zero coupling with every $v\in V$. The variational suprema over increasing
finite spaces are increasing; density and continuity show that their
supremum is the one over $Y_U$. Bounded monotone Hermitian matrices on the
fixed finite-dimensional $V$ converge in norm. Their limit represents the
displayed quadratic form.
\end{proof}

No assertion that the high operator has a bounded inverse on $\mathscr H$ is
needed here. No arbitrary block-$\ell^2$ series is admitted into $\mathscr X$.
If a fixed series does converge in $\mathscr X$, then the energy inequalities
pass to its sum; its increasing $Z^2$ sums are bounded by the finite
sparse inequalities and hence converge.

\subsection{Centering, positivity, and the missing coverage}

Set $c(x)=\cosh(x/2)$, $m(u)=\int c(x)u(x)\,dx$, and
$q[u]=q_0[u]+2|m(u)|^2$. The mean is continuous on $\mathscr H$ because
$\int c^2W<\infty$. If each high block satisfies the \emph{exact}
condition $m(u_j)=0$, then $m(h)=0$ also on $Y_U$. Therefore
\begin{equation}\label{ls:eq:mean}
 q[v+h]\ge q[v]-\epsilon_V(U)\norm v_{\mathscr H}^2.
\end{equation}
The mixed mean term is identically zero under this explicit condition;
the low mean correction remains. If both spaces are centered, then
$q=q_0$ on their sum. For actual zero jets the required mean identity must
be checked for the exact permitted derivative orders. Subtracting a mean
representer changes the frequency packet, so the packet estimates cannot
simply be reapplied after that operation.

If an independently certified low bound is
$q_0[v]\ge a_V\norm v_{\mathscr H}^2$ with $a_V>0$, then
$\epsilon_V(U)\le a_V$ proves nonnegativity on $V+Y_U$.
For $\epsilon_V(U)\le a_V/4$, the stronger finite bound is
\[
 q_0[v+h]\ge\tfrac12a_V\norm v_{\mathscr H}^2+\tfrac14Z^2.
\]
A lower bound $q_0[v]\ge-b\norm v_{\mathscr H}^2$ with $b>0$ does not supply
this positive low margin. If the low form has a null vector, positivity of
the combined form requires its exact mixed coupling to vanish; a small
absolute bound alone cannot pay a nonzero coupling from a zero margin.

The restriction to sparse windows is substantive. It cannot be removed
just by using a fixed finite number of sparse classes.

\begin{proposition}[A counting obstruction to boundedly many sparse classes]\label{ls:prop:counting}
Suppose selected unit windows cover every nontrivial zeta zero with
ordinate in $[T,2T]$, and their starts are partitioned into $Q(T)$ classes,
each satisfying \eqref{sw:eq:spacing}. Then for sufficiently large $T$,
\[
 Q(T)\ge c(\log T)^8
\]
for an absolute $c>0$. This concerns coverage by windows, not a claim about
the closed linear span of their associated functions.
\end{proposition}
\begin{proof}
The zero-counting formula with an $O(\log T)$ remainder implies
$N_\zeta(2T)-N_\zeta(T)\ge aT\log T$ for large $T$, while every interval
of length one in this range contains at most $b\log(2T+1)$ zeros with
multiplicity. These consequences also follow from the explicit bound in
\cite[Corollary 1.2]{HSW}. Enlarge $b$ to include endpoints.
Windows that meet $[T,2T]$ start in $[T-1,2T]$. Each sparse class contains
at most
\[
 1+\Phi(2T)-\Phi(T-1)
 \le1+\frac{T+1}{(\log(T-1))^8}
\]
such starts. Counting zeros with possible overcounting therefore gives
\[
 aT\log T\le Q(T)b\log(2T+1)
 \left(1+\frac{T+1}{(\log(T-1))^8}\right).
\]
Rearrangement proves the assertion. Positivity within each class would
still leave all interactions between different classes to be controlled.
\end{proof}

For a sequence of growing low spaces $V_j$, the present response bound is
useful only with a joint estimate on $K_{V_j}$, $U_j$, and a genuine low
margin $a_j$. For example,
$2C_*^2K_{V_j}^2\Theta(U_j)\le a_j$ pays this sparse response.
It does not pay interactions with the unselected windows below or above
$U_j$. Taking $U_j$ extremely large for each fixed $V_j$ does not make
those windows disappear. The form-core exhaustion and all remaining
couplings must be proved on the same domain and in the same Gram metric.

Thus the two concrete operator statements obtained here are
\[
 C_F\succeq\tfrac12\operatorname{diag}_j((\log T_j)H_j),
 \qquad B_F^*C_F^{-1}B_F\preceq\epsilon_V(U)G_V,
\]
where $H_j$ is the ordinary $L^2$ Gram within block $j$.
The first uses the stated sparse coercivity input; the second is the new
mixed-response consequence. They apply to these actual blocks and all
their finite selections, not to an arbitrary full complement or a different
spectral projection. Neither turns a fixed-stage numerical global bound
into a proof at the one-half endpoint.


\paragraph{Separation from the finite response-gate spaces.}
The low space in this section is not automatically the cut-centered
cosine source space $\operatorname{span}\{C(W1_Ie_n)\}$ of
Section~\ref{vrg:sec:geometry}, a piecewise-linear trial space, or an
optimized response-image space. None of these is identified here with
the finite global theta--exponential-polynomial class
\eqref{ls:eq:low-decomp}. Thus an existing finite response gate or source
Gram cannot simply be reused as $A_V$ or $G_V$ in this theorem: an exact
identification or a proved transfer is required. Even for a matching
space, a global variance coefficient strictly below one half does not
supply the positive unshifted $q_0$ margin $a_V$ used above.

In fact, direct equality fails for every nonzero centered sharp-cut
source of this type. If $C(W1_Ip)=WQ$ for an entire exponential polynomial
$Q$, division by $W>0$ gives
$Q=1_Ip-\beta\cosh(x/2)$ for a constant $\beta$.
On an open exterior interval, $Q=-\beta\cosh(x/2)$; the identity theorem
then gives this equality everywhere. It follows that $p=0$ on $I$,
so the cut source and its centering are zero. This argument blocks direct
identification only. It does not exclude approximation or a transfer
with separately proved and paid form and Gram errors.


% Body-only inclusion. All source theorem and proof content is retained.
% Existing macros required: \Hh and \norm. X is the common Gamma form domain.
\section{A certified four zero theta subspace and sparse attachment}
\label{sec:four-zero-certificate}

This section supplies a concrete positive low space for the preceding
low-to-sparse response theorem. The finite computed result uses the actual
theta weight, full prime form, and a positive lower part of the complete
Gamma form with an explicit interpolation debit. Its sparse infinite
attachment has a separate, explicitly quantified analytical premise.

\subsection{The actual form and the certified low space}

For $r=|x|$ and $\tau_j=\pi j^2e^{2r}$, set
\[
 W(x)=\frac{\sum_{j\ge1}(4\tau_j^2-6\tau_j)e^{-\tau_j}}{1+e^{-r}},
 \qquad \Hh=L^2_{\rm even}(\mathbb R,dx/W).
\]
All unsubscripted physical integrals below use Lebesgue measure. Put
$\rho(s)=e^{-s/2}/(1-e^{-2s})$, $c_\Gamma=\log\pi-\psi(1/4)$,
$a_n=\Lambda(n)/\sqrt n$, and
\begin{align}
 t_\Gamma[u]&=\int_0^\infty\rho(s)\norm{u(\cdot+s)-u}_2^2\,ds,\\
 p[u]&=2\Re\sum_{n\ge2}a_n\int_{\mathbb R}u(x+\log n)\overline{u(x)}\,dx,\\
 q_0[u]&=t_\Gamma[u]-c_\Gamma\norm u_2^2-p[u].\label{fz:eq:q}
\end{align}
Polarization defines the corresponding Hermitian forms. The form domain is
$X=\{u\in\Hh:t_\Gamma[u]<\infty\}$ with
$\norm u_{X}^2=\norm u_{\Hh}^2+t_\Gamma[u]$.

Let $\gamma_1,\ldots,\gamma_4$ be the first four positive zero ordinates:
\[
 (14.1347251417\ldots,\ 21.0220396388\ldots,
   25.0108575801\ldots,\ 30.4248761259\ldots).
\]
They are rigorously isolated as simple critical-line zeros by the FLINT
zero routine \cite{fzFlintZeros}; no global RH assumption is used. The
certificate evaluates functions with enclosing balls for the exact
ordinates, not rounded point ordinates. More explicitly, the following outward decimal intervals enclose the exact
ordinates:
\begin{align*}
\gamma_1&\in[14.134725141734693790457251983562470270784257116
                 \mathbin{\pm}3.13\cdot10^{-46}],\\
\gamma_2&\in[21.022039638771554992628479593896902777334340525
                 \mathbin{\pm}1.33\cdot10^{-46}],\\
\gamma_3&\in[25.010857580145688763213790992562821818659549673
                 \mathbin{\pm}4.69\cdot10^{-46}],\\
\gamma_4&\in[30.424876125859513210311897530584091320181560024
                 \mathbin{\pm}3.23\cdot10^{-46}].
\end{align*}
Here $[a\mathbin{\pm}b]$ denotes $[a-b,a+b]$. Each of these four zeros
is simple, so this certified space uses their zeroth jets only. Define
\[
 v_j(x)=W(x)\cos(\gamma_jx),\qquad V=\operatorname{span}\{v_1,v_2,v_3,v_4\},
 \qquad G_{ij}=\int W(x)\cos(\gamma_ix)\cos(\gamma_jx)\,dx.
\]
Thus $c^*Gc=\norm{\sum c_jv_j}_{\Hh}^2$ is the actual function norm.

\begin{theorem}[Finite computed certificate]\label{fz:thm:finite}
The supplied enclosing computation and exact rational sign verifier give
\begin{equation}\label{fz:eq:main}
 \boxed{q_0[v]\ge\frac1{32}\norm v_{\Hh}^2\quad(v\in V).}
\end{equation}
Moreover, for
\begin{equation}\label{fz:eq:KV}
 K_V=\norm W_2\sum_{j=1}^4\sqrt{(G^{-1})_{jj}},
 \qquad K_V^2<32.
\end{equation}
The proof is the analytic error enclosure and positive sufficient matrix
constructed in the following three subsections, with the executable arithmetic supplied.
\end{theorem}

The mean is exactly zero, independently of numerical quadrature. Set
$\kappa=2\cosh(x/2)W$ and $m(u)=\int\cosh(x/2)u(x)\,dx$.
The classical theta Fourier representation \cite[(1.2)]{fzLagarias} gives
\begin{equation}\label{fz:eq:mean}
 \int\kappa(x)\cos(zx)\,dx=\Xi(z),\qquad
 m(v_j)=\tfrac12\Xi(\gamma_j)=0.
\end{equation}
Consequently $q=q_0+2|m|^2$ equals $q_0$ on $V$. No approximate
centering projection or deletion of a mean correction is involved.

\subsection{Theta, spatial, and prime enclosures}

For $r\ge0$, direct differentiation of the defining series gives
\begin{equation}\label{fz:eq:envelope}
 |W^{(k)}(r)|\le C_ke^{(4+2k)r-\pi e^{2r}},\qquad
 (C_0,C_1,C_2)=(42,700,10000),\quad k=0,1,2.
\end{equation}
Here $C_0$ in this section is a derivative-envelope constant, not a
frequency-complexity budget. To verify the constants, write
$s=(1+e^{-r})^{-1}$ and $N=\sum(4\tau^2-6\tau)e^{-\tau}$.
The first two derivatives of a summand of $N$ have polynomials
\[
 -8\tau^3+28\tau^2-12\tau,\qquad
 16\tau^4-104\tau^3+136\tau^2-24\tau.
\]
Use $|s|\le1$, $|s'|,|s''|\le1/4$, and $\tau\ge\pi$.
Sufficient constants are respectively
$4\pi^2\sum j^4e^{-\pi(j^2-1)}$,
$22\pi^3\sum j^6e^{-\pi(j^2-1)}$,
and $80\pi^4\sum j^8e^{-\pi(j^2-1)}$; the code encloses these below
the displayed integers. Smoothness at zero follows from the theta
transformation, also underlying \eqref{fz:eq:mean}.

Let $W_8$ retain $j\le8$ in the theta series. Uniformly for $r\ge0$,
\begin{equation}\label{fz:eq:theta-omission}
 |W^{(k)}(r)-W_8^{(k)}(r)|\le10^{-90},\quad k\le2.
\end{equation}
For example the omitted derivatives are bounded by
$10000\sum_{j\ge9}j^8e^{-3j^2}$. The ratio of successive terms is
at most $(10/9)^8e^{-57}<1/2$, and
$20000\,9^8e^{-243}<10^{-90}$. These deliberately loose estimates
avoid treating a finite theta sum as the exact weight.

We use physical radius $R=2$. Put $E=e^{2R}$ and
\[
 J_m(a,E)=\int_E^\infty t^me^{-at}\,dt,
 \quad J_0=e^{-aE}/a,\quad
 J_m=e^{-aE}E^m/a+(m/a)J_{m-1}.
\]
From \eqref{fz:eq:envelope},
\begin{equation}\label{fz:eq:spatial}
 \norm{W^{(k)}}_{L^2(|x|>R)}\le C_k\sqrt{J_{3+2k}(2\pi,E)},
 \qquad \int_{|x|>R}W\le42J_1(\pi,E).
\end{equation}
The integrals defining $G$, the physical Gram $H_{ij}=\int v_iv_j$,
and the norms $\norm{W^{(k)}}_2$ are evaluated on $[0,2]$ with
\eqref{fz:eq:theta-omission} and \eqref{fz:eq:spatial} added outwards.

Every prime power $n\le32$ is included. For $s=\frac12\log n$, evenness
rewrites the correlation as an integral on $y\ge0$:
\begin{align*}
 C_{ij}(\log n)=\int_0^\infty&W(y+s)W(|y-s|)\big[
 \cos(\gamma_i(y+s))\cos(\gamma_j(y-s))\\
 &+\cos(\gamma_j(y+s))\cos(\gamma_i(y-s))\big],dy.
\end{align*}
The two finite intervals $[0,s]$ and $[s,2]$ use fixed analytic branches
of $W_8$. The omitted $y$ tail is bounded by
$2\norm W_2\norm W_{L^2(|x|>2)}$; the finite theta error is also added.
This computes the full head matrix $P_{\le32}=2\sum_{n\le32}a_nC(\log n)$.
Ball integration uses meromorphic finite expressions, so the analytic
callback condition of \cite{fzFlintIntegral} is satisfied on its evaluation
disks or yields a nonfinite ball for subdivision.

For the entire remaining prime series, the weight envelope gives
\[
 |p_{>M}(u,v)|\le r_M\norm u_{\Hh}\norm v_{\Hh},\qquad
 r_M=84\sum_{n>M}\sqrt n\log n\,e^{-\pi n/2}.
\]
Indeed, at least one of $|x|,|x+\log n|$ is at least
$\frac12\log n$, giving
$\sup_x\sqrt{W(x+\log n)W(x)}\le42n e^{-\pi n/2}$.
Both orientations have been included. For $z=e^{-\pi/2}$ and $m=M+1$,
\begin{equation}\label{fz:eq:prime-tail}
 r_M\le84z^m\left[\frac{m^2}{1-z}
             +\frac{2mz}{(1-z)^2}+\frac{z(1+z)}{(1-z)^3}\right],
\end{equation}
using $\sqrt n\log n\le n^2$. At $M=32$ the certified upper bound
is less than $3.609\times10^{-18}$. In matrix order the tail costs
$r_{32}G$, not $r_{32}I$ in unrelated coordinates.

\subsection{Full Gamma energy and a certified interpolation debit}

Expand $\rho(s)=\sum_{\ell\ge0}e^{-a_\ell s}$,
$a_\ell=2\ell+1/2$. Each resulting form $t_a$ is positive. Its Fourier
multiplier is $2\xi^2/[a(a^2+\xi^2)]$, so
\begin{equation}\label{fz:eq:gamma-upper}
 0\le d_\Gamma(\xi)\le4+\tfrac12\log(1+4\xi^2).
\end{equation}
For the upper bound, the first summand is at most $4$ and the rest are
bounded by their decreasing integral in the label variable. The integral
is exactly $\frac12\log(1+4\xi^2)$.
Thus, for $g\in H^1$ with $m_0=\norm g_2^2$, $m_1=\norm{g'}_2^2$,
\begin{equation}\label{fz:eq:jensen}
 t_\Gamma[g]\le F(m_0,m_1):=
 4m_0+\frac{m_0}{2}\log(1+4m_1/m_0).
\end{equation}
This follows by Jensen's inequality in Fourier space. The right side is
increasing in both arguments; $g=0$ is handled by continuity.

On $[-2,2]$ use $8192$ uniform cells of width $h=1/2048$.
Let $f_j$ be the piecewise linear interpolant of the exact nodal values
$v_j$, with its two endpoint values replaced by zero, extended by zero.
All exact nodal values are enclosed with theta and zero errors included.
Put $e_j=v_j-f_j$. An upper bound for $\norm{v_j''}_2$ is
\[
 D_j=\norm{W''}_2+2\gamma_j\norm{W'}_2+\gamma_j^2\norm W_2.
\]
Poincare's inequality twice on each cell, followed by the endpoint ramp
correction, gives
\begin{align}
 \norm{e_j}_2&\le\frac{h^2}{\pi^2}D_j
       +\sqrt{2h/3}\,W(R)+\norm W_{L^2(|x|>R)}=:b_{0j},\\
 \norm{e_j'}_2&\le\frac h\pi D_j+\sqrt{2/h}\,W(R)
       +\norm{W'}_{L^2(|x|>R)}
       +\gamma_j\norm W_{L^2(|x|>R)}=:b_{1j}.\label{fz:eq:interp}
\end{align}
For the derivative estimate the interpolation error has mean-zero
derivative on each cell; for the function estimate it has zero endpoint
values. The endpoint correction makes $f_j\in H^1(\mathbb R)$.

Set $E_{\rm err}=\sum_jF(b_{0j}^2,b_{1j}^2)$ and let
$E_F$ bound $\sum_jt_\Gamma[f_j]$. The code uses
$F(\sum_j\norm{f_j}_2^2,\sum_j\norm{f_j'}_2^2)$ for $E_F$,
which is an upper bound by the same concavity argument. Cauchy--Schwarz
in the full Gamma energy then gives, for all coefficient vectors $c$,
\begin{equation}\label{fz:eq:debit}
 t_\Gamma\left[\sum_jc_jv_j\right]
 \ge t_\Gamma\left[\sum_jc_jf_j\right]
       -\Delta_\Gamma\norm c_2^2,
 \qquad \Delta_\Gamma=2\sqrt{E_FE_{\rm err}}.
\end{equation}
The nonnegative error energy has only been dropped in a lower bound.
The computed enclosures give
\[
 E_F<1.248327,\qquad E_{\rm err}<1.493\times10^{-9},\qquad
 \Delta_\Gamma<8.633\times10^{-5}.
\]

To enclose a positive part of $t_\Gamma[f]$, keep labels $0\le\ell<128$.
If $d_i$ is the derivative vector on cell $i$, then
\[
 t_a(f_j,f_k)=a^{-2}\iint e^{-a|x-y|}f_j'(x)f_k'(y)\,dx\,dy.
\]
With $r_a=e^{-ah}$, the same-cell and off-cell weights are
\[
 D_a=2\left(\frac h a-\frac{1-r_a}{a^2}\right),\qquad
 O_a=\frac{(1-r_a)^2}{a^2}.
\]
The exact matrix is
\begin{equation}\label{fz:eq:recurrence}
 T_a=a^{-2}\left[D_a\sum_i d_id_i^T
 +O_a\sum_{i>j}r_a^{i-j-1}(d_id_j^T+d_jd_i^T)\right].
\end{equation}
The second sum is accumulated with
$w_i=\sum_{j<i}r_a^{i-j-1}d_j$, $w_{i+1}=r_aw_i+d_i$.
This is an exact finite-element integral, with ball arithmetic enclosing
its coefficients. Let $\Gamma_{128}=\sum_{\ell<128}T_{a_\ell}$.
The entire omitted Gamma sum remains a nonnegative form, so
$t_\Gamma[f_c]\ge c^*\Gamma_{128}c$. No assertion that the full Gamma
energy equals its finite truncation is made.

\subsection{The sufficient matrix and its independent rational sign check}

Combining the preceding bounds proves that the following exact sufficient
matrix is below $q_0|_V-G/32$ in common coefficient coordinates:
\begin{equation}\label{fz:eq:certificate}
 M=\Gamma_{128}-c_\Gamma H-P_{\le32}
       -\Delta_\Gamma I-(r_{32}+1/32)G.
\end{equation}
The calculation uses 160-bit balls and $10^{-23}$ requested finite
integration tolerances; the returned enclosing radii, not a tolerance
assumption, enter every subsequent operation. For readability, the
matrix is approximately
\[
 M\approx\begin{pmatrix}
 .02977055&.01936909&.00937675&.00169132\\
 .01936909&.04462486&.04171459&.01736883\\
 .00937675&.04171459&.05025826&.03680333\\
 .00169132&.01736883&.03680333&.05998650
 \end{pmatrix}.
\]
These displayed decimals are not the sign certificate. The exact export
encloses every entry in $[L_{ij},L_{ij}+1]/10^{18}$, where
{\setlength{\arraycolsep}{4pt}
\[
L=\begin{pmatrix}
29770549382006073&19369088169309656&9376745804446878&1691319957294064\\
19369088169309656&44624860215294801&41714594770375288&17368833986541499\\
9376745804446878&41714594770375288&50258260922383360&36803330101386192\\
1691319957294064&17368833986541499&36803330101386192&59986503767699879
\end{pmatrix}.
\]
}
The same integer endpoints are included in the machine-readable ledger.

Let $M_c=(L+\frac12\boldsymbol1\boldsymbol1^T)/10^{18}$ and
$M_*=M_c-2\cdot10^{-18}I$. Every symmetric matrix in the entry boxes
is at least $M_*$ by the absolute row-sum bound on its error. The
standard-library-only verifier performs exact fraction $LDL^T$ and obtains
\begin{equation}\label{fz:eq:pivots}
 D_{11}>\frac{29}{1000},\quad D_{22}>\frac{31}{1000},\quad
 D_{33}>\frac7{1000},\quad D_{44}>\frac8{1000}.
\end{equation}
Thus $M\succeq M_*\succ0$, proving \eqref{fz:eq:main}.
These are coordinate factorization pivots, not weighted eigenvalues.

The actual weighted Gram is also certified positive. Its enclosing inverse
and the complete $\norm W_2$ bound give
\[
 K_V^2\le21.849544666284<32.
\]
Indeed $v_j=(e^{i\gamma_jx}W+e^{-i\gamma_jx}W)/2$, so its two envelope
norms sum to $\norm W_2$. For every $v=\sum c_jv_j$,
\begin{equation}\label{fz:eq:gram-coords}
 \norm W_2\sum_j|c_j|\le K_V\norm v_{\Hh},
\end{equation}
by the coordinate-functional norms $\sqrt{(G^{-1})_{jj}}$.
This completes the proof of Theorem~\ref{fz:thm:finite}.

\subsection{Attachment to infinitely many sparse higher windows}

The rest of the result has a separate analytical input. It is not needed
for the finite certificate.

\begin{hypothesis}[Sparse high-family estimates]\label{fz:hyp:sparse}
For fixed frequency-prefix lists, let
$u_T=(e^{iTx}A(x)+e^{-iTx}A(-x))/2$, $A=WQ$, with offsets in
$[0,1]+i[-1/2,1/2]$ and total paired prefix complexity $N\le8\log T$.
At a specified threshold $T_0$, assume that
\begin{align}
 q_0[\sum_j u_j]&\ge\tfrac12\sum_j\log T_j\norm{u_j}_2^2,\label{fz:eq:sparse-pos}\\
 |q_0(v_i,u_T)|&\le C_*\frac{(\log T)^6}{T}\norm W_2\norm{u_T}_2\label{fz:eq:mixed}
\end{align}
whenever $T_j\ge T_0$ and successive starts satisfy
$\Phi(T_{j+1})-\Phi(T_j)\ge1$, where
$\Phi'(T)=(\log T)^{-8}$. The mixed estimate is for the complete form
\eqref{fz:eq:q}, with every prime power and the unbounded Gamma multiplier.
Assume the coefficient-free theta envelope and derivative estimates admit
the constant recipe below. Prefix degree $k$ costs $k+1$; each block has
one fixed list, not an unrestricted union of varying lists. This application
uses the coefficient $8$ in the total paired budget. It does not assert
that every full actual unit-window zero-jet space satisfies that smaller
budget; actual permitted zero jets must separately meet the stated count.
\end{hypothesis}

The numerical threshold below is conditional on the stated analytical
estimates with the displayed constant recipe and its complete sufficient
threshold conditions. Equivalently, one may take
\eqref{fz:eq:sparse-pos}--\eqref{fz:eq:mixed} themselves as assumptions for
every admissible family above $U$, with the displayed $C_*$. Mere existence
of unspecified finite constants and an unspecified threshold would not
justify the numerical value of $U$.

We verify a conservative explicit choice for that recipe. Set
\begin{align*}
 c_W&=(4\pi^2-6\pi)/2,&m_0&=c_We^{-\pi e^2},&L_W&=700/c_W,\\
 D_0&=10\,18^{3/2},&B&=2\cdot42D_0e/m_0,&K&=2^{16},\\
 K_1&=12\sqrt K+4e^{-1}/(1-e^{-1})^2,&&
 K_2=K(D_0e^{L_We^2}+L_W+2).
\end{align*}
The chosen $B$ dominates $1,4e,2\cdot42e^{1/2}/m_0,2\cdot700e^{1/2}/m_0$.
The absorption test is
\[
 \frac1K\left[\frac{\log B}{2\log3}+\frac72+
 \frac{\log(2B)}{\log3}+\frac{15}2
       \left(1+\frac{\log K}{2\log3}\right)\right]\le\frac\pi2.
\]
Its enclosing value is less than $0.001442$. This recipe uses the
complex-exponent derivative tools represented by \cite{Erdelyi}; the
arithmetic check does not independently prove the sparse analytical input.

Put $A_0=9$, $B_0=10$, $L_0=1+\log10$,
$B_1=K_2A_0^4L_0$, $C_X=KA_0L_0+5+2\log3$,
$G_0=6+\log9$, $E_0=2G_0+|c_\Gamma|+12\sqrt{C_X}$,
$C_*=8E_0B_1+16$. The scalar verifier gives
\[
 K_1<2^{12},\quad B_1<2^{800},\quad C_*<2^{800},\quad |c_\Gamma|<6.
\]
For $t_0=2^{512}$ and $T_0=U=e^{t_0}$, it verifies the whole-ray tests
\begin{align*}
 16B_1^2t_0^{10}e^{-2t_0}&\le1/8,\\
 K_1\sqrt{9\log(10t_0)/t_0}+|c_\Gamma|/t_0&\le1/8,\\
 C_*\left[\frac{2^{17}+4/3}{t_0^2}+\frac{2^{19/2}}{t_0^3}
       +(2+2^{11/2})t_0^5e^{-t_0}\right]&\le1/4.
\end{align*}
The low-carrier compatibility conditions are also checked:
$\gamma_4<31$, the envelope $W$ has budget $N_V=1$,
$b_1=16K_2\log3$, and
\[
 t_0>\log62,\qquad
 t_0>\log(32K_2\log3),\qquad 1<8t_0.
\]
Thus $T\ge U$ implies $T\ge2\gamma_4$, $b_1\le T/2$, and the required
low-envelope budget condition. These are needed for the low-to-high mixed
estimate, in addition to the high-family tests.

All left sides decrease for $t_0\ge18$. The first two are the diagonal
packet tests and the last pays the complete sparse row sum. Computation
is in logarithmic coordinates; $\exp(2^{512})$ is never constructed as a
floating number. This is an effective but entirely impractical height.

For $U\ge e^{18}$ define
\[
 \Theta(U)=\frac{(\log U)^{11}}{U^2}+
 \frac{(\log U)^3+3(\log U)^2+6\log U+6}{U}.
\]
For sparse starts $T_j\ge U$, the packing inequality gives
\[
 \sum_j\frac{(\log T_j)^{11}}{T_j^2}\le\Theta(U).
\]
It follows by summing a decreasing function in the separated coordinate
$\Phi$, then evaluating $\int_U^\infty (\log s)^3s^{-2}ds$.
From \eqref{fz:eq:mixed} and \eqref{fz:eq:gram-coords}, Cauchy--Schwarz gives
\[
 |q_0(v,h)|\le C_*K_V\sqrt{\Theta(U)}\norm v_{\Hh}Z,
 \quad h=\sum_j u_j,\quad Z^2=\sum_j\log T_j\norm{u_j}_2^2.
\]
The exact Schur response therefore satisfies
\begin{equation}\label{fz:eq:response}
 B_F^*C_F^{-1}B_F\preceq\epsilon_V(U)G,
 \qquad \epsilon_V(U)=2C_*^2K_V^2\Theta(U),
\end{equation}
for every finite high selection. This follows either by completing the
square or by the supremum of $|q_0(v,h)|^2/q_0[h]$.

At the specified $U$, $\epsilon_V(U)\le1/128$. For a simple explicit
check, $\Theta(e^t)\le17t^{11}e^{-t}$ for $t\ge1$, so
\[
 \epsilon_V(e^{t_0})<2^{1611}t_0^{11}e^{-t_0}
 =\exp(7243\log2-2^{512})<1/128.
\]

\begin{theorem}[Conditional infinite sparse attachment]\label{fz:thm:attachment}
Under Hypothesis~\ref{fz:hyp:sparse}, with the certified four-zero space $V$
and the indicated threshold, the exact Schur matrix for every finite
admissible high selection satisfies
\[
 A_V-B_F^*C_F^{-1}B_F\succeq\frac3{128}G.
\]
Every finite admissible high sum also obeys
\begin{equation}\label{fz:eq:attached}
 \boxed{q_0[v+h]\ge\frac1{64}\norm v_{\Hh}^2
                  +\frac14\sum_j\log T_j\norm{u_j}_2^2.}
\end{equation}
The conclusion passes to fixed series convergent in $X$, and
nonnegativity holds on the $X$-closure of $V$ plus the algebraic sum of
the fixed high block spaces. If those high blocks are actual permitted
zero jets, the mean correction is retained exactly.
\end{theorem}
\begin{proof}
Write $k=C_*K_V\sqrt{\Theta(U)}\norm v_{\Hh}$. The sparse input and
the finite certificate give
\[
 q_0[v+h]\ge\tfrac1{32}\norm v_{\Hh}^2+\tfrac12 Z^2-2kZ
 \ge(\tfrac1{32}-2\epsilon_V(U))\norm v_{\Hh}^2+\tfrac14Z^2,
\]
which proves \eqref{fz:eq:attached}. Continuity of the Gamma form in $X$
and the bounded prime form prove the closure claim. A bounded high inverse
in $\Hh$ or convergence of arbitrary abstract block sequences is not assumed.
Finally, for an actual zero $w$ of order $\nu$, differentiation of
\eqref{fz:eq:mean} gives
$m(W\partial_z^k\cos(zx)|_{z=w})=\Xi^{(k)}(w)/2=0$ for $k<\nu$.
Thus $q=q_0$ throughout the centered sum and its form closure.
\end{proof}

\subsection{Scope and reproduction}

The finite result concerns four exact functions and all their complex
linear combinations. It is not a floating eigenvalue sample or a local
radius certificate. The sparse extension is conditional on its stated
analytical estimates, whose constants have now been instantiated.
It omits the other low and intermediate zero windows and the unselected
high windows. Neither a uniform bound for growing $V$ nor a full-domain
form-core exhaustion is proved. A $1/32$ margin for $q_0$ on this subspace
cannot be added to an unrelated global energy coefficient.

The bundle contains the full regeneration program, finite-integral
checkpoint, complete enclosing matrices, scalar-constant verifier,
integer interval matrix, and a standard-library rational verifier.
With Python and \texttt{python-flint==0.9.0}, run in order:
\begin{verbatim}
python3 certify_four_zeros.py
python3 check_sparse_constants.py
python3 export_rational_matrix.py
python3 check_rational_matrix.py
\end{verbatim}
The last command checks only the exported matrix sign; it does not
regenerate its analytic enclosures. The complete run reported here uses
trusted FLINT/Arb ball arithmetic and rigorous quadrature
\cite{fzArbIntegration,fzFlintIntegral}. It is a reproducible computation
with an accompanying proof of its error bounds, not a formal proof-assistant
verification or an independent second-implementation audit. A separate
source review and full reexecution with Python-FLINT 0.9.0 and FLINT 3.6.0
reproduced the complete finite-integral checkpoint, rational matrix,
rational sign certificate, and sparse-scalar ledger. A 256-bit rerun gave
the same rational matrix. A live-ball serialization check verified all
sixteen entries against the exported rational boxes; it retained both
endpoints and the full printed ball radius.


% Full six-section source body retained, with namespaced labels and audit additions.
% Uses the existing \norm and \Hh macros and the existing fz bibliography keys.
\section{A derivative Gram bound and a five zero theta certificate}
\label{sec:five-zero-certificate}

The first five actual zero functions admit a stronger finite-space lower
bound than the scalar interpolation estimate supplies at the same target.
The improvement comes from preserving the second-derivative cross terms
inside the Gamma interpolation error. The result below concerns this one
five-dimensional space. It does not upgrade the preceding conditional
sparse attachment to a new low space without rechecking its constants.

\subsection{The form, normalization, and result}
For $r=|x|$, $\tau_j=\pi j^2e^{2r}$, define
\[
 W(x)=\frac{\sum_{j\ge1}(4\tau_j^2-6\tau_j)e^{-\tau_j}}{1+e^{-r}},
 \qquad \Hh=L^2_{\rm even}(\mathbb R,dx/W).
\]
Let $\rho(s)=e^{-s/2}/(1-e^{-2s})$, $c_\Gamma=\log\pi-\psi(1/4)$,
and $a_n=\Lambda(n)/\sqrt n$. Physical norms and integrals use Lebesgue
measure. The full forms are
\begin{align}
 t_\Gamma[u]&=\int_0^\infty\rho(s)\norm{u(\cdot+s)-u}_2^2\,ds,\\
 p[u]&=2\Re\sum_{n\ge2}a_n\int_{\mathbb R}u(x+\log n)\overline{u(x)}\,dx,\\
 q_0[u]&=t_\Gamma[u]-c_\Gamma\norm u_2^2-p[u].\label{fz5:eq:q}
\end{align}
Set $X=\{u\in\Hh:t_\Gamma[u]<\infty\}$, with
$\norm u_X^2=\norm u_\Hh^2+t_\Gamma[u]$.
The functions used here are smooth and belong to $X\cap H^2(\mathbb R)$.

Let $\gamma_1,\ldots,\gamma_5$ be the first five positive zero ordinates:
\[
14.13472514\ldots,\quad21.02203964\ldots,\quad25.01085758\ldots,
\quad30.42487613\ldots,\quad32.93506159\ldots.
\]
Their exact values are enclosed by the rigorous FLINT zero routine
\cite{fzFlintZeros}. The verification for these indices does not assume
RH at arbitrary height. Put
\[
 v_j=W\cos(\gamma_jx),\qquad V_5=\operatorname{span}_{\mathbb C}\{v_1,\ldots,v_5\},
 \qquad G_{ij}=\int W\cos(\gamma_ix)\cos(\gamma_jx).
\]
Thus $c^*Gc=\norm{\sum c_jv_j}_\Hh^2$ is the actual function norm.

\begin{theorem}[Computed finite-space certificate]\label{fz5:thm:main}
The supplied enclosing computation and rational sign verification give
\begin{equation}\label{fz5:eq:main}
 \boxed{q_0[v]\ge\frac1{25}\norm v_\Hh^2\qquad(v\in V_5).}
\end{equation}
The exact Schur complement of $q_0-\norm{\cdot}_\Hh^2/25$ with respect to
$V_4=\operatorname{span}\{v_1,\ldots,v_4\}$ is greater than $1/2000$
in the coefficient of $v_5$.
\end{theorem}

The mean is handled analytically. With $\kappa=2\cosh(x/2)W$ and
$m(u)=\int\cosh(x/2)u(x)\,dx$, the classical theta representation
\cite[(1.2)]{fzLagarias} gives
\[
 \int\kappa(x)\cos(zx)\,dx=\Xi(z),\qquad
 m(v_j)=\Xi(\gamma_j)/2=0.
\]
Consequently $q=q_0+2|m|^2$ equals $q_0$ on $V_5$. There is no numerical
centering projection or deletion of a nonzero mean term.

\subsection{An interpolation inequality retaining derivative correlations}
The expansion $\rho(s)=\sum_{\ell\ge0}e^{-(2\ell+1/2)s}$ gives positive
forms $t_a$ with multiplier $2\xi^2/[a(a^2+\xi^2)]$. Bounding the first
term by $4$ and the rest by the decreasing integral in the label gives
\[
0\le d_\Gamma(\xi)\le4+\tfrac12\log(1+4\xi^2).
\]
For $g\in H^1$ this implies, by Jensen's inequality,
\begin{equation}\label{fz5:eq:F}
 t_\Gamma[g]\le F(\norm g_2^2,\norm{g'}_2^2),\qquad
 F(a,b)=4a+\tfrac a2\log(1+4b/a).
\end{equation}
The right side is increasing in both arguments and homogeneous of degree
one. The zero case is defined by continuity.

Take any finite family $v_1,\ldots,v_N\in H^2(\mathbb R)$, and a uniform
mesh of width $h$ on $[-R,R]$. Let $f_j$ be its linear interpolant with
both endpoint values set to zero, extended by zero outside the interval.
For $v_c=\sum c_jv_j$, write $v_c-f_c=a_c+r_c$: inside the interval,
$a_c$ is the interpolation error using the \emph{original} endpoint
values, and it is zero outside. The remainder $r_c$ is the sum of the
endpoint ramps inside and $v_c$ outside. Both pieces are in $H^1$;
the traces of $r_c$ match at the two boundaries.

Let $D_{ij}=\int\overline{v_i''}v_j''$, so that
$c^*Dc=\norm{v_c''}_2^2$.
The two cellwise Poincare inequalities yield
\[
 \norm{a_c}_2\le\frac{h^2}{\pi^2}\sqrt{c^*Dc},\qquad
 \norm{a_c'}_2\le\frac h\pi\sqrt{c^*Dc}.
\]
Define
\begin{equation}\label{fz5:eq:chi}
 \chi(h)=\frac{h^4}{\pi^4}
          \left[4+\frac12\log(1+4\pi^2/h^2)\right].
\end{equation}
Homogeneity in \eqref{fz5:eq:F} gives $t_\Gamma[a_c]\le\chi(h)c^*Dc$.
Choose any $\tau\ge\sum_jt_\Gamma[r_j]$. Then
$t_\Gamma[r_c]\le\tau\norm c_2^2$.

\begin{theorem}[Matrix interpolation bound]\label{fz5:thm:interp}
For $\eta>0$, set
\[
 E=(1+\eta)\chi(h)D+(1+\eta^{-1})\tau I.
\]
Then $t_\Gamma[v_c-f_c]\le c^*Ec$. If $\Gamma_L$ is the matrix of any
positive finite label sum of $t_\Gamma$ on the interpolants $f_j$, then,
for every $0<\varepsilon<1$,
\begin{equation}\label{fz5:eq:matrix-error}
 \boxed{t_\Gamma[v_c]\ge
 c^*\left[(1-\varepsilon)\Gamma_L
                -(\varepsilon^{-1}-1)E\right]c.}
\end{equation}
The numerical constants in this statement do not depend on $N$.
\end{theorem}
\begin{proof}
The first statement is the Hilbert seminorm inequality
$\norm{a+r}_\Gamma^2\le(1+\eta)\norm a_\Gamma^2+
(1+\eta^{-1})\norm r_\Gamma^2$.
Apply Young's inequality to $v_c=f_c+(v_c-f_c)$ to obtain
\[
 t_\Gamma[v_c]\ge(1-\varepsilon)t_\Gamma[f_c]
             -(\varepsilon^{-1}-1)t_\Gamma[v_c-f_c].
\]
Since $1-\varepsilon>0$, the complete positive label sum may be
replaced by $\Gamma_L$ in this lower bound.
\end{proof}

The distinction from a scalar debit matters: small values of
$\norm{v_c''}_2$ caused by cancellation are retained in $c^*Dc$.
The theorem does \emph{not} assert $D\preceq C G$ with a constant uniform
in source dimension. For example, if $D\preceq L_NG$ and
$G\succeq g_NI$, the error cost is bounded by
\[
 E\preceq\left[(1+\eta)\chi(h)L_N+
                   (1+\eta^{-1})\tau/g_N\right]G.
\]
Growth of $L_N$, loss of $g_N$, and the tail trace $\tau$ must still be
controlled when $N$ increases. Refining the mesh improves an approximation
bound; it does not decide the sign of the underlying exact form.

\subsection{Complete enclosures for the five theta functions}
All computations use 160-bit Arb balls through Python-FLINT 0.9.0
\cite{fzFlintIntegral,fzArbIntegration}. The returned quadrature radii enter
the ledger; a requested tolerance is not treated as a proven error bound.
The parameters are
\[
 R=2,\quad h=1/2048,\quad 8192\text{ cells},\quad L=128,
 \quad\varepsilon=1/16384,\quad\eta=1/1024.
\]

For $k=0,1,2$, differentiation of the theta series gives
\begin{equation}\label{fz5:eq:env}
 |W^{(k)}(r)|\le C_ke^{(4+2k)r-\pi e^{2r}},\qquad
 (C_0,C_1,C_2)=(42,700,10000).
\end{equation}
Indeed the numerator derivative polynomials are
$-8t^3+28t^2-12t$ and $16t^4-104t^3+136t^2-24t$.
Use $s=(1+e^{-r})^{-1}$, $|s|\le1$, $|s'|,|s''|\le1/4$, and $t\ge\pi$.
The sufficient constants are
$4\pi^2\sum j^4e^{-\pi(j^2-1)}$,
$22\pi^3\sum j^6e^{-\pi(j^2-1)}$, and
$80\pi^4\sum j^8e^{-\pi(j^2-1)}$, checked below the displayed integers.
Retaining $j\le8$ costs at most $10^{-90}$ in each derivative through
order two, since the omitted terms are bounded by
$10000\sum_{j\ge9}j^8e^{-3j^2}$ and
$20000\,9^8e^{-243}<10^{-90}$.

For $E_R=e^{2R}$ and $J_m(a,E_R)=\int_{E_R}^\infty t^me^{-at}\,dt$,
\[
 \norm{W^{(k)}}_{L^2(|x|>R)}\le C_k\sqrt{J_{3+2k}(2\pi,E_R)},
 \qquad \int_{|x|>R}W\le42J_1(\pi,E_R).
\]
The $J_m$ are evaluated by their elementary integration-by-parts recurrence.
The true physical Gram $H_{ij}=\int v_iv_j$, weighted Gram $G$, prime
head, and $D$ use these spatial tails and the theta omission bound.
For the new derivative matrix we integrate
\[
 v_j''=W''\cos(\gamma_jx)-2\gamma_jW'\sin(\gamma_jx)
                         -\gamma_j^2W\cos(\gamma_jx).
\]
If $d_j$ bounds $\norm{v_j''}_2$, the finite-interval approximation error
has norm at most $e_j=2\cdot10^{-90}(1+\gamma_j)^2$.
An outward entry error is $e_id_j+e_jd_i+e_ie_j$, plus the product of
the two derivative-tail norm bounds. This encloses every cross term.

For the endpoint remainder in Theorem~\ref{fz5:thm:interp}, valid bounds are
\begin{align*}
 b_{0j}&=\sqrt{2h/3}\,W(R)+\norm W_{L^2(|x|>R)},\\
 b_{1j}&=\sqrt{2/h}\,W(R)+\norm{W'}_{L^2(|x|>R)}
                           +\gamma_j\norm W_{L^2(|x|>R)}.
\end{align*}
We use $\tau=\sum_j F(b_{0j}^2,b_{1j}^2)$. The certified upper values are
\[
 \chi(h)<7.857\times10^{-15},\qquad \tau<5.114\times10^{-140}.
\]

The exact piecewise linear Gamma integrals have a linear recurrence.
For one label $a=2\ell+1/2$, let $d_i$ be the vector of slopes on cell $i$,
$r_a=e^{-ah}$, $D_a=2[h/a-(1-r_a)/a^2]$, $O_a=(1-r_a)^2/a^2$.
Then
\[
 T_a=a^{-2}\left[D_a\sum_i d_id_i^T+
 O_a\sum_{i>j}r_a^{i-j-1}(d_id_j^T+d_jd_i^T)\right].
\]
The second sum uses $w_{i+1}=r_aw_i+d_i$. Summing $0\le\ell<128$
produces $\Gamma_L$; all omitted labels remain nonnegative.

All prime powers $n\le32$ are integrated, with both orientations.
Writing $s=\frac12\log n$, the even correlation integrates
$W(y+s)W(|y-s|)$ times the two swapped cosine products on $y\ge0$.
We split at $s$ and $R=2$, using fixed analytic branches. Its $y>R$
error is at most $2\norm W_2\norm W_{|x|>R}$.
The whole remaining prime operator satisfies
\begin{align}
 |p_{>M}(u,v)|&\le r_M\norm u_\Hh\norm v_\Hh,\\
 r_M&\le84z^{M+1}\left[\frac{(M+1)^2}{1-z}
       +\frac{2(M+1)z}{(1-z)^2}+\frac{z(1+z)}{(1-z)^3}\right],
 \quad z=e^{-\pi/2}.\label{fz5:eq:prime}
\end{align}
For the proof, \eqref{fz5:eq:env} implies
$\sup_x\sqrt{W(x+\log n)W(x)}\le42n e^{-\pi n/2}$;
use $\Lambda(n)\le\log n$ and $\sqrt n\log n\le n^2$ and sum the
complete integer tail. Here $r_{32}<3.609\times10^{-18}$, charged as
$r_{32}G$, not as a matrix in an unrelated norm.

\subsection{The Schur certificate and the scalar comparison}
The preceding proof gives the common-coordinate lower matrix
\begin{equation}\label{fz5:eq:Qlower}
 \underline Q=(1-\varepsilon)\Gamma_L-(\varepsilon^{-1}-1)E
                -c_\Gamma H-P_{\le32}-r_{32}G,
 \qquad q_0|_{V_5}\succeq\underline Q.
\end{equation}
Set $\alpha=1/25$ and partition
\[
 M=\underline Q-\alpha G=\begin{pmatrix}A&b\\b^*&c\end{pmatrix},
 \qquad \underline s=c-b^*A^{-1}b.
\]
The following decimals are descriptive. The proofs use the enclosing
balls and rational interval ledger supplied with the paper.
\begin{center}
\begin{tabular}{lr}\toprule
Quantity & Enclosed value, approximately\\\midrule
$c$ & $0.0622228365822704$\\
$b^*A^{-1}b$ & $0.0616945688507312$\\
$\underline s$ & $0.000528267731539247$\\
$\norm{v_5-P_{V_4}^{\Hh}v_5}_\Hh^2$ & $0.0194662457002191$\\
$\underline s/\norm{v_5-P_{V_4}^{\Hh}v_5}_\Hh^2$ & $0.027137627854625$\\
\bottomrule\end{tabular}
\end{center}

Every entry of $M$ is rounded outwards to an integer interval on the
$10^{-18}$ grid. Let $M_c$ be the midpoint matrix and let $\rho$ be the
largest row sum of interval radii. Then every enclosed symmetric matrix
is at least $M_c-\rho I$. A separate standard-library rational
$LDL^T$ calculation on this latter matrix gives successive pivots larger
than
\[
 \frac{28}{1000},\quad\frac{31}{1000},\quad\frac7{1000},
 \quad\frac7{1000},\quad\frac1{2000}.
\]
In particular $M\succ0$, proving \eqref{fz5:eq:main}. These are coordinate
pivots, not weighted eigenvalues.

To transfer the Schur bound to the actual full form, observe that for
fixed fifth coefficient $z$,
\begin{align*}
 \inf_{x\in\mathbb C^4}(q_0-\alpha\norm{\cdot}_\Hh^2)
          \left[\sum_{j\le4}x_jv_j+zv_5\right]
 &\ge \inf_x (x,z)^*M(x,z)\\
 &=\underline s\,|z|^2>\frac{|z|^2}{2000}.
\end{align*}
Both low blocks are strictly positive, so these infima are their exact
Schur complements. The monotonicity argument retains the actual full
Gamma and prime forms; it does not identify their response with the
response of the sufficient matrix.

Orthogonalizing the new function in the \emph{actual} $\Hh$ norm makes
the same calculation more interpretable. Put $r=v_5-P_{V_4}^{\Hh}v_5$.
The sufficient shifted diagonal and response become approximately
$0.00323356827784684$ and $0.00270530054630759$. Their difference is
unchanged, because this basis change is triangular. Consequently the
actual shifted Schur remainder is at least
\begin{equation}\label{fz5:eq:intrinsic}
 \frac{27}{1000}\norm r_\Hh^2\,|z|^2.
\end{equation}
This is a reserve in the new quotient direction, not a global eigenvalue
bound to be added to $1/25$.

For comparison, put
\begin{align*}
 E_F&=F(\sum_j\norm{f_j}_2^2,\sum_j\norm{f_j'}_2^2),\\
 E_e&=\sum_jF((h^2d_j/\pi^2+b_{0j})^2,(hd_j/\pi+b_{1j})^2),
 \qquad \Delta=2\sqrt{E_FE_e}.
\end{align*}
Cauchy--Schwarz in the full Gamma energy gives
$t_\Gamma[v_c]\ge t_\Gamma[f_c]-\Delta\norm c_2^2$.
The enclosing code uses outward upper bounds in these expressions and
obtains $\Delta<0.000127358$. The resulting scalar sufficient matrix is
\[
 M_{\rm scalar}=\Gamma_L-c_\Gamma H-P_{\le32}-r_{32}G
                        -\Delta I-\tfrac1{25}G.
\]
Its leading four-dimensional block is positive, but its Schur remainder
is enclosed near $-0.0009495689556652$. Thus this particular sufficient
matrix is indefinite. It is not a negative direction for $q_0-\norm{\cdot}_\Hh^2/25$;
Theorem~\ref{fz5:thm:main} proves that exact form positive on the same space.
The successful comparison preserves the second-derivative correlations
instead of charging the worst coefficient loss to every direction.

\subsection{What an unlimited extension would require}
The following exact linear algebra states the continuing obligation.
It is included to separate an approximation improvement from a sign proof.

\begin{lemma}[One-step extension]
On a finite direct sum, let the matrix of $q_0-\alpha\norm{\cdot}_\Hh^2$
be $\left(\begin{smallmatrix}A&B\\B^*&C\end{smallmatrix}\right)$ with
$A\succ0$. It is nonnegative if and only if
$C-B^*A^{-1}B\succeq0$.
If $A$ is only nonnegative, the necessary and sufficient conditions are
$\operatorname{ran}B\subseteq\operatorname{ran}A$ and
$C-B^*A^\dagger B\succeq0$.
\end{lemma}
\begin{proof}
For $A\succ0$, complete the square in $x+A^{-1}By$.
For a singular $A$, a nonzero component of $By$ in $\ker A$ produces
an arbitrarily signed mixed term along that kernel, so the range
condition is necessary. On its orthogonal complement the same square
completion uses $A^\dagger$ and proves sufficiency and necessity.
\end{proof}

For nested finite spaces exhausting a form core, actual nonnegativity
at every stage would pass to the closure by form continuity. Positive
Schur remainders can tend to zero; a fixed positive reserve like $1/25$
is not required to prove nonnegativity of $q_0$. However, neither the
first successful extension nor a dimension-independent interpolation
formula proves the signs of all subsequent Schur remainders.

For theta, the unresolved tasks include every missing lower and
intermediate direction, admissible higher zero jets, complete tail
couplings, and the precise full-domain decomposition/core needed by a
global one-half bridge. Positivity of a new diagonal alone is insufficient.
The actual new block must pay its inverse-response term in common
coordinates, with any kernel compatibility preserved. No all-dimension
estimate establishing those facts is supplied here.

\subsection{Reproduction and trust boundary}
The incremental generator reuses the hash-fixed four-function integral
block and computes the new row, column, and full five-function errors.
The derivative-Gram computation and final response test are new.
The bundle also includes a generator for the frozen input in a separate
subdirectory. A normal incremental reproduction is
\begin{verbatim}
python3 certify_five_zeros.py
python3 certify_matrix_error.py
python3 export_rational_matrix.py
python3 check_rational_matrix.py
\end{verbatim}
It requires \texttt{python-flint==0.9.0}; the last verifier uses only
Python's standard library. The latter checks the exported matrix sign,
not the analytic quadrature or its source bounds. The computation relies
on FLINT/Arb and the explicit error arguments above. It is not a
proof-assistant formalization, peer review, or a second independent
implementation of the complete certificate.


\subsection{Enclosure interpretation and independent arithmetic checks}

The derivative Gram in this argument is the exact Hermitian matrix
$D_{ij}=\langle v_i'',v_j''\rangle_{L^2}$, not a matrix formed by taking
entrywise upper endpoints. An interval enclosure of $D$ is not, by itself,
a single Loewner upper matrix. If $D_c$ is its midpoint and $R_D$ its
symmetric matrix of radii, then
\[
 D\preceq D_c+\delta_D I,\qquad
 \delta_D=\max_i\sum_j(R_D)_{ij},
\]
would give one explicit operator upper bound. The supplied calculation
instead propagates all enclosing balls through the affine sufficient-matrix
formula and proves positivity for every real symmetric matrix in the
resulting final box. This certifies the particular matrix constructed from
the exact $D$, $G$, $H$, prime head and interpolant Gamma forms, with the
outward scalar upper bounds used in the debit. No entrywise ordering is
substituted for operator ordering.

For completeness, put $d=10^{18}$. The final enclosure is
$M_{ij}\in[L_{ij},L_{ij}+1]/d$, where $L$ is symmetric and its upper
triangle is the following integer table.
\begin{center}
\begin{tabular}{rrr}\toprule
$i$&$j$&$L_{ij}$\\\midrule
1&1&28760435990307785\\
1&2&19011187197637525\\
1&3&9323663727615164\\
1&4&1692331546476943\\
1&5&450954585499171\\
2&2&43612876726907555\\
2&3&40952089771902003\\
2&4&17243682628951712\\
2&5&8610340444149185\\
3&3&49244789402245086\\
3&4&36245526582564298\\
3&5&26052393710586973\\
4&4&58970078116844913\\
4&5&57484313167688249\\
5&5&62222836582270409\\
\bottomrule\end{tabular}
\end{center}
Thus $M\succeq B$, where
\[
 B=\frac{L+\tfrac12\mathbf1\mathbf1^T-\tfrac52I}{10^{18}}.
\]
An exact rational determinant computation, independent of the supplied
$LDL^T$ implementation, confirms all five positive leading principal
minors of $B$. Its successive determinant ratios exceed the five pivot
floors stated above, and its last ratio satisfies
\[
 \frac{\det B}{\det B_{1:4,1:4}}>\frac{528}{10^6}>\frac1{2000}.
\]
Applying the same operator enclosure method to the actual weighted Gram
gives an upper bound
$\norm{v_5-P_{V_4}^{\Hh}v_5}_\Hh^2<19467/10^6$.
Consequently the normalized lower reserve is greater than $27/1000$,
using exact rational comparisons rather than the printed decimals.
The scalar comparison has a positive leading block but a Schur upper
bound smaller than $-949/10^6$, also by exact rational arithmetic.

The five-function envelope constant is computed separately:
\[
 K_{V_5}=\norm W_2\sum_{j=1}^5\sqrt{(G^{-1})_{jj}},\qquad
 K_{V_5}^2<75.
\]
The enclosing numerical upper value is approximately $74.271700130409$.
This follows from the coordinate-functional inequality
$|c_j|\le\sqrt{(G^{-1})_{jj}}\norm{\sum c_i v_i}_\Hh$ and the freshly
computed five-dimensional Gram. In particular, the earlier bound
$K_{V_4}^2<32$ is not assigned to $V_5$.

A complete source replay regenerated the frozen four-function input and
reproduced its mathematical fields exactly; only elapsed-time metadata
changed. The five-function checkpoint, rational matrix and rational
sign certificate were reproduced byte for byte. Both stages also passed
at 256 bits while retaining the original certified four-function block,
with the same final rational matrix. The earlier section records a
separate complete 256-bit regeneration of that four-function block.
A live-object serialization check verifies all 225 matrix entries across
the two five-function stages and all 25 final rational boxes. These checks
use Python-FLINT 0.9.0 and FLINT 3.6.0. They confirm reproduction and
outward serialization; they do not constitute a second implementation of
the entire analytic quadrature or a proof-assistant formalization.

% Graph-Schur source integration: only headings, labels, common macros and citation keys mapped.
\section{Graph bounds for a projected finite source and complete high sector}
\label{gs:sec:transfer}
The finite certificate from Section~\ref{sec:five-zero-certificate} can be
transferred to its Gamma spectral projection only after controlling the
complete operator row. This section proves that transfer under its
explicit admitted-sector and row hypotheses. Its low space remains the
projection of the five-dimensional source, not the entire low spectral
sector of Section~\ref{gb:sec:bridge}. In this section $L$ denotes the
physical representative $T_\Gamma-D_b$, and $K$ denotes the chosen closed
sector. These local symbols do not change the earlier form domains.

\subsection{Framework and precise inputs}
Let $t$ be a nonnegative densely defined closed form on a Hilbert space
$\Hh$, with representative $T$, and let $D_b=D_b^*$ be bounded.
Write $L=T-D_b$ and $q=t-\langle D_b\,\cdot,\cdot\rangle$.
Let $K$ be a closed subspace for which $X_K=\operatorname{Dom}(t)\cap K$
is dense in $K$. The restricted form is closed; let $\Gamma_K$ represent
it, put $D_K=P_KD_b|_K$, and set
\[
 L_K=\Gamma_K-D_K,\qquad\operatorname{Dom}L_K=\operatorname{Dom}\Gamma_K.
\]
These constructions use the closed-form representation and spectral
calculus framework; see \cite[Sections 2.3, 3.1]{gbTeschl}.
No reduction of $K$ by $T$ or $D_b$ separately is assumed.
For $v\in\operatorname{Dom}T\cap K$, form testing gives
\begin{equation}\label{gs:eq:compression}
 v\in\operatorname{Dom}L_K,\qquad L_Kv=P_KLv,\qquad
 \norm{L_Kv}\le\norm{Lv}.
\end{equation}
If the actual $L$ reduces $K$, this agrees with its operator restriction.
Reduction and a physical bridge are additional requirements when drawing
a global statement outside this sector.

For $R>d\ge\norm{D_b}$ use the \emph{Gamma spectral} projections
\[
 P=\mathbf1_{[0,R)}(\Gamma_K),\quad Q=I-P,
 \quad G_+=\Gamma_K|_{QK},\quad g=R-d.
\]
The exact block forms/operators of $L_K$ are
\[
 A=P\Gamma_KP-PD_KP,\qquad B=-QD_KP,\qquad
 C=G_+-QD_KQ\succeq gI.
\]
Here $A$ and $B$ are bounded, $\operatorname{Dom}C=\operatorname{Dom}G_+$,
and $S=A-B^*C^{-1}B$ is bounded on $PK$.
The low sector $PK$ is not assumed finite dimensional.

\subsection{The graph-to-Schur transfer identity}
\begin{theorem}\label{gs:thm:transfer}
For every $v\in\operatorname{Dom}L_K$, write $\ell=Pv$, $h=Qv$.
Then the complete high residual is $r=QL_Kv=B\ell+Ch$, and
\begin{equation}\label{gs:eq:identity}
 \boxed{S[\ell]=q[v]-\langle QL_Kv,C^{-1}QL_Kv\rangle.}
\end{equation}
If a linear space $V\subseteq\operatorname{Dom}L_K$ satisfies
\[
 q[v]\ge\alpha\norm v^2,\qquad\norm{L_Kv}^2\le\beta\norm v^2,
 \qquad v\in V,
\]
then
\begin{equation}\label{gs:eq:transfer}
 S[Pv]\ge\left(\alpha-\frac\beta{R-d}\right)\norm v^2.
\end{equation}
If the parenthesis is positive, $P|_V$ is injective and the same constant
is a lower bound relative to $\norm{Pv}^2$ on $PV$.
\end{theorem}
\begin{proof}
Spectral projections preserve $\operatorname{Dom}\Gamma_K$, so
$h\in\operatorname{Dom}C$. Expanding the full form gives
\[
 q[\ell+h]=\langle A\ell,\ell\rangle
 +2\Re\langle B\ell,h\rangle+\langle Ch,h\rangle.
\]
Since $r=B\ell+Ch$ and $C^{-1}r=h+C^{-1}B\ell$, subtraction gives
\eqref{gs:eq:identity}. The bound $C^{-1}\preceq g^{-1}I$ proves
\eqref{gs:eq:transfer}. If $Pv=0$ and the coefficient is positive, the
left side is zero and $v=0$. Finally $\norm{Pv}\le\norm v$.
\end{proof}

This identity subtracts the entire inverse-response energy of the trial
high component $Qv$. It is not an estimate of a compressed finite row.
It explains why positivity on $V$ alone cannot be copied to $PV$.
For example, $\Gamma=\operatorname{diag}(1,4)$ and
$D_b=\left(\begin{smallmatrix}2&1\\1&0\end{smallmatrix}\right)$, with
cutoff $R=3$, give $A=-1$, $B=-1$, $C=4$.
For $v=(1,1)$, $q[v]=1$ but $S[Pv]=-5/4$; indeed $QLv=3$ and
$1-3^2/4=-5/4$. This synthetic example is not a theta counterexample.

\subsection{Complete theta graph bounds from the derivative Gram}
For $r=|x|$, $\tau_j=\pi j^2e^{2r}$, let
\[
 W(x)=\frac{\sum_{j\ge1}(4\tau_j^2-6\tau_j)e^{-\tau_j}}{1+e^{-r}},
 \qquad\Hh=L^2_{\rm even}(\mathbb R,dx/W).
\]
Use $\rho(s)=e^{-s/2}/(1-e^{-2s})$ and the physical Gamma form
\[
 t_\Gamma[u]=\int_0^\infty\rho(s)\norm{u(\cdot+s)-u}_2^2\,ds.
\]
The full correction is
\[
 \langle D_bu,u\rangle_\Hh=c_\Gamma\norm u_2^2+
 2\Re\sum_{n\ge2}\frac{\Lambda(n)}{\sqrt n}
 \int u(x+\log n)\overline{u(x)}\,dx,
 \quad c_\Gamma=\log\pi-\psi(1/4).
\]
Thus $q_0=t_\Gamma-\langle D_b\,\cdot,\cdot\rangle_\Hh$.

\begin{hypothesis}[Admitted sector and finite input]\label{gs:hyp:inputs}
The chosen closed zero-jet sector $K$ has densely defined restricted
Gamma form. The five exact functions
\[
 v_j=W\cos(\gamma_jx),\quad 1\le j\le5,\qquad V_5=\operatorname{span}\{v_j\},
\]
where $\gamma_j$ are the first five positive zero ordinates, belong to
$K\cap\operatorname{Dom}T_\Gamma$. For these functions the actual row is
$T_\Gamma v=W d_\Gamma(D)v$, where $d_\Gamma$ is the ordinary Fourier
multiplier of $t_\Gamma$. The supplied finite certificate establishes
\[
 q_0[v]\ge\frac1{25}\norm v_\Hh^2\quad(v\in V_5),
\]
and supplies enclosures of the positive weighted Gram $G_5$ and the
physical derivative Gram $(D_2)_{ij}=\int\overline{v_i''}v_j''$.
\end{hypothesis}

These are explicit inputs, not conclusions from form or graph density
alone. The admitted-row identity follows by form testing when its
weighted row is admitted; membership in $\Hh\cap H^2$ makes the norm
estimates below finite. No $H^2$ regularity of \emph{projected} jets is
assumed. The projection in Theorem~\ref{gs:thm:transfer} is applied only after
bounding the full row of the original smooth functions.

\subsubsection*{The new finite derivative comparison}
From the frozen enclosing $G_5,D_2$ data, ball arithmetic gives
\begin{equation}\label{gs:eq:D2}
 D_2\preceq2^{20}G_5,\qquad
 \norm{v''}_2\le2^{10}\norm v_\Hh\quad(v\in V_5).
\end{equation}
We verify the stronger positive definiteness of $2^{20}G_5-D_2$.
Its entries are exported as outward rational intervals; subtracting the
maximal interval-radius row sum from the midpoint diagonal and applying
exact rational $LDL^T$ proves the comparison. The trace
$\operatorname{tr}(G_5^{-1}D_2)$ is about $690204.364$, but this decimal
is not used as the sign certificate. No theta integral is rerun here.

\subsubsection*{Bounds including all prime and Gamma rows}
The theta envelope
\[
 W(r)\le42e^{4r-\pi e^{2r}}
\]
gives $M=\norm W_\infty\le42e^{-\pi}<2$, since its exponent decreases
for $r\ge0$. One derivation of $42$ is
$4\pi^2\sum_{j\ge1}j^4e^{-\pi(j^2-1)}<42$; the supplied theta ledger
bounds the first eight terms and the remaining series.
In particular $\norm v_2\le\sqrt2\norm v_\Hh$.
The same envelope gives
$\sup_x\sqrt{W(x+\log n)W(x)}\le42n e^{-\pi n/2}$.
With $z=e^{-\pi/2}$ and $\Lambda(n)\le\log n$, the complete prime
operator has norm at most
\[
 p_*:=84\sum_{n\ge2}n^2z^n
 =84z^2\left[\frac4{1-z}+\frac{4z}{(1-z)^2}
                        +\frac{z(1+z)}{(1-z)^3}\right].
\]
All prime powers and both orientations are covered. Therefore
\begin{equation}\label{gs:eq:d}
 \norm{D_b}\le2c_\Gamma+p_*<35.719135<\frac{143}{4}<36.
\end{equation}

Expanding $\rho$ in exponential labels and bounding the decreasing sum
by its integral gives
\[
0\le d_\Gamma(\xi)\le4+\tfrac12\log(1+4\xi^2).
\]
The tangent to the concave logarithm at $k^2$, with $k=40$, gives
\[
 d_\Gamma(\xi)\le a+b\xi^2,\quad
 a=4+\tfrac12\log(1+4k^2)-\frac{2k^2}{1+4k^2},\quad
 b=\frac2{1+4k^2}.
\]
Plancherel and the triangle inequality, in physical $L^2$, now yield
\begin{align}
 \norm{T_\Gamma v}_\Hh
 &\le\sqrt M\norm{d_\Gamma(D)v}_2
 \le\sqrt2\bigl(a\norm v_2+b\norm{v''}_2\bigr)\notag\\
 &\le\sqrt2\bigl(a\sqrt2+b\,2^{10}\bigr)\norm v_\Hh
 \le\frac{65}{4}\norm v_\Hh.\label{gs:eq:graph}
\end{align}
The checked scalar value before the last rounding is less than
$16.216844$. Combining \eqref{gs:eq:d}, \eqref{gs:eq:graph}, and
\eqref{gs:eq:compression} proves the full graph bound
\begin{equation}\label{gs:eq:L}
 \boxed{\norm{L_Kv}_\Hh\le\norm{Lv}_\Hh
                   \le52\norm v_\Hh\qquad(v\in V_5).}
\end{equation}
In matrix terms the complete row Gram is at most $2704G_5$.
It is not replaced by a Gram of finitely projected rows.

\subsection{Projected low positivity and the entire high sector}
Take
\[
 R=2^{18}=262144,\qquad d=36,\qquad g=262108,
 \qquad\mathcal L_5=P_RV_5.
\]
Here $R$ is a spectral value of $\Gamma_K$, not a zeta-zero height.
No spacing or sparse-window assumption is imposed on $Q_RK$.

\begin{theorem}[Finite projected source with complete high attachment]
Under Hypothesis~\ref{gs:hyp:inputs}, $\dim\mathcal L_5=5$ and
\begin{equation}\label{gs:eq:schur}
 S_R[\ell]\ge\frac1{35}\norm\ell_\Hh^2
 \qquad(\ell\in\mathcal L_5).
\end{equation}
On the form domain
$\mathcal Y_R=\mathcal L_5\oplus Q_RX_K$ one has
\begin{equation}\label{gs:eq:attach}
 \boxed{q_0[u]\ge\frac1{36}\norm u_\Hh^2\qquad(u\in\mathcal Y_R).}
\end{equation}
\end{theorem}
\begin{proof}
Theorem~\ref{gs:thm:transfer} and \eqref{gs:eq:L} give
\[
 S_R[P_Rv]\ge\left(\frac1{25}-\frac{2704}{262108}\right)\norm v_\Hh^2
 =\frac{48627}{1638175}\norm v_\Hh^2\ge
 \frac1{35}\norm v_\Hh^2.
\]
This proves injectivity and \eqref{gs:eq:schur}. For arbitrary
$\ell\in\mathcal L_5$ and $h\in Q_RX_K$, complete the exact square:
\[
 q_0[\ell+h]=S_R[\ell]+\norm{C^{1/2}(h+C^{-1}B\ell)}^2.
\]
Put $y=h+C^{-1}B\ell$. Since $\norm{C^{-1}B}\le d/g$,
\[
 \norm{\ell+h}_\Hh^2\le[1+2(d/g)^2]\norm\ell_\Hh^2+2\norm y_\Hh^2.
\]
Meanwhile $q_0[\ell+h]\ge\norm\ell_\Hh^2/35+g\norm y_\Hh^2$.
The exact rational comparisons
\[
 \frac{1/35}{1+2(36/262108)^2}>\frac1{36},\qquad
 \frac{262108}{2}>\frac1{36}
\]
prove \eqref{gs:eq:attach} on the whole indicated form domain.
\end{proof}

\subsubsection*{The two relative operator inequalities, with restricted low source}
Let $E=-G_+^{-1/2}QD_KQG_+^{-1/2}$. On the \emph{entire} high space,
\begin{equation}\label{gs:eq:E}
 E\succeq-\kappa I,\qquad\kappa=\frac{36}{262144}=\frac9{65536}<1.
\end{equation}
Let $A_5$ denote the form $A$ restricted to $\mathcal L_5$. From
$A\succeq S_R$ and \eqref{gs:eq:schur}, $A_5\succeq I/35$ there.
For $\ell\in\mathcal L_5$, the \emph{complete} normalized coupling obeys
\begin{align}
 \langle B\ell,G_+^{-1}B\ell\rangle
 &\le\frac{36^2}{262144}\norm\ell_\Hh^2\notag\\
 &\le\frac{2835}{16384}A_5[\ell]
 \le(1-\kappa)A_5[\ell].\label{gs:eq:relative}
\end{align}
For any orthogonal decomposition reducing $G_+$, its left side equals
$\sum_j\langle B_j\ell,G_j^{-1}B_j\ell\rangle$, as an increasing
positive form sum. Thus all such high blocks are included together.
The relative estimate is proved only on $\mathcal L_5$, not on all of
$P_RK$, and not for an unrelated zero-height decomposition.
It follows from the newly established Schur lower bound; no unknown
low positivity was assumed to derive it.

\subsection{What this transfers, and what remains open}
Physical $L^2$ derivative data are converted to the reciprocal-theta norm
by \eqref{gs:eq:D2} and $M<2$. The full high residual $QL_Kv$ is bounded;
positivity is not simply copied through a spectral projection.

The operator projections and $C^{-1}$ were not constructed numerically.
The proof applies to the exact projections defined by spectral calculus.
It neither classifies $P_RK$ as finite dimensional nor certifies
$S_R$ outside the five-dimensional space $\mathcal L_5$.
The estimate covers every high Gamma vector while leaving the other low
directions untouched.

A global endpoint still requires signs on a family exhausting the needed
form core, or the entire low Schur operator at a fixed split. Raising
$R$ with source dimension changes both projections and Schur operators;
nesting and completeness require proof. Graph density gives neither an
effective derivative bound nor a low sign. Our finite derivative estimate
does not imply $H^2$ regularity of all projected jets. The actual reducing
decomposition and physical bridge are further inputs to any global
one-half equivalence. The remaining signs, that endpoint, and RH are
not proved here.

\subsection{Verification and reproduction}
The supplied numerical inputs are frozen copies of already enclosed
five-function Gram data. Their hashes are checked before use. The new
steps are the $D_2\preceq2^{20}G_5$ comparison, complete scalar operator
bounds, and the rational transfer inequalities. No prime or Gamma
integral, spectral projection, or response inverse is recomputed.
Ball arithmetic uses Python-FLINT/Arb \cite{fzArbIntegration}; a separate
standard-library rational verifier checks the exported derivative matrix
and all final rational inequalities. Run
\begin{verbatim}
python3 verify_graph_bound.py
python3 check_rational_transfer.py
\end{verbatim}
The first command requires \texttt{python-flint==0.9.0}. The second
also checks the synthetic projection example; that example is not
theta evidence. The finite positivity input and analytic admission/
sector hypotheses remain explicit. This is not a full independent
reproof of those inputs, proof-assistant formalization, or peer review.
No originality claim is made for the block identity or spectral calculus.


% Fixed-split six-zero source: full body retained after declared notation/label/heading maps.
\section{A sixth low direction at a fixed Gamma split}
\label{sx:sec:six}
The preceding five-mode result can be enlarged while keeping the same
Gamma spectral projection. The new calculation retains the physical and
second-derivative Gram correlations in the complete row cost. It gives
an incremental six-mode certificate, under the explicit sector and row
admission below, rather than a sign on the entire low spectral space.
The frozen five-function integral blocks come from
Section~\ref{sec:five-zero-certificate}; their current error debits are
recomputed after adding the sixth row. The notation $\mathcal L_m=PV_m$
is local to this section and denotes a projected finite source space.

\subsection{Normalization and scope}
For $r=|x|$ and $\tau_j=\pi j^2e^{2r}$, put
\[
 W(x)=\frac{\sum_{j\ge1}(4\tau_j^2-6\tau_j)e^{-\tau_j}}{1+e^{-r}},
 \qquad\Hh=L^2_{\mathrm{even}}(\mathbb R,dx/W).
\]
Physical norms $\norm\cdot_2$ use Lebesgue measure. Set
$\rho(s)=e^{-s/2}/(1-e^{-2s})$ and $c_\Gamma=\log\pi-\psi(1/4)$.
The complete forms are
\begin{align}
 t_\Gamma[u]&=\int_0^\infty\rho(s)\norm{u(\cdot+s)-u}_2^2\,ds,\label{sx:eq:t}\\
 p[u]&=2\Re\sum_{n\ge2}\frac{\Lambda(n)}{\sqrt n}
          \int_{\mathbb R}u(x+\log n)\overline{u(x)}\,dx,\label{sx:eq:p}\\
 q_0[u]&=t_\Gamma[u]-c_\Gamma\norm u_2^2-p[u].\label{sx:eq:q}
\end{align}
Write $X=\operatorname{Dom}(t_\Gamma)$ in $\Hh$, let $T_\Gamma$
represent the closed Gamma form, and let $D_b$ be the bounded correction
in \eqref{sx:eq:q}, so $L=T_\Gamma-D_b$.

Let $\gamma_j$ denote the first six positive ordinates of zeta zeros and
$v_j=W\cos(\gamma_jx)$, $V_m=\operatorname{span}\{v_1,\ldots,v_m\}$.
The added ordinate is
\[
 \gamma_6=37.586178158825671257217763480705332821405597351\ldots.
\]
Exact zeros, enclosed by rigorous intervals, are used throughout; this
decimal is not substituted for an exact zero. FLINT verifies the
critical-line location and simplicity for the indices used here
\cite{fzFlintZeros}. No assumption about zeros at arbitrary heights is made.

For $v_c=\sum_{j=1}^6c_jv_j$, the three Gram matrices are
\[
 c^*Gc=\norm{v_c}_\Hh^2,\qquad
 c^*Hc=\norm{v_c}_2^2,\qquad
 c^*D_2c=\norm{v_c''}_2^2.
\]
They refer to different norms; the conversions below are proved explicitly.
With $\kappa=2\cosh(x/2)W$ and
$m(u)=\int\cosh(x/2)u(x)\,dx$, the theta Fourier representation
\cite[(1.2)]{fzLagarias} gives
\[
 \int\kappa(x)\cos(zx)\,dx=\Xi(z),\qquad m(v_j)=\Xi(\gamma_j)/2=0.
\]
Thus the mean-corrected form $q=q_0+2|m|^2$ equals $q_0$ on $V_6$.

For the concrete sector
$K=\Hh_c\ominus\overline{\operatorname{span}\{r_\ell:\ell\ge1\}}$,
where $\Hh_c=\ker m$ and
$r_\ell=\kappa^{(2\ell)}-4^{-\ell}\kappa$, membership of these raw
functions follows exactly. Integration by parts gives
\[
 \langle v_j,r_\ell\rangle_\Hh
 =\bigl[(-1)^\ell\gamma_j^{2\ell}-4^{-\ell}\bigr]\Xi(\gamma_j)=0.
\]
All boundary terms vanish by theta decay. This verifies the finite
membership; it is not a proof of density of the entire zero-jet family.

\begin{hypothesis}[Sector and row admission]\label{sx:hyp:sector}
The chosen closed zero-jet sector $K\subset\Hh$ contains these six exact
functions, and $X_K=X\cap K$ is dense in $K$. The restricted Gamma
form has representative $\Gamma_K$. For the raw functions used here,
the admitted operator row is $T_\Gamma v=W d_\Gamma(D)v$, where
$d_\Gamma$ is the physical Fourier multiplier of \eqref{sx:eq:t}.
\end{hypothesis}
The raw functions are in $\Hh\cap H^2$. Since $W$ is bounded and
$d_\Gamma$ grows at most logarithmically, the displayed row has finite
$\Hh$ norm; form testing gives its operator admission. No regularity of
spectrally projected functions is inferred from this observation.
We use the closed-form representation and spectral calculus described
in \cite[Sections 2.3 and 3.1]{gbTeschl}.

Put $D_K=P_KD_b|_K$ and $L_K=\Gamma_K-D_K$. Form testing on $X_K$
gives $L_Kv=P_KLv$ for raw $v\in\operatorname{Dom}T_\Gamma\cap K$.
In particular $\norm{L_Kv}\le\norm{Lv}$. Separate invariance of $K$
under $T_\Gamma$ or $D_b$ is not needed. A global physical bridge and
the appropriate full reducing decomposition are additional inputs to
any conclusion outside the present sector.

\subsection{One fixed split and the exact response identity}
Throughout this section, keep
\[
 R=2^{18}=262144,\qquad d=36,\qquad g=R-d=262108,
\]
and use $P=\mathbf1_{[0,R)}(\Gamma_K)$ and $Q=I-P$.
This $R$ is a Gamma spectral value, not a zero height or a spatial radius.
The blocks of $L_K$ are
\[
 A=P\Gamma_KP-PD_KP,\quad B=-QD_KP,\quad
 C=\Gamma_K|_{QK}-QD_KQ\succeq gI.
\]
The complete high block $C$ is not truncated. The Schur operator
$S=A-B^*C^{-1}B$ is bounded on $PK$; we do not assume $PK$ is finite
dimensional. Define $\mathcal L_m=PV_m$. Fixed $P$ implies
$\mathcal L_5\subseteq\mathcal L_6$.

\begin{lemma}[Complete graph response]\label{sx:lem:graph}
For $v\in\operatorname{Dom}L_K$, one has
\begin{equation}\label{sx:eq:identity}
 S[Pv]=q_0[v]-\langle QL_Kv,C^{-1}QL_Kv\rangle.
\end{equation}
If $Q_m^{\rm lo}$ is a lower matrix for $q_0$ on a finite source
$v_c$, and $H_m^{\rm row}$ is an upper matrix for
$\norm{L_Kv_c}^2$, then
\begin{equation}\label{sx:eq:matrix-transfer}
 S[Pv_c]\ge c^*\bigl(Q_m^{\rm lo}-g^{-1}H_m^{\rm row}\bigr)c.
\end{equation}
\end{lemma}
\begin{proof}
Write $v=\ell+h$ with $\ell=Pv$, $h=Qv$. The complete high residual is
$r=QL_Kv=B\ell+Ch$. Expanding $q_0[\ell+h]$ and subtracting
$\langle r,C^{-1}r\rangle$ leaves precisely $S[\ell]$.
Spectral projections preserve the operator domain, so these expressions
are admitted. Now $C^{-1}\preceq g^{-1}I$ and
$\norm{QL_Kv}\le\norm{L_Kv}$ prove \eqref{sx:eq:matrix-transfer}.
\end{proof}
The estimate bounds every high response row. It is not a bound on a
finite compression of $QL_K$.

\subsection{A matrix cost retaining derivative correlations}
The theta envelope $W(r)\le42e^{4r-\pi e^{2r}}$ gives
$M=\norm W_\infty\le42e^{-\pi}<2$.
It also bounds a weighted shift by
$\sup_x\sqrt{W(x+\log n)W(x)}\le42ne^{-\pi n/2}$.
Using $z=e^{-\pi/2}$, $\Lambda(n)\le\log n$, and
$\sqrt n\log n\le n^2$ for $n\ge2$, all prime powers and both
orientations are bounded by
\[
 p_* =84z^2\left[\frac4{1-z}+\frac{4z}{(1-z)^2}
                      +\frac{z(1+z)}{(1-z)^3}\right].
\]
Consequently $\norm{D_b}\le2c_\Gamma+p_*<35.719135<36$.
The positive exponential-label expansion of $\rho$ implies
\[
 0\le d_\Gamma(\xi)\le4+\tfrac12\log(1+4\xi^2)\le a+b\xi^2,
\]
where the last bound is the tangent at $k=40$:
\[
 a=4+\tfrac12\log(1+4k^2)-\frac{2k^2}{1+4k^2},\qquad
 b=\frac2{1+4k^2}.
\]

\begin{proposition}[Dimension-independent formula, matrix-dependent cost]
For any admitted finite family of raw functions and any $\theta,\zeta>0$,
the complete row Gram is bounded by
\begin{equation}\label{sx:eq:general-row}
 H_m^{\rm row}=(1+\theta)M\left[(1+\zeta)a^2H_m
                  +(1+\zeta^{-1})b^2D_{2,m}\right]
                +(1+\theta^{-1})d^2G_m.
\end{equation}
In particular, $M=2$, $d=36$, $\theta=8$ and $\zeta=1/16$ give
\begin{equation}\label{sx:eq:row}
 \boxed{H_m^{\rm row}=\frac{153}{8}a^2H_m+306b^2D_{2,m}+1458G_m.}
\end{equation}
\end{proposition}
\begin{proof}
The weighted row identity and Plancherel give
\[
 \norm{T_\Gamma v}_\Hh^2
 \le M\norm{d_\Gamma(D)v}_2^2
 \le M\left[(1+\zeta)a^2\norm v_2^2
                +(1+\zeta^{-1})b^2\norm{v''}_2^2\right].
\]
Here the second inequality applies Young's inequality pointwise to
$(a+b\xi^2)^2$. Next,
\[
 \norm{Lv}_\Hh^2\le(1+\theta)\norm{T_\Gamma v}_\Hh^2
                   +(1+\theta^{-1})d^2\norm v_\Hh^2.
\]
Finally $L_Kv=P_KLv$ is a contraction of the full row. This proves the
matrix inequality for every coefficient vector simultaneously.
\end{proof}
No arbitrary addition of constants from different norms is made.
The actual cross terms in $H_m$ and $D_{2,m}$ remain. The coefficients
in \eqref{sx:eq:row} do not grow with $m$, but this does not bound
$D_{2,m}$ relative to $G_m$ uniformly in $m$.

\subsection{The six-function enclosure and fixed-split certificate}
The numerical input consists of frozen enclosing five-function integral
blocks with checked SHA-256 hashes. The added row and column of $G,H$,
the prime head, the Gamma interpolant matrix, and $D_2$ are calculated
at 160-bit precision using Arb ball arithmetic \cite{fzArbIntegration}.
The returned integration radii, not just requested tolerances, enter
the enclosures. Details of the infinite terms are in Section~\ref{sx:app:tails}.

Let $Q_6^{\rm lo}$ be the resulting matrix lower bound for the full
$q_0$. Exact rational signs of outward interval matrices establish
\begin{align}
 Q_6^{\rm lo}-G/40&\succ0,\label{sx:eq:finite}\\
 2^{20}G-D_2&\succ0,\label{sx:eq:derivative}\\
 M_6:=Q_6^{\rm lo}-H_6^{\rm row}/262108-G/50&\succ0.\label{sx:eq:M6}
\end{align}
The last comparison is the relevant new projected-source certificate.
Partition $M_6$ into its first five coefficient directions and its last
direction, writing $M_6=\left(\begin{smallmatrix}F&b_6\\b_6^*&c_6\end{smallmatrix}\right)$.
The enclosing values are
\begin{align*}
 c_6&=0.06539682658459459\ldots,\\
 b_6^*F^{-1}b_6&=0.065169161193679\ldots,\\
 c_6-b_6^*F^{-1}b_6&=0.000227665390916\ldots>\frac1{5000}.
\end{align*}
These are the diagonal, response cost, and remainder of the
\emph{sufficient matrix}, not a numerical construction of $C^{-1}$.
The last lower bound also bounds the corresponding exact shifted
projected-form Schur remainder, by minimizing the Loewner inequality.
Precisely, if $\widetilde S$ is the pullback of $S$ through $c\mapsto Pv_c$,
then $\widetilde S-G/50\succeq M_6$. The projected Gram is at most
$G$, so the same statement holds when shifting by that projected Gram.

The actual Gram innovation of the added raw function is
\[
 G_{66}-G_{6,1:5}G_{1:5,1:5}^{-1}G_{1:5,6}
 =0.04273614502025719\ldots>0.
\]
All four matrices in the rational ledger are handled by a midpoint
matrix minus its maximum interval-radius row sum times the identity,
followed by exact rational $LDL^T$. For $M_6$, the six strict pivot
floors are respectively
\[
 0.030,\quad0.032,\quad0.008,\quad0.010,\quad0.0014,\quad0.0002.
\]
They are rational thresholds in the checker, not floating eigenvalues.

\begin{theorem}[Six low directions with all high directions]\label{sx:thm:main}
Under Hypothesis~\ref{sx:hyp:sector} and the stated enclosing inputs,
$\dim\mathcal L_6=6$, $\mathcal L_5\subset\mathcal L_6$, and
\begin{equation}\label{sx:eq:main-low}
 \boxed{S[\ell]\ge\frac1{50}\norm\ell_\Hh^2\qquad(\ell\in\mathcal L_6).}
\end{equation}
On the full attached form domain $\mathcal Y_6=\mathcal L_6\oplus QX_K$,
\begin{equation}\label{sx:eq:main-full}
 \boxed{q_0[u]\ge\frac1{51}\norm u_\Hh^2\qquad(u\in\mathcal Y_6).}
\end{equation}
\end{theorem}
\begin{proof}
Equations \eqref{sx:eq:matrix-transfer} and \eqref{sx:eq:M6} give
$S[Pv]\ge\norm v_\Hh^2/50$. If $Pv=0$, then $v=0$; hence $P|_{V_6}$
is injective. Contractivity of $P$ proves \eqref{sx:eq:main-low}.
For $\ell\in\mathcal L_6$, $h\in QX_K$, put $y=h+C^{-1}B\ell$.
The exact square completion yields
\[
 q_0[\ell+h]=S[\ell]+\norm{C^{1/2}y}^2
             \ge\norm\ell^2/50+g\norm y^2.
\]
Since $\norm{C^{-1}B}\le d/g$ and $\ell\perp h$,
\[
 \norm{\ell+h}^2\le[1+2(d/g)^2]\norm\ell^2+2\norm y^2.
\]
The rational checks
\[
 \frac{1/50}{1+2(36/262108)^2}
 =\frac{4293787729}{214689394550}>\frac1{51},
 \qquad \frac{262108}{2}>\frac1{51}
\]
prove \eqref{sx:eq:main-full}.
\end{proof}

For comparison, \eqref{sx:eq:derivative} and the tangent bound also yield
$\norm{L_Kv}\le\norm{Lv}\le52\norm v$ on $V_6$. The simpler scalar
transfer would give only
\[
 S[Pv]\ge\left(\frac1{40}-\frac{2704}{262108}\right)\norm v^2
 =\frac{38487}{2621080}\norm v^2>\frac1{70}\norm v^2.
\]
The improved $1/50$ uses the matrix cost and is not obtained by adding
these two lower bounds. A stronger source comparison
$Q_6^{\rm lo}-G/25$ has a negative last Schur value; this only means
that this sufficient comparison does not certify $1/25$ on $V_6$.
It does not prove a negative direction of $q_0$.

\subsubsection*{Complete relative coupling on the enlarged source}
Set $G_+=\Gamma_K|_{QK}$ and
$E=-G_+^{-1/2}QD_KQG_+^{-1/2}$. On the whole high space,
\[
 E\succeq-\kappa I,\qquad\kappa=\frac9{65536}<1.
\]
On $\mathcal L_6$, the restriction $A_6$ satisfies $A_6\succeq I/50$, since
$A\succeq S$. Therefore the complete normalized coupling obeys
\[
 \langle B\ell,G_+^{-1}B\ell\rangle
 \le\frac{36^2}{262144}\norm\ell^2
 \le\frac{2025}{8192}A_6[\ell]
 \le(1-\kappa)A_6[\ell].
\]
For $\ell\ne0$, the last inequality is strict. For any orthogonal decomposition reducing $G_+$, the left side is the
sum of all corresponding block costs. This is not a claim about an
unrelated decomposition by zero height, or about all of $PK$.

\subsection{The exact next requirement for exhaustion}
The certificate just established is one finite extension. It does not
inductively determine the signs of further extensions.

\begin{proposition}[Semidefinite extension criterion]\label{sx:prop:extension}
A finite Hermitian block matrix
$\left(\begin{smallmatrix}F&b\\b^*&c\end{smallmatrix}\right)$ is nonnegative
if and only if
\[
 F\succeq0,\qquad b\in\Ran F,\qquad c-b^*F^\dagger b\ge0.
\]
Equivalently, its coupling annihilates $\ker F$, and the remaining
diagonal pays the exact inverse cost on $\Ran F$.
\end{proposition}
\begin{proof}
Testing vectors $(tz,1)$ with $z\in\ker F$ and varying complex $t$
forces $b^*z=0$. On $\Ran F$, complete the square with $F^\dagger b$.
This gives both necessity and sufficiency.
\end{proof}
In an enclosing calculation $F$ must be the leading block of the
\emph{current} lower matrix. When tail or interpolation costs change,
one cannot silently use a previously larger lower matrix in its place.
The present calculation reuses the integral blocks, but recomputes
their current error debits. This matters even though the boundary debit
here is extremely small.

\begin{theorem}[Fixed-split cofinal sign transfer]\label{sx:thm:cofinal}
Keep $R>d$ fixed. Let $V_m\subset\operatorname{Dom}L_K$ be nested
finite source spaces with $\overline{\bigcup_mPV_m}=PK$ in Hilbert norm.
Suppose, on every $V_m$, complete enclosing matrices give
\[
 S[Pv_c]\ge c^*F_mc,\qquad
 F_m\succeq-\varepsilon_mG_m,\qquad\varepsilon_m\downarrow0,
\]
where $G_m$ is the actual raw-source Gram. Then $S\succeq0$ on $PK$
and $q_0\ge0$ on $X_K$.
\end{theorem}
\begin{proof}
Fix $v\in V_k$. For every $m\ge k$ the inequality holds for the same
function, giving $S[Pv]\ge-\varepsilon_m\norm v^2$. Let $m\to\infty$.
Since $S$ is bounded, its nonnegativity extends from the dense union
of projected spaces to all of $PK$. Finally, for any
$u=\ell+h\in X_K$, use the complete square from
Theorem~\ref{sx:thm:main} and $C\succeq gI$.
\end{proof}

This proof does not need a uniform strictly positive gap. It does need
every stage (or a nested cofinal subsequence), the actual Gram metric,
complete response costs, and the indicated density. Density of an
appropriate full zero-jet form core would imply the required projected
density because $P$ is bounded. The six real simple-zero functions do
not establish such a core. A continuation must include all required
exact zero jets, including nonreal frequencies or multiplicity jets
where required; restricting forever to a convenient real-frequency
list is not automatically exhaustive.

The usable certificate template is
\begin{equation}\label{sx:eq:target}
 \boxed{Q_m^{\rm lo}-g^{-1}H_m^{\rm row}
                  +\varepsilon_mG_m\succeq0.}
\end{equation}
Formula \eqref{sx:eq:row} provides a rigorously admissible cost but does
not prove \eqref{sx:eq:target} for arbitrary $m$. If it becomes too coarse,
the exact identity suggests bounding the genuinely projected residual
$QL_Kv$, with complete response approximations, rather than the whole
$L_Kv$. Increasing approximation order alone never proves the final
sign. Changing $R$ would change the operator in Theorem~\ref{sx:thm:cofinal}
and would require a different limiting argument.


\subsection{Enclosures, infinite tails, and reproduction}\label{sx:app:tails}
The computational parameters are spatial radius $2$, $8192$ linear
cells of width $h=1/2048$, $128$ positive Gamma labels, theta indices
$j\le8$, and prime powers $n\le32$. Only new row/column integrals are
evaluated. All existing entries are frozen enclosing data, not rounded
point estimates. Their validity is an inherited input, not independently
reproved by a hash comparison.

\subsubsection*{Theta and spatial tails}
For $k=0,1,2$ one has
\[
 |W^{(k)}(r)|\le C_ke^{(4+2k)r-\pi e^{2r}},\qquad
 (C_0,C_1,C_2)=(42,700,10000).
\]
These follow by differentiating the numerator and the factor
$(1+e^{-r})^{-1}$; sufficient series constants are
$4\pi^2\sum j^4e^{-\pi(j^2-1)}$,
$22\pi^3\sum j^6e^{-\pi(j^2-1)}$, and
$80\pi^4\sum j^8e^{-\pi(j^2-1)}$.
Their interval checks are below the displayed integers.
The omission in each derivative through order two is at most
$10^{-90}$, by
$20000\,9^8e^{-243}<10^{-90}$ and a geometric ratio below $1/2$.
For $E=e^4$ and $J_m(a,E)=\int_E^\infty t^me^{-at}\,dt$,
\[
 \norm{W^{(k)}}_{L^2(|x|>2)}\le C_k\sqrt{J_{3+2k}(2\pi,E)},
 \qquad \int_{|x|>2}W\le42J_1(\pi,E).
\]
These are evaluated by the exact integration-by-parts recurrence.
For derivative entries the integrated function is
$v_j''=W''\cos(\gamma_jx)-2\gamma_jW'\sin(\gamma_jx)
-\gamma_j^2W\cos(\gamma_jx)$.
The finite-interval error norm is bounded by
$e_j=2\cdot10^{-90}(1+2\gamma_j+\gamma_j^2)$.
Entry errors include $e_id_j+e_jd_i+e_ie_j$, with $d_j$ an upper
bound for $\norm{v_j''}_2$, and the product of spatial tail norms.

\subsubsection*{Prime powers and the complete remainder}
Every nonzero von Mangoldt term through $32$ is included, with its
cross correlations. For $m=33$ and $z=e^{-\pi/2}$, the whole omitted
prime operator has norm at most
\[
 r_{32}=84z^m\left[\frac{m^2}{1-z}+\frac{2mz}{(1-z)^2}
                          +\frac{z(1+z)}{(1-z)^3}\right]
 <3.609\cdot10^{-18}.
\]
It costs $r_{32}G$, not $r_{32}I$ in coefficient coordinates.
Each retained correlation uses the analytic pieces separated at
$\log n/2$, with both spatial tails bounded by Cauchy--Schwarz.

\subsubsection*{Complete Gamma energy from a positive finite sum}
The label $a_\ell=2\ell+1/2$ has multiplier
$2\xi^2/[a_\ell(a_\ell^2+\xi^2)]$; every omitted label is nonnegative.
Let $f_j$ be the linear interpolant, with endpoints set to zero and
extended by zero. For $F(a,b)=4a+(a/2)\log(1+4b/a)$, Plancherel and
Jensen give $t_\Gamma[e]\le F(\norm e_2^2,\norm{e'}_2^2)$.
Split the interpolation error into the ordinary cell error and the
endpoint/tail remainder. The cell Poincare estimates give
\[
 \chi(h)=\frac{h^4}{\pi^4}\left[4+\tfrac12\log(1+4\pi^2/h^2)\right],
 \quad E_6=(1+\eta)\chi(h)D_2+(1+\eta^{-1})\tau I.
\]
Here $\eta=1/1024$ and $\tau$ bounds the sum of the six Gamma energies
of the endpoint/tail remainders. The enclosing values satisfy
$\chi(h)<7.857\cdot10^{-15}$ and $\tau<6.137\cdot10^{-140}$.
For $\epsilon=1/16384$, Young's inequality yields
\[
 t_\Gamma[v_c]\ge
 c^*\bigl[(1-\epsilon)\Gamma_{128}
                    -(\epsilon^{-1}-1)E_6\bigr]c.
\]
The exact label integrals on piecewise-linear functions are evaluated
by exponential recurrences for the cell slopes. Thus the used lower
matrix is
\[
 Q_6^{\rm lo}=(1-\epsilon)\Gamma_{128}
 - (\epsilon^{-1}-1)E_6-c_\Gamma H-P_{32}-r_{32}G.
\]
The complete Gamma energy and the prime tail have therefore not been
discarded or identified with a finite model.

\subsubsection*{Reproduction and limits of the verification}
Use Python-FLINT 0.9.0 for the first three commands; the last uses only
Python's standard library:
\begin{verbatim}
python3 certify_six_zeros.py
python3 certify_matrix_error.py
python3 verify_fixed_split.py
python3 check_exact.py
\end{verbatim}
The enclosing five-function JSON inputs are part of the bundle.
The workflow reproduces this incremental calculation from those inputs;
it does not rederive every earlier integral from scratch. The last
command checks four $6\times6$ interval matrix signs and the final
rational transfer inequalities. It is not an independent validation of
the quadrature algorithms or analytic tail lemmas.

The exact Gamma spectral projections and the inverse $C^{-1}$ are
used analytically and have not been numerically constructed. Positivity
outside $\mathcal L_6$ in $PK$, uniform control in increasing source order,
and full cofinal exhaustion have not been established. The global
one-half lower bound and RH are not proved. This work has not undergone
proof-assistant formalization or independent peer review. No novelty
claim is made for the standard Schur and spectral-calculus identities.


% Source body preserved under the declared heading, label, notation, and citation maps.
\section{Uniform response brackets and fixed graph debit obstruction}
\label{ur:sec:response}
The preceding five- and six-direction certificates remain valid, but their
fixed graph debit cannot certify an exhaustive nested family of raw zero
functions. This section proves that obstruction and retains the complete
inverse-energy structure in a dimension-independent response approximation.
The bounds below measure Schur approximation error, not a new theta positivity
margin. Each operator polynomial still contains the full Gamma spectral
projection and complete high operator; the projected exhaustion statement
retains the whole zero-jet form-core hypothesis.

\subsection{Fixed sector, definitions, and inputs}
Work on an admitted closed theta sector $K$ of the reciprocal-weight
Hilbert space $\Hh=L^2_{\rm even}(\mathbb R,dx/W)$, where
\[
 W(x)=\frac{\sum_{j\ge1}(4\tau_j^2-6\tau_j)e^{-\tau_j}}{1+e^{-|x|}},
 \qquad\tau_j=\pi j^2e^{2|x|}.
\]
We use its positive, even, smooth theta weight with rapid decay of its
derivatives. Physical $L^2$ norms use Lebesgue measure. The Gamma form is
\[
 t_\Gamma[u]=\int_0^\infty
 \frac{e^{-s/2}}{1-e^{-2s}}\norm{u(\cdot+s)-u}_2^2\,ds.
\]
Assume ${X}_K=\operatorname{Dom}(t_\Gamma)\cap K$ is dense in $K$,
and let $\Gamma$ represent this closed restricted form.
Let $D=D^*$ be the compression to $K$ of the complete correction
\[
 c_\Gamma\norm u_2^2+
 2\Re\sum_{n\ge2}\frac{\Lambda(n)}{\sqrt n}
 \int u(x+\log n)\overline{u(x)}\,dx,
 \qquad c_\Gamma=\log\pi-\psi(1/4).
\]
Let $D_b$ denote the correction on the full weighted space,
$q_0[u]=t_\Gamma[u]-\langle D_bu,u\rangle$, and $D=P_KD_b|_K$.
Put $L=\Gamma-D$ and $q=q_0|_{{X}_K}$, and suppose
$\norm{D_b}\le d$. These closed-form and operator-domain assumptions
are retained explicitly; the analysis uses the representation and
spectral calculus of \cite[Sections 2.3 and 3.1]{gbTeschl}.
All prime powers and the complete Gamma form enter these definitions.

Keep $R>d$ fixed and define
\[
 P=\mathbf1_{[0,R)}(\Gamma),\quad Q=I-P,\quad
 G=\Gamma|_{QK},\quad g=R-d.
\]
Then
\[
 A=P\Gamma P-PDP,\quad B=-QDP,\quad C=G-QDQ\succeq gI,
 \qquad S=A-B^*C^{-1}B.
\]
Here $G$ is the high Gamma operator, not a coefficient Gram matrix.
The low operator $S$ is bounded on $PK$, which need not be finite
dimensional. No separate reduction of $K$ by the original unrestricted
Gamma and prime operators is assumed. A physical bridge and reduction
outside $K$ remain additional inputs to a global one-half conclusion.

The numerical specialization is
\begin{equation}\label{ur:eq:constants}
 R=262144,\qquad d=36,\qquad g=262108.
\end{equation}
One explicit bound for this $d$ uses $W(r)\le42e^{4r-\pi e^{2r}}$,
$\norm W_\infty<2$, and the weighted shift bound
$\sup_x\sqrt{W(x+\log n)W(x)}\le42ne^{-\pi n/2}$.
With $z=e^{-\pi/2}$ it gives
\[
 \norm D\le2c_\Gamma+
 84z^2\left[\frac4{1-z}+\frac{4z}{(1-z)^2}
                     +\frac{z(1+z)}{(1-z)^3}\right]<36.
\]
The same bound applies to every finite prime-head correction. We use
this complete analytic bound as an input; no prime quadrature is rerun.

\subsection{Why fixed graph debits cannot exhaust the raw sources}
For real $T$ put $u_T=W\cos(Tx)$. These raw functions are in the full
operator domain: $u_T\in H^2$, boundedness of $W$, and the Gamma
multiplier bound below put the weighted row $W d_\Gamma(D)u_T$ in
$\Hh$; form testing admits its compression when $u_T\in K$.
Write $m(u)=\int\cosh(x/2)u(x)\,dx$ and
$\Xi(z)=\xi(1/2+iz)$. In particular, for a real
zero ordinate $T$, exact mean and derivative orthogonality give their
membership in the standard sector
\[
 K=\ker m\ominus\overline{\operatorname{span}\{r_\ell:\ell\ge1\}},
 \quad r_\ell=\kappa^{(2\ell)}-4^{-\ell}\kappa,
 \quad\kappa=2\cosh(x/2)W.
\]
Indeed $m(u_T)=\Xi(T)/2$ and
$\langle u_T,r_\ell\rangle_\Hh=
[(-1)^\ell T^{2\ell}-4^{-\ell}]\Xi(T)$ by the theta Fourier
formula \cite[(1.2)]{fzLagarias} and integration by parts.
There are unbounded real critical-line zero ordinates without assuming
RH; see \cite[Section 25.10(i)]{DLMF}.

\begin{lemma}\label{ur:lem:asympt}
Let $w_1=\int W>0$ and $w_2=\int W^2>0$. Then, as $T\to\infty$,
\[
 \norm{u_T}_\Hh^2\longrightarrow w_1/2,\qquad
 \norm{u_T}_2^2\longrightarrow w_2/2,\qquad
 T^{-4}\norm{u_T''}_2^2\longrightarrow w_2/2.
\]
Moreover $q_0[u_T]=O(1+\log T)$ in absolute value, while
$q_0[u_T]/\norm{u_T}_\Hh^2\to+\infty$.
For admitted zero functions in $K$, this is also the restricted form $q$.
\end{lemma}
\begin{proof}
The first two limits follow from $\cos^2(Tx)=(1+\cos(2Tx))/2$ and
Riemann--Lebesgue. The identity
$u_T''=-T^2W\cos(Tx)-2TW'\sin(Tx)+W''\cos(Tx)$ proves the third.
The decreasing positive exponential-label sum of the Gamma multiplier
gives
\begin{equation}\label{ur:eq:multiplier}
 \tfrac12\log(1+4\xi^2)\le d_\Gamma(\xi)
 \le4+\tfrac12\log(1+4\xi^2).
\end{equation}
For example, compare the sum of
$2\xi^2/[a(a^2+\xi^2)]$ over $a=2\ell+1/2$ with its decreasing
integral; the first term costs at most $4$ in the upper bound.
Jensen and Plancherel imply
\[
 t_\Gamma[u]\le4\norm u_2^2+
 \frac{\norm u_2^2}{2}
 \log\left(1+\frac{4\norm{u'}_2^2}{\norm u_2^2}\right).
\]
Since $\norm{u_T'}_2=O(T)$, this is $O(1+\log T)$.
The bounded correction costs at most $d\norm{u_T}_\Hh^2$.

For the divergence, the unitary Fourier transform is
$\widehat u_T(\xi)=\tfrac12[\widehat W(\xi-T)+\widehat W(\xi+T)]$.
Its mass on every fixed bounded frequency interval tends to zero by
$L^2$ translation tails, while its total mass tends to $w_2/2>0$.
The lower multiplier in \eqref{ur:eq:multiplier} tends to infinity off
large intervals. Therefore $t_\Gamma[u_T]\to\infty$; the bounded
correction and the first norm limit prove the last assertion.
\end{proof}

\begin{theorem}[Failure of a fixed second-derivative debit]\label{ur:thm:obstruction}
For fixed $c_0,c_1\ge0$, $c_2>0$ and $g>0$, define
\[
 \mathcal F[u]=q_0[u]-\frac1g
  \left(c_0\norm u_\Hh^2+c_1\norm u_2^2+c_2\norm{u''}_2^2\right).
\]
Then
\begin{equation}\label{ur:eq:quartic}
 \lim_{T\to\infty}\frac{\mathcal F[u_T]}
                   {T^4\norm{u_T}_\Hh^2}
 =-\frac{c_2}{g}\frac{w_2}{w_1}<0.
\end{equation}
In particular this applies to the admissible row cost
\[
 c_0=1458,\quad c_1=\tfrac{153}{8}a^2,\quad
 c_2=306b^2=\frac{1224}{40972801},
\]
where $a=4+\tfrac12\log(6401)-3200/6401$, $b=2/6401$.
\end{theorem}
\begin{proof}
Divide by $T^4\norm{u_T}_\Hh^2$ and apply Lemma~\ref{ur:lem:asympt}.
Only the derivative debit survives. Restricting $T$ to unbounded real
zero ordinates gives the same conclusion on actual theta zero functions.
\end{proof}
Thus any nested raw-source family eventually containing all these
functions cannot satisfy $\mathcal F\ge-\varepsilon_j\norm\cdot^2$
with $\varepsilon_j\downarrow0$ at every stage. Fix one sufficiently
high zero function with a strictly negative value and then let $j$
increase. An even smaller enclosing lower matrix cannot repair this.
This is a failure of a sufficient form, not a negative direction of $q$;
in fact $q_0[u_T]$ is eventually positive. Finite certificates already
obtained with this debit are unaffected.

\begin{proposition}[The exact graph norm does not cure the fixed-gap loss]
Put $r(v)=q[v]/\norm v^2$. For $v\in\operatorname{Dom}L$,
\[
 \frac{q[v]-g^{-1}\norm{Lv}^2}{\norm v^2}
 \le r(v)-r(v)^2/g.
\]
Even retaining the high projection gives
\[
 \frac{q[v]-g^{-1}\norm{QLv}^2}{\norm v^2}
 \le r(v)-r(v)^2/g+(R+2d)^2/g.
\]
Both right sides tend to $-\infty$ along the admitted high zero functions.
\end{proposition}
\begin{proof}
Cauchy--Schwarz gives $\norm{Lv}\norm v\ge|q[v]|$.
Also $PLv=APv+B^*Qv$ and
$\norm{PLv}\le(R+2d)\norm v$. Use
$\norm{QLv}^2=\norm{Lv}^2-\norm{PLv}^2$ and the divergent Rayleigh
quotient of Lemma~\ref{ur:lem:asympt}.
\end{proof}
The exact identity
$S[Pv]=q[v]-\langle QLv,C^{-1}QLv\rangle$ remains valid.
The obstruction concerns replacing $C^{-1}$ by the single worst-case
number $g^{-1}$ on these raw functions. It does not exclude projected
low sources, different lifts, or the actual signed Schur operator.

\subsection{Monotone brackets for the complete high inverse}
Define bounded operators
\[
 {\mathsf K}=G^{-1/2}QDQG^{-1/2},\qquad
 {\mathsf T}=G^{-1/2}B,\qquad \kappa_0=d/R<1.
\]
Then ${\mathsf K}={\mathsf K}^*$, $\norm{{\mathsf K}}\le\kappa_0$, and
$\norm{{\mathsf T}}^2\le d^2/R$. The form factorization and its inverse are
\begin{equation}\label{ur:eq:factor}
 C=G^{1/2}(I-{\mathsf K})G^{1/2},\qquad
 C^{-1}=G^{-1/2}(I-{\mathsf K})^{-1}G^{-1/2}.
\end{equation}
These are form/representation identities for the possibly unbounded
$G$. No commutation between $G$ and $D$ is assumed.
Consequently $S=A-{\mathsf T}^*(I-{\mathsf K})^{-1}{\mathsf T}$ on the entire low sector.

For $n\ge0$ set
\[
 p_n(x)=\sum_{k=0}^{2n-1}x^k\quad(p_0=0),\qquad
 t_n=\frac{\kappa_0^{2n}}{1-\kappa_0},
\]
and define the upper and lower operators
\begin{equation}\label{ur:eq:brackets}
 U_n=A-{\mathsf T}^*p_n({\mathsf K}){\mathsf T},\qquad
 F_n=U_n-t_n{\mathsf T}^*{\mathsf T}.
\end{equation}

\begin{theorem}[Dimension-independent complete response brackets]\label{ur:thm:brackets}
On all of $PK$,
\begin{align}
 F_n&\preceq F_{n+1}\preceq S\preceq U_{n+1}\preceq U_n,\label{ur:eq:monotone}\\
 0\preceq U_n-S&\preceq t_n{\mathsf T}^*{\mathsf T},\qquad
 0\preceq S-F_n\preceq t_n{\mathsf T}^*{\mathsf T}.\label{ur:eq:relative-error}
\end{align}
Both sequences converge to $S$ in operator norm, with
\begin{equation}\label{ur:eq:absolute}
 \norm{U_n-S},\ \norm{S-F_n}
 \le e_n:=\frac{d^2}{R}\frac{(d/R)^{2n}}{1-d/R}.
\end{equation}
The constants do not depend on a source dimension, zero multiplicity,
frequency spacing, or raw coefficient condition number.
\end{theorem}
\begin{proof}
For $x\in[-\kappa_0,\kappa_0]$, the exact geometric identity gives
\[
 (1-x)^{-1}-p_n(x)=\frac{x^{2n}}{1-x},\qquad
 0\le\frac{x^{2n}}{1-x}\le t_n.
\]
Apply the spectral calculus to ${\mathsf K}$ and take the congruence by ${\mathsf T}$.
This proves \eqref{ur:eq:relative-error}. Moreover
\[
 p_{n+1}(x)-p_n(x)=x^{2n}(1+x)\ge0,
\]
and
\[
 (t_n-t_{n+1})-x^{2n}(1+x)
 =\kappa_0^{2n}(1+\kappa_0)-x^{2n}(1+x)\ge0.
\]
These prove both monotonicities. Finally use
$\norm{{\mathsf T}}^2\le d^2/R$.
\end{proof}
The signed terms ${\mathsf T}^*{\mathsf K}^k{\mathsf T}$ are retained. In particular, $U_1$
contains both ${\mathsf T}^*{\mathsf T}$ and the possibly signed ${\mathsf T}^*{\mathsf K}{\mathsf T}$.
The scalar tail bounds only the remaining infinite series.

For \eqref{ur:eq:constants}, $\kappa_0=9/65536$ and
\[
 e_n=\frac{324}{65527}\left(\frac9{65536}\right)^{2n}.
\]
Exact rational arithmetic gives the following outward bounds.
\begin{center}
\begin{tabular}{ccc}\toprule
$n$ & Degree of $p_n$ & Uniform Schur error $e_n$ \\\midrule
0 & empty sum & $<4.945\cdot10^{-3}$\\
1 & 1 & $<9.33\cdot10^{-11}$\\
2 & 3 & $<1.76\cdot10^{-18}$\\
3 & 5 & $<3.32\cdot10^{-26}$\\\bottomrule
\end{tabular}
\end{center}
For example, on the entire low sector,
\[
 S\succeq A-{\mathsf T}^*(I+{\mathsf K}){\mathsf T}-e_1I.
\]
This is an approximation bound, not proof that its right side is
nonnegative. The index $n$ is a polynomial order in a complete bounded
operator. It is \emph{not} a finite response-space dimension: each
application of ${\mathsf K}$ still involves the full high sector and exact Gamma
spectral projections.

\subsection{A complete prime-tail error in the same Schur norm}
The fixed split depends on $\Gamma$, so it stays unchanged when only
the prime correction is truncated. This permits a sharper perturbation
bound than termwise subtraction of inverse formulas.

\begin{theorem}[Variational Schur stability]\label{ur:thm:stability}
Let $D_1,D_2$ be bounded self-adjoint corrections on the same sector,
with $\norm{D_i}\le d<R$ and $\norm{D_1-D_2}\le r$. Use the same
$\Gamma,P,Q$ for both Schur operators $S_i$. Then
\begin{equation}\label{ur:eq:stability}
 \boxed{\norm{S_1-S_2}\le r\left[1+\left(\frac d{R-d}\right)^2\right].}
\end{equation}
\end{theorem}
\begin{proof}
For fixed low $\ell$, the unique minimizing high vector for $q_i$ is
$h_i=-C_i^{-1}B_i\ell$, and
$\norm{\ell+h_i}^2\le[1+(d/g)^2]\norm\ell^2$.
Evaluate $q_1$ on $\ell+h_2$ to obtain
\[
 S_1[\ell]\le S_2[\ell]+r[1+(d/g)^2]\norm\ell^2.
\]
Interchanging the indices gives the opposite inequality. Both minimizers
are admitted in the common high form domain because $C_i\succeq gI$.
\end{proof}

For the complete prime tail beyond $M$, the theta shift envelope gives
\[
 r_M\le84z^{M+1}\left[
 \frac{(M+1)^2}{1-z}+\frac{2(M+1)z}{(1-z)^2}
                     +\frac{z(1+z)}{(1-z)^3}\right],\qquad z=e^{-\pi/2}.
\]
This bounds the entire integer tail, hence all prime powers. In
particular $r_{32}<3.609\cdot10^{-18}$. The two complete correction
bounds remain below $36$, so \eqref{ur:eq:stability} yields
\[
 \norm{S(D)-S(D_{32})}<3.610\cdot10^{-18}.
\]
Combining this with the cubic response bound on $D_{32}$ gives
\begin{equation}\label{ur:eq:combined}
 \norm{S(D)-U_2(D_{32})}<5.37\cdot10^{-18}.
\end{equation}
The constants are added here only after both have been converted into
the same operator norm on the same exact low space. Gamma projection,
spatial integration, and finite matrix evaluation errors are not
included in \eqref{ur:eq:combined}. Holding $M=32$ fixed leaves a nonzero
certified error and does not prove an exact endpoint.

\subsection{Cofinal low signs and the remaining arithmetic problem}
\begin{lemma}[Projected exhaustion from an admitted form core]\label{ur:lem:core}
Assume the algebraic span $\mathcal J_{\rm fin}$ of all admitted exact
zero jets is a form core for ${X}_K$. Enumerate every required generator
and let $V_j$ be the spans of the initial finite sets. Then
${\mathcal L}_j=PV_j$ is nested and $\overline{\bigcup_j{\mathcal L}_j}=PK$.
\end{lemma}
\begin{proof}
Every $\ell\in PK$ lies in ${X}_K$ because its Gamma spectral support
is bounded by $R$. Choose $v_n\in\mathcal J_{\rm fin}$ tending to
$\ell$ in form norm. The spectral projection $P$ is a contraction in
both Hilbert and Gamma form norms, so $Pv_n\to\ell$. Each $v_n$
belongs to some $V_j$.
\end{proof}
This lemma transfers an explicit core input; it is not a new independent
proof of that input. It also requires the entire list of generators,
not just the real simple-zero functions that have been tested finitely.

Let ${\mathcal L}_j\subset PK$ be nested finite-dimensional spaces with dense
union. Choose a basis map $J_j$ with actual low Gram $G_j=J_j^*J_j$.
Let $M_j\to\infty$ be prime cutoffs and $n_j\to\infty$ response orders.
Suppose a computed Hermitian matrix $\widehat U_j$ encloses
$J_j^*U_{n_j}(D_{M_j})J_j$ in the relative sense
\[
 -\eta_jG_j\preceq
 \widehat U_j-J_j^*U_{n_j}(D_{M_j})J_j
 \preceq\eta_jG_j.
\]
The coefficient norm cannot replace $G_j$ without a verified comparison.

\begin{theorem}[Two independent approximation indices]\label{ur:thm:cofinal}
If
\[
 \widehat U_j\succeq-\delta_jG_j,\qquad
 \delta_j+\eta_j+e_{n_j}
 +r_{M_j}[1+(d/g)^2]\longrightarrow0,
\]
then $S\succeq0$ on $PK$ and $q\ge0$ on ${X}_K$.
\end{theorem}
\begin{proof}
Theorems~\ref{ur:thm:brackets} and \ref{ur:thm:stability} imply
\[
 J_j^*S(D)J_j\succeq-
 \bigl[\delta_j+\eta_j+e_{n_j}
           +r_{M_j}(1+(d/g)^2)\bigr]G_j.
\]
Apply this to a fixed vector in the nested union and let $j\to\infty$.
The boundedness of $S$ and density give its nonnegativity everywhere.
For $u=\ell+h\in{X}_K$, exact square completion gives
$q[u]=S[\ell]+\norm{C^{1/2}(h+C^{-1}B\ell)}^2\ge0$.
\end{proof}

The analytic response error $e_n$ is now uniform in source dimension.
It can be assigned independently by choosing
\[
 n\ge\max\left\{0,\left\lceil
 \frac{\log(d^2/[R(1-d/R)\epsilon])}{2\log(R/d)}
 \right\rceil\right\}
\]
for a desired error $\epsilon>0$ when $d>0$. The case $d=0$ is exact
without a response cost. This does not provide an effective algorithm
for evaluating the projected operator moments, nor the signs of
$\widehat U_j$ as $j$ increases. An actual zero-jet form-core theorem,
if admitted, supplies the projected density; finitely many real simple
zeros do not establish it. All required frequencies and multiplicity
jets must be represented.

The signed operator whose finite restrictions now require evaluation is
\[
 U_n=A-\sum_{k=0}^{2n-1}{\mathsf T}^*{\mathsf K}^k{\mathsf T}.
\]
Its arithmetic signs remain unknown. No general theorem can infer those
signs from $\norm{{\mathsf K}}<1$ alone: keeping $B,C$ fixed and changing $A$
can make $S$ positive, zero, or negative. At an exact zero direction,
the lower bracket $F_n$ may stay negative for every finite $n$ even
though $S=0$. This is why vanishing, explicitly paid errors are a valid
endpoint target; a uniformly strict positive gap is not necessary.

\subsection{Verification boundary and reproduction}
The asymptotic obstruction, inverse brackets, monotonicity, and Schur
stability are proved above. Their application uses the stated theta
regularity, sector-domain, and bounded-correction inputs. No new
low-frequency theta matrix, spectral projection, Gamma inverse, or
response moment is computed. No previously successful finite integral
certificate is rerun or invalidated by the obstruction.

The standard-library checker
\begin{verbatim}
python3 check_response.py
\end{verbatim}
verifies the displayed rational constants and three exact rational
noncommuting block examples, including positive, zero, and negative
Schur operators. It checks both bracket monotonicities and the inverse
identity. Those examples validate algebraic implementation, not actual
theta positivity. The proofs do not depend on a floating eigenvalue
test. A second script, \texttt{check\_theta\_constants.py}, uses
Python-FLINT 0.9.0 to enclose only the elementary exponential/digamma
bounds for $\norm D<36$ and $r_{32}<3.609\cdot10^{-18}$; it performs no
theta matrix integration. This note is not a proof-assistant formalization or independent
peer review, and no novelty claim is made for Neumann series or
variational Schur identities. The global one-half bound and RH remain
unproved.


\section{Fixed lattice positivity and signed infinite block coupling}\label{sec:fixed-lattice}
A fixed real frequency lattice permits a synthesis estimate independent of mode count. Its geometry differs from the actual complex zero family. The signed block result then describes what is still needed to combine positive subspaces.

\subsection{The complete form and elementary theta bounds}
The same actual weight and complete form are used on the full-line space $H=L^2(\R,dx/W)$ and $X=\{u\in H:t_\Gamma[u]<\infty\}$ in this section. Put $r=|x|$ and $\tau_j=\pi j^2e^{2r}$. The full-line symbols are local to this section; the physical even sector is obtained by restricting to even functions. The Fourier transform is the common unitary angular-frequency transform. This extension does not change any complete prime or Gamma term.

Define positive finite constants
\begin{align*}
 c_W&=(4\pi^2-6\pi)/2,\qquad
 C_W=4\pi^2\sum_{j\ge1}j^4e^{-\pi(j^2-1)},\\
 D_W&=71\sum_{j\ge1}e^{-\pi j^2}
       \big((\pi j^2)^3+3(\pi j^2)^2+6\pi j^2+6\big),\\
 D_0&=C_W\left(c_\Gamma e^{-\pi}
            +2\sum_{n\ge2}(\log n)n^{3/2}e^{-\pi n}\right).
\end{align*}

\begin{lemma}\label{lat:lem:theta}
$W$ is smooth, positive and even. Moreover
\[
 c_We^{4r-\pi e^{2r}}\le W(x)\le C_We^{4r-\pi e^{2r}},\qquad
 \norm{W''}_1\le D_W,
\]
and, on the entire space $H$,
\begin{equation}\label{lat:eq:perturbation}
 \big|c_\Gamma\norm u_2^2+p[u]\big|\le D_0\norm u_H^2.
\end{equation}
The Gamma multiplier obeys
\begin{equation}\label{lat:eq:gamma}
 t_\Gamma[u]=\int d_\Gamma(\xi)|\widehat u(\xi)|^2\,d\xi,
 \qquad d_\Gamma(\xi)\ge\tfrac12\log(1+4\xi^2).
\end{equation}
\end{lemma}
\begin{proof}
The two-sided envelope and smoothness are the actual-theta estimates of Appendix~\ref{sw:sec:theta}. We give the additional explicit derivative constant required for the lattice leakage bound.

For $r>0$, write $h=(1+e^{-r})^{-1}$ and
$F(\tau)=(4\tau^2-6\tau)e^{-\tau}$. Differentiation with
$\tau'=2\tau$ gives the polynomials
\[
 F_r=(-8\tau^3+28\tau^2-12\tau)e^{-\tau},\qquad
 F_{rr}=(16\tau^4-104\tau^3+136\tau^2-24\tau)e^{-\tau}.
\]
Since $\tau\ge\pi>3$, their absolute values are at most
$6\tau^2e^{-\tau},19\tau^3e^{-\tau},67\tau^4e^{-\tau}$, respectively.
Using $|h'|,|h''|\le1/4$ bounds $|W''|$ by
$71\sum_j\tau_j^4e^{-\tau_j}$. Integrating on both half-lines yields
$\norm{W''}_1\le71\sum_j\int_{\pi j^2}^\infty t^3e^{-t}\,dt=D_W$.
These envelopes also justify $W,W'\in L^1$ and termwise differentiation.

For $s=\log n$, put $r=|x|$, $t=|x+s|$, so $r+t\ge s$.
The arithmetic--geometric mean inequality and monotonicity of
$2z-\pi e^z$ for $z\ge0$ give
\[
 \sqrt{W(x+s)W(x)}\le C_We^{2(r+t)-\pi e^{r+t}}
                         \le C_Wn^2e^{-\pi n}.
\]
Substitute $v=u/\sqrt W$ and apply Cauchy--Schwarz to the translated
$v$'s. Sum with $\Lambda(n)\le\log n$ and use
$\sup W\le C_We^{-\pi}$ to obtain~\eqref{lat:eq:perturbation}.
This also proves absolute convergence of the complete prime form on $H$.

The Gamma multiplier identity and its lower bound are proved in Appendix~\ref{sw:sec:theta}; they use the full unbounded Gamma integral and hold with extended nonnegative values.
\end{proof}

\subsection{A whole infinite-dimensional lattice subspace}
Fix a spacing $d>0$ and let $L=2\pi/d$. Define periodizations on
$[-L/2,L/2]$ by
\[
 w_1(s)=\sum_{n\in\mathbb Z}W(s+nL),\qquad
 w_2(s)=\sum_{n\in\mathbb Z}W(s+nL)^2,
\]
and let $a_d=\min w_2>0$, $b_d=\max w_1<\infty$.
The series converge uniformly, by Lemma~\ref{lat:lem:theta}.
These are fixed-spacing constants; there is no mode-count parameter.

Let $P\in L^2(\R/L\mathbb Z)$ have Fourier coefficients $c_k=0$ for
$|k|<K$, and set $u=WP$. Take $u\in X$ when evaluating $q_0$.
Thus infinitely many Fourier coefficients are permitted.
If even profiles are desired, take $P$ even.

\begin{theorem}[No bound on the number of lattice modes]\label{lat:thm:lattice}
For every integer $K\ge2$, define
\[
 \eta_K=\frac{4D_W^2}{3\pi^2a_dd^2K^2}\left(1+\frac6K\right),
 \qquad
 L_d(K)=\frac{a_d}{2b_d}\log(1+d^2K^2)(1-\eta_K)_+-D_0.
\]
Every profile just specified satisfies
\begin{equation}\label{lat:eq:lattice}
 \boxed{q_0[WP]\ge L_d(K)\norm{WP}_H^2.}
\end{equation}
For fixed $d$, $L_d(K)\to+\infty$. In particular one sufficiently high
cutoff works simultaneously for every finite number of modes and every
profile in this entire periodic energy class.
\end{theorem}
\begin{proof}
Periodization and Parseval give, with no coefficient conditioning loss,
\begin{align*}
 \norm{WP}_H^2&=\int_{-L/2}^{L/2}w_1|P|^2
                       \le b_dL\sum_k|c_k|^2,\\
 \norm{WP}_2^2&=\int_{-L/2}^{L/2}w_2|P|^2
                       \ge a_dL\sum_k|c_k|^2.
\end{align*}
Thus $\norm u_2^2\ge(a_d/b_d)\norm u_H^2$.
For finite sums, $\widehat u(\xi)=\sum_kc_k\widehat W(\xi-kd)$.
Integration by parts twice yields
$|\widehat W(\zeta)|\le D_W/(\sqrt{2\pi}|\zeta|^2)$.
On $|\xi|\le Kd/2$ we have $|\xi-kd|\ge d(|k|-K/2)$.
Cauchy--Schwarz in the coefficient index therefore gives
\begin{align*}
 \norm{\mathbf1_{[-Kd/2,Kd/2]}\widehat u}_2^2
 &\le\sum_k|c_k|^2\,
 \frac{D_W^2K}{\pi d^3}\sum_{k=K}^\infty(k-K/2)^{-4}\\
 &\le\sum_k|c_k|^2\,
 \frac{D_W^2K}{\pi d^3}
              \left(\frac{16}{K^4}+\frac8{3K^3}\right)
 \le\eta_K\norm u_2^2.
\end{align*}
This is an operator estimate on all coefficient vectors, not a modewise
estimate added with coefficient absolute values.

For general $P\in L^2$ of a period, its finite Fourier sums converge
after multiplication by $W$ in both $H$ and ordinary $L^2$, since both
periodized weights are bounded. The low-band estimate extends by continuity.
No convergence in the Gamma graph norm is needed: for the given $u\in X$,
apply~\eqref{lat:eq:gamma} directly outside the band. It gives
\[
 t_\Gamma[u]\ge\tfrac12\log(1+d^2K^2)
                         (1-\eta_K)_+\norm u_2^2.
\]
The coefficient is nonnegative, so the proved norm conversion has the
correct direction. Subtract~\eqref{lat:eq:perturbation} to finish.
\end{proof}

All constants can be bounded without a numerical spectral certificate.
For example
\begin{align*}
 a_d&\ge c_W^2e^{4L-2\pi e^L},\\
 b_d&\le C_W\left(e^{-\pi}+2\sum_{n\ge1}
       e^{4(n-1/2)L-\pi e^{2(n-1/2)L}}\right).
\end{align*}
Set $C_\eta=16D_W^2/(3\pi^2a_dd^2)$. For a desired $h>0$, it suffices to take
\[
 K\ge\left\lceil\max\left\{2,\sqrt{2C_\eta},
 d^{-1}\exp\!\left(\frac{2b_d(D_0+h)}{a_d}\right)\right\}\right\rceil.
\]
This is an existence prescription, potentially extremely large. It is
not a practical threshold calculation. Its dependence on $d$ is essential.

\paragraph{Centering without dropping a mean term.}
In the physical even sector, the normalized $\kappa$ satisfies the full-domain source
identity $q_0(\kappa,u)=-\frac12\ip\kappa u_H$, let
$\mathcal C=I-\kappa\otimes\kappa$. Then
\[
 q_0[\mathcal Cu]=q_0[u]+\tfrac12|\ip u\kappa_H|^2.
\]
When $L_d(K)\ge0$,~\eqref{lat:eq:lattice} implies
$q_0[\mathcal Cu]\ge L_d(K)\norm{\mathcal Cu}_H^2$.
This centering corollary retains the exact mean correction; its source
identity follows on the full domain from Proposition~\ref{vrg:hyp:realization}; it is not a consequence of periodicity.

\subsection{An extension theorem for infinitely many signed blocks}
The lattice result supplies one example of a stable high subspace.
To add it or other subspaces to a low part, coupling still matters.
Here is a criterion whose constants do not count the number of blocks.

Let $V_0$ be finite dimensional with a positive semidefinite matrix $A$.
For each high block $V_j$, choose a positive definite energy metric $P_j$,
and put $y_j=P_j^{1/2}v_j$. Suppose the high part of the actual form
obeys, on all finite block sums,
\[
 q_{\rm high}[v]\ge\ip{(I+E)y}{y},\qquad
 E=E^*,\quad E\succeq-\kappa I,\quad0\le\kappa<1,
\]
where $E$ is one bounded operator on $\bigoplus_{j\ge1}V_j$.
It contains all signed high--high couplings. Write the low--high term as
$2\operatorname{Re}\ip{b(v_0)}y$, with
$b(v_0)=(P_j^{-1/2}B_jv_0)_j$.

\begin{theorem}[Uniform relative coupling]\label{lat:thm:blocks}
Assume the low form is at least $\ip{Av_0}{v_0}$ and
\begin{equation}\label{lat:eq:coupling}
 \sum_{j\ge1}\norm{P_j^{-1/2}B_jv_0}^2
                  \le\eta^2\ip{Av_0}{v_0},\qquad
 \eta^2\le1-\kappa.
\end{equation}
Then the full form is nonnegative on every finite block sum, with no
restriction on its complexity. If the inequality is strict, it also obeys
\[
 q[v_0+v]\ge\mu\left(\norm{A^{1/2}v_0}^2+\sum_j\norm{y_j}^2\right),
 \quad
 \mu=\frac{2-\kappa-\sqrt{\kappa^2+4\eta^2}}2>0.
\]
For a closed semibounded realization this extends to the form-norm closure
of the algebraic sum. An independent form-core premise is needed to
identify that closure with a desired full domain.
\end{theorem}
\begin{proof}
Condition~\eqref{lat:eq:coupling} implies $b(\ker A)=0$ and a factorization
$b=K A^{1/2}$ with $\norm K\le\eta$ on $\operatorname{ran}A^{1/2}$.
Set $x=A^{1/2}v_0$. The entire signed form is bounded below by
\begin{align*}
 \norm x^2+2\operatorname{Re}\ip{Kx}y+(1-\kappa)\norm y^2
 ={}&(1-\kappa)\norm{y+(1-\kappa)^{-1}Kx}^2\\
 &+\norm x^2-(1-\kappa)^{-1}\norm{Kx}^2\ge0.
\end{align*}
Alternatively its lower bound is the scalar two-by-two form with diagonal
$1,1-\kappa$ and off-diagonal $-\eta$, whose minimum eigenvalue is $\mu$.
All estimates use the complete infinite coupling norm, hence survive every
finite truncation. Form-norm continuity gives the closure assertion.
\end{proof}

This is a sufficient theorem, not a claim that its premises have been proved
for the actual zero-jet decomposition. The kernel condition is essential:
a positive semidefinite full form with $q[v_0]=0$ must satisfy
$q(v_0,w)=0$ for all $w$, by varying a real or complex scalar in
$q[v_0+tw]$. Arbitrarily small nonzero coupling to an exact zero direction
already prevents positivity. Pseudoinverses must therefore not be used
without this compatibility. See~\cite{Boyd} for the finite Schur framework.

An absolute row-sum bound is one possible way to estimate $E$, but is
not required. The one-sided signed operator bound $E\succeq-\kappa I$
can retain cancellation discarded by absolute values.

\subsection{Two examples separating accumulation from cancellation}
Positive diagonal blocks do not suffice. For
\[
 q_N(x,y)=x^2+\sum_{j=1}^Ny_j^2+\frac{x}{5}\sum_{j=1}^Ny_j
 =\sum_{j=1}^N(y_j+x/10)^2+(1-N/100)x^2,
\]
every individual low--high pair is positive definite, yet $q_{101}$ has a
negative direction. The collective coupling has norm $\sqrt N/10$.
This example is not a theta counterexample.

There is also a real model in which signs prevent accumulation even though
absolute row sums diverge. Let $x,y_j\in\R^2$, $j\ge1$, and
\[
 J=\begin{pmatrix}0&-1\\1&0\end{pmatrix},\quad
 E_{ij}=\frac{J}{4(i-j)}\ (i\ne j),\ E_{ii}=0,
 \quad (Kx)_j=\frac{x}{4j}.
\]
The block matrix $E$ is self-adjoint. The discrete convolution
$H_0a(i)=\sum_{j\ne i}a(j)/(i-j)$ has norm at most $\pi$ on $\ell^2(\mathbb Z)$.
Indeed its Fourier multiplier is $i(\theta-\pi)$ on $0<\theta<2\pi$,
up to the Fourier sign convention, by the sawtooth Fourier series.
Restriction to the positive integers is a compression. Thus
\[
 \norm E\le\pi/4,\qquad K^*K=(\pi^2/96)I,
 \qquad \pi^2/96<1-\pi/4.
\]
This $\ell^2$ argument is elementary; broader sharp discrete Hilbert
transform bounds are established in~\cite{Hilbert}.
Theorem~\ref{lat:thm:blocks} applies to
$q=\norm x^2+2\ip{Kx}y+\ip{(I+E)y}y$.
Using $\pi<22/7$, the scalar comparison matrix has determinant at least
$131/1176$ and trace $17/14$ when its entries are conservatively fixed at
$1$, $3/14$, and $-\sqrt{121/1176}$. Its smallest eigenvalue is at least
determinant divided by trace, giving
\[
 \boxed{q_N\ge\frac{131}{1428}
       \left(\norm x^2+\sum_{j=1}^N\norm{y_j}^2\right)
       \quad\text{for every }N.}
\]
Meanwhile $\sum_{j\ne i}\norm{E_{ij}}=\infty$.
This supplies a precise cancellation mechanism, not an identification of
the theta interaction matrix with this model.

\subsection{What remains for actual zero jets and the half bound}
Theorem~\ref{lat:thm:lattice} is a proved statement for the explicit complete
theta form on a restricted infinite-dimensional class. It is stronger than
a fixed-mode-count statement on that class, but its arithmetic geometry is
different from that of actual zero jets. Such jets are not supplied as
real frequencies on one fixed lattice. Polynomial factors from multiplicity
and complex frequency offsets are also outside the periodic class.
Letting $d\downarrow0$ does not repair the issue: $a_d,b_d$ and the cutoff
may deteriorate, and no uniform estimate in $d$ has been proved.

For an actual jet family, the missing substitute is an estimate on whole
linear combinations. Either prove stable synthesis Gram bounds and a
uniform low-frequency leakage bound, or establish the complete signed
interaction bound and relative coupling inequality of
Theorem~\ref{lat:thm:blocks} in the actual block metrics. An upper bound on the
number of zeros in each window gives neither estimate. Individual-window
positivity does not control the matrix of interactions between windows.

At the unshifted endpoint, positivity of the low block and exact
annihilation of its kernel must also be established. A certificate for
$q_0+a\norm{\cdot}_H^2$ with $a>0$ is not that input. To use a sequence of
shifted certificates instead, all assumptions must hold on an appropriate
nested form-core exhaustion with shifts tending to zero. No fixed-stage
margin or fixed-complexity threshold supplies this automatically.

No new global numerical coefficient, universal actual-jet positivity,
or RH proof is asserted. Exact symbolic checks of the constants and synthetic examples are recorded with the validation materials. The infinite-dimensional conclusions rest on the proofs above, not those finite tests.


\subsection{A later growing-band extension}
Appendix~\ref{app:growing-zero-bands} gives a coefficient-independent
strip and full-prefix extrapolation argument. It proves the two-sided
ordinary-$L^2$ asymptotic on the entire actual zero-jet space in a single
band of width $o(\log T\sqrt{\log\log T})$, including every internal
coupling. A conservative positive result holds for fixed boundary coefficient
$d<\pi\sqrt2/7200$. The exact inverse-retaining Schur consequence applies
inside a union already contained in a covered band. It does not glue
arbitrary bands or settle the global sign. The optional explicit count uses
the cited zero-count theorem, including its finite-height numerical
verification dependency; the qualitative conclusion uses the classical
asymptotic zero count.

\section{The remaining sign}\label{sec:remaining}
All reductions retain the complete Gamma integral, every prime power, the original mean and the relevant minimum domains. The historical local-sign and response-contract branches remain valid at their stated premises. The newer response gate improves the conditional finite-stage comparisons, but each certified or reported stage still has strictly positive loss. The evidence classes are summarized in Appendix~\ref{app:current-evidence}.

Exact half-eigenstates and essential spectral bottom $1/2$ leave open a lower discrete spectrum. Delta-scaled complements control infinite blocks but do not assign every finite Schur sign as $\delta\downarrow0$. Full zero-jet cores close an approximation issue for the actual operator; they do not prove the signs of the resulting finite Hermitian forms. Fixed-complexity coercivity, arbitrarily many sparse windows, a single full band of width $11\log T$, sparse families of entire growing bands under exponent-$12$ spacing, and a fixed real lattice solve distinct restricted problems. They do not control arbitrary unions of adjacent actual-zero bands, unrestricted multiplicity budgets, or the general full-complement low--high interaction operator. Section~\ref{sec:low-sparse} does control the exact finite and infinite variational response of the declared sparse family against a fixed admissible low space, with a cost tending to zero in its actual weighted Gram. The general response estimate does not by itself supply positivity of that low space. Section~\ref{sec:four-zero-certificate} supplies a positive margin for one specific four-zero space and pays its selected sparse tail; it supplies neither uniform constants for growing low spaces nor the omitted adjacent-window couplings. Boundedly many sparse classes cannot cover all high zero windows.

The signed block theorem makes the missing collective conditions explicit: the full signed high interaction, complete relative low--high coupling, and compatibility with true null directions must all be justified in the actual metrics. A collection of individually positive blocks or a finite truncation cannot replace those conditions. No uniform threshold as lattice spacing tends to zero is supplied.

The response toolkit in Appendix~\ref{rtc:sec:tools} separates source spaces from inverse source images, keeps the full residual row Gram, and states what fixed or ordered-varying cofinal certificates would imply. Fixed inverse conditioning, exact finite optimization and monotone refinement do not themselves certify those signs. Appendix~\ref{ext:sec:exterior} explains why infinitely many exact equality directions also prevent a uniformly strict endpoint overcharge.

The remaining assertion is nonnegativity of every finite full-zero-jet form in Theorem~\ref{thm:actual-sign}. That assertion, the global variance half bound, and the Riemann hypothesis are not proved here. No zero-simplicity, historical-priority, journal-acceptance, human-peer-review or proof-assistant-verification claim is made.


The graph-transfer results in Sections~\ref{gs:sec:transfer}
and~\ref{sx:sec:six} address a different restricted domain: a five- or
six-dimensional projected low source with the entire complementary
Gamma-high form space at one fixed spectral split. The sixth-direction
extension retains the complete matrix row cost and proves a conditional
fixed-split exhaustion criterion, but supplies no further low signs or
cofinal exhaustion. It neither changes the global half boundary nor
licenses a zero-height interpretation of the Gamma spectral projection.


Section~\ref{ur:sec:response} separates a further analytic issue from
that remaining sign problem. A fixed worst-gap graph debit fails along
high raw zero functions, even if its graph norm is exact or retains the
high projection. Complete inverse-energy brackets instead give a uniform
geometric response error, with a separate complete prime-tail debit in
the same Schur norm. The source dimension and operator-polynomial order
are independent indices; a polynomial step still uses the full high
operator and Gamma projection. No response-accuracy bound alone supplies
the low signs or an evaluated cofinal certificate.

\appendix
\section{Complete scalar and interface estimates}\label{app:costs}
The estimates below are part of the proofs of Section~\ref{sec:bounds}. They retain the full theta series, every prime tail and the endpoint singularities of the Gamma kernels.


\subsection{Elementary full-series bounds}\label{c:app:scalar}
For $r>0$ set $S_m(r)=\sum_{k=0}^m r^k/k!$.
The positive exponential series and its geometric tail give
\[
 S_m(r)<e^r\le S_m(r)+\frac{r^{m+1}}{(m+1)!}
                  \frac1{1-r/(m+2)}\qquad(0<r<m+2).
\]
The positive logarithm identity is
\begin{equation}\label{c:eq:atanh}
 \log\frac{1+z}{1-z}=2\sum_{k\ge0}\frac{z^{2k+1}}{2k+1},
 \qquad0<z<1.
\end{equation}
Omitted denominators may be bounded below and powers summed geometrically.
For example
\begin{gather*}
 693/1000<842/1215<\log2<25/36<7/10,\\
 1098/1000<7379/6720<\log3<1+1/12+1/60=11/10,\\
 1609/1000<\log5<13724/8505<81/50.
\end{gather*}
The lower bound for $\log5$ follows from the first eight terms at $z=2/3$,
which exceed respectively
$1.33333,.19753,.05267,.01672,.00578,.00210,.00079,.00030$.
These are exact finite decimals summing to $1.60922$.
The upper bound uses the first three terms and their geometric remainder.
Squaring proves
$\sqrt2>140/99$, $\sqrt3>55/32$, $\sqrt5>324/145$;
these imply \eqref{c:eq:coefficientcaps}.
Also $\log n/\sqrt n\le2/e<1$, so $a_n\le1$.
The inequality $\Lambda(n)\le n/2$ follows from $\log n<n/2$ for $n\ge2$.

The arctangent addition formula and alternating series give
\[
 \pi>4\sum_{d\in\{2,3\}}\sum_{k=0}^3
       \frac{(-1)^k}{(2k+1)d^{2k+1}}
 =\frac{1538665}{489888}>\frac{157}{50},
\]
where the last difference is $10417/12247200$.
Also
$\frac{22}{7}-\pi=\int_0^1z^4(1-z)^4/(1+z^2)\,dz>0$.
Thus $\pi^2<10$.
Useful Taylor consequences are
\[
 e^{2/5}<3/2,\quad e^{1/5}<5/4,\quad
 e^{8/5}<5,\quad e^3>401/20>20,\quad e^2>73/10.
\]
The bound at $8/5$ follows from the displayed tail formula with $m=6$.

\subsubsection{The exterior theta envelope}
For $g(t)=(4t^2-6t)e^{-t}$ and $t\ge3$,
\begin{equation}\label{c:eq:thetatotal}
 \sum_{j\ge2}g(j^2t)<128t^2e^{-4t}.
\end{equation}
Indeed bound each term by $4t^2j^4e^{-tj^2}$.
Consecutive ratios for $j\ge2$ are at most
$(3/2)^4e^{-5t}<1/2$, so twice the first majorant pays the complete tail.
Both $g(t)$ and $t^2e^{-4t}$ decrease for $t>3$.

At ${B_0}=4/5$, $S_7(8/5)>99/20$ gives
$t_{B_0}>T_0=15543/1000>31/2$.
The logarithm bounds imply $e^{B_0}<\sqrt5<9/4$, so $1+e^{-{B_0}}>13/9$.
The degree-nine exponential sum at $3$ is
$7280604/362880>401/20$; the degree-four sum at $543/1000$ is
$41297535960801/24000000000000>43/25$.
Hence $e^{15}>20^5(1+5/400)=3240000$ and $e^{T_0}>5572800$.
Since $4T_0^2-6T_0<8731/10$, the leading part of $A({B_0})$ is below
$78579/724464000$.
The entire tail is below $32768/20^{20}<10^{-12}$, while
\[
 \frac1{9216}-\frac{78579}{724464000}
 =\frac{279936}{9216\cdot724464000}>10^{-8}.
\]
This proves $A<1/9216$.
The logarithmic derivatives of its two summands are at most
\[
 \frac1{1+e^x}+\frac{2(8t-6)}{4t-6}-2t\le6-2t<-25,
 \qquad 4-8t<-25
\]
on $x\ge {B_0}$.
Therefore $W({B_0}+h)\le A({B_0}+h)\le Ae^{-25h}$ and
$\int_{B_0}^\infty W\le A/25$.

\subsubsection{A global theta cap and the exterior probability}
Termwise differentiation is justified by locally uniform summable majorants, and gives
\[
 \kappa'(x)=e^{x/2}\sum_{j\ge1}t_jP(t_j)e^{-t_j},
 \qquad P(t)=-8t^2+30t-15.
\]
At zero the first theta term is below $46/49$:
$4\pi^2-6\pi<4(22/7)^2-6(22/7)$, while
$e^\pi>e^{31/10}>22$.
The remaining terms are bounded by $40j^4e^{-3j^2}$,
whose first term is below $1/250$ and successive ratio below $1/10000$.
Thus $\kappa(0)<46/49+40/9999$.
The polynomial $P$ decreases on $[3,\infty)$, with
$P(31/10)=28/25$ and $P(16/5)=-23/25$.
All $j\ge2$ derivative terms are negative.
A positive derivative is therefore possible only on
$0<x<\frac12\log(32/31)<1/62$.
There,
\[
 \kappa'_+(x)<\frac{60}{59}\frac{16}{5}\frac{28}{25}\frac1{22}
 =\frac{2688}{16225}<\frac16.
\]
Hence
\[
 0<\kappa(x)<\frac{46}{49}+\frac{40}{9999}+\frac1{372}
 =\frac{57440653}{60753924}<\frac{19}{20}.
\]
Evenness gives the global assertion and $\sup W<1/2$.
The normalization proved in Section~\ref{a:sec:energy} gives
\[
 \int_0^\infty W
 =\int_0^\infty\frac{\kappa}{2c}\le\frac12\int_0^\infty c\kappa=\frac18.
\]

For completeness, the endpoint cap $\kappa({B_0})<1/4000$ follows from
$t_{B_0}>31/2$, $e^{{B_0}/2}<3/2$, $e^{31/2}>5300000$:
\[
 \kappa({B_0})<\frac{1302}{5300000}
 +\frac{23064}{5300000^4}\frac{10000}{9999}
 <\frac{1302}{5300000}+\frac1{10^9}<\frac1{4000}.
\]
The exponential lower bound is supplied by $S_{24}(31/2)$;
its excess over $5300000$ is
\[
 \frac{37076883949467729399970752217216753}
 {10409396852733332453861621760000}>0.
\]
The tail ratio in this estimate is below $1/10000$.
Furthermore
\[
 -\kappa'-25\kappa
 =e^{x/2}\sum_{j\ge1}t_j(8t_j^2-130t_j+165)e^{-t_j}>0
 \quad(x\ge {B_0}),
\]
since at $t=31/2+u$ the polynomial is $8u^2+118u+72$.
It follows that
\[
 \int_{B_0}^\infty e^{x/2}\kappa(x)\,dx
 <\frac{(3/2)(1/4000)}{49/2}=\frac3{196000},
 \qquad \mu(O)<\frac3{49000}.
\]

\subsubsection{All needed interior point estimates}
Since $t(1/10)>19/5>16/5$, the differentiated polynomial proves that
$\kappa$ and $W$ decrease to the right of $1/10$.
Every tail below is the complete bound \eqref{c:eq:thetatotal}, with the factor $e^{x/2}$ when estimating $\kappa$.

At $x=1/10$,
$t>(157/50)(61/50)>19/5$,
$1+e^{-x}>19/10$, and $e^{19/5}>20(11/5)=44$.
The leading bound for $W$ is $437/1045$; its tail is below
$46208/96000000<1/1000$.
The margin $21/50-437/1045=19/10450$ pays this tail.

At $x=1/4$, $t>(157/50)(79/48)>103/20$,
$e^{1/4}<9/7$ gives $1+e^{-x}>16/9$, and
$e^{103/20}>20(73/10)(29/25)=4234/25$.
The leading bound is $67671/270976$, the full tail is below
$3456/(20^6\cdot7)<1/100000$, and the margin to $1/4$ is
$73/270976$.

At $x=9/20$,
$(157/50)S_3(9/10)=761921/100000>15/2$,
$1+e^{-9/20}>31/20$, and $e^{15/2}>20^2\cdot4=1600$.
The leading bound is $9/124$, the tail is below $7200/20^{10}<1/1000$,
and the margin to $3/40$ is $3/1240$.

For $\kappa(2/5)$, use $t>69/10$, $e^{1/5}<5/4$,
$e^{69/10}>900$.
The leading bound is $207/1000$, and the tail is below
$\frac{19044}{5\cdot900^4}\frac{10000}{9999}<10^{-6}$,
proving $\kappa(2/5)<1/4$.
The first inequality follows from $(157/50)S_3(4/5)>69/10$.

At $x=1/2$ take $t>17/2$, since $(157/50)S_5(1)>17/2$.
Use $e^{1/4}<9/7<4/3$ and $e^{17/2}>20^2\cdot7\cdot(79/48)$.
The leading bound is $952/13825$ and the tail is below
$36992/(3\cdot20^{11})<1/1000$.
The margin $7/100-952/13825=63/55300$ pays the tail.

For $\kappa(3/5)$,
$S_6(6/5)-520/157=88083/12265625>0$ and
$S_5(7/5)-4=1979/46875>0$ give $t>52/5$ and
$e^{52/5}>32000$; also $e^{3/10}<10/7$.
The leading bound is $(10/7)(9256/25)/32000$.
The full tail is below $19840/20^{13}$, by the majorants
$620j^4e^{-10j^2}$ and their ratio below $1/2$.
Their sum is below $1/50$.

At $x=7/10$, take $t>25/2$, using $e^{7/5}>4$.
Also $e^{7/20}<47/33<10/7$ and
$e^{25/2}>20^4(79/48)$.
The leading bound is $33/11060$ and the tail is below
$200000/(7\cdot20^{16})<10^{-6}$.
The margin to $3/1000$ is $9/553000$.

At $x=3/4$,
$(157/50)S_5(3/2)=896627/64000>14$,
$e^{3/8}<3/2$, and $e^{14}>20^4(73/10)=1168000$.
The leading bound is $21/23360$ and the tail is below
$37632/20^{18}<10^{-5}$.
The margin to $1/1000$ is $59/584000$.
All exponential comparisons above follow from the stated positive sums or their geometric remainder formula; no evaluated theta value enters.

\subsubsection{Gamma constants and all overlap attenuations}
Since $\rho(s)=e^{s/2}/(2\sinh s)$, the inequality
$\rho(s)\le1/(2s)+1/4$ is equivalent to
$2se^{s/2}\le(2+s)\sinh s$.
For odd degree $n$ coefficient comparison reduces to $n\le2^{n-1}$;
for even $n\ge2$ it reduces to $2^{2-n}\le1$.
Thus the power-series difference is nonnegative.
The lower bound $\rho({B_0})>5/6$ uses $e^{2/5}<3/2$, $e^{8/5}<5$;
the upper bound $\rho({B_0})<9/10$ uses $e^{2/5}>10/7$, $e^{8/5}>9/2$.

The harmonic error formula gives
\[
 H_n-\log n-\gamma_E
 =\sum_{k=n}^\infty\int_k^{k+1}\left(\frac1x-\frac1{k+1}\right)\,dx
 >\frac1{2(n+1)}.
\]
At $n=5$ and $\log5>8/5$ this yields $\gamma_E<3/5$.
Together with $\log\pi<6/5$, $\pi/2<11/7$ and $\log2<7/10$,
\[
 \cg<\frac65+\frac35+\frac{11}{7}+\frac{21}{10}
 =\frac{383}{70}<\frac{11}{2}.
\]

For the six overlap caps let $\sigma=25/2$.
The reflection exponents exceed $19/8,47/8,35/4$, from the stated logarithm lower bounds.
The positive exponential sums give
\begin{align*}
 e^{19/8}&>7(185/128)>10,\\
 e^{47/8}&>20\cdot7(289/128)>300,\\
 e^{35/4}&>400(73/10)(21/10)=6132>6000.
\end{align*}
For the translations, use
$\sqrt{3/4}<7/8$, $\sqrt{3/5}<4/5$, $\sqrt{4/5}<9/10$, and
\[
 28\cdot3^{12}<4^{12},\qquad
 2000\cdot3^{12}<5^{13},\qquad
 27\cdot4^{12}<2\cdot5^{12}.
\]
These give $(3/4)^\sigma<1/32$, $(3/5)^\sigma<1/500$ and
$(4/5)^\sigma<1/15$, paying every listed pair.


\subsection{Coercivity of the rank 64 lift}\label{c:sec:lift}
Let
\[
 e_0=(2{B_0})^{-1/2},\qquad e_n={B_0}^{-1/2}\cos(k_nx),\quad
 k_n=n\pi/{B_0}\quad(n\ge1),
\]
be the physical even Neumann basis. Let $P_{64}$ project onto
$e_0,\ldots,e_{63}$ and $H_{64}=L_I+9P_{64}$. Write $H=H_{64}$ in this appendix. We prove the lift estimate
without using the stored response data.

\subsubsection{The entire Gamma form dominates a diagonal form}
For $a>0$ the restricted whole-line Green kernel
$e^{-a\abs{x-y}}/(2a)$ is the inverse $R_{R,a}$ of
$-\partial_x^2+a^2$ with Robin form boundary contribution
$a(\abs{v(-{B_0})}^2+\abs{v({B_0})}^2)$. Removing this nonnegative term gives
the Neumann inverse $R_{N,a}$, so $R_{R,a}\preceq R_{N,a}$.
For every zero-extended local vector,
\[
 \int_0^\infty e^{-as}\norm{v(\cdot+s)-v}_2^2\dd s
 =\frac2a\norm v_2^2-2a\ip v{R_{R,a}v}
 \ge\frac2a\norm v_2^2-2a\ip v{R_{N,a}v}.
\]
Using $\rho(s)=\sum_{j\ge0}e^{-\lambda_js}$, where
$\lambda_j=2j+1/2$, and monotone summation of nonnegative forms, yields
\begin{equation}\label{c:eq:DN}
 t_\Gamma\succeq D_N,\qquad
 D_Ne_n=d_ne_n,\qquad
 d_n=2\sum_{j\ge0}\frac{k_n^2}{\lambda_j(\lambda_j^2+k_n^2)}.
\end{equation}
This is a form inequality on the whole local minimum domain. It does
not impose Neumann boundary conditions on an original vector. Lower
semicontinuity extends the core comparison, with inclusion in the
form domain of $D_N$.

The numbers $d_n$ are nonnegative and increasing. We give a short
rational proof that $d_{64}>9$. Set
\[
 f_t(a)=\frac{2t^2}{a(a^2+t^2)},\qquad G(t)=\sum_{j\ge0}f_t(2j+1/2).
\]
Each summand increases in $t$, and $k_{64}=80\pi>240$.
Differentiation gives
\[
 f_t''(a)=\frac{4t^2(6a^4+3a^2t^2+t^4)}{a^3(a^2+t^2)^3}>0.
\]
The mesh-two trapezoidal inequality for this convex decreasing function
therefore gives
\begin{align*}
 G(240)\ge{}&f_{240}(1/2)+f_{240}(5/2)+f_{240}(9/2)
 +\tfrac12f_{240}(13/2)\\
 &+\tfrac12\log\bigl(1+(480/13)^2\bigr),
\end{align*}
since $\int_A^\infty f_t(a)\dd a=\log(1+t^2/A^2)$.
The inequality $(1+x)^{-1}>1-x$ proves that the four rational terms
exceed, respectively,
\[
 \frac{3999}{1000},\qquad\frac{799}{1000},\qquad
 \frac{444}{1000},\qquad\frac{153}{1000}.
\]
For example they exceed $4-1/57600$, $4/5-5/57600$,
$4/9-1/6400$ and $2/13-13/115200$, respectively.
Also \eqref{c:eq:atanh} gives
\[
 \log2>\frac{842}{1215}>\frac{693}{1000},\qquad
 \log(15/13)>\frac17.
\]
It follows that
\[
 \tfrac12\log\bigl(1+(480/13)^2\bigr)
 >\log(480/13)=5\log2+\log(15/13)
 >\frac{693}{200}+\frac17.
\]
Consequently
\begin{equation}\label{c:eq:d64}
 d_{64}>G(240)>\frac{3151}{350}=9+\frac1{350}>9.
\end{equation}

\subsubsection{A whole operator bound for the local prime term}
Write $s=\log2$, $t=\log3$, $d={B_0}-s$ and
$a=s/\sqrt2$, $b=s/2$, $a_3=t/\sqrt3$.
The elementary log bounds give $.69<s<.70$, $1.09<t<1.10$ and
$\log5>2{B_0}$. The lower bound for $s$ follows from $842/1215>.693$;
for $t$, the first three terms in \eqref{c:eq:atanh} give
$\log3>1+1/12+1/80>1.09$.
Under the unitary fold on $(0,{B_0})$ partition that interval
(up to endpoints) into
\[
 C=(0,d),\quad L_-=(s-d,s),\quad L_+=(s,s+d),\quad
 J_0=(d,s-d).
\]
Prime 2 connects $y\in C$ to $s-y$ and $s+y$, with weight $a$,
and reflects $J_0$. Prime 4 connects the two leaves $L_-,L_+$
with weight $b$. Prime 3 reflects $S=(t-{B_0},{B_0})$ with weight $a_3$.
The inequalities
$2s>t$, $t+s>2{B_0}$, $t+2{B_0}<4s$ show that the leaves lie in $S$,
that their prime-3 images lie in $J_0$, and that $C$ does not meet $S$.

Define the positive number
\[
 R_p=\frac{b+a_3+\sqrt{(b+a_3)^2+8a^2}}2,
 \qquad\eta_0=2a/R_p,
 \qquad w_0=\eta_0 1_C+1_{J_0\cup L_-\cup L_+}.
\]
The weighted row ratio $(\PP_Iw_0)/w_0$ equals $R_p$ on the center
and leaves and is at most $a+a_3$ on $J_0$. Here
$\PP_I$ is the complete local prime operator in the definition of $L_I$.
One has $R_p>5/4$: its defining quadratic at $5/4$ is negative using
$b>1/3$, $a_3>4/7$, $2a^2>4/9$.
Moreover $a<1/2$, $a_3<121/190$ and hence $a+a_3<5/4$.
Every edge map is measure-preserving and every edge weight is
nonnegative. Applying
$2\abs{uv}\le(w_y/w_x)u^2+(w_x/w_y)v^2$ on each symmetric edge,
and summing, proves
$\abs{\ip v{\PP_Iv}}\le R_p\norm v_2^2$.
This retains all low/high interactions of the full operator.
Finally
\[
 b+a_3<\frac7{20}+\frac{121}{190}=\frac{75}{76},\quad
 2a^2<\frac{49}{100},\quad
 \left(\frac{27}{20}\right)^2-\frac{75}{76}\frac{27}{20}
 -\frac{49}{100}=\frac1{3800}>0.
\]
Thus
\begin{equation}\label{c:eq:primeop}
 \norm{\PP_I}<\frac{27}{20}.
\end{equation}

\subsubsection{The lift constant and a residual majorant}
The elementary estimate
\begin{equation}\label{c:eq:cgamma}
 c_\Gamma<\frac65+\frac35+\frac{11}{7}+\frac{21}{10}
 =\frac{383}{70}
\end{equation}
uses $e^{6/5}>401/125>22/7$ and $\gamma_E<3/5$.
To check the latter without evaluating Euler's constant, the harmonic
error identity gives
\[
 H_n-\log n-\gamma_E
 =\sum_{k=n}^\infty\int_k^{k+1}
       \left(\frac1x-\frac1{k+1}\right)\dd x
 >\frac1{2(n+1)}.
\]
Take $n=5$ and $\log5>8/5$ to get
$\gamma_E<H_5-8/5-1/12=3/5$.

\begin{theorem}[Rank-64 lift coercivity]\label{c:thm:lift}
On the complete mean-inclusive local minimum domain,
\begin{equation}\label{c:eq:lift}
 H_{64}=L_I+9P_{64}\succeq\frac{61}{28}I,
 \qquad H^{-1}\preceq\frac{28}{61}I.
\end{equation}
This is a coercivity statement for $H$, not for the unlifted $L_I$.
\end{theorem}
\begin{proof}
Dropping only the nonnegative pole in a lower estimate, equations
\eqref{c:eq:DN} and \eqref{c:eq:primeop} give
\[
 L_I\succeq D_N-(c_\Gamma+27/20)I.
\]
Below mode 64 the lifted diagonal is at least $9-c_\Gamma-27/20$;
from mode 64 onward the same bound follows from \eqref{c:eq:d64}.
Since $9-383/70-27/20=61/28$, \eqref{c:eq:lift} follows.
The pole has not been removed from the operator being inverted.
\end{proof}
For the residual contract define the larger diagonal majorant
\begin{equation}\label{c:eq:Dold}
 D_{\rm old}e_n=\beta_ne_n,\qquad
 \beta_n=\begin{cases}(d_n+1/2)^{-1},&n<64,\\
 (d_n-17/2)^{-1},&n\ge64.
 \end{cases}
\end{equation}
Translation Cauchy--Schwarz gives
$\norm{\PP_I}\le2(a_2+a_3+a_4)<3$, and
$c_\Gamma<11/2$, so $H\succeq D_N-17I/2+9P_{64}$.
Equation~\eqref{c:eq:d64} proves
$0<H^{-1}\preceq D_{\rm old}\preceq2I$.
Contract (a) uses this larger majorant. The sharper scalar bound for $H^{-1}$ is used separately in the inverse-cost estimate.

\subsection{An exact signed exterior prime identity}\label{c:sec:signed}
In this appendix write $A_{B_0}=A$.
This section is analytic and requires no stored numerical input or
exterior moment constraint. Set $g=u/\alpha$ and
$E_n=\{x:x,x+s_n\in O\}$. Define the nonnegative full-prime slack
\begin{equation}\label{c:eq:slack}
 D_p(g)=\sum_{n\ge2}a_n\int_{E_n}
 \alpha(x)\alpha(x+s_n)[g(x+s_n)-g(x)]^2\dd x.
\end{equation}
The sum includes all prime powers. The two row orientations arise from
the two diagonal terms in each edge square, as the proof below shows.

\subsubsection{The complete exterior row and its one reflection}
On the positive exterior fold $f=\sqrt2u$, the full prime expression is
\begin{equation}\label{c:eq:exteriorfold}
 (\PP_Of)(x)=\sum_{n\ge2}a_n\big[
 f(x+s_n)+1_{x>{B_0}+s_n}f(x-s_n)
 +1_{{B_0}<x<s_n-{B_0}}f(s_n-x)\big],\quad x>{B_0}.
\end{equation}
The first two terms are same-side orientations. The third is one
reflected opposite-side input, with one coefficient $a_n$. The cases
for $x-s_n$ are disjoint: positive exterior, killed interval and
negative exterior. Since $\log4<2{B_0}<\log5$, no reflection occurs for
$n=2,3,4$. The unitary fold already absorbs the factor two from the
two half-lines.

The relation $p(u,u)=\int_{B_0}^\infty f\PP_Of$ is understood as an
absolutely convergent form or weighted-duality pairing, first for
truncated sums. It is not an assertion that $\PP_Of$ belongs to
physical $L^2$ for every finite-$\wo$ vector. Equivalently it is the
$L^2$ pairing with the bounded conjugate
$M_\alpha\PP_OM_\alpha$, justified by the following row bound.
Define the nonnegative potential on all of $O$ by
\begin{equation}\label{c:eq:primepotential}
 V_p(x)=\alpha(x)\sum_{n\ge2}a_n
 \big[\alpha(x+s_n)1_O(x+s_n)
      +\alpha(x-s_n)1_O(x-s_n)\big].
\end{equation}
For $x>{B_0}$, either existing same-side product is at most
$A_{B_0}n^{-25/2}$: the sum of the two distances from ${B_0}$ is at least
$s_n$. On the reflected branch it is at most
$A_{B_0}e^{-(25/2)(s_n-2{B_0})}$, because those distances sum exactly to
$s_n-2{B_0}$. There is only one such branch in a row.

The two same-side budgets together satisfy
\[
 2\sum_{n\ge2}a_nn^{-25/2}\le2\sum_{n\ge2}n^{-12}
 \le2\left(2^{-12}+\int_2^\infty t^{-12}\dd t\right)
 =\frac{13}{22528}<\frac1{1000}.
\]
At $n=5$, equations \eqref{c:eq:coefficientcaps} and
$e^{-(25/2)(\log5-2{B_0})}<9/10$ give cost less than $261/400$.
There is no $n=6$ coefficient. Every remaining reflection is covered by
\eqref{c:eq:bothtails}, since
$a_ne^{-(25/2)(\log n-2{B_0})}\le(5/n)^{12}$ for $n\ge7$.
Consequently, for all $x\in O$,
\begin{equation}\label{c:eq:primecap}
 V_p(x)<\left(\frac1{1000}+\frac{261}{400}+\frac3{100}\right)A_{B_0}
 =\frac{1367}{2000}A_{B_0}<\frac7{10}A_{B_0}.
\end{equation}
Evenness extends the pointwise bound to the negative half-line without
doubling it. The summable majorants give uniform row convergence and
boundedness of the conjugated operator, including every prime tail.

\subsubsection{The signed identity and the analytic cap}
\begin{theorem}[Full-prime slack and mass bound]\label{c:thm:signed}
For every real even exterior $u$ of finite $\wo$, one has
\begin{equation}\label{c:eq:slackidentity}
 p(u,u)+D_p(u/\alpha)
 =\int_O V_p(x)|u(x)/\alpha(x)|^2\dd x,
\end{equation}
with $0\le D_p(u/\alpha)\le(7/5)A_{B_0}\wo(u)$.
Put
\[
 V(x)=c_\Gamma\alpha(x)^2+V_p(x),\qquad N=\frac{31}{46080}.
\]
Then $0<V(x)<N$ and
\begin{equation}\label{c:eq:negativeanalytic}
 c_\Gamma\norm u_2^2+p(u,u)
 \le N \wo(u)-D_p(u/\alpha).
\end{equation}
Whenever $t_\Gamma(u)$ is finite, the exact spatial identity is
\begin{equation}\label{c:eq:signedq0}
 q_0(u)=t_\Gamma(u)+D_p(u/\alpha)
              -\int_OV(x)|u(x)/\alpha(x)|^2\dd x.
\end{equation}
The cap is an analytic conclusion, not a numerical premise.
\end{theorem}
\begin{proof}
For a finite sum, expansion of one edge square in \eqref{c:eq:slack}
gives cross term $-2a_nu(x)u(x+s_n)$. The two diagonal terms,
after a change of variable in one, give the two orientations in
\eqref{c:eq:primepotential}. This proves \eqref{c:eq:slackidentity} for
truncated sums with exactly the factor in the definition of $p$.
The inequality $(a-b)^2\le2(a^2+b^2)$ gives
\[
 0\le D_p(g)\le2\int_OV_pg^2\le\frac75 A_{B_0}\norm g_2^2.
\]
The same row bound controls the absolute prime bilinear form by
weighted Cauchy--Schwarz. Monotone convergence of the slack series and
absolute convergence of the prime pairing extend the identity to all
$g\in L^2(O)$. Thus negative prime correlation increases the retained
slack; its sign has not been replaced by an absolute correlation.

Equation~\eqref{c:eq:cgamma} gives $c_\Gamma<11/2$. Together with
\eqref{c:eq:envelope} and \eqref{c:eq:primecap},
\[
 0<V(x)<\left(\frac{11}{2}+\frac7{10}\right)A_{B_0}
 =\frac{31}{5}A_{B_0}<\frac{31}{46080}.
\]
Adding the mass term proves \eqref{c:eq:negativeanalytic}, and
subtracting the exact identity from $t_\Gamma$ proves
\eqref{c:eq:signedq0}. The slack and spatial potential are $\wo$-bounded,
so the same identities extend in the original form topology, in
particular on the fixed minimum domain $\D_O^0$.
\end{proof}

The slack has trivial kernel on the even weighted exterior space.
Indeed its $n=2$ same-side edges have strictly positive weight, so
zero slack forces $g(x+\log2)=g(x)$ almost everywhere on $({B_0},\infty)$.
The $L^2$ norm is then an infinite sum of equal nonnegative period
integrals and can be finite only if each vanishes. Evenness supplies
the other half-line. This trivial-kernel statement does not assert a
uniform inequality $D_p\ge\delta\wo$ with $\delta>0$.


\subsection{Sharp cuts and the centered exterior minimum domain}\label{c:app:cuts}
We record the domain argument, rather than identifying a minimum
closure with an arbitrary maximal finite-energy class.
Write $W_\R(h)=\int_\R |h|^2/\alpha^2$. Under $h=\kappa f$ it equals
$\norm f_{L^2(\mu)}^2$. Proposition~\ref{a:prop:domain}
shows that the original topology and
$(t_\Gamma(h)+W_\R(h))^{1/2}$ are equivalent on numerators.

For a smooth numerator $h$, expansion gives
\[
 t_\Gamma(h1_I)+t_\Gamma(h1_O)
 =t_\Gamma(h)+2\ip{h1_I}{\J_\Gamma(h1_O)}.
\]
For $0<s\le2$, \eqref{c:eq:rhoelementary} gives
$\rho(s)\le1/(2s)+1/4\le1/s$. For $s\ge2$, the function
$se^{-s/2}$ decreases, so
$s\rho(s)\le(2/e)/(1-e^{-4})<1$; the last inequality follows already
from $e>8/3$. Thus the unweighted Carleman bound for the Gamma crossing
is finite: $\rho(s)\le1/s$ and the two boundary Carleman maps give,
for example, $\norm{\J_\Gamma}_{L^2(O)\to L^2(I)}\le2\pi$.
Thus the right-hand correction is bounded by a constant times
$\norm h_2^2$, and $W_\R$ splits exactly. Since $\alpha^2$ is bounded,
both cut maps extend continuously in the numerator form topology.

For a fixed smooth compact numerator, approximate either step cut from
its own side by smooth killed-core cutoffs. The error $e_\delta$ is
supported in strips of width $\delta$ at $\pm {B_0}$ and satisfies
\[
 \norm{e_\delta}_\infty+\operatorname{TV}(e_\delta)\le C,
 \qquad\norm{e_\delta}_2^2+W_\R(e_\delta)\le C\delta.
\]
The bounded-variation translation estimate gives
\[
 \norm{e_\delta(\cdot+s)-e_\delta}_2^2\le C\min(s,\delta).
\]
Integrating against $\rho(s)=O(1/s)$ near zero and its integrable
large-shift tail proves
$t_\Gamma(e_\delta)=O(\delta(1+|\log\delta|))\to0$.
Consequently the cuts are in the corresponding killed minimum closures.
Apply continuity to a general minimum-domain numerator to obtain the
same conclusion. On the bounded local interval $W_\R$ is equivalent
to physical $L^2$, because $\alpha$ is continuous and strictly positive.

The constant function belongs to $\D$. Indeed even smooth cutoffs
$\chi_R$ of the constant satisfy $\kappa\chi_R\to\kappa$ in physical
$H^1$ and in $W_\R$ by theta decay. The topology equivalence and
$t_\Gamma(h)\le C\norm h_{H^1}^2$ make them Cauchy in the original
energy topology. Subtracting a constant is an admissible gauge change,
preserving energy and variance.

For an exterior limit with zero moment, choose even compact smooth
$\psi$ in $O$ with $m_O(\psi)=1$. If $u_n$ are killed-core approximants,
replace them by $u_n-m_O(u_n)\psi$. Moment continuity makes the
correction vanish in form norm. Hence the centered cut lies in
precisely $\D_O^0$, not just in the moment kernel of a larger domain.


\section{Complete response quantities and contracts}\label{app:response}
These quantities specify the additional response branch. Their definitions and analytic error propagation do not constitute a new verification of the stated numerical enclosures.


\subsection{The exact centered-row quantities}\label{c:app:rows}
This appendix fixes the analytic interpretation of contract (b).
Use $K=64$, $J_0=512$, $d_0=6$ and $x_*=2$.
Let $\varphi_0={B_0}^{-1/2}$ and
$\varphi_m=(2/{B_0})^{1/2}\cos(k_mx)$ for $m>0$, where $k_m=m\pi/{B_0}$.
They are the unitary folds of the physical cosines.
All rows below are formed from the same matrix $Z$ and the same fixed rational vector $a$ in Hypothesis~\ref{c:hyp:response}.
The displayed formulas define quantities to be enclosed; their small numerical bounds remain assumptions.
For the physical columns $w_i$ of $\V$, let $Q_i=\J^*w_i$ denote
their weighted-dual kernels. This notation means
$\ip{\J u}{w_i}=\int_OuQ_i$ and
$\alpha Q_i\in L^2(O)$; it does not assert $Q_i\in L^2(O,\dd x)$.
If $m_O(u)=0$, then for every fixed vector $a$,
\begin{equation}\label{c:eq:centering}
 \V^*\J u=\int_Ou(x)[Q(x)-ac(x)]\dd x,\qquad
 \norm{\V^*\J u}\le\sqrt{\wo(u)}\norm{\alpha(Q-ac)}_{L^2(O;\R^{64})}.
\end{equation}
This is exact moment cancellation, not an approximate mean repair.
The choice of $a$ need not be accurate as a numerical minimizer.
The direction $c$ is subtracted in the weighted dual, not by physical
$L^2$ projection; $c$ is not a physical exterior $L^2$ vector.

Split the trial response into cosine rows below $K=64$ and rows
64 through 1023. Duality and Theorem~\ref{c:thm:source} give
\begin{equation}\label{c:eq:highrow}
 \norm{\alpha\J^*\V_{\ge64}}_{L^2(O;\R^{64})}
 \le j\norm{Z_{\ge64}}_{\rm HS}.
\end{equation}
This pays every discarded row of the finite trial response. The error
of the trial response as an approximation to $H_{64}^{-1}E$ belongs instead
to \eqref{c:eq:contracta}; it is not silently removed.

Let $b_m=\varphi_m({B_0})$ and $\lambda_j=2j+1/2$.
For $x>{B_0}$, direct cosine integration of the two folded exponential
kernels gives the low-response Gamma adjoint
\begin{equation}\label{c:eq:Qgamma}
 Q_{\Gamma,<K}(x)=\sum_{j\ge0}g_je^{-\lambda_j(x-{B_0})},\qquad
 g_j=\frac{1-e^{-2\lambda_jB_0}}{\sqrt2}
       \sum_{m<K}\frac{b_m\lambda_j}{\lambda_j^2+k_m^2}Z_m.
\end{equation}
The factor $1/\sqrt2$ converts folded coefficients to physical
full-line columns. Put $A_J=2J_0+1/2$, and define signed moments
\[
 M_r=\sum_{m<K}b_m(-k_m^2)^rZ_m\quad(0\le r<d_0),\qquad
 H_{d_0}=\sum_{m<K}|b_m|\norm{Z_m}k_m^{2d_0}.
\]
The exact finite geometric identity for
$\lambda/(\lambda^2+k^2)$ has remainder of modulus at most
$k^{2d_0}/\lambda^{2d_0+1}$, without requiring $k<\lambda$.
Therefore
\[
 \norm{\sum_{m<K}\frac{b_m\lambda}{\lambda^2+k_m^2}Z_m}
 \le\sum_{r<d_0}\frac{\norm{M_r}}{\lambda^{2r+1}}
       +\frac{H_{d_0}}{\lambda^{2d_0+1}}.
\]
With $\sigma_p(A)=A^{-p}+A^{1-p}/[2(p-1)]$, decreasing integral
comparison gives $\sum_{j\ge J_0}\lambda_j^{-p}\le\sigma_p(A_J)$.
Minkowski's inequality and
$\norm{e^{-\lambda h}}_{L^2(0,\infty)}=(2\lambda)^{-1/2}$ prove
\begin{equation}\label{c:eq:egamma}
 e_\Gamma=\sqrt{A/2}\left[
 \sum_{r<d_0}\norm{M_r}\sigma_{2r+3/2}(A_J)
 +H_{d_0}\sigma_{2d_0+3/2}(A_J)\right]
\end{equation}
as a bound for the complete omitted weighted Gamma norm. Both
exterior halves are included. The signed moments are summed before
their vector norms are taken.

The remaining finite kernel is
\begin{equation}\label{c:eq:K0}
 \begin{split}
 K_0(x)={}&\sum_{j<J_0}g_je^{-\lambda_j(x-{B_0})}-ac(x)\\
 &+\sum_{\substack{2\le n\le16\\|x-\log n|<{B_0}}}
  \frac{a_n}{\sqrt2}\sum_{m<K}|b_m|Z_m\cos(k_m(x-\log n)).
 \end{split}
\end{equation}
All ten prime powers through 16 are included:
$2,3,4,5,7,8,9,11,13,16$. Their Gamma, prime and mean contributions
are combined before vector norms. Boundary values at support switches
have measure zero.

The rough bound
\begin{equation}\label{c:eq:roughalpha}
 \alpha(x)^2\le80e^{4|x|-\pi e^{2|x|}}\quad(x\in O)
\end{equation}
follows already from the envelope: $4\pi^2<40$ and the tail-to-leading
coarse ratio $32e^{-3t}<1$ permit the constant 80.
Define
\[
 C_X=\sum_{j<J_0}\norm{g_j}
 +\left(\sum_{n\le16}a_n\right)
      \sum_{m<K}\frac{|b_m|\norm{Z_m}}{\sqrt2}+\norm a.
\]
For $x\ge x_*=2$, $\norm{K_0(x)}\le C_Xe^{x/2}$, including any finite
prime supports continuing past $x_*$. Since
$5-2\pi e^{2x}\le5-2\pi e^{2x_*}<0$, the full spatial tail is bounded by
\begin{equation}\label{c:eq:eX}
 e_X^2=\frac{160C_X^2e^{5x_*-\pi e^{2x_*}}}{2\pi e^{2x_*}-5}.
\end{equation}
The factor 160 includes both exterior half-lines.

On $[{B_0},2]$ the nine support switches, in increasing order, are
\[
 \begin{gathered}\log5-{B_0},\ \log7-{B_0},\ \log8-{B_0},\ \log9-{B_0},\ \log2+{B_0},\\
 \log11-{B_0},\ \log13-{B_0},\ \log3+{B_0},\ \log16-{B_0}.\end{gathered}
\]
The ten intervals have the fixed subdivision counts
\[
 (5,169,67,59,48,53,84,67,37,14),
\]
totaling 603.
On an exact coarse interval $[L_c,R_c]$ with count $N$, the exact finer
endpoints are $l_i=L_c+(R_c-L_c)i/N$, $r_i=L_c+(R_c-L_c)(i+1)/N$.
This algebraic identity proves adjacency and coverage. Overlap of
arbitrary interval boxes alone would not do so. Each open cell lies
between consecutive support switches and has a constant active-prime
set.

For a cell $[l,r]$, center $x_0$ and half-width $h$, define
$v_s=K_0^{(s)}(x_0)/s!$, $0\le s<16$. Taylor's theorem gives the
following explicit uniform vector remainder bounds:
\begin{align*}
 R_\Gamma&=\sum_{j<J_0}\norm{g_j}e^{-\lambda_j(l-{B_0})}
                   \frac{(\lambda_jh)^{16}}{16!},\\
 R_p&=\left(\sum_{n\ {\rm active}}a_n\right)
       \sum_{m<K}\frac{|b_m|\norm{Z_m}}{\sqrt2}
                   \frac{(k_mh)^{16}}{16!},\\
 R_c&=\norm a\cosh(r/2)\frac{(h/2)^{16}}{16!}.
\end{align*}
Thus $M_{[l,r]}=\sum_{s=0}^{15}\norm{v_s}h^s+R_\Gamma+R_p+R_c$
bounds $\norm{K_0(x)}$ on the whole cell, and
\begin{equation}\label{c:eq:Ifin}
 I_{\rm fin}=\sum_{[l,r]}2(r-l)A(l)M_{[l,r]}^2
\end{equation}
bounds its full even weighted integral up to $x_*$.
Contract (b) requires outward enclosures of these
quantities; the analytical formula is a supremum-and-remainder bound,
not a quadrature accuracy inferred from mesh size. The subdivision counts do not replace these exact endpoint definitions.

\subsubsection{The complete prime tail in the response budget}
For $y\in I$, $|y\pm\log n|\ge\log n-{B_0}$ and each of the two
translated restrictions has weighted norm at most
$\sqrt{80}e^{-2{B_0}}n^2e^{-\zeta n^2}$, with
$\zeta=(\pi/2)e^{-2{B_0}}>3/10$. For $N\ge16$, integration by parts and
$\int_N^\infty e^{-\zeta x^2}\dd x\le
 e^{-\zeta N^2}/(2\zeta N)$ give
\[
 \norm{\J_{p,>N}}_{\wo\to L^2(I)}
 \le2\sqrt{80}e^{-2{B_0}}
 \left(\frac{N}{2\zeta}+\frac1{4\zeta^2N}\right)e^{-\zeta N^2}.
\]
At $N=16$, $2\sqrt{80}e^{-2{B_0}}<18$, the parenthesis is less than 27,
and $\zeta N^2>75$; hence the bound is smaller than
$b_{16}=500/20^{25}$. These entire integer tails dominate all omitted
prime powers and retain both orientations.


The finite part has weighted norm at most $\sqrt{I_{\rm fin}}+e_X$,
the complete omitted Gamma part has norm at most $e_\Gamma$,
and the complete omitted prime part has norm at most
$b_{16}\norm{Z_{<64}}_{\mathrm{HS}}$.
The high trial rows have norm at most $j\norm{Z_{64:1023}}_{\mathrm{HS}}$ by
\eqref{c:eq:highrow}. Exact exterior centering therefore gives
\[
 \norm{V^*Ju}\le
 \left(\sqrt{I_{\rm fin}}+e_X+e_\Gamma+
 b_{16}\norm{Z_{<64}}_{\mathrm{HS}}
 +j\norm{Z_{64:1023}}_{\mathrm{HS}}\right)\sqrt{\wo(u)}.
\]
This proves the decomposition used in \eqref{c:eq:responsebound};
the numerical caps themselves remain contract (b).


\section{Actual weight source recovery and the Abel endpoint}\label{app:abel}
This appendix proves the weighted estimates invoked by the synthesis theorem. In particular, the source is recovered in the actual weight before the endpoint conclusion is drawn. Carrier growth alone is not used as a weighted admission criterion.


\subsection{A scalar Borel source and exact theta reduction}\label{d:sec:scalar-source}

The coarse estimate~\eqref{d:eq:Abel-norm} uses quotient growth but not the full strength of actual weighted admission of $u$. We now use that admission to recover an auxiliary Laplace source in its correct weighted norm. This is the additional input that makes the Abel bound uniform.

The exponential coordinate $q=e^{2b}$ maps $L$ bijectively onto
\begin{equation}\label{d:eq:q-domain}
 \mathcal D_q=\{q:\operatorname{Re}q\ge1/2,\ |q-1|\le1\}.
\end{equation}
Indeed the two lens inequalities are
$\operatorname{Re}(e^{2b})\ge1/2$ and
$\operatorname{Re}(e^{-2b})\ge1/2$. The latter becomes
$\operatorname{Re}q/|q|^2\ge1/2$, equivalent to $|q-1|\le1$.
Conversely $\mathcal D_q$ lies in the right half-plane, and
$b=\frac12\operatorname{Log}q$ satisfies the two inequalities in the principal theta strip. Fractional powers of $q$ below always use this principal branch.

Fix once and for all a positively oriented, centrally symmetric rectifiable Jordan contour $\Gamma\subset\mathbb C\setminus L$ with $L\subset\operatorname{int}\Gamma$, winding number one about every point of $L$, with its interior in $|\operatorname{Im}b|<\pi/4$, and such that
\begin{equation}\label{d:eq:fixed-contour-margin}
 \operatorname{Re}(e^{2b})\ge\frac14\qquad(b\in\Gamma).
\end{equation}
A sufficiently small exterior parallel contour has these properties by compactness and the strict margin between $L$ and the strip boundary. Symmetry also gives the same margin for $e^{-2b}$. The contour remains fixed independently of $u$ and $r$.

For $u\in\Ii$, with $P(w)=Q(-iw)$ and Borel transform $B_P$, define the scalar analytic functional
\begin{equation}\label{d:eq:mu-functional}
 \mu(h)=\frac1{2\pi i}\int_\Gamma B_P(b)h(b)\,db.
\end{equation}
All functions $h$ used below are holomorphic on a neighborhood of the contour and its interior. Theorem~\ref{d:thm:Borel} gives
\begin{equation}\label{d:eq:mu-bound}
 |\mu(h)|\le C_\Gamma\norm{u}_{\Hh}\sup_{b\in\Gamma}|h(b)|.
\end{equation}
The Borel series of an even $P$ is odd, and continuation gives
$B_P(-b)=-B_P(b)$. A change of variable on the symmetric contour therefore yields
\begin{equation}\label{d:eq:mu-symmetry}
 \mu(h(-b))=\mu(h(b)),\qquad \mu(e^{bw})=P(w).
\end{equation}

\begin{lemma}[Scalar source identification]\label{d:lem:scalar-source}
The function
\begin{equation}\label{d:eq:scalar-v}
 v(x)=\mu(\kappa(x+b))
\end{equation}
is even, belongs to $L^1(\R,e^{B|x|}dx)$ for every finite $B$, and equals $u$ almost everywhere.
\end{lemma}
\begin{proof}
The contour lies in a compact subset of the theta strip. The two-ended theta bounds therefore give uniform ordinary exponential moments for $\kappa(\cdot+b)$, $b\in\Gamma$. Equation~\eqref{d:eq:mu-bound} and ordinary dominated integration prove the asserted $L^1$ properties. Symmetry~\eqref{d:eq:mu-symmetry} and evenness of $\kappa$ give evenness of $v$.

For each $b$ on the contour, a shift within the strip gives
\[
 \int_\R\kappa(x+b)\cos(zx)\,dx=\Xi(z)\cos(bz).
\]
The vertical sides vanish by the theta tails. Uniform exponential moments justify passing the scalar contour through the Fourier integral, so
\[
 F_v(z)=\Xi(z)\mu(\cos(bz))=\Xi(z)P(iz)=\Xi(z)Q(z)=F_u(z).
\]
The given $u\in\Hh$ also lies in ordinary $L^1$ by~\eqref{d:eq:eval-bound}. Fourier uniqueness there gives $v=u$ almost everywhere.
\end{proof}

The integral~\eqref{d:eq:scalar-v} is a scalar contour integral, not an $\Hh$-valued Bochner integral over an exterior or boundary contour. No such weighted contour admission is being assumed.

Define for $t\ge0$
\begin{equation}\label{d:eq:fg-source}
 f(t)=\mu(q^{9/4}e^{-tq}),\qquad
 g(t)=\mu(q^{5/4}e^{-tq}),\qquad q=e^{2b}.
\end{equation}
These are continuous, and the fixed contour gives
\begin{equation}\label{d:eq:fg-crude}
 |f(t)|+|g(t)|\le C\norm{u}_{\Hh}e^{-t/4}.
\end{equation}
Ordinary integration on the contour, justified by~\eqref{d:eq:fixed-contour-margin}, gives
\begin{equation}\label{d:eq:g-tail}
 g(t)=\int_t^\infty f(s)\,ds.
\end{equation}

For $x\ge0$, put
\begin{equation}\label{d:eq:t-U-H}
 t=\pi e^{2x},\qquad
 U(t)=\frac{u(x)}{4\pi^2e^{9x/2}},\qquad
 H_f(t)=\sum_{n\ge1}n^4f(n^2t).
\end{equation}
Here $u$ is represented by the continuous scalar source of Lemma~\ref{d:lem:scalar-source}. The precise theta identity on the strip is
\begin{align}
 \kappa(x+b)
 &=4\pi^2e^{9x/2}\sum_{n\ge1}n^4q^{9/4}e^{-n^2tq}\notag\\
 &\quad-6\pi e^{5x/2}\sum_{n\ge1}n^2q^{5/4}e^{-n^2tq}.
 \label{d:eq:exact-theta-q}
\end{align}
It follows by expanding~\eqref{a:eq:kappa}, with analytic continuation as in Lemma~\ref{d:lem:theta-tails}. The uniform exponential bound on the fixed contour justifies all sums and integrals.

\begin{proposition}[Exact reduction in the actual weight]\label{d:prop:exact-reduction}
For $t\ge\pi$,
\begin{equation}\label{d:eq:U-H-identity}
 U(t)=H_f(t)-\frac3{2t}\int_t^\infty H_f(s)\,ds.
\end{equation}
Writing
\begin{equation}\label{d:eq:E-tail}
 E_T=L^2([T,\infty),t^{3/2}e^t\,dt),\qquad T\ge\pi,
\end{equation}
one has two-sided norm comparability
\begin{equation}\label{d:eq:actual-weight-comparison}
 \norm{u}_{\Hh}\asymp\norm{U}_{E_\pi}.
\end{equation}
The comparison holds as an identity of extended finite or infinite integrals for every even scalar source with the displayed normalization.
\end{proposition}
\begin{proof}
Insert~\eqref{d:eq:exact-theta-q} into~\eqref{d:eq:mu-functional} and divide by the normalization in~\eqref{d:eq:t-U-H}. The lower-degree coefficient is exactly
$6\pi e^{5x/2}/(4\pi^2e^{9x/2})=3/(2t)$. Thus
\[
 U(t)=\sum_{n\ge1}n^4f(n^2t)
       -\frac3{2t}\sum_{n\ge1}n^2g(n^2t).
\]
By~\eqref{d:eq:fg-crude} and~\eqref{d:eq:g-tail}, the latter sum equals
$\int_t^\infty H_f(s)\,ds$, proving~\eqref{d:eq:U-H-identity}. No theta remainder has been discarded.

Finally use~\eqref{d:eq:W-two-sided}, $u(x)=4\pi^2e^{9x/2}U(t)$, and $dx=dt/(2t)$. The squared source normalization contributes $t^{9/2}$, the weight contributes $t^{-2}e^t$, and the change of variables contributes $t^{-1}$. Their product is $t^{3/2}e^t$. Evenness covers both physical half-lines and proves~\eqref{d:eq:actual-weight-comparison}.
\end{proof}

\subsection{Noncircular recovery of the Laplace source}\label{d:sec:source-recovery}

The contour bound alone does not imply the weighted membership needed for the endpoint. The actual membership $u\in\Hh$, through~\eqref{d:eq:actual-weight-comparison}, supplies it.

\begin{theorem}[Recovery from actual weighted admission]\label{d:thm:recovery}
For every $u\in\Ii$, the function $f$ in~\eqref{d:eq:fg-source} belongs to
\begin{equation}\label{d:eq:E-global}
 E=L^2((0,\infty),(1+t)^{3/2}e^t\,dt),
\end{equation}
and
\begin{equation}\label{d:eq:recovery-bound}
 \norm{f}_E\le C\norm{u}_{\Hh}.
\end{equation}
\end{theorem}
\begin{proof}
First recover $H_f$ without assuming it belongs to $E_\pi$. The crude bound~\eqref{d:eq:fg-crude} gives
$H_f(t)=O(\norm{u}_{\Hh}e^{-t/4})$ for $t\ge\pi$. Set
$J_0(t)=\int_t^\infty H_f(s)\,ds$. Equation~\eqref{d:eq:U-H-identity} becomes
\[
 J_0'(t)+\frac3{2t}J_0(t)=-U(t).
\]
The known exponential bound already gives $t^{3/2}J_0(t)\to0$. Multiplying by $t^{3/2}$ and integrating to infinity yields the pointwise identities
\begin{align}
 J_0(t)&=t^{-3/2}\int_t^\infty s^{3/2}U(s)\,ds,\notag\\
 H_f(t)&=U(t)+\frac32t^{-5/2}\int_t^\infty s^{3/2}U(s)\,ds.
 \label{d:eq:Volterra-recovery}
\end{align}
This derivation has not used weighted membership of $H_f$ or $f$.

The second operator in~\eqref{d:eq:Volterra-recovery} is bounded on $E_\pi$. Conjugate by the isometry
$V(s)=s^{3/4}e^{s/2}U(s)$ to ordinary $L^2$. Apart from its constant $3/2$, the resulting kernel is
\[
 \boldsymbol1_{s\ge t}\,t^{-1}(s/t)^{3/4}e^{-(s-t)/2}.
\]
For $t\ge\pi$ it is bounded by the integrable function
$\pi^{-1}(1+(s-t)/\pi)^{3/4}e^{-(s-t)/2}$ of $s-t\ge0$.
The elementary $L^1$ convolution bound, after extension by zero, proves boundedness. Therefore
\begin{equation}\label{d:eq:Hf-weighted}
 \norm{H_f}_{E_\pi}\le C\norm{U}_{E_\pi}
 \le C'\norm{u}_{\Hh}.
\end{equation}

To remove the higher theta terms, define
\[
 (Sh)(t)=\sum_{n\ge2}n^4h(n^2t).
\]
On $E_T$, a change of variables gives
\begin{equation}\label{d:eq:dilation-E}
 \norm{n^4h(n^2\,\cdot)}_{E_T}
 \le n^{3/2}e^{-(n^2-1)T/2}\norm{h}_{E_T}.
\end{equation}
Indeed the squared polynomial prefactor is $n^8n^{-3}n^{-2}=n^3$, and on $s\ge n^2T$ the exponential ratio
$e^{s/n^2-s}$ is at most $e^{-(n^2-1)T}$. Hence the norm of $S$ on $E_T$ is less than one for all sufficiently large fixed $T$.

One must not apply this inverse to $f$ before proving $f\in E_T$. Use simultaneously the Banach space
\[
 \mathcal B_T=\{h\in C([T,\infty)):
       \norm{h}_{\mathcal B_T}:=\sup_{t\ge T}e^{t/4}|h(t)|<\infty\}.
\]
There
\begin{equation}\label{d:eq:dilation-B}
 \norm{S}_{\mathcal B_T\to\mathcal B_T}
 \le\sum_{n\ge2}n^4e^{-(n^2-1)T/4}.
\end{equation}
The sums in~\eqref{d:eq:dilation-E} and~\eqref{d:eq:dilation-B} tend to zero as $T$ tends to infinity, by domination by any fixed convergent Gaussian series. Choose one $T$ for which both operator norms are less than one. This is an analytic choice; no numerical tail calculation is needed.

The known contour bound places $f$ and $H_f=(I+S)f$ in $\mathcal B_T$, while~\eqref{d:eq:Hf-weighted} places $H_f$ in $E_T$. The same finite Neumann sums
\[
 f_M=\sum_{k=0}^M(-S)^kH_f
\]
converge to $f$ in $\mathcal B_T$ and to a weighted inverse in $E_T$. The first convergence is uniform on compact subintervals, while the second implies local $L^2$ convergence. Their limits therefore agree almost everywhere. Consequently $f\in E_T$, with
\[
 \norm{f}_{E_T}\le
 (1-\norm{S}_{E_T\to E_T})^{-1}\norm{H_f}_{E_T}
 \le C\norm{u}_{\Hh}.
\]
Finally~\eqref{d:eq:fg-crude} controls the fixed interval $[0,T]$ with a constant times $\norm{u}_{\Hh}$, and the weights $t^{3/2}$ and $(1+t)^{3/2}$ are comparable on $[T,\infty)$. This proves~\eqref{d:eq:recovery-bound}.
\end{proof}

\begin{lemma}[Forward weighted reconstruction]\label{d:lem:forward-reconstruction}
Suppose an even scalar theta source has the exact identity~\eqref{d:eq:U-H-identity} with $H_f=(I+S)f$, and $f\in E_\pi$. Then its physical numerator is admitted to $\Hh$, with
\begin{equation}\label{d:eq:forward-reconstruction}
 \norm{u}_{\Hh}\le C\norm{f}_{E_\pi}.
\end{equation}
\end{lemma}
\begin{proof}
The dilation estimate~\eqref{d:eq:dilation-E} at $T=\pi$ proves boundedness of $I+S$, without requiring contraction. The tail operator
$Jh(t)=t^{-1}\int_t^\infty h(s)\,ds$ is bounded on $E_\pi$: after the same isometry, its kernel is
\[
 \boldsymbol1_{s\ge t}\,t^{-1}(t/s)^{3/4}e^{-(s-t)/2}
 \le\pi^{-1}e^{-(s-t)/2}.
\]
Its norm is at most $2/\pi$. Thus~\eqref{d:eq:U-H-identity} bounds $U$ in $E_\pi$, and the extended-integral comparison~\eqref{d:eq:actual-weight-comparison} proves admission and~\eqref{d:eq:forward-reconstruction}. The scalar source is already defined; the argument does not assume the weighted conclusion.
\end{proof}

\subsection{Construction of the positive fractional power kernel}\label{d:sec:positive-kernel}

For the endpoint analysis restrict
\begin{equation}\label{d:eq:kernel-parameters}
 \frac34\le r<1,\qquad
 \eta=\eta_r=\frac94(1-r),\qquad
 c_0=\frac12,\qquad d_r=c_0^{1/r}.
\end{equation}
Thus $0<\eta\le9/16$ and $2^{-4/3}\le d_r\le c_0$. The symbol $a=i/2$ remains the centering point and is not reused as a fractional-order parameter.

Let $I_\eta$ denote the positive measure
$s^{\eta-1}ds/\Gamma(\eta)$ on $(0,\infty)$, with $I_0=\delta_0$ if the limiting parameter is used. Its Laplace transform is $q^{-\eta}$ in the right half-plane~\cite[Eq.~5.9.1]{DLMF}.

\begin{lemma}[Positive probability and convolution kernels]\label{d:lem:positive-kernel}
For $0<r<1$ and $t\ge0$ there is a positive probability measure $\rho_{r,t}$ on $[0,\infty)$ with
\begin{equation}\label{d:eq:rho-transform}
 \int_0^\infty e^{-qs}\rho_{r,t}(ds)=e^{-tq^r},
 \qquad \operatorname{Re}q>0.
\end{equation}
The family is weakly continuous in $t$. The positive locally finite kernel
\begin{equation}\label{d:eq:k-definition}
 k_r(t,ds)=(I_\eta*\rho_{r,t})(ds)
\end{equation}
is measurable in $t$ and satisfies
\begin{equation}\label{d:eq:k-Laplace}
 \int_0^\infty e^{-qs}k_r(t,ds)=q^{-\eta}e^{-tq^r}.
\end{equation}
\end{lemma}
\begin{proof}
The Gamma integral and integration by parts give, first for positive $q$,
\begin{equation}\label{d:eq:Levy-power}
 q^r=\frac r{\Gamma(1-r)}
      \int_0^\infty(1-e^{-qs})s^{-1-r}\,ds.
\end{equation}
Both endpoints converge for $0<r<1$. For $\epsilon>0$ let
\[
 \nu_\epsilon(ds)=\frac r{\Gamma(1-r)}
           \boldsymbol1_{s\ge\epsilon}s^{-1-r}\,ds.
\]
This measure is finite. Its compound-Poisson probability measure
\[
 e^{-t\nu_\epsilon([0,\infty))}
 \sum_{j=0}^\infty\frac{t^j}{j!}\nu_\epsilon^{*j},
 \qquad \nu_\epsilon^{*0}=\delta_0,
\]
has Laplace transform
$\exp[-t\int(1-e^{-qs})\nu_\epsilon(ds)]$.

These probabilities are tight as $\epsilon\downarrow0$. If $L_\epsilon$ denotes their Laplace transform, their mass above $R>0$ is at most
\[
 (1-e^{-1})^{-1}(1-L_\epsilon(1/R))
 \le C tR^{-r},
\]
because the truncated exponent is at most $R^{-r}$ by~\eqref{d:eq:Levy-power}. Every weak subsequential limit therefore has transform $e^{-tq^r}$ for positive $q$. Uniqueness of Laplace transforms identifies all such limits. For clarity, uniqueness for finite measures follows by the change of variable $y=e^{-s}$ and polynomial approximation on $[0,1]$; for positive locally finite measures with a finite exponential tilt, first apply that tilt. This also gives the uniqueness used below.

The same tightness argument, locally uniformly in $t$, and continuity of $e^{-tq^r}$ prove weak continuity in $t$. The right-half-plane identity~\eqref{d:eq:rho-transform} follows by holomorphy and the identity theorem. Convolution with $I_\eta$ gives~\eqref{d:eq:k-Laplace}. Its finite exponentially weighted mass proves local finiteness. Weak continuity of the probability kernel and integration against the locally finite gamma measure give measurability of the convolution kernel, first on bounded continuous compactly supported tests and then on Borel sets. Positivity permits all these convolution integrations by Tonelli.
\end{proof}

\begin{proposition}[Fractional contraction of the actual Laplace source]\label{d:prop:fractional-source}
For the $f$ of~\eqref{d:eq:fg-source}, define
\begin{equation}\label{d:eq:fr-definition}
 f_r(t)=\mu(q^{9r/4}e^{-tq^r}).
\end{equation}
Then
\begin{equation}\label{d:eq:fr-kernel}
 f_r(t)=\int_0^\infty f(s)k_r(t,ds).
\end{equation}
The integral is absolutely convergent. This identity applies to arbitrary complex sources, without positivity of $f$ or $\mu$.
\end{proposition}
\begin{proof}
The contour lies in $\operatorname{Re}q\ge1/4$, so
$|q^{9/4}e^{-sq}|\le C e^{-s/4}$ there. The integral of this bound against $k_r(t,ds)$ is finite by~\eqref{d:eq:k-Laplace} at $q=1/4$. Thus absolute Fubini permits interchange with the scalar contour. The result is
\[
 \mu\left(q^{9/4}\int_0^\infty e^{-sq}k_r(t,ds)\right)
 =\mu(q^{9/4-\eta}e^{-tq^r})=f_r(t).
\]
The branches agree because $\operatorname{Log}q=2b$ in the strip, so $q^r=e^{2rb}$ and $9/4-\eta=9r/4$ exactly.
\end{proof}

\subsection{A positive resolvent and uniform weighted Schur control}\label{d:sec:weighted-Schur}

We require a bound uniform as $r\uparrow1$, including small values of the input variable. It can be obtained without a pointwise density formula for the probability measures in Lemma~\ref{d:lem:positive-kernel}.

\begin{lemma}[Keyhole resolvent formula]\label{d:lem:keyhole}
For the parameters~\eqref{d:eq:kernel-parameters} and real $q>d_r$,
\begin{equation}\label{d:eq:resolvent-Stieltjes}
 \frac1{q^r-c_0}
 =\frac{d_r^{1-r}/r}{q-d_r}
   +\int_0^\infty\frac{\varrho_r(v)}{q+v}\,dv,
\end{equation}
where
\begin{equation}\label{d:eq:cut-density}
 \varrho_r(v)=\frac{\sin(\pi r)}{\pi}
 \frac{v^r}{v^{2r}-2c_0v^r\cos(\pi r)+c_0^2}\ge0.
\end{equation}
The positive function
\begin{equation}\label{d:eq:pr-positive}
 p_r(s)=\frac{d_r^{1-r}}r e^{d_rs}
             +\int_0^\infty e^{-vs}\varrho_r(v)\,dv
\end{equation}
has Laplace transform $(q^r-c_0)^{-1}$ for $q>d_r$, and satisfies uniformly for $s>0$ and $3/4\le r<1$,
\begin{equation}\label{d:eq:pr-bound}
 p_r(s)\le C e^{d_rs}(1+s^{-1/4}).
\end{equation}
\end{lemma}
\begin{proof}
Use the principal branch of $z^r$ on $\C\setminus(-\infty,0]$. The function $(z^r-c_0)^{-1}$ has one simple pole there at $z=d_r$, of residue $d_r^{1-r}/r$. Apply a keyhole Cauchy contour to
$(z^r-c_0)^{-1}/(z-q)$. The large-circle integral tends to zero as $O(R^{-r})$, while the small circle tends to zero because $(z^r-c_0)^{-1}$ remains bounded at zero. On the upper bank of the negative axis its negative imaginary part is
\[
 \frac{v^r\sin(\pi r)}
 {v^{2r}-2c_0v^r\cos(\pi r)+c_0^2}.
\]
The lower bank is its conjugate. Their oriented difference and the simple pole give~\eqref{d:eq:resolvent-Stieltjes}, including its positive sign. The cut integral converges at both endpoints. Integrating the nonnegative expression~\eqref{d:eq:pr-positive} against $e^{-qs}$ and using Tonelli proves its claimed Laplace transform.

Since $\cos(\pi r)\le0$ in the specified range, the denominator in~\eqref{d:eq:cut-density} is at least $v^{2r}+c_0^2$. For $0<v\le1$, $\varrho_r(v)\le C v^r$; for $v\ge1$,
$\varrho_r(v)\le C\sin(\pi r)v^{-r}$. Thus the cut term in~\eqref{d:eq:pr-positive} is bounded by
\[
 C+C\sin(\pi r)\Gamma(1-r)s^{r-1}
 \le C'(1+s^{-1/4}).
\]
The last step uses
$\sin(\pi r)\Gamma(1-r)=\pi/\Gamma(r)$, uniformly bounded on $[3/4,1]$, by the Gamma reflection formula~\cite[Eq.~5.5.3]{DLMF}. For $s\le1$, $s^{r-1}\le s^{-1/4}$, and for $s\ge1$ it is at most one. The residue coefficient is uniformly bounded and $e^{d_rs}\ge1$, proving~\eqref{d:eq:pr-bound}.
\end{proof}

\begin{lemma}[Integrated potentials and their densities]\label{d:lem:potentials}
Uniformly in the same parameter range,
\begin{equation}\label{d:eq:cubic-potential}
 p_r^{*3}(s)\le C e^{d_rs}(s^{5/4}+s^2).
\end{equation}
For every integer $j\ge0$, the positive locally finite measure
\[
 M_j(ds)=\int_0^\infty t^je^{c_0t}k_r(t,ds)\,dt
\]
has the density
\begin{equation}\label{d:eq:potential-identity}
 M_j(ds)=j!\,[I_\eta*p_r^{*(j+1)}](s)\,ds.
\end{equation}
In particular these measures have no atom at zero.
\end{lemma}
\begin{proof}
For~\eqref{d:eq:cubic-potential}, expand the triple convolution of the bound
$C e^{d_rs}(1+s^{-1/4})$. If $k$ of the three factors are singular, the beta integral gives a constant times $e^{d_rs}s^{2-k/4}$, for $0\le k\le3$. The beta parameters are positive; explicitly
\[
 \int_0^s y^{\alpha-1}(s-y)^{\beta-1}\,dy
 =s^{\alpha+\beta-1}
       \frac{\Gamma(\alpha)\Gamma(\beta)}{\Gamma(\alpha+\beta)}
 \quad(\alpha,\beta>0),
\]
as in~\cite[Eq.~5.12.1]{DLMF}. The intermediate powers are bounded by a constant times $s^{5/4}+s^2$.

For real $q>d_r$, Tonelli and~\eqref{d:eq:k-Laplace} give
\[
 \int_0^\infty e^{-qs}M_j(ds)
 =q^{-\eta}\int_0^\infty t^je^{-t(q^r-c_0)}\,dt
 =\frac{j!\,q^{-\eta}}{(q^r-c_0)^{j+1}}.
\]
This is finite, so $M_j$ is locally finite. The right side is also the Laplace transform of the measure on the right of~\eqref{d:eq:potential-identity}, by Lemma~\ref{d:lem:keyhole}. Laplace uniqueness after a fixed finite exponential tilt proves the measure identity. The convolution on its right is locally integrable and is a density, proving the no-atom assertion.
\end{proof}

We will also use, for $\beta\ge0$,
\begin{equation}\label{d:eq:gamma-convolution-bound}
 [I_\eta*(e^{d_r\,\cdot}\,\cdot^\beta)](s)
 \le e^{d_rs}s^\beta\int_0^\infty e^{-d_ry}I_\eta(dy)
 =d_r^{-\eta}e^{d_rs}s^\beta.
\end{equation}
The factor $d_r^{-\eta}$ is uniformly bounded under~\eqref{d:eq:kernel-parameters}. This formulation preserves the cancellation between the singular gamma integral and $1/\Gamma(\eta)$ as $\eta\downarrow0$.

\begin{theorem}[Uniform weighted kernel estimate]\label{d:thm:weighted-kernel}
There is a constant independent of $3/4\le r<1$ such that
\begin{align}
 &\int_\pi^\infty(1+t)^{3/2}e^t
       \left|\int_0^\infty f(s)k_r(t,ds)\right|^2dt\notag\\
 &\hspace{35mm}\le C\int_0^\infty(1+s)^{3/2}e^s|f(s)|^2\,ds.
 \label{d:eq:uniform-kernel}
\end{align}
For general $E$ elements the integral is interpreted by the bounded extension; for the continuous source in~\eqref{d:eq:fg-source} it agrees with the absolutely convergent integral~\eqref{d:eq:fr-kernel}.
\end{theorem}
\begin{proof}
Set
\[
 K_r(t,ds)=e^{c_0t-c_0s}k_r(t,ds).
\]
Its mass and first moment follow from~\eqref{d:eq:k-Laplace} and one derivative at $q=c_0$:
\begin{align}
 \int K_r(t,ds)&=c_0^{-\eta}e^{-t(c_0^r-c_0)},\notag\\
 \int sK_r(t,ds)&=
 c_0^{-\eta}e^{-t(c_0^r-c_0)}
        [\eta/c_0+trc_0^{r-1}].
 \label{d:eq:kernel-moments}
\end{align}
The derivative is legitimate by exponential integrability. Since $c_0^r\ge c_0$ and the remaining parameters stay in bounded ranges, Cauchy--Schwarz for the positive finite measure gives
\begin{equation}\label{d:eq:Schur-row}
 \int(1+s)^{1/2}K_r(t,ds)\le C(1+t)^{1/2}.
\end{equation}

The required column estimate is a measure domination on the whole input half-line:
\begin{equation}\label{d:eq:Schur-column}
 \int_\pi^\infty(1+t)^2K_r(t,ds)\,dt
 \le C(1+s)^2\,ds.
\end{equation}
To prove it, use $(1+t)^2\le Ct^2$ for $t\ge\pi$, and then the $j=2$ case of~\eqref{d:eq:potential-identity}. By~\eqref{d:eq:cubic-potential} and~\eqref{d:eq:gamma-convolution-bound}, the resulting density is bounded by
\[
 C e^{-c_0s}[I_\eta*p_r^{*3}](s)
 \le C'e^{-(c_0-d_r)s}(s^{5/4}+s^2)
 \le C''(1+s)^2.
\]
This holds also for $0<s<1$, and the potential measure has no atom at zero. Thus no small-input region has been discarded. The restriction $t\ge\pi$ matches the complete physical half-line $x\ge0$, not an omitted physical tail.

Write $G(s)=e^{c_0s}f(s)$ and $G_r(t)=e^{c_0t}f_r(t)$. Weighted Cauchy--Schwarz gives
\[
 |G_r(t)|^2
 \le\left[\int(1+s)^{1/2}K_r(t,ds)\right]
       \left[\int |G(s)|^2(1+s)^{-1/2}K_r(t,ds)\right].
\]
Multiply by $(1+t)^{3/2}$, apply~\eqref{d:eq:Schur-row}, integrate over $t\ge\pi$, and use Tonelli and~\eqref{d:eq:Schur-column}. The result is
\[
 \int_\pi^\infty(1+t)^{3/2}|G_r(t)|^2dt
 \le C\int_0^\infty(1+s)^{3/2}|G(s)|^2ds,
\]
which is~\eqref{d:eq:uniform-kernel}. The argument starts with bounded compactly supported measurable $G$ and extends by completion. The column domination ensures that input null sets are ignored for almost every output parameter. For the actual continuous exponentially bounded $f$, absolute convergence from Proposition~\ref{d:prop:fractional-source} identifies the extension with its ordinary kernel integral. Cauchy--Schwarz applies without change to complex inputs.
\end{proof}


\section{Compact division and supporting spectral distinctions}\label{app:compactmodel}
The following compact-model results describe exact root spaces and explain why a compact spectral list alone would not establish synthesis. They are supporting results, not substitutes for Appendix~\ref{app:abel}.


\subsection{Compact division and the origin boundary condition}\label{d:sec:division}

For $w\in\C$, write $S_w(s)=\sin(ws)/w$ with $S_0(s)=s$. On the positive half-line set
\begin{equation}\label{d:eq:V-definition}
 (V_wu)(x)=\int_x^\infty S_w(t-x)u(t)\,dt,\qquad x\ge0,
\end{equation}
and extend evenly.

\begin{theorem}[Hilbert Schmidt division]\label{d:thm:division}
For every $w$, $V_w$ is a Hilbert--Schmidt operator on $\Hh$. The map $w\mapsto V_w$ is entire with values in $\HS(\Hh)$ and is bounded there on each horizontal strip $|\operatorname{Im}w|\le B$. If $F_u(w)=0$, then $g=V_wu$ satisfies, distributionally on $\R$,
\begin{equation}\label{d:eq:division-equations}
 g''+w^2g=u,\qquad
 F_g(z)=\frac{F_u(z)}{w^2-z^2},
\end{equation}
where all apparent singularities are removable. Thus division in~\eqref{d:eq:division-equations} is bounded and compact on the closed subspace $\{u:F_u(w)=0\}$. Without that condition the even extension has a derivative jump at zero and the displayed division identity need not hold.
\end{theorem}
\begin{proof}
The isometry $u\mapsto\sqrt2u/\sqrt W$ from $\Hh$ to $L^2(0,\infty)$ transforms $V_w$ into the kernel
\[
 \boldsymbol1_{t\ge x}S_w(t-x)\sqrt{W(t)/W(x)}.
\]
For $x,s\ge0$,~\eqref{d:eq:W-two-sided} gives
\begin{equation}\label{d:eq:weight-ratio}
 \frac{W(x+s)}{W(x)}\le C\exp\bigl(4s-\pi e^{2x}(e^{2s}-1)\bigr).
\end{equation}
Using $S_w(s)=\int_0^s\cos(wr)\,dr$, we have $|S_w(s)|\le s e^{Bs}$ whenever $|\operatorname{Im}w|\le B$. Thus the squared kernel is bounded by
\begin{equation}\label{d:eq:HS-majorant}
 Cs^2\exp\bigl((4+2B)s-\pi e^{2x}(e^{2s}-1)\bigr).
\end{equation}
For $x$ in a bounded interval, the integral in $s$ is finite by double-exponential decay. Choose $X$ so that $\pi e^{2X}\ge4+2B$. For $x\ge X$, the inequality $e^{2s}-1\ge2s$ makes the exponent at most $-\pi e^{2x}s$. The $s^2$ integral is then $2\pi^{-3}e^{-6x}$, integrable in $x$. This proves both the Hilbert--Schmidt property and the uniform strip bound. Parameter derivatives introduce only additional fixed powers of $s$; the same estimates justify holomorphy in the Hilbert--Schmidt norm.

The exponential integrability following~\eqref{d:eq:eval-bound} justifies differentiating the tail integral in the weak sense. On $x>0$,
$g''+w^2g=u$, while
\begin{equation}\label{d:eq:jump}
 g'(0+)=-\int_0^\infty\cos(wt)u(t)\,dt=-\tfrac12F_u(w).
\end{equation}
The even extension therefore satisfies, on all of $\R$,
\begin{equation}\label{d:eq:jump-distribution}
 g''+w^2g=u-F_u(w)\delta_0.
\end{equation}
Exactly the stipulated vanishing removes the delta term. Since both $g,u\in\Hh$ have all exponential moments, their distributional Fourier identity extends from the real axis to the entire functions in~\eqref{d:eq:division-equations}. This proves the quotient identity and removability. At $w=0$, evenness makes $F_u(0)=0$ a zero of order at least two, as required by the denominator $-z^2$.

There is no nonzero even solution of $g''+w^2g=0$ in $\Hh$: such a solution is a scalar cosine, or a constant when $w=0$, and cannot be square integrable against the reciprocal double-exponential weight. Thus the admitted solution is unique. Restricting a Hilbert--Schmidt operator to a closed subspace is compact, completing the proof.
\end{proof}

Every fixed finite sequence of divisions is bounded whenever the remaining zero orders permit it. No bound uniform in the number of divisions, no uniform spectral-projection bound, and no frame estimate has been established. Compact division is not a polynomial-density theorem.

\subsection{A centered compact backward operator}\label{d:sec:model}

Put
\begin{equation}\label{d:eq:d-delta-T}
 d=\frac14,\qquad a=\frac i2,\qquad
 \Delta(z)=a^2-z^2=-\frac14-z^2,\qquad
 T=CV_a|_{\Hh_c}.
\end{equation}
Because $F_u(a)=0$ for $u\in\Hh_c$, $V_a$ is genuine division on this whole subspace. Write
\begin{equation}\label{d:eq:sj}
 s_j=(D^2-d)^j\kappa\qquad(j\ge1).
\end{equation}

\begin{theorem}[Dense range and all generalized nullspaces]\label{d:thm:backward}
The operator $T$ is compact on $\Hh_c$, its range is dense, and for every $n\ge1$,
\begin{equation}\label{d:eq:kernel-exact}
 \ker T^n=\operatorname{span}\{s_1,\ldots,s_n\}
 =\operatorname{span}\{r_1,\ldots,r_n\}.
\end{equation}
These spaces have dimension $n$. The chain action is
\begin{equation}\label{d:eq:shift-chain}
 Ts_1=0,\qquad Ts_j=s_{j-1}\quad(j\ge2),
\end{equation}
and consequently
\begin{equation}\label{d:eq:N-generalized}
 \Nn=\overline{\bigcup_{n\ge1}\ker T^n}^{\Hh}.
\end{equation}
\end{theorem}
\begin{proof}
Compactness follows from Theorem~\ref{d:thm:division} and boundedness of $C$. Since $F_{s_j}=\Delta^j\Xi$, all $s_j$ are centered. Division and centering give~\eqref{d:eq:shift-chain}. The polynomial $\Delta^j$ has leading term $(-z^2)^j$ and vanishes at $a$, so a triangular change of basis gives the two finite spans in~\eqref{d:eq:kernel-exact}. Linear independence follows by dividing any transform relation by the nonzero entire function $\Xi$.

If $Tu=0$, then $V_au=\alpha\kappa$ for some scalar $\alpha$. Applying $D^2-d$ gives $u=\alpha s_1$. More generally, if $Tu$ lies in $\operatorname{span}\{s_1,\ldots,s_n\}$, then $V_au$ is that vector plus a multiple of $\kappa$. Applying $D^2-d$ puts $u$ in $\operatorname{span}\{s_1,\ldots,s_{n+1}\}$. Induction, together with~\eqref{d:eq:shift-chain}, proves the exact kernel formula rather than only one inclusion.

For every even $g\in C_c^\infty(\R)$, let $u=(D^2-d)g$. Integration by parts against $c$ gives $m(u)=0$. The uniqueness part of Theorem~\ref{d:thm:division} gives $V_au=g$ and hence $Tu=Cg$. Even compact smooth functions are dense in $\Hh$: first truncate in the weighted norm, then approximate on a compact set where $W$ is positive and continuous, and symmetrize. Their images under $C$ are therefore dense in $\Hh_c$. This proves dense range and the last assertion.
\end{proof}

For later use the transform action is
\begin{equation}\label{d:eq:T-transform}
 F_{Tu}(z)=\frac{F_u(z)}{\Delta(z)}-\beta(u)\Xi(z),\qquad
 \beta(u)=2F_{V_au}(a).
\end{equation}
In particular $T$ is not self-adjoint: its nontrivial kernel and dense range would be incompatible with self-adjointness.

\begin{proposition}[The invariant divisible ideal]\label{d:prop:T-ideal}
The spaces $\Nn$ and $\Ii$ are $T$-invariant. If $u\in\Ii$ and $F_u=\Xi Q$, then
\begin{equation}\label{d:eq:T-quotient}
 F_{Tu}=\Xi\bigl(R(z)-R(a)\bigr),\qquad
 R(z)=\frac{Q(z)}{\Delta(z)}.
\end{equation}
The function $R$ is even entire.
\end{proposition}
\begin{proof}
Invariance of $\Nn$ follows from~\eqref{d:eq:shift-chain} and continuity. For $\Ii$, the condition $Q(a)=0$ and evenness give zeros at both $a$ and $-a$, so division by $\Delta$ is removable. Since $2\Xi(a)=1$,~\eqref{d:eq:T-transform} becomes~\eqref{d:eq:T-quotient}, whose quotient again vanishes at $a$.
\end{proof}

\subsection{Exact nonzero root spaces and multiplicities}\label{d:sec:roots}

A root space of a bounded operator at $\lambda$ means the union of the kernels of all powers of that operator minus $\lambda$. For a nonzero eigenvalue of a compact operator this is finite dimensional and the union stabilizes. Put
\begin{equation}\label{d:eq:lambda-w}
 \lambda(w)=\frac1{\Delta(w)}=-\frac1{w^2+1/4}.
\end{equation}
It is defined at every zero of $\Xi$ by~\eqref{d:eq:Xi-xi}.

\begin{theorem}[All nonzero adjoint root spaces]\label{d:thm:adjoint}
For each distinct zero pair $\{w,-w\}$ of $\Xi$, the root space of $T^*$ at $\lambda(w)$ is exactly
\begin{equation}\label{d:eq:adjoint-root}
 \operatorname{span}\{J_{w,k}:0\le k<\nu(w)\}.
\end{equation}
It has dimension $\nu(w)$ and consists of one Jordan block. These are all the nonzero root spaces of $T^*$. For the Riesz waves $E_{w,k}$ the corresponding eigenvalue is $\overline{\lambda(w)}$. The sign pair yields the same space, not two copies. Consequently $\Zz$ is the closed span of all nonzero adjoint root spaces.
\end{theorem}
\begin{proof}
For $k<\nu(w)$, differentiate~\eqref{d:eq:T-transform}. All derivatives of the correction term that occur vanish at $w$. Using~\eqref{d:eq:Riesz} and the linear-first convention gives the raw-wave formula
\begin{equation}\label{d:eq:adjoint-action}
 T^*J_{w,k}=\sum_{\ell=0}^k\binom{k}{\ell}
 \lambda^{(\ell)}(w)J_{w,k-\ell}.
\end{equation}
The vectors in~\eqref{d:eq:adjoint-root} are linearly independent. Indeed their exponential-polynomial expressions have distinct leading polynomial degrees at the two distinct frequencies $iw$ and $-iw$; such expressions are independent for $w\ne0$. One may verify this by successively applying powers of $D-iw$ and $D+iw$ to isolate each frequency. Here $w\ne0$ by~\eqref{d:eq:Xi-xi}. Also
\[
 \lambda'(w)=\frac{2w}{\Delta(w)^2}\ne0.
\]
Thus the first adjacent coefficient in the triangular action~\eqref{d:eq:adjoint-action} never vanishes. It is one block of length $\nu(w)$. The identity $J_{-w,k}=(-1)^kJ_{w,k}$ handles the sign pair.

It remains to exhaust the root spaces. Let $h$ be a root vector of $T^*$ at some $\lambda\ne0$, and put $v=h/W$. This is an even locally square-integrable distribution. For even compact smooth $g$, Theorem~\ref{d:thm:backward} gives $T(D^2-d)g=Cg$. Since $h\in\Hh_c$,
\begin{equation}\label{d:eq:root-test}
 \ip{(D^2-d)g}{T^*h}_{\Hh}=\ip{g}{h}_{\Hh}.
\end{equation}
For a chain $T^*h=\lambda h+h_-$, conjugating the tested distributional identity gives
\[
 (D^2-d)(\lambda v+v_-)=v,
 \qquad (D^2-d-\lambda^{-1})v=-\lambda^{-1}(D^2-d)v_-.
\]
The constant-coefficient factors commute. Iterating along a chain proves
\begin{equation}\label{d:eq:root-ODE}
 (D^2-d-\lambda^{-1})^n v=0
\end{equation}
for some $n$. Even test functions suffice because each distribution involved is even.

Choose $w$ such that $\Delta(w)=\lambda^{-1}$. Then~\eqref{d:eq:root-ODE} is $(D^2+w^2)^nv=0$. Distributional solutions of a constant-coefficient ordinary differential equation are the usual exponential polynomials: factor the operator into first-order factors and solve each distributional equation after multiplication by its integrating exponential. For $w\ne0$, its even solutions are therefore finite linear combinations of $\partial_z^k\cos(zx)|_{z=w}$. For $w=0$, they are even polynomials. It remains to determine which of these possibilities survive the Hilbert-space orthogonality constraints.

Every nonzero adjoint root vector annihilates $\ker T^m$ for every $m$. To see this, $T^*$ is invertible on its finite root space, so write $h=(T^*)^mh_m$ within that space. Then $\ip{u}{h}=\ip{T^mu}{h_m}=0$ for $u\in\ker T^m$. In particular $h\perp\Nn$.

Suppose $w\ne0$ and write $h=\sum_{k=0}^M c_kJ_{w,k}$ with $c_M\ne0$. Let $\nu$ be the zero order of $\Xi$ at $w$, using $\nu=0$ if it is not a zero. If $M\ge\nu$, choose the even polynomial
\[
 p(z)=\Delta(z)(z^2-\overline{w}^{\,2})^{M-\nu}.
\]
It vanishes at $a$, so the source $u_p$ with transform $p\Xi$ is a finite radical combination by~\eqref{d:eq:derivative-transforms}. Since $\Delta(\overline w)\ne0$, this transform has exact order $M$ at $\overline w$. But
\[
 \ip{u_p}{h}_{\Hh}
 =\sum_{k=0}^M\overline{c_k}(p\Xi)^{(k)}(\overline w)
 =\overline{c_M}(p\Xi)^{(M)}(\overline w)\ne0,
\]
contradicting $h\perp\Nn$. Hence $M<\nu$. This proves that $w$ must be a zero and that no jet above the asserted order is present.

If $w=0$, write $v=\sum_{j=0}^M c_jx^{2j}$ with $c_M\ne0$. The centered polynomial $p(z)=\Delta(z)z^{2M}$ makes $p\Xi$ have exact zero order $2M$ at zero, since $\Xi(0)\ne0$. The same pairing, using $F_u^{(2j)}(0)=(-1)^j\int x^{2j}u(x)\,dx$, is nonzero. This excludes $w=0$. The case $\Delta(w)=0$ is impossible because $\lambda$ is finite and nonzero. Exhaustion, multiplicity, and the final span statement follow.
\end{proof}

\begin{corollary}[Primal divided chains]\label{d:cor:primal}
For a zero $w$ of multiplicity $\nu$, set $s=w^2+1/4$ and define $h_j\in\Hh$ by
\begin{equation}\label{d:eq:primal-source}
 F_{h_j}(z)=\frac{\Xi(z)}{(w^2-z^2)^j},\qquad
 \phi_j=h_j-s^{-j}\kappa,\qquad 1\le j\le\nu.
\end{equation}
The vectors $h_j$ exist, and the vectors $\phi_j$ are centered and constitute the full root space of $T$ at $-s^{-1}$, of dimension $\nu$. Their action is
\begin{equation}\label{d:eq:primal-action}
 T\phi_j=-\sum_{r=1}^j s^{-(j-r+1)}\phi_r.
\end{equation}
If $j<\nu$, $F_{\phi_j}$ vanishes in value at every zero of $\Xi$, but $\phi_j\notin\Ii$.
\end{corollary}
\begin{proof}
Repeated application of Theorem~\ref{d:thm:division}, starting with $\kappa$, gives $h_j$ for $j\le\nu$. At $a$, its transform is $s^{-j}/2$, giving centering. With $y=w^2-z^2$ and $\Delta=y-s$, the identity
\[
 \frac{y^{-j}-s^{-j}}{y-s}
 =-\sum_{r=1}^j s^{-(j-r+1)}y^{-r}
\]
and centering prove~\eqref{d:eq:primal-action}. The vectors are independent by the resulting rational-function identity after division by $\Xi$. Their adjacent block coefficients are nonzero. For a compact operator the nonzero algebraic multiplicities of $T$ and $T^*$ at conjugate eigenvalues agree; this follows by restricting the two adjoint Riesz projections to their finite-dimensional ranges. Theorem~\ref{d:thm:adjoint}, applied at $\overline w$, therefore shows there is no larger primal root space.

At the pair $\{w,-w\}$, the first term in~\eqref{d:eq:primal-source} has zero order $\nu-j$, whereas the subtracted multiple of $\Xi$ has order $\nu$. The resulting order is exactly $\nu-j$. For $j<\nu$ this is positive but less than $\nu$; all other zeros retain their values. This proves the last assertion without identifying the values-only kernel with $\Ii$.
\end{proof}

The root-space theorems use only multiplicities that actually occur. They neither assume nor prove simplicity. Their algebraic description alone does not establish completeness; that conclusion follows from Theorem~\ref{d:thm:full-synthesis}.

\subsection{The compact quasinilpotent model}\label{d:sec:gap}

\begin{theorem}[Quasinilpotent restriction and synthesis criterion]\label{d:thm:quasinilpotent}
The compact operator $T|_{\Ii}$ is quasinilpotent. Its generalized nullspace is precisely the algebraic span of all the $r_j$. Hence global synthesis~\eqref{d:eq:ideal-synthesis} is equivalent to
\begin{equation}\label{d:eq:remaining-density}
 \overline{\bigcup_{n\ge1}\ker(T|_{\Ii})^n}^{\Hh}=\Ii.
\end{equation}
Before using synthesis, the quotient $\Ii/\Nn$ carries a compact quasinilpotent induced operator and identifies as a Hilbert space with $\Ii\cap\Nn^\perp$. Theorem~\ref{d:thm:full-synthesis} proves this quotient is zero; compactness alone does not.
\end{theorem}
\begin{proof}
Suppose $Tu=\lambda u$ with $u\in\Ii$ and $\lambda\ne0$. For some scalar $\alpha$,
$V_au=\lambda u+\alpha\kappa$. Applying $D^2-d$ and transforming gives
\[
 [1-\lambda\Delta(z)]F_u(z)=\alpha\Delta(z)\Xi(z).
\]
Write $F_u=\Xi Q$ and cancel $\Xi$ away from its discrete zero set, then continue analytically:
\begin{equation}\label{d:eq:no-eigen}
 [1-\lambda\Delta(z)]Q(z)=\alpha\Delta(z).
\end{equation}
At either root of $1-\lambda\Delta(z)$, the right side is $\alpha/\lambda$. Thus $\alpha=0$, including when the root is the double point zero. Equation~\eqref{d:eq:no-eigen} then forces $Q=0$. There are no nonzero eigenvalues of the restriction. A compact operator has no nonzero spectral points except eigenvalues; hence its spectrum is $\{0\}$. The space $\Ii$ is infinite dimensional because it contains the independent radical chains, so zero does belong to this spectrum.

The exact kernels follow from Theorem~\ref{d:thm:backward}, since every $s_j$ belongs to $\Ii$. Their closed span is $\Nn$, proving the equivalence. Both $\Nn$ and $\Ii$ are invariant, so the quotient operator is defined and compact. Its $n$th power has norm at most $\norm{(T|_{\Ii})^n}$, proving spectral radius zero. The usual orthogonal representative identifies the quotient Hilbert space with $\Ii\ominus\Nn$ and represents its action by compression. These structural facts do not by themselves prove that the space is zero; the uniform endpoint estimate will do so.
\end{proof}

Compactness, dense range, and the exact dimensions in~\eqref{d:eq:kernel-exact} do not imply~\eqref{d:eq:remaining-density}. The following elementary example records the precise logical obstruction.

\begin{proposition}[A compact model with an invisible residual part]\label{d:prop:countermodel}
There is a compact quasinilpotent operator with dense range, a one-dimensional kernel, and an $n$-dimensional $n$th kernel for every $n$, whose generalized nullspace is not dense.
\end{proposition}
\begin{proof}
On $\ell^2(\mathbb N_0)$ define $Be_0=0$ and $Be_n=n^{-2}e_{n-1}$ for $n\ge1$. The weights tend to zero, so $B$ is compact, and every basis vector lies in its range. Moreover
\[
 \ker B^n=\operatorname{span}\{e_0,\ldots,e_{n-1}\},\qquad
 \norm{B^n}=(n!)^{-2}.
\]
Thus it has dense range and spectral radius zero.

On $L^2(0,1)$ let $(Vf)(x)=\int_0^x f(t)\,dt$. Its triangular kernel is square integrable, so it is Hilbert--Schmidt. It is injective. Its adjoint is $V^*f(t)=\int_t^1f(x)\,dx$, also injective, so $V$ has dense range. The iterated kernel is
$\boldsymbol1_{t<x}(x-t)^{n-1}/(n-1)!$. Its row and column integrals are bounded by $1/n!$, so $\norm{V^n}\le1/n!$ and $V$ is quasinilpotent.

The direct sum $B\oplus V$ has all the claimed properties, but
\[
 \overline{\bigcup_{n\ge1}\ker(B\oplus V)^n}
 =\ell^2(\mathbb N_0)\oplus\{0\}
 \ne\ell^2(\mathbb N_0)\oplus L^2(0,1).
\]
\end{proof}

This is not a counterexample for the theta space. It shows that the compact-model properties alone cannot exclude a residual quotient. The actual theta-space conclusion uses the additional uniform weighted endpoint estimate of Theorem~\ref{d:thm:uniform-endpoint}, rather than this invalid generic inference.

\subsection{Failure of the inherited de Branges norm axiom}\label{d:sec:debranges}

Every even entire transform can be written uniquely as $f(\lambda)=F_u(z)$ with $\lambda=z^2$. Transport the $\Hh$ norm to this entire-function space in $\lambda$. A necessary de Branges axiom says that, at a nonreal zero $\alpha$ of $f$, the entire function
\begin{equation}\label{d:eq:zero-flip}
 f(\lambda)\frac{\lambda-\overline\alpha}{\lambda-\alpha}
\end{equation}
has the same norm as $f$. This is axiom (H1) in de Branges's original formulation~\cite{deBranges1961}. Bounded division alone does not provide this isometric statement.

\begin{theorem}[Explicit zero flip obstruction]\label{d:thm:debranges}
The Fourier image of $\Hh$, written in the variable $\lambda=z^2$ with its inherited norm, fails the axiom~\eqref{d:eq:zero-flip}. The centered image of $\Hh_c$ fails it as well. The counterexample can be chosen so that both physical numerators are even and compactly supported smooth functions.
\end{theorem}
\begin{proof}
Write $\omega=1/W$ and take even $g\in C_c^\infty(\R;\C)$. For nonreal $\alpha$, define
\[
 u=g''+\alpha g,\qquad \widetilde u=g''+\overline\alpha g.
\]
If $G(\lambda)=F_g(z)$, their transforms are $(\alpha-\lambda)G(\lambda)$ and $(\overline\alpha-\lambda)G(\lambda)$, exactly related by~\eqref{d:eq:zero-flip}. Both are in the space. Direct expansion in the linear-first convention gives
\begin{align}
 \norm{u}_{\Hh}^2-\norm{\widetilde u}_{\Hh}^2
 &=4\operatorname{Im}\alpha\,\operatorname{Im}\int_\R g''\overline g\,\omega\,dx\notag\\
 &=-4\operatorname{Im}\alpha\int_\R\omega'(x)
           \operatorname{Im}(g'(x)\overline{g(x)})\,dx.
 \label{d:eq:flip-norm}
\end{align}
The integration by parts has no endpoint term.

Differentiating the theta asymptotic gives
$(\log\omega)'(x)=2\pi e^{2x}-4+o(1)>0$ for sufficiently large positive $x$. Choose two reflected compact intervals in these tails. On their union $x\omega'(x)>0$. Let $\eta$ be a nonzero real even compact smooth function supported there and put $g(x)=\eta(x)e^{itx^2}$ with real $t\ne0$. Then
\[
 \operatorname{Im}(g'\overline g)=2tx\eta^2,
\]
so~\eqref{d:eq:flip-norm} becomes
\begin{equation}\label{d:eq:flip-nonzero}
 \norm{u}_{\Hh}^2-\norm{\widetilde u}_{\Hh}^2
 =-8t\operatorname{Im}\alpha\int_\R x\omega'(x)\eta(x)^2\,dx\ne0.
\end{equation}
This disproves the axiom for the full image.

For the centered image, choose $\eta$ additionally in the common kernel of the two real linear functionals
\[
 \eta\longmapsto\int_\R c(x)\eta(x)\cos(tx^2)\,dx,
 \qquad
 \eta\longmapsto\int_\R c(x)\eta(x)\sin(tx^2)\,dx.
\]
Such a nonzero choice exists because the real even compact test space on the chosen intervals is infinite dimensional and these impose at most two linear constraints. Now $m(g)=0$, and integration by parts gives $m(g'')=dm(g)=0$. Hence $u,\widetilde u\in\Hh_c$, while the strict sign in~\eqref{d:eq:flip-nonzero} persists. The transform is nonzero by Fourier injectivity and the absence of a nonzero compactly supported solution of $g''+\alpha g=0$. This supplies a genuine nonzero function with a zero at $\alpha$ in the centered space as well.
\end{proof}

The conclusion concerns the actual inherited norm. It does not exclude every equivalent norm, auxiliary entire-function space, nonisometric embedding, or other completeness theorem. Any such approach must establish its own identification and hypotheses before drawing a synthesis conclusion.


\begin{proposition}[Semigroup structure alone is insufficient]\label{d:prop:semigroup-model}
On $\ell^2(\mathbb N)\oplus\C$, define
\[
 B_r(x,c)=\bigl((r^{2n}(x_n+c))_{n\ge1},0\bigr),\qquad 0<r<1.
\]
Each $B_r$ is bounded, compact, and injective, and $B_rB_s=B_{rs}$. Its range closure is $\ell^2(\mathbb N)\oplus\{0\}$, while $B_re_n=r^{2n}e_n$ there. Nevertheless
\[
 \norm{B_r(0,1)}^2=\frac{r^4}{1-r^4}\longrightarrow\infty.
\]
\end{proposition}
\begin{proof}
The map is a compact diagonal plus a rank-one operator. A zero image requires the constant sequence $x_n=-c$, which is square summable only when $c=0$. Direct substitution gives the composition law. Finite sequences in the first summand belong to the range, while the second output coordinate is always zero. The geometric series gives the displayed norm.
\end{proof}

This abstract example is not a theta counterexample: it has no analogue of the proved recovery estimate~\eqref{d:eq:recovery-bound} and uniform kernel bound~\eqref{d:eq:uniform-kernel}. It explains why those estimates, rather than semigroup structure alone, are the decisive steps.


\section{Finite clusters and the fixed complexity proof}\label{app:clusters}


\subsection{Uniform envelopes in a bounded frequency window}\label{e:sec:ode}

\begin{lemma}[Compact ODE estimate]\label{e:lem:ode}
Fix $N\in\N$ and $0\le M<\infty$. If $P$ has complexity at most $N$ and all
its frequencies obey $|\lambda_j|\le M$, then, with $A=WP$,
\begin{equation}\label{e:eq:envelope}
 \norm A_W+\norm{A'}_2\le C_{N,M,W}\norm A_2.
\end{equation}
The estimate is uniform when distinct frequencies approach one another
and when polynomial coefficients vary arbitrarily.
\end{lemma}
\begin{proof}
Fix the total order $m=\sum_j(\deg p_j+1)\le N$. The function $P$
satisfies the monic differential equation
\begin{equation}\label{e:eq:annihilator}
 \prod_j(\partial_x-i\lambda_j)^{\deg p_j+1}P=0.
\end{equation}
The coefficients of the monic polynomial on the left lie in a compact set:
they are continuous symmetric polynomials in $m$ roots from a closed disk.
Repeated roots are included in that set. Let $B_a$ be the corresponding
companion matrix and set $Y=(P,P',\ldots,P^{(m-1)})^{\mathsf T}$. Then
$Y'=B_aY$, and $\norm{B_a}\le L$ uniformly in this compact family.

We claim the uniform observability estimate
\begin{equation}\label{e:eq:observable}
 |Y(0)|^2\le C_m\int_{-1}^1|P(x)|^2\,dx.
\end{equation}
Otherwise there would be coefficient vectors $a_n$ in the compact family
and unit initial vectors $v_n$ for which the integrals tend to zero.
Extract $a_n\to a$ and $v_n\to v$, with $|v|=1$. The power series of the
matrix exponential converges uniformly for $|x|\le1$ and bounded matrices,
so the first components of $e^{xB_{a_n}}v_n$ converge uniformly to that of
$e^{xB_a}v$. The latter component vanishes identically on $[-1,1]$.
Its derivatives of orders $0,\ldots,m-1$ at zero are the components of $v$,
so $v=0$, a contradiction.

The matrix exponential estimate gives
$|Y(x)|\le e^{L|x|}|Y(0)|$ for all real $x$. The same type of bound holds
for $P'$, including when $m=1$ by its differential equation. Since
$w_*:=\min_{[-1,1]}W>0$,
\[
 \int_{-1}^1|P|^2\,dx\le w_*^{-2}\norm{WP}_2^2.
\]
Combining this with~\eqref{e:eq:observable} yields
\begin{equation}\label{e:eq:growth}
 |P(x)|+|P'(x)|\le C e^{L|x|}\norm{WP}_2.
\end{equation}
The integrals of $e^{2L|x|}W$, $e^{2L|x|}W^2$, and
$e^{2L|x|}|W'|^2$ are finite by~\eqref{e:eq:weight}.
Integrating~\eqref{e:eq:growth} and using $(WP)'=W'P+WP'$ proves the lemma
for order $m$. Take the maximum of the finitely many constants for
$1\le m\le N$.
\end{proof}

The compact family is formed only after frequencies are placed in a
bounded window. The companion matrices of an original configuration with
unbounded real frequencies need not be bounded. No inverse Vandermonde
matrix or lower bound on root spacing is used.

\subsection{Finite clusters at arbitrary real separations}\label{e:sec:clusters}

\begin{lemma}[Subsequential cluster decomposition]\label{e:lem:clusters}
Let $P_n\in\EE$ be nonzero. There is a subsequence, still indexed by $n$,
and a fixed finite collection of clusters $\mathcal C$ such that
\begin{equation}\label{e:eq:cluster}
 WP_n=\sum_{C\in\mathcal C}e^{it_{C,n}x}A_{C,n},\qquad
 A_{C,n}=WQ_{C,n},
\end{equation}
where each $t_{C,n}$ is the real part of an original frequency,
all frequencies of each $Q_{C,n}$ lie in a fixed bounded disk depending
on $C$, and $|t_{C,n}-t_{D,n}|\to\infty$ for $C\ne D$.
There are constants $L_C<\infty$, independent of $n$, such that
\begin{equation}\label{e:eq:cluster-envelope}
 \norm{A_{C,n}}_W+\norm{A_{C,n}'}_2\le L_C\norm{A_{C,n}}_2.
\end{equation}
Moreover,
\begin{equation}\label{e:eq:decouple}
 \norm{WP_n}_2^2=(1+o(1))
 \sum_{C\in\mathcal C}\norm{A_{C,n}}_2^2,
\end{equation}
where the relative error is independent of the cluster amplitudes.
\end{lemma}
\begin{proof}
Extract a subsequence on which the number of distinct frequencies and all
polynomial degrees are fixed; there are only finitely many such choices.
Label the frequencies along this subsequence. For every pair $j,k$, extract
further so that $\Rea\lambda_{j,n}-\Rea\lambda_{k,n}$ converges in the
extended real line. There are finitely many pairs. Declare $j\sim k$ when
this limit is finite. Reflexivity and symmetry are immediate; transitivity
follows by adding the two convergent real differences. The equivalence
classes are the clusters.

Choose a member $j_C$ in each cluster and put
\[
 t_{C,n}=\Rea\lambda_{j_C,n},\qquad
 Q_{C,n}(x)=\sum_{j\in C}p_{j,n}(x)
 e^{i(\lambda_{j,n}-t_{C,n})x}.
\]
Within a cluster the real offsets are bounded, while the imaginary parts
remain in $[-\beta,\beta]$. Between distinct clusters the carrier differences
tend in modulus to infinity. This proves the stated decomposition and,
by Lemma~\ref{e:lem:ode},~\eqref{e:eq:cluster-envelope}.

For distinct clusters, integration by parts gives
\begin{align}\label{e:eq:cross}
 \left|\ip{e^{it_{C,n}x}A_{C,n}}{e^{it_{D,n}x}A_{D,n}}_2\right|
 &\le\frac{\norm{A_{C,n}'}_2\norm{A_{D,n}}_2+
 \norm{A_{C,n}}_2\norm{A_{D,n}'}_2}{|t_{C,n}-t_{D,n}|}\notag\\
 &\le\frac{L_C+L_D}{|t_{C,n}-t_{D,n}|}
 \norm{A_{C,n}}_2\norm{A_{D,n}}_2.
\end{align}
The product $A_{C,n}\overline{A_{D,n}}$ belongs to $W^{1,1}(\R)$ by
Cauchy--Schwarz. Its endpoint limits are zero: the derivative is integrable,
so the limits exist, and its integrability forces those limits to vanish.
Thus the integration by parts is valid.

Write $s=|\mathcal C|\le N$ and $h_{C,n}=\norm{A_{C,n}}_2$. If $s>1$,
let $\delta_n$ be the largest coefficient on the right of~\eqref{e:eq:cross}.
Then $\delta_n\to0$ and
\[
 \left|\norm{WP_n}_2^2-\sum_C h_{C,n}^2\right|
 \le2\delta_n\sum_{C<D}h_{C,n}h_{D,n}
 \le(s-1)\delta_n\sum_C h_{C,n}^2.
\]
For $s=1$ the identity is exact. This proves~\eqref{e:eq:decouple}; in
particular, for all sufficiently large $n$,
\begin{equation}\label{e:eq:amplitude}
 \sum_C\norm{A_{C,n}}_2^2\le2\norm{WP_n}_2^2.
\end{equation}
\end{proof}

The constants $L_C$ depend on the bounded offset windows found along the
subsequence. They are not yet universal constants for the original class.
This is sufficient for the contradiction arguments that follow. Crucially,
\eqref{e:eq:amplitude} is established without any prior bound on the separate
envelope norms or on the raw coefficients.

\subsection{Proof of uniform comparison and coercivity}\label{e:sec:proof}

\begin{proof}[Proof of Theorem~\ref{e:thm:main}\textup{(i)}]
Suppose no constant in~\eqref{e:eq:comparison} existed. Choose nonzero
$P_n\in\EE$ with $\norm{WP_n}_W^2>n\norm{WP_n}_2^2$. Apply
Lemma~\ref{e:lem:clusters}. Multiplication by a real modulation is an
isometry in both weighted and ordinary norms. If $s=|\mathcal C|$ and
$L=\max_C L_C$, the triangle inequality and Cauchy--Schwarz give
\[
 \norm{WP_n}_W^2\le s\sum_C\norm{A_{C,n}}_W^2
 \le sL^2\sum_C\norm{A_{C,n}}_2^2
 \le2sL^2\norm{WP_n}_2^2
\]
for all sufficiently large $n$, a contradiction. Thus a single finite
constant exists for the entire class.
\end{proof}

\begin{proof}[Proof of Theorem~\ref{e:thm:main}\textup{(ii)}]
If~\eqref{e:eq:escape} failed, there would be $T_n\to\infty$ and
$P_n\in\EE(T_n)$ such that $\norm{WP_n}_2=1$ and the low-band Fourier
norm remained bounded below by a positive constant. Extract the clusters
of Lemma~\ref{e:lem:clusters}. Each $t_{C,n}$ is the real part of an original
frequency, so $|t_{C,n}|\ge T_n\to\infty$.

For $|t_{C,n}|>R$, Fourier translation and Plancherel imply
\begin{align}\label{e:eq:low-band}
 \norm{\mathbf1_{[-R,R]}
 \widehat{e^{it_{C,n}x}A_{C,n}}}_2^2
 &=\int_{|t|\le R}|\widehat A_{C,n}(t-t_{C,n})|^2\,dt\notag\\
 &\le\frac{\norm{A_{C,n}'}_2^2}{(|t_{C,n}|-R)^2}
 \le\frac{L_C^2\norm{A_{C,n}}_2^2}{(|t_{C,n}|-R)^2}.
\end{align}
Equation~\eqref{e:eq:amplitude} bounds every envelope norm by $\sqrt2$.
There are finitely many clusters, so the triangle inequality makes the
low-band norm of their sum tend to zero, a contradiction.
\end{proof}

\begin{proof}[Proof of Theorem~\ref{e:thm:main}\textup{(iii)}]
Let $C$ be the constant in~\eqref{e:eq:comparison} and put
$\varepsilon(T;R)=\varepsilon_{N,\beta,W}(T;R)$. For every fixed $R$,
Lemma~\ref{e:lem:gamma} and the bounded perturbation estimate give
\begin{equation}\label{e:eq:master-lower}
 \qD[WP]\ge\left\{
 \frac12\log(1+4R^2)\bigl(1-\varepsilon(T;R)^2\bigr)
 -\norm D C\right\}\norm{WP}_2^2
 \qquad(P\in\EE(T)).
\end{equation}
All Gamma frequencies remain in the original form; its nonnegative
low-frequency part was only discarded in this lower estimate. Given
$L>0$, first choose $R$ so that
$\tfrac14\log(1+4R^2)-\norm D C\ge L$. Then choose $T$ large enough
that $\varepsilon(T;R)^2\le1/2$. This proves the first divergence in
\eqref{e:eq:coercivity}. For a prescribed weighted coefficient $L_W>0$,
apply that result with $L=L_W C$ and use~\eqref{e:eq:comparison} to get
$\qD[WP]\ge L_W\norm{WP}_W^2$.
\end{proof}

The order of choices is important: first fix a Fourier band $R$, then
send the minimum real carrier height $T$ to infinity. The extraction proof
does not yield a logarithmic-in-$T$ rate for arbitrary configurations.


\subsection{An entire two frequency formula}\label{e:sec:kernel}

For the theta weight and $D_\theta$ of Section~\ref{e:sec:theta}, define
\begin{equation}\label{e:eq:overlap}
 e_z(x)=W(x)\cos(zx),\qquad
 H_s(t)=\int_\R W(x+s/2)W(x-s/2)\cos(tx)\,dx\quad(s\ge0).
\end{equation}
The following exact formula records the complete form on cosine vectors.
It is not a zero-sum representation and makes no zero-location assumption.

\begin{proposition}[Entire bilinear kernel]\label{e:prop:kernel}
The function
$\mathcal K(z,v)=\qth(e_z,e_{\overline v})$ is entire on $\C^2$ and is
symmetric in $z,v$, even in each variable, and real under simultaneous
complex conjugation. It satisfies
\begin{align}\label{e:eq:kernel}
 \mathcal K(z,v)
 &=\int_0^\infty\rho(s)\Bigl[
 H_0(z-v)+H_0(z+v)\notag\\[-2pt]
 &\hspace{3.5em}-\cos\bigl((z+v)s/2\bigr)H_s(z-v)
 -\cos\bigl((z-v)s/2\bigr)H_s(z+v)\Bigr]\,ds\notag\\
 &\quad-\frac{c_\Gamma}{2}\bigl[H_0(z-v)+H_0(z+v)\bigr]\notag\\
 &\quad-\sum_{n=2}^\infty a_n\Bigl[
 \cos\bigl((z+v)s_n/2\bigr)H_{s_n}(z-v)
 +\cos\bigl((z-v)s_n/2\bigr)H_{s_n}(z+v)\Bigr].
\end{align}
The Gamma bracket is kept combined. The integral and series converge
locally uniformly, and every fixed parameter derivative can be taken under
them. In particular,
\begin{equation}\label{e:eq:jet-matrix}
 \qth(J_{w,k},J_{v,\ell})
 =\left.\partial_z^k\partial_\eta^\ell\mathcal K(z,\eta)
 \right|_{z=w,\,\eta=\overline v}.
\end{equation}
\end{proposition}
\begin{proof}
For a compact parameter set, $e_z$ and each of its fixed parameter
derivatives have uniformly bounded $\HH$ and $H^1$ norms by the theta
tails. The same estimates, with one additional derivative, justify
complex differentiation by dominated convergence in both norms. Thus
$z\mapsto e_z$ is entire into $\XX$ by~\eqref{e:eq:h1bound}.

Set $C_s(z,v)=\int_\R e_z(x+s)e_v(x)\,dx$. With $y=x+s/2$, the weight
product is even in $y$. The product-to-sum identity, whose odd sine terms
integrate to zero, yields
\begin{equation}\label{e:eq:Cs}
 C_s(z,v)=\frac12\Bigl[
 \cos\bigl((z+v)s/2\bigr)H_s(z-v)
 +\cos\bigl((z-v)s/2\bigr)H_s(z+v)\Bigr].
\end{equation}
In particular $C_s$ is symmetric in $z,v$. Translation invariance gives
\[
 \int_\R(e_z(x+s)-e_z(x))(e_v(x+s)-e_v(x))\,dx
 =2C_0(z,v)-2C_s(z,v).
\]
The two prime orientations also both equal $C_{s_n}(z,v)$.
Substitution into the exact forms proves~\eqref{e:eq:kernel}.

On compact parameter sets, the Gamma bracket is $O(s^2)$ near zero,
bounded by $s^2\norm{e_z'}_2\norm{e_v'}_2$; for large $s$ it is bounded
by $4\norm{e_z}_2\norm{e_v}_2$. Both bounds are integrable against
$\rho(s)\,ds$. They remain valid for fixed parameter derivatives, with
the corresponding differentiated vectors. For the prime part,
Lemma~\ref{a:lem:prime} and its norm-summable weighted translation proof give absolute
locally uniform convergence, also after every fixed parameter derivative.
Each $H_s$ is entire by the same theta domination. This proves all
convergence and analyticity assertions. Evenness, symmetry and the real
conjugation identity follow from the formulas; differentiation gives
\eqref{e:eq:jet-matrix}. Notice that the sesquilinear kernel is
$\qth(e_z,e_v)=\mathcal K(z,\overline v)$, not $\mathcal K(z,v)$.
\end{proof}





\section{Quantitative theta and Gamma estimates for growing budgets}\label{sw:sec:theta}
We first verify the regularity and decay required throughout the proof. Set
\begin{equation}\label{sw:eq:theta-constants}
 \begin{split}
 c_W&=\frac{4\pi^2-6\pi}{2},\qquad
 C_4=4\pi^2\sum_{j\ge1}j^4e^{-\pi(j^2-1)},\\
 C_6&=22\pi^3\sum_{j\ge1}j^6e^{-\pi(j^2-1)},\qquad L_W=C_6/c_W.
 \end{split}
\end{equation}
These are positive finite constants. Their defining series have elementary geometric tail majorants, so they introduce no ineffective constants.

\begin{lemma}\label{sw:lem:theta}
The weight $W$ is smooth, positive, and even. For $r=|x|$,
\begin{align}
 c_We^{4r-\pi e^{2r}}&\le W(x)\le C_4e^{4r-\pi e^{2r}},\label{sw:eq:W-bounds}\\
 |W'(x)|&\le C_6e^{6r-\pi e^{2r}},\qquad
 |(\log W)'(x)|\le L_We^{2r}.\label{sw:eq:Wprime}
\end{align}
In particular, on $[-1,1]$ we may use the explicit positive lower bound
\begin{equation}\label{sw:eq:m}
 W(x)\ge m_0:=c_We^{-\pi e^2}.
\end{equation}
\end{lemma}
\begin{proof}
All differentiated series converge locally uniformly for $r\ge0$. Each summand of~\eqref{sw:eq:W} is positive, since $\tau_j\ge\pi>3/2$. Keeping $j=1$ and using $1+e^{-r}\le2$ proves the lower bound. Dropping the negative part of the polynomial and using $e^{2r}\ge1$ proves the upper bound.

Smooth evenness at zero needs justification. Let
\[
 H(x)=e^{x/2}\sum_{j\ge1}e^{-\pi j^2e^{2x}},\qquad x\in\R.
\]
The Gaussian Poisson identity
\[
 \sum_{j\in\mathbb Z}e^{-\pi t j^2}
   =t^{-1/2}\sum_{j\in\mathbb Z}e^{-\pi j^2/t}\quad(t>0)
\]
follows by taking Fourier coefficients of the periodized Gaussian: its $k$th coefficient is $t^{-1/2}e^{-\pi k^2/t}$, and both series converge absolutely. It gives $H(-x)=H(x)+\sinh(x/2)$. Hence $\kappa=H''-H/4$ is smooth and even. Direct differentiation on $x\ge0$ gives
\[
 \kappa(x)=e^{x/2}\sum_{j\ge1}(4\tau_j^2-6\tau_j)e^{-\tau_j},
 \qquad W(x)=\frac{\kappa(x)}{2\cosh(x/2)}.
\]
Thus this smooth even quotient agrees with~\eqref{sw:eq:W} everywhere.

For $r>0$, the derivative of $(4\tau^2-6\tau)e^{-\tau}$ is
$(-8\tau^3+28\tau^2-12\tau)e^{-\tau}$. Including the derivative of the denominator in~\eqref{sw:eq:W} bounds its absolute contribution by
\[
 (8\tau^3+32\tau^2+18\tau)e^{-\tau}\le22\tau^3e^{-\tau},
 \qquad \tau\ge\pi>3.
\]
Summing proves the first inequality in~\eqref{sw:eq:Wprime}; smooth evenness extends it to zero. Division by the lower bound for $W$ proves the logarithmic derivative estimate. Equation~\eqref{sw:eq:m} follows immediately from~\eqref{sw:eq:W-bounds}.
\end{proof}

Use the unitary Fourier convention~\eqref{sw:eq:fourier}.
\begin{lemma}\label{sw:lem:gamma}
For $u\in L^2(\R)$, including possibly infinite values,
\[\tg[u]=\int_\R d_\Gamma(\xi)|\widehat u(\xi)|^2\,d\xi,\quad d_\Gamma\text{ as in~\eqref{sw:eq:gamma-symbol}}.\]
The exact multiplier satisfies
\begin{equation}\label{sw:eq:gamma-lower}
 d_\Gamma(\xi)\ge\tfrac12\log(1+4\xi^2).
\end{equation}
Every finite theta-weighted exponential polynomial belongs to $H^1(\R)$ and has finite Gamma form.
\end{lemma}
\begin{proof}
Plancherel for each shift and Tonelli give~\eqref{sw:eq:gamma-symbol}. Because $1-e^{-2s}\le2s$, its right-hand symbol is at least
\[
 I(\xi)=\int_0^\infty\frac{e^{-s/2}(1-\cos(\xi s))}{s}\,ds.
\]
Differentiating on compact $\xi$ intervals gives
$I'(\xi)=\xi/(\xi^2+1/4)$. Since $I(0)=0$, this is~\eqref{sw:eq:gamma-lower}. Also $\rho_\Gamma(s)=O(s^{-1})$ at zero and $O(e^{-s/2})$ at infinity. The estimate
$\norm{u(\cdot+s)-u}_2\le s\norm{u'}_2$ gives
\[
 \tg[u]\le\left(\int_0^\infty s^2\rho_\Gamma(s)\,ds\right)\norm{u'}_2^2.
\]
The theta bounds make $WP$ and $(WP)'$ square integrable, since $P,P'$ grow at most polynomially times an exponential. This proves the assertion.
\end{proof}
The factor $2$ in~\eqref{sw:eq:gamma-symbol} and the full density~\eqref{eq:common-shifts} are retained throughout. The logarithm is only a lower estimate for the complete multiplier.

\section{Exponential growth and coefficient-free localization}\label{sw:sec:localization}
We use two explicit inequalities for exponential sums from Erd\'elyi~\cite{Erdelyi}. Write $\mathcal E_n$ for sums of at most $n$ functions $a_je^{z_jt}$ with arbitrary complex $a_j,z_j$, and $\mathcal E_n^-$ for those with $\Rea z_j\le0$. Lemma 3.2 of that source gives
\begin{equation}\label{sw:eq:erd-growth}
 |F(t)|\le(2et/d)^n\norm{F}_{L^\infty[a,a+d]},
 \quad F\in\mathcal E_n^-,\quad a,d>0,\quad t\ge a+d.
\end{equation}
Its Lemma 3.5 gives, for arbitrary complex exponents,
\begin{equation}\label{sw:eq:erd-interior}
 |F(y)|\le\sqrt{n/\delta}\,\norm{F}_{L^2[y-\delta,y+\delta]},
 \quad F\in\mathcal E_n,\quad \delta>0.
\end{equation}
In particular, this is not a real-exponent Markov inequality.

\begin{lemma}\label{sw:lem:growth}
If $F(t)=\sum_j e^{z_jt}p_j(t)$ has prefix complexity at most $N$ and $|\Rea z_j|\le\beta$, where $\beta\ge0$, then for $t\ge1$,
\begin{equation}\label{sw:eq:growth}
 |F(t)|\le2\sqrt N\,e^{\beta t+\beta}(4et)^N
                         \norm{F}_{L^2[0,1]}.
\end{equation}
The corresponding negative-tail bound uses $F(-t)$ and its anchor $[-1,0]$.
\end{lemma}
\begin{proof}
First take distinct exponential terms. Apply~\eqref{sw:eq:erd-growth} to $e^{-\beta t}F(t)$ with $a=1/4,d=1/2$. For $y\in[1/4,3/4]$, equation~\eqref{sw:eq:erd-interior} with $\delta=1/4$ gives
$|F(y)|\le2\sqrt N\norm{F}_{L^2[0,1]}$. Damping on this anchor is nonexpansive. This proves even the slightly stronger version without $e^\beta$.

For completeness, confluence preserves the complexity count. Replace each polynomial $p_j(t)$ by $p_j(D_\varepsilon(t))$, where
\begin{equation}\label{sw:eq:D}
 D_\varepsilon(t)=\frac{e^{i\varepsilon t}-1}{i\varepsilon}.
\end{equation}
A cluster of prefix length $m_j$ becomes at most $m_j$ exponentials with roots $z_j+ik\varepsilon$, $0\le k<m_j$. Their real parts are unchanged. Choose a sequence of nonzero real $\varepsilon\to0$ avoiding intercluster collisions, or merge coincidences. The approximants converge uniformly on every compact real interval. Passing to the limit proves~\eqref{sw:eq:growth}. Reflection gives the negative-tail statement.
\end{proof}

We make the uniform absorption of complexity explicit. This will also cover the later derivative tails.
\begin{lemma}\label{sw:lem:absorb}
Fix $B\ge1$ and $b,d\ge0$. There is an effectively specifiable $K\ge e^2$, depending only on $B,b,d$, such that, for all integers $N\ge1$, with
\begin{equation}\label{sw:eq:MR}
 M_N=K(N+1)\log(N+2),\qquad R_N=\tfrac12\log M_N,
\end{equation}
one has, for every $r\ge R_N$,
\begin{equation}\label{sw:eq:absorb}
 B(N+1)^d\exp\{br+N\log(B(r+2))-\pi e^{2r}\}
      \le e^{-(\pi/2)e^{2r}}.
\end{equation}
\end{lemma}
\begin{proof}
Since $r+2\le2e^r$ for $r\ge0$, put $D=\log(2B)$ and bound the positive part of the logarithm on the left by
\[
 \log B+d\log(N+1)+ND+(N+b)r.
\]
Each term divided by $e^{2r}$ is decreasing for $r\ge1$. Also
$R_N\le\frac12\log K+\log(N+2)$, since
$(N+1)\log(N+2)\le(N+2)^2$. At $r=R_N$ the ratio is at most
\begin{equation}\label{sw:eq:K-rule}
 \frac1K\left\{\frac{\log B}{2\log3}+d+\frac{D}{\log3}
       +(b+1)\left(1+\frac{\log K}{2\log3}\right)\right\}.
\end{equation}
Choose $K\ge e^2$ so that this expression is at most $\pi/2$. Such a choice exists because $(1+\log K)/K\to0$, and is specified by an elementary scalar inequality. Then $R_N\ge1$, and monotonicity proves~\eqref{sw:eq:absorb} for the entire ray.
\end{proof}

Apply Lemma~\ref{sw:lem:growth} to $P$ in~\eqref{sw:eq:prefix} with $z_j=i\lambda_j$, so $|\Rea z_j|\le1/2$ independently of the real carriers. The anchor norms are at most $m_0^{-1}\norm{WP}_2$. Lemma~\ref{sw:lem:theta} therefore gives, for $r\ge1$,
\begin{equation}\label{sw:eq:u-envelope}
 \frac{|WP(\pm r)|}{\norm{WP}_2}
 \le \frac{2C_4e^{1/2}}{m_0}\sqrt N\,
             e^{(9/2)r-\pi e^{2r}}(4er)^N
 \qquad(P\ne0).
\end{equation}
The absorption lemma applies with fixed $B,b,d$. After choosing $K$ once, and enlarging it later if necessary, it follows that
\begin{equation}\label{sw:eq:u-tail}
 \norm{\one_{\{|x|>r\}}WP}_2\le e^{-e^{2r}}\norm{WP}_2,
 \qquad r\ge R_N.
\end{equation}
Indeed, squaring~\eqref{sw:eq:absorb}, integrating both tails, and substituting $v=e^{2x}$ bounds the squared ratio by
$e^{-\pi e^{2r}}/(\pi e^{2r})\le e^{-2e^{2r}}$.
No exponential coefficient has been bounded.

\section{A bound for the full prime form}\label{sw:sec:prime}
We include the elementary counting bound that controls the head of the series. Write $\Psi(X)=\sum_{n\le X}\Lambda(n)$.

\begin{lemma}\label{sw:lem:cheb}
For every $X\ge1$,
\begin{equation}\label{sw:eq:cheb}
 \Psi(X)\le3X,\qquad
 \sum_{2\le n\le X}\frac{\Lambda(n)}{\sqrt n}\le6\sqrt X.
\end{equation}
\end{lemma}
\begin{proof}
For an integer $m\ge1$, every prime power in $(m,2m]$ contributes its prime at least once to $\binom{2m}{m}$. This follows from the valuation formula
\[
 v_p\binom{2m}{m}=\sum_{k\ge1}
 \left(\left\lfloor\frac{2m}{p^k}\right\rfloor
       -2\left\lfloor\frac{m}{p^k}\right\rfloor\right),
\]
whose summands are nonnegative and equal one when $m<p^k\le2m$. Thus
$\Psi(2m)-\Psi(m)\le\log\binom{2m}{m}\le2m\log2$.
Telescoping at powers of two gives $\Psi(2^k)<2^{k+1}\log2$. For $2^{k-1}<X\le2^k$ this implies
$\Psi(X)<4X\log2<3X$. The case $X=1$ is immediate; the last strict inequality follows, for example, from $1+3/4+(3/4)^2/2>2$. Abel summation gives
\[
 \sum_{n\le X}\frac{\Lambda(n)}{\sqrt n}
 =\frac{\Psi(X)}{\sqrt X}+\frac12\int_1^X\frac{\Psi(t)}{t^{3/2}}\,dt
 \le6\sqrt X.
\]
\end{proof}

For $u=WP$, denote its correlation by
$C_u(s)=\int_\R u(x+s)\overline{u(x)}\,dx$. For $s\ge0$, the two sets
$\{|x|>s/2\}$ and $\{|x+s|>s/2\}$ cover the line except a null set. Split the integral into the first set and its complement, and use Cauchy--Schwarz, changing variable in the second part. This gives
\begin{equation}\label{sw:eq:correlation}
 |C_u(s)|\le2\norm u_2\norm{\one_{\{|x|>s/2\}}u}_2.
\end{equation}
Always $|C_u(s)|\le\norm u_2^2$. Let $J_N=\lfloor M_N\rfloor$. For the head $2\le n\le J_N$, Lemma~\ref{sw:lem:cheb} gives a contribution to $|p[u]|$ of at most $12\sqrt{M_N}\norm u_2^2$. For every remaining integer $n>M_N$, equations~\eqref{sw:eq:u-tail} and~\eqref{sw:eq:correlation} give
\[
 |C_u(\log n)|\le2e^{-n}\norm u_2^2.
\]
The entire tail is therefore at most $C_P\norm u_2^2$, where a sufficient fixed constant is
\begin{equation}\label{sw:eq:CP}
 C_P=4\frac{e^{-1}}{(1-e^{-1})^2};
\end{equation}
we used $\Lambda(n)/\sqrt n\le\log n/\sqrt n\le n$ and the geometric derivative sum. This also proves absolute convergence of~\eqref{sw:eq:p}. In view of~\eqref{sw:eq:MR}, one may take
\begin{equation}\label{sw:eq:K1}
 K_1=12\sqrt K+C_P.
\end{equation}
This proves Proposition~\ref{sw:prop:analytic}(i), including all prime powers and the infinite remainder.

\section{A local derivative bound for complex confluent sums}\label{sw:sec:derivative}
The bounded-offset derivative estimate must be uniform as roots coalesce. We derive a quantitative local bound from a complex-exponent half-line theorem, rather than infer a rate from finite-dimensional compactness.

For a function $f$ on $[0,\infty)$, let
$\norm f_{2,e^{-t}}^2=\int_0^\infty|f(t)|^2e^{-t}\,dt$.
Erd\'elyi's Theorem 2.9.1~\cite{Erdelyi}, for complex roots $\eta_j$ with $\Rea\eta_j<1/2$, states
\begin{equation}\label{sw:eq:source-derivative}
 \norm{f'}_{2,e^{-t}}\le
 \left[\max_j|\eta_j|+
 \sqrt{\sum_{j<k}(1-2\Rea\eta_j)(1-2\Rea\eta_k)}\right]
 \norm f_{2,e^{-t}}.
\end{equation}
His Lemma 3.4 gives
\begin{equation}\label{sw:eq:source-eval}
 |f(0)|\le\sqrt{\sum_j(1-2\Rea\eta_j)}\,
                         \norm f_{2,e^{-t}}.
\end{equation}
We use these weighted statements directly. Unweighted norms below satisfy $\norm f_2^2=\int|f|^2$; no derivative norm is used.

\begin{lemma}\label{sw:lem:local-derivative}
Suppose $Q(t)=\sum_j e^{z_jt}p_j(t)$ has prefix complexity at most $N$ and $|z_j|\le2$. With
\begin{equation}\label{sw:eq:D0}
 D_0=10\,18^{3/2},
\end{equation}
one has, for every $x\in\R$ and $0<\delta\le1$,
\begin{equation}\label{sw:eq:local-derivative}
 |Q'(x)|\le D_0(N+1)^3\delta^{-3/2}
                            \norm Q_{L^2[x,x+\delta]}.
\end{equation}
\end{lemma}
\begin{proof}
First consider sums of $n\le N$ ordinary exponentials. Put $a=16(N+1)$ and $g(t)=e^{-at}Q(t)$. Lemma~\ref{sw:lem:growth}, with $\beta=2$, bounds the squared tail integrand relative to $\norm Q_{L^2[0,1]}^2$ by
\[
 4N e^{4t+4}(4et)^{2N}e^{-32(N+1)t},\qquad t\ge1.
\]
For $t\ge1$ and $N\ge1$, use
$\log(4N)\le2N$, $\log(4et)\le3t$, $4\le4Nt$, and $4t\le4Nt$.
The logarithm of this expression is at most $-16Nt-32t$, so its integral on $[1,\infty)$ is less than one. Damping is at most one on $[0,1]$, and hence
\begin{equation}\label{sw:eq:damped-norm}
 \norm g_{L^2[0,\infty)}\le\sqrt2\norm Q_{L^2[0,1]}.
\end{equation}
Write $\mu_j=z_j-a$. Apply~\eqref{sw:eq:source-derivative} to $f=e^{t/2}g$, whose roots are $\eta_j=\mu_j+1/2$. Its left side is $\norm{g'+g/2}_2$. The triangle inequality yields
\begin{equation}\label{sw:eq:Bmu}
 \norm{g'}_2\le B(\mu)\norm g_2,
 \quad B(\mu)=\tfrac12+\max_j|\mu_j+\tfrac12|
                +2\sqrt{\sum_{j<k}\Rea\mu_j\Rea\mu_k}.
\end{equation}
All these norms are on the half-line. Since $|z_j|\le2$ and $\Rea\mu_j<0$,
$B(\mu)\le4N(a+2)$. Applying~\eqref{sw:eq:source-eval} to $e^{t/2}h$, for $h=g$ and $h=g'$, gives
\begin{equation}\label{sw:eq:endpoint}
 |h(0)|\le\sqrt{-2\sum_j\Rea\mu_j}\,\norm h_2
          \le\sqrt{2N(a+2)}\,\norm h_2.
\end{equation}
Differentiation preserves the same exponential span. Therefore, from $Q'(0)=g'(0)+ag(0)$,
\begin{align*}
 |Q'(0)|
 &\le\sqrt{2N(a+2)}\,[4N(a+2)+a]\sqrt2\norm Q_{L^2[0,1]}\\
 &\le10[N(a+2)]^{3/2}\norm Q_{L^2[0,1]}
 \le D_0(N+1)^3\norm Q_{L^2[0,1]}.
\end{align*}

We justify passage to polynomial prefixes also on the unbounded interval. Use~\eqref{sw:eq:D} in each cluster. For real $t\ge0$,
\[
 |D_\varepsilon(t)|\le t,\qquad |D_\varepsilon'(t)|=1.
\]
For a fixed confluent $Q$ and its prefixes, the damped approximants and their first derivatives consequently have an $\varepsilon$-independent polynomial-times-$e^{-(a-2)t}$ majorant. Its square is integrable. Their endpoint values and derivatives converge, as do their half-line norms, by dominated convergence. The shifted roots preserve real parts; their moduli tend to those of the original repeated roots. Apply the exact continuous multipliers in~\eqref{sw:eq:Bmu}--\eqref{sw:eq:endpoint} before taking the limit, and then use the displayed bounds with $|z_j|\le2$. The damping estimate follows from the already established confluent Lemma~\ref{sw:lem:growth}. Thus the same constants pass to the limit. The majorant may depend on the fixed $Q$, but the resulting inequality does not depend on its coefficients or root separations.

Finally apply the endpoint result to $F(t)=Q(x+\delta t)$. Its roots are $\delta z_j$, still of modulus at most two, while
$F'(0)=\delta Q'(x)$ and
$\norm F_{L^2[0,1]}=\delta^{-1/2}\norm Q_{L^2[x,x+\delta]}$.
This proves~\eqref{sw:eq:local-derivative}.
\end{proof}

\section{The weighted envelope derivative}\label{sw:sec:weighted}
Let $Q$ satisfy the hypotheses of Proposition~\ref{sw:prop:analytic}(ii) and set $A=WQ$. The claim is trivial for $Q=0$, so assume $Q\ne0$. Its exponential roots $z_j=i\lambda_j$ have modulus at most $\sqrt{5}/2<2$. Use~\eqref{sw:eq:MR} and put $\delta=M_N^{-1}$. If $|x|\le R_N$ and $x\le y\le x+\delta$, then~\eqref{sw:eq:Wprime} gives
\[
 |\log W(x)-\log W(y)|\le L_W\delta e^{2(R_N+\delta)}
 \le L_We^2.
\]
Write $C_{\rm osc}=e^{L_We^2}$. Multiplying~\eqref{sw:eq:local-derivative} by $W(x)$ and using $W(x)/W(y)\le C_{\rm osc}$ yields
\begin{align}\label{sw:eq:interior}
 \int_{|x|\le R_N}|W(x)Q'(x)|^2\,dx
 &\le D_0^2C_{\rm osc}^2(N+1)^6\delta^{-3}
       \int_{|x|\le R_N}\int_x^{x+\delta}|A(y)|^2\,dy\,dx\\
 &\le D_0^2C_{\rm osc}^2(N+1)^6M_N^2\norm A_2^2.\nonumber
\end{align}
Fubini supplies the factor $\delta$ in the second line: each fixed $y$ is counted over a set of $x$ of length at most $\delta$. Also,
\begin{equation}\label{sw:eq:interior-Wprime}
 \norm{\one_{\{|x|\le R_N\}}W'Q}_2\le L_WM_N\norm A_2.
\end{equation}

The exterior estimate must not use an unknown global derivative norm. Instead, apply~\eqref{sw:eq:local-derivative} with $\delta=1$ for $x\in[-1,0]$. Since $[x,x+1]\subset[-1,1]$,
\begin{equation}\label{sw:eq:derivative-anchor}
 \norm{Q'}_{L^2[-1,0]}
 \le D_0(N+1)^3\norm Q_{L^2[-1,1]}
 \le D_0m_0^{-1}(N+1)^3\norm A_2.
\end{equation}
The function $Q'$ has complexity at most $N$ and the same root bounds. For its positive tail use Lemma~\ref{sw:lem:growth} on $F(t)=Q'(t-1)$ at $t=r+1$; for its negative tail use $F(t)=Q'(-t)$ at $t=r$. Both anchors are exactly $[-1,0]$. Combining with~\eqref{sw:eq:derivative-anchor} and~\eqref{sw:eq:W-bounds} bounds $|WQ'|/\norm A_2$ by an expression of the form on the left of~\eqref{sw:eq:absorb}, with fixed $B,b,d$; the extra factors are at most $(N+1)^{7/2}$ and the distance factor is at most $(r+2)^N$. For $W'Q$, use~\eqref{sw:eq:growth} and~\eqref{sw:eq:Wprime}; only the fixed linear exponent changes from $9r/2$ to $13r/2$.

Choose one $B$ and one pair $b,d$ dominating these two envelopes and~\eqref{sw:eq:u-envelope}, and increase the single $K$ using~\eqref{sw:eq:K-rule}. Integrating the absorbed exterior envelopes, exactly as for~\eqref{sw:eq:u-tail}, gives
\begin{equation}\label{sw:eq:exterior}
 \norm{\one_{\{|x|>R_N\}}WQ'}_2\le\norm A_2,
 \qquad
 \norm{\one_{\{|x|>R_N\}}W'Q}_2\le\norm A_2.
\end{equation}
A sufficient common choice is $b=13/2$, $d=7/2$, and
\begin{equation}\label{sw:eq:common-B}
 B=\max\left\{1,4e,\frac{2C_4e^{1/2}}{m_0},
                  \frac{2C_4D_0e}{m_0},\frac{2C_6e^{1/2}}{m_0}\right\}.
\end{equation}
This makes the common choice of $K$ in~\eqref{sw:eq:K-rule} completely specified by elementary inequalities and the fixed theta series.

Combining~\eqref{sw:eq:interior}, \eqref{sw:eq:interior-Wprime}, and~\eqref{sw:eq:exterior} proves
\[
 \norm{A'}_2\le (D_0C_{\rm osc}+L_W+2)(N+1)^3M_N\norm A_2.
\]
Thus a sufficient symbolic choice is
\begin{equation}\label{sw:eq:K2}
 K_2=K(D_0C_{\rm osc}+L_W+2).
\end{equation}
Together with~\eqref{sw:eq:K1}, this proves the two quantitative estimates with a single consistent $K$. No compactness-to-polynomial inference is used.

\section{The single window lower estimate}\label{sw:sec:carrier}
We prove Proposition~\ref{sw:prop:analytic}(iii), completing its proof. The evenness of $W$ is essential to the representation~\eqref{sw:eq:packet}: $A(-x)=W(x)Q(-x)$, so $u_T$ is $W$ times the paired exponential polynomial and is even. Reflection preserves both ordinary norms $\norm A_2$ and $\norm{A'}_2$.

The cross term in its squared norm is
\[
 I_T=\int_\R e^{2iTx}A(x)\overline{A(-x)}\,dx.
\]
The derivative of $A(x)\overline{A(-x)}$ has $L^1$ norm at most
$2\norm{A'}_2\norm A_2$. The theta decay makes its boundary terms vanish. Integration by parts therefore gives
\[
 |I_T|\le\frac{b_N}{T}\norm A_2^2,
 \qquad
 \norm{u_T}_2^2\ge\frac12\left(1-\frac{b_N}{T}\right)\norm A_2^2.
\]
For $T\ge2b_N$, in particular,
\begin{equation}\label{sw:eq:norm-compare}
 \norm A_2\le2\norm{u_T}_2.
\end{equation}
Plancherel gives
$\int_{|\xi|\ge T/2}|\widehat A(\xi)|^2d\xi
 \le4\norm{A'}_2^2/T^2$.
The two translated Fourier envelopes in~\eqref{sw:eq:packet}, including their factors $1/2$, consequently satisfy
\begin{equation}\label{sw:eq:leakage}
 \norm{\one_{[-T/2,T/2]}\widehat u_T}_2
 \le\frac{2b_N}{T}\norm A_2
 \le\frac{4b_N}{T}\norm{u_T}_2.
\end{equation}
On the complementary band, Lemma~\ref{sw:lem:gamma} gives
$d_\Gamma(\xi)\ge\frac12\log(1+T^2)$. Subtract the constant term and use~\eqref{sw:eq:prime-main} for the \emph{total} paired prefix budget. Equation~\eqref{sw:eq:lower} follows.



% Consolidated response tools. No additional macros are required.
\section{Complete response certificates and their limiting obligations}\label{rtc:sec:tools}
This appendix develops response tests for the same centered weighted operator,
not for a local pencil or a projected operator. Throughout,
$\Hh=L^2_{\mathrm{even}}(\R,dx/W)$, $\Hh_c=\ker m$,
$X_c=X\cap\Hh_c$, and $B$ denotes the restriction to $\Hh_c$ of the
actual representative of $q_0$. Write $q=q_0|_{X_c}$.
The realization and full-domain bridge are those of
Section~\ref{sec:weighted-physical} and Proposition~\ref{a:prop:domain}.
In particular $B\ge-I/2$, since the original energy is nonnegative.
All physical inner products and norms below are weighted ones; coefficient
matrix norms, when used, are Euclidean operator norms in the stated basis.
The convention is linearity in the first inner-product argument and matrix
entries $M_{ij}=\ip{v_j}{v_i}_{\Hh}$, so quadratic forms are $x^*Mx$.

Set
\begin{equation}\label{rtc:eq:setup}
 A=B+I,\qquad R=A^{-1},\qquad
 q_\sigma=q+\sigma\norm{\cdot}_{\Hh}^{2},\qquad k=1-\sigma.
\end{equation}
When a sharper prior is used it is an explicit premise
$B\ge-dI$ for this same operator, with $0<d<1/2$ and $a=1-d$.
The unconditional endpoint is $d=1/2$, $a=1/2$.
Thus $0<R\le a^{-1}I$: $R$ is injective and positive on each nonzero
vector, but need not be uniformly bounded below. The letter $a$ here is
an inverse-coercivity constant, not a spatial cutoff radius.
No rounded numerical coefficient is silently substituted for $d$.
These constructions add no actual theta matrix or numerical sign evidence.

\subsection{Divided sources and the full physical kernel}\label{rtc:null:section}
The full zero-jet sector of Theorem~\ref{thm:actual-core} is the orthogonal
complement of the derivative radical closure $\mathcal N$, not necessarily
the orthogonal complement of the entire physical kernel.  The following
identification separates these spaces without a simplicity assumption.
Only in this subsection, $B=T_\Gamma-D_b$ is the ambient weighted operator from
\eqref{eq:Db-physical}, and $B_c$ is its self-adjoint restriction to
$\Hh_c$.  We retain $X_c=X\cap\Hh_c$, $q_0$, and the complete prime and
Gamma terms.  The normalization and reduction already proved in
Theorem~\ref{thm:actual-core} give
\[
 \norm{\kappa}_{\Hh}=1,\qquad
 \ip{u}{\kappa}_{\Hh}=2m(u),\qquad
 Cu=u-2m(u)\kappa,\qquad B\kappa=-\kappa/2.
\]
In particular, $C$ preserves the actual operator domain and is bounded in
its graph norm.  No ordinary-$L^2$ projection is substituted for $C$.

\begin{lemma}[Division admits actual operator rows]\label{rtc:null:division}
The division operator $V_w$ of~\eqref{d:eq:V-definition} maps $\Hh$
boundedly into $\operatorname{Dom}B$.  Its weighted first derivative is
bounded, uniformly for $|\operatorname{Im}w|\le b$ with fixed $b$.
The quotient identity of Theorem~\ref{d:thm:division} still requires
$F_u(w)=0$; operator-domain admission does not require that vanishing.
\end{lemma}
\begin{proof}
Theorem~\ref{d:thm:division} proves the weighted bound for $V_w$ itself.
Differentiating its half-line tail integral gives
\[
 (V_wu)'(x)=-\int_x^\infty\cos(w(t-x))u(t)\,dt\qquad(x>0).
\]
After the half-line isometry $u\mapsto\sqrt2u/\sqrt W$, the square of
this derivative kernel is bounded, by~\eqref{d:eq:weight-ratio}, by
\[
 C\exp\{(4+2b)s-\pi e^{2x}(e^{2s}-1)\},\qquad x,s\ge0.
\]
For bounded $x$ its $s$ integral is finite by double-exponential decay.
For sufficiently large $x$, its exponent is at most $-\pi e^{2x}s$;
the $s$ integral is then at most $\pi^{-1}e^{-2x}$, which is integrable
in $x$.  Thus the derivative operator is Hilbert--Schmidt into weighted
$L^2$, with a bound uniform on the stated strip.  The even extension is
continuous at zero, so its first distributional derivative contains no
point mass.  Boundedness of $W$ gives $V_wu\in\Hh\cap H^1(\R)$.
Lemma~\ref{d:lem:Gamma-H1}, with $\mathsf D=D_b$, now supplies actual
$B$-domain admission and a graph bound.  The second derivative can still
have the point mass in~\eqref{d:eq:jump-distribution}; hence the distinct
vanishing condition for transform division has not been discarded.
\end{proof}

For each zero $w$ of $\Xi$, let $\nu(w)$ be its multiplicity and put
$J_w(x)=W(x)\cos(wx)$.  The exponential weight moments give
\begin{equation}\label{rtc:null:values}
 J_w\in\Hh_c,\qquad
 \ip{J_w}{h}_{\Hh}=\overline{F_h(\overline w)},\qquad
 \sup_{|\operatorname{Im}w|\le1/2}\norm{J_w}_{\Hh}<\infty.
\end{equation}
Indeed $m(J_w)=\Xi(w)/2=0$.  Notice the conjugation: the Riesz vector
for evaluation at $w$ is $J_{\overline w}$.

\begin{lemma}[An absolutely convergent full operator row]\label{rtc:null:row}
For every even compact smooth $g$,
\begin{equation}\label{rtc:null:row-identity}
 Bg=\sum_{\Xi(w)=0}\nu(w)F_g(w)J_w-m(g)\kappa,\qquad
 B(Cg)=\sum_{\Xi(w)=0}\nu(w)F_g(w)J_w.
\end{equation}
The sums run over distinct zeros, including both signs, and converge
absolutely in $\Hh$.
For $u\in\operatorname{Dom}B$ and even compact smooth $g$ one also has
\begin{equation}\label{rtc:null:mixed}
 q_0(u,g)+2m(u)\overline{m(g)}
 =\sum_{\Xi(w)=0}\nu(w)F_u(w)\overline{F_g(\overline w)},
\end{equation}
again absolutely convergently.
\end{lemma}
\begin{proof}
Repeated integration by parts makes $F_g(w)$ decrease faster than any
fixed inverse power of $1+|\operatorname{Re}w|$, uniformly in the zero
strip.  Combining~\eqref{rtc:null:values} with the zero count
\eqref{a:eq:zerocount} proves absolute convergence in $\Hh$.
Pair the first proposed row with a compact smooth $h$.  Its pairing is
the complex version of Lemma~\ref{a:lem:bilinear}, since
$\ip{m(g)\kappa}{h}_{\Hh}=2m(g)\overline{m(h)}$.
These tests are dense in $\Hh$, and $g$ is in the actual operator domain
by Theorem~\ref{d:thm:Gamma-graph-core}; consequently the row is $Bg$.
Using $B\kappa=-\kappa/2$ proves the centered identity.
For the mixed identity, pair $u$ with the actual row for $Bg$, use
self-adjointness, and reindex $w\mapsto\overline w$.
The function $F_u$ is uniformly bounded on the zero strip by
\eqref{d:eq:eval-bound}, while $F_g$ has the rapid decay just established.
This proves absolute convergence and~\eqref{rtc:null:mixed} without
asserting a zero-sum formula for two arbitrary rough form vectors.
\end{proof}

\begin{theorem}[Normalized divided sources]\label{rtc:null:source}
For a zero $w$ of multiplicity $\nu$, use the divided chain $h_j,\phi_j$
from Corollary~\ref{d:cor:primal}, and set
\begin{equation}\label{rtc:null:normalization}
 s=w^2+\tfrac14,\qquad
 a_w=\frac{\Xi^{(\nu)}(w)}{\nu!(-2w)^\nu},\qquad
 \psi_w=a_w^{-1}\phi_\nu.
\end{equation}
Then $s\ne0$, $a_w\ne0$, every $h_j,\phi_j,\psi_w$ is in
$\operatorname{Dom}B$, and
\begin{equation}\label{rtc:null:source-row}
 B\phi_j=0\quad(1\le j<\nu),\qquad
 B\psi_w=2\nu(w)J_w.
\end{equation}
The centered transform $F_{\psi_w}$ equals one at $w$ and $-w$, and
vanishes at every other zero of $\Xi$.
\end{theorem}
\begin{proof}
The existence, quotient transforms, and exact centering
$\phi_j=h_j-s^{-j}\kappa$ are already proved in
Corollary~\ref{d:cor:primal}.  Its construction is $h_j=V_w^j\kappa$.
Lemma~\ref{rtc:null:division} admits every such vector to the actual
operator domain; subtracting the admitted vector $s^{-j}\kappa$
preserves this admission.  Also $w\ne0$ because $\Xi(0)=\int_\R\kappa>0$, and $s\ne0$ because
$\Xi(\pm i/2)=1/2$ in~\eqref{d:eq:Xi-xi}.  The removable value at $w$ of
$\Xi(z)/(w^2-z^2)^\nu$ is $a_w\ne0$; evenness gives the same value
at $-w$.  At all other zeros the quotient, and therefore the centered
transform, vanishes.  For $j<\nu$ the remaining zero order at $w$ is
positive, so $F_{\phi_j}$ vanishes at every zero.
Apply~\eqref{rtc:null:mixed} to these admitted centered vectors.
The sum is zero for $\phi_j$, and consists of the two terms $w,-w$
for $\psi_w$.  Those terms equal
$2\nu(w)\overline{F_g(\overline w)}=\ip{2\nu(w)J_w}{g}_{\Hh}$.
Density of the compact tests identifies the full operator rows in
\eqref{rtc:null:source-row}.
\end{proof}

Define two closed subspaces of the centered weighted space by
\begin{equation}\label{rtc:null:spaces}
 \begin{split}
 \mathcal V_0&=\{u\in\Hh_c:F_u(w)=0\text{ at every zero of }\Xi\},\\
 \mathcal Z_0&=\overline{\operatorname{span}\{J_w:\Xi(w)=0\}}^{\Hh}.
 \end{split}
\end{equation}
Here $\mathcal Z_0$ is the span of value waves only; it is not the
full-jet span $\mathcal Z$ from~\eqref{d:eq:N-intro}.

\begin{theorem}[Full physical kernel and range closure]\label{rtc:null:kernel}
In the original centered weighted realization,
\begin{equation}\label{rtc:null:kernel-identity}
 \ker B_c=\mathcal V_0,\qquad
 \overline{\operatorname{Ran}B_c}^{\Hh}=\mathcal Z_0,\qquad
 \Hh_c=\mathcal V_0\oplus\mathcal Z_0.
\end{equation}
Both spaces reduce $B_c$.  In particular every member of
$\mathcal V_0$, although defined only by Hilbert-space evaluations,
automatically belongs to $\operatorname{Dom}B$ and has zero actual row.
\end{theorem}
\begin{proof}
The ambient compact smooth graph core in
Theorem~\ref{d:thm:Gamma-graph-core}, followed by the graph-bounded
projection $C$, is a graph core for $B_c$.
Its image under $B_c$ is in $\mathcal Z_0$ by
\eqref{rtc:null:row-identity}.  Graph approximation and closedness of
$\mathcal Z_0$ therefore imply
$\operatorname{Ran}B_c\subseteq\mathcal Z_0$.
Conversely~\eqref{rtc:null:source-row} puts every $J_w$ in
$\operatorname{Ran}B_c$, giving the asserted range closure.
Equation~\eqref{rtc:null:values} and conjugation symmetry of the zero
set identify $\mathcal V_0=\mathcal Z_0^\perp$ inside $\Hh_c$.
For a densely defined self-adjoint operator,
$(\operatorname{Ran}B_c)^\perp=\ker B_c$: orthogonality to the range
is precisely membership in the domain of its adjoint with zero adjoint
row.  This proves the kernel identity and the stated domain admission.
The kernel projection is a spectral projection, so it and its orthogonal
complement reduce $B_c$.
\end{proof}

\begin{corollary}[Multiplicity can enlarge the physical kernel]\label{rtc:null:multiplicity}
One has $\mathcal N\subseteq\mathcal V_0$.  If a zero has multiplicity
$\nu>1$, its vectors $\phi_1,\ldots,\phi_{\nu-1}$ are in
$\ker B_c$ and are linearly independent modulo $\mathcal N$.
Consequently no identification $\ker B_c=\mathcal N$ follows without
excluding these additional directions.
\end{corollary}
\begin{proof}
The transforms of the derivative radicals are polynomial multiples of
$\Xi$.  Every derivative evaluation is continuous by
\eqref{d:eq:eval-bound}; hence their Hilbert closure $\mathcal N$
retains vanishing to the full multiplicity at every zero.  Thus
$\mathcal N\subseteq\mathcal V_0$.
Theorem~\ref{rtc:null:source} puts each indicated $\phi_j$ in the kernel.
Its transform has exact order $\nu-j$ at $w$, as established in
Corollary~\ref{d:cor:primal}.  In a nonzero combination choose the
largest $j$ with nonzero coefficient.  Its strictly smallest zero order
cannot be canceled by any other summand, so the combination fails
full-multiplicity vanishing and cannot lie in $\mathcal N$.
This is a conditional statement about a multiple zero, not an assertion
that such a zero exists.
\end{proof}

\begin{proposition}[Exact conditional sign blocks]\label{rtc:null:sign}
Choose $\psi_{-w}=\psi_w$ and
$\psi_{\overline w}=\overline{\psi_w}$.
For a real zero $w>0$,
$q_0[\psi_w]=2\nu(w)$.
For a nonreal zero with $\operatorname{Re}w>0$,
\begin{equation}\label{rtc:null:sign-block}
 \bigl(q_0(\psi_\alpha,\psi_\beta)\bigr)_{
             \alpha,\beta\in\{w,\overline w\}}
 =2\nu(w)\begin{pmatrix}0&1\\1&0\end{pmatrix},\qquad
 q_0[\operatorname{Re}\psi_w]=\nu(w),\quad
 q_0[\operatorname{Im}\psi_w]=-\nu(w).
\end{equation}
Different sign-and-conjugation orbits have zero cross-pairing.
If such a nonreal zero exists, its real centered admitted profile
$f=\operatorname{Im}\psi_w/\kappa$ satisfies
\begin{equation}\label{rtc:null:rayleigh}
 \frac{\mathcal E[f]}{\operatorname{Var}_\mu(f)}
 =\frac12-\frac{\nu(w)}{
                \norm{\operatorname{Im}\psi_w}_{\Hh}^2}.
\end{equation}
\end{proposition}
\begin{proof}
Reality and evenness of the construction give the chosen symmetries.
There are no purely imaginary zeros, since
$\Xi(iy)=\int_\R\kappa(x)\cosh(yx)\,dx>0$.
For any two normalized sources, the actual row gives
\[
 q_0(\psi_w,\psi_v)
 =2\nu(w)\overline{F_{\psi_v}(\overline w)}.
\]
The values in Theorem~\ref{rtc:null:source} yield each matrix entry and
the zero cross-pairings.  Expansion of
$\operatorname{Re}\psi_w=(\psi_w+\psi_{\overline w})/2$ and
$\operatorname{Im}\psi_w=(\psi_w-\psi_{\overline w})/(2i)$ gives the
two signed values.  The negative value ensures that the imaginary part
is nonzero.  The full-domain bridge in
Proposition~\ref{a:prop:domain}, applied to this centered operator-domain
vector, proves both admission of $f$ to the original energy domain and
\eqref{rtc:null:rayleigh}.
\end{proof}
A certified prior $q_0[u]\ge-\varepsilon\norm{u}_{\Hh}^2$ on $X_c$,
with $\varepsilon>0$, would therefore impose
$\norm{\operatorname{Im}\psi_w}_{\Hh}^2\ge\nu(w)/\varepsilon$
on a hypothetical nonreal zero.  This is not an exclusion: the factor
$a_w^{-1}$ and increasing division order have no established uniform
bound here.  Likewise, positivity on individual real-zero sources is
not a proof of positivity on arbitrary form-domain vectors.

\begin{proposition}[Full-kernel deflation in the actual norms]\label{rtc:null:deflation}
Let $P_0$ be the orthogonal projection of $\Hh_c$ onto $\mathcal V_0$
and $P_*=I-P_0$.  These projections preserve the form and operator
domains and are bounded in their respective norms.  For $u\in X_c$
and $\delta\ge0$,
\begin{equation}\label{rtc:null:deflation-identity}
 \begin{split}
 q_0[u]&=q_0[P_*u],\\
 q_0[u]+\delta\norm{u}_{\Hh}^2
 &=\delta\norm{P_0u}_{\Hh}^2+
   q_0[P_*u]+\delta\norm{P_*u}_{\Hh}^2.
 \end{split}
\end{equation}
The projections under $P_*$ of the full zero-jet family from
Theorem~\ref{thm:actual-core} form a graph core on
$\operatorname{Dom}B\cap\mathcal Z_0$ and a form core on
$X\cap\mathcal Z_0$.
\end{proposition}
\begin{proof}
For $n\in\mathcal V_0$, Theorem~\ref{rtc:null:kernel} gives $Bn=0$,
whence
\[
 \tg[n]=\ip{D_bn}{n}_{\Hh}\le\norm{D_b}\norm{n}_{\Hh}^2.
\]
Thus $P_0:\Hh_c\to X$ is bounded.  Its actual row is zero, so it is
bounded into the graph domain as well.  Subtraction proves the asserted
preservation and boundedness for $P_*$.  Since
$q_0(n,u)=\ip{Bn}{u}_{\Hh}=0$ for every $u\in X_c$, orthogonality
proves~\eqref{rtc:null:deflation-identity}.
Finally $\mathcal N\subseteq\mathcal V_0$ gives
$\mathcal Z_0\subseteq\mathcal K$.
Approximate a vector in either stated domain by finite full-jet sums in
the corresponding norm, using Theorem~\ref{thm:actual-core}, and apply
$P_*$.  Its boundedness in that norm and $P_*u=u$ prove the two core
claims.
\end{proof}
For a finite full-jet list $e_1,\ldots,e_n$, the standard coefficient
matrices for these projected vectors are
\[
 M_{ij}=q_0(e_j,e_i),\qquad
 G^*_{ij}=\ip{P_*e_j}{P_*e_i}_{\Hh}.
\]
The form matrix stays $M$, whereas the norm matrix changes to $G^*$.
Thus, for $\delta\ge0$, positivity of $M+\delta G^*$ for every finite
list is equivalent, by the proved form core, to $B|_{\mathcal Z_0}
+\delta I\ge0$; on $\mathcal V_0$ the shifted form is already
$\delta\norm{\cdot}_{\Hh}^2$.
This requires all finite lists, not a finite sample.  The projection has
not been made effective, and removing the kernel gives neither a
positive sign nor a positive spectral gap.  The resolvent constructions
below avoid requiring an explicit formula for this full-kernel
projection.

We return to the centered notation $B=B_c$ and $q=q_0|_{X_c}$ for
the remaining response constructions.

\subsection{Fixed inverses and kernel removing filters}\label{rtc:sec:cores}
\begin{proposition}[Fixed inverse cores and residual costs]\label{rtc:prop:fixed-core}
Let $Z$ be a reducing subspace of $B$, and let $E$ have dense linear span
in $Z$ in the Hilbert norm. If $\alpha>0$ and
$B+\alpha I\ge c_\alpha I>0$, then
$R_\alpha E$, where $R_\alpha=(B+\alpha I)^{-1}$, spans a graph core
for $B|_Z$ and a form core for $X_c\cap Z$.
For $r=R_\alpha y$ and an admitted trial $\widetilde r\in\operatorname{Dom}B$,
put $e=y-(B+\alpha)\widetilde r$. Then
\begin{align}\label{rtc:eq:fixed-error}
 |q_\sigma[r]-q_\sigma[\widetilde r]|
 &\le {2\over c_\alpha}\norm{(B+\sigma)\widetilde r}_{\Hh}
                 \norm e_{\Hh}
 +\left({1\over c_\alpha}+{|\alpha-\sigma|\over c_\alpha^2}\right)
                 \norm e_{\Hh}^2,\\
 q_\sigma[r]&\ge q_\sigma[\widetilde r]
 -t\norm{(B+\sigma)\widetilde r}_{\Hh}^2
 -{t^{-1}+\alpha-\sigma\over c_\alpha^2}\norm e_{\Hh}^2
 \quad(0\le\sigma\le\alpha,\ t>0).\label{rtc:eq:fixed-lower}
\end{align}
For responses $r_i=R_\alpha y_i$, the exact form matrix is
$M_{ij}=\ip{y_j}{r_i}_{\Hh}-(\alpha-\sigma)\ip{r_j}{r_i}_{\Hh}$.
\end{proposition}
\begin{proof}
For $v\in\operatorname{Dom}B\cap Z$, approximate $(B+\alpha)v$ by
$y_n\in\operatorname{span}E$. Boundedness of $R_\alpha$ and
$BR_\alpha=I-\alpha R_\alpha$ proves graph convergence of $R_\alpha y_n$
to $v$. Bounded spectral truncations make the operator domain a form core;
graph convergence implies shifted-form convergence, proving the second core
claim by a diagonal approximation. For the residual statements,
$h=r-\widetilde r=R_\alpha e$ satisfies
$\norm h_{\Hh}\le c_\alpha^{-1}\norm e_{\Hh}$ and
$0\le q_\alpha[h]=\ip{e}{R_\alpha e}_{\Hh}\le c_\alpha^{-1}\norm e_{\Hh}^2$.
Expand $q_\sigma[\widetilde r+h]$ and use
$q_\sigma[h]=q_\alpha[h]-(\alpha-\sigma)\norm h_{\Hh}^2$.
Cauchy--Schwarz gives the absolute estimate; discarding $q_\alpha[h]$
and applying $2xy\le tx^2+t^{-1}y^2$ gives the lower estimate.
The matrix formula follows from $Br_j=y_j-\alpha r_j$.
\end{proof}
For $\alpha=1$, the unconditional choice is $c_1=1/2$.
The constants remain bounded as $\sigma\downarrow0$; this removes a
solve-conditioning loss but supplies no signed margin.
In particular, Theorem~\ref{rtc:null:kernel} makes the value waves $J_w$
a Hilbert-total family in $\mathcal Z_0$; their $R_\alpha$ images therefore
form an actual graph core there without requiring derivative jets.

\begin{proposition}[Bounded row alternatives]\label{rtc:prop:row-free}
For the unconditional prior $B\ge-I/2$, $0\le\sigma\le1$, put
$b_\sigma=\max\{2-4\sigma,0\}$. A single fixed-shift residual
$e=y-A\widetilde r$ satisfies, for every $t>0$,
\begin{align}\label{rtc:eq:row-free-single}
 q_\sigma[Ry]&\ge q_\sigma[\widetilde r]
 -t\norm{(B+\sigma)\widetilde r}_{\Hh}^2
 -(4/t+b_\sigma)\norm e_{\Hh}^2,\notag\\
 q_\sigma[Ry]&\ge q_\sigma[\widetilde r]
 -t\norm{\widetilde r}_{\Hh}^2
 -(1/t+b_\sigma)\norm e_{\Hh}^2,\\
 |q_\sigma[Ry]-q_\sigma[\widetilde r]|
 &\le2\norm{\widetilde r}_{\Hh}\norm e_{\Hh}
               +2\norm e_{\Hh}^2.\notag
\end{align}
The lower constant $b_\sigma$ for $q_\sigma[Re]$ is optimal using
only $0<R\le2I$, over that abstract class.
\end{proposition}
\begin{proof}
The error is $Re$, and $q_\sigma[Re]$ has multiplier
$p_\sigma(r)=r-(1-\sigma)r^2$ on $[0,2]$.
Concavity gives its minimum $\min\{0,4\sigma-2\}=-b_\sigma$.
The endpoint $r=2$ or scalar operators with $r\downarrow0$ give sharpness.
Its absolute value is at most $2$. Furthermore
$(B+\sigma)R=I-(1-\sigma)R$ has norm at most $1$.
The form cross term is bounded either by
$4\norm{(B+\sigma)\widetilde r}_{\Hh}\norm e_{\Hh}$
or, after moving the bounded resolvent, by
$2\norm{\widetilde r}_{\Hh}\norm e_{\Hh}$.
Young's inequality gives both lower bounds, and the latter cross estimate
with $|p_\sigma|\le2$ gives the absolute bound.
\end{proof}
The first lower bound improves the corresponding unconditional row-based
coefficient in~\eqref{rtc:eq:fixed-lower}. The second is an alternative,
not a uniform comparison, because its trial Gram replaces the row Gram.

\begin{theorem}[A zero independent filtered core]\label{rtc:thm:filter}
Assume $B\ge-dI$, $0\le d\le1/2$, and put $\beta=(1-d)^{-1}$,
$Z=(\ker B)^\perp$, and $K_f=R(I-R)$.
Then
\begin{gather}\label{rtc:eq:filter}
 \ker K_f=\ker B,\qquad BK_f=(I-R)^2,\qquad
 \norm{K_f}\le\max\{1/4,\beta(\beta-1)\},\quad\norm{BK_f}\le1.
\end{gather}
For any Hilbert-total source family $(e_n)$ in $\Hh_c$, the span of
$K_fe_n$ is a graph core on $Z$ and a form core on $X_c\cap Z$.
The bounded operators
\begin{equation}\label{rtc:eq:filter-matrices}
 \mathsf G=R^2(I-R)^2,\qquad
 \mathsf M_\sigma=R(I-R)^2[I-(1-\sigma)R]
\end{equation}
satisfy $q_\sigma(K_fy,K_fz)=\ip{\mathsf M_\sigma y}{z}_{\Hh}$ and
$\ip{K_fy}{K_fz}_{\Hh}=\ip{\mathsf Gy}{z}_{\Hh}$.
For $\sigma\ge0$, $\mathsf M_\sigma\ge0$ if and only if $B+\sigma I\ge0$.
\end{theorem}
\begin{proof}
At $\lambda\ge-d$, set $r=(1+\lambda)^{-1}\in(0,\beta]$.
The filter and its row have multipliers $r(1-r)$ and $(1-r)^2$.
These give all bounds and actual operator-domain admission, and the first
multiplier vanishes only at $\lambda=0$. For
$v\in\operatorname{Dom}B\cap Z$, truncate its spectrum to
$\eta\le|\lambda|\le M$. These truncations converge in graph norm as
$\eta\downarrow0$, $M\uparrow\infty$. On each retained set $1/[r(1-r)]$
is bounded; hence the truncation is $K_fy$ for some $y\in\Hh_c$.
Approximate that $y$ by source combinations and use the graph bounds in
\eqref{rtc:eq:filter}. Diagonal approximation proves graph density; the
spectral-truncation form-core argument of Proposition~\ref{rtc:prop:fixed-core}
then proves form density.
The two bounded multipliers in \eqref{rtc:eq:filter-matrices} are
$\lambda^2/(1+\lambda)^4$ and
$\lambda^2(\lambda+\sigma)/(1+\lambda)^4$.
This proves the identities. Positivity on the filtered core extends to
$X_c\cap Z$; on $\ker B$ the remaining form is
$\sigma\norm{\cdot}_{\Hh}^2$ with zero cross terms. This proves the
reverse sign implication, and the forward one is immediate.
\end{proof}
A concrete total source family is
$e_n=C(\sqrt W h_{2n})$, with normalized even Hermite functions $h_{2n}$
and the exact projection $Cu=u-2m(u)\kappa$.
Multiplication by $\sqrt W$ is unitary from ordinary even $L^2$ onto
$\Hh$, and $C$ maps a total family to a total family in $\Hh_c$.
For completeness of the Hermites, orthogonality to all of them makes every
derivative at zero of
$\int f(x)e^{-x^2/2}e^{zx}\,dx$ vanish. Gaussian Cauchy--Schwarz makes
this function entire; Fourier uniqueness then gives $f=0$.
Centering need not preserve independence or orthogonality. It is not kernel
removal: in the ambient realization $B\kappa=-\kappa/2$ and
$K_f\kappa=-2\kappa$.

\begin{proposition}[Two complete residuals for the filter]\label{rtc:prop:filter-error}
For $y\in\Hh_c$ and $z,v\in\operatorname{Dom}B$, let
$e_1=y-Az$, $e_2=z-Av$, and $g=z-v$. Then
\begin{align}\label{rtc:eq:filter-error}
 h:=K_fy-g&=K_fe_1-Re_2,&
 Bh&=(I-R)^2e_1-(I-R)e_2,\\
 \norm h_{\Hh}^2&\le L_d(\norm{e_1}_{\Hh}^2+\norm{e_2}_{\Hh}^2),&
 \norm{Bh}_{\Hh}&\le\norm{e_1}_{\Hh}+\norm{e_2}_{\Hh},\notag
\end{align}
where $L_d=\beta^2[1+(\beta-1)^2]$.
For $0\le\sigma\le d$ and $t>0$,
\begin{equation}\label{rtc:eq:filter-gate}
 q_\sigma[K_fy]\ge q_\sigma[g]-t\norm{(B+\sigma)g}_{\Hh}^2
 -L_d(t^{-1}+d-\sigma)(\norm{e_1}_{\Hh}^2+\norm{e_2}_{\Hh}^2).
\end{equation}
\end{proposition}
\begin{proof}
The identities $Ry=z+Re_1$ and $R^2y=v+Re_2+R^2e_1$ give the errors.
The squared row norm of $[K_f,-R]$ is $\norm{K_f^2+R^2}$, bounded by
$L_d$ because $r^2[(1-r)^2+1]$ has derivative
$2r(2r^2-3r+2)>0$ for $r>0$. The graph bound follows from
$\norm{I-R}\le1$. Expand the form at $g$, retain
$q_\sigma[h]\ge-(d-\sigma)\norm h_{\Hh}^2$, and use Young's inequality.
\end{proof}
For linear trial maps these are matrix inequalities for every coefficient
combination, including residual correlations. If the trial form has margin
$\gamma\norm g_{\Hh}^2$, the trial row is bounded by
$\Lambda\norm g_{\Hh}$, and the sum of residual squares by
$\tau^2\norm g_{\Hh}^2$, minimizing over $t$ gives the sufficient condition
$\gamma\ge2\Lambda\sqrt{L_d}\tau+L_d(d-\sigma)\tau^2$.
The minimum includes zero-factor limiting cases. None of these relative
premises follows merely from small absolute errors.
There is also a bounded-row filter alternative. For the same residuals,
$0\le\sigma\le d\le1/2$, and $\mathfrak r=\norm{e_1}_{\Hh}^2+
\norm{e_2}_{\Hh}^2$, define
\[
 C_\sigma=\max_{0\le r\le\beta}
 [1-(1-\sigma)r]^2[(1-r)^2+1]\le2.
\]
The exact row in~\eqref{rtc:eq:filter-error} gives
$\norm{(B+\sigma)h}_{\Hh}^2\le C_\sigma\mathfrak r$ by its squared
row-operator norm. Move the form cross term to this error row and use
$q_\sigma[h]\ge-(d-\sigma)L_d\mathfrak r$. Young's inequality then proves
\begin{equation}\label{rtc:eq:row-free-filter}
 q_\sigma[K_fy]\ge q_\sigma[g]-t\norm g_{\Hh}^2
       -[C_\sigma/t+L_d(d-\sigma)]\mathfrak r.
\end{equation}
The bound $C_\sigma\le2$ uses $\beta\le2$ and
$|1-r|,|1-(1-\sigma)r|\le1$ on the indicated interval.
If $q_\sigma[g]\ge\gamma\norm g_{\Hh}^2$ and
$\mathfrak r\le\tau^2\norm g_{\Hh}^2$, optimization gives the sufficient
condition $\gamma\ge2\sqrt{C_\sigma}\tau+L_d(d-\sigma)\tau^2$,
without a separate trial-row bound. All these statements apply as full
coefficient matrix inequalities, not just columnwise estimates.
At $\lambda=-b<0$, the unshifted filtered form and norm multipliers are
$-b^3/(1-b)^4$ and $b^2/(1-b)^4$; their quotient remains $-b$.
More generally an analytic filter with leading term $c\lambda^s$ gives
form order $2s+1$ and norm order $2s$, not a smaller normalized deficit.
For nested source prefixes $E_m$, the optimal tested deficit
\[
 d_m=\max\left\{0,-\inf_{y\in E_m,\ K_fy\ne0}
                   {q[K_fy]\over\norm{K_fy}_{\Hh}^2}\right\}
\]
(with value zero if $K_fE_m=\{0\}$) is nondecreasing, since the infimum
is over an expanding set. Enlarging tested sources does not by itself
make a deficit decrease.

\subsection{A single endpoint theorem for two response penalties}\label{rtc:sec:penalties}
For admitted responses $z,v$ define
\begin{equation}\label{rtc:eq:two-base}
 e_1=y-Az,\quad e_2=z-Av,\qquad
 J_\sigma=q_\sigma[z]+2\operatorname{Re}\ip{z-kv}{e_1}_{\Hh}.
\end{equation}
Throughout this subsection $a\in[1/2,1]$, $A\ge aI$, and
$0\le\sigma\le1/2$.
\begin{theorem}[All constant residual penalties]\label{rtc:thm:endpoint}
For a real symmetric $2\times2$ matrix $C_0$, the bound
\begin{equation}\label{rtc:eq:general-penalty}
 q_\sigma[Ry]\ge J_\sigma-
 \sum_{i,j=1}^2(C_0)_{ij}\ip{e_j}{e_i}_{\Hh}
\end{equation}
holds for every self-adjoint $A\ge aI$ and every admitted source and
response pair if and only if
\begin{equation}\label{rtc:eq:endpoints}
 C_0\ge0,\qquad C_0+F_k(a^{-1})\ge0,\qquad
 F_k(r)=\begin{pmatrix}r-kr^2&-kr\\-kr&0\end{pmatrix}.
\end{equation}
\end{theorem}
\begin{proof}
Expanding $Ry=z+Re_1$ and $Rz=v+Re_2$ gives the exact identity
\begin{equation}\label{rtc:eq:signed-remainder}
 q_\sigma[Ry]-J_\sigma
 =\ip{(R-kR^2)e_1}{e_1}_{\Hh}
       -2k\operatorname{Re}\ip{Re_2}{e_1}_{\Hh}.
\end{equation}
No second operator row is needed. Set $M=a^{-1}$. For $0\le r\le M$,
\[
 C_0+F_k(r)=(1-r/M)C_0+(r/M)[C_0+F_k(M)]
                  +kr(M-r)\begin{pmatrix}1&0\\0&0\end{pmatrix}.
\]
The two endpoints therefore suffice by spectral calculus on the pair of
residuals. Conversely, $A=r^{-1}$ on a one-dimensional space realizes
arbitrary residual pairs by $z=0$, $y=e_1$, $v=-re_2$.
Thus universal validity forces $C_0+F_k(r)\ge0$ for $0<r\le M$;
taking $r\downarrow0$ and $r=M$ yields both endpoints.
\end{proof}
The optimality assertions here are over this abstract operator class, not
over the spectrum of the theta operator. A least-squares residual is
orthogonal to $AF$, not generally to $F$; hence the signed correction in
$J_\sigma$ cannot be dropped. For a first Galerkin response and both trials
in the same $F$, $e_1\perp F$ does justify dropping it.

\begin{corollary}[Sharp rank one and unequal penalties]\label{rtc:cor:penalties}
For $t>0$ such that $D_t=t^2(a-k)+2akt>0$, put
$\alpha_t=k^2/D_t$. Then the smallest universal nonnegative coefficient
in front of $\norm{e_1+te_2}_{\Hh}^2$ is $\alpha_t$.
In particular the sum coefficient is
\begin{equation}\label{rtc:eq:chi}
 \alpha_1={k^2\over a(1+2k)-k}
 ={(1-\sigma)^2\over 2-3d-(1-2d)\sigma}.
\end{equation}
If $0<d<1/2$ and $0\le\sigma<d$, write $\epsilon=d-\sigma>0$.
The choice $t_*=ak/\epsilon$ gives the valid penalty
\begin{equation}\label{rtc:eq:unequal}
 {1\over\epsilon}\norm{(\epsilon/a)e_1+ke_2}_{\Hh}^2.
\end{equation}
Every valid $C_0$ has $(C_0)_{11}\ge\epsilon/a^2$; equality forces
$(C_0)_{12}=k/a$ and $(C_0)_{22}\ge k^2/\epsilon$.
\end{corollary}
\begin{proof}
For $C_0=\alpha(1,t)^T(1,t)$, the second endpoint determinant is
$M\{\alpha[t^2(1-kM)+2kt]-k^2M\}$.
Its positive lower-right diagonal makes positivity equivalent to
$\alpha\ge k^2/D_t$. The necessity argument in
Theorem~\ref{rtc:thm:endpoint} proves sharpness. Substitution of $t_*$
gives \eqref{rtc:eq:unequal}. The first diagonal endpoint gives the lower
bound on $(C_0)_{11}$. If equality holds, a positive semidefinite matrix
with zero first diagonal has zero first row, forcing $(C_0)_{12}=k/a$.
Finally $\det C_0\ge0$ gives the last inequality, attained by
\eqref{rtc:eq:unequal}.
\end{proof}
The unequal penalty minimizes the first cost, not both costs. It has a
$k^2/\epsilon$ second cost and is used only for $\sigma<d$.
If $\norm{e_2}_{\Hh}\le\theta\epsilon(ak)^{-1}\norm{e_1}_{\Hh}$,
its cost is at most $(1+\theta)^2\epsilon a^{-2}\norm{e_1}_{\Hh}^2$;
for source spaces this premise must hold for every combination.
At the unconditional endpoint $a=1/2$, \eqref{rtc:eq:chi} is $2k^2$.
There is then the useful exact positive remainder
\begin{align}\label{rtc:eq:coarse-remainder}
 q_\sigma[Ry]-J_\sigma+2k^2\norm{e_1+e_2}_{\Hh}^2
 &=k^2\ip{(2I-R)(e_1+e_2)}{e_1+e_2}_{\Hh}\notag\\
 &\quad+\ip{R(\sigma e_1-ke_2)}{\sigma e_1-ke_2}_{\Hh}
 +k\ip{R(2I-R)e_1}{e_1}_{\Hh}.
\end{align}
It follows by expanding the right side and using $k+\sigma=1$.
For general $d$, set $b=k\beta$, $c=(d-\sigma)\beta^2$ and
$D_*=2b-c$ in the range $0\le\sigma\le d$. The sharp sum remainder is
\[
 \alpha_1\ip{(I-R/\beta)(e_1+e_2)}{e_1+e_2}_{\Hh}
 +{1\over\beta D_*}\ip{Rw}{w}_{\Hh}
 +k\ip{R(\beta I-R)e_1}{e_1}_{\Hh},\qquad
 w=(b-c)e_1-be_2.
\]
Expansion proves this equivalent factorization, with
$\alpha_1=b^2/D_*$. It exhibits equality at $B=-d$ and
$e_2=(b-c)e_1/b$, so the original coupled-residual sharpness is retained.

\subsection{Complete row matrices and monotone response optimization}\label{rtc:sec:optimization}
Let $\Phi:\C^n\to\operatorname{Dom}B$ have independent centered columns.
Set
\begin{equation}\label{rtc:eq:full-grams}
 G=\Phi^*\Phi,\qquad H=\Phi^*A\Phi,\qquad
 L=(A\Phi)^*(A\Phi),\qquad
 D=L-HG^{-1}H.
\end{equation}
The adjoints use the weighted inner product. Only first rows are required;
$\operatorname{Dom}B^2$ is not assumed. If $P_F$ projects onto
$F=\operatorname{Ran}\Phi$, then
\begin{equation}\label{rtc:eq:full-leakage}
 D=((I-P_F)A\Phi)^*((I-P_F)A\Phi)\ge0.
\end{equation}
Indeed $P_F=\Phi G^{-1}\Phi^*$ and expansion gives
\eqref{rtc:eq:full-leakage}. It contains the complete omitted row,
including weighted spatial tails, mean corrections, compensated Gamma rows,
all prime powers and every cross term. Replacing $L$ by $HG^{-1}H$
sets this genuine cost to zero without justification.

For sources $y=\Phi Sx$ and responses $z=\Phi Ux$, $v=\Phi Vx$, let
$X=\binom UV$, $G_s=S^*GS$ and use the rank-one penalty
$\alpha\norm{e_1+te_2}_{\Hh}^2$, with $\alpha=\alpha_t$.
Expansion yields the lower matrix
\begin{align}\label{rtc:eq:two-matrices}
 \Gamma(X)&=-\alpha G_s+X^*b+b^*X-X^*KX,\notag\\
 K_{11}&=H+kG+\alpha(L-2tH+t^2G),\notag\\
 K_{12}&=\alpha tL-(\alpha t^2+k)H,\qquad K_{22}=\alpha t^2L,\\
 b&=\binom{[(1-\alpha t)G+\alpha H]S}{[-kG+\alpha tH]S},\qquad
 K=\begin{pmatrix}K_{11}&K_{12}\\K_{12}^*&K_{22}\end{pmatrix}.\notag
\end{align}
Equivalently its unpenalized term is
$U^*(H-kG)U+(U-kV)^*(GS-HU)+(GS-HU)^*(U-kV)$, and its residual Gram is
\[
 (S+tU)^*G(S+tU)-(S+tU)^*H(U+tV)
 -(U+tV)^*H(S+tU)+(U+tV)^*L(U+tV).
\]
These expansions verify every signed term of \eqref{rtc:eq:two-matrices}.
For an independent source map $Y:\C^j\to\Hh_c$, no inclusion in
$\operatorname{Ran}\Phi$ is needed: take $G_s=Y^*Y$, keep $K$ unchanged,
and replace the coupling by
\begin{equation}\label{rtc:eq:two-rectangular}
 b=\binom{(1-\alpha t)\Phi^*Y+\alpha(A\Phi)^*Y}
              {-k\Phi^*Y+\alpha t(A\Phi)^*Y}.
\end{equation}
Expansion of the same complete residual objective gives
$\Gamma(X)=-\alpha G_s+X^*b+b^*X-X^*KX$ again.
Consequently the optimum, range and refinement arguments below also hold
for arbitrary Hilbert source columns; no operator-domain admission of those
source columns is inferred or required.

\begin{theorem}[Finite optimum and refinement]\label{rtc:thm:opt}
The matrix $K$ is positive semidefinite. With $h=1+k/(\alpha t^2)$,
\begin{equation}\label{rtc:eq:two-factor}
 K[(z,v)]=\alpha t^2\norm{Av+(A/t-hI)z}_{\Hh}^2
 +(1+2k/t)\ip{(A-aI)z}{z}_{\Hh}.
\end{equation}
The columns of $b$ lie in $\operatorname{Ran}K$, and
\begin{equation}\label{rtc:eq:two-opt}
 \widehat\Gamma=-\alpha G_s+b^*K^\dagger b,\qquad
 \widehat\Gamma-\Gamma(X)=(X-K^\dagger b)^*K(X-K^\dagger b)\ge0.
\end{equation}
At fixed operator, source coordinates, shift and penalty, exact optima
increase under nested response refinement. If enlarged positive definite
matrices have common coordinates
$K_m=\left(\begin{smallmatrix}K_n&J\\J^*&D_0\end{smallmatrix}\right)$,
$b_m=\binom{b_n}{f}$, put
$W_0=D_0-J^*K_n^{-1}J$ and $r=f-J^*K_n^{-1}b_n$. Then
\begin{equation}\label{rtc:eq:gain}
 \widehat\Gamma_m-\widehat\Gamma_n=r^*W_0^{-1}r\ge0.
\end{equation}
\end{theorem}
\begin{proof}
The square in \eqref{rtc:eq:two-factor} has the stated off-diagonal
block. Subtracting its upper-left block from $K_{11}$ leaves
$(1+2k/t)H-(k+k^2/(\alpha t^2))G$.
The identity $k+k^2/(\alpha t^2)=a(1+2k/t)$ proves the factorization.
The residual bound puts a finite upper bound on $\Gamma(X)$ for each
source. A nonzero component of $bx$ in $\ker K$ would give an unbounded
linear term along that component. Thus $bx\in\operatorname{Ran}K$.
Finite square completion proves \eqref{rtc:eq:two-opt}, simultaneously
for all sources. The old response pairs remain feasible in the larger
space, proving monotonicity. Finally the triangular substitution
$w_{\rm old}\mapsto w_{\rm old}+K_n^{-1}Jw_{\rm new}$ diagonalizes
the block quadratic into $K_n$ and $W_0$; completing both squares gives
\eqref{rtc:eq:gain}. Positive definiteness of $K_m$ gives $W_0>0$.
\end{proof}
Ordinary inverses in this theorem require positive definiteness; the prior
$A\ge aI$ alone gives only semidefiniteness. At the coarse choice
$a=1/2$, $t=1$, $\alpha=2k^2$, eliminating $K_{22}$ gives
\begin{equation}\label{rtc:eq:reduced-hessian}
 K_{11}-K_{12}K_{22}^{-1}K_{12}^*
 =(1+2k)\{E+(k+1/2)\Delta\},\quad
 E=H-G/2,\quad\Delta=G-HL^{-1}H\ge0.
\end{equation}
Here $L>0$ because $A\ge I/2$, and
$\Delta=((I-P_{AF})\Phi)^*((I-P_{AF})\Phi)$ by the projection formula
$P_{AF}=A\Phi L^{-1}(A\Phi)^*$. This proves its sign and, by expanding
the eliminated product, \eqref{rtc:eq:reduced-hessian}.
It is a different omitted-component Gram from $D$ in
\eqref{rtc:eq:full-leakage}. A strictly stronger prior $B\ge-dI$,
$d<1/2$, makes $E\ge(1/2-d)G>0$ and hence makes this unchanged coarse
Hessian positive definite. That strictness does not automatically carry
over after changing its penalty to the sharp prior-dependent coefficient.
For example the sharp objective may have null pairs when $A=aI$.

If $0<W_0\le\overline W_0$, inverse order gives the lower gain
$r^*W_0^{-1}r\ge r^*\overline W_0^{-1}r$.
Thus an upper Schur bound, not a lower one, can pay a signed deficit.
A stronger prior with $a>1/2$ improves the retained coarse sum objective
by exactly $(2k^2-\alpha_1)$ times the complete Gram of $e_1+e_2$.
Its coefficient is nonnegative by \eqref{rtc:eq:chi}; whether the gain
pays the old deficit is a separate full-matrix inequality.

For the coarse objective only, write $\Gamma_\sigma$ for $t=1$,
$\alpha=2k^2$, set $E_1=GS-HU$, and let $W_s$ be the complete Gram of
$e_1+e_2$. Holding the responses fixed and lowering the shift by
$0\le s\le\sigma$ gives the exact identity
\begin{equation}\label{rtc:eq:two-shift}
 \Gamma_{\sigma-s}(X)=\Gamma_\sigma(X)
 -s\{U^*GU+V^*E_1+E_1^*V+(4k+2s)W_s\}.
\end{equation}
It follows by replacing $k$ by $k+s$ in the signed correction and penalty.
The middle terms are signed; the brace is not asserted positive.
An upper bound $\nu I$ on that brace and a lower margin $\mu I$ on
$\Gamma_\sigma$ certify the smaller shift if $\mu-s\nu\ge0$.
Source enlargement is independent: an enlarged Hermitian block
$\left(\begin{smallmatrix}P&C_1\\C_1^*&D_1\end{smallmatrix}\right)$ with
$P>0$ is nonnegative exactly when $D_1-C_1^*P^{-1}C_1\ge0$, by
completion of its first square. This block must be evaluated at the new
shift. No monotonicity theorem supplies its sign.

For a fixed rational $X_0$, the direct trial matrix avoids certifying an
exact inverse. If Hermitian approximations to $G_s,b,K$ have errors
$\eta_G,\eta_b,\eta_K$, then
\begin{equation}\label{rtc:eq:two-round}
 \Gamma(X_0)\ge\widetilde\Gamma(X_0)-\rho I,\qquad
 \rho=\alpha\eta_G+2\norm{X_0}\eta_b+\norm{X_0}^2\eta_K.
\end{equation}
Expand the difference in \eqref{rtc:eq:two-matrices} and apply the matrix
operator-norm triangle inequality. If $K>0$, its exact optimization loss
is $r_0^*K^{-1}r_0$, $r_0=b-KX_0$; a lower $0<K_*\le K$ gives an
upper loss $r_0^*K_*^{-1}r_0$. These inverse directions serve different
purposes. Exact-optimum monotonicity need not hold for independently
rounded enclosures or for successive, nonoptimal least-squares choices.
Old responses and valid enclosures remain usable by zero extension.

\subsection{Exact null response repair}\label{rtc:sec:null-repair}
In this subsection $0<d<1/2$, $0\le\sigma<d$ and
$\epsilon=d-\sigma>0$. Write $g_\sigma$ for $J_\sigma$ minus the unequal
penalty \eqref{rtc:eq:unequal}, and
$E_a[w]=\ip{(A-aI)w}{w}_{\Hh}$ on $\operatorname{Dom}A$.
\begin{proposition}[The exact loss and its removable null part]\label{rtc:prop:null-loss}
For $u=Ry$,
\begin{equation}\label{rtc:eq:unequal-loss}
 q_\sigma[u]-g_\sigma(y;z,v)
 ={2k-a\over a}E_a[u-z]
 +{k^2\over\epsilon}
 \norm{e_2+{\epsilon\over k}(a^{-1}I-R)e_1}_{\Hh}^2.
\end{equation}
For a finite-rank weighted orthogonal projection $P$ into $\ker B$, set
$Q_P=I-P$ and
$z^+=Q_Pz+Py$, $v^+=Q_Pv+Py$. Then
\begin{align}\label{rtc:eq:null-gain}
 g_\sigma(y;z^+,v^+)-g_\sigma(y;z,v)
 &={(2k-a)d\over a}\norm{Pe_1}_{\Hh}^2\notag\\
 &\quad+{k^2\over\epsilon}
 \norm{Pe_2+{\epsilon d\over ak}Pe_1}_{\Hh}^2\ge0.
\end{align}
Only overlaps are needed: $Pe_1=P(y-z)$, $Pe_2=P(z-v)$.
\end{proposition}
\begin{proof}
Add the unequal penalty to \eqref{rtc:eq:signed-remainder} and complete
the square in $e_2$. The remaining first-residual term is
$(2k-a)\ip{R(a^{-1}I-R)e_1}{e_1}_{\Hh}$.
Since $Re_1=u-z$, it equals $(2k-a)E_a[u-z]/a$, proving the loss.
Its second square vanishes at the admitted vector
$v_*=Rz+(\epsilon/k)R(a^{-1}I-R)e_1$; thus the best possible value for
fixed $z$ is exactly $q_\sigma[u]-(2k-a)E_a[u-z]/a$.
The kernel projection reduces $A$ and $R$ and satisfies $AP=RP=P$.
It preserves the domains, and the repaired residuals are $Q_Pe_1,Q_Pe_2$.
On its range, $R=I$ and $A-aI=dI$. Orthogonal splitting of
\eqref{rtc:eq:unequal-loss} therefore removes exactly
\eqref{rtc:eq:null-gain}. The overlap identities follow from $PA=P$.
\end{proof}
At $\sigma=0$ the gain is
$d(1+d)a^{-1}\norm{Pe_1}_{\Hh}^2+
 d^{-1}\norm{Pe_2+(d^2/a)Pe_1}_{\Hh}^2$.
It removes approximation loss but supplies no positive energy on a true
null direction. The actual derivative radicals from
Lemma~\ref{a:lem:radical} and Proposition~\ref{f:prop:halfeigen} are
admitted choices for $P$; their proofs are not repeated here.
The vector $\kappa$ is not one of these centered null vectors.

More generally, for any valid penalty matrix $C_0$, set
$\Omega=C_0+F_k(1)\ge0$, using Theorem~\ref{rtc:thm:endpoint}.
The same repair gain is
$\sum\Omega_{ij}\ip{Pe_j}{Pe_i}_{\Hh}$: the exact remainder
\eqref{rtc:eq:signed-remainder} on the kernel has $R=I$.
For the unequal and sharp balanced penalties respectively,
\[
 \Omega=\begin{pmatrix}\epsilon/a^2+\sigma&kd/a\\kd/a&k^2/\epsilon\end{pmatrix},
 \qquad
 \Omega=\begin{pmatrix}\alpha_1+\sigma&\alpha_1-k\\
                         \alpha_1-k&\alpha_1\end{pmatrix}.
\]
Let $N$ have independent known null columns, $G_N=N^*N>0$,
$P=NG_N^{-1}N^*$, and let $Y,Z,V$ be the source and two response maps.
With $T_1=N^*(Y-Z)$, $T_2=N^*(Z-V)$ and $T=\binom{T_1}{T_2}$, the
complete gain matrix is
\begin{equation}\label{rtc:eq:null-matrix}
 \mathcal G=T^*(\Omega\otimes G_N^{-1})T.
\end{equation}
This follows by substituting the orthogonal projection formula in the
preceding quadratic. For nested null projections, the new orthogonal
projection increment gives another nonnegative quadratic, proving exact
gain monotonicity for fixed responses and source coordinates.

There is an overlap-only augmentation of the complete matrices. Put
$C_N=\Phi^*N$. Since $AN=N$, all three augmented matrices have the same
new blocks:
\begin{equation}\label{rtc:eq:null-augmentation}
 G_{\rm aug}=\begin{pmatrix}G&C_N\\C_N^*&G_N\end{pmatrix},\quad
 H_{\rm aug}=\begin{pmatrix}H&C_N\\C_N^*&G_N\end{pmatrix},\quad
 L_{\rm aug}=\begin{pmatrix}L&C_N\\C_N^*&G_N\end{pmatrix}.
\end{equation}
Pairing the known rows with old columns proves these formulas, without
new high-derivative operator rows. Projecting the old trials by $I-P$
instead subtracts $C_NG_N^{-1}C_N^*$ from each of $G,H,L$; hence $H-G$
is unchanged. Dependent augmented columns must be removed or treated at
their exact rank. The full null Gram and overlaps still require enclosures.

For a rigorous lower gain take $G_N\le\overline G_N$, with
$\overline G_N>0$, and $K_0=\Omega\otimes\overline G_N^{-1}$.
If $\norm{K_0^{1/2}(T-\widetilde T)}\le\tau$, then for $0<\eta<1$,
\begin{equation}\label{rtc:eq:null-robust}
 \mathcal G\ge(1-\eta)\widetilde T^*K_0\widetilde T
                  -(\eta^{-1}-1)\tau^2I.
\end{equation}
Inverse order first gives $\mathcal G\ge T^*K_0T$.
Expanding this square at $\widetilde T$ and applying Young's inequality
to its cross term proves \eqref{rtc:eq:null-robust}.
The upper Gram bound is essential for this lower-gain direction.

If a same-domain approximation has
$0\le D_j=B_j^+-B\le h_jI$, null preservation can be restored by
\begin{equation}\label{rtc:eq:null-head}
 B_{j,P}^{\sharp}=(I-P)B_j^+(I-P)
                  =B+(I-P)D_j(I-P),\qquad
 0\le B_{j,P}^{\sharp}-B\le h_jI.
\end{equation}
The equality uses reduction by $P$; the bound follows by testing the
positive bounded $D_j$ on $(I-P)u$. Bounded perturbation preserves the
operator and form domains and every existing graph core.
This is not an assertion that $B_{j,P}^{\sharp}\le B_j^+$.
Its own entries must be enclosed: raw head rows need not obey $A_jN=N$,
and repaired compact columns may acquire noncompact radical tails.

\subsection{A rectangular one response inverse certificate}\label{rtc:sec:single}
Now assume $0<d<1/2$, $a=1-d$, and $0\le\sigma<d$.
The following construction tests $R^{1/2}y$, not $Ry$ or $K_fy$.
The spectral theorem gives
\begin{equation}\label{rtc:eq:inverse-target}
 q_\sigma[R^{1/2}y]=\ip{(I-kR)y}{y}_{\Hh},\qquad
 B+\sigma I\ge0\ \Longleftrightarrow\ I-kR\ge0.
\end{equation}
The map $R^{1/2}$ is onto $\operatorname{Dom}A^{1/2}=X_c$ and
$\norm{A^{1/2}R^{1/2}y}_{\Hh}=\norm y_{\Hh}$.
For a trial $z\in\operatorname{Dom}A$, put $e=y-Az$ and
\begin{equation}\label{rtc:eq:single-upper}
 U(y;z)=2\operatorname{Re}\ip y z_{\Hh}-\ip{Az}z_{\Hh}
                      +a^{-1}\norm e_{\Hh}^2,
 \qquad h_\sigma(y;z)=\norm y_{\Hh}^2-kU(y;z).
\end{equation}
\begin{theorem}[One response loss and finite optimum]\label{rtc:thm:single}
Let $C_A=[A(A-aI)]^{1/2}$, with
$\operatorname{Dom}C_A=\operatorname{Dom}A$. Then
\begin{equation}\label{rtc:eq:single-loss}
 U(y;z)-\ip{Ry}y_{\Hh}
 =\ip{(a^{-1}I-R)e}e_{\Hh}
 ={1\over a}\norm{C_A(Ry-z)}_{\Hh}^2\ge0.
\end{equation}
Let $Y:\C^j\to\Hh_c$ have independent source columns and
$\Phi:\C^n\to\operatorname{Dom}A$ independent response columns.
Use $G_s=Y^*Y$, the full $H,L$ of \eqref{rtc:eq:full-grams}, and set
$K_1=L-aH$, $b_1=((A-aI)\Phi)^*Y$. For $z=\Phi Xx$, $y=Yx$,
\begin{align}\label{rtc:eq:single-finite}
 U_X&=a^{-1}(G_s-b_1^*X-X^*b_1+X^*K_1X),\notag\\
 \widehat U&=a^{-1}(G_s-b_1^*K_1^\dagger b_1),\qquad
 \widehat\Theta_\sigma=G_s-k\widehat U,\\
 U_X-\widehat U&=a^{-1}(X-K_1^\dagger b_1)^*K_1
                            (X-K_1^\dagger b_1)\ge0.\notag
\end{align}
Here $K_1\ge0$, $\operatorname{Ran}b_1\subset\operatorname{Ran}K_1$,
and $Y^*RY\le\widehat U\le a^{-1}G_s$.
The optimal response is independent of the target shift, and its exact
finite sign criterion is
\begin{equation}\label{rtc:eq:single-payment}
 b_1^*K_1^\dagger b_1\ge{d-\sigma\over1-\sigma}G_s.
\end{equation}
\end{theorem}
\begin{proof}
Expand $\ip{A(Ry-z)}{Ry-z}_{\Hh}$ to obtain
$\ip{Ry}y_{\Hh}=2\operatorname{Re}\ip y z_{\Hh}-\ip{Az}z_{\Hh}
+\ip{Re}e_{\Hh}$. Subtraction gives the first equality.
For $e=A(Ry-z)$, spectral calculus gives the second. The graph weights
$1+\lambda^2$ and $1+\lambda(\lambda-a)$ are comparable on
$[a,\infty)$, proving the domain equality without assuming $A^2z$.
Expanding the complete residual gives the finite expression.
Moreover $K_1=(C_A\Phi)^*C_A\Phi$. A vector in its kernel gives a trial
with spectral support at $a$, so $(A-aI)\Phi c=0$, implying $c^*b_1=0$.
This proves the range condition. Square completion yields the optimum.
The loss identity gives the lower bound on $\widehat U$ and $X=0$ gives
the upper bound. Expanding $G_s-k\widehat U$ proves the payment condition.
\end{proof}
In particular source columns need not lie in $\operatorname{Dom}A$.
The projection interpretation is exact: for
$T_A=(I-aR)^{1/2}$ and the Hilbert projection $P_{C_AF}$ onto
$C_AF$, $F=\operatorname{Ran}\Phi$,
\begin{equation}\label{rtc:eq:captured}
 b_1^*K_1^\dagger b_1=Y^*T_AP_{C_AF}T_AY.
\end{equation}
Indeed $(C_A\Phi)^*T_AY=b_1$ by spectral calculus, and the finite
projection formula proves the identity. It does not supply the required
lower bound on the captured norm.

Nested response spaces decrease $\widehat U$ and increase
$\widehat\Theta_\sigma$, by inclusion of admissible trials. The block
square completion of Theorem~\ref{rtc:thm:opt} gives the exact gain
$(k/a)r^*W_0^{-1}r$ when the corresponding $K_1$ blocks are positive
definite. Holding responses fixed, or optimizing them, gives
\begin{equation}\label{rtc:eq:single-shift}
 \Theta_{\sigma-s}=\Theta_\sigma-sU\qquad(0\le s\le\sigma).
\end{equation}
The cost is positive definite on each independent finite source space
because $U\ge Y^*RY>0$. If $u_*$ is the largest generalized eigenvalue
of $(\widehat U,G_s)$, the smallest nonnegative certified shift is
$\max\{0,1-u_*^{-1}\}\le d$. This follows by diagonalizing the
Gram-normalized $\widehat U$; the endpoint $d$ merely recovers the prior.

Null correction $z^+=(I-P)z+Py$ has $e^+=(I-P)e$, so
\eqref{rtc:eq:single-loss} gives
\begin{equation}\label{rtc:eq:single-null}
 h_\sigma(y;z^+)-h_\sigma(y;z)
 ={kd\over a}\norm{P(y-z)}_{\Hh}^2.
\end{equation}
It uses the same full null Grams and overlap logic as
\eqref{rtc:eq:null-matrix}. For fixed rational $X$, finite-entry errors
$\eta_G,\eta_b,\eta_K$ in $G_s,b_1,K_1$ give the lower debit
\begin{equation}\label{rtc:eq:single-round}
 \Theta_\sigma(X)\ge\widetilde\Theta_\sigma(X)-\rho_1I,
 \quad \rho_1=|1-k/a|\eta_G+{k\over a}
                    (2\norm X\eta_b+\norm X^2\eta_K).
\end{equation}
Expansion of \eqref{rtc:eq:single-finite} proves this by the triangle
inequality. These errors include analytic tails and evaluation rounding.

\subsection{Internal energy versus complete leakage}\label{rtc:sec:leakage}
Throughout this subsection retain $0<d<1/2$, $0\le\sigma<d$,
$a=1-d$ and $k=1-\sigma$. Take the common source and response family
$Y=\Phi$ and use $G,H,L,D$ from \eqref{rtc:eq:full-grams}. Put
$Q=H-G$ and $V_0=H-aG=Q+dG$. Assume $V_0>0$; then
$K_1\ge aV_0>0$, since $(A-aI)A\ge a(A-aI)$.
\begin{proposition}[Exact finite leakage payment]\label{rtc:prop:payment}
For $0\le\sigma<d$, the common-space certificate is nonnegative exactly
when
\begin{equation}\label{rtc:eq:leak-payment}
 (d-\sigma)D\le a(Q+\sigma G)G^{-1}(Q+dG).
\end{equation}
\end{proposition}
\begin{proof}
Here $b_1=V_0$ and \eqref{rtc:eq:single-payment} says
$V_0K_1^{-1}V_0\ge((d-\sigma)/k)G$.
All matrices being inverted are positive definite, and $d-\sigma>0$.
Inverse order and congruence make this equivalent to
$(d-\sigma)K_1\le kV_0G^{-1}V_0$.
Use $K_1=HG^{-1}H+D-aH$ and $k-(d-\sigma)=a$ to subtract the two
sides, obtaining \eqref{rtc:eq:leak-payment}.
Its right product is Hermitian because it equals
$a[QG^{-1}Q+(d+\sigma)Q+\sigma dG]$.
\end{proof}
If $V_0$ is singular, a nonzero null coefficient gives a true
$A$-eigenvector at $a$, since $A-aI\ge0$. Its exact inverse target is
$(1-k/a)\norm{\Phi x}_{\Hh}^2<0$ for $\sigma<d$.
Thus singular endpoint sources cannot be discarded with a pseudoinverse
in the common-space sign test. The general rectangular theorem remains
valid, and enlarging responses need not enlarge the source space.
Nor does failure of this finite common-space test establish a negative
operator direction. For $A=\operatorname{diag}(1,100)$, $a=9/10$,
and $\Phi=(100,1)^T$, the complete data are
$G=10001$, $H=10100$, $L=20000$.
Substituting them into~\eqref{rtc:eq:single-finite} gives
$\widehat\Theta_0=-107811891/109100<0$, whereas the exact source target
$\Phi^*(I-R)\Phi=99/100>0$ and $B\ge0$.
This follows by a scalar inverse with $V_0=10991/10$, $K_1=10910$;
an independent response direction removes this approximation obstruction.

\begin{theorem}[Sharp explicit source constants]\label{rtc:thm:scalar}
Suppose $Q\ge\gamma G$, $0\le D\le\ell G$, where
$\gamma>-d$ and $\ell\ge0$. Then
\begin{align}\label{rtc:eq:mu}
 \widehat\Theta_\sigma&\ge\mu_\sigma G,&
 \mu_\sigma&={a(\gamma+\sigma)(\gamma+d)-(d-\sigma)\ell
                  \over a[(1+\gamma)(\gamma+d)+\ell]}.
\end{align}
On the inverse-source space $R^{1/2}\operatorname{Ran}\Phi$, and not
on the direct source space by this argument, one obtains
\begin{equation}\label{rtc:eq:beta}
 q[u]\ge\beta(\gamma,\ell)\norm u_{\Hh}^2,\qquad
 \beta={a\gamma(\gamma+d)-d\ell\over a(\gamma+d)+\ell}
       =\gamma-{(\gamma+d)\ell\over a(\gamma+d)+\ell}.
\end{equation}
Both constants are best possible given just this prior, internal lower
bound and complete-leakage upper bound over the abstract operator class.
\end{theorem}
\begin{proof}
Normalize by $G^{-1/2}$ and let
$J=G^{-1/2}HG^{-1/2}\ge(1+\gamma)I>aI$,
$D_0=G^{-1/2}DG^{-1/2}\le\ell I$.
Then the normalized $K_1$ is
$J(J-aI)+D_0\le J(J-aI)+\ell I$.
Inverse order followed by congruence by $J-aI$ gives the lower bound
$f(J)$ for $(J-aI)K_1^{-1}(J-aI)$, where
\[
 f(t)={(t-a)^2\over t(t-a)+\ell},\qquad
 f'(t)={(t-a)[a(t-a)+2\ell]\over[t(t-a)+\ell]^2}>0\quad(t>a).
\]
Only the upper comparison matrix commutes with $J$; no commutation of
$D_0$ and $J$ is used. Scalar spectral calculus gives
$f(J)\ge f(1+\gamma)I$. Substitution into
$\widehat\Theta_\sigma=a^{-1}[kV_0K_1^{-1}V_0-(d-\sigma)G]$
proves \eqref{rtc:eq:mu}.
For $u=R^{1/2}\Phi x$ the zero-shift bound is
$q[u]\ge\mu_0\norm{A^{1/2}u}_{\Hh}^2
=\mu_0(q[u]+\norm u_{\Hh}^2)$.
The denominator minus numerator of $\mu_0$ is
$a(\gamma+d)+\ell>0$, so $\mu_0<1$, even if it is negative.
Division by $1-\mu_0$ proves \eqref{rtc:eq:beta}.

For sharpness set $v=\gamma+d>0$, $h=1+\gamma$ and
\[
 A_* =\begin{pmatrix}h&\sqrt\ell\\\sqrt\ell&a+\ell/v\end{pmatrix},
 \qquad \Phi=e_1.
\]
Here $A_*-aI=ww^*/v\ge0$, $w=(v,\sqrt\ell)^T$,
and $G=1$, $H=h$, $D=\ell$, $K_1=hv+\ell$.
Direct inversion gives
$\ip{A_*^{-1}e_1}{e_1}=(av+\ell)/[a(hv+\ell)]$.
Thus the source target equals $\mu_\sigma$ and its transformed-source
Rayleigh quotient equals $\beta$. Equality holds also when $\ell=0$.
\end{proof}
The extremizer retains an unobserved $B_*=A_*-I$ eigenvalue $-d$.
Finite-source positivity therefore does not imply global positivity.
The formula gives $-d<\beta\le\gamma$;
nonnegativity is equivalent to $d\ell\le a\gamma(\gamma+d)$.
For $-d<\gamma<0$ this cannot hold, and for $\gamma=0$ it requires
$\ell=0$. At zero leakage $\beta=\gamma$.
The excluded endpoint $\gamma=-d$ can make the finite matrices $V_0$
and $K_1$ singular; the full operator $A$ remains invertible.

\begin{theorem}[Relative leakage preserving zero directions]\label{rtc:thm:relative}
Let $T=G^{-1/2}QG^{-1/2}\ge0$, $D_0=G^{-1/2}DG^{-1/2}$, and assume the
full matrix inequality
\begin{equation}\label{rtc:eq:relative-premise}
 D_0\le{a\theta\over d}T(T+dI)+\tau(T+dI)^2,
 \qquad 0\le\theta\le1,\quad\tau\ge0.
\end{equation}
Define $\alpha=1+\tau d$, $c=1+a\theta/d+\tau$, and
$r=\tau d^2/[a(1+\tau d)]$. Then
\begin{equation}\label{rtc:eq:relative-bound}
 G^{-1/2}\widehat\Theta_\sigma G^{-1/2}
 \ge(\sigma-kr)I+{k(1-\theta)\over\alpha}
                      T(\alpha I+cT)^{-1}.
\end{equation}
For $0<\sigma<d$, the sufficient scalar allowance is
$\tau\le a\sigma/[d(d-\sigma)]$.
Moreover on $R^{1/2}\operatorname{Ran}\Phi$,
\begin{equation}\label{rtc:eq:relative-physical}
 q[u]\ge-{\tau d^2\over a+\tau d}\norm u_{\Hh}^2.
\end{equation}
\end{theorem}
\begin{proof}
The normalized $K_1$ satisfies
$(I+T)(T+dI)+D_0\le(T+dI)(\alpha I+cT)$.
Both factors on the right are positive definite functions of $T$.
Inverse order and congruence by $T+dI$ yield, at zero shift, the lower
function
\[
 f_0(T)=[(1-\theta-\tau d/a)T-\tau d^2I/a](\alpha I+cT)^{-1}.
\]
The scalar identity
$f_0(t)=-r+(1-\theta)t/[\alpha(\alpha+ct)]$
proves the bound for $\sigma=0$.
Since $\widehat\Theta_\sigma=\sigma G+k\widehat\Theta_0$, it proves
\eqref{rtc:eq:relative-bound}. The scalar part is nonnegative exactly
when $a\sigma\ge\tau d(d-\sigma)$; division uses $\sigma<d$.
At zero shift, converting the lower value $-r$ as in the preceding proof
gives $-r/(1+r)=-\tau d^2/(a+\tau d)$.
\end{proof}
For $\tau=0$ the matrix term is
$k(1-\theta)T[I+(1+a\theta/d)T]^{-1}$ and retains the kernel of $T$.
Indeed domination by the right side of \eqref{rtc:eq:relative-premise}
then forces $D_0$ to kill that kernel. A weighted remainder can spend
$\tau d^2$ even there. An absolute remainder $\zeta I$ can be covered
by $\tau=\zeta/d^2$ because $(T+dI)^2\ge d^2I$; the factor $d^2$
cannot then be counted a second time as free improvement.
The premises are complete matrix inequalities, not diagonal tests.
If $Q\ge0$ were proved on a form-dense cofinal family, it would already
establish the desired sign, so that premise cannot conceal the missing
arithmetic argument. For sources $\Phi S$, all bounds restrict by
congruence with $S$, but imposing them on the entire response space can
be stronger than the rectangular criterion actually needs.

\subsection{Cofinal logic and independent response refinement}\label{rtc:sec:cofinal}
The logical limits are different from finite response convergence.
Let $(E_j)$ be nested finite-dimensional source spaces with Hilbert-dense
union in $\Hh_c$, and let $(F_n)\subset\operatorname{Dom}B$ be nested
finite-dimensional response spaces with graph-dense union.
Independent response refinement is allowed; neither $n=j$ nor a prescribed
complexity growth rate is assumed.

\begin{theorem}[Fixed operator completeness]\label{rtc:thm:cofinal}
The following statements hold.
\begin{enumerate}
\item If $\sigma_j\ge0$ tends to zero and $(e_n)$ is a fixed total source
family, then $B\ge0$ is equivalent to nonnegativity of the exact matrices
of $q_{\sigma_j}$ on $K_f\operatorname{span}\{e_0,\ldots,e_{m_j}\}$
for every $j$, whenever $m_j$ eventually exceeds every fixed index.
\item For each fixed source space and shift, the optimized two-response
matrices tend in matrix norm to the exact matrix of $q_\sigma[R\,\cdot]$,
and $\widehat U$ tends in matrix norm to the compression of $R$.
Here the two-response penalty is any fixed valid rank-one penalty from
Corollary~\ref{rtc:cor:penalties}, including the coarse endpoint.
\item For positive shifts $\sigma_j\to0$ in their respective admitted
ranges, $B\ge0$ is equivalent to
$\forall j\ \exists n_j$ such that the optimized lower matrix is
nonnegative on $E_j$, for either of these response schemes.
\item If $\varepsilon_j\ge0$ tends to zero and
$q[R^{1/2}y]\ge-\varepsilon_j\norm{R^{1/2}y}_{\Hh}^2$ for every
$y\in E_j$, then $q\ge0$ on $X_c$.
Thus finite bounds \eqref{rtc:eq:beta} suffice if their negative parts
tend to zero on such a nested dense source family.
\end{enumerate}
\end{theorem}
\begin{proof}
For the first assertion, fix a finite source combination $y$.
It belongs to every sufficiently late prefix, so letting $j\to\infty$
gives $q[K_fy]\ge0$. The filtered form core of
Theorem~\ref{rtc:thm:filter} extends this to $X_c\cap(\ker B)^\perp$;
the kernel contributes zero. The converse follows directly from $q\ge0$.
Monotone decrease of shifts is unnecessary, but nonnegative convergence
and eventual inclusion of each fixed source are necessary to this argument.

For the second assertion graph density approximates each fixed $Ry$ and
$R^2y$ by admitted trials. Their complete residuals tend to zero, their
trial forms converge, and the signed correction tends to zero. Consequently
the trial lower values converge to $q_\sigma[Ry]$; the optimum lies between
them and the exact value. Similarly \eqref{rtc:eq:single-loss} makes the
one-response trial upper values tend to $\ip{Ry}y_{\Hh}$ and squeezes
the optimum. Choose these approximations simultaneously for a finite source
basis. Linearity and polarization give convergence of all matrix entries,
hence finite-dimensional operator-norm convergence.

For the third assertion, if $B\ge0$, then
$q_{\sigma_j}[Ry]\ge\sigma_j\norm{Ry}_{\Hh}^2$ and
$I-(1-\sigma_j)R\ge\sigma_jR$.
Injectivity and positivity of $R$ give a strictly positive minimum on each
fixed finite source unit sphere. The second assertion therefore gives a
finite $n_j$ with a strictly positive lower matrix.
Conversely the certified lower matrices give the appropriate exact signs
on each $E_j$. On every retained source let $j\to\infty$.
For the two-response scheme this proves $q[Ry]\ge0$ on the dense source
union; Proposition~\ref{rtc:prop:fixed-core} makes its resolvent image a
graph core and then a form core. For the one-response scheme it proves
$\ip{(I-R)y}y_{\Hh}\ge0$ on that union. Boundedness of $I-R$ extends the
inequality to all $\Hh_c$, and \eqref{rtc:eq:inverse-target} proves $q\ge0$.
The fourth assertion uses the same last argument: for fixed $y$, the
right-hand deficit tends to zero, and $R^{1/2}$ maps the Hilbert-total
source union to a form core by its displayed isometry.
\end{proof}
A constructive choice in a common containing family is successive complete
least squares:
$U=L^{-1}HS$, $V=L^{-1}HU$. Both inverses are legitimate because
$H\ge aG>0$ and $L\ge a^2G>0$.
Since $A$ maps its graph domain onto $\Hh_c$, the union of $AF_n$ is
dense. The first residual tends to zero; writing its trial as
$z_n=Ry-Re_{1,n}$ gives graph convergence to $Ry$.
The second residual is bounded by
$\norm{z_n-Ry}_{\Hh}+\operatorname{dist}(Ry,AF_n)\to0$.
Thus these nonoptimal choices also prove the existence direction in the
two-index theorem. Their values need not be monotone.
For example let $A=I+ZZ^*/20$, $\sigma=1/100$, source $e_1$, and standard
prefix trials, where
\[
 Z=\begin{pmatrix}
 2&1&1&1&4\\3&-4&-1&-4&-1\\-1&-3&-2&-1&3\\
 -4&-3&4&1&-3\\0&-1&0&4&-4
 \end{pmatrix}.
\]
Substitution of $G=I_n$, $H=A_{1:n,1:n}$,
$L=(A^2)_{1:n,1:n}$ and the stated least-squares coefficients into
\eqref{rtc:eq:two-matrices}, with its coarse penalty, gives at $n=1,2$
respectively
\[
 -{20290748111431\over34863282857500},\qquad
 -{116578365688419098736203\over198931329099749364778600}.
\]
Cross multiplication makes the second strictly smaller, although $B\ge0$.
This exact finite example separates convergence of these responses from
monotonicity of the jointly optimized value.

For comparison, form Galerkin convergence alone is weaker than the graph
input just used. If $D_b$ is the bounded perturbation in the actual
realization and $b=\norm{D_b}$, then $q_1\ge\norm{\cdot}_{\Hh}^2/2$
and $\norm u_X^2=\norm u_{\Hh}^2+\tg[u]$ imply
\[
 (1+2b)^{-1}\norm u_X^2\le q_1[u]\le(1+b)\norm u_X^2.
\]
The lower estimate follows from
$\norm u_X^2=q_1[u]+\ip{D_bu}u_{\Hh}$ and
$\norm u_{\Hh}^2\le2q_1[u]$; the upper follows directly from the
form expression. Hence conforming dense form Galerkin trials converge in
the form norm by the energy projection property. This does not itself
make their complete Hilbert residuals converge; graph approximation is the
separate input for that conclusion.

\begin{proposition}[Cofinal transfer for ordered varying operators]\label{rtc:prop:varying}
Suppose $B_j=B+D_j$ on the same operator domain, with
$0\le D_j\le h_jI$, $h_j\to0$, and $R_j=(B_j+I)^{-1}$.
For any admitted prior $A\ge aI>0$ and $0\le\sigma_j<1$,
\begin{equation}\label{rtc:eq:inverse-transfer}
 0\le[I-(1-\sigma_j)R_j]-(I-R)
 \le\left({h_j\over a^2}+{\sigma_j\over a}\right)I.
\end{equation}
If $\sigma_j>0$ tends to zero, the two-index characterization in
Theorem~\ref{rtc:thm:cofinal} remains valid for lower matrices formed
using $B_j$ at stage $j$, with fixed nested dense source spaces and a
common graph-core response family. For the two-response scheme additionally require $\sigma_j\le1/2$ and
use its coarse penalty; sharper valid priors may also be used stagewise
within their stated parameter ranges.
\end{proposition}
\begin{proof}
Inverse order gives $0\le R-R_j$, and the resolvent identity
$R_j-R=-RD_jR_j$ gives $\norm{R-R_j}\le h_j/a^2$.
Thus the difference is $(R-R_j)+\sigma_jR_j$, proving the ordered bound.
This argument does not treat a product of noncommuting positive operators
as positive; positivity comes from inverse order.
If $B\ge0$ then $B_j\ge0$. At each fixed $j$, bounded perturbation
preserves the graph core, so the fixed-operator theorem supplies a finite
response dimension with positive margin.
Conversely, for fixed retained $y$, one-response certificates and
\eqref{rtc:eq:inverse-transfer} give $\ip{(I-R)y}y_{\Hh}\ge0$ in the
limit. For two responses, $R_jy\to Ry$ in graph norm because
$B(R_j-R)=-BRD_jR_j$ and $BR=I-R$ is bounded.
Also
$q_j[R_jy]=q[R_jy]+\ip{D_jR_jy}{R_jy}_{\Hh}$,
whose second term tends to zero; graph continuity of the first gives
$q[Ry]\ge0$. The same core or bounded-form extension as before completes
the proof. Null-preserving heads \eqref{rtc:eq:null-head} satisfy exactly
the same hypotheses.
\end{proof}
All these conclusions require genuinely cofinal certificates, not any
finite list of successful stages. A diagonal abstract operator with one
negative entry, infinitely many zero entries and unbounded positive entries
satisfies the fixed-shift and filtered-core inputs; placing the negative
entry beyond a prescribed finite source prefix defeats every inference
from that prefix to the global sign. Likewise $B=0$ on a known null sector
and $B=-d/2$ on its orthogonal complement shows that exact null repair
cannot determine the complement's sign. These are abstract limitations,
not assertions about a theta counterexample. The remaining theta task is
a fully enclosed signed payment on a cofinal source family, or a proved
uniform full-space inequality implying those payments.

\subsection{Ordered theta tail transfer on the same domain}
\label{rtc:tail-section}

The fixed-operator response tests can be coupled to explicit tail budgets
without changing the weighted space or discarding the omitted row components.
In this subsection let $\mathcal B=T_\Gamma-D_b$ denote the ambient
weighted representative in~\eqref{eq:Db-physical}, and let
$B=\mathcal B|_{\Hh_c}$, where $\Hh_c=\ker m$ and
$Cu=u-2m(u)\kappa$.  The full-domain bridge and
\eqref{eq:actual-unitary} give $B\ge-I/2$; they also give
$\mathcal B\kappa=-\kappa/2$, so $\Hh_c$ reduces $\mathcal B$.
Write $q=q_0|_{X_c}$.  Every norm and pairing below is weighted unless
its subscript is $2$.

Lemma~\ref{a:lem:explicitW} already proves the full-series estimates
$W(x)\le42e^{4|x|-\pi e^{2|x|}}$ and $\norm{W}_\infty<1/2$.
Lemma~\ref{a:lem:prime} already proves the complete prime-power bound
\begin{equation}\label{rtc:tail-prime-budget}
 \norm{\mathsf P-\mathsf P_N}\le r_N,
 \qquad r_N=96(N+1)^3(27/512)^{N+1},\quad N\ge2,
\end{equation}
where $\mathsf P$ and $\mathsf P_N$ are the weighted representatives
of $p$ and $p_N$.  In that proof the translation norms are taken first
on ordinary ambient $L^2$; reflection invariance of the paired shifts
then permits restriction to even functions.  Neither an individual
translation on even functions nor a bounded unweighted infinite prime
operator is being asserted here.

\begin{lemma}[Complete Gamma displacement tail]\label{rtc:tail-gamma}
For $T>0$, the form
\[
 t_{>T}[u]=\int_T^\infty\rho(s)
               \norm{u(\cdot+s)-u}_2^2\,ds
\]
has a bounded positive representative $\mathsf T_{>T}$ on $\Hh$, with
\begin{equation}\label{rtc:tail-gamma-budget}
 0\preceq\mathsf T_{>T}\preceq\gamma_T I,
 \qquad \gamma_T=\frac{4e^{-T/2}}{1-e^{-2T}}.
\end{equation}
\end{lemma}
\begin{proof}
For ordinary translations, $\norm{u(\cdot+s)-u}_2\le2\norm{u}_2$.
Also $\norm{u}_2^2\le\norm{W}_\infty\norm{u}_{\Hh}^2$ and
\[
 \int_T^\infty\rho(s)\,ds
 \le\frac{1}{1-e^{-2T}}\int_T^\infty e^{-s/2}\,ds
 =\frac{2e^{-T/2}}{1-e^{-2T}}.
\]
Consequently $0\le t_{>T}[u]\le\gamma_T\norm{u}_{\Hh}^2$,
using $\norm{W}_\infty<1/2$.  Polarization and bounded-form
representation prove the assertion.
\end{proof}
Here $T$ cuts off displacement, not physical position or Fourier
frequency.  The compensated small-displacement Gamma integral remains
intact; its divergent diagonal rate is never integrated separately.

Put $\mathsf P_{>N}=\mathsf P-\mathsf P_N$.  On the unchanged domain
$\operatorname{Dom}B\subset\Hh_c$, define
\begin{align}
 B_{N,T}^{\mathrm{hd}}
 &=B-C\mathsf T_{>T}C+C\mathsf P_{>N}C,
 \label{rtc:tail-head}\\
 B_{N,T}^{+}
 &=B_{N,T}^{\mathrm{hd}}+(\gamma_T+r_N)I,
 & h_{N,T}&=\gamma_T+2r_N.
 \label{rtc:tail-upper}
\end{align}
The two compressions are necessary: the raw retained head need not
preserve the centered subspace.

\begin{proposition}[Ordered transfer]\label{rtc:tail-order}
These self-adjoint operators have the same operator domain, form domain
$X_c$, and graph cores as $B$.  Their bounded difference satisfies
\begin{equation}\label{rtc:tail-sandwich}
 0\preceq B_{N,T}^{+}-B\preceq h_{N,T}I.
\end{equation}
If $\delta\ge h_{N,T}$ and $\sigma=\delta-h_{N,T}$, then
\begin{equation}\label{rtc:tail-shift-transfer}
 q[u]+\delta\norm{u}_{\Hh}^2
 \ge q_{N,T}^{+}[u]+\sigma\norm{u}_{\Hh}^2
 \qquad(u\in X_c).
\end{equation}
In particular $B_{N,T}^{+}\ge-I/2$.
\end{proposition}
\begin{proof}
The difference in question is
\[
 (\gamma_T I-C\mathsf T_{>T}C)
       +(r_N I+C\mathsf P_{>N}C).
\]
By~\eqref{rtc:tail-gamma-budget} its first summand lies between $0$
and $\gamma_T I$.  By~\eqref{rtc:tail-prime-budget} its second
summand lies between $0$ and $2r_NI$; no positivity of the prime tail
has been used.  This proves~\eqref{rtc:tail-sandwich}.
Adding a bounded self-adjoint operator preserves self-adjointness on the
same operator domain.  The graph norms are equivalent by the two
triangle inequalities, and the closed forms differ by the corresponding
bounded form on the same form domain.  Thus their graph cores also
coincide.  The upper side of~\eqref{rtc:tail-sandwich} gives
\eqref{rtc:tail-shift-transfer}, and the lower side preserves $B\ge-I/2$.
\end{proof}
This ordered head is not the positive jump energy with selected edges
removed.  It retains the full theta weight, the constant term $\cg$, and
the exact mean projection.  It is still an infinite-dimensional operator.

\begin{proposition}[An all-stage dyadic budget]\label{rtc:tail-schedule}
For every integer $j\ge1$, set
\begin{equation}\label{rtc:tail-dyadic}
 \delta_j=2^{-j},\qquad N_j=j+2,\qquad
 T_j=2(j+6)\log2.
\end{equation}
Write $r_j=r_{N_j}$, $\gamma_j=\gamma_{T_j}$,
$h_j=\gamma_j+2r_j$, $\sigma_j=\delta_j-h_j$, and
$B_j^+=B_{N_j,T_j}^+$.  Then
\begin{equation}\label{rtc:tail-all-stage}
 r_j<\delta_j/8,\qquad \gamma_j<\delta_j/8,\qquad
 h_j<3\delta_j/8,\qquad \sigma_j>5\delta_j/8>0,
\end{equation}
and $2h_j<1$.  In particular,
\begin{equation}\label{rtc:tail-stage-transfer}
 B+\delta_jI\succeq B_j^++\sigma_jI
 \succeq B_j^++\tfrac12\delta_jI+\tfrac18\delta_jI.
\end{equation}
\end{proposition}
\begin{proof}
For $j=1$ one has
\[
 r_3=\frac{1594323}{33554432}<\frac1{16}=\frac{\delta_1}{8}.
\]
For $N\ge3$,
\[
 \frac{r_{N+1}}{r_N}
 =\left(\frac{N+2}{N+1}\right)^3\frac{27}{512}
 \le\frac{3375}{32768}<\frac12.
\]
Induction proves the first bound at every stage.  Direct substitution gives
\[
 \gamma_j=\frac{\delta_j}{16(1-2^{-4(j+6)})}<\frac{\delta_j}{8},
\]
since $2^{-4(j+6)}<1/2$.  The definitions give the remaining bounds;
$\delta_j\le1/2$ implies $2h_j<3/8<1$.
Proposition~\ref{rtc:tail-order} proves~\eqref{rtc:tail-stage-transfer}.
\end{proof}
The schedule begins at $j=1$; its stated bounds are not claimed for
$j=0$.  They hold uniformly over all centered vectors, independently of
trial dimension, but give no sign for $B_j^++\sigma_jI$.

\subsection{Compact centered rows and complete Gram transfer}
\label{rtc:tail-rows-section}

\begin{lemma}[Compact centered graph core]\label{rtc:tail-core}
The space $C_{c,\mathrm{even}}^\infty(\R)\cap\ker m$ is a graph core
for $B$ and every $B_j^+$.
\end{lemma}
\begin{proof}
Theorem~\ref{d:thm:Gamma-graph-core}, applied to $D_b$, supplies the
ambient graph core for $\mathcal B$.  Choose a fixed even compact
smooth $\eta$ with $m(\eta)=1$; for example normalize a nonzero
nonnegative even bump.  For $u\in\operatorname{Dom}B$, choose
$g_n\in C_{c,\mathrm{even}}^\infty(\R)$ tending to $u$ in the
ambient graph norm.  Normalization gives
\[
 \ip{v}{\kappa}_{\Hh}=2m(v),\qquad
 |m(v)|\le\tfrac12\norm{v}_{\Hh},
\]
so $m(g_n)\to m(u)=0$.  Therefore
$g_n-m(g_n)\eta$ is compact, smooth, exactly centered and converges to
$u$ in that graph norm.  As $\Hh_c$ reduces $\mathcal B$, this is the
graph norm of $B$.  Proposition~\ref{rtc:tail-order} transfers the core
to $B_j^+$.  In particular, no cutoff of the noncompact vector $Cg_n$
and no uncharged cutoff error is needed.
\end{proof}

For an even compact smooth $u$ define the raw retained row
\begin{align}
 \mathcal Q_{N,T}u(x)=W(x)\Bigg\{&
 \int_0^T\rho(s)[2u(x)-u(x+s)-u(x-s)]\,ds
 -\cg u(x)\notag\\
 &-\sum_{2\le n\le N}a_n
      [u(x+\log n)+u(x-\log n)]\Bigg\}.
 \label{rtc:tail-raw-row}
\end{align}
The combined Gamma bracket is $O(s^2)$ near zero, locally uniformly in
$x$, and has compact support after the retained shifts.  It therefore
defines the actual weighted row on these inputs.  Indeed pairing with a
compact smooth test, changing variables in the two shifts and then
using the form identity identifies its Gamma and prime contributions;
density of those tests identifies the representative.  Bounded removal
of the tails agrees with~\eqref{rtc:tail-head}.  For centered $u$,
\begin{equation}\label{rtc:tail-centered-row}
 A_ju:=(B_j^++I)u
 =C\bigl((1+\gamma_j+r_j)u+\mathcal Q_{N_j,T_j}u\bigr).
\end{equation}
This row identity extends by graph convergence to its full operator
domain, with the raw row interpreted as its graph-continuous extension.

\begin{proposition}[Mean subtraction and the exact exterior row]
\label{rtc:tail-gram}
Let $\phi_1,\ldots,\phi_n$ be centered compact smooth trial vectors and
put $d_a=(1+\gamma_j+r_j)\phi_a+\mathcal Q_{N_j,T_j}\phi_a$.
Use the linear-first pairing convention and define
\begin{align}
 G_{ab}&=\ip{\phi_b}{\phi_a}_{\Hh},&
 H^{(j)}_{ab}&=\ip{A_j\phi_b}{\phi_a}_{\Hh},&
 L^{(j)}_{ab}&=\ip{A_j\phi_b}{A_j\phi_a}_{\Hh}.
 \label{rtc:tail-GHL}
\end{align}
Their exact raw-row formulas are
\begin{equation}\label{rtc:tail-mean-Gram}
 H^{(j)}_{ab}=\ip{d_b}{\phi_a}_{\Hh},\qquad
 L^{(j)}_{ab}=\ip{d_b}{d_a}_{\Hh}
                  -4m(d_b)\overline{m(d_a)}.
\end{equation}
If the trial supports lie in $[-a_0,a_0]$, then, for
$R_*=a_0+\max\{T_j,\log N_j\}$,
\begin{equation}\label{rtc:tail-exterior-row}
 \mathbf1_{\{|x|>R_*\}}A_j\phi_a
 =-2m(d_a)\mathbf1_{\{|x|>R_*\}}\kappa.
\end{equation}
The exterior contribution to $L^{(j)}$ is exactly
\begin{equation}\label{rtc:tail-exterior-Gram}
 4\mu(\{|x|>R_*\})
       \bigl(m(d_b)\overline{m(d_a)}\bigr)_{ab},
\end{equation}
a positive semidefinite matrix of rank at most one.
\end{proposition}
\begin{proof}
The centered trial vectors are orthogonal to $\kappa$, so
$\ip{Cd_b}{\phi_a}_{\Hh}=\ip{d_b}{\phi_a}_{\Hh}$.
As $C$ is the orthogonal projection removing the unit vector $\kappa$,
\[
 \ip{Cd_b}{Cd_a}_{\Hh}
 =\ip{d_b}{d_a}_{\Hh}
  -\ip{d_b}{\kappa}_{\Hh}\ip{\kappa}{d_a}_{\Hh},
\]
which gives~\eqref{rtc:tail-mean-Gram}.  Every displacement in
\eqref{rtc:tail-raw-row} has magnitude at most
$\max\{T_j,\log N_j\}$, so $d_a$ vanishes outside $[-R_*,R_*]$.
Equation~\eqref{rtc:tail-centered-row} then gives
\eqref{rtc:tail-exterior-row}.  Finally
$\kappa^2/W=2\cosh(x/2)\kappa$; integration of the mixed exterior
row pairings gives~\eqref{rtc:tail-exterior-Gram}.
\end{proof}
The negative rank-one term in~\eqref{rtc:tail-mean-Gram} already
includes the exterior positive mass together with the interior cross
terms.  The exterior matrix must not be added a second time.  Full theta
weights, finite-interval raw-row pairings and moments, and the
compensated Gamma integral still require rigorous evaluation.

If the $\phi_a$ are independent and $P_F$ is their weighted orthogonal
projection, the complete omitted-row cost is
\begin{equation}\label{rtc:tail-leakage}
 D^{(j)}=L^{(j)}-H^{(j)}G^{-1}H^{(j)}
 =\bigl(\ip{(I-P_F)A_j\phi_b}{(I-P_F)A_j\phi_a}_{\Hh}\bigr)_{ab}
 \succeq0.
\end{equation}
To check this, the coordinates of $P_FA_j\phi_b$ are the $b$th
column of $G^{-1}H^{(j)}$, so its Gram is
$H^{(j)}G^{-1}H^{(j)}$; weighted orthogonality gives the identity.
No tail cutoff permits replacing $L^{(j)}$ by that projected Gram.

\begin{proposition}[Complete row and graph-resolvent transfer]
\label{rtc:tail-relative}
Let $A=B+I$, $R=A^{-1}$, and $R_j=A_j^{-1}$.
For any finite admitted trial family in their common operator domain,
let $G,H,L$ be the trial Gram, the trial--row pairing matrix, and
the full row Gram, respectively, as in~\eqref{rtc:tail-GHL}.
In particular, $L$ requires no $\operatorname{Dom}A^2$ assumption.  Then
\begin{equation}\label{rtc:tail-relative-Grams}
 0\preceq H^{(j)}-H\preceq h_jG,\qquad
 (1-2h_j)^2L^{(j)}\preceq L\preceq(1+2h_j)^2L^{(j)}.
\end{equation}
Moreover
\begin{equation}\label{rtc:tail-resolvent-graph}
 \norm{R_j-R}\le4h_j,\qquad
 \norm{B(R_j-R)}\le2h_j.
\end{equation}
\end{proposition}
\begin{proof}
Put $\mathcal D_j=B_j^+-B$, so
$0\preceq\mathcal D_j\preceq h_jI$.  Testing this order on each
trial combination gives the first matrix inequality.  Since $A_j\ge I/2$,
\[
 \norm{(A_j-A)u}_{\Hh}\le h_j\norm{u}_{\Hh}
                     \le2h_j\norm{A_ju}_{\Hh}.
\]
The triangle and reverse triangle inequalities, together with $2h_j<1$,
give
\[
 (1-2h_j)\norm{A_ju}_{\Hh}
 \le\norm{Au}_{\Hh}
 \le(1+2h_j)\norm{A_ju}_{\Hh}.
\]
Squaring for every coefficient combination proves the second pair in
\eqref{rtc:tail-relative-Grams}.  This argument does not commute
$A$ and $\mathcal D_j$.

Both resolvents have norm at most $2$ and map into the common operator
domain.  Their bounded resolvent identity is
$R_j-R=-R\mathcal D_jR_j$.  Hence the first bound in
\eqref{rtc:tail-resolvent-graph} follows.  Spectral calculus for
$B\ge-I/2$ gives $BR=I-R$ and $\norm{BR}\le1$; applying this
bounded operator to the identity gives the second bound.  Thus the
resolvents converge from the source Hilbert norm to the actual graph
norm, not merely in Hilbert norm.
\end{proof}
These constants are dimension independent in the physical weighted
space, not condition-number bounds in arbitrary finite coordinates.
The order, row and resolvent estimates allow a separately proved
cofinal sign certificate to be transferred between the displayed
operators.  They do not compare optimized response matrices from
different stages without that transfer.  In particular, if an admitted
$u$ has $q[u]=-d\norm{u}_{\Hh}^2$ with $d>0$, then
\[
 q_j^+[u]+\sigma_j\norm{u}_{\Hh}^2
 \le q[u]+\delta_j\norm{u}_{\Hh}^2
 =(\delta_j-d)\norm{u}_{\Hh}^2<0
 \qquad(\delta_j<d).
\]
Thus even these vanishing, all-stage tail budgets cannot remove an
actual negative direction.  No such theta direction is asserted, and
no additional numerical sign certificate is supplied by this subsection.


\section{Exterior equality directions and exact response deflation}\label{ext:sec:exterior}
The derivative radicals also constrain the exact exterior Schur problem. The source identities, original domains and sharp-cut admission are established in Sections~\ref{a:sec:radical} and~\ref{c:sec:limits} and Appendix~\ref{c:app:cuts}. This section first works with real even vectors, so the finite projection formulas below use real bilinear forms and transposes. The corresponding complex inequalities follow by applying the real result to real and imaginary parts; no change of inner-product convention is implicit.

Fix $B>0$, $I=(-B,B)$ and $O=\R\setminus[-B,B]$. Here $\D_I$ and $\D_O$ denote the zero-extended compact-smooth minimum closures for this radius, in the actual $X$ topology, and $\D_O^0=\{u\in\D_O:m_O(u)=0\}$. These symbols are local to this section rather than fixed to the earlier radius $B_0$. Write $q=q_0+2|m|^2$. Assume that its mean-inclusive local representative $L$ satisfies $L\succeq\ell I>0$ on physical $L^2(I,dx)$. This is an explicit local positivity premise, not a consequence of the global radicals. Let $J$ be the complete physical exterior source, with every prime power and Gamma interface retained. For each fixed $B$, the Gamma crossing has the same ordinary-$L^2$ Carleman bound as in Appendix~\ref{c:app:cuts}. The complete prime form is bounded on $\Hh$, and $\min_I W>0$ bounds the weighted norm of a local test by its physical $L^2(I)$ norm. Together with $\|u\|_2\le\|W\|_\infty^{1/2}\|u\|_{\Hh}$, these facts give a bounded complete row $J:L^2(O,dx/W)\to L^2(I,dx)$. No numerical source constant specific to $B_0$ is transferred. This bound and the strict local floor admit $L^{-1}Ju$ for $u$ in the original centered exterior minimum domain $\D_O^0$.

For $v\in\D_I$ and $u\in\D_O^0$, the exact block identity is
\begin{equation}\label{ext:eq:S}
 q[v+u]=q_I[v-L^{-1}Ju]+S[u],\qquad
 S[u]=q_0[u]-\ip{Ju}{L^{-1}Ju}.
\end{equation}
The local mean pole has not been deleted. Put $\alpha=\sqrt W$, $\mathcal W(u)=\int_O|u|^2/W$, and use the complete nonnegative prime difference form $D_p$ and full row potential $V_p$ of Section~\ref{c:sec:signed}. Thus, with $V_0=c_\Gamma W+V_p$,
\begin{align}
 \mathcal P[u]&=t_\Gamma[u]+D_p(u/\alpha),\label{ext:eq:P}\\
 \mathcal K[u]&=\int_O V_0|u/\alpha|^2+\ip{Ju}{L^{-1}Ju},\label{ext:eq:K}\\
 S&=\mathcal P-\mathcal K.\label{ext:eq:PK}
\end{align}
These are the original-domain identities, not declarations of an unproved completion. The Gamma form is strictly positive on a nonzero exterior numerator extended by zero on $I$: zero energy makes its translations equal almost everywhere and hence gives a constant ordinary-$L^2$ function, necessarily zero. Consequently $\mathcal P$ is an inner product on $\D_O^0$. Proposition~\ref{c:prop:nogap} shows that it does not dominate a fixed positive multiple of $\mathcal W$.

For $d_j=\kappa^{(2j)}$, the full-domain identities already proved give $m(d_j)=4^{-j}/2$ and $q_0(d_j,z)=-4^{-j}m(z)$ for real $z\in X$. Therefore $q(d_j,z)=0$. These are bilinear radical identities, not an assumption that the whole form is nonnegative.

\subsection{Exact exterior equality directions}
Define, for $j\ge1$,
\begin{equation}\label{ext:eq:uj}
 M_j=m_O(d_j),\quad \beta_j=M_j/M_0,\quad
 w_j=d_j-\beta_j\kappa,\quad
 u_j=\one_Ow_j,\quad v_j=\one_Iw_j.
\end{equation}
Here $M_0=m_O(\kappa)=\mu(O)/2>0$.

\begin{theorem}\label{ext:thm:null}
Under the strict-local-positivity premise and the exact identities stated above, each $u_j$ belongs to $\D_O^0$ and
\begin{equation}\label{ext:eq:exactresponse}
 v_j\in\operatorname{Dom}(L),\qquad Lv_j=Ju_j,\qquad
 S(u_j,u)=0\quad(u\in\D_O^0).
\end{equation}
The exterior vectors $(u_j)_{j\ge1}$ are linearly independent. In particular
\begin{equation}\label{ext:eq:equality}
 \mathcal P(u_j,u)=\mathcal K(u_j,u),\qquad
 \mathcal P[u_j]=\mathcal K[u_j]>0.
\end{equation}
These conclusions do not require $S\ge0$.
\end{theorem}
\begin{proof}
The moment definition and sharp-cut input give $u_j\in\D_O^0$ and $v_j\in\D_I$. Equation \eqref{a:eq:radicalidentity} gives $q(w_j,h)=0$ for every $h\in X$, since both $d_j$ and $\kappa$ annihilate $q$ bilinearly. Testing with an arbitrary local form vector yields
\[
 q_I(v_j,h)=\ip{Ju_j}{h}.
\]
The right side is an $L^2(I)$ representing row. The representation theorem for the closed local form implies $v_j\in\operatorname{Dom}(L)$ and $Lv_j=Ju_j$. Testing instead with an arbitrary $u\in\D_O^0$ gives
$q_0(u_j,u)-\ip{v_j}{Ju}=0$, which is the polarized Schur identity in \eqref{ext:eq:exactresponse}.

If $\sum_{j=1}^N b_ju_j=0$, the analytic function
$\sum_{j=1}^N b_j(\kappa^{(2j)}-\beta_j\kappa)$ vanishes on the open interval $(B,\infty)$, hence on $\R$. Fourier transformation gives
\[
 \left[\sum_{j=1}^N b_j\bigl((-t^2)^j-\beta_j\bigr)\right]\widehat\kappa(t)=0.
\]
Since $\widehat\kappa(0)=\int\kappa>0$, the polynomial vanishes identically. Its successive highest-degree coefficients give $b_j=0$ for every $j$. Finally \eqref{ext:eq:PK} and strict positivity of $\mathcal P$ give \eqref{ext:eq:equality}.
\end{proof}

The centering coefficients admit boundary formulas. Twice integrating by parts on $(B,\infty)$ gives
\begin{equation}\label{ext:eq:recurrence}
 M_j=\frac14 M_{j-1}
 -2\bigl[c(B)\kappa^{(2j-1)}(B)-c'(B)\kappa^{(2j-2)}(B)\bigr].
\end{equation}
For example
\[
 \beta_1=\frac14-
 \frac{2\{c(B)\kappa'(B)-c'(B)\kappa(B)\}}{M_0}.
\]
All derivatives at a specified finite order and the single exterior mass $M_0$ must be enclosed if these formulas are used numerically. No uniform derivative-tail assertion is made as $j$ increases.

\subsection{An endpoint obstruction stronger than absence of a weighted gap}
\begin{corollary}\label{ext:cor:critical}
There is no $\varepsilon>0$ such that $\mathcal K\le(1-\varepsilon)\mathcal P$ on the whole exterior domain. The same statement holds after imposing any finite number of linear constraints: their common kernel still contains a nonzero finite combination of the $u_j$.

If $\mathcal K$ extends to a bounded form on the abstract Hilbert completion of $(\D_O^0,\mathcal P)$, its representing positive operator $T$ has eigenvalue $1$ with infinite multiplicity. It is not compact. Moreover
\begin{equation}\label{ext:eq:essential}
 \norm{T-F}\ge1\quad\hbox{for every finite-rank operator }F.
\end{equation}
\end{corollary}
\begin{proof}
Substitute a nonzero combination of $u_j$ into \eqref{ext:eq:equality}. More unknown coefficients than linear constraints provide such a combination satisfying those constraints. If the bounded representation exists, density and polarization give $Tu_j=u_j$. An infinite orthonormal sequence in their $\mathcal P$-span proves noncompactness. For finite-rank $F$, its restriction to this infinite-dimensional eigenspace has a nontrivial kernel. A unit vector $x$ in that kernel satisfies $(T-F)x=x$.
\end{proof}
The completion in this corollary is used only as an abstract form Hilbert space. No identification of all of its elements with physical numerator functions, or unproved boundedness of $\mathcal K$, is needed. If the desired $\mathcal K\le\mathcal P$ holds on the original domain, that bounded extension would follow and would still be noncompact.

This obstruction concerns approximating the \emph{whole relative cost} by a finite-rank operator with norm error below one. It does not obstruct a finite Schur reduction with a separately proved infinite complement or a representation retaining an infinite reference operator.
It also shows why a positive definite additive error is incompatible with a direct endpoint upper certificate: if $\widehat{\mathcal K}\succeq\mathcal K$ and $\widehat{\mathcal K}-\mathcal K$ is positive on some $u_j$, then $\widehat{\mathcal K}\npreceq\mathcal P$. In particular $\mathcal K+\varepsilon\mathcal W\npreceq\mathcal P$ for any $\varepsilon>0$.

\subsection{Removing known null directions without changing the defect}
Fix $n\ge1$ and set $N_n=\operatorname{span}\{u_1,\ldots,u_n\}$.
Let $G_{ij}=\mathcal P(u_i,u_j)$, $b_i(u)=\mathcal P(u_i,u)$, and define
\begin{equation}\label{ext:eq:projection}
 a(u)=G^{-1}b(u),\qquad
 \Pi_nu=\sum_{j=1}^n a_j(u)u_j,\qquad Q_n=I-\Pi_n.
\end{equation}
The matrix $G$ is positive definite. These maps preserve the original domain; only finite combinations of admitted vectors are added.

\begin{theorem}\label{ext:thm:deflate}
For every $u\in\D_O^0$,
\begin{align}
 S[u]&=S[Q_nu],\label{ext:eq:Sdeflate}\\
 \mathcal P[Q_nu]&=\mathcal P[u]-b(u)^{\mathsf T}G^{-1}b(u),\label{ext:eq:Pdeflate}\\
 \mathcal K[Q_nu]&=\mathcal K[u]-b(u)^{\mathsf T}G^{-1}b(u).\label{ext:eq:Kdeflate}
\end{align}
The required exact local response of the removed part is
\begin{equation}\label{ext:eq:responseN}
 L^{-1}J\Pi_nu=\sum_{j=1}^n a_j(u)v_j,
 \qquad
 q_I[L^{-1}J\Pi_nu]\le\mathcal P[\Pi_nu]\le\mathcal P[u].
\end{equation}
Thus this prescribed response map has norm at most one from the $\mathcal P$ energy to the local $q_I$ energy, independently of $n$.
\end{theorem}
\begin{proof}
$\Pi_n$ is the $\mathcal P$-orthogonal projection. Theorem \ref{ext:thm:null} gives $S(N_n,\D_O^0)=0$ and $\mathcal K(N_n,\D_O^0)=\mathcal P(N_n,\D_O^0)$. These yield the three identities by expansion. The response identity follows from $Lv_j=Ju_j$. For $a=\Pi_nu$, equality of $\mathcal K[a]$ and $\mathcal P[a]$ gives
\[
 q_I[L^{-1}Ja]=\mathcal P[a]-\int_O V_0|a/\alpha|^2\le\mathcal P[a].
\]
\end{proof}
There is no claim of a uniformly bounded map between the \emph{weighted} spaces used in derivative interpolation. The norms in \eqref{ext:eq:responseN} are different, and the absent inequality $\mathcal P\ge\varepsilon\mathcal W$ cannot be inserted. A poorly scaled derivative Gram matrix can still be numerically ill-conditioned.

The projection data also have useful inverse-free expressions:
\begin{align}
 G_{ij}&=\int_O V_0\frac{\overline{u_i}u_j}{W}+q_I(v_i,v_j),\label{ext:eq:Gsource}\\
 b_i(u)&=\int_O V_0\frac{\overline{u_i}u}{W}+\ip{v_i}{Ju}.
 \label{ext:eq:bsource}
\end{align}
Thus no unknown inverse response for the prescribed null columns is needed. These expressions retain the full source, potential, and prime tail; they are not finite numerical certifications by themselves.

\subsection{An exact correction of approximate responses}
Let $T_0:\D_O^0\to\D_I$ be any linear trial-response map. Define
\begin{equation}\label{ext:eq:Tcorr}
 T_nu=T_0Q_nu+\sum_{j=1}^n a_j(u)v_j.
\end{equation}
It interpolates the true local response on $N_n$. Use the form-dual residual
\[
 r_n(u)(h)=\ip{Ju}{h}-q_I(T_nu,h),\qquad h\in\D_I,
\]
and its full energy-dual squared norm
\begin{equation}\label{ext:eq:resnorm}
 \mathcal R_n[u]=\sup_{h\ne0}\frac{|r_n(u)(h)|^2}{q_I[h]}.
\end{equation}
The supremum is over the whole local form domain, not just the trial modes. It includes every residual row and tail.

\begin{theorem}\label{ext:thm:residual}
The corrected response satisfies
\begin{align}
 r_n(u)&=r_0(Q_nu),\label{ext:eq:rcorrect}\\
 \mathcal R_n[u]&=q_I[L^{-1}Ju-T_nu].\label{ext:eq:resenergy}
\end{align}
If
\begin{equation}\label{ext:eq:Strial}
 \widetilde S_n[u]=q_0[u]-2\operatorname{Re}\ip{Ju}{T_nu}+q_I[T_nu],
\end{equation}
then the exact signed identity is
\begin{equation}\label{ext:eq:exacterror}
 \boxed{S[u]=\widetilde S_n[u]-\mathcal R_n[u].}
\end{equation}
Both $\widetilde S_n$ and $\mathcal R_n$ annihilate $N_n$ bilinearly. No global weighted residual bound or existence of an ordinary operator residual row is assumed.
\end{theorem}
\begin{proof}
The contribution of $\Pi_nu$ to the residual vanishes by \eqref{ext:eq:responseN}. The Riesz representative of $r_n(u)$ in the local $q_I$ Hilbert space is $L^{-1}Ju-T_nu$, which proves \eqref{ext:eq:resenergy}. Completing the square in \eqref{ext:eq:Strial} gives \eqref{ext:eq:exacterror}. The residual is zero on $N_n$, while $S$ annihilates it by Theorem \ref{ext:thm:null}; hence both polarized forms vanish there.
\end{proof}

For example, certified forms $\underline S_n\le\widetilde S_n$ and $\overline R_n\ge\mathcal R_n$ give the sufficient test
\begin{equation}\label{ext:eq:reggate}
 \underline S_n[w]+\eta\mathcal P[w]-\overline R_n[w]\ge0
 \quad\hbox{for every }w\in Q_n\D_O^0.
\end{equation}
Together with
\begin{equation}\label{ext:eq:regblock}
 (S+\eta\mathcal P)[u]=\eta\mathcal P[\Pi_nu]+(S+\eta\mathcal P)[Q_nu],
\end{equation}
this certifies $S+\eta\mathcal P\ge0$ on the whole original exterior domain. A finite calculation for \eqref{ext:eq:reggate} still needs a proved infinite complement and all enclosure errors. Finite trial positivity alone is not a proof of \eqref{ext:eq:reggate}.

At $\eta=0$, the lower certificate must preserve every equality direction it encounters. With $\eta>0$, each exact null direction has margin exactly $\eta\mathcal P[u_j]$. Any overestimate $E$ of the cost must therefore satisfy $E[u_j]\le\eta\mathcal P[u_j]$ on these directions. This bounds the allowable \emph{overcharge}; it is not a claim that the entire response itself is small.

Finite deflation cannot remove the infinitely many equality directions: there are always nonzero combinations of the $u_j$ in $Q_n\D_O^0$. Its purpose is exact cancellation in selected computed columns, not an automatic strictly positive gap on the remainder.

\subsection{What the next certificate must establish}
If one has the separate global estimate $S\ge\mathcal P-\varepsilon_8\mathcal W$, then a certificate $S+\eta\mathcal P\ge0$ yields
\begin{equation}\label{ext:eq:epsilon}
 S\ge-\frac{\eta}{1+\eta}\varepsilon_8\mathcal W.
\end{equation}
Indeed add $S+\eta\mathcal P\ge0$ to $\eta(S-\mathcal P+\varepsilon_8\mathcal W)\ge0$. For any cofinal sequence $\eta_j\downarrow0$ with certificates on the same original domain, the pointwise limit gives $S\ge0$, and \eqref{ext:eq:S} gives the global half bound. A single positive-$\eta$ certificate cannot be iterated without proving a new one.

Any implementation of this criterion must form admitted centered columns \eqref{ext:eq:uj}, enclose their $\mathcal P$ Gram matrix through \eqref{ext:eq:Gsource}, correct trial local responses by \eqref{ext:eq:Tcorr}, and certify \eqref{ext:eq:reggate} with a full complement and complete tails. The mean condition, the complete prime sum, and the Gamma boundary form are retained at every stage. Increasing $n$ also requires new high-derivative and Gram-conditioning control.

The positive-parameter weighted-pencil route of Section~\ref{sec:schur-unified} remains compatible with these conclusions. Adding a positive weighted mass changes the equality directions into positive directions. Its finite Schur form and the exterior $S$ here are different objects; constants and residual norms cannot be transferred without the corresponding form identities.

\paragraph{A finite block witness for the remaining limitation.}
The null identities alone cannot supply the remaining sign. For an explicit model take
\[
 L=\begin{pmatrix}3&1\\1&2\end{pmatrix},\qquad
 J=\frac1{10}\begin{pmatrix}1&2&-1&1\\2&-1&1&3\end{pmatrix},
\]
\[
 K=\operatorname{diag}(1,2,3,4),\quad
 P=\operatorname{diag}(1,2,4,7/2),\quad
 V=K-J^*L^{-1}J.
\]
Here $L\succ I$ and $\norm J^2\le\norm J_{\rm HS}^2=11/50$, so $V\succeq39I/50>0$.
The block form with exterior block $P-V$ has Schur complement
\[
 S=P-V-J^*L^{-1}J=P-K=\operatorname{diag}(0,0,1,-1/2).
\]
Its first two exterior coordinate vectors are exact null directions, with local responses $L^{-1}Je_1$ and $L^{-1}Je_2$. Every projection and response-correction identity above holds, but the fourth coordinate is still negative. This is an abstract model, not a counterexample to the theta half bound.


The exact finite example checks the projection and response algebra only. It does not compute an actual theta Schur matrix or certify its remaining sign. The local coercivity premise, complete source and row identities, original form domain and full residual norm remain essential. Finite deflation supplies no new global coefficient and no proof of the half inequality.


\section{Computational evidence and public reproducibility}\label{app:provenance-unified}
The analytic arguments are proved in this single source file. The following evidence identifies computational premises without treating a hash as a proof, and without citing an earlier manuscript in place of an analytic argument.


\subsection{Computational input identity and verification boundary}\label{c:app:provenance}
The proof has two different computational input classes.
Hypothesis~\ref{c:hyp:local} is an accepted whole-domain local theorem at ${B_0}=4/5$.
Its supplied packet gives $q_I[u]\ge10^{-22}\norm u_2^2$.
The records below identify that input and its computational evidence.
They are data-provenance identifiers, not citations to an earlier manuscript as a substitute for an analytic proof.
A cryptographic identity does not establish an inequality.

\begingroup\small
\begin{longtable}{@{}p{0.96\textwidth}@{}}
\toprule Local input artifact and SHA-256\\\midrule\endhead
Local theorem source\\*
\nolinkurl{ac04a393e8b6112d81bfb598c2dd00450b2a22eca36c298e3bbe82f7ef78b4c6}\\[3pt]
Local certificate archive version 1.1\\*
\nolinkurl{b805a431c4c77d955d5a6d4ed7397a01faf12cd727c1051af607ed14f77ab74b}\\[3pt]
Radius-four-fifths source diagnostic\\*
\nolinkurl{4c67f0c3f37bb8ce3c4f7c6da166050f2a1760c62bd47227c7709fb3a0b2ea05}\\[3pt]
Radius-four-fifths producer success record\\*
\nolinkurl{19156f6635bb25e853b9a1aecd650672043113f85d279c562e1a090269b0ca91}\\[3pt]
Exact-rational replay source\\*
\nolinkurl{3c5e006df82fbc985ce9b7ed73e26d62a9679086cc399dcb022224325c9fde13}\\[3pt]
Exact-rational replay result\\*
\nolinkurl{f2d4eef6e3efcd632ee252b2cca56be120e5ce5b3c59cdcc499277bc4ee528e2}\\[3pt]
Separate MPFR input archive\\*
\nolinkurl{4e57e261db7c8cfebdc5dc82b23d8aecfab3b639aa7fa5c95901b90f3aa18f8b}\\[3pt]
Separate MPFR evidence\\*
\nolinkurl{2a6fe3e8dfa04308ed555271bb19441b7d3ce783b6cc427419840a2d39a231d8}\\
\bottomrule
\end{longtable}\endgroup
The full local minimum domain, ordinary physical $L^2$ norm, positive mean pole, exact local prime powers $2,3,4$, and complete Gamma remainder are part of the input statement.
Finite matrix positivity alone is not the input theorem.
The supplied exact-rational replay and separately supplied source-enclosure evidence represent different validation layers.
This historical input record is not a new source-enclosure verification. Exact scalar and selected stored-matrix checks performed during consolidation are recorded separately below; the local certificate was not rerun for this document build.

The second class consists of Hypothesis~\ref{c:hyp:response}(a),(b).
One fixed data source for those contracts is the archive
\texttt{centered\_exterior\_cost\_sources.zip}, SHA-256
\begin{center}\small
\nolinkurl{4984450bc9f3174d354fdfe9731919dfb4c134f2db1ef0b55c40bc701b04792e}.
\end{center}
The relevant member identifiers are
\begingroup\small
\begin{longtable}{@{}p{0.96\textwidth}@{}}
\toprule Response input member and SHA-256\\\midrule\endhead
\texttt{centering\_vector.json}\\*
\nolinkurl{c04a7d7ac467971ae101290adbaec719826af343d7c0f589bdb7789c8ca8b4ed}\\[3pt]
\texttt{run/evidence.json}\\*
\nolinkurl{6c404a912f5986ae08080f2811a5d797d9aed22fc5683e24c9f2241525e4d3cd}\\[3pt]
\texttt{run/cells.json}\\*
\nolinkurl{d93657d28a399250f87f0d12f5421ced5c72178ff874c67ec5a3087a387ea6c7}\\[3pt]
\texttt{local\_response/run1024/response\_coefficients.json}\\*
\nolinkurl{b40fff3249a60dbb1224c9099360c96333d19fe22b7bd10bda126e9ece44f7a9}\\[3pt]
\texttt{local\_response/run1024/exact\_matrices.json}\\*
\nolinkurl{792718e3d299d74ca5757e5f17942608ab005468759673bd2f36388cf102dfba}\\[3pt]
\texttt{local\_response/run1024/evidence.json}\\*
\nolinkurl{62232dc8890083bfd967ad7fa3e854fe4ca00925b7b493a45dc7a45ecf234cbb}\\
\bottomrule
\end{longtable}\endgroup
The archive also identifies 60 finite local Gram blocks covering rows 1024 through 16383 and the infinite-tail enclosure.
Its local parameters are 1024 trial modes, cutoff 16384, 192-bit precision, 64 Gamma labels and tail order six.
The exterior parameters are fixed in Appendix~\ref{c:app:rows}.
These identify the declared inputs, not a new certification.
File identity, finite rational replay, and full source-enclosure validation remain distinct.
The single manuscript source requires no certificate data or executable research code to compile. Any separately distributed evidence retains its own manifest and trust scope.


\subsection{Public supporting archive and verification boundary}
The immutable public supporting archive is
\begin{center}\small
\href{https://github.com/serabing-hash/theta-weighted-certificate/tree/f539e52fe08379c0639ff94a812ec73a40d29a0e}{\texttt{serabing-hash/theta-weighted-certificate}}\\
commit \nolinkurl{f539e52fe08379c0639ff94a812ec73a40d29a0e}.
\end{center}
The accepted radius-$4/5$ local certificate package is available in that commit at
\begin{center}\small
\href{https://github.com/serabing-hash/theta-weighted-certificate/tree/f539e52fe08379c0639ff94a812ec73a40d29a0e/research/2026-10-28-four-fifths-local-certificate}{\texttt{research/2026-10-28-four-fifths-local-certificate}}.
\end{center}
Its complete local-domain theorem, source-enclosure layer and exact-rational replay layer must be distinguished. The original local package and later response stages have distinct verification histories. The manuscript build itself performs no prime or zero enumeration, source enclosure or response solve.
The exact source archive with SHA-256 beginning \texttt{b805a431} in the table above is directly linked as
\href{https://github.com/serabing-hash/theta-weighted-certificate/blob/f539e52fe08379c0639ff94a812ec73a40d29a0e/research/2026-10-28-four-fifths-local-certificate/prime_gamma_local_certificate_sources.zip}{\texttt{prime\_gamma\_local\_certificate\_sources.zip}} in the same immutable commit.

The historical response archive identified above is not asserted to be publicly downloadable from that older commit. The response theorem is an implication from the explicitly displayed complete-operator contracts, not a claim that their enclosure premises were newly certified or made public. This distinction does not affect the analytic implication, but it limits independent computational reproduction of that branch until the corresponding source archive is published.

This manuscript builds from one \texttt{.tex} file with standard packages. Every proof appendix and the bibliography are inline; no figures, bibliography database, external input file or certificate data are required to compile it. Document compilation, PDF rendering and file hashes verify artifact production only, not the mathematical correctness of computational input statements.

\section{Current dependency and verification map}\label{app:current-evidence}
\subsection{Analytical and computational boundaries}
The mathematical arguments are integrated in this source. The foundational and spectral parts retain their original analytic hypotheses and classical inputs; their inclusion is not a claim that every line has undergone a new independent audit. The response-gate derivation has scoped analytical review of normalization, weak source, minimum--maximal domains, cut and trial admission, complement estimates and matrix-to-energy composition. The sparse and contiguous-band arguments, the fixed-low-space sparse response theorem, the four- and five-zero computations, the four-mode conditional attachment, the actual Gamma-sector bridge, and the fixed-lattice argument have separate restricted-scope reviews. These reviews are not formal verification and do not certify the unresolved global sign.

The numerical gates are conditional on full residual and source enclosures for the actual operator. Exact rational or integer acceptance proves a finite statement about the supplied matrices; it does not by itself prove that those matrices enclose the complete infinite operator. File hashes establish identity only.

\subsection{Third and fourth response stages}
For the third stage $a=1/1024$, the completed historical separate-environment
replay used the same implementation and libraries. It regenerated the
full residual/sign pipeline and supplied fresh finite reconstructions
and exact integer witnesses. This differs from the earlier selected-file
checks and alternate-margin calculation. The later exact-rational
counterexample in Section~\ref{vrg:sec:review} identifies shared
integer-power implementation exposure. The earlier replay therefore
remains a reproducibility result, not a repaired source-enclosure
validation. No claim is made that an actual operand of that historical
run was shown to fail, or that its conditional analytic implication was
refuted. The third-stage enclosure obligation has not been discharged by
the corrected fourth-stage run.

For the fourth stage $a=1/2048$, the original producer replay, independent
analytical/static review, exact scalar composition, and independent
stored-matrix reconstruction remain separate historical evidence. A
subsequent complete recomputation now uses the directed-integer-power
repair from the unchanged authenticated physical candidate. It completed
source geometry, $T$, $G$, the complete Gamma calculation, the residual
model, all $5120$ compact residual cells, and both sufficient gates.
The exact-penalty gate has all $256$ corrected shifted LDL lower pivots
positive. Separate Python integer/Fraction checks reconstruct all
$65,536$ entries of each gate and verify all $256$ strict congruence
margins, for the exact penalty and for the conservative penalty $26240$.
This is a full corrected generation using the same numerical algorithms
and libraries with the stated arithmetic repair and storage adapters,
not a second complete numerical implementation or a formal proof assistant audit.

The corrected run does not reproduce every old intermediate endpoint:
cell and error-ledger enclosures and shifted-pivot intervals genuinely
widened. There are $142$ outward scientific endpoint changes: two lower
weight bounds in $G$, five upper bounds in the model error ledger, and
$135$ endpoints across $68$ shifted LDL pivot intervals. All $256$ corrected
lower pivots are positive. In particular, changing source-code and provenance hashes is
not the only difference. The final sufficient matrices and their exact
sign acceptance were retained. These are distinct facts: fresh execution,
corrected directed bounds, mathematical-field equality, and exact finite
acceptance must not be replaced by a single hash-match claim.

The retained-margin consequence $\lambda_{4,\mathrm{ret}}$ in
Corollary~\ref{vrg:cor:retained4} was also checked afresh against corrected
$M,G,K,F$. The historical integer witness supplied rows only; its old
acceptance result was not reused. Complete exact products verify all
$256$ strict retained-congruence margins and all $256$ symmetric row
bounds for $\mathsf U\prec4060I$. The same fixed
$\beta=10508819/40000000000$, continuation parameter and exact coefficient
therefore remain accepted. This is a consequence of the same fourth-stage
source computation, not a new source stage or a positive unshifted
low-space margin.

These numerical checks retain the full analytic source, residual, domain,
complement and energy-bridge premises, as well as the explicitly stated
implementation trust. They do not revalidate earlier historical stages
or prove the entire monograph formally. The separate local radius-$4/5$
input keeps its own exact-rational and Arb/MPFR verification history.
Every response coefficient remains below one half. The fixed-complement
continuation ceiling limits that method, not the true operator gap.
The global half inequality and the Riemann hypothesis remain unproved.

\subsection{The four zero low-space computation}\label{app:four-zero-evidence}
Theorem~\ref{fz:thm:finite} uses a separate Python-FLINT 0.9.0 and
FLINT 3.6.0 Arb/Acb computation. Its full source was reviewed and rerun:
the complete finite-integral checkpoint, rational matrix, rational sign
record and sparse-scalar record reproduced; elapsed-time metadata is not
part of the mathematical comparison. A higher-precision run gave the
same exported rational sufficient matrix. A separate integer Bareiss
calculation confirms all four positive leading principal minors of its
uniformly lowered midpoint matrix and the stated rational pivot floors.
These finite signs are distinct from the source-enclosure proof, which
uses the displayed theta, prime and Gamma bounds and trusted ball
integration. This is a same-code source replay with independent exact
finite acceptance, not a second complete source implementation or formal
verification.

This computation uses no source or stored matrix from the Decimal-based
response stages. Its actual weighted low Gram supplies the coupling
constant for this specific global theta--exponential-polynomial space.
The explicitly checked sparse constant recipe and height remain under
Hypothesis~\ref{fz:hyp:sparse}; the numerical checks do not themselves
prove that analytical hypothesis. The selected tail has total paired
budget $8\log T$, not an asserted covering of all actual unit-window jet
spaces. Neither its low margin nor its success repairs a separate
response-stage arithmetic premise or yields a global half bound.


\subsection{The fifth zero and the actual Gamma-sector bridge}
\label{app:five-gamma-evidence}
The fifth-zero certificate in Section~\ref{sec:five-zero-certificate}
retains the derivative Gram through the full Gamma interpolation debit.
Its complete source replay regenerated the four-function input and the
new five-function data. The final rational matrix and acceptance record
reproduced exactly, and the higher-precision and live-ball serialization
checks stated there passed. Separate rational determinant comparisons
verify the five positive leading minors, the shifted coordinate Schur
margin, and its quotient-normalized reserve. These exact signs concern
the exported sufficient matrices; the source enclosure still relies on
the explicit full-tail proofs and trusted Arb/Acb integration. The new
weighted Gram gives $K_{V_5}^2<75$, not the earlier four-mode bound $32$.
No sparse attachment with inherited four-mode constants, all-dimensional
sign, or improvement of a global comparison is asserted.

Section~\ref{gb:sec:bridge} is an analytic addition, under the existing
actual reducing-sector and full-zero-jet graph-core results. Its
restricted-form representative, projected graph cores, high bounds,
exact Schur equivalence, complete residual debit, and all-stage criterion
were checked at source level. The accompanying exact rational examples
check structural identities only. No actual theta Schur sign is supplied,
and the low Gamma spectral projection is neither assumed finite rank nor
identified with zero-height windows or finite response-source coordinates.
Neither this bridge nor the independent fifth-zero ball computation
repairs an arithmetic premise of the Decimal-based response stages.


The flexible-anchor improvement in Section~\ref{bl:sec:anchor} uses the
stronger published exponential-sum factor, with fixed parameters before
taking the budget limit. Its localization proof, positive-envelope count,
internal-cut residual inequality, complete growing-band cross bound and
three infinite row estimates were reviewed. The separate scalar checks
verify explicit rational margins and sufficient constants; they are not
actual theta-matrix computations or evaluated final PNT heights. The
ordinary-$L^2$ margins are not uniform reciprocal-theta margins, and the
new growing-band geometry does not inherit the separate unit-window
low-space attachment without an additional mixed estimate.


\subsection{Five-mode graph transfer and complete Gamma high attachment}
Section~\ref{gs:sec:transfer} retains the explicit sector admission and
operator-row hypotheses. Its new finite comparison bounds the complete
second-derivative Gram by $2^{20}$ times the actual weighted Gram, using
frozen data from the independently replayed five-mode computation.
The complete prime and Gamma estimates give a graph bound of $52$ and
then an exact inverse-response debit. At Gamma spectral cutoff $2^{18}$,
the projected five-dimensional source has Schur margin $1/35$; attaching
the entire complementary Gamma-high form domain leaves margin $1/36$.
The audited small producer and separate rational checker reproduce the
stated derivative-matrix and scalar comparisons. They do not recompute
theta integrals, construct the spectral projection or high inverse, or
establish the admitted-sector hypothesis from graph density. The other
low spectral directions, cofinal sign and global half bound remain open.


\subsection{Incremental sixth direction at the same Gamma split}
Section~\ref{sx:sec:six} evaluates the new sixth row and column while
retaining frozen enclosing five-function integral blocks, then recomputes
the current Gamma interpolation and complete-tail debits. A fresh
incremental source replay, started without a sixth-row checkpoint,
regenerated the new Gram, prime, Gamma and derivative entries. The four
interval matrix signs and rational transfer inequalities were checked
separately. This is not a fresh regeneration of every inherited
five-function integral, and hashes alone do not prove those inherited
enclosures. The original raw-sector and admitted-row inputs remain
explicit. The fixed-split exhaustion theorem is a conditional criterion
for every required nested stage and projected density; six functions do
not supply that cofinal family or certify the other low directions.


\subsection{Uniform response accuracy and its sign boundary}
Section~\ref{ur:sec:response} preserves the complete inverse-energy
structure at one fixed Gamma split. Its paired polynomial brackets
converge monotonically in operator norm with a geometric error independent
of finite source dimension. The full prime-tail comparison uses that same
low-space operator norm. These are analytic approximation statements,
not a new actual theta low-matrix or positivity computation. The original
rational and elementary ball-arithmetic checkers were replayed, and a
separately written exact check tested the constants, paired bracket
identities, monotonicity and noncommuting synthetic examples. The source
proofs retain the restricted-domain and complete-correction hypotheses;
the projected-core step assumes the whole admitted exact-zero-jet form
core. Its fixed graph-debit obstruction does not invalidate earlier
successful five- or six-mode certificates. It explains why those same
fixed debits cannot certify an exhaustive raw family. Evaluated projected
moments, actual finite low signs and all required cofinal stages remain
missing.

\subsection{Restricted positivity and infinite extensions}
The sparse theorem permits an unbounded number of windows only under its declared spacing, local prefix budgets and actual-form convergence. The contiguous theorem includes every actual zero jet and internal coupling in one band. Its flexible-anchor refinement reaches an eleven-logarithmic band with an ordinary-$L^2$ margin, and the complete exponent-$12$ row estimate permits its declared sparse growing-band families. Neither permits an arbitrary adjacent covering. Its prime-number-theorem refinement is qualitative and has no effective final height here. Fixed-lattice positivity permits infinitely many real lattice modes only for profiles already admitted to the actual form domain, with one fixed positive spacing. The signed model demonstrates a cancellation mechanism; it is not an identification of actual theta interactions with that model. The low-to-sparse response estimate has a fixed finite global theta--exponential-polynomial low-space hypothesis and uses its true weighted Gram. Its closure statement is variational in the actual form norm, not a bounded inverse on the reciprocal-theta Hilbert space. Cut-centered cosine sources, piecewise-linear trials and response images from a finite numerical gate are not automatically that low space; neither their matrices nor a sub-half global comparison supplies its required positive low margin. The four-zero space in Section~\ref{sec:four-zero-certificate} is a different, explicitly admissible low space with its own certified positive margin and weighted Gram. It does not identify the numerical response spaces with that low space or include omitted windows.

The public repository for this manuscript and its reproducibility materials is \url{https://github.com/serabing-hash/theta-weighted-certificate}. The integrated manuscript, source and selected verification files are available in the \href{https://github.com/serabing-hash/theta-weighted-certificate/tree/main/papers/2026-11-03-unified-theta-energy-v2}{integrated release folder}.

The public companion discovery path is the \href{https://github.com/serabing-hash/theta-weighted-certificate/tree/main/papers/2026-11-03-unified-theta-energy-v2/reproducibility}{reproducibility folder of the integrated release}. This is a mutable repository path, not an immutable certificate identifier. The release manifest and file SHA-256 values identify the exact artifacts; the published release record supplies the immutable commit permalink separately.

The assembled source is standalone. All proofs and external bibliography entries are inline, and the compile command has no scientific data dependency. Build success, extracted-text checks, page rendering and SHA-256 manifests validate artifact production only. Current public file availability is established by the accompanying release manifest, not by private storage links or historical transfer records.

\section{Short extrapolation and growing zero jet bands}\label{app:growing-zero-bands}
This appendix incorporates the complete analytic argument of the separately
published growing zero jet bands manuscript~\cite{ChoiGrowingBands}. It strengthens the earlier
single-band result of Section~\ref{sec:sparse-bands}; its ordinary-$L^2$
comparison should not be read as a uniform reciprocal-theta weighted gap.
Here the forms are defined on the displayed functions without imposing evenness
until the actual even zero jets are selected. All internal band couplings,
full multiplicity prefixes, the complete Gamma term, and every prime power are
retained. The fixed-boundary result is a lower bound under its stated strict
constant, not the two-sided asymptotic of the smaller-width result.

\subsection{The form and the main conclusions}
All logarithms are natural. Norms with subscript $2$ are ordinary,
unnormalized Lebesgue $L^2$ norms, on the interval indicated or on $\R$.
Inner products are linear in the first argument.
For $r=|x|$ define the positive even weight
\begin{equation}\label{gzb:eq:weight}
 W(x)=\frac{1}{1+e^{-r}}
 \sum_{j=1}^{\infty}(4\tau_j^2-6\tau_j)e^{-\tau_j},
 \qquad \tau_j=\pi j^2e^{2r}.
\end{equation}
Let $\Lambda$ be the von Mangoldt function and put
\begin{equation}\label{gzb:eq:forms}
\begin{split}
 \rho(s)&=\frac{e^{-s/2}}{1-e^{-2s}},\qquad
 a_n=\frac{\Lambda(n)}{\sqrt n},\qquad
 \cg=\log\pi-\psi(1/4),\quad \psi=\Gamma'/\Gamma,\\
 \tg[u]&=\int_0^\infty\rho(s)
               \norm{u(\cdot+s)-u}_2^2\,ds,\\
 p[u]&=2\Rea\sum_{n=2}^{\infty}a_n
                  \int_\R u(x+\log n)\overline{u(x)}\,dx,\\
 q_0[u]&=\tg[u]-\cg\norm u_2^2-p[u].
\end{split}
\end{equation}
The sum runs over all integers, so it includes every prime power with its
correct von Mangoldt coefficient and both translation orientations.
We use these full forms, not finite heads. They are well defined on every
$u=WP$ considered below. Such functions belong to
$L^2(\R,dx/W)\cap H^1(\R)$; the estimates below prove absolute convergence
of the prime sum and finiteness of the Gamma integral.

An exponential polynomial has prefix budget at most $N$ if it has a
representation
\begin{equation}\label{gzb:eq:prefix}
 P(x)=\sum_{j=1}^{J}p_j(x)e^{z_jx},\qquad
 \deg p_j<m_j,\qquad \sum_jm_j\le N.
\end{equation}
Equal $z_j$ may be merged. A selected degree $k$ costs its entire prefix
of length $k+1$ even if some coefficients vanish. No coefficient norm,
separation condition or inverse Vandermonde estimate is imposed.

To state the actual zero consequence, write
\[
 \xi(s)=\tfrac12s(s-1)\pi^{-s/2}\Gamma(s/2)\zeta(s),\qquad
 \Xi(w)=\xi(1/2+iw).
\]
For each zero $w$ of $\Xi$ with $\Rea w>0$, let $\nu(w)$ be its order and set
\begin{equation}\label{gzb:eq:jets}
 J_{w,k}(x)=W(x)\left.\partial_z^k\cos(zx)\right|_{z=w},
 \qquad 0\le k<\nu(w).
\end{equation}
Let $V_{T,L}$ be the complex span of all these vectors with
$T\le\Rea w\le T+L$, including every prefix in \eqref{gzb:eq:jets}.
Zeros off the critical line are included: the classical critical strip
only gives $|\Ima w|<1/2$ \cite[Sections 25.4 and 25.10]{DLMF}.

\begin{theorem}[Entire growing actual bands]\label{gzb:thm:main}
Let $L(T)\ge0$ satisfy
\begin{equation}\label{gzb:eq:width}
 L(T)=o\bigl((\log T)\sqrt{\log\log T}\bigr)
 \qquad (T\longrightarrow\infty).
\end{equation}
There is a scalar $\varepsilon_T\to0$, independent of the coefficient
vector in $V_{T,L(T)}$, such that
\begin{equation}\label{gzb:eq:main}
 (1-\varepsilon_T)\log T\,\norm u_2^2
 \le q_0[u]\le
 (1+\varepsilon_T)\log T\,\norm u_2^2,\qquad u\in V_{T,L(T)}.
\end{equation}
In particular, this holds for $L(T)=(\log T)(\log\log T)^{1/4}$.
The norm is that of the entire represented sum, with all correlations
inside the band. Empty spaces cause no difficulty.
\end{theorem}

The analytic estimates hold for arbitrary exponents in the strip.
Zeta specific information enters only through the strip location and
the zero count with multiplicity.
\begin{corollary}[A fixed small boundary coefficient]\label{gzb:cor:boundary}
For every fixed $d>0$,
\begin{equation}\label{gzb:eq:boundary}
 \liminf_{T\to\infty}\ \inf_{0\ne u\in
 V_{T,d(\log T)\sqrt{\log\log T}}}
 \frac{q_0[u]}{\log T\,\norm u_2^2}
 \ge 1-\frac{120\sqrt d}{2^{3/4}\sqrt\pi}.
\end{equation}
Thus every fixed $0<d<\pi\sqrt2/7200$ gives eventual strict positivity
on the whole band. This constant is conservative; no assertion is made
at equality.
\end{corollary}

\subsubsection{Relation to prior methods and the scope of the result}
Zhang's real frequency sparse Fourier extrapolation theorem
\cite[Theorem 1.1 and Proposition 3.7]{Zhang} supplies the starting
methodological comparison. His weighted frequency elimination and energy
concentration argument is the model for Lemma~\ref{gzb:lem:energy}.
The stated frequency class in that paper has real Fourier frequencies,
equivalently purely imaginary exponential roots, and does not itself
state the strip and polynomial prefix conclusion needed here.
We give this extension directly, and combine it with an orthonormal
Volterra construction. At zero strip width our explicit bound recovers
the constant in Zhang's Proposition 3.7 after translation and rescaling;
see Remark~\ref{gzb:rem:zhang}.
Erd\'elyi's work \cite{Erdelyi} provides an earlier general setting for
coefficient independent inequalities for exponential sums. Our derivative
estimate is proved below from the explicit basis, rather than being
inferred from a compactness argument.

The theta form normalization and the carrier comparison were used in
preceding stages of the author's research. This paper reproduces their
definitions and the estimates it needs; no earlier manuscript of the
author is an external hypothesis. The new localization scale improves
the $O(\sqrt N)$ prime loss available from a linear physical cutoff to
$O(\sqrt N/(\log N)^{1/4})$. This is the gain that allows
\eqref{gzb:eq:width}. We do not assert optimality of that scale or its constants.

No local positivity result can simply be summed across arbitrary bands:
the mixed terms of $q_0$ remain. Section~\ref{gzb:sec:schur} retains the full
inverse response within a union already covered by the whole band
estimate. It does not provide an exhaustive gluing theorem.
In particular, \eqref{gzb:eq:main} is an ordinary $L^2$ comparison on growing
finite dimensional spaces, not a uniform gap in $L^2(dx/W)$.
Neither global nonnegativity nor the Riemann hypothesis follows.

\subsection{A finite interval energy comparison}
\begin{lemma}[Weighted half line concentration]\label{gzb:lem:energy}
Suppose $P$ has prefix budget at most $N\ge1$ and $|\Rea z_j|\le a$,
where $a\ge0$. For $\ell>0$ set
\[
 \alpha=8N/\ell+2a,\qquad w(t)=e^{\alpha(t-\ell)}.
\]
Then
\begin{equation}\label{gzb:eq:energy}
 E(P):=\int_{-\infty}^{\ell}w(t)|P(t)|^2\,dt
       \le2\norm P_{L^2[0,\ell]}^2.
\end{equation}
The reflected version is
\begin{equation}\label{gzb:eq:energy-forward}
 \int_0^\infty e^{-\alpha t}|P(t)|^2\,dt
       \le2\norm P_{L^2[0,\ell]}^2.
\end{equation}
\end{lemma}
\begin{proof}
At an exponent $z$ of a current exponential polynomial $f$, put
$c=\alpha+2\Rea z>0$ and define
\[
 b(t)=c\int_0^\infty e^{-(c-z)s}f(t-s)\,ds,\qquad h=f-b.
\]
The denominator associated with another root $z_j$ has real part
$\alpha+\Rea z+\Rea z_j\ge\alpha-2a>0$.
All integrals therefore converge, including polynomial factors.
For the selected root, convolution sends $e^{zt}p(t)$ to
$e^{zt}(p(t)-p'(t)/c+\cdots)$, so $h$ decreases that root's polynomial
prefix by at least one. Other prefixes do not increase.
Successive elimination gives $f_0=P$, $f_m=0$ with $m\le N$ and
$b_j=f_j-f_{j+1}$, $0\le j<m$.

Differentiation gives $b'=zb+ch$ and the exact identity
\begin{equation}\label{gzb:eq:energy-identity}
 (w|b|^2)'=c\,w(|f|^2-|h|^2).
\end{equation}
The boundary at $-\infty$ vanishes because $\alpha>2a$.
Integration yields
\[
 E(f_j)-E(f_{j+1})=|b_j(\ell)|^2/c_j\ge0.
\]
Multiplication by $\ell-t$ and integration by parts yield
\[
 \int_{-\infty}^{\ell}(\ell-t)w(t)|P(t)|^2\,dt
   =\sum_{j=0}^{m-1}\frac{E(b_j)}{c_j}
   \le\frac{4m}{\alpha-2a}E(P)\le\frac\ell2E(P).
\]
Here $E(b_j)\le2E(f_j)+2E(f_{j+1})\le4E(P)$.
For $t<0$, $\ell-t\ge\ell$, so the weighted energy there is at most
$E(P)/2$. On $[0,\ell]$, $w\le1$. This proves \eqref{gzb:eq:energy}.
Apply it to $P(\ell-t)$ and change variables to obtain
\eqref{gzb:eq:energy-forward}.
The argument includes repeated roots directly, without taking a limit
of a singular coefficient Gram matrix.
\end{proof}

\subsection{The stable Volterra basis and short extrapolation}
The orthonormal construction below is the classical right half plane
Takenaka--Malmquist system, written as a product of stable Volterra
all pass operators. Zhang \cite[Appendix A, especially Remark A.7]{Zhang}
discusses the all pass formulation. We include its proof, endpoint
identities and continuation estimate to cover varying positive real
parts and every repeated root prefix explicitly. The construction
itself is not claimed as new.

\begin{lemma}[An orthonormal stable basis]\label{gzb:lem:basis}
Let $\zeta_1,\ldots,\zeta_m$ have positive real parts $b_j$, with
repetitions allowed. On $L^2(0,\infty)$ define
\begin{equation}\label{gzb:eq:volterra}
 (T_\zeta h)(t)=h(t)-2\Rea\zeta
                    \int_0^t e^{-\zeta(t-v)}h(v)\,dv.
\end{equation}
The functions
\begin{equation}\label{gzb:eq:basis}
 e_j(t)=\sqrt{2b_j}\,e^{-\zeta_jt},\qquad
 \phi_j=T_{\zeta_1}\cdots T_{\zeta_{j-1}}e_j
\end{equation}
are an orthonormal basis of the exponential polynomial space with root
list $-\zeta_1,\ldots,-\zeta_m$, including its repeated root prefixes.
If $0<b_j\le B$, then for $d\ge0$,
\begin{equation}\label{gzb:eq:basis-growth}
 |\phi_j(-d)|\le\sqrt{2B}
   \exp\bigl(Bd+2\sqrt{2B(j-1)d}\bigr).
\end{equation}
Also
\begin{equation}\label{gzb:eq:basis-endpoints}
 \phi_j(0)=\sqrt{2b_j},\qquad
 \phi_j'(0)=-\Bigl(\zeta_j+2\sum_{k<j}b_k\Bigr)\sqrt{2b_j}.
\end{equation}
\end{lemma}
\begin{proof}
Extend $h$ by zero to the negative half line. The full line Fourier
multiplier of \eqref{gzb:eq:volterra} is
$(i\xi-\overline\zeta)/(i\xi+\zeta)$, of modulus one for real $\xi$.
Causality makes it an isometry on the half line.
Direct Fubini gives
$\ip{T_\zeta h}{e^{-\zeta t}}_{L^2(0,\infty)}=0$:
the inner convolution integral has denominator
$\zeta+\overline\zeta=2\Rea\zeta$.
These assertions first hold for compactly supported functions and extend
by boundedness.
For $i<j$, cancel the first $i-1$ common isometries in the inner product
of $\phi_i$ and $\phi_j$, and use this orthogonality for $T_{\zeta_i}$.
Each $\phi_j$ belongs to the prefix space of the first $j$ roots.
That space has dimension $j$, including repetitions. Thus they are a basis.

Their unique exponential polynomial continuations to negative real $t$
satisfy, on reversing the finite integral in \eqref{gzb:eq:volterra},
\[
 |(T_\zeta h)(-d)|\le |h(-d)|+
       2B\int_0^d e^{B(d-v)}|h(-v)|\,dv.
\]
After multiplying by $e^{-Bd}$, iterate $I+2BV$, where
$(Vg)(d)=\int_0^d g(v)\,dv$. The resulting majorant is
\[
 \sqrt{2B}\,e^{Bd}\sum_{k=0}^{j-1}
          \binom{j-1}{k}\frac{(2Bd)^k}{k!}.
\]
Use $\binom{j-1}{k}\le(j-1)^k/k!$ and
$\sum_{k\ge0}y^k/(k!)^2\le e^{2\sqrt y}$ for $y\ge0$.
The latter follows from the even terms of the exponential series, since
$\binom{2k}{k}\le4^k$. This proves \eqref{gzb:eq:basis-growth}.
At zero, $T_\zeta h(0)=h(0)$ and
$(T_\zeta h)'(0)=h'(0)-2\Rea\zeta\,h(0)$.
Iteration proves \eqref{gzb:eq:basis-endpoints}.
\end{proof}

\begin{theorem}[Complex strip short extrapolation]\label{gzb:thm:extrapolate}
Suppose \eqref{gzb:eq:prefix} holds, $|\Rea z_j|\le a$, $\ell>0$, and
$N\ge1$. Put
\begin{equation}\label{gzb:eq:C}
 \alpha=8N/\ell+2a,\qquad C=\alpha+2a=8N/\ell+4a.
\end{equation}
For every $d\ge0$,
\begin{equation}\label{gzb:eq:extrapolate}
 |P(\ell+d)|\le\sqrt{2NC}
     \exp\bigl(ad+2\sqrt{NCd}\bigr)\norm P_{L^2[0,\ell]}.
\end{equation}
Constants are independent of imaginary parts, coefficient conditioning,
collisions and multiplicity distribution.
\end{theorem}
\begin{proof}
Set $H(u)=e^{-\alpha u/2}P(\ell-u)$, $u\ge0$.
Its stable roots are $-\zeta_j$, with
$\zeta_j=\alpha/2+z_j$ and $0<\Rea\zeta_j\le C/2$.
Lemma~\ref{gzb:lem:energy} gives
$\norm H_{L^2(0,\infty)}\le\sqrt2\norm P_{L^2[0,\ell]}$.
Expanding in Lemma~\ref{gzb:lem:basis} and applying Cauchy--Schwarz gives
\[
 |H(-d)|\le\sqrt{NC}\,e^{Cd/2+2\sqrt{NCd}}
                                      \norm H_{L^2(0,\infty)}.
\]
Since $P(\ell+d)=e^{-\alpha d/2}H(-d)$ and $C-\alpha=2a$,
\eqref{gzb:eq:extrapolate} follows. A shorter list only decreases the
displayed upper bounds.
\end{proof}
\begin{remark}\label{gzb:rem:zhang}
At $a=0$, $\ell=2$, the prefactor in \eqref{gzb:eq:extrapolate} is
$2\sqrt2\,N$ and the exponential factor is $e^{4N\sqrt d}$.
Translation to $[-1,1]$ recovers Zhang's Proposition 3.7 \cite{Zhang}.
The strip and full prefix cases are separately proved extensions.
No sharpness at large extrapolation distances is asserted.
\end{remark}

\subsection{The full theta tail}
\begin{lemma}[Theta envelopes]\label{gzb:lem:theta}
For $r=|x|$,
\begin{equation}\label{gzb:eq:theta-envelope}
 c_W e^{4r-\pi e^{2r}}\le W(x)\le C_4 e^{4r-\pi e^{2r}},
 \qquad c_W=2\pi^2-3\pi>9,\quad C_4<42.
\end{equation}
The weight is locally absolutely continuous and, almost everywhere,
\begin{equation}\label{gzb:eq:logW}
 |(\log W)'(x)|\le112e^{2|x|}.
\end{equation}
\end{lemma}
\begin{proof}
Every term in \eqref{gzb:eq:weight} is positive. The $j=1$ term and
$1\le1+e^{-r}\le2$ give the lower bound.
For the upper bound take
$C_4=4\pi^2\sum_{j\ge1}j^4e^{-\pi(j^2-1)}$.
The first term beyond $j=1$ is less than $1/32$, using $\pi>3$, $e>2$.
Consecutive tail terms have ratio at most $(3/2)^4e^{-15}<1/2$.
Thus the full sum is less than $17/16$, and $\pi<22/7$ gives
$C_4<2057/49<42$. Also $c_W>9$ follows from $\pi>3$.

For $r>0$, differentiation is locally uniform. Write
$g(r)=(1+e^{-r})^{-1}$, so $0<g\le1$, $0\le g'\le1$.
The absolute derivative of a summand times $g$ is at most
\[
 (8\tau_j^3+32\tau_j^2+18\tau_j)e^{-\tau_j}
       \le21\tau_j^3e^{-\tau_j},\qquad \tau_j\ge\pi>3.
\]
Hence $|W'|\le C_6e^{6r-\pi e^{2r}}$, where
$C_6=21\pi^3\sum_{j\ge1}j^6e^{-\pi(j^2-1)}<1000$.
Indeed the first tail term is less than $1/8$, consecutive tail ratios
are at most $(3/2)^6e^{-15}<1/2$, and
$21(22/7)^3(5/4)<1000$.
The same bound holds almost everywhere for the continuous even
continuation, giving local absolute continuity at zero.
Divide by the lower bound and use $1000/9<112$.
\end{proof}

\begin{theorem}[An explicit sublinear mass cutoff]\label{gzb:thm:tail}
Suppose $P$ has prefix budget at most $N\ge4096$ and
$|\Rea z_j|\le1/2$. Set
\begin{equation}\label{gzb:eq:cutoff}
 s=\log N,\quad M_N=\frac{100N}{\sqrt s},\quad
 R_N=\tfrac12\log M_N,\quad U=\norm{WP}_2.
\end{equation}
For every $r\ge R_N$, writing $v=e^{2r}$,
\begin{align}
 |W(r)P(\pm r)|^2&\le e^{-4v}U^2,\label{gzb:eq:point-tail}\\
 \norm{\one_{\{|x|>r\}}WP}_2^2
     &\le\frac{e^{-4v}}{4v}U^2\le e^{-4v}U^2.\label{gzb:eq:mass-tail}
\end{align}
Moreover,
\begin{equation}\label{gzb:eq:moment}
 \int_\R e^{4|x|}|W(x)P(x)|^2\,dx\le(M_N^2+1)U^2.
\end{equation}
\end{theorem}
\begin{proof}
The theta envelope gives square integrability before normalization.
Assume $U>0$. For $r>1$ apply Theorem~\ref{gzb:thm:extrapolate} with
$a=1/2$, $\ell=r-1$ and $d=1$, and reflect for the negative ray.
The lower comparison function is decreasing on $x\ge0$, so
\[
 \norm P_{L^2[0,r-1]}
       \le c_W^{-1}e^{-4(r-1)+\pi e^{2(r-1)}}U.
\]
Multiplication by the upper envelope gives
\begin{equation}\label{gzb:eq:point-majorant}
 \frac{|W(r)P(\pm r)|^2}{U^2}
 \le\left(\frac{C_4}{c_W}\right)^2e^8(2NC_r)
  \exp\bigl(1+4\sqrt{NC_r}-2\pi(1-e^{-2})e^{2r}\bigr),
 \quad C_r=\frac{8N}{r-1}+2.
\end{equation}
Here $s>8$, $R_N\ge s/4$ and $R_N-1\ge s/8$.
On the entire ray $r\ge R_N$, $C_r\le64N/s+2\le66N/s$.
The logarithm of the prefactor in \eqref{gzb:eq:point-majorant}, including
its factor $e^1$, is at most
\[
 \log25+9+\log132+2s-\log s<18+2s\le3N/\sqrt s.
\]
For the last whole ray inequality,
$(18+2s)\sqrt s/e^s$ is decreasing for $s\ge8$ and is less than $3$ at $8$.
Also $4\sqrt{NC_r}<33N/\sqrt s$ and $2\pi(1-e^{-2})>5$.
Since $e^{2r}\ge100N/\sqrt s$,
\[
 \log\frac{|W(r)P(\pm r)|^2}{U^2}
       \le36N/\sqrt s-5e^{2r}\le-4e^{2r}.
\]
This proves \eqref{gzb:eq:point-tail} everywhere on the ray.
Integrating both rays gives
\[
 2\int_r^\infty e^{-4e^{2x}}\,dx
       =\int_v^\infty\frac{e^{-4t}}t\,dt\le\frac{e^{-4v}}{4v},
\]
proving \eqref{gzb:eq:mass-tail}.
Finally the layer cake identity gives
\[
 \int e^{4|x|}|WP|^2
 \le M_N^2U^2+4\int_{R_N}^\infty
              e^{4r}\norm{\one_{\{|x|>r\}}WP}_2^2\,dr
 \le\left(M_N^2+\frac{e^{-4M_N}}8\right)U^2.
\]
This implies \eqref{gzb:eq:moment}.
\end{proof}

\subsection{Every prime power and the complete prime estimate}
\begin{lemma}[Elementary prime head]\label{gzb:lem:head}
For every $X\ge1$,
\begin{equation}\label{gzb:eq:head}
 \sum_{2\le n\le X}\frac{\Lambda(n)}{\sqrt n}\le6\sqrt X,
 \qquad |p_{n\le X}[u]|\le12\sqrt X\,\norm u_2^2.
\end{equation}
Also $0\le a_n\le1$.
\end{lemma}
\begin{proof}
Write $\Psi(x)=\sum_{n\le x}\Lambda(n)$.
The factorial valuation formula gives
$\Psi(2m)-\Psi(m)\le\log\binom{2m}{m}\le2m\log2$.
Every prime power in $(m,2m]$ contributes one to the corresponding
valuation, whose other summands are nonnegative.
At powers of two, summation gives
$\Psi(2^k)\le(2^{k+1}-2)\log2$.
For $2^{k-1}<x\le2^k$ it follows that
$\Psi(x)<4(\log2)x<3x$, with the same weak bound at endpoints.
Partial summation gives
\[
 \sum_{n\le X}\frac{\Lambda(n)}{\sqrt n}
    =\frac{\Psi(X)}{\sqrt X}
       +\frac12\int_1^X\Psi(t)t^{-3/2}\,dt\le6\sqrt X.
\]
Each translation correlation is at most $\norm u_2^2$ in modulus,
which proves the two orientation head bound.
Finally $a_n\le(\log n)/\sqrt n\le2/e<1$.
\end{proof}

\begin{theorem}[Complete strip prime cap]\label{gzb:thm:prime}
Under the hypotheses of Theorem~\ref{gzb:thm:tail}, the full prime series
is absolutely convergent and
\begin{equation}\label{gzb:eq:primecap}
 |p[WP]|\le P_N\norm{WP}_2^2,\qquad
 P_N=12\sqrt{\left\lceil\frac{100N}{\sqrt{\log N}}\right\rceil}+1.
\end{equation}
Thus $P_N=O(\sqrt N/(\log N)^{1/4})$.
\end{theorem}
\begin{proof}
For $C_u(s)=\int u(x+s)\overline{u(x)}\,dx$, splitting into
$|x|>s/2$ and its complement yields
\begin{equation}\label{gzb:eq:correlation}
 |C_u(s)|\le2\norm u_2\norm{\one_{\{|x|>s/2\}}u}_2,\qquad s>0.
\end{equation}
On the complement $|x+s|\ge s/2$, justifying the second term after
a change of variables.
Set $Y=\lceil M_N\rceil$ and use Lemma~\ref{gzb:lem:head} for the head.
For every integer $n>Y$, \eqref{gzb:eq:mass-tail} at $r=(\log n)/2$ gives
$|C_u(\log n)|\le2e^{-2n}\norm u_2^2$.
The remaining series, with both orientations, costs at most
$4\sum_{n>Y}e^{-2n}\norm u_2^2<\norm u_2^2$.
This proves absolute convergence and \eqref{gzb:eq:primecap}.
No cancellation of carrier phases is used.
\end{proof}

\subsection{A coefficient independent derivative estimate}
The same basis gives an explicit polynomial derivative bound, useful
because both offset width and prefix count grow.
\begin{lemma}[A local derivative bound]\label{gzb:lem:local-derivative}
Suppose $Q$ has prefix budget at most $N\ge1$, with
$|\Rea z_j|\le1/2$ and $|\Ima z_j|\le N$.
For every $x\in\R$ and $0<\delta\le1$,
\begin{equation}\label{gzb:eq:local-derivative}
 |Q'(x)|\le100N^3\delta^{-3/2}
                           \norm Q_{L^2[x,x+\delta]}.
\end{equation}
\end{lemma}
\begin{proof}
Translate $x$ to zero. Put $\alpha=8N/\delta+1$,
$H(t)=e^{-\alpha t/2}Q(t)$.
By \eqref{gzb:eq:energy-forward},
$\norm H_{L^2(0,\infty)}\le\sqrt2\norm Q_{L^2[0,\delta]}$.
The stable roots $-\zeta_j$, $\zeta_j=\alpha/2-z_j$, satisfy
\[
 b_j=\Rea\zeta_j\le5N/\delta,\quad
 S:=\sum_jb_j\le5N^2/\delta,\quad
 \max_j|\zeta_j|\le6N/\delta,\quad \alpha/2\le\tfrac92N/\delta.
\]
Expanding in Lemma~\ref{gzb:lem:basis} and using \eqref{gzb:eq:basis-endpoints} gives
\[
 |H(0)|\le\sqrt{2S}\norm H_2,\qquad
 |H'(0)|\le\sqrt{2S}(\max_j|\zeta_j|+2S)\norm H_2.
\]
Since $Q'(0)=H'(0)+(\alpha/2)H(0)$,
\[
 |Q'(0)|\le\frac{41\sqrt{20}}2N^3\delta^{-3/2}
                   \norm Q_{L^2[0,\delta]}
          <100N^3\delta^{-3/2}\norm Q_{L^2[0,\delta]}.
\]
Repeated roots are included directly.
\end{proof}

\begin{proposition}[A global weighted derivative bound]\label{gzb:prop:derivative}
Under Lemma~\ref{gzb:lem:local-derivative}, for $N\ge4096$,
\begin{equation}\label{gzb:eq:derivative}
 \norm{(WQ)'}_2\le b_N\norm{WQ}_2,\qquad
 b_N=KN^4,\qquad K=101(100e^{118}+112).
\end{equation}
The constant is deliberately loose and independent of root separation.
\end{proposition}
\begin{proof}
Write $A=WQ$ and choose $\delta(x)=e^{-2(|x|+1)}$.
For $y\in[x,x+\delta(x)]$, \eqref{gzb:eq:logW} gives
$|\log W(y)-\log W(x)|\le112\delta(x)e^{2(|x|+1)}=112$.
Thus \eqref{gzb:eq:local-derivative} gives
\[
 |W(x)Q'(x)|^2
   \le100^2N^6e^{224}\delta(x)^{-3}
                         \int_x^{x+\delta(x)}|A(y)|^2\,dy.
\]
When $x\le y\le x+\delta(x)$, one has $|x-y|\le1$ and
$e^{-2}\delta(y)\le\delta(x)\le e^2\delta(y)$.
For each $y$ the set of such $x$ has length at most $e^2\delta(y)$.
Fubini gives
\[
 \norm{WQ'}_2^2
   \le100^2N^6e^{224+8}\int\delta(y)^{-2}|A(y)|^2\,dy
    =100^2N^6e^{236}\int e^{4|y|}|A(y)|^2\,dy.
\]
Also $\norm{W'Q}_2\le112\norm{e^{2|x|}A}_2$.
The moment bound \eqref{gzb:eq:moment} now gives
\[
 \norm{A'}_2
 \le(100e^{118}N^3+112)\sqrt{M_N^2+1}\norm A_2.
\]
Since $M_N\le100N$ and $\sqrt{M_N^2+1}\le101N$ for $N\ge4096$,
\eqref{gzb:eq:derivative} follows. Integrability before normalization
follows from the theta envelope, so no unknown derivative norm was
used to estimate itself.
\end{proof}

\subsection{Carrier separation and the complete Gamma comparison}
For the unitary Fourier transform
$\widehat u(\xi)=(2\pi)^{-1/2}\int u(x)e^{-i\xi x}\,dx$,
Plancherel and the nonnegative difference integral give
\begin{equation}\label{gzb:eq:gamma-symbol}
 \tg[u]=\int_\R d_\Gamma(\xi)|\widehat u(\xi)|^2\,d\xi,\qquad
 d_\Gamma(\xi)=2\int_0^\infty\rho(s)(1-\cos(\xi s))\,ds.
\end{equation}
Since $(2s)^{-1}\le(1-e^{-2s})^{-1}\le1+(2s)^{-1}$,
\begin{equation}\label{gzb:eq:gamma-bounds}
 \tfrac12\log(1+4\xi^2)\le d_\Gamma(\xi)
                           \le4+\tfrac12\log(1+4\xi^2).
\end{equation}
Here
$\int_0^\infty e^{-s/2}(1-\cos(\xi s))\,ds/s
=\tfrac12\log(1+4\xi^2)$, by differentiation in $\xi$ and integration
back from zero. The remaining upper integral is
$2\int_0^\infty e^{-s/2}(1-\cos(\xi s))\,ds\le4$.
No Gamma displacement or frequency is truncated.

Let $Q$ have budget at most $N\ge4096$,
$|\Rea z_j|\le1/2$ and $0\le\Ima z_j\le L\le N$. Set
\begin{equation}\label{gzb:eq:packet}
 A=WQ,\qquad
 u_T(x)=\frac{e^{iTx}A(x)+e^{-iTx}A(-x)}2.
\end{equation}
Only the positive envelope budget is $N$; the paired representation
may have budget $2N$. Reflection is $A(-x)$, not $\overline{A(x)}$.

\begin{proposition}[Complete whole packet estimate]\label{gzb:prop:packet}
Suppose $T\ge\max(1,2b_N)$. Set
\begin{equation}\label{gzb:eq:packet-constants}
 h=\log T,\quad h_T=\tfrac12\log(1+T^2),\quad
 \eta_{T,N}=16b_N^2/T^2,\quad C_\Gamma=4+\tfrac12\log37.
\end{equation}
Then
\begin{align}
 |p[u_T]|&\le P_N\norm{u_T}_2^2,\label{gzb:eq:packet-prime}\\
 h_T(1-\eta_{T,N})\norm{u_T}_2^2
   &\le\tg[u_T]\le(h+C_\Gamma)\norm{u_T}_2^2,\label{gzb:eq:packet-gamma}\\
 |q_0[u_T]-h\norm{u_T}_2^2|
   &\le\left(|\cg|+P_N+
      \max(C_\Gamma,\eta_{T,N}h_T)\right)\norm{u_T}_2^2.
       \label{gzb:eq:packet-error}
\end{align}
The prime cap uses the positive envelope budget $N$, not $2N$.
\end{proposition}
\begin{proof}
The reflected carrier correlation is
$I_T=\int e^{2iTx}A(x)\overline{A(-x)}\,dx$.
The nonoscillatory factor has derivative $L^1$ norm at most
$2\norm{A'}_2\norm A_2$ and vanishing endpoint terms.
Integration by parts gives
\begin{equation}\label{gzb:eq:carrier-norm}
 |I_T|\le(b_N/T)\norm A_2^2,\quad
 \norm{u_T}_2^2\ge\tfrac12(1-b_N/T)\norm A_2^2,\quad
 \norm A_2\le2\norm{u_T}_2.
\end{equation}
Triangle inequality and reflection transfer \eqref{gzb:eq:mass-tail} to
\[
 \norm{\one_{\{|x|>r\}}u_T}_2
       \le2e^{-2e^{2r}}\norm{u_T}_2\quad(r\ge R_N).
\]
The prime head is charged directly to $\norm{u_T}_2^2$.
For $n>\lceil M_N\rceil$, \eqref{gzb:eq:correlation} now gives
$|C_{u_T}(\log n)|\le4e^{-2n}\norm{u_T}_2^2$.
The remaining series costs at most
$8\sum_{n>\lceil M_N\rceil}e^{-2n}\norm{u_T}_2^2
 <\norm{u_T}_2^2$, proving \eqref{gzb:eq:packet-prime}.
The factor two in \eqref{gzb:eq:carrier-norm} affects only the
exponentially small prime tail, not the head.

Plancherel and \eqref{gzb:eq:derivative} imply
\[
 \norm{\one_{[-T/2,T/2]}\widehat{u_T}}_2
      \le(2b_N/T)\norm A_2\le(4b_N/T)\norm{u_T}_2.
\]
On the complementary Fourier region, \eqref{gzb:eq:gamma-bounds} is at least
$h_T$, proving the lower Gamma bound.
Differentiation of \eqref{gzb:eq:packet} gives
$\norm{u_T'}_2\le(T+b_N)\norm A_2
                  \le2(T+b_N)\norm{u_T}_2$.
For $u_T\ne0$, Jensen's inequality gives
\[
 \frac{\tg[u_T]}{\norm{u_T}_2^2}
 \le4+\tfrac12\log\left(1+4\frac{\norm{u_T'}_2^2}{\norm{u_T}_2^2}\right)
 \le4+\tfrac12\log(1+16(T+b_N)^2)\le h+C_\Gamma.
\]
The last step uses $T\ge1$ and $b_N/T\le1/2$.
Since $h_T\ge h$, the lower Gamma error from $h$ is at most
$\eta_{T,N}h_T$. Add the absolute scalar and prime costs to obtain
\eqref{gzb:eq:packet-error}. The zero function needs no division.
\end{proof}

\subsection{The complete actual band and its zero count}
The asymptotic conclusions need only the classical Riemann--von Mangoldt
formula $N_\zeta(t)=F(t)+O(\log t)$, with multiplicities
\cite{vonMangoldt,HSW}. To declare a concrete scalar upper budget without
enumerating the zeros, we invoke the published explicit bound
\cite[Corollary 1.2]{HSW}: for $t\ge e$,
\begin{equation}\label{gzb:eq:count}
 |N_\zeta(t)-F(t)|\le R(t),\quad
 F(t)=\frac{t}{2\pi}\log\frac{t}{2\pi e},\quad
 R(t)=0.1038\log t+0.2573\log\log t+9.3675.
\end{equation}
Here $N_\zeta$ counts positive ordinates with multiplicity. The proof of
this explicit refinement in \cite{HSW} uses finite height numerical zero
verification. We perform no such enumeration here; the asymptotic band
conclusions do not depend on the particular numerical constants in this
refinement. Endpoint conventions will not matter because we widen both endpoints.
For $T\ge e+1$ put
\begin{equation}\label{gzb:eq:countbudget}
 \mathcal M(T,L)=F(T+L+1)-F(T-1)+R(T+L+1)+R(T-1).
\end{equation}
Thus the full multiplicity count $M(T,L)$ in the closed band satisfies
$M(T,L)\le\mathcal M(T,L)$.
Uniformly for $0\le L=o(T)$, the mean value theorem gives
\begin{equation}\label{gzb:eq:count-asym}
 M(T,L)\le\frac{L+2}{2\pi}\log T+
       O\left(\log T+\frac{(L+2)^2}{T}\right)
       =O((L+1)\log T).
\end{equation}
This is an upper bound, not a relative short interval counting asymptotic.

Every $u\in V_{T,L}$ has exactly the representation \eqref{gzb:eq:packet}, with
\begin{equation}\label{gzb:eq:actual-envelope}
 Q(x)=\sum_{\substack{w:\ T\le\Rea w\le T+L}}
       \ \sum_{k<\nu(w)}c_{w,k}(ix)^k e^{i(w-T)x}.
\end{equation}
Its positive envelope budget is at most $M(T,L)$ and its roots have
imaginary parts in $[0,L]$ and real parts in $[-1/2,1/2]$.
Differentiate both exponentials in the cosine to verify the identity;
it is valid for complex $w$ and complex $c_{w,k}$.

\begin{proof}[Proof of Theorem~\ref{gzb:thm:main}]
Write $h=\log T$. Choose the common padded integer budget
\begin{equation}\label{gzb:eq:common-budget}
 n(T)=\left\lceil\max\{4096,h,L(T),\mathcal M(T,L(T))\}\right\rceil.
\end{equation}
No zero enumeration is required.
By \eqref{gzb:eq:count-asym} and \eqref{gzb:eq:width},
\[
 n(T)=O((L(T)+1)h),\qquad n(T)\ge h,\qquad
 n(T)=O(h^2\sqrt{\log h}).
\]
Consequently $b_{n(T)}/T\to0$ and $h_Tb_{n(T)}^2/T^2\to0$,
so Proposition~\ref{gzb:prop:packet} applies eventually.
The complete prime cap gives
\begin{equation}\label{gzb:eq:relative-prime}
 \frac{P_{n(T)}}h
 =O\left(\frac{\sqrt{n(T)}}{h(\log n(T))^{1/4}}\right)
 =O\left(\sqrt{\frac{L(T)+1}{h\sqrt{\log h}}}\right)=o(1).
\end{equation}
Now \eqref{gzb:eq:packet-error}, divided by $h$, tends to zero uniformly
over \eqref{gzb:eq:actual-envelope}, proving \eqref{gzb:eq:main}.
For $L=h(\log h)^{1/4}$ the prime contribution is
$O((\log h)^{-1/8})$.
\end{proof}

\begin{proof}[Proof of Corollary~\ref{gzb:cor:boundary}]
Fix $d>0$ and set $L=dh\sqrt{\log h}$.
The explicit upper budget \eqref{gzb:eq:common-budget} satisfies
\[
 n(T)=\left(\frac d{2\pi}+o(1)\right)h^2\sqrt{\log h},
 \qquad\log n(T)=2\log h+o(\log h).
\]
This is an asymptotic for the declared upper budget using $\mathcal M$,
not an unsupported relative asymptotic for the zeros actually present.
It follows that
$\limsup_{T\to\infty}P_{n(T)}/h
 \le120\sqrt d/(2^{3/4}\sqrt\pi)$.
The complete Gamma coefficient tends to one, and $\cg/h\to0$.
Subtracting the full absolute prime cap proves \eqref{gzb:eq:boundary},
whose lower bound is positive in the stated strict range of $d$.
\end{proof}

For any spanning list $v_1,\ldots,v_m$, including a redundant list,
let $G_{ij}=\ip{v_j}{v_i}_2$ and $Q_{ij}=q_0(v_j,v_i)$.
The uniform result is equivalently
\[
 (1-\varepsilon_T)hG\preceq Q\preceq(1+\varepsilon_T)hG.
\]
No raw coefficient Gram eigenvalue, frequency gap or condition number
is divided out. A represented zero vector lies in the kernel of both sides.

\subsection{An exact inverse retaining Schur consequence}\label{gzb:sec:schur}
The whole band estimate controls screening on every subspace of that same
band. The following version preserves the full prime error and its
inverse response debit.
On even finite theta exponential polynomials define
\begin{equation}\label{gzb:eq:reference}
\begin{split}
 \mathcal K(u,v)&=\int_\R e^{-|s|/2}
                     \int_\R u(x+s)\overline{v(x)}\,dx\,ds,\\
 B&=\tg-\cg\ip{\cdot}{\cdot}_2+\mathcal K,\qquad
 E=q_0-B=-p-\mathcal K.
\end{split}
\end{equation}
The Fourier multiplier of $\mathcal K$ is $(\xi^2+1/4)^{-1}$, so
$0\preceq\mathcal K\preceq4G$.
On the band in Theorem~\ref{gzb:thm:main}, eventually
\begin{equation}\label{gzb:eq:reference-bound}
 (h-C_B)G\preceq B\preceq(h+C_B)G,\qquad
 C_B=C_\Gamma+|\cg|+5,
\end{equation}
by \eqref{gzb:eq:packet-gamma} with $\eta_{T,n}h_T\le1$.
This gives a positive reference on the entire represented band without
a separately assumed sign of $q_0$.
The full arithmetic cap gives
\begin{equation}\label{gzb:eq:E-bound}
 -(P_n+4)G\preceq E\preceq P_nG.
\end{equation}

\begin{corollary}[Signed screened determinant]\label{gzb:cor:schur}
Let a finite represented union $R+V$ lie in one band satisfying
\eqref{gzb:eq:reference-bound}, with $h>C_B$. Set
$r_T=(P_n+4)/(h-C_B)<1$.
Choose residual sources spanning $R$ after $B$ orthogonal projection
off $V$, and let $H$ be their $B$ Gram.
Let $Q_{VV}$ be the actual complete $q_0$ matrix on represented support
of $V$. It is positive definite there. Exact screening gives
\begin{equation}\label{gzb:eq:Schur}
 S=H+E_{RR}-E_{RV}Q_{VV}^{-1}E_{VR},\qquad
 (1-r_T)H\preceq S\preceq(1+r_T)H.
\end{equation}
For two residual sources, put
$F=E_{RR}-E_{RV}Q_{VV}^{-1}E_{VR}$. Then
\begin{equation}\label{gzb:eq:det}
 \operatorname{tr}(\adj(H)F)+\det F
                   \ge(-2r_T+r_T^2)\det H.
\end{equation}
Under \eqref{gzb:eq:width}, $r_T\to0$.
All inverses act on represented support, not redundant coordinates.
\end{corollary}
\begin{proof}
From \eqref{gzb:eq:reference-bound} and \eqref{gzb:eq:E-bound},
$(1-r_T)B\preceq q_0\preceq(1+r_T)B$ on the complete union.
Thus $Q_{VV}$ is positive on represented support.
For residual $x\in R$, minimization over $v\in V$ gives
\[
 (1-r_T)B[x]\le\inf_{v\in V}q_0[x+v]\le(1+r_T)B[x],
\]
since $B(x,V)=0$.
Completing the square in the actual $V$ block gives \eqref{gzb:eq:Schur},
including the nonnegative inverse response debit.
When $H$ is positive definite, normalization by $H^{-1/2}$ implies
$\det S\ge(1-r_T)^2\det H$.
Use the exact two by two identity
$\det(H+F)=\det H+\operatorname{tr}(\adj(H)F)+\det F$.
When $H$ is singular, \eqref{gzb:eq:Schur} gives the same kernel for $H$
and $S$, so both determinants vanish and the identity still proves
\eqref{gzb:eq:det}. Finally \eqref{gzb:eq:relative-prime} gives $r_T\to0$.
\end{proof}

The debit has not been replaced by an ordinary Gram inverse, bounded
entry by entry, or omitted. Control comes from the joint all coefficients
estimate on the complete union. This is not a gluing argument from two
neighboring positive bands. A union wider than the proved range is not
covered.

\subsection{Verification and remaining limitations}
The estimates proved here are analytic. We perform no numerical theta
integration, zero enumeration or matrix sign computation, and we do not
use the prime number theorem. The explicit zero counting theorem
\eqref{gzb:eq:count} is invoked as published, including its finite height
numerical verification dependency. The qualitative band conclusions need
only the classical asymptotic zero count, so no claim that the complete
dependency chain of the explicit refinement is nonnumerical is needed.
The constants $4096$, $100$, $12$ and the boundary coefficient are
conservative consequences of the displayed inequalities.
The very large $K$ in \eqref{gzb:eq:derivative} is explicit but is not
advertised as an efficient numerical height bound.
For specified $T,L$, \eqref{gzb:eq:common-budget}, the carrier condition
$T\ge2Kn(T)^4$ and \eqref{gzb:eq:packet-error} provide a scalar sufficient
test. No final numerical starting height is evaluated here.

The accompanying growing-band check script checks exact scalar inequalities, the basis
endpoint formulas in simple configurations, and the two by two determinant
algebra. Floating point sample checks are labelled as illustrations,
not substitutes for whole ray proofs.
The standalone source, PDF and reproducibility files are available at
\url{https://github.com/serabing-hash/theta-weighted-certificate/tree/main/publications/growing-zero-jet-bands-2026-11-03}.
This source includes all arguments except the classical zeta strip
location and the cited zero counting theorem. Compilation uses standard \LaTeX\ and requires no
private data or computer algebra system.

Three restrictions remain essential. The prime estimate controls bounded
prefix complexity, not arbitrary infinite exponential sums.
The actual zero theorem covers an entire single contiguous band, not
the mixed terms in an arbitrary covering of all positive ordinates.
Ordinary $L^2$ and reciprocal theta weighted norms cannot be interchanged
with a budget independent constant. The Schur corollary does not remove
these restrictions. The unresolved global sign, and any deduction of
the Riemann hypothesis through that sign, remains outside this result.


\section*{License}
Copyright \textcopyright\ 2026 Jongmin Choi. This work is licensed under the Creative Commons Attribution-NonCommercial 4.0 International License (CC BY-NC 4.0): \url{https://creativecommons.org/licenses/by-nc/4.0/}. External references retain their respective rights.

\begin{thebibliography}{99}\small
\bibitem{Guinand}
A.~P. Guinand, A summation formula in the theory of prime numbers,
\emph{Proceedings of the London Mathematical Society} (2) \textbf{50} (1948), 107--119.
\url{https://doi.org/10.1112/plms/s2-50.2.107}.
\bibitem{vonMangoldt}
H.~von Mangoldt, Zur Verteilung der Nullstellen der Riemannschen Funktion,
\emph{Mathematische Annalen} \textbf{60} (1905), 1--19.
\url{https://doi.org/10.1007/BF01447494}.
\bibitem{Kato}
T.~Kato, \emph{Perturbation Theory for Linear Operators}, second edition,
Classics in Mathematics, Springer, Berlin, 1995, Chapter~VI.
\url{https://doi.org/10.1007/978-3-642-66282-9}.
\bibitem{deBranges1961}
L.~de Branges, Some Hilbert spaces of entire functions.~II,
\emph{Transactions of the American Mathematical Society} \textbf{99} (1961), 118--152.
\url{https://doi.org/10.1090/S0002-9947-1961-0133456-2}.
\bibitem{DLMF}
National Institute of Standards and Technology, \emph{Digital Library of Mathematical Functions}.
Jacobi transformation, Eq.~20.7.32; Gamma integral, reflection, beta integral and uniform complex Stirling formulas, Sections~5.5, 5.9, 5.11 and 5.12; completed zeta function and Euler product, Sections~25.2 and 25.4; zeros of the Riemann zeta function, Section~25.10; all-prime-power Chebyshev function and qualitative prime number theorem, Eqs.~25.16.1 and 25.16.3.
\url{https://dlmf.nist.gov/20.7}; \url{https://dlmf.nist.gov/5.9};
\url{https://dlmf.nist.gov/5.11}; \url{https://dlmf.nist.gov/5.12};
\url{https://dlmf.nist.gov/25.2}; \url{https://dlmf.nist.gov/25.4}; \url{https://dlmf.nist.gov/25.10}; \url{https://dlmf.nist.gov/25.16.i}.
\bibitem{ChoiGrowingBands}
J.~Choi, \emph{Short extrapolation and growing zero jet bands for a theta
prime Gamma form}, 3 November 2026, repository publication at commit
\texttt{f2a5e6e390f213149f5f9d1d023aa2dd81d33518}.
\url{https://github.com/serabing-hash/theta-weighted-certificate/tree/f2a5e6e390f213149f5f9d1d023aa2dd81d33518/publications/growing-zero-jet-bands-2026-11-03}.
\bibitem{Zhang}
R.~Zhang, \emph{Optimal extrapolation bounds for sparse Fourier sums},
arXiv:2607.10501v2, 13 October 2026.
\url{https://arxiv.org/html/2607.10501v2}.
\bibitem{Erdelyi}
T. Erd\'elyi, Inequalities for exponential sums, \emph{Sbornik: Mathematics} \textbf{208} (2017), no.~3, 433--464. \url{https://doi.org/10.1070/SM8670}. Author manuscript, November 8, 2016.
\url{https://people.tamu.edu/~terdelyi/papers-online/sbornik150R.pdf}.
\bibitem{HSW}
E. Hasanalizade, Q. Shen, and P.-J. Wong, \emph{Counting zeros of the Riemann zeta function}, arXiv:2107.06506v1, July 14, 2021.
\url{https://arxiv.org/abs/2107.06506v1}.
\bibitem{Boyd}
S. Boyd and L. Vandenberghe, \emph{Convex Optimization}, Cambridge University Press, 2004, Appendix A.5.
\url{https://web.stanford.edu/~boyd/cvxbook/}.
\bibitem{Hilbert}
R. Ba\~nuelos and M. Kwa\'snicki, On the $\ell^p$-norm of the discrete Hilbert transform,
\emph{Duke Mathematical Journal} \textbf{168}, no.~3 (2019), 471--504.
\url{https://doi.org/10.1215/00127094-2018-0047}.
\url{https://arxiv.org/abs/1709.07427v2}.
% Append these four entries to the common bibliography. Reuse existing Erdelyi.

\bibitem{fzLagarias} J. C. Lagarias and D. Montague,
\emph{The integral of the Riemann xi-function}, arXiv:1106.4348v3, 2011,
equation (1.2). \url{https://arxiv.org/abs/1106.4348}.
\bibitem{fzFlintZeros} FLINT authors, \emph{Zeros of the Riemann zeta function},
\texttt{acb\_dirichlet} documentation, accessed November 3, 2026.
\url{https://flintlib.org/doc/acb_dirichlet/index.html}.
\bibitem{fzFlintIntegral} Python-FLINT authors, \emph{acb: complex numbers},
version 0.9.0 documentation, integral callback contract.
\url{https://python-flint.readthedocs.io/en/latest/acb.html}.
\bibitem{fzArbIntegration} F. Johansson, \emph{Numerical integration in
arbitrary-precision ball arithmetic}, arXiv:1802.07942, 2018.
\url{https://arxiv.org/abs/1802.07942}.

\bibitem{gbTeschl}
G. Teschl, \emph{Mathematical Methods in Quantum Mechanics},
Graduate Studies in Mathematics 99, American Mathematical Society, 2009,
Theorem 2.13, Section 3.1, and Theorem 6.4.
\url{https://www.mat.univie.ac.at/~gerald/ftp/book-schroe/schroe.pdf}.

\end{thebibliography}
\end{document}
